% 自动生成；修改模板或打印配置后运行 _台账/生成LaTeX.py。
% 独立文件：XeLaTeX 编译两次即可，不需要 input 或外部代码文件。
\documentclass[UTF8,fontset=fandol,a4paper,twoside]{ctexart}
\usepackage[inner=11mm,outer=9mm,top=11mm,bottom=12mm,
    headheight=11pt,headsep=3mm,footskip=6mm]{geometry}
\usepackage{amsmath,amssymb,multicol,listings,xcolor,enumitem,fancyhdr,titlesec,tabularx}
\usepackage[hidelinks,unicode]{hyperref}
\setmonofont{lmmono10-regular.otf}[BoldFont=lmmonolt10-bold.otf]
\setCJKmonofont{FandolFang-Regular}[BoldFont=FandolHei-Regular]
\definecolor{ink}{gray}{0.18}
\setlength{\columnsep}{6mm}
\setlength{\columnseprule}{0.2pt}
\setlength{\parindent}{0pt}
\setlength{\parskip}{1pt}
\setlength{\multicolsep}{3pt}
\setlength{\abovedisplayskip}{3pt}
\setlength{\belowdisplayskip}{3pt}
\setlength{\abovedisplayshortskip}{2pt}
\setlength{\belowdisplayshortskip}{2pt}
\setlist[itemize]{nosep,leftmargin=1.1em,topsep=1pt}
\titleformat{\section}{\fontsize{11}{12}\selectfont\bfseries}{\thesection}{0.4em}{}
\titleformat{\subsection}{\fontsize{9}{10}\selectfont\bfseries}{\thesubsection}{0.4em}{}
\titlespacing*{\section}{0pt}{7pt}{3pt}
\titlespacing*{\subsection}{0pt}{5pt}{2pt}
\setcounter{tocdepth}{2}
\makeatletter
\renewcommand{\@pnumwidth}{2em}
\renewcommand{\l@section}{\@dottedtocline{1}{0em}{1.6em}}
\renewcommand{\l@subsection}{\@dottedtocline{2}{0.8em}{2.7em}}
\makeatother
\pagestyle{fancy}\fancyhf{}
\fancyhead[LE,RO]{\fontsize{7.5}{9}\selectfont\thepage}
\fancyhead[LO,RE]{\fontsize{7.5}{9}\selectfont ICPC 算法手册\quad\nouppercase{\leftmark}}
\renewcommand{\headrulewidth}{0.2pt}
\renewcommand{\sectionmark}[1]{\markboth{\thesection\ #1}{}}
\lstset{language=C++,basicstyle=\ttfamily\fontsize{7.8}{8.8}\selectfont,
    keywordstyle=\bfseries,commentstyle=\color{ink},stringstyle=\color{ink},
    numbers=none,frame=none,columns=fullflexible,keepspaces=true,
    tabsize=4,showstringspaces=false,showspaces=false,escapeinside={(*@}{@*)},
    breaklines=true,breakatwhitespace=false,breakindent=1em,
    aboveskip=2pt,belowskip=3pt,
    literate={Σ}{{$\Sigma$}}1 {∑}{{$\sum$}}1 {Δ}{{$\Delta$}}1
    {α}{{$\alpha$}}1 {π}{{$\pi$}}1 {φ}{{$\varphi$}}1 {θ}{{$\theta$}}1
    {μ}{{$\mu$}}1 {σ}{{$\sigma$}}1 {Π}{{$\Pi$}}1 {∏}{{$\prod$}}1
    {≤}{{$\le$}}1 {≥}{{$\ge$}}1 {≠}{{$\ne$}}1 {≈}{{$\approx$}}1
    {×}{{$\times$}}1 {∈}{{$\in$}}1 {∞}{{$\infty$}}1
    {≡}{{$\equiv$}}1 {⊆}{{$\subseteq$}}1
    {→}{{$\to$}}1 {⇒}{{$\Rightarrow$}}1 {⊕}{{$\oplus$}}1
    {²}{{$^2$}}1 {³}{{$^3$}}1 {√}{{$\sqrt{\ }$}}1
}
\newcommand{\topic}[1]{\par\smallskip\textbf{#1}\par\nobreak}
\newcommand{\note}[1]{{\fontsize{8}{9.5}\selectfont #1\par}}
\begin{document}
\fontsize{8.3}{9.8}\selectfont

\begin{center}{\fontsize{15}{17}\selectfont\bfseries ICPC 算法手册}\quad\small 2026-10-05\end{center}
\note{双栏完整实现版\quad 197 份代码模板、1 份速查表。各模板独立使用；同名全局量和函数按题目取舍，不将整本直接拼接编译。公共头文件与 \texttt{ll} 定义仅在此列出，其他容量、类型和依赖保留在各模板中。}
\begin{lstlisting}
#include<bits/stdc++.h>
using namespace std;
typedef long long ll;
\end{lstlisting}
\begin{multicols}{2}\tableofcontents\end{multicols}
\clearpage\begin{multicols}{2}
\section{基础与技巧}
% SOURCE: 01-基础与技巧/STL容器速查.md
\subsection{STL容器速查}
\note{下标从 0 起，迭代器区间为 \texttt{[\allowbreak{}l,\allowbreak{}r)\allowbreak{}}；读取端点或弹出前保证非空，不解引用 \texttt{end(\allowbreak{})\allowbreak{}}。各容器共有的 \texttt{size(\allowbreak{})\allowbreak{}}、\texttt{empty(\allowbreak{})\allowbreak{}} 不逐一重复。下文 \texttt{n} 为元素数，\texttt{k} 为所取数量。}
\topic{vector 与区间修改}
\begingroup\fontsize{7.8}{9}\selectfont
\renewcommand{\arraystretch}{1.05}
\begin{tabularx}{\linewidth}{@{}>{\raggedright\arraybackslash}p{0.45\linewidth}@{\hspace{4pt}}>{\raggedright\arraybackslash}X@{}}\hline
\textbf{写法} & \textbf{作用、条件与复杂度}\\\hline
\texttt{v.\allowbreak{}reserve(\allowbreak{}n)\allowbreak{}} /\allowbreak{} \texttt{v.\allowbreak{}resize(\allowbreak{}n,\allowbreak{}x)\allowbreak{}} & 前者只预留容量，不增加 size；后者改变 size，新增元素填 x。\\[2pt]
\texttt{v.\allowbreak{}erase(\allowbreak{}l,\allowbreak{}r)\allowbreak{}} /\allowbreak{} \texttt{v.\allowbreak{}insert(\allowbreak{}it,\allowbreak{}x)\allowbreak{}} & 删除半开迭代器区间 /\allowbreak{} 在 it 前插入。移动后续元素，O(\allowbreak{}n)\allowbreak{}。\\[2pt]
\texttt{v.\allowbreak{}erase(\allowbreak{}remove(\allowbreak{}v.\allowbreak{}begin(\allowbreak{})\allowbreak{},\allowbreak{}v.\allowbreak{}end(\allowbreak{})\allowbreak{},\allowbreak{}x)\allowbreak{},\allowbreak{}v.\allowbreak{}end(\allowbreak{})\allowbreak{})\allowbreak{}} & 删除所有等于 x 的项，O(\allowbreak{}n)\allowbreak{}；remove 只移动元素，须配合 erase。\\[2pt]
迭代器失效 & vector 扩容使全部迭代器、指针、引用失效；未扩容的插入及删除使操作位置及之后失效。clear 不释放容量。\\[2pt]
\hline\end{tabularx}\par\endgroup
\topic{排序、查找与数值算法}
\begingroup\fontsize{7.8}{9}\selectfont
\renewcommand{\arraystretch}{1.05}
\begin{tabularx}{\linewidth}{@{}>{\raggedright\arraybackslash}p{0.45\linewidth}@{\hspace{4pt}}>{\raggedright\arraybackslash}X@{}}\hline
\textbf{写法} & \textbf{作用、条件与复杂度}\\\hline
\texttt{sort(\allowbreak{}l,\allowbreak{}r,\allowbreak{}cmp)\allowbreak{}} /\allowbreak{} \texttt{stable\_sort(\allowbreak{}l,\allowbreak{}r,\allowbreak{}cmp)\allowbreak{}} & 排序 /\allowbreak{} 保留等价元素原顺序。sort 为 O(\allowbreak{}n log n)\allowbreak{}；stable\_sort 有足够辅助内存时相同，否则 O(\allowbreak{}n log\ensuremath{^2} n)\allowbreak{}。\\[2pt]
\texttt{v.\allowbreak{}erase(\allowbreak{}unique(\allowbreak{}v.\allowbreak{}begin(\allowbreak{})\allowbreak{},\allowbreak{}v.\allowbreak{}end(\allowbreak{})\allowbreak{})\allowbreak{},\allowbreak{}v.\allowbreak{}end(\allowbreak{})\allowbreak{})\allowbreak{}} & 去相邻重复，O(\allowbreak{}n)\allowbreak{}。全局去重须先排序；unique 不改变 size。\\[2pt]
\texttt{lower\_bound(\allowbreak{}l,\allowbreak{}r,\allowbreak{}x)\allowbreak{}} /\allowbreak{} \texttt{upper\_bound(\allowbreak{}l,\allowbreak{}r,\allowbreak{}x)\allowbreak{}} & 非降序区间中首个 \ensuremath{\ge}x /\allowbreak{} >\allowbreak{}x 的迭代器；随机访问时 O(\allowbreak{}log n)\allowbreak{}。降序数据须同时传匹配比较器。\\[2pt]
\texttt{equal\_range(\allowbreak{}l,\allowbreak{}r,\allowbreak{}x)\allowbreak{}} & 返回等于 x 的半开区间；vector 中两端相减就是出现次数。\\[2pt]
\texttt{nth\_element(\allowbreak{}l,\allowbreak{}l+\allowbreak{}k-\allowbreak{}1,\allowbreak{}r)\allowbreak{}} & 1\ensuremath{\le}k\ensuremath{\le}n；第 k 小就位，两侧分别不大于 /\allowbreak{} 不小于它，内部无序，平均 O(\allowbreak{}n)\allowbreak{}。\\[2pt]
\texttt{partial\_sort(\allowbreak{}l,\allowbreak{}l+\allowbreak{}k,\allowbreak{}r)\allowbreak{}} & 最小 k 项排好序，其余无序；0\ensuremath{\le}k\ensuremath{\le}n，O(\allowbreak{}n log k)\allowbreak{}，k\ensuremath{\le}1 时 O(\allowbreak{}n)\allowbreak{}。\\[2pt]
\texttt{next\_permutation(\allowbreak{}l,\allowbreak{}r)\allowbreak{}} & 下一个字典序排列，O(\allowbreak{}n)\allowbreak{}；枚举全部排列须先升序。重复元素自然去重，返回 false 时重置为最小排列。\\[2pt]
\texttt{accumulate(\allowbreak{}l,\allowbreak{}r,\allowbreak{}0LL)\allowbreak{}} & 累加，O(\allowbreak{}n)\allowbreak{}；初值决定累加类型，写 0 会按 int 累加。\\[2pt]
\texttt{rotate(\allowbreak{}l,\allowbreak{}m,\allowbreak{}r)\allowbreak{}} & 把 \texttt{[\allowbreak{}m,\allowbreak{}r)\allowbreak{}} 移到 \texttt{[\allowbreak{}l,\allowbreak{}m)\allowbreak{}} 前，两段内部顺序不变，O(\allowbreak{}n)\allowbreak{}。\\[2pt]
\hline\end{tabularx}\par\endgroup
\topic{set / multiset / map}
\begingroup\fontsize{7.8}{9}\selectfont
\renewcommand{\arraystretch}{1.05}
\begin{tabularx}{\linewidth}{@{}>{\raggedright\arraybackslash}p{0.45\linewidth}@{\hspace{4pt}}>{\raggedright\arraybackslash}X@{}}\hline
\textbf{写法} & \textbf{作用、条件与复杂度}\\\hline
\texttt{s.\allowbreak{}lower\_bound(\allowbreak{}x)\allowbreak{}} /\allowbreak{} \texttt{s.\allowbreak{}upper\_bound(\allowbreak{}x)\allowbreak{}} & 首个 \ensuremath{\ge}x /\allowbreak{} >\allowbreak{}x；O(\allowbreak{}log n)\allowbreak{}。必须用成员函数，通用二分在树迭代器上移动为 O(\allowbreak{}n)\allowbreak{}。\\[2pt]
\texttt{it=\allowbreak{}s.\allowbreak{}lower\_bound(\allowbreak{}x)\allowbreak{}}；\texttt{prev(\allowbreak{}it)\allowbreak{}} & it 是 \ensuremath{\ge}x 的后继（需 it!=\allowbreak{}end）；prev(\allowbreak{}it)\allowbreak{} 是 <\allowbreak{}x 的前驱（需 it!=\allowbreak{}begin），包括 it=\allowbreak{}end 的情况。\\[2pt]
\texttt{ms.\allowbreak{}erase(\allowbreak{}x)\allowbreak{}} /\allowbreak{} \texttt{ms.\allowbreak{}erase(\allowbreak{}it)\allowbreak{}} & 前者删除所有相等项，O(\allowbreak{}log n+\allowbreak{}个数)\allowbreak{}；后者只删一项，均摊 O(\allowbreak{}1)\allowbreak{}，it 须有效。\\[2pt]
\texttt{ms.\allowbreak{}count(\allowbreak{}x)\allowbreak{}} /\allowbreak{} \texttt{ms.\allowbreak{}equal\_range(\allowbreak{}x)\allowbreak{}} & 计数 /\allowbreak{} 相等项区间，O(\allowbreak{}log n+\allowbreak{}个数)\allowbreak{} /\allowbreak{} O(\allowbreak{}log n)\allowbreak{}。遍历或 distance 该区间仍需 O(\allowbreak{}个数)\allowbreak{}。\\[2pt]
\texttt{mp[\allowbreak{}x]\allowbreak{}} /\allowbreak{} \texttt{mp.\allowbreak{}find(\allowbreak{}x)\allowbreak{}} & 前者在缺键时插入默认值；后者只查找，缺键返回 end。均为 O(\allowbreak{}log n)\allowbreak{}。\\[2pt]
\texttt{it=\allowbreak{}s.\allowbreak{}erase(\allowbreak{}it)\allowbreak{}} & 删除时取返回的下一迭代器；树容器插入不使旧迭代器失效，删除仅使被删元素失效。\\[2pt]
\hline\end{tabularx}\par\endgroup
\topic{比较与堆}
\begingroup\fontsize{7.8}{9}\selectfont
\renewcommand{\arraystretch}{1.05}
\begin{tabularx}{\linewidth}{@{}>{\raggedright\arraybackslash}p{0.45\linewidth}@{\hspace{4pt}}>{\raggedright\arraybackslash}X@{}}\hline
\textbf{写法} & \textbf{作用、条件与复杂度}\\\hline
\texttt{pair} /\allowbreak{} \texttt{tuple} /\allowbreak{} \texttt{tie(\allowbreak{}a.\allowbreak{}x,\allowbreak{}a.\allowbreak{}y)\allowbreak{}} & 默认字典序；tie 可省掉手写多关键字比较。\texttt{auto [\allowbreak{}x,\allowbreak{}y]\allowbreak{}=\allowbreak{}p} 是复制，\texttt{auto \&[\allowbreak{}x,\allowbreak{}y]\allowbreak{}=\allowbreak{}p} 可修改原对象。\\[2pt]
\texttt{return tie(\allowbreak{}a.\allowbreak{}x,\allowbreak{}a.\allowbreak{}y)\allowbreak{}<\allowbreak{}tie(\allowbreak{}b.\allowbreak{}x,\allowbreak{}b.\allowbreak{}y)\allowbreak{};} & sort /\allowbreak{} set 按 (\allowbreak{}x,\allowbreak{}y)\allowbreak{} 升序；比较器必须严格弱序，相等时返回 false，不能用 \ensuremath{\le}。set 以 cmp(\allowbreak{}a,\allowbreak{}b)\allowbreak{} 与 cmp(\allowbreak{}b,\allowbreak{}a)\allowbreak{} 均为 false 判等。\\[2pt]
\texttt{priority\_queue<\allowbreak{}T>\allowbreak{} q} & 默认大根堆；top 为 O(\allowbreak{}1)\allowbreak{}，push /\allowbreak{} pop 为 O(\allowbreak{}log n)\allowbreak{}，pop 不返回值。\\[2pt]
\texttt{priority\_queue<\allowbreak{}T,\allowbreak{}vector<\allowbreak{}T>\allowbreak{},\allowbreak{}greater<\allowbreak{}T>\allowbreak{}>\allowbreak{} q} & 小根堆；T 为 pair 时，先取 first 小，再取 second 小。\\[2pt]
\texttt{return tie(\allowbreak{}a.\allowbreak{}x,\allowbreak{}a.\allowbreak{}y)\allowbreak{}>\allowbreak{}tie(\allowbreak{}b.\allowbreak{}x,\allowbreak{}b.\allowbreak{}y)\allowbreak{};} & 堆的比较器表示 a 优先级低于 b；此式使 (\allowbreak{}x,\allowbreak{}y)\allowbreak{} 最小者在顶。堆不能按值删除或直接修改；可另记有效性并懒删除。\\[2pt]
\hline\end{tabularx}\par\endgroup
\topic{哈希、字符串与其他容器}
\begingroup\fontsize{7.8}{9}\selectfont
\renewcommand{\arraystretch}{1.05}
\begin{tabularx}{\linewidth}{@{}>{\raggedright\arraybackslash}p{0.45\linewidth}@{\hspace{4pt}}>{\raggedright\arraybackslash}X@{}}\hline
\textbf{写法} & \textbf{作用、条件与复杂度}\\\hline
\texttt{h.\allowbreak{}max\_load\_factor(\allowbreak{}0.\allowbreak{}7)\allowbreak{}}；\texttt{h.\allowbreak{}reserve(\allowbreak{}n)\allowbreak{}} & unordered 容器先设负载因子，再按元素数预留；查增删平均 O(\allowbreak{}1)\allowbreak{}、最坏 O(\allowbreak{}n)\allowbreak{}，无顺序。不是抗碰撞保证。\\[2pt]
rehash 的失效规则 & rehash 使迭代器失效，元素引用 /\allowbreak{} 指针仍有效；插入可能触发 rehash，删除仅使被删元素失效。\\[2pt]
\texttt{s.\allowbreak{}substr(\allowbreak{}pos,\allowbreak{}len)\allowbreak{}} /\allowbreak{} \texttt{s.\allowbreak{}erase(\allowbreak{}pos,\allowbreak{}len)\allowbreak{}} & 第二参数是长度；到末尾为止，不是右端点。substr 复制结果，erase 移动后续字符。\\[2pt]
\texttt{s.\allowbreak{}find(\allowbreak{}t,\allowbreak{}pos)\allowbreak{}} & 从 pos 开始查找子串，失败用 \texttt{string::npos} 判断，不存入普通 int 后与 -\allowbreak{}1 混用。\\[2pt]
\texttt{deque} /\allowbreak{} \texttt{stack} /\allowbreak{} \texttt{queue} & deque 两端增删 O(\allowbreak{}1)\allowbreak{}，支持下标；stack 用 top，queue 用 front /\allowbreak{} back，push /\allowbreak{} pop 不返回元素。deque 两端插入使迭代器失效，旧元素引用仍有效；中间插删使两者均失效。\\[2pt]
\texttt{bitset<\allowbreak{}N>\allowbreak{}} & 位编号从低位 0 起，容量是编译期常量；\texttt{to\_ullong(\allowbreak{})\allowbreak{}} 的数值装不下会抛异常。批量位运算与 DP 用法见 bitset 技巧。\\[2pt]
\hline\end{tabularx}\par\endgroup
% SOURCE: 01-基础与技巧/快读快写.cpp
\subsection{快读快写}

\note{输入输出对象维护缓冲区、当前指针与有效长度；read把结果写到传入变量，write输出一个值。普通批量题按块缓冲，交互题每轮显式flush；混用其他IO时须考虑各自缓冲。}
\note{按已知数量读取合法整数，输入须在 ll 范围内；EOF 返回 0，不能据此区分文件结束与整数零。}
\begin{lstlisting}
inline ll readd()
{
    int c;
    bool flag=false;
    while((c=getchar())>'9'||c<'0')
    {
        if(c==EOF)return 0;
        if(c=='-')flag=true;
    }
    unsigned long long res=c-'0';
    while((c=getchar())>='0'&&c<='9')res=(res<<3)+(res<<1)+c-'0';
    if(flag)return res==(1ULL<<63)?LLONG_MIN:-(ll)res;
    return (ll)res;
}
void print(ll x)
{
    if(x<0)
    {
        putchar('-');
        if(x<-9)print(-(x/10)); // 先除再取负，兼容 LLONG_MIN
        putchar('0'-x%10);
    }
    else
    {
        if(x>9)print(x/10);
        putchar(x%10+'0');
    }
}
\end{lstlisting}
\note{用法：n=\allowbreak{}readd(\allowbreak{})\allowbreak{}; print(\allowbreak{}ans)\allowbreak{}; putchar(\allowbreak{}'\textbackslash{}n')\allowbreak{}; 不自动输出空格或换行。 无锁版本可按平台替换 getchar/\allowbreak{}putchar，参见 D:\textbackslash{}code\textbackslash{}luogu\textbackslash{}P10815.\allowbreak{}cpp。}
% SOURCE: 01-基础与技巧/离散化.cpp
\subsection{离散化}

\note{a[\allowbreak{}1.\allowbreak{}.\allowbreak{}n]\allowbreak{}是原值，b是升序去重值表，id[\allowbreak{}i]\allowbreak{}是从1起的排名；排名r对应b[\allowbreak{}r-\allowbreak{}1]\allowbreak{}。调用discrete(\allowbreak{})\allowbreak{}重建b和id，a不变；离散化只保留顺序，不保留数值距离。}
\note{排序去重 O(\allowbreak{}n log n)\allowbreak{}，映射 O(\allowbreak{}log n)\allowbreak{}；保留大小关系，不保留实际距离。}
\begin{lstlisting}
const int N=1000005;
int n,id[N];
ll a[N];
vector<ll> b;
void discrete()
{
    b.clear();
    for(int i=1;i<=n;i++)b.push_back(a[i]);
    sort(b.begin(),b.end());
    b.erase(unique(b.begin(),b.end()),b.end());
    for(int i=1;i<=n;i++)
        id[i]=lower_bound(b.begin(),b.end(),a[i])-b.begin()+1;
}
\end{lstlisting}
\note{id[\allowbreak{}i]\allowbreak{} 是 1-\allowbreak{}based 排名，原数组 a 不变；排名 r 的原值为 b[\allowbreak{}r-\allowbreak{}1]\allowbreak{}。}
% SOURCE: 01-基础与技巧/双指针.cpp
\subsection{双指针}

\note{a从0起且非降序，two\_sum(\allowbreak{}target)\allowbreak{}返回两个不同位置；s为文本，t为要覆盖的字符多重集，cnt按字节记窗口缺额。min\_cover(\allowbreak{})\allowbreak{}每次重置cnt，返回最短覆盖长度；输入不变。}
\note{适用：有序数组两数和、满足条件的最短连续窗口；各指针只向一个方向走。 vector/\allowbreak{}string 都为 0-\allowbreak{}indexed；指针单调保证总扫描线性，窗口计数允许为负表示多余。}
\begin{lstlisting}
vector<int> a;
string s,t;
int cnt[256];
\end{lstlisting}
\note{O(\allowbreak{}n)\allowbreak{}、额外空间 O(\allowbreak{}1)\allowbreak{}；a 必须非降序，target 为目标和，返回两个不同的 0-\allowbreak{}based 下标。}
\begin{lstlisting}
pair<int,int> two_sum(int target)
{
    int l=0,r=(int)a.size()-1;
    while(l<r)
    {
        long long sum=(long long)a[l]+a[r];
        if(sum==target)return {l,r};
        if(sum<target)l++;
        else r--;
    }
    return {-1,-1};
}
\end{lstlisting}
\note{O(\allowbreak{}|s|+\allowbreak{}|t|)\allowbreak{}、空间 O(\allowbreak{}256)\allowbreak{}；s 是文本、t 是需求多重集，窗口 [\allowbreak{}l,\allowbreak{}r]\allowbreak{} 两端均包含。 unsigned char 避免高位字节作为负下标；先扩大至满足，再收缩直到缺少一个所需字符。}
\begin{lstlisting}
int min_cover()
{
    memset(cnt,0,sizeof cnt); // 每次调用重置窗口计数
    if(t.empty())return 0;
    int need=t.size(),l=0,ans=INT_MAX;
    for(unsigned char c:t)cnt[c]++;
    for(int r=0;r<(int)s.size();r++)
    {
        unsigned char c=s[r];
        if(cnt[c]-->0)need--;
        while(need==0)
        {
            ans=min(ans,r-l+1);
            unsigned char d=s[l++];
            if(++cnt[d]>0)need++;
        }
    }
    return ans==INT_MAX?-1:ans;
}
\end{lstlisting}
% SOURCE: 01-基础与技巧/单调栈.cpp
\subsection{单调栈}

\note{a[\allowbreak{}1.\allowbreak{}.\allowbreak{}n]\allowbreak{}是输入，st[\allowbreak{}1.\allowbreak{}.\allowbreak{}top]\allowbreak{}存候选下标，ans[\allowbreak{}i]\allowbreak{}存相邻更大/\allowbreak{}更小位置；不存在时用代码对应的0或n+\allowbreak{}1哨兵。每次扫描重置top；改比较符前先确定相等元素是否阻挡。}
\begin{lstlisting}
const int N=1000005;
int n,a[N],st[N],top,ans[N];
\end{lstlisting}
\note{O(\allowbreak{}n)\allowbreak{}，右边第一个严格更大的下标，不存在为 0。}
\begin{lstlisting}
void right_greater()
{
    top=0;
    fill(ans+1,ans+n+1,0);
    for(int i=1;i<=n;i++)
    {
        while(top&&a[st[top]]<a[i])ans[st[top--]]=i;
        st[++top]=i;
    }
}
\end{lstlisting}
\note{O(\allowbreak{}n)\allowbreak{}，左边第一个严格更小的下标，不存在为 0。}
\begin{lstlisting}
void left_less()
{
    top=0;
    for(int i=1;i<=n;i++)
    {
        while(top&&a[st[top]]>=a[i])top--;
        ans[i]=top?st[top]:0;
        st[++top]=i;
    }
}
\end{lstlisting}
\note{右侧写法：弹栈时给旧元素赋答案；左侧写法：弹栈后给当前元素赋答案。 修改比较方向可求较大/\allowbreak{}较小；是否包含等号由题意决定。}
\note{柱状图最大矩形：a[\allowbreak{}1.\allowbreak{}.\allowbreak{}n]\allowbreak{} 为非负高度，柱宽为 1，O(\allowbreak{}n)\allowbreak{}。 弹出 j 时，右边界 i、左边界为弹出后的栈顶 L，可覆盖 [\allowbreak{}L+\allowbreak{}1,\allowbreak{}i-\allowbreak{}1]\allowbreak{}。}
\begin{lstlisting}
ll histogram()
{
    top=0;
    ll res=0;
    for(int i=1;i<=n+1;i++)
    {
        while(top&&(i==n+1||a[st[top]]>=a[i]))
        {
            int j=st[top--],L=top?st[top]:0;
            res=max(res,1LL*a[j]*(i-L-1));
        }
        if(i<=n)st[++top]=i;
    }
    return res;
}
\end{lstlisting}
\note{n+\allowbreak{}1 只负责清空栈，不访问 a[\allowbreak{}n+\allowbreak{}1]\allowbreak{}；等高柱合并后由更靠右的柱继续扩展。}
\note{所有非空子数组最小值之和，$O(n)$。对位置 j，L 是左侧首个严格更小位置，R 是右侧首个小于等于的位置。两端分别有 j-L 与 R-j 种选择：}\[\operatorname{contrib}_j=a_j(j-L)(R-j).\]\note{乘积及总和须在 \texttt{ll} 范围内。}
\begin{lstlisting}
ll sum_min()
{
    top=0;
    ll res=0;
    for(int i=1;i<=n+1;i++)
    {
        while(top&&(i==n+1||a[st[top]]>=a[i]))
        {
            int j=st[top--],L=top?st[top]:0;
            res+=1LL*a[j]*(j-L)*(i-j);
        }
        if(i<=n)st[++top]=i;
    }
    return res;
}
\end{lstlisting}
\note{一侧严格、一侧非严格，把相同最小值归给最右位置，避免重复或遗漏。 求最大值之和：>\allowbreak{}=\allowbreak{} 改为 <\allowbreak{}=\allowbreak{}；子数组极差之和 =\allowbreak{} 最大值之和 -\allowbreak{} 最小值之和。}
% SOURCE: 01-基础与技巧/单调队列.cpp
\subsection{单调队列}

\topic{等长区间平移：差集配对}
\note{旧区间 $[s,s+L)$ 右移 $d$，$0\le d\le L$。重叠部分不变；移出 $[s,s+d)$，移入 $[s+L,s+L+d)$。记移出、移入边际收益为 $u_i,v_i$，则}
\[\Delta(s,d)=\sum_{i=s}^{s+d-1}(u_i+v_{i+L}).\]
\note{例如最大化恰被覆盖 $k$ 次的位置数，原覆盖次数为 $c_i$：$u_i=[c_i=k+1]-[c_i=k]$，$v_i=[c_i=k-1]-[c_i=k]$。令 $z_i=u_i+v_{i+L}$，允许 $d\le L$ 时变成长度至多 $L$ 的最大子段和，用前缀最小值队列。保留不移动的零收益；$d>L$ 为不交情形，须另算（2025 成都 K）。}

\note{a[\allowbreak{}1.\allowbreak{}.\allowbreak{}n]\allowbreak{}存数据，k是窗口长度，ans存窗口答案；DP示例f[\allowbreak{}i]\allowbreak{}是以i结束的状态。队列保存候选下标，失效条件由允许的转移范围决定，淘汰条件由候选比较决定；换题替换这两个条件。}
\note{q 存候选下标；每个候选最多入队、出队一次，总 O(\allowbreak{}n)\allowbreak{}。 适用：候选按入队顺序失效，过期后不再有效；队首始终是当前最优候选。 新候选在旧候选剩余有效期内始终有效且不劣，才能永久删除旧候选。 若候选优劣随查询改变、可能反转，不能直接套这个循环。}
\begin{lstlisting}
const int N=1000005;
int n,k,ans[N];
ll a[N],f[N];
\end{lstlisting}
\note{窗口最小值，下标从 1 开始，1<\allowbreak{}=\allowbreak{}k<\allowbreak{}=\allowbreak{}n；先加入 i，再查询包含 i 的窗口。}
\begin{lstlisting}
void solve()
{
    deque<int> q;
    for(int i=1;i<=n;i++)
    {
        while(!q.empty()&&q.front()<=i-k)q.pop_front();
        while(!q.empty()&&a[q.back()]>=a[i])q.pop_back();
        q.push_back(i);
        ans[i]=q.front(); // 记录当前最优候选下标；具体取值/转移在这里改
    }
}
\end{lstlisting}
\note{窗口 [\allowbreak{}max(\allowbreak{}1,\allowbreak{}i-\allowbreak{}k+\allowbreak{}1)\allowbreak{},\allowbreak{}i]\allowbreak{}；求最大值把队尾比较 >\allowbreak{}=\allowbreak{} 改为 <\allowbreak{}=\allowbreak{}。 上述非严格比较保留相等值中最新的下标；改成 >\allowbreak{} /\allowbreak{} <\allowbreak{} 则保留最早的下标。 ans[\allowbreak{}i]\allowbreak{} 存下标，窗口值为 a[\allowbreak{}ans[\allowbreak{}i]\allowbreak{}]\allowbreak{}；只要完整窗口时，i>\allowbreak{}=\allowbreak{}k 才记录答案。 自定义题目：将队首条件换成 expired(\allowbreak{}id,\allowbreak{}i)\allowbreak{}，队尾条件换成 check(\allowbreak{}old,\allowbreak{}now)\allowbreak{}。 expired 判断是否失效，check 判断新候选能否永久淘汰旧候选，须满足上面的适用条件。}
\topic{有界前驱 DP}\note{a 从 1 起，k 至少为 1；先过期、再转移、最后加入 i。初始候选 0 不可漏，不能让 i 转移到自己。}\[f_0=0,\qquad f_i=a_i+\max_{\max(0,i-k)\le j<i}f_j.\]
\begin{lstlisting}
void dp()
{
    f[0]=0; // 每次调用从初始状态重算 f[1..n]
    deque<int> q{0};
    for(int i=1;i<=n;i++)
    {
        while(!q.empty()&&q.front()<i-k)q.pop_front();
        f[i]=f[q.front()]+a[i];
        while(!q.empty()&&f[q.back()]<=f[i])q.pop_back();
        q.push_back(i);
    }
}
\end{lstlisting}
\note{按题目改转移值、有效区间及队尾比较；只把可达状态入队，空队时当前状态不可达。}
% SOURCE: 01-基础与技巧/位运算技巧.cpp
\subsection{位运算技巧}

\note{mask表示选中位的集合，sub是其子集，n为有效位数；宽度按uint32/\allowbreak{}uint64区分。枚举函数的参数是当前掩码或位数等标量；零值、最高位及移位等于字长的边界见各模块正文。}
\note{整数掩码的枚举与计数、异或区间和。布尔集合批量操作见bitset与字并行，所有子集聚合见SOS，异或空间见线性基，避免同一功能散放多处。}
\topic{类型与位宽}
\note{掩码用ull，编号从0开始。64位掩码只能左移0.\allowbreak{}.\allowbreak{}63位；构造全集的枚举接口限定bits\ensuremath{\le}63，不能写1ULL<\allowbreak{}<\allowbreak{}64。clz/\allowbreak{}ctz的输入必须非零；下方bit\_length显式处理0。有符号负数不参与这些掩码移位。}
\begin{lstlisting}
typedef unsigned long long ull;
\end{lstlisting}
\topic{置位查询}
\note{lowbit取最低置位的权值，bit\_count计置位数，bit\_length取有效位宽，is\_pow2判正的2次幂。置位编号用ctz；从低到高枚举时反复清掉最低置位，时间与置位数相同。这些短操作集中在此，不在卡常里重复。}
\begin{lstlisting}
ull lowbit(ull x){return x&(-x);}
bool is_pow2(ull x){return x>0&&(x&(x-1))==0;}
int bit_count(ull x){return __builtin_popcountll(x);}
int bit_length(ull x){return x?64-__builtin_clzll(x):0;}
\end{lstlisting}
\begin{lstlisting}
vector<int> set_bits(ull x)
{
    vector<int> res;
    for(;x;x&=x-1)res.push_back(__builtin_ctzll(x));
    return res;
}
\end{lstlisting}
\topic{子集、超集与枚举规模}
\note{submasks含mask和空集，共 $2^{\operatorname{popcount}(mask)}$ 项；mask=\allowbreak{}0也只产生一个0。更新式在0后会回到mask，所以必须先处理再判断结束。supermasks在bits位全集内枚举，要求 $0\le bits\le63$ 且mask不超出全集。两种函数的push\_back可直接替换为题目操作，不必存结果。}
\begin{lstlisting}
vector<ull> submasks(ull mask)
{
    vector<ull> res;
    for(ull s=mask;;s=(s-1)&mask)
    {
        res.push_back(s);
        if(s==0)break;
    }
    return res;
}
\end{lstlisting}
\begin{lstlisting}
vector<ull> supermasks(ull mask,int bits)
{
    ull all=(1ULL<<bits)-1;
    vector<ull> res;
    for(ull s=mask;;s=(s+1)|mask)
    {
        res.push_back(s);
        if(s==all)break;
    }
    return res;
}
\end{lstlisting}
\note{对所有mask再枚举其子集，每位有“不在mask、只在mask、也在子集”三种状态，总次数 $3^{bits}$。若只需要mask的子集中同时包含于allow的数量，每个交集位自由选，答案为 $2^{\operatorname{popcount}(mask\&allow)}$；下方接口使用32位unsigned，结果放ull。}
\begin{lstlisting}
unsigned long long count_subset(unsigned mask,unsigned allow){return 1ULL<<__builtin_popcount(mask&allow);}
\end{lstlisting}
\topic{固定置位数：Gosper}
\note{next\_combination返回下一更大的同置位数掩码，0或无法继续用ull表示时返回0。combinations枚举bits位中恰选k位， $0\le k\le bits\le63$ ，时间为组合数；k=\allowbreak{}0单独返回空集，避免除以lowbit(\allowbreak{}0)\allowbreak{}。}
\begin{lstlisting}
ull next_combination(ull x)
{
    ull t=lowbit(x),y=x+t;
    return y?y|(((x^y)>>2)/t):0;
}
vector<ull> combinations(int bits,int k)
{
    if(k==0)return {0};
    vector<ull> res;
    ull limit=1ULL<<bits;
    for(ull s=(1ULL<<k)-1;s&&s<limit;s=next_combination(s))res.push_back(s);
    return res;
}
\end{lstlisting}
\topic{连续整数的异或}
\note{0到n的异或按n模4循环为n、1、n+\allowbreak{}1、0。闭区间[\allowbreak{}l,\allowbreak{}r]\allowbreak{}用两个前缀异或；l=\allowbreak{}0直接取右前缀，避免无符号下溢。}
\begin{lstlisting}
ull xor_prefix(ull n){return n%4==0?n:(n%4==1?1:(n%4==2?n+1:0));}
\end{lstlisting}
\topic{对齐块异或后的算术和}
\note{块 $[b,b+L)$ 满足 $L=2^k$、b是L的倍数，整个块在ull值域内。低k位异或后仍是一次全排列，清掉v的低k位后令 $H=b\mathbin{\mathtt{xor}}(v\mathbin{\mathtt{\&}}\mathord{\sim}(L-1))$ ，则总和为 $LH+L(L-1)/2$ 。乘积与答案放128位；这项是算术和，与上一项区间异或不同。}
\begin{lstlisting}
unsigned __int128 xor_block_sum(ull b,ull len,ull v)
{
    assert(is_pow2(len)&&b%len==0&&b<=ULLONG_MAX-(len-1));
    return (unsigned __int128)len*(b^(v&~(len-1)))+(unsigned __int128)len*(len-1)/2;
}
\end{lstlisting}
\topic{关联模型}
\note{SOS是在全部掩码上逐位聚合的DP，不放在枚举工具中；子集和、超集和及逆变换统一使用第六章SOS接口。GF(\allowbreak{}2)\allowbreak{}线性空间使用第五章线性基；固定掩码AND/\allowbreak{}OR/\allowbreak{}XOR与循环移位的仿射复合保留在赛事建模专题。}
% SOURCE: 01-基础与技巧/逆序对.cpp
\subsection{逆序对}

\note{a是待统计数组，tmp为同容量归并缓冲；count\_inv(\allowbreak{}l,\allowbreak{}r)\allowbreak{}使用闭区间，返回严格逆序对数并把a[\allowbreak{}l.\allowbreak{}.\allowbreak{}r]\allowbreak{}排序。需要再次统计原序列时先恢复a，答案用ll。}
\note{O(\allowbreak{}n log n)\allowbreak{}，统计严格逆序对；会将 a[\allowbreak{}l.\allowbreak{}.\allowbreak{}r]\allowbreak{} 排好序。}
\begin{lstlisting}
const int N=1000005;
int a[N],tmp[N];
ll count_inv(int l,int r)// 相等不算逆序对
{
    if(l>=r)return 0;
    int m=(l+r)>>1;
    ll res=count_inv(l,m)+count_inv(m+1,r);
    int i=l,j=m+1,cnt=l;
    while(i<=m&&j<=r)
    {
        if(a[i]<=a[j])tmp[cnt++]=a[i++];
        else res+=m-i+1,tmp[cnt++]=a[j++];
    }
    while(i<=m)tmp[cnt++]=a[i++];
    while(j<=r)tmp[cnt++]=a[j++];
    for(int i=l;i<=r;i++)a[i]=tmp[i];
    return res;
}
\end{lstlisting}
% SOURCE: 01-基础与技巧/反悔贪心.cpp
\subsection{反悔贪心}

\note{jobs存Job\{d,\allowbreak{}p\}：截止时刻与单位时长任务收益；tasks存Task\{t,\allowbreak{}d\}：耗时与截止时刻。分别调用job\_schedule(\allowbreak{})\allowbreak{}、task\_schedule(\allowbreak{})\allowbreak{}，返回最大收益/\allowbreak{}最大件数；均原地按截止时刻排序，堆每次重新建立。}
\note{重复权值哈夫曼：最小 w 有 c 个时，批量 floor(\allowbreak{}c/\allowbreak{}2)\allowbreak{} 次合并，代价 2*w*floor(\allowbreak{}c/\allowbreak{}2)\allowbreak{}，生成 2*w。 c 为奇数时，余下一个 w 与当前次小配对；原权排序后与新生成的单调队列取两边最小。 若每种符号必须编码，零权符号也不能扔掉；计数与总代价用足够宽的整数。}
\note{按截止时间排序，堆保存已选任务；放不下时撤销最不值得保留的。 收益最大：踢最低收益，用小根堆；件数最多：踢最长耗时，用大根堆。 两种模型均为 O(\allowbreak{}n log n)\allowbreak{}，不能推广为任意耗时、任意收益的调度。}
\note{1.\allowbreak{} 单位时长，最大化收益（P2949）：d>\allowbreak{}=\allowbreak{}0、p>\allowbreak{}=\allowbreak{}0。 到截止时刻 d 最多选 d 件，小根堆维护已选任务的收益。}
\begin{lstlisting}
struct Job
{
    int d,p;
};
vector<Job> jobs;
ll job_schedule()
{
    sort(jobs.begin(),jobs.end(),[](const Job &a,const Job &b){return a.d<b.d;});
    priority_queue<int,vector<int>,greater<int> > q;
    ll ans=0;
    for(Job x:jobs)
    {
        q.push(x.p),ans+=x.p;
        if((int)q.size()>x.d)ans-=q.top(),q.pop(); // 连新任务一起比较，撤销最低收益
    }
    return ans;
}
\end{lstlisting}
\note{为什么能撤销：此前任务已可行，本轮只可能多出一件；删除任意一件都恢复可行。 同样腾出一个时段，删最低收益损失最小；后续截止时间不会更早。}
\note{2.\allowbreak{} 不同耗时，最大化件数（P4053）：t>\allowbreak{}0、d>\allowbreak{}=\allowbreak{}0，从时刻 0 起可安排任务。 tot 为已选总耗时；任务按截止时间顺序执行，大根堆保存耗时。}
\begin{lstlisting}
struct Task
{
    ll t,d;
};
vector<Task> tasks;
int task_schedule()
{
    sort(tasks.begin(),tasks.end(),[](const Task &a,const Task &b){return a.d<b.d;});
    priority_queue<ll> q;
    ll tot=0;
    for(Task x:tasks)
    {
        if(tot+x.t<=x.d)
        {
            tot+=x.t;
            q.push(x.t);
        }
        else if(!q.empty()&&q.top()>x.t)
        {
            tot+=x.t-q.top(); // 数量不变，总耗时变短，为后续任务留下空间
            q.pop();
            q.push(x.t);
        }
    }
    return (int)q.size();
}
\end{lstlisting}
\note{替换为何可行：移除旧任务后，其余任务只会提前；新任务放在当前末尾。 新耗时更短，所以新 tot 小于旧 tot，而旧 tot<\allowbreak{}=\allowbreak{}当前截止时间。 同样的件数，删最长任务能省下最多时间；新任务不更短时直接跳过。 例：(\allowbreak{}耗时,\allowbreak{}截止)\allowbreak{}=\allowbreak{}(\allowbreak{}4,\allowbreak{}4)\allowbreak{},\allowbreak{}(\allowbreak{}2,\allowbreak{}5)\allowbreak{},\allowbreak{}(\allowbreak{}3,\allowbreak{}6)\allowbreak{}：先选 4，再换成 2，最后可选 2 和 3。 不同耗时还要求收益最大时，以上两种堆策略都不保证最优。}
% SOURCE: 01-基础与技巧/记忆化搜索.cpp
\subsection{记忆化搜索}

\note{h[\allowbreak{}1.\allowbreak{}.\allowbreak{}n]\allowbreak{}[\allowbreak{}1.\allowbreak{}.\allowbreak{}m]\allowbreak{}是高度，f[\allowbreak{}i]\allowbreak{}[\allowbreak{}j]\allowbreak{}是从该格出发的最长下降路径；0表示尚未计算。dx、dy为四邻方向。高度严格下降保证无环，换地图清空f，搜索结果包含起点本身。}
\note{滑雪：从 (\allowbreak{}x,\allowbreak{}y)\allowbreak{} 出发只能走向严格更低的相邻格子，求最长路径。 状态：dfs(\allowbreak{}x,\allowbreak{}y)\allowbreak{} 为从该格出发的最长路径长度，包含当前格；转移取下一格答案 +\allowbreak{}1。 缓存下标须包含影响后续结果的全部信息；剩余次数、已选集合等若有影响也要入状态。 f[\allowbreak{}x]\allowbreak{}[\allowbreak{}y]\allowbreak{}=\allowbreak{}0 表示未计算；每个状态算一次，O(\allowbreak{}nm)\allowbreak{}，最坏递归深度 O(\allowbreak{}nm)\allowbreak{}。 本题答案至少为 1；答案可能为 0/\allowbreak{}负数时，用独立 vis 标记或不会与答案重合的哨兵。 高度严格下降保证无环；存在循环依赖时，缓存或访问标记本身不能解决状态转移。}
\begin{lstlisting}
const int N=1005;
int n,m,h[N][N],f[N][N];
int dx[4]={0,0,1,-1},dy[4]={1,-1,0,0};
int dfs(int x,int y)
{
    if(f[x][y])return f[x][y];
    f[x][y]=1;
    for(int i=0;i<4;i++)
    {
        int nx=x+dx[i],ny=y+dy[i];
        if(nx<1||nx>n||ny<1||ny>m||h[nx][ny]>=h[x][y])continue;
        f[x][y]=max(f[x][y],dfs(nx,ny)+1);
    }
    return f[x][y];
}
int solve()
{
    memset(f,0,sizeof f);
    int ans=0;
    for(int i=1;i<=n;i++)
        for(int j=1;j<=m;j++)ans=max(ans,dfs(i,j));
    return ans;
}
\end{lstlisting}
\note{多测或修改输入后须清缓存；长链可能爆栈，可按高度从低到高排序后迭代转移。}
% SOURCE: 01-基础与技巧/二维差分.cpp
\subsection{二维差分}

\note{n、m是矩阵行列数；a、d从1起，分别存原值和差分。第0行列为0，容量覆盖n+\allowbreak{}1、m+\allowbreak{}1；区间修改写d，还原后d变成实际值。重新处理数据前清空使用范围。}
\begin{lstlisting}
const int N=1005;
int n,m;
ll a[N][N],d[N][N]; // 1-indexed，预留第 0 行/列和 n+1、m+1
\end{lstlisting}
\note{原数组非零时先建差分；原数组全零时直接清 d，再做修改。}
\begin{lstlisting}
void build()
{
    memset(d,0,sizeof d);
    for(int i=1;i<=n;i++)
        for(int j=1;j<=m;j++)
            d[i][j]=a[i][j]-a[i-1][j]-a[i][j-1]+a[i-1][j-1];
}
\end{lstlisting}
\note{O(\allowbreak{}1)\allowbreak{}，闭矩形 [\allowbreak{}x1,\allowbreak{}x2]\allowbreak{} \ensuremath{\times} [\allowbreak{}y1,\allowbreak{}y2]\allowbreak{} 加 v。}
\begin{lstlisting}
void add(int x1,int y1,int x2,int y2,ll v)
{
    d[x1][y1]+=v;
    d[x2+1][y1]-=v;
    d[x1][y2+1]-=v;
    d[x2+1][y2+1]+=v;
}
\end{lstlisting}
\note{O(\allowbreak{}nm)\allowbreak{}，所有修改结束后还原一次，结果写入 a。}
\begin{lstlisting}
void restore()
{
    for(int i=1;i<=n;i++)
        for(int j=1;j<=m;j++)
        {
            d[i][j]+=d[i-1][j]+d[i][j-1]-d[i-1][j-1];
            a[i][j]=d[i][j];
        }
}
\end{lstlisting}
\note{流程：build -\allowbreak{}>\allowbreak{} 多次 add -\allowbreak{}>\allowbreak{} restore；不要重复 restore。}
% SOURCE: 01-基础与技巧/高维差分.cpp
\subsection{高维差分}

\note{令 $x=(x_1,\ldots,x_D)$，$e_i$ 为单位向量，$b\in\{0,1\}^D$，
$|b|=\sum b_i$。任一坐标为 0 时，数组及中间阶段的值均为 0。}
\textbf{差分与直接还原}
\begin{align*}
d(x)&=\sum_{b\in\{0,1\}^D}(-1)^{|b|}a(x-b),\\
a(x)&=d(x)+\sum_{b\ne0}(-1)^{|b|+1}a(x-b).
\end{align*}
\textbf{逐维递推：$O(DS)$，$S$ 为格数}
\begin{align*}
G_0(x)&=a(x),\\
G_i(x)&=G_{i-1}(x)-G_{i-1}(x-e_i),\quad G_D=d,\\
F_0(x)&=d(x),\\
F_i(x)&=F_{i-1}(x)+F_i(x-e_i),\quad F_D=a.
\end{align*}
\note{原地建表沿当前维倒序，还原正序。直接容斥还原为 $O(2^D S)$。}
\textbf{闭区域加 $v$：$2^D$ 个角点}
\[
p_i=\begin{cases}l_i,&b_i=0,\\r_i+1,&b_i=1,\end{cases}
\qquad d(p)\mathrel{+}=(-1)^{|b|}v.
\]
\note{各维预留第 0 层及右端点加 1。全零初始可直接修改；结果只还原一次。
下方三维代码是逐维递推的具体写法。}

\note{n、m、k是三维长度，d是差分表，坐标从1起；每维预留0和长度+\allowbreak{}1。区间加写8个角点，逐维累加后d变成实际值；D维公式中的S是选取坐标轴的集合，并非数组下标。}
\begin{lstlisting}
const int N=105;
int n,m,k;
ll d[N][N][N]; // 1-indexed；各维预留 0 和右端点+1，按题目调整大小
\end{lstlisting}
\note{非零初始值直接填 d，再建差分；全零起始清 d 后跳过 build。 建差分沿每一维倒序做相邻相减，避免覆盖尚未使用的原值。}
\begin{lstlisting}
void build()
{
    for(int x=n;x>=1;x--)for(int y=1;y<=m;y++)for(int z=1;z<=k;z++)
        d[x][y][z]-=d[x-1][y][z];
    for(int x=1;x<=n;x++)for(int y=m;y>=1;y--)for(int z=1;z<=k;z++)
        d[x][y][z]-=d[x][y-1][z];
    for(int x=1;x<=n;x++)for(int y=1;y<=m;y++)for(int z=k;z>=1;z--)
        d[x][y][z]-=d[x][y][z-1];
}
\end{lstlisting}
\note{闭长方体 [\allowbreak{}x1,\allowbreak{}x2]\allowbreak{} \ensuremath{\times} [\allowbreak{}y1,\allowbreak{}y2]\allowbreak{} \ensuremath{\times} [\allowbreak{}z1,\allowbreak{}z2]\allowbreak{} 加 v，修改八个角点。 每选一个右端点+\allowbreak{}1，符号翻转一次：0/\allowbreak{}2 个取正，1/\allowbreak{}3 个取负。}
\begin{lstlisting}
void add(int x1,int y1,int z1,int x2,int y2,int z2,ll v)
{
    ++x2,++y2,++z2;
    d[x1][y1][z1]+=v;
    d[x2][y1][z1]-=v; d[x1][y2][z1]-=v; d[x1][y1][z2]-=v;
    d[x2][y2][z1]+=v; d[x2][y1][z2]+=v; d[x1][y2][z2]+=v;
    d[x2][y2][z2]-=v;
}
\end{lstlisting}
\note{修改结束后，沿每一维正序做前缀和；结果仍在 d，只还原一次。}
\begin{lstlisting}
void restore()
{
    for(int x=1;x<=n;x++)for(int y=1;y<=m;y++)for(int z=1;z<=k;z++)
        d[x][y][z]+=d[x-1][y][z];
    for(int x=1;x<=n;x++)for(int y=1;y<=m;y++)for(int z=1;z<=k;z++)
        d[x][y][z]+=d[x][y-1][z];
    for(int x=1;x<=n;x++)for(int y=1;y<=m;y++)for(int z=1;z<=k;z++)
        d[x][y][z]+=d[x][y][z-1];
}
\end{lstlisting}
\note{三维建表/\allowbreak{}还原 O(\allowbreak{}nmk)\allowbreak{}，修改 O(\allowbreak{}1)\allowbreak{}；D 维建表/\allowbreak{}还原 O(\allowbreak{}D*S)\allowbreak{}，S 为格数。 空间随各维长度相乘；这里静态数组约 8.\allowbreak{}8 MiB，须保证 n+\allowbreak{}1,\allowbreak{}m+\allowbreak{}1,\allowbreak{}k+\allowbreak{}1<\allowbreak{}N。}
% SOURCE: 01-基础与技巧/三分法.cpp
\subsection{三分法}

\note{f/\allowbreak{}g是题目目标函数，l/\allowbreak{}r是搜索边界，返回最优位置而非函数值。实数与整数各留一个入口；调用前证明单峰方向，整数版最后枚举剩余位置。}
\note{求最小值：最优点/\allowbreak{}最优区间左侧严格下降，右侧严格上升；求最大值则反过来。 多个局部极值、非最优处的平台不能直接套；单调函数也可用，最优点在端点。 均返回极值位置；函数值用 f/\allowbreak{}g 再算。若能直接求导或推出最优点，优先直接求。}
\begin{lstlisting}
double f(double x);
double ternary_min(double l,double r)
{
    for(int i=1;i<=100;i++)
    {
        double x=l+(r-l)/3,y=r-(r-l)/3;
        if(f(x)<f(y))r=y;
        else l=x;
    }
    return (l+r)/2;
}
\end{lstlisting}
\note{100 轮后区间约为原来的 (\allowbreak{}2/\allowbreak{}3)\allowbreak{}\textasciicircum{}100，按初始范围及要求调整次数。 每轮调用两次 f；函数很贵时可用黄金分割复用一个点的函数值。}
\begin{lstlisting}
ll g(ll x); // 整数自变量的目标函数；值类型按题意改，r-l 须不溢出。
ll ternary_int(ll l,ll r)
{
    while(r-l>3)
    {
        ll x=l+(r-l)/3,y=r-(r-l)/3;
        if(g(x)<g(y))r=y;
        else l=x;
    }
    ll pos=l;
    for(ll i=l;i<r;i++)if(g(i+1)<g(pos))pos=i+1;
    return pos;
}
\end{lstlisting}
\note{整数区间最后剩至多 4 个点，必须枚举；不能照搬固定轮数的实数版。 求最大值：实数版的 <\allowbreak{} 改成 >\allowbreak{}；整数版两处 <\allowbreak{} 都改成 >\allowbreak{}（循环条件不改）。}
% SOURCE: 01-基础与技巧/分数规划.cpp
\subsection{分数规划}

\note{n为物品数，k为必须选取件数；a[\allowbreak{}i]\allowbreak{}是收益，b[\allowbreak{}i]\allowbreak{}>\allowbreak{}0是成本，求选出k件后的总收益/\allowbreak{}总成本。c[\allowbreak{}i]\allowbreak{}=\allowbreak{}a[\allowbreak{}i]\allowbreak{}-\allowbreak{}x*b[\allowbreak{}i]\allowbreak{}是判定临时权值，排序会修改c，a、b保持不变；下标1.\allowbreak{}.\allowbreak{}n。}
\note{n 件物品恰选 k 件，使总分子与总分母的比最大。a 可表示价值，b 表示重量，要求每个 b 为正。这是总和之比，不能按各物品比值直接取最大的 k 件。}\[\frac{\sum a_i}{\sum b_i}\ge x\quad\Longleftrightarrow\quad\sum(a_i-xb_i)\ge0.\]\note{猜比值 x 后，取最大的 k 个新权值 a-xb：和非负表示有方案达到 x。二分范围取单项比值的最小值与最大值；100 轮排序判定为 $O(100n\log n)$。}
\begin{lstlisting}
const int N=100005;
int n,k; // 物品数、必须选取的件数，1<=k<=n；数组下标从 1 开始。
double a[N],b[N],c[N]; // a：分子属性（收益）；b：分母属性（成本，>0）；c：判定时的新权值。
bool check(double x)
{
    for(int i=1;i<=n;i++)c[i]=a[i]-x*b[i];
    sort(c+1,c+n+1,greater<double>());
    double sum=0;
    for(int i=1;i<=k;i++)sum+=c[i];
    return sum>=0;
}
double fractional()
{
    double l=a[1]/b[1],r=l;
    for(int i=2;i<=n;i++)l=min(l,a[i]/b[i]),r=max(r,a[i]/b[i]);
    for(int i=1;i<=100;i++)
    {
        double mid=(l+r)/2;
        if(check(mid))l=mid;
        else r=mid;
    }
    return l;
}
\end{lstlisting}
\note{其他约束只改 check 内的选取方式，不可直接套“取最大的 k 项”。}
% SOURCE: 01-基础与技巧/高精度BigInt.cpp
\subsection{高精度BigInt}

\note{BigInt对象保存符号和按低位在前排列的数位；运算参数表示两个独立大整数，保留对象传递。按公共表示、比较、加减、乘除分别摘取，结果规范化去掉高位零，零的符号统一。}
\note{按区复制：非负加法 0+\allowbreak{}1，用 add\_abs；普通 a+\allowbreak{}b 复制 0+\allowbreak{}1+\allowbreak{}2+\allowbreak{}3。 仅乘法：0+\allowbreak{}4 或 0+\allowbreak{}5；除法/\allowbreak{}取模：0+\allowbreak{}2+\allowbreak{}4+\allowbreak{}6；比较：0+\allowbreak{}7。 各区放在同一个 struct BigInt 内，不需要的区整块省略。}
\topic{0 公共：存储、构造、去零、绝对值比较、输出}
\begin{lstlisting}
struct BigInt
{
    vector<int> d;
    int len,sign; // d[0] 为最低位，十进制；sign=1 正，-1 负，零统一为正
    BigInt():d(1,0),len(1),sign(1){}
    BigInt(ll v)
    {
        len=0,sign=1;
        unsigned long long u=v;
        if(v<0)sign=-1,u=0-u;
        while(u>0)
        {
            d.push_back((int)(u%10));
            u/=10;
            len++;
        }
        if(len==0)
        {
            len=1;
            d.push_back(0);// v=0 时存一个 0
        }
    }
    BigInt(const string& s)
    {
        len=0,sign=1;
        if(s.empty())
        {
            len=1,d.push_back(0);
            return;
        }
        int st=0;
        if(s[0]=='-')sign=-1,st=1;
        for(int i=(int)s.size()-1;i>=st;i--)d.push_back(s[i]-'0');
        len=(int)d.size();
        if(len==0)len=1,d.push_back(0);
        trim();
    }
    void trim()// 去掉高位 0，保证 0 的符号是正
    {
        while(len>1&&d[len-1]==0)len--;
        if(len<=0)// 空对象归一成「一个 0」
        {
            len=1;
            d.assign(1,0);
        }
        if(len==1&&d[0]==0)sign=1;
    }
    bool is_zero()const{return len<=0||(len==1&&d[0]==0);}
    bool abs_less(const BigInt& b)const// 比绝对值
    {
        if(len!=b.len)return len<b.len;
        for(int i=len-1;i>=0;i--)
            if(d[i]!=b.d[i])return d[i]<b.d[i];
        return false;
    }
    string to_string()const// 输出用
    {
        string s;
        if(len<=0)return "0";// 空对象当 0
        if(sign==-1&&!is_zero())s+='-';
        for(int i=len-1;i>=0;i--)s+=(char)('0'+d[i]);
        return s;
    }
\end{lstlisting}
\topic{1 绝对值加法：仅需 0，O(\allowbreak{}n)\allowbreak{}}
\begin{lstlisting}
    static BigInt add_abs(const BigInt& a,const BigInt& b)// 只用绝对值相加
    {
        BigInt res;
        res.d.assign(max(a.len,b.len)+1,0);
        int carry=0;
        for(int i=0;i<a.len||i<b.len||carry;i++)
        {
            int sum=carry+(i<a.len?a.d[i]:0)+(i<b.len?b.d[i]:0);
            res.d[i]=sum%10;
            carry=sum/10;
        }
        res.len=(int)res.d.size();
        res.trim();
        return res;
    }
\end{lstlisting}
\topic{2 绝对值减法：仅需 0，要求 |a|>\allowbreak{}=\allowbreak{}|b|，O(\allowbreak{}n)\allowbreak{}}
\begin{lstlisting}
    static BigInt sub_abs(const BigInt& a,const BigInt& b)// 只用绝对值，返回 |a|-|b|（要求 |a|>=|b|）
    {
        BigInt res;
        res.d.assign(max(a.len,b.len)+1,0);// 多一位防借位越界
        int borrow=0;
        for(int i=0;i<a.len||i<b.len||borrow;i++)
        {
            int cur=(i<a.len?a.d[i]:0)-borrow-(i<b.len?b.d[i]:0);
            if(cur<0)cur+=10,borrow=1;
            else borrow=0;
            res.d[i]=cur;
        }
        res.len=(int)res.d.size();
        res.trim();
        return res;
    }
\end{lstlisting}
\topic{3 带符号加减：需 0+\allowbreak{}1+\allowbreak{}2，O(\allowbreak{}n)\allowbreak{}}
\begin{lstlisting}
    BigInt operator-()const// 取相反数
    {
        BigInt res=*this;
        if(!res.is_zero())res.sign=-res.sign;
        return res;
    }
    BigInt operator+(const BigInt& b)const
    {
        BigInt res;
        if(sign==b.sign)
        {
            res=add_abs(*this,b);
            if(!res.is_zero())res.sign=sign;
        }
        else
        {
            bool less=abs_less(b);
            res=less?sub_abs(b,*this):sub_abs(*this,b);
            if(!res.is_zero())res.sign=less?b.sign:sign;
        }
        return res;
    }
    BigInt operator-(const BigInt& b)const
    {
        BigInt t=b;// a-b = a+(-b)，符号分情况走加法
        if(!t.is_zero())t.sign=-b.sign;
        return *this+t;
    }
\end{lstlisting}
\topic{4 高精乘 int：仅需 0，O(\allowbreak{}n)\allowbreak{}}
\begin{lstlisting}
    BigInt operator*(int v)const// 高精乘低精，除法试商时用
    {
        BigInt res;
        if(v==0)return res;// 已经是 0
        res.d.clear();
        res.sign=sign;
        ll u=v;
        if(u<0)res.sign=-res.sign,u=-u;
        ll carry=0;
        int i=0;
        for(;i<len||carry;i++)
        {
            ll cur=carry+(i<len?(ll)d[i]*u:0);
            res.d.push_back((int)(cur%10));
            carry=cur/10;
        }
        res.len=i>0?i:1;
        res.trim();
        return res;
    }
\end{lstlisting}
\topic{5 高精乘高精：仅需 0，O(\allowbreak{}nm)\allowbreak{}}
\begin{lstlisting}
    BigInt operator*(const BigInt& b)const// 高精乘高精
    {
        BigInt res;
        if(is_zero()||b.is_zero())return res;
        int n=len+b.len+1;// 多一位放最高进位
        res.d.assign(n,0);
        for(int i=0;i<len;i++)
        {
            int carry=0;
            for(int j=0;j<b.len||carry;j++)
            {
                int cur=res.d[i+j]+carry+(j<b.len?d[i]*b.d[j]:0);
                res.d[i+j]=cur%10;
                carry=cur/10;
            }
        }
        res.len=n;
        res.sign=sign*b.sign;
        res.trim();
        return res;
    }
\end{lstlisting}
\topic{6 除法与取模：需 0+\allowbreak{}2+\allowbreak{}4，无需加法区和高精乘高精区，O(\allowbreak{}nm)\allowbreak{}}
\begin{lstlisting}
    static pair<BigInt,BigInt> divmod_abs(const BigInt& a,const BigInt& b)// 只用绝对值，返回 {商,余数}
    {
        if(a.abs_less(b))return make_pair(BigInt(0),a);
        BigInt q(0),r(0);
        q.d.assign(a.len,0);
        q.len=a.len;
        for(int i=a.len-1;i>=0;i--)// 逐位试商，每位二分 0..9
        {
            r.d.insert(r.d.begin(),a.d[i]); // r=r*10+当前位，无需加法区
            r.len=(int)r.d.size();
            r.trim();
            int lo=0,hi=9,dig=0;
            while(lo<=hi)
            {
                int mid=(lo+hi)>>1;
                if(!r.abs_less(b*mid))dig=mid,lo=mid+1;
                else hi=mid-1;
            }
            q.d[i]=dig;
            r=sub_abs(r,b*dig);
        }
        q.trim();
        r.trim();
        return make_pair(q,r);
    }
    static pair<BigInt,BigInt> divmod(const BigInt& a,const BigInt& b)// 返回 {商,余数}，符号同 C++
    {
        assert(!b.is_zero()); // 除数必须非零
        BigInt x=a,y=b;
        int sa=x.sign,sb=y.sign;
        x.sign=1,y.sign=1;
        pair<BigInt,BigInt> pr=divmod_abs(x,y);
        BigInt q=pr.first,r=pr.second;
        if(!q.is_zero())q.sign=sa*sb;
        if(!r.is_zero())r.sign=sa;// 余数跟被除数同号
        return make_pair(q,r);
    }
    BigInt operator/(const BigInt& b)const{return divmod(*this,b).first;}
    BigInt operator%(const BigInt& b)const{return divmod(*this,b).second;}
\end{lstlisting}
\topic{7 比较重载：仅需 0，按需选用}
\begin{lstlisting}
    bool operator<(const BigInt& b)const
    {
        if(sign!=b.sign)return sign<b.sign;
        if(sign==1)return abs_less(b);
        return b.abs_less(*this);
    }
    bool operator>(const BigInt& b)const{return b<*this;}
    bool operator==(const BigInt& b)const
    {
        if(len!=b.len||sign!=b.sign)return false;
        for(int i=0;i<len;i++)
            if(d[i]!=b.d[i])return false;
        return true;
    }
    bool operator!=(const BigInt& b)const{return !(*this==b);}
    bool operator<=(const BigInt& b)const{return !(b<*this);}
    bool operator>=(const BigInt& b)const{return !(*this<b);}
};
\end{lstlisting}
\section{数据结构}
% SOURCE: 02-数据结构/树状数组.cpp
\subsection{树状数组}

\note{BIT对象的tr存区间累计量，n是有效长度；单点更新与前缀查询从1起，不能在0上执行add。a/\allowbreak{}tmp和bit\_inv用于逆序对应用；多测清空树状数组，差分及二维模块另按正文说明初始化。}
\note{一维、二维点加求和，差分区间加，以及频数顺序统计与逆序对计数。}
\topic{公共容量}
\note{下标从1开始；容量模板参数包含第0格，须满足 $1\le n<NMAX$。二维同样要求 $1\le n,m<NMAX$。默认N用于一维、M用于二维；按题目容量调整，不必同时摘取全部模块。}
\begin{lstlisting}
const int N=100005;
const int M=1005;//二维树状数组的默认边长
\end{lstlisting}
\topic{单点加、区间和与第k个}
\note{BIT：先init(\allowbreak{}n)\allowbreak{}，add(\allowbreak{}x,\allowbreak{}v)\allowbreak{}加到单点，sum(\allowbreak{}x)\allowbreak{}求[\allowbreak{}1,\allowbreak{}x]\allowbreak{}，query(\allowbreak{}l,\allowbreak{}r)\allowbreak{}求闭区间和。单次操作 $O(\log n)$、初始化 $O(n)$，空间 $O(n)$。add的x必须至少1，否则lowbit为0；区间查询允许l>\allowbreak{}r返回0，其他下标须在范围内。T须容纳各前缀和。}
\note{kth(\allowbreak{}k)\allowbreak{}返回首个前缀和至少k的位置，不存在返回n+\allowbreak{}1，要求k>\allowbreak{}0且每个位置的当前值非负，使前缀和单调。普通区间和可以存负数；用了负数后不能继续套kth。频数树上的第k小直接用这一接口。}
\begin{lstlisting}
template<typename T,int NMAX=N>
struct BIT
{
    T tr[NMAX];
    int n;
    void init(int n_)//下标 1..n_
    {
        n=n_;
        for(int i=1;i<=n;i++)tr[i]=0;
    }
    void add(int x,T v)
    {
        for(;x<=n;x+=x&-x)tr[x]+=v;
    }
    T sum(int x)//前缀和 [1,x]
    {
        T s=0;
        for(;x>0;x-=x&-x)s+=tr[x];
        return s;
    }
    T query(int l,int r)//区间和 [l,r]
    {
        if(l>r)return 0;
        return sum(r)-sum(l-1);
    }
    // 前缀二分；使用条件见对应说明。
    int kth(T k)
    {
        int p=0,lg=1;
        T s=0;
        while((lg<<1)<=n)lg<<=1;//不超过 n 的最大 2 的幂
        for(int j=lg;j;j>>=1)
            if(p+j<=n&&s+tr[p+j]<k)s+=tr[p+j],p+=j;
        return p+1;
    }
};
\end{lstlisting}
\topic{区间加、单点查}
\note{DiffBIT依赖BIT。以差分存储，update(\allowbreak{}l,\allowbreak{}r,\allowbreak{}v)\allowbreak{}修改两端点，query(\allowbreak{}x)\allowbreak{}还原单点；两者均 $O(\log n)$。从全零开始，r=\allowbreak{}n时不修改n+\allowbreak{}1。它不提供区间和接口；需要区间加与区间和时采用两棵树状数组或线段树。}
\begin{lstlisting}
template<typename T,int NMAX=N>
struct DiffBIT
{
    BIT<T,NMAX> c;
    int n;
    void init(int n_)
    {
        n=n_,c.init(n_);
    }
    void update(int l,int r,T v)
    {
        c.add(l,v);
        if(r+1<=n)c.add(r+1,-v);//边界：r=n 时右端点不越界
    }
    T query(int x)
    {
        return c.sum(x);
    }
};
\end{lstlisting}
\topic{二维单点加、子矩阵和}
\note{BIT2D：add(\allowbreak{}x,\allowbreak{}y,\allowbreak{}v)\allowbreak{}单点加，sum(\allowbreak{}x,\allowbreak{}y)\allowbreak{}求左上前缀，query(\allowbreak{}x1,\allowbreak{}y1,\allowbreak{}x2,\allowbreak{}y2)\allowbreak{}用四项容斥求闭子矩阵。单次 $O(\log n\log m)$，初始化和空间 $O(nm)$；两个更新坐标都必须大于0。大网格先检查平方容量与内存。}
\begin{lstlisting}
template<typename T,int NMAX=M>
struct BIT2D
{
    T tr[NMAX][NMAX];
    int n,m;
    void init(int n_,int m_)
    {
        n=n_,m=m_;
        for(int i=1;i<=n;i++)
            for(int j=1;j<=m;j++)tr[i][j]=0;
    }
    void add(int x,int y,T v)
    {
        for(int i=x;i<=n;i+=i&-i)
            for(int j=y;j<=m;j+=j&-j)tr[i][j]+=v;
    }
    T sum(int x,int y)//左上角前缀和
    {
        T s=0;
        for(int i=x;i>0;i-=i&-i)
            for(int j=y;j>0;j-=j&-j)s+=tr[i][j];
        return s;
    }
    T query(int x1,int y1,int x2,int y2)//子矩阵和
    {
        return sum(x2,y2)-sum(x1-1,y2)-sum(x2,y1-1)+sum(x1-1,y1-1);
    }
};
\end{lstlisting}
\topic{离散化计逆序对}
\note{依赖BIT。把输入放在a[\allowbreak{}1.\allowbreak{}.\allowbreak{}len]\allowbreak{}，调用calc\_inv(\allowbreak{}len)\allowbreak{}，返回严格逆序对数：只计 $i<j$ 且 $a_i>a_j$，相等值不计。先压缩值域，再从左到右查询之前大于当前值的数量；时间 $O(n\log n)$，额外空间 $O(n)$，答案为ll。tmp和bit\_inv会被重置，原数组a不变。}
\begin{lstlisting}
int n;
int a[N],tmp[N];
BIT<int> bit_inv;
ll calc_inv(int len)
{
    for(int i=1;i<=len;i++)tmp[i]=a[i];
    sort(tmp+1,tmp+len+1);
    int tot=unique(tmp+1,tmp+len+1)-tmp-1;
    bit_inv.init(tot);
    ll ans=0;
    for(int i=1;i<=len;i++)
    {
        int rk=lower_bound(tmp+1,tmp+tot+1,a[i])-tmp;
        ans+=i-1-bit_inv.sum(rk);//前面比 a[i] 严格大的个数
        bit_inv.add(rk,1);
    }
    return ans;
}
\end{lstlisting}
% SOURCE: 02-数据结构/线段树(区间加区间和).cpp
\subsection{线段树(区间加区间和)}

\note{a[\allowbreak{}1.\allowbreak{}.\allowbreak{}n]\allowbreak{}为初值，seg保存区间和与待下传加量；br仅供对照。修改区间闭合，sum增加长度乘增量；每次换数据重新build，不能只清根懒标记。}
\begin{lstlisting}
const int N=100005;
\end{lstlisting}
\note{线段树：区间加 +\allowbreak{} 区间和，懒标记下推，每次操作 O(\allowbreak{}log n)\allowbreak{}}
\begin{lstlisting}
template<typename T,int NMAX=N>
struct SegmentTree
{
    T tr[NMAX*4],lazy[NMAX*4];
    T *src;//建树用的原数组(1-indexed)
    #define lp (p<<1)
    #define rp ((p<<1)|1)
    #define mid ((l+r)>>1)
    void build(int l,int r,int p)
    {
        lazy[p]=0;
        if(l==r)
        {
            tr[p]=src[l];
            return;
        }
        build(l,mid,lp);
        build(mid+1,r,rp);
        tr[p]=tr[lp]+tr[rp];
    }
    void push_down(int p,int l,int r)//把 p 的标记发给两个儿子
    {
        if(lazy[p]==0)return;
        lazy[lp]+=lazy[p],tr[lp]+=lazy[p]*(mid-l+1);
        lazy[rp]+=lazy[p],tr[rp]+=lazy[p]*(r-mid);
        lazy[p]=0;
    }
    void update(int L,int R,T v,int l,int r,int p)
    {
        if(L<=l&&r<=R)//(l,r)<=(L,R)
        {
            lazy[p]+=v,tr[p]+=v*(r-l+1);
            return;
        }
        push_down(p,l,r);
        if(L<=mid)update(L,R,v,l,mid,lp);
        if(R>mid)update(L,R,v,mid+1,r,rp);
        tr[p]=tr[lp]+tr[rp];
    }
    T query(int L,int R,int l,int r,int p)
    {
        if(L<=l&&r<=R)return tr[p];
        push_down(p,l,r);
        T res=0;
        if(L<=mid)res+=query(L,R,l,mid,lp);
        if(R>mid)res+=query(L,R,mid+1,r,rp);
        return res;
    }
};
int n;
ll a[205],br[205];
SegmentTree<ll,205> seg;
\end{lstlisting}
% SOURCE: 02-数据结构/线段树(区间加乘).cpp
\subsection{线段树(区间加乘)}

\note{a[\allowbreak{}1.\allowbreak{}.\allowbreak{}n]\allowbreak{}是输入，seg保存区间和、乘法与加法懒标记；br仅用于朴素对照。乘法标记初值1，加法初值0；复合新变换时同时更新已有加法，不能交换执行顺序。}
\begin{lstlisting}
const int N=100005;
\end{lstlisting}
\note{线段树：区间乘 +\allowbreak{} 区间加 +\allowbreak{} 区间和，每次操作 O(\allowbreak{}log n)\allowbreak{} T 为不超过 64 位的整数，mod>\allowbreak{}=\allowbreak{}0；不取模时结果须能放入 T。 标记约定：先乘后加，儿子值 =\allowbreak{} 儿子值*mul +\allowbreak{} add*区间长度}
\begin{lstlisting}
template<typename T,int NMAX=N>
struct SegmentTree
{
    T tr[NMAX*4],add[NMAX*4],mul[NMAX*4],mod;
    T *src;//建树用的原数组(1-indexed)
    #define lp (p<<1)
    #define rp ((p<<1)|1)
    #define mid ((l+r)>>1)
    T mo(__int128 x){if(!mod)return (T)x; x%=mod; return (T)(x<0?x+mod:x);}
    void build(int l,int r,int p)
    {
        add[p]=0,mul[p]=1;//加标记 0，乘标记 1
        if(l==r)
        {
            tr[p]=mo(src[l]);
            return;
        }
        build(l,mid,lp);
        build(mid+1,r,rp);
        tr[p]=mo((__int128)tr[lp]+tr[rp]);
    }
    void push_down(int p,int l,int r)
    {
        if(mul[p]==1&&add[p]==0)return;
        tr[lp]=mo((__int128)tr[lp]*mul[p]+(__int128)add[p]*(T)(mid-l+1));
        mul[lp]=mo((__int128)mul[lp]*mul[p]);
        add[lp]=mo((__int128)add[lp]*mul[p]+add[p]);
        tr[rp]=mo((__int128)tr[rp]*mul[p]+(__int128)add[p]*(T)(r-mid));
        mul[rp]=mo((__int128)mul[rp]*mul[p]);
        add[rp]=mo((__int128)add[rp]*mul[p]+add[p]);
        mul[p]=1,add[p]=0;
    }
    void update_add(int L,int R,T v,int l,int r,int p)
    {
        if(L<=l&&r<=R)//(l,r)<=(L,R)
        {
            tr[p]=mo((__int128)tr[p]+(__int128)v*(T)(r-l+1));
            add[p]=mo((__int128)add[p]+v);
            return;
        }
        push_down(p,l,r);
        if(L<=mid)update_add(L,R,v,l,mid,lp);
        if(R>mid)update_add(L,R,v,mid+1,r,rp);
        tr[p]=mo((__int128)tr[lp]+tr[rp]);
    }
    void update_mul(int L,int R,T v,int l,int r,int p)
    {
        if(L<=l&&r<=R)
        {
            tr[p]=mo((__int128)tr[p]*v);
            mul[p]=mo((__int128)mul[p]*v);
            add[p]=mo((__int128)add[p]*v);//加法标记也要乘
            return;
        }
        push_down(p,l,r);
        if(L<=mid)update_mul(L,R,v,l,mid,lp);
        if(R>mid)update_mul(L,R,v,mid+1,r,rp);
        tr[p]=mo((__int128)tr[lp]+tr[rp]);
    }
    T query(int L,int R,int l,int r,int p)
    {
        if(L<=l&&r<=R)return tr[p];
        push_down(p,l,r);
        T res=0;
        if(L<=mid)res=mo((__int128)res+query(L,R,l,mid,lp));
        if(R>mid)res=mo((__int128)res+query(L,R,mid+1,r,rp));
        return res;
    }
};
int n;
ll a[205],br[205];
SegmentTree<ll,205> seg;
\end{lstlisting}
% SOURCE: 02-数据结构/并查集.cpp
\subsection{并查集}

\note{d保存父亲与集合大小，n为元素数，下标1.\allowbreak{}.\allowbreak{}n；init(\allowbreak{}n)\allowbreak{}使每点独立。find返回代表元，merge合并集合并维护大小；路径压缩会改父亲，需回滚时改用回滚并查集。}
\begin{lstlisting}
const int N=100005;
\end{lstlisting}
\note{并查集：路径压缩 +\allowbreak{} 按大小合并，单次操作近似 O(\allowbreak{}alpha(\allowbreak{}n)\allowbreak{})\allowbreak{}}
\begin{lstlisting}
template<int NMAX=N>
struct DSU
{
    int fa[NMAX],sz[NMAX];
    void init(int n)//1..n 各自成一个集合
    {
        for(int i=1;i<=n;i++)fa[i]=i,sz[i]=1;
    }
    int find(int x)
    {
        return fa[x]==x?x:fa[x]=find(fa[x]);
    }
    bool same(int x,int y)
    {
        return find(x)==find(y);
    }
    void unionn(int x,int y)
    {
        int fx=find(x),fy=find(y);
        if(fx==fy)return;
        if(sz[fx]<sz[fy])swap(fx,fy);//小的树挂到大树下面
        fa[fy]=fx,sz[fx]+=sz[fy];
    }
    int size(int x)//x 所在集合大小
    {
        return sz[find(x)];
    }
};
int n;
DSU<305> d;
\end{lstlisting}
% SOURCE: 02-数据结构/带权并查集.cpp
\subsection{带权并查集}

\note{父亲数组记录代表元，权值记录节点与父亲的势能差，find压缩时累加差值。合并参数是两点及它们的关系；若已同根先核对关系，换实例初始化父亲、大小和零差。}
\begin{lstlisting}
const int N=100005;
const int INF=0x3f3f3f3f;
\end{lstlisting}
\note{种类并查集(\allowbreak{}3n 拆点)\allowbreak{}：x 同类域 /\allowbreak{} x+\allowbreak{}n x 的猎物 /\allowbreak{} x+\allowbreak{}n*2 x 的天敌，n 个点要开 3n 单次操作近似 O(\allowbreak{}alpha(\allowbreak{}n)\allowbreak{})\allowbreak{}，只能表示"同类/\allowbreak{}吃"这种循环关系}
\begin{lstlisting}
template<int NMAX=N>
struct KindDSU
{
    int fa[NMAX*3],n;
    void init(int n_)
    {
        n=n_;
        for(int i=1;i<=n*3;i++)fa[i]=i;
    }
    int find(int x)
    {
        return fa[x]==x?x:fa[x]=find(fa[x]);
    }
    void unionn(int x,int y)
    {
        fa[find(x)]=find(y);
    }
    // 声明 x,y 同类，返回 false 表示与已有关系矛盾
    bool same(int x,int y)
    {
        if(find(x)==find(y+n)||find(x)==find(y+2*n))return false;//已经互相吃
        unionn(x,y),unionn(x+n,y+n),unionn(x+2*n,y+2*n);
        return true;
    }
    // 声明 x 吃 y，返回 false 表示矛盾
    bool eat(int x,int y)
    {
        if(find(x)==find(y)||find(x)==find(y+n))return false;//同类或已被 y 吃
        unionn(x+n,y),unionn(x,y+2*n),unionn(x+2*n,y+n);
        return true;
    }
};
\end{lstlisting}
\note{边权并查集：d[\allowbreak{}x]\allowbreak{} =\allowbreak{} val[\allowbreak{}x]\allowbreak{}-\allowbreak{}val[\allowbreak{}fa[\allowbreak{}x]\allowbreak{}]\allowbreak{}，find 时先递归到根再累加权值 merge(\allowbreak{}x,\allowbreak{}y,\allowbreak{}w)\allowbreak{} 表示 val[\allowbreak{}y]\allowbreak{}-\allowbreak{}val[\allowbreak{}x]\allowbreak{}=\allowbreak{}w；mod>\allowbreak{}0 时权值对 mod 取模(\allowbreak{}种类并查集用 mod=\allowbreak{}3，食物链里 2 表示 x 吃 y)\allowbreak{}}
\begin{lstlisting}
template<int NMAX=N>
struct WeightDSU
{
    int fa[NMAX],mod;
    ll d[NMAX];
    ll norm(ll x)
    {
        if(mod==0)return x;
        x%=mod;return x<0?x+mod:x;
    }
    void init(int n,int m=0)
    {
        mod=m;
        for(int i=1;i<=n;i++)fa[i]=i,d[i]=0;
    }
    int find(int x)
    {
        if(fa[x]==x)return x;
        int f=fa[x];
        fa[x]=find(f);
        d[x]=norm(d[x]+d[f]);//val[x]-val[根]
        return fa[x];
    }
    // 加入关系 val[y]-val[x]=w，返回 false 表示与已有关系矛盾
    bool merge(int x,int y,ll w)
    {
        int fx=find(x),fy=find(y);
        if(fx==fy)return norm(d[y]-d[x])==norm(w);
        fa[fy]=fx;
        d[fy]=norm(d[x]+w-d[y]);//解出 val[fy]-val[fx]
        return true;
    }
    // 同集合返回 val[y]-val[x]，否则返回 INF
    ll query(int x,int y)
    {
        if(find(x)!=find(y))return INF;
        return norm(d[y]-d[x]);
    }
};
\end{lstlisting}
% SOURCE: 02-数据结构/ST表.cpp
\subsection{ST表}

\note{a[\allowbreak{}1.\allowbreak{}.\allowbreak{}n]\allowbreak{}是静态输入，stx/\allowbreak{}stn分别为最大/\allowbreak{}最小值ST表；对象中的f[\allowbreak{}i]\allowbreak{}[\allowbreak{}k]\allowbreak{}覆盖从i开始的2\textasciicircum{}k项，lg为向下取整对数。build(\allowbreak{}n,\allowbreak{}a)\allowbreak{}复制输入到表，修改原数组不会更新查询；容量模板参数按实际规模设置。}
\begin{lstlisting}
const int N=100005;
\end{lstlisting}
\note{ST 表：静态区间最值(\allowbreak{}RMQ)\allowbreak{}，n log n 预处理，O(\allowbreak{}1)\allowbreak{} 查询，不支持修改 查询用两段长度为 2\textasciicircum{}k 的区间覆盖 [\allowbreak{}l,\allowbreak{}r]\allowbreak{}，可重叠}
\begin{lstlisting}
template<typename T,int NMAX=N,int LOG=18>
struct STmax
{
    T f[NMAX][LOG];
    int lg[NMAX],n;
    void build(int n_,T *src)
    {
        n=n_;
        lg[1]=0;
        for(int i=2;i<=n;i++)lg[i]=lg[i>>1]+1;//预处理 log2 向下取整
        for(int i=1;i<=n;i++)f[i][0]=src[i];
        for(int j=1;(1<<j)<=n;j++)
            for(int i=1;i+(1<<j)-1<=n;i++)
                f[i][j]=max(f[i][j-1],f[i+(1<<(j-1))][j-1]);
    }
    T query(int l,int r)
    {
        int k=lg[r-l+1];
        return max(f[l][k],f[r-(1<<k)+1][k]);
    }
};
template<typename T,int NMAX=N,int LOG=18>
struct STmin
{
    T f[NMAX][LOG];
    int lg[NMAX],n;
    void build(int n_,T *src)
    {
        n=n_;
        lg[1]=0;
        for(int i=2;i<=n;i++)lg[i]=lg[i>>1]+1;
        for(int i=1;i<=n;i++)f[i][0]=src[i];
        for(int j=1;(1<<j)<=n;j++)
            for(int i=1;i+(1<<j)-1<=n;i++)
                f[i][j]=min(f[i][j-1],f[i+(1<<(j-1))][j-1]);
    }
    T query(int l,int r)
    {
        int k=lg[r-l+1];
        return min(f[l][k],f[r-(1<<k)+1][k]);
    }
};
int n;
int a[105];
STmax<int,105,8> stx;
STmin<int,105,8> stn;
\end{lstlisting}
% SOURCE: 02-数据结构/树链剖分.cpp
\subsection{树链剖分}

\note{adj存无向树；val是原顶点编号下的点权，a为预留数组，建线段树时读取val[\allowbreak{}rnk[\allowbreak{}dfn]\allowbreak{}]\allowbreak{}。fa/\allowbreak{}dep/\allowbreak{}sz/\allowbreak{}son为父亲、深度、大小、重儿子，top是链顶，dfn/\allowbreak{}rnk互为映射；rt为根，stk/\allowbreak{}order供迭代预处理。换树清空adj，重新预处理并建线段树。}
\note{树链剖分（重链剖分）+\allowbreak{} 线段树：路径加/\allowbreak{}路径和 +\allowbreak{} 子树加/\allowbreak{}子树和 建树 O(\allowbreak{}n)\allowbreak{}，路径操作 O(\allowbreak{}log\textasciicircum{}2 n)\allowbreak{}，子树操作 O(\allowbreak{}log n)\allowbreak{} 关键性质：子树在 dfs 序上是一段连续区间 [\allowbreak{}dfn[\allowbreak{}x]\allowbreak{},\allowbreak{} dfn[\allowbreak{}x]\allowbreak{}+\allowbreak{}sz[\allowbreak{}x]\allowbreak{}-\allowbreak{}1]\allowbreak{} 两遍 dfs 都写成迭代版（显式栈），链状数据也不会爆栈}
\begin{lstlisting}
const int N=100005;
int n,m,rt,cnt,op;
int a[N],val[N];
vector<int> adj[N];
int fa[N],dep[N],sz[N],son[N],top[N],dfn[N],rnk[N];
int stk[N],stk2[N];                     // 迭代 dfs 用的栈
void add_edge(int u,int v)
{
    adj[u].push_back(v);
    adj[v].push_back(u);
}
\end{lstlisting}
\note{第一遍（迭代）：求 fa /\allowbreak{} dep /\allowbreak{} sz /\allowbreak{} 重儿子}
\begin{lstlisting}
int order[N],ocnt;
void dfs1_iter(int root)
{
    int tp=0;
    fa[root]=0,dep[root]=1;
    stk[++tp]=root,ocnt=0;
    while(tp)
    {
        int u=stk[tp--];
        order[++ocnt]=u;
        for(int v:adj[u])
        {
            if(v==fa[u])continue;
            fa[v]=u,dep[v]=dep[u]+1;
            stk[++tp]=v;
        }
    }
    for(int i=ocnt;i>=1;i--)            // 逆序：先算完儿子再算父亲
    {
        int u=order[i];
        sz[u]=1,son[u]=0;
        for(int v:adj[u])
        {
            if(v==fa[u])continue;
            sz[u]+=sz[v];
            if(sz[v]>sz[son[u]])son[u]=v;
        }
    }
}
\end{lstlisting}
\note{第二遍（迭代）：编 dfs 序，重儿子优先 → 同一条重链上 dfn 连续}
\begin{lstlisting}
void dfs2(int root,int tproot)
{
    cnt=0;
    int tp=0;
    stk[tp]=root,stk2[tp]=tproot,tp++;
    while(tp)
    {
        tp--;
        int u=stk[tp],tpv=stk2[tp];
        top[u]=tpv,dfn[u]=++cnt,rnk[cnt]=u;
        for(int v:adj[u])   // 轻儿子各自开新链
        {
            if(v==fa[u]||v==son[u])continue;
            stk[tp]=v,stk2[tp]=v,tp++;
        }
        if(son[u])stk[tp]=son[u],stk2[tp]=tpv,tp++;   // 重儿子接着这条链
    }
}
struct SegmentTree{
    ll tr[N<<2],lazy[N<<2];
    #define lp (p<<1)
    #define rp ((p<<1)|1)
    #define mid ((l+r)>>1)
    void build(int l,int r,int p)
    {
        lazy[p]=0;
        if(l==r){tr[p]=val[rnk[l]];return;}
        build(l,mid,lp),build(mid+1,r,rp);
        tr[p]=tr[lp]+tr[rp];
    }
    void push_down(int l,int r,int p)
    {
        if(!lazy[p])return;
        tr[lp]+=lazy[p]*(mid-l+1),lazy[lp]+=lazy[p];
        tr[rp]+=lazy[p]*(r-mid),lazy[rp]+=lazy[p];
        lazy[p]=0;
    }
    void update(int L,int R,ll k,int l,int r,int p)
    {
        if(L<=l&&r<=R){tr[p]+=k*(r-l+1),lazy[p]+=k;return;}
        push_down(l,r,p);
        if(L<=mid)update(L,R,k,l,mid,lp);
        if(R>mid)update(L,R,k,mid+1,r,rp);
        tr[p]=tr[lp]+tr[rp];
    }
    ll query(int L,int R,int l,int r,int p)
    {
        if(L<=l&&r<=R)return tr[p];
        push_down(l,r,p);
        ll res=0;
        if(L<=mid)res+=query(L,R,l,mid,lp);
        if(R>mid)res+=query(L,R,mid+1,r,rp);
        return res;
    }
}seg;
\end{lstlisting}
\note{路径加：u,\allowbreak{}v 往上跳到同一条重链}
\begin{lstlisting}
void path_add(int u,int v,ll k)
{
    while(top[u]!=top[v])
    {
        if(dep[top[u]]<dep[top[v]])swap(u,v);
        seg.update(dfn[top[u]],dfn[u],k,1,n,1);
        u=fa[top[u]];
    }
    if(dep[u]>dep[v])swap(u,v);
    seg.update(dfn[u],dfn[v],k,1,n,1);
}
ll path_sum(int u,int v)
{
    ll res=0;
    while(top[u]!=top[v])
    {
        if(dep[top[u]]<dep[top[v]])swap(u,v);
        res+=seg.query(dfn[top[u]],dfn[u],1,n,1);
        u=fa[top[u]];
    }
    if(dep[u]>dep[v])swap(u,v);
    res+=seg.query(dfn[u],dfn[v],1,n,1);
    return res;
}
void subtree_add(int u,ll k){seg.update(dfn[u],dfn[u]+sz[u]-1,k,1,n,1);}
ll subtree_sum(int u){return seg.query(dfn[u],dfn[u]+sz[u]-1,1,n,1);}
\end{lstlisting}
% SOURCE: 02-数据结构/线段树(最大子段和).cpp
\subsection{线段树(最大子段和)}

\note{a[\allowbreak{}1.\allowbreak{}.\allowbreak{}n]\allowbreak{}为序列；seg节点保存sum、最大前缀、最大后缀和最大子段和。合并按左到右顺序，查询结果不能交换左右；br为对照缓冲，空段是否允许以本模板的叶子初始化为准。}
\begin{lstlisting}
const int N=500005;
\end{lstlisting}
\note{线段树：单点修改 +\allowbreak{} 区间查询最大子段和(\allowbreak{}不能为空)\allowbreak{} 每个结点维护 sum /\allowbreak{} lmax(\allowbreak{}左端起)\allowbreak{} /\allowbreak{} rmax(\allowbreak{}右端起)\allowbreak{} /\allowbreak{} mx，merge O(\allowbreak{}1)\allowbreak{}，每次操作 O(\allowbreak{}log n)\allowbreak{}}
\begin{lstlisting}
template<typename T,int NMAX=N>
struct SegmentTree
{
    struct Node
    {
        T sum,lmax,rmax,mx;
        Node(){}
        Node(T x){sum=lmax=rmax=mx=x;}
    }tr[NMAX*4];
    T *src;//建树用的原数组(1-indexed)
    #define lp (p<<1)
    #define rp ((p<<1)|1)
    #define mid ((l+r)>>1)
    Node merge(Node a,Node b)
    {
        Node c;
        c.sum=a.sum+b.sum;
        c.lmax=max(a.lmax,a.sum+b.lmax);
        c.rmax=max(b.rmax,b.sum+a.rmax);
        c.mx=max(max(a.mx,b.mx),a.rmax+b.lmax);//跨过中点的那一段
        return c;
    }
    void build(int l,int r,int p)
    {
        if(l==r)
        {
            tr[p]=Node(src[l]);
            return;
        }
        build(l,mid,lp);
        build(mid+1,r,rp);
        tr[p]=merge(tr[lp],tr[rp]);
    }
    void update(int L,T v,int l,int r,int p)//把 a[L] 改成 v
    {
        if(l==r)
        {
            tr[p]=Node(v);
            return;
        }
        if(L<=mid)update(L,v,l,mid,lp);
        else update(L,v,mid+1,r,rp);
        tr[p]=merge(tr[lp],tr[rp]);
    }
    Node query(int L,int R,int l,int r,int p)
    {
        if(L<=l&&r<=R)return tr[p];
        if(R<=mid)return query(L,R,l,mid,lp);
        if(L>mid)return query(L,R,mid+1,r,rp);
        return merge(query(L,R,l,mid,lp),query(L,R,mid+1,r,rp));
    }
};
int n;
ll a[205],br[205];
SegmentTree<ll,205> seg;
\end{lstlisting}
% SOURCE: 02-数据结构/主席树(可持久化线段树).cpp
\subsection{主席树(可持久化线段树)}

\note{a[\allowbreak{}1.\allowbreak{}.\allowbreak{}n]\allowbreak{}是原数组，b为升序去重值，sz为值域大小；第i个根代表前i项的频数版本。节点存左右儿子与计数，区间第k小用两个前缀根之差；多测重置节点池和根，不修改旧版本。}
\note{主席树（可持久化线段树）：静态区间第 k 小 /\allowbreak{} 区间内某值域的数个数 建树 O(\allowbreak{}n log n)\allowbreak{}，单次询问 O(\allowbreak{}log n)\allowbreak{}，空间 O(\allowbreak{}n log n)\allowbreak{} 思路：对每个前缀 [\allowbreak{}1.\allowbreak{}.\allowbreak{}i]\allowbreak{} 建一棵权值线段树，区间 [\allowbreak{}l,\allowbreak{}r]\allowbreak{} 的信息就是 root[\allowbreak{}r]\allowbreak{}-\allowbreak{}root[\allowbreak{}l-\allowbreak{}1]\allowbreak{} 调用约定：kth(\allowbreak{}root[\allowbreak{}l-\allowbreak{}1]\allowbreak{},\allowbreak{}root[\allowbreak{}r]\allowbreak{},\allowbreak{}1,\allowbreak{}sz,\allowbreak{}k)\allowbreak{}、cnt(\allowbreak{}root[\allowbreak{}l-\allowbreak{}1]\allowbreak{},\allowbreak{}root[\allowbreak{}r]\allowbreak{},\allowbreak{}1,\allowbreak{}sz,\allowbreak{}x)\allowbreak{} 第一个参数是「小版本」，写反了会算出 0 或负数 每次插入 +\allowbreak{}1，可先 sum[\allowbreak{}old]\allowbreak{}+\allowbreak{}1 或递归后重算；kth 要求 1<\allowbreak{}=\allowbreak{}k<\allowbreak{}=\allowbreak{}区间元素数。}
\begin{lstlisting}
const int N=200005;
const int LOG=20;       // log2(N)+2
int n,sz;               // n 元素个数，sz 离散化后互不相同的值个数
int a[N],b[N];          // a 原数组，b 离散化后的值（升序去重）
struct ChairTree{
    int ls[N*LOG],rs[N*LOG],sum[N*LOG];
    int root[N],tot;
    void init()
    {
        tot=0,root[0]=0;
        ls[0]=rs[0]=sum[0]=0;
    }
    // 在 old 版本上插入 pos，返回新版本根，O(log n)
    int insert(int old,int l,int r,int pos)
    {
        int p=++tot;
        ls[p]=ls[old],rs[p]=rs[old];
        if(l==r){sum[p]=sum[old]+1;return p;}
        int mid=(l+r)>>1;
        if(pos<=mid)ls[p]=insert(ls[old],l,mid,pos);
        else rs[p]=insert(rs[old],mid+1,r,pos);
        sum[p]=sum[ls[p]]+sum[rs[p]];       // 必须放在递归之后
        return p;
    }
    // [l-1,r] 版本相减，求第 k 小，返回离散化下标，O(log n)
    int kth(int u,int v,int l,int r,int k)
    {
        while(l<r)
        {
            int mid=(l+r)>>1;
            int cnt=sum[ls[v]]-sum[ls[u]];      // 值落在左半边的有几个
            if(k<=cnt)u=ls[u],v=ls[v],r=mid;
            else k-=cnt,u=rs[u],v=rs[v],l=mid+1;
        }
        return l;
    }
    // [l-1,r] 里值 <= x（离散化下的位置 x）的个数，O(log n)
    int count_le(int u,int v,int l,int r,int x)
    {
        if(x<=0)return 0;                   // 比最小值还小
        if(r<=x)return sum[v]-sum[u];       // 整段都 <= x
        int mid=(l+r)>>1;
        int res=count_le(ls[u],ls[v],l,mid,x);
        if(x>mid)res+=count_le(rs[u],rs[v],mid+1,r,x);
        return res;
    }
    // 区间 [l,r] 里值等于 x（离散化位置）的个数
    int cnt_value(int u,int v,int x){return count_le(u,v,1,sz,x)-count_le(u,v,1,sz,x-1);}
}ct;
\end{lstlisting}
\note{建树：把 a[\allowbreak{}1.\allowbreak{}.\allowbreak{}n]\allowbreak{} 离散化后建成主席树，去重后的值个数写回全局 sz}
\begin{lstlisting}
void build()
{
    for(int i=1;i<=n;i++)b[i]=a[i];
    sort(b+1,b+n+1);
    sz=unique(b+1,b+n+1)-b-1;
    ct.init();
    for(int i=1;i<=n;i++)
    {
        int pos=lower_bound(b+1,b+sz+1,a[i])-b;
        ct.root[i]=ct.insert(ct.root[i-1],1,sz,pos);
    }
}
\end{lstlisting}
\note{值域 [\allowbreak{}lo,\allowbreak{}hi]\allowbreak{}（真实值）内的数个数：把 lo-\allowbreak{}1、hi 转到离散化位置再做差}
\begin{lstlisting}
int count_range(int L,int R,int lo,int hi)
{
    if(lo>hi||L>R||!sz)return 0;
    int p=lower_bound(b+1,b+sz+1,lo)-b-1;   // 最后一个位置，它的值 < lo
    int q=upper_bound(b+1,b+sz+1,hi)-b-1;   // 最后一个位置，它的值 <= hi
    if(q<=0)return 0;
    if(p<0)p=0;
    return ct.count_le(ct.root[L-1],ct.root[R],1,sz,q)-ct.count_le(ct.root[L-1],ct.root[R],1,sz,p);
}
\end{lstlisting}
% SOURCE: 02-数据结构/动态开点线段树.cpp
\subsection{动态开点线段树}

\note{节点池按需分配左右儿子、区间信息与懒标记，0为空节点；root为根，l/\allowbreak{}r为完整值域，L/\allowbreak{}R为操作区间。空节点的默认信息须符合合并单位元，节点池重置后旧根编号失效。}
\note{动态开点线段树：值域很大（1e9 级）时只开用到的结点 支持区间加 /\allowbreak{} 区间和，单次操作 O(\allowbreak{}log V)\allowbreak{}，V 是值域 空间：每次修改最多新建 O(\allowbreak{}log V)\allowbreak{} 个点，m 次修改共 O(\allowbreak{}m log V)\allowbreak{} 结点用 vector 按需 push\_back（也可以开静态数组池，见注释）}
\begin{lstlisting}
struct DynSeg{
    struct Node{
        int ls,rs;      // 左右儿子下标，0 表示没建出来（等价于整棵空树）
        ll sum,lazy;    // 区间和、加的懒标记
    };
    vector<Node> t;     // t[0] 是空结点哨兵，sum/lazy 都是 0
    ll lo,hi;           // 值域（闭区间），如 1..1e9
    DynSeg(ll lo_=1,ll hi_=1000000000){init(lo_,hi_);}
    void init(ll lo_,ll hi_)
    {
        lo=lo_,hi=hi_;
        t.clear();
        t.push_back(Node{0,0,0,0});//哨兵 t[0]，空结点
        t.push_back(Node{0,0,0,0});//根 t[1]，整棵树一开始只有一个根
    }
    int new_node()
    {
        t.push_back(Node{0,0,0,0});
        return (int)t.size()-1;
    }
    // p 的标记发给两个儿子，儿子不存在就建出来
    // 注意：t 是 vector，push_back 会重新分配内存，不能提前把 t[p] 取成引用
    void push_down(int p,ll l,ll r)
    {
        if(t[p].lazy==0||l==r)return;
        ll v=t[p].lazy,mid=l+(ll)(((__int128)r-l)/2);
        if(!t[p].ls)t[p].ls=new_node();
        if(!t[p].rs)t[p].rs=new_node();
        int L=t[p].ls,R=t[p].rs;        // 建完点之后再取，下标不会失效
        t[L].lazy+=v,t[L].sum+=v*(mid-l+1);
        t[R].lazy+=v,t[R].sum+=v*(r-mid);
        t[p].lazy=0;
    }
    void add(ll L,ll R,ll v){add(L,R,v,lo,hi,1);}
    // 区间 [L,R] 加 v，O(log V)
    void add(ll L,ll R,ll v,ll l,ll r,int p)
    {
        if(R<l||r<L)return;
        if(L<=l&&r<=R)//(l,r)<=(L,R)
        {
            t[p].sum+=v*(r-l+1),t[p].lazy+=v;
            return;
        }
        push_down(p,l,r);
        ll mid=l+(ll)(((__int128)r-l)/2);
        if(L<=mid)
        {
            if(!t[p].ls)t[p].ls=new_node();
            add(L,R,v,l,mid,t[p].ls);
        }
        if(R>mid)
        {
            if(!t[p].rs)t[p].rs=new_node();
            add(L,R,v,mid+1,r,t[p].rs);
        }
        t[p].sum=t[t[p].ls].sum+t[t[p].rs].sum;
    }
    ll sum(ll L,ll R){return sum(L,R,lo,hi,1);}
    // 区间 [L,R] 求和，O(log V)；结点没建出来说明这段全是 0
    ll sum(ll L,ll R,ll l,ll r,int p)
    {
        if(!p||R<l||r<L)return 0;             // 空结点直接返回 0，不用再往下走
        if(L<=l&&r<=R)return t[p].sum;
        push_down(p,l,r);
        ll mid=l+(ll)(((__int128)r-l)/2),res=0;
        if(L<=mid)res+=sum(L,R,l,mid,t[p].ls);
        if(R>mid)res+=sum(L,R,mid+1,r,t[p].rs);
        return res;
    }
    // 单点加就是 add(x,x,v)
    void point_add(ll x,ll v){add(x,x,v);}
    int nodes(){return (int)t.size()-1;}//已用结点数
};
\end{lstlisting}
\note{静态数组池写法（赛场上省 vector 开销），容量按 m*log V 估： const int MAXNODE=\allowbreak{}2000005; int ls[\allowbreak{}MAXNODE]\allowbreak{},\allowbreak{}rs[\allowbreak{}MAXNODE]\allowbreak{};ll sum[\allowbreak{}MAXNODE]\allowbreak{},\allowbreak{}lazy[\allowbreak{}MAXNODE]\allowbreak{};int tot; int new\_node(\allowbreak{})\allowbreak{}\{+\allowbreak{}+\allowbreak{}tot;ls[\allowbreak{}tot]\allowbreak{}=\allowbreak{}rs[\allowbreak{}tot]\allowbreak{}=\allowbreak{}sum[\allowbreak{}tot]\allowbreak{}=\allowbreak{}lazy[\allowbreak{}tot]\allowbreak{}=\allowbreak{}0;return tot;\}}
\begin{lstlisting}
const ll V=1000000000;//值域 1..1e9
const int XN=4005;    // 对拍用到的坐标个数
\end{lstlisting}
% SOURCE: 02-数据结构/线段树合并.cpp
\subsection{线段树合并}

\note{adj存树的邻点，add\_edge只加一个方向，无向树须调用两次。col[\allowbreak{}u]\allowbreak{}为颜色，root[\allowbreak{}u]\allowbreak{}为子树权值线段树根，ans[\allowbreak{}u]\allowbreak{}为不同颜色数；子树合并进父根后，旧根不可再当独立版本使用。换树清空adj，重置节点池、根和答案。}
\note{线段树合并（权值线段树合并）：求每个子树内不同颜色数 /\allowbreak{} 每个权值出现次数 建树 O(\allowbreak{}n log n)\allowbreak{}；合并总复杂度 O(\allowbreak{}结点总数)\allowbreak{} =\allowbreak{} O(\allowbreak{}n log n)\allowbreak{} 关键性质：merge 里只要有一边是空就直接返回另一边， 否则两边都不空 -\allowbreak{}>\allowbreak{} 当前这个结点被丢弃，每个结点一生只会被丢一次， 所以所有 merge 加起来是 O(\allowbreak{}总结点数)\allowbreak{}，均摊到一次合并是 O(\allowbreak{}log n)\allowbreak{} 坑：合并是「破坏性」的，merge(\allowbreak{}a,\allowbreak{}b)\allowbreak{} 之后 b 的子结点已经被挂到 a 上，b 不能再单独用}
\begin{lstlisting}
const int N=100005;
const int LOG=18;      // log2(N)+1
const int MAXNODE=N*(LOG+1);//n 个叶子各建一条到根的链
int n,col[N];
vector<int> adj[N];
void add_edge(int u,int v)
{
    adj[u].push_back(v);
}
struct MergeSeg{
    int ls[MAXNODE],rs[MAXNODE],sum[MAXNODE],cnt[MAXNODE],tot;//cnt 是这段里被占用的叶子数
    void init(){tot=0;}
    int new_node()
    {
        ++tot;
        ls[tot]=rs[tot]=sum[tot]=cnt[tot]=0;
        return tot;
    }
    // 在位置 pos 插入 cnt 个（叶子），返回根，O(log n)
    int insert(int p,int l,int r,int pos,int c)
    {
        if(!p)p=new_node();
        if(l==r)
        {
            if(sum[p]==0)cnt[p]=1;//这个叶子原来没有颜色，现在有了，占用数 +1
            sum[p]+=c;
            return p;
        }
        int mid=(l+r)>>1;
        if(pos<=mid)ls[p]=insert(ls[p],l,mid,pos,c);
        else rs[p]=insert(rs[p],mid+1,r,pos,c);
        sum[p]=sum[ls[p]]+sum[rs[p]];
        cnt[p]=cnt[ls[p]]+cnt[rs[p]];
        return p;
    }
    // 把 b 合并进 a，返回合并后的根，均摊 O(log n)；要求 a、b 表示的颜色集合不交
    // cnt 是「这段里被占用的叶子数（不同颜色数）」
    // 关键：递归会把子树的 cnt 改掉，所以先把 ca、cb 和两个儿子的原值存下来，
    //       再用「合并前两边之和 - 合并后的交集叶子数」推出共有几种颜色
    int merge(int a,int b)
    {
        if(!a||!b)return a|b;       // 有一边空，直接接过去；b 空则只剩 a
        int ca=cnt[a],cb=cnt[b];
        int la=ls[a],lb=ls[b],ra=rs[a],rb=rs[b];
        if(!la&&!lb&&!ra&&!rb)      // 两个都是叶子结点（同一个颜色），并成一个
        {
            sum[a]+=sum[b],cnt[a]=1;
            return a;
        }
        int cla=cnt[la],clb=cnt[lb],cra=cnt[ra],crb=cnt[rb];
        int mla=merge(la,lb),mra=merge(ra,rb);
        ls[a]=mla,rs[a]=mra;
        int dup=(cla+clb-cnt[mla])+(cra+crb-cnt[mra]);//两边都有的颜色数
        sum[a]=sum[mla]+sum[mra];
        cnt[a]=ca+cb-dup;
        return a;
    }
    // 不同颜色数 = 权值线段树里 sum>0 的叶子个数，直接返回维护好的 cnt
    int count_distinct(int a){return cnt[a];}
}seg;
int ans[N];//每个子树内不同颜色数
int root[N];//每个点对应的权值线段树根
\end{lstlisting}
\note{自底向上：把自己的颜色建树，再把每个儿子的树合并进来，最后统计根上的 sum}
\begin{lstlisting}
void dfs(int u,int fa)
{
    root[u]=seg.insert(0,1,n,col[u],1);
    for(int v:adj[u])
    {
        if(v==fa)continue;
        dfs(v,u);
        root[u]=seg.merge(root[u],root[v]);
    }
    ans[u]=seg.count_distinct(root[u]);
}
\end{lstlisting}
% SOURCE: 02-数据结构/线段树分裂与合并.cpp
\subsection{线段树分裂与合并}

\note{节点池t存左右儿子和频数和，0为空根；多棵树用不同根编号共享池。split通过引用修改原根并返回抽出的区间根，merge消耗两棵旧树；非持久化，旧根不能继续独立操作，计数不可减成负数。}
\note{动态权值线段树，多重集频数非负；拆出 [\allowbreak{}L,\allowbreak{}R]\allowbreak{} 返回新根，原根按引用更新。 split O(\allowbreak{}log \ensuremath{\sigma})\allowbreak{} 边界节点，merge 摊还 O(\allowbreak{}总节点数)\allowbreak{}；所有操作是破坏性的，根不能共享节点。 节点用 vector 扩容；递归调用不得持有 vector 元素引用。值域 0.\allowbreak{}.\allowbreak{}\ensuremath{\sigma}-\allowbreak{}1。}
\begin{lstlisting}
struct SplitMergeSeg
{
    struct Node{int l=0,r=0;long long sum=0;};vector<Node> t={{}};int sigma;
    SplitMergeSeg(int sigma):sigma(sigma){assert(sigma>0);}
    int node(){t.push_back({});return t.size()-1;}
    void pull(int p){t[p].sum=t[t[p].l].sum+t[t[p].r].sum;}
    int add(int p,int l,int r,int x,int delta){if(!p)p=node();if(l==r){t[p].sum+=delta;assert(t[p].sum>=0);return p;}int m=(l+r)/2;if(x<=m){int q=add(t[p].l,l,m,x,delta);t[p].l=q;}else{int q=add(t[p].r,m+1,r,x,delta);t[p].r=q;}pull(p);return p;}
    int split(int &p,int l,int r,int L,int R)
    {
        if(!p||R<l||r<L)return 0;if(L<=l&&r<=R){int q=p;p=0;return q;}
        int q=node(),m=(l+r)/2,a=t[p].l,b=t[p].r;
        int x=split(a,l,m,L,R),y=split(b,m+1,r,L,R);
        t[p].l=a;t[p].r=b;t[q].l=x;t[q].r=y;pull(p);pull(q);return q;
    }
    int merge(int a,int b,int l,int r){if(!a||!b)return a?a:b;if(l==r)t[a].sum+=t[b].sum;else{int m=(l+r)/2;t[a].l=merge(t[a].l,t[b].l,l,m);t[a].r=merge(t[a].r,t[b].r,m+1,r);pull(a);}return a;}
};
\end{lstlisting}
% SOURCE: 02-数据结构/扫描线(矩形面积并).cpp
\subsection{扫描线(矩形面积并)}

\note{rx1/\allowbreak{}ry1、rx2/\allowbreak{}ry2是各矩形的两端，ys为离散化纵坐标；覆盖长度按相邻坐标的实际差计算。矩形用半开边界解释面积，退化矩形贡献0；g是小范围验证用网格，实际扫描线不依赖它。}
\note{扫描线求矩形面积并、周长并，O(\allowbreak{}n log n)\allowbreak{} 周长：两个方向分别累计覆盖长度变化，同坐标先加入后删除}
\begin{lstlisting}
const int N=100005;
int n;
int rx1[N],ry1[N],rx2[N],ry2[N],ys[N<<1];
struct ScanTree{
    ll cov[N<<3],len[N<<3],num[N<<3];
    bool lc[N<<3],rc[N<<3];
    int *ys;
    void init(int *y){ys=y;}
    void clear(int cnt)
    {
        for(int i=1;i<=(cnt<<2);i++)cov[i]=len[i]=num[i]=0,lc[i]=rc[i]=0;
    }
    void push_up(int p,int l,int r)
    {
        if(cov[p]>0)
        {
            len[p]=(ll)ys[r+1]-ys[l];
            num[p]=1,lc[p]=rc[p]=1;
        }
        else if(l==r)len[p]=num[p]=0,lc[p]=rc[p]=0;
        else
        {
            int lp=p<<1,rp=p<<1|1;
            len[p]=len[lp]+len[rp];
            num[p]=num[lp]+num[rp]-(rc[lp]&&lc[rp]);
            lc[p]=lc[lp],rc[p]=rc[rp];
        }
    }
    void update(int L,int R,ll v,int l,int r,int p)
    {
        if(L<=l&&r<=R)
        {
            cov[p]+=v;
            push_up(p,l,r);
            return;
        }
        int mid=(l+r)>>1;
        if(L<=mid)update(L,R,v,l,mid,p<<1);
        if(R>mid)update(L,R,v,mid+1,r,p<<1|1);
        push_up(p,l,r);
    }
}seg;
struct Ev{
    int x,lo,hi,typ;
    bool operator<(const Ev &o)const{return x!=o.x?x<o.x:typ>o.typ;}
}ev[N<<1];
struct SweepRes{ll area,diff;};
SweepRes sweep(int *a1,int *b1,int *a2,int *b2)
{
    int ycnt=0,ec=0;
    for(int i=1;i<=n;i++)
    {
        if(a1[i]>=b1[i]||a2[i]>=b2[i])continue;
        ys[++ycnt]=a2[i],ys[++ycnt]=b2[i];
        ev[++ec]={a1[i],a2[i],b2[i],1};
        ev[++ec]={b1[i],a2[i],b2[i],-1};
    }
    SweepRes ans{0,0};
    if(!ec)return ans;
    sort(ys+1,ys+ycnt+1);
    ycnt=unique(ys+1,ys+ycnt+1)-ys-1;
    sort(ev+1,ev+ec+1);
    seg.init(ys),seg.clear(ycnt);
    for(int i=1;i<=ec;i++)
    {
        if(i>1)ans.area+=seg.len[1]*((ll)ev[i].x-ev[i-1].x);
        int l=lower_bound(ys+1,ys+ycnt+1,ev[i].lo)-ys;
        int r=lower_bound(ys+1,ys+ycnt+1,ev[i].hi)-ys-1;
        ll last=seg.len[1];
        seg.update(l,r,ev[i].typ,1,ycnt-1,1);
        ans.diff+=llabs(seg.len[1]-last);
    }
    return ans;
}
ll union_area()
{
    return sweep(rx1,rx2,ry1,ry2).area;
}
ll union_perimeter()
{
    SweepRes a=sweep(rx1,rx2,ry1,ry2),b=sweep(ry1,ry2,rx1,rx2);
    return a.diff+b.diff;
}
const int TESTN=100005;
int g[64][64];
\end{lstlisting}
% SOURCE: 02-数据结构/Trie字典树.cpp
\subsection{Trie字典树}

\note{t是字典树，s为待操作串；节点转移按字符映射，词尾次数区分完整词与前缀经过次数。根为0，容量按总串长设置；换词典调用clear，单个插入/\allowbreak{}查询的字符串参数用于区分不同操作。}
\note{字典树 Trie：插入 /\allowbreak{} 查询是否出现 /\allowbreak{} 出现次数 /\allowbreak{} 前缀计数 结点数 N 按「所有串长度之和」开，别只按串个数开 设 f[\allowbreak{}j]\allowbreak{} 为某 j 个已插入串能共享的最长前缀：节点深 h、经过数 cnt 增 1 时，只更新 f[\allowbreak{}cnt]\allowbreak{}。 更小 j 的答案此前已>\allowbreak{}=\allowbreak{}h；每次插入只访问该串路径，总 O(\allowbreak{}字符串总长度)\allowbreak{}，不必区间 chmax。}
\begin{lstlisting}
const int N=200005;
\end{lstlisting}
\note{输入仅允许字母数字；其他字符需修改 id。Trie 实例须零初始化（例如全局 t）。 字符映射：小写 0\textasciitilde{}25，大写 26\textasciitilde{}51，数字 52\textasciitilde{}61（别的字符集改这里）}
\begin{lstlisting}
int id(char c)
{
    if(c>='a'&&c<='z')return c-'a';
    if(c>='A'&&c<='Z')return c-'A'+26;
    if(c>='0'&&c<='9')return c-'0'+52;
    return 0;
}
struct Trie{
    int tr[N][62];      // tr[p][c] 儿子编号，0 表示没有
    int word[N];        // 以 p 结尾的串个数
    int pre[N];         // 经过 p 的串个数（含以 p 结尾的）
    int tot=0;            // 已用结点数，根固定为 0，所以从 1 开始编号
    void clear()        // 清空全部：O(结点数*62)，只有一个测试点时可省
    {
        for(int i=0;i<=tot;i++)
        {
            word[i]=0,pre[i]=0;
            for(int j=0;j<62;j++)tr[i][j]=0;
        }
        tot=0;
    }
    void insert(const string &s)              // O(|s|)
    {
        int p=0,len=s.length();
        pre[0]++;
        for(int i=0;i<len;i++)
        {
            int c=id(s[i]);
            if(!tr[p][c])tr[p][c]=++tot;
            p=tr[p][c];
            pre[p]++;
        }
        word[p]++;
    }
    bool exist(const string &s)               // 是否作为完整串出现过，O(|s|)
    {
        int p=0,len=s.length();
        for(int i=0;i<len;i++)
        {
            int c=id(s[i]);
            if(!tr[p][c])return false;
            p=tr[p][c];
        }
        return word[p]>0;
    }
    int count_word(const string &s)           // 完整串出现次数，O(|s|)
    {
        int p=0,len=s.length();
        for(int i=0;i<len;i++)
        {
            int c=id(s[i]);
            if(!tr[p][c])return 0;
            p=tr[p][c];
        }
        return word[p];
    }
    int count_pre(const string &s)            // 以 s 为前缀的串个数（含 s 本身），O(|s|)
    {
        int p=0,len=s.length();
        for(int i=0;i<len;i++)
        {
            int c=id(s[i]);
            if(!tr[p][c])return 0;
            p=tr[p][c];
        }
        return pre[p];
    }
};
Trie t;
string s;
int n;
\end{lstlisting}
% SOURCE: 02-数据结构/01Trie(最大异或).cpp
\subsection{01Trie(最大异或)}

\note{BinaryTrie的bits是有效位数，节点计数支持重复数和删除；输入为无符号整数。max\_xor返回异或值，匹配的原数等于输入异或返回值；查询要求非空，erase删除一个副本。}
\note{无符号整数多重集：插入、删一个副本、查询 max(\allowbreak{}x\textasciicircum{}y)\allowbreak{}，O(\allowbreak{}B)\allowbreak{}，B=\allowbreak{}32 含最高位。 修改为较小 B 时所有输入必须小于 2\textasciicircum{}B；空集合不能查最大异或。 删除不回收节点，空间 O(\allowbreak{}插入过的不同数*B)\allowbreak{}；要输出 y，用 x\textasciicircum{}max\_xor(\allowbreak{}x)\allowbreak{}。}
\begin{lstlisting}
struct BinaryTrie
{
    struct Node{int ch[2]={0,0},cnt=0;};
    vector<Node> tr=vector<Node>(1);
    int B;
    BinaryTrie(int bits=32):B(bits){assert(1<=B&&B<=32);}
    void clear(){tr.assign(1,Node{});}
    void insert(uint32_t x)
    {
        assert(B==32||(x>>B)==0);
        int p=0;
        ++tr[p].cnt;
        for(int i=B-1;i>=0;i--)
        {
            int b=x>>i&1;
            if(!tr[p].ch[b])
            {
                tr[p].ch[b]=tr.size();
                tr.push_back(Node{});
            }
            p=tr[p].ch[b],++tr[p].cnt;
        }
    }
    bool erase(uint32_t x)
    {
        assert(B==32||(x>>B)==0);
        int p=0;
        for(int i=B-1;i>=0;i--)
        {
            p=tr[p].ch[x>>i&1];
            if(!p||!tr[p].cnt)return false;
        }
        p=0,--tr[p].cnt;
        for(int i=B-1;i>=0;i--)p=tr[p].ch[x>>i&1],--tr[p].cnt;
        return true;
    }
    uint32_t max_xor(uint32_t x)
    {
        assert(tr[0].cnt>0&&(B==32||(x>>B)==0));
        uint32_t ans=0;
        int p=0;
        for(int i=B-1;i>=0;i--)
        {
            int b=x>>i&1,q=tr[p].ch[b^1];
            if(q&&tr[q].cnt)ans|=uint32_t(1)<<i,p=q;
            else p=tr[p].ch[b];
        }
        return ans;
    }
};
\end{lstlisting}
% SOURCE: 02-数据结构/bitset与字并行.cpp
\subsection{bitset与字并行}

\note{a存非负物品重量；f[\allowbreak{}s]\allowbreak{}是和s的可达性，每次背包清空f。g[\allowbreak{}i]\allowbreak{}[\allowbreak{}j]\allowbreak{}是0起有向邻接关系，closure(\allowbreak{})\allowbreak{}原地改成可达关系并加入零长路径；换图重新设置g。N按查询和上限或点数设置。}
\note{把布尔状态或集合压成位向量，每机器字同时处理多位。整数掩码枚举与计数见第一章位运算；这里保留长布尔向量的两种实际应用。}
\topic{容量与基本操作}
\note{bitset<\allowbreak{}N>\allowbreak{}容量在编译期固定。与、或、异或、取反、移位用于集合交并、差和状态平移；any/\allowbreak{}none检查空集，count计位数，test读取单个位。集合差用a\&\textasciitilde{}b；取反只针对N位全集。复杂度按N除以机器字宽W估算，不能因写成一行就当O(\allowbreak{}1)\allowbreak{}。}
\begin{lstlisting}
const int N=513;
\end{lstlisting}
\topic{非负01背包可达性}
\note{填全局a后调用subset\_sum(\allowbreak{})\allowbreak{}，结果在全局f，也作为返回值给出。f[\allowbreak{}s]\allowbreak{}表示能否取出和s，每次清空后只有f[\allowbreak{}0]\allowbreak{}。每件物品执行f|=\allowbreak{}f<\allowbreak{}<\allowbreak{}x：右侧使用更新前状态，所以每件只用一次。非负输入允许截断到N-\allowbreak{}1，高位和不会再减回查询范围；负数不能套。时间 $O(|a|N/W)$，空间N位，按查询上限设置N。}
\begin{lstlisting}
vector<int> a;
bitset<N> f;
bitset<N> subset_sum()
{
    f.reset();
    f[0]=1;
    for(int x:a)
    {
        assert(x>=0);
        f|=f<<x;
    }
    return f;
}
\end{lstlisting}
\topic{布尔传递闭包}
\note{填全局g的邻接位向量后调用closure(\allowbreak{})\allowbreak{}，g原地变成可达关系。以第k点为中转，若i能到k，就并入k的整行可达集合。k必须在最外层，顶点编号0.\allowbreak{}.\allowbreak{}n-\allowbreak{}1且n\ensuremath{\le}N；代码加入长度0的路径。输入n以上的位须清零，换图重新设置g；只回答可达性，不统计路径数。时间 $O(n^2N/W)$，空间nN位。}
\begin{lstlisting}
vector<bitset<N>> g;
void closure()
{
    int n=g.size();
    assert(n<=N);
    for(int i=0;i<n;i++)g[i][i]=1;
    for(int k=0;k<n;k++)
        for(int i=0;i<n;i++)if(g[i][k])g[i]|=g[k];
}
\end{lstlisting}
% SOURCE: 02-数据结构/分块.cpp
\subsection{分块}

\note{a[\allowbreak{}1.\allowbreak{}.\allowbreak{}n]\allowbreak{}是输入，blk维护块和与懒标记；br仅是独立朴素对照缓冲，比赛代码可省去。初始化传n与输入数组，操作传区间及增量；同一题的块状态可放全局对象中，换数组重新init。}
\begin{lstlisting}
const int N=100005;
\end{lstlisting}
\note{分块：区间加 +\allowbreak{} 区间和，块长 sqrt(\allowbreak{}n)\allowbreak{} update/\allowbreak{}query 复杂度 O(\allowbreak{}sqrt(\allowbreak{}n)\allowbreak{})\allowbreak{}：散块暴力，整块用 tag/\allowbreak{}sum}
\begin{lstlisting}
template<typename T,int NMAX=N>
struct Block
{
    T a[NMAX],tag[NMAX],sum[NMAX];
    int id[NMAX],L[NMAX],R[NMAX],B,cnt,n;
    void init(int n_,T *src)//用 1..n_ 的原数组建块
    {
        n=n_;
        B=max(1,(int)sqrt((double)n));//块长
        cnt=(n-1)/B+1;
        for(int i=1;i<=cnt;i++)
        {
            L[i]=(i-1)*B+1;
            R[i]=min(i*B,n);
            tag[i]=0,sum[i]=0;
        }
        for(int i=1;i<=n;i++)
        {
            a[i]=src[i];
            id[i]=(i-1)/B+1;
            sum[id[i]]+=a[i];
        }
    }
    void update(int l,int r,T c)
    {
        if(id[l]==id[r])//同块
        {
            for(int i=l;i<=r;i++)a[i]+=c,sum[id[l]]+=c;
            return;
        }
        for(int i=l;i<=R[id[l]];i++)a[i]+=c,sum[id[l]]+=c;
        for(int i=L[id[r]];i<=r;i++)a[i]+=c,sum[id[r]]+=c;
        for(int i=id[l]+1;i<id[r];i++)tag[i]+=c,sum[i]+=(T)(R[i]-L[i]+1)*c;
    }
    T query(int l,int r)
    {
        T s=0;
        if(id[l]==id[r])
        {
            for(int i=l;i<=r;i++)s+=a[i]+tag[id[l]];
            return s;
        }
        for(int i=l;i<=R[id[l]];i++)s+=a[i]+tag[id[l]];
        for(int i=L[id[r]];i<=r;i++)s+=a[i]+tag[id[r]];
        for(int i=id[l]+1;i<id[r];i++)s+=sum[i];
        return s;
    }
};
int n;
ll a[205],br[205];
Block<ll,205> blk;
\end{lstlisting}
% SOURCE: 02-数据结构/莫队算法.cpp
\subsection{莫队算法}

\note{a[\allowbreak{}1.\allowbreak{}.\allowbreak{}n]\allowbreak{}为离散化值，m为查询数，block为分块长度；cnt[\allowbreak{}x]\allowbreak{}是当前区间中值x的频数，now为当前统计量，ans按原编号存结果。指针移动前从空区间开始，换查询组清空频数。}
\note{莫队算法（普通莫队）：离线处理区间询问，本例求「区间不同数个数」 分块 +\allowbreak{} 排序后左右指针总移动 O(\allowbreak{}n*sqrt(\allowbreak{}n)\allowbreak{})\allowbreak{}，单次移动 O(\allowbreak{}1)\allowbreak{} 排序时奇数块右端点反过来排，常数更小（这行别删）}
\begin{lstlisting}
const int N=100005;
int n,m,block;
int a[N],ans[N],cnt[N],now;
struct Query{
    int l,r,id;
}q[N];
bool cmp(const Query &x,const Query &y)
{
    int bx=x.l/block,by=y.l/block;
    if(bx!=by)return bx<by;
    if(bx&1)return x.r>y.r;         // 奇数块右端点递减
    return x.r<y.r;
}
void add(int pos)
{
    if(cnt[a[pos]]++==0)now++;
}
void del(int pos)
{
    if(--cnt[a[pos]]==0)now--;
}
\end{lstlisting}
\note{把询问排好序后统一回答，结果写回 ans[\allowbreak{}id]\allowbreak{}}
\begin{lstlisting}
void solve()
{
    block=max(1,(int)sqrt(n));
    sort(q+1,q+m+1,cmp);
    int l=1,r=0;
    memset(cnt,0,sizeof(cnt));
    now=0;
    for(int i=1;i<=m;i++)
    {
        while(l>q[i].l)add(--l);
        while(r<q[i].r)add(++r);
        while(l<q[i].l)del(l++);
        while(r>q[i].r)del(r--);
        ans[q[i].id]=now;
    }
}
\end{lstlisting}
% SOURCE: 02-数据结构/莫队变体.cpp
\subsection{莫队变体}

\note{所有下标0起，值和颜色已离散化到0.\allowbreak{}.\allowbreak{}sigma-\allowbreak{}1。带修改版读a、change和query\{l,\allowbreak{}r,\allowbreak{}time\}，执行后a停在最后处理的时间状态；树上版读g、color和queries点对；回滚版读a和range\_query闭区间。均按原查询顺序返回，计数、排序和回滚记录每次重建。}
\note{三种莫队：输入值已经离散化到 0.\allowbreak{}.\allowbreak{}sigma-\allowbreak{}1；例题指标供替换 add/\allowbreak{}remove/\allowbreak{}check。 带修改：区间不同数，记录修改前后值；排序块长 n\textasciicircum{}(\allowbreak{}2/\allowbreak{}3)\allowbreak{}，典型移动量 O(\allowbreak{}n\textasciicircum{}(\allowbreak{}5/\allowbreak{}3)\allowbreak{})\allowbreak{}（q\textasciitilde{}n）。 树上：路径不同色，Euler 每点出现两次，进出时翻转 active；LCA 额外加入一次。 回滚：区间同值最远距离，只增加右端、临时增加左端后撤销，O(\allowbreak{}(\allowbreak{}n+\allowbreak{}q)\allowbreak{}sqrt n)\allowbreak{}。}
\begin{lstlisting}
struct Modification{int pos,before,after;};
struct TimedQuery{int l,r,time;};
vector<int> a,color;
vector<Modification> change;
vector<TimedQuery> query;
vector<vector<int>> g;
vector<pair<int,int>> queries,range_query;
vector<int> modified_mo(int sigma)
{
    int n=a.size(),block=max(1,(int)pow(max(1,n),2.0/3));vector<int> order(query.size());iota(order.begin(),order.end(),0);
    sort(order.begin(),order.end(),[&](int x,int y){auto a=query[x],b=query[y];return tuple(a.l/block,a.r/block,a.time)<tuple(b.l/block,b.r/block,b.time);});
    vector<int> freq(sigma),ans(query.size());int l=0,r=-1,t=0,different=0;
    auto add=[&](int x){if(!freq[x]++)different++;};auto del=[&](int x){if(!--freq[x])different--;};
    auto apply=[&](int i,bool forward){auto c=change[i];if(l<=c.pos&&c.pos<=r)del(a[c.pos]);a[c.pos]=forward?c.after:c.before;if(l<=c.pos&&c.pos<=r)add(a[c.pos]);};
    for(int id:order){auto q=query[id];while(t<q.time)apply(t++,true);while(t>q.time)apply(--t,false);while(l>q.l)add(a[--l]);while(r<q.r)add(a[++r]);while(l<q.l)del(a[l++]);while(r>q.r)del(a[r--]);ans[id]=different;}
    return ans;
}
vector<int> tree_mo(int sigma)
{
    int n=g.size();if(!n){assert(queries.empty());return {};}
    int log=1;while((1LL<<log)<=n)log++;vector<vector<int>> up(log,vector<int>(n));vector<int> tin(n),tout(n),depth(n),euler,idx(n),parent(n,-1),stack={0};
    parent[0]=0;
    while(!stack.empty())
    {int u=stack.back();if(idx[u]==0){tin[u]=euler.size();euler.push_back(u);up[0][u]=parent[u];for(int k=1;k<log;k++)up[k][u]=up[k-1][up[k-1][u]];}
        if(idx[u]<(int)g[u].size()){int v=g[u][idx[u]++];if(v==parent[u])continue;parent[v]=u;depth[v]=depth[u]+1;stack.push_back(v);}else{tout[u]=euler.size();euler.push_back(u);stack.pop_back();}}
    auto lca=[&](int a,int b){if(depth[a]<depth[b])swap(a,b);int d=depth[a]-depth[b];for(int k=0;k<log;k++)if(d>>k&1)a=up[k][a];if(a==b)return a;for(int k=log-1;k>=0;k--)if(up[k][a]!=up[k][b])a=up[k][a],b=up[k][b];return up[0][a];};
    struct Q{int l,r,extra,id;};vector<Q> q;
    for(int i=0;i<(int)queries.size();i++){auto [u,v]=queries[i];if(tin[u]>tin[v])swap(u,v);int a=lca(u,v);if(a==u)q.push_back({tin[u],tin[v],-1,i});else q.push_back({tout[u],tin[v],a,i});}
    int block=max(1,(int)sqrt(2*n));sort(q.begin(),q.end(),[&](Q a,Q b){if(a.l/block!=b.l/block)return a.l<b.l;return (a.l/block&1)?a.r>b.r:a.r<b.r;});
    vector<int> freq(sigma),ans(q.size());vector<bool> active(n);int distinct=0,l=0,r=-1;
    auto flip=[&](int u){int c=color[u];if(active[u]){if(!--freq[c])distinct--;}else if(!freq[c]++)distinct++;active[u]=!active[u];};
    for(auto a:q){while(l>a.l)flip(euler[--l]);while(r<a.r)flip(euler[++r]);while(l<a.l)flip(euler[l++]);while(r>a.r)flip(euler[r--]);if(a.extra>=0)flip(a.extra);ans[a.id]=distinct;if(a.extra>=0)flip(a.extra);}
    return ans;
}
vector<int> rollback_mo(int sigma)
{
    int n=a.size(),block=max(1,(int)sqrt(max(1,n)));vector<int> order(range_query.size()),ans(range_query.size());iota(order.begin(),order.end(),0);
    sort(order.begin(),order.end(),[&](int i,int j){return pair(range_query[i].first/block,range_query[i].second)<pair(range_query[j].first/block,range_query[j].second);});
    vector<int> first(sigma,n),last(sigma,-1);struct Change{int v,first,last,best;};vector<Change> log;int best=0;
    auto add=[&](int pos){int v=a[pos];log.push_back({v,first[v],last[v],best});first[v]=min(first[v],pos);last[v]=max(last[v],pos);best=max(best,last[v]-first[v]);};
    auto rollback=[&](int size){while((int)log.size()>size){auto c=log.back();log.pop_back();first[c.v]=c.first;last[c.v]=c.last;best=c.best;}};
    int current=-1,right=-1;
    for(int id:order)
    {
        auto [l,r]=range_query[id];int b=l/block,end=min(n,(b+1)*block);
        if(b!=current){rollback(0);current=b;right=end-1;}
        if(r<end){int save=log.size();for(int p=l;p<=r;p++)add(p);ans[id]=best;rollback(save);}
        else{while(right<r)add(++right);int save=log.size();for(int p=end-1;p>=l;p--)add(p);ans[id]=best;rollback(save);}
    }
    return ans;
}
\end{lstlisting}
% SOURCE: 02-数据结构/CDQ分治(三维偏序).cpp
\subsection{CDQ分治(三维偏序)}

\note{a是0起三维点集，每点为\{x,\allowbreak{}y,\allowbreak{}z\}；dominance(\allowbreak{})\allowbreak{}返回按原输入顺序的支配计数，相同点计入其他副本，不计自己。每次重建排序、离散化、BIT及计数。}
\note{若 q[\allowbreak{}j]\allowbreak{}=\allowbreak{}sum(\allowbreak{}k<\allowbreak{}=\allowbreak{}j)\allowbreak{}(\allowbreak{}j-\allowbreak{}k+\allowbreak{}1)\allowbreak{}*f[\allowbreak{}k]\allowbreak{}，则 f[\allowbreak{}j]\allowbreak{}=\allowbreak{}q[\allowbreak{}j]\allowbreak{}-\allowbreak{}2*q[\allowbreak{}j-\allowbreak{}1]\allowbreak{}+\allowbreak{}q[\allowbreak{}j-\allowbreak{}2]\allowbreak{}；二阶差分消前缀转移。 区间截断需补边界项；再判断剩余转移是否是卷积，才组合 CDQ+\allowbreak{}NTT。}
\note{每个点统计其他满足 x<\allowbreak{}=\allowbreak{}x\_i,\allowbreak{}y<\allowbreak{}=\allowbreak{}y\_i,\allowbreak{}z<\allowbreak{}=\allowbreak{}z\_i 的点，按原输入顺序返回。 相同三元组先合并，组内另加 w-\allowbreak{}1；不能直接拿重复点做 CDQ。 x 排序后分治，跨左右统计 y，树状数组统计 z；O(\allowbreak{}n log\ensuremath{^2}n)\allowbreak{}，空间 O(\allowbreak{}n)\allowbreak{}。}
\begin{lstlisting}
vector<array<int,3>> a;
vector<ll> dominance()
{
    struct Node{int x,y,z,w,id;ll ans;};
    int n=a.size();
    vector<int> order(n),group(n),zs;
    iota(order.begin(),order.end(),0);
    sort(order.begin(),order.end(),[&](int i,int j){return a[i]<a[j];});
    vector<Node> p;
    for(int i:order)
    {
        if(p.empty()||array<int,3>{p.back().x,p.back().y,p.back().z}!=a[i])
            p.push_back({a[i][0],a[i][1],a[i][2],0,(int)p.size(),0});
        ++p.back().w,group[i]=p.back().id;
    }
    for(auto v:p)zs.push_back(v.z);
    sort(zs.begin(),zs.end());
    zs.erase(unique(zs.begin(),zs.end()),zs.end());
    for(auto &v:p)v.z=lower_bound(zs.begin(),zs.end(),v.z)-zs.begin()+1;
    vector<int> bit(zs.size()+1);
    auto add=[&](int x,int v){for(;x<(int)bit.size();x+=x&-x)bit[x]+=v;};
    auto sum=[&](int x){int s=0;for(;x;x-=x&-x)s+=bit[x];return s;};
    auto cdq=[&](auto &&self,int l,int r)->void
    {
        if(r-l<=1)return;
        int m=(l+r)/2;
        self(self,l,m),self(self,m,r);
        int i=l;
        for(int j=m;j<r;j++)
        {
            while(i<m&&p[i].y<=p[j].y)add(p[i].z,p[i].w),++i;
            p[j].ans+=sum(p[j].z);
        }
        for(int j=l;j<i;j++)add(p[j].z,-p[j].w);
        inplace_merge(p.begin()+l,p.begin()+m,p.begin()+r,
            [](Node a,Node b){return a.y<b.y;});
    };
    cdq(cdq,0,p.size());
    vector<ll> count(p.size()),ans(n);
    for(auto v:p)count[v.id]=v.ans+v.w-1;
    for(int i=0;i<n;i++)ans[i]=count[group[i]];
    return ans;
}
\end{lstlisting}
% SOURCE: 02-数据结构/整体二分(带修改区间第k小).cpp
\subsection{整体二分(带修改区间第k小)}

\note{a有效下标1.\allowbreak{}.\allowbreak{}n；ops按时间存Operation，type=\allowbreak{}0时\{x=\allowbreak{}位置,\allowbreak{}y=\allowbreak{}新值\}，type=\allowbreak{}1时\{x=\allowbreak{}l,\allowbreak{}y=\allowbreak{}r,\allowbreak{}k=\allowbreak{}排名\}。函数会把a更新到最后的修改值，返回查询顺序的答案；每次重建值域、事件和BIT。递归分组的事件容器分别属于不同分支，保留移动传递。}
\note{a[\allowbreak{}1.\allowbreak{}.\allowbreak{}n]\allowbreak{}；0 操作：\{0,\allowbreak{}位置,\allowbreak{}新值,\allowbreak{}0\}；1 操作：\{1,\allowbreak{}l,\allowbreak{}r,\allowbreak{}k\}。}
\begin{lstlisting}
struct Operation{int type,x,y,k;};
\end{lstlisting}
\note{按值域二分：<\allowbreak{}=\allowbreak{}mid 的修改进入 BIT，查询数出区间内这一半有几个。 查询去右半时 k 减掉左半数量；分组必须保持原时间顺序，递归前撤销 BIT。 修改拆成旧值 -\allowbreak{}1、新值 +\allowbreak{}1；O(\allowbreak{}(\allowbreak{}n+\allowbreak{}q)\allowbreak{} log V log n)\allowbreak{}，V 是不同值数。 空间 O(\allowbreak{}n+\allowbreak{}q)\allowbreak{}：分组后释放父层事件，避免同一批事件在多层递归中重复保留。}
\begin{lstlisting}
vector<int> a;
vector<Operation> ops;
vector<int> range_kth()
{
    int n=(int)a.size()-1;
    assert(n>=0);
    vector<int> value(a.begin()+1,a.end());
    for(auto e:ops)if(e.type==0)value.push_back(e.y);
    sort(value.begin(),value.end());
    value.erase(unique(value.begin(),value.end()),value.end());
    auto id=[&](int x){return lower_bound(value.begin(),value.end(),x)-value.begin();};
    struct Event{int type,x,y,k,delta,id;};
    vector<Event> events;
    for(int i=1;i<=n;i++)events.push_back({0,i,(int)id(a[i]),0,1,0});
    int queries=0;
    for(auto e:ops)
    {
        if(e.type==0)
        {
            assert(1<=e.x&&e.x<=n);
            events.push_back({0,e.x,(int)id(a[e.x]),0,-1,0});
            events.push_back({0,e.x,(int)id(e.y),0,1,0});
            a[e.x]=e.y;
        }
        else
        {
            assert(e.type==1&&1<=e.x&&e.x<=e.y&&e.y<=n&&1<=e.k&&e.k<=e.y-e.x+1);
            events.push_back({1,e.x,e.y,e.k,0,queries++});
        }
    }
    vector<int> ans(queries),bit(n+1);
    if(!queries)return ans;
    auto add=[&](int x,int v){for(;x<=n;x+=x&-x)bit[x]+=v;};
    auto sum=[&](int x){int s=0;for(;x;x-=x&-x)s+=bit[x];return s;};
    auto solve=[&](auto &&self,int l,int r,vector<Event> e)->void
    {
        if(e.empty())return;
        if(l==r){for(auto t:e)if(t.type==1)ans[t.id]=value[l];return;}
        int m=(l+r)/2;
        vector<Event> left,right;
        for(auto t:e)
        {
            if(t.type==0)
            {
                if(t.y<=m)add(t.x,t.delta),left.push_back(t);
                else right.push_back(t);
            }
            else
            {
                int count=sum(t.y)-sum(t.x-1);
                if(t.k<=count)left.push_back(t);
                else t.k-=count,right.push_back(t);
            }
        }
        for(auto t:left)if(t.type==0)add(t.x,-t.delta);
        vector<Event>().swap(e);
        self(self,l,m,move(left)),self(self,m+1,r,move(right));
    };
    solve(solve,0,value.size()-1,move(events));
    return ans;
}
\end{lstlisting}
% SOURCE: 02-数据结构/整体二分(前缀判定).cpp
\subsection{整体二分(前缀判定)}

\note{点1.\allowbreak{}.\allowbreak{}n；e按加入顺序存无向边，query存要问的点对，两个向量从0起。返回每问最早连通的前缀长度，0表示初始已连通，e.\allowbreak{}size(\allowbreak{})\allowbreak{}+\allowbreak{}1表示永不连通。每次重建二分范围、分桶和DSU。}
\note{给定 m 条按时间加入的无向边，求每对点最早连通时间：0 初始连通，m+\allowbreak{}1 表示始终不连通。 每轮把查询按 mid 分桶，清空 DSU 后扫所有前缀；整体二分也适用于其他单调前缀判定。 O(\allowbreak{}(\allowbreak{}n+\allowbreak{}m+\allowbreak{}q)\allowbreak{} log(\allowbreak{}m+\allowbreak{}2)\allowbreak{} * alpha(\allowbreak{}n)\allowbreak{})\allowbreak{}，空间 O(\allowbreak{}n+\allowbreak{}m+\allowbreak{}q)\allowbreak{}；边/\allowbreak{}查询向量从 0 起，点 1.\allowbreak{}.\allowbreak{}n。}
\begin{lstlisting}
int n;
vector<pair<int,int>> e,query;
vector<int> first_connected()
{
    int m=e.size(),q=query.size();
    vector<int> l(q),r(q,m+1),fa(n+1),sz(n+1);
    auto find=[&](int x){while(x!=fa[x])x=fa[x]=fa[fa[x]];return x;};
    while(1)
    {
        vector<vector<int>> bucket(m+1);bool pending=false;
        for(int i=0;i<q;++i)if(l[i]<r[i])bucket[l[i]+(r[i]-l[i])/2].push_back(i),pending=true;
        if(!pending)return l;
        iota(fa.begin(),fa.end(),0),fill(sz.begin(),sz.end(),1);
        for(int t=0;t<=m;++t)
        {
            if(t)
            {
                auto [u,v]=e[t-1];u=find(u),v=find(v);
                if(u!=v){if(sz[u]<sz[v])swap(u,v);fa[v]=u,sz[u]+=sz[v];}
            }
            for(int i:bucket[t])
                if(find(query[i].first)==find(query[i].second))r[i]=t;
                else l[i]=t+1;
        }
    }
}
\end{lstlisting}
\note{换成 BIT/\allowbreak{}AC fail 树统计时，只替换重置、插入和 check；check 必须随前缀长度单调。 此处二分“时间/\allowbreak{}前缀”，不使用第 k 小模板里的“去右边减 k”；允许删除会破坏本例的单调性。}
% SOURCE: 02-数据结构/可回滚并查集与时间线段树.cpp
\subsection{可回滚并查集与时间线段树}

\note{n为点数，编号1.\allowbreak{}.\allowbreak{}n；ops按时间存\{type,\allowbreak{}u,\allowbreak{}v\}，0加边、1删边、2问连通。返回查询顺序的0/\allowbreak{}1；每次重建时间线段树、活动边记录和回滚并查集。RollbackDSU保留独立实例及快照接口。}
\note{按大小合并，禁止路径压缩；snapshot 记栈长，rollback 撤回到该时刻。 find 最坏 O(\allowbreak{}log n)\allowbreak{}，一次撤销 O(\allowbreak{}1)\allowbreak{}；无需为没有实际合并的操作入栈。}
\begin{lstlisting}
struct RollbackDSU
{
    vector<int> fa,sz;
    vector<pair<int,int>> history;// 被挂的根、挂接前父根大小
    RollbackDSU(int n):fa(n+1),sz(n+1,1){iota(fa.begin(),fa.end(),0);}
    int find(int x){while(fa[x]!=x)x=fa[x];return x;}
    int snapshot(){return history.size();}
    bool merge(int x,int y)
    {
        x=find(x),y=find(y);
        if(x==y)return false;
        if(sz[x]<sz[y])swap(x,y);
        history.push_back({y,sz[x]});
        fa[y]=x,sz[x]+=sz[y];
        return true;
    }
    void rollback(int t)
    {
        while((int)history.size()>t)
        {
            auto [y,old]=history.back();history.pop_back();
            sz[fa[y]]=old,fa[y]=y;
        }
    }
};
\end{lstlisting}
\note{点 1.\allowbreak{}.\allowbreak{}n，初始无边；操作 \{type,\allowbreak{}u,\allowbreak{}v\}：0 加边，1 删边，2 查询连通。 无向边允许重复加入：全部副本删完才失效；删除必须有对应的活动副本。 每条边生效区间 [\allowbreak{}l,\allowbreak{}r)\allowbreak{} 挂到时间线段树；进入结点合并，退出撤销。 q 次操作 O(\allowbreak{}q log q log n)\allowbreak{}，存边 O(\allowbreak{}q log q)\allowbreak{}，按查询出现顺序返回 0/\allowbreak{}1。}
\begin{lstlisting}
int n;
vector<array<int,3>> ops;
vector<int> dynamic_connectivity()
{
    int q=ops.size();
    if(!q)return {};
    vector<vector<pair<int,int>>> seg(4*q);
    auto put=[&](auto &&self,int p,int l,int r,int L,int R,pair<int,int> e)->void
    {
        if(R<=l||r<=L)return;
        if(L<=l&&r<=R){seg[p].push_back(e);return;}
        int m=(l+r)/2;
        self(self,p*2,l,m,L,R,e),self(self,p*2+1,m,r,L,R,e);
    };
    map<pair<int,int>,pair<int,int>> active;// 边 -> {副本数,起始时刻}
    for(int t=0;t<q;t++)
    {
        auto [type,u,v]=ops[t];
        assert(0<=type&&type<=2&&1<=u&&u<=n&&1<=v&&v<=n);
        if(type==2)continue;
        if(u>v)swap(u,v);
        auto &a=active[{u,v}];
        if(type==0){if(a.first++==0)a.second=t;}
        else
        {
            assert(a.first>0);
            if(--a.first==0)put(put,1,0,q,a.second,t,{u,v});
        }
    }
    for(auto [e,a]:active)if(a.first)put(put,1,0,q,a.second,q,e);
    RollbackDSU d(n);
    vector<int> ans;
    auto dfs=[&](auto &&self,int p,int l,int r)->void
    {
        int old=d.snapshot();
        for(auto [u,v]:seg[p])d.merge(u,v);
        if(r-l==1)
        {
            if(ops[l][0]==2)ans.push_back(d.find(ops[l][1])==d.find(ops[l][2]));
        }
        else{int m=(l+r)/2;self(self,p*2,l,m),self(self,p*2+1,m,r);}
        d.rollback(old);
    };
    dfs(dfs,1,0,q);
    return ans;
}
\end{lstlisting}
% SOURCE: 02-数据结构/笛卡尔树.cpp
\subsection{笛卡尔树}

\note{a[\allowbreak{}1.\allowbreak{}.\allowbreak{}n]\allowbreak{}为键，lc/\allowbreak{}rc/\allowbreak{}fa记录左右儿子和父亲，root为根，0为空；stk是建树单调栈。中序按原下标、堆序按a，比较符决定等值的归属；重新建树清空儿子和父亲。}
\note{笛卡尔树：O(\allowbreak{}n)\allowbreak{} 单调栈建树 性质：中序遍历就是原数组顺序；每个结点是「小根堆」（默认最小值在根）， 所以区间 [\allowbreak{}l,\allowbreak{}r]\allowbreak{} 的最小值就是 lca(\allowbreak{}l,\allowbreak{}r)\allowbreak{}，配合 O(\allowbreak{}1)\allowbreak{} LCA 就能做 RMQ 反过来：把「下标」当 BST 键、「值」当堆键，笛卡尔树唯一确定 单调栈过程：新点 x 从栈顶往下弹掉所有 val >\allowbreak{} a[\allowbreak{}x]\allowbreak{} 的点， 最后弹出来的那个点当 x 的左儿子（它右边的点都比它大，全在左边）}
\begin{lstlisting}
const int N=100005;
const int LOG=18;
int n,a[N];
int lc[N],rc[N],fa[N],stk[N];
int root;
\end{lstlisting}
\note{O(\allowbreak{}n)\allowbreak{} 建笛卡尔树（小根堆）；返回根；相等时取靠前的当根（用 >\allowbreak{} 判断）}
\begin{lstlisting}
int build_cart()
{
    int top=0,rt=0;
    for(int i=1;i<=n;i++)lc[i]=rc[i]=fa[i]=0;
    for(int i=1;i<=n;i++)
    {
        int last=0;
        while(top&&a[stk[top]]>a[i])last=stk[top--];//弹掉比 a[i] 大的
        lc[i]=last;
        if(last)fa[last]=i;
        if(top)rc[stk[top]]=i,fa[i]=stk[top];
        else fa[i]=0,rt=i;//栈空了，i 就是当前的根
        stk[++top]=i;
    }
    return rt;
}
\end{lstlisting}
% SOURCE: 02-数据结构/平衡树Treap.cpp
\subsection{平衡树Treap}

\note{t保存Treap节点，root为根，seed为随机优先级发生器状态；x是键。重复次数与子树大小决定排名和第k小；插删更新根，换集合清空节点池，随机优先级仅负责平衡。}
\note{平衡树 Treap（带旋）：插入 /\allowbreak{} 删除 /\allowbreak{} 排名 /\allowbreak{} 第k小 /\allowbreak{} 前驱后继 每个结点随机一个优先值，小根堆性质 → 期望 O(\allowbreak{}log n)\allowbreak{}，最坏 O(\allowbreak{}n)\allowbreak{} 重复插入忽略（集合），排名约定：get\_rank(\allowbreak{}x)\allowbreak{} 返回「比 x 小的数个数」（0 起） 关键坑：旋转必须先用 int \&q=\allowbreak{}ls[\allowbreak{}p]\allowbreak{} 绑好引用，递归返回后 ls[\allowbreak{}p]\allowbreak{} 可能已经不是原来那个结点了}
\begin{lstlisting}
const int N=200005;
const int INF=0x3f3f3f3f;
int n,x;
\end{lstlisting}
\note{xorshift 伪随机，比 rand(\allowbreak{})\allowbreak{} 快}
\begin{lstlisting}
unsigned long long seed=20240513;
int rand_int()
{
    seed^=seed<<7;
    seed^=seed>>9;
    return (int)(seed%1000000007);
}
struct Treap{
    int val[N],pri[N],sz[N],ls[N],rs[N];   // 0 号是空结点
    int tot;
    void clear(){tot=sz[0]=ls[0]=rs[0]=0;}
    int new_node(int x)
    {
        tot++;
        val[tot]=x,pri[tot]=rand_int(),sz[tot]=1,ls[tot]=0,rs[tot]=0;
        return tot;
    }
    void push_up(int p){sz[p]=sz[ls[p]]+sz[rs[p]]+1;}
    // 左旋：p 的右儿子转上来
    void zig(int &p)
    {
        int q=rs[p];
        rs[p]=ls[q],ls[q]=p,p=q;
        push_up(ls[p]),push_up(p);
    }
    // 右旋：p 的左儿子转上来
    void zag(int &p)
    {
        int q=ls[p];
        ls[p]=rs[q],rs[q]=p,p=q;
        push_up(rs[p]),push_up(p);
    }
    void insert(int &p,int x)
    {
        if(p==0){p=new_node(x);return;}
        if(x<val[p])
        {
            int &q=ls[p];                  // 先绑引用，旋转后 ls[p] 会变
            insert(q,x);
            if(pri[q]<pri[p])zag(p);
        }
        else if(x>val[p])
        {
            int &q=rs[p];
            insert(q,x);
            if(pri[q]<pri[p])zig(p);
        }
        else return;                       // 已有相同值（多重集合应另加出现次数，不要只改比较号）
        push_up(p);
    }
    // 把 p 一路旋转到叶子再摘掉
    void erase(int &p,int x)
    {
        if(p==0)return;
        if(x<val[p])erase(ls[p],x);
        else if(x>val[p])erase(rs[p],x);
        else
        {
            if(ls[p]==0)p=rs[p];
            else if(rs[p]==0)p=ls[p];
            else if(pri[ls[p]]<pri[rs[p]])zag(p),erase(rs[p],x);
            else zig(p),erase(ls[p],x);
        }
        if(p)push_up(p);
    }
    bool find(int p,int x)
    {
        while(p)
        {
            if(x<val[p])p=ls[p];
            else if(x>val[p])p=rs[p];
            else return true;
        }
        return false;
    }
    // 比 x 小的数个数（0 起）
    int get_rank(int p,int x)
    {
        int cnt=0;
        while(p)
        {
            if(x<=val[p])p=ls[p];
            else cnt+=sz[ls[p]]+1,p=rs[p];
        }
        return cnt;
    }
    // 第 k 小（k 从 1 开始），不存在返回 0
    int kth(int p,int k)
    {
        if(k<1||k>sz[p])return 0;
        while(p)
        {
            if(k<=sz[ls[p]])p=ls[p];
            else if(k==sz[ls[p]]+1)return val[p];
            else k-=sz[ls[p]]+1,p=rs[p];
        }
        return 0;
    }
    int get_pre(int p,int x)               // 小于 x 的最大值
    {
        int ans=-INF;
        while(p)
        {
            if(val[p]<x)ans=val[p],p=rs[p];
            else p=ls[p];
        }
        return ans;
    }
    int get_next(int p,int x)              // 大于 x 的最小值
    {
        int ans=INF;
        while(p)
        {
            if(val[p]>x)ans=val[p],p=ls[p];
            else p=rs[p];
        }
        return ans;
    }
    int get_min(int p)
    {
        while(ls[p])p=ls[p];
        return val[p];
    }
    int get_max(int p)
    {
        while(rs[p])p=rs[p];
        return val[p];
    }
    void print_inorder(int p)              // 中序，调试用
    {
        if(!p)return;
        print_inorder(ls[p]);
        printf(" %d",val[p]);
        print_inorder(rs[p]);
    }
};
Treap t;
int root;
\end{lstlisting}
% SOURCE: 02-数据结构/平衡树Splay.cpp
\subsection{平衡树Splay}

\note{sp为Splay对象，节点维护键、重复次数、子树大小、父亲和两个儿子；根在对象内。排名与第k小按重复数计，哨兵参与排名时按接口扣除；换集合重置对象，删除后同步维护大小。}
\note{平衡树 Splay（伸展树）：插入 /\allowbreak{} 删除 /\allowbreak{} 排名 /\allowbreak{} 第k小 /\allowbreak{} 前驱后继 /\allowbreak{} 区间翻转 单次操作均摊 O(\allowbreak{}log n)\allowbreak{}，常数比 Treap 大，但能支持区间操作（文艺平衡树） 本模板每个结点存一个互不相同的值 +\allowbreak{} 出现次数 num，所以同一份代码既能当普通平衡树， 也能当文艺平衡树（此时把 val 看成原序列元素、num 全为 1，忽略 val 的 BST 性质） 约定：get\_rank(\allowbreak{}x)\allowbreak{} 返回严格小于 x 的个数（0 起），kth 的 k 从 1 开始 文艺平衡树用法：init(\allowbreak{}n)\allowbreak{} 建出 1.\allowbreak{}.\allowbreak{}n 并带两个哨兵，位置 i 对应中序第 i+\allowbreak{}1 个结点 坑 1：splay/\allowbreak{}查找/\allowbreak{}旋转前都要 push\_down，否则 rev 标记会让左右儿子对不上 坑 2：序列模式哨兵也计入 sz，位置查询含哨兵；输出时跳过哨兵}
\begin{lstlisting}
const int N=200005;
const int INF=0x3f3f3f3f;
int n,m,x;
struct Splay{
    int val[N],num[N],sz[N],fa[N],ch[N][2];
    bool rev[N];
    int tot,root;
    void clear()
    {
        tot=0,root=0;
        for(int i=0;i<2;i++)ch[0][i]=0;
        fa[0]=sz[0]=num[0]=0,rev[0]=0,val[0]=0;
    }
    int new_node(int v,int f)
    {
        tot++;
        val[tot]=v,num[tot]=1,sz[tot]=1,fa[tot]=f,rev[tot]=0;
        ch[tot][0]=ch[tot][1]=0;
        return tot;
    }
    int get(int p){return ch[fa[p]][1]==p;}          // p 是父结点的右儿子吗
    void push_up(int p){sz[p]=sz[ch[p][0]]+sz[ch[p][1]]+num[p];}
    void push_down(int p)
    {
        if(!p||!rev[p])return;
        swap(ch[p][0],ch[p][1]);                     // 翻转：先换儿子，再下传标记
        if(ch[p][0])rev[ch[p][0]]^=1;
        if(ch[p][1])rev[ch[p][1]]^=1;
        rev[p]=0;
    }
    // 单旋：p 往上走一层
    void rotate(int p)
    {
        int f=fa[p],g=fa[f],k=get(p),w=ch[p][k^1];
        push_down(f),push_down(p);                   // 旋转前先把标记放下去
        ch[f][k]=w;
        if(w)fa[w]=f;
        ch[p][k^1]=f,fa[f]=p,fa[p]=g;
        if(g)ch[g][ch[g][1]==f]=p;
        push_up(f),push_up(p);
    }
    // 把 p 旋到 goal 的儿子处，goal=0 表示旋到根
    void splay(int p,int goal=0)
    {
        vector<int> path;
        for(int u=p;u;u=fa[u])path.push_back(u);
        for(int i=(int)path.size()-1;i>=0;i--)push_down(path[i]);
        while(fa[p]!=goal)
        {
            int f=fa[p],g=fa[f];
            if(g!=goal)
            {
                if(get(p)==get(f))rotate(f);         // 同向：先转父亲（zig-zig）
                else rotate(p);                      // 异向：先转自己（zig-zag）
            }
            rotate(p);
        }
        if(!goal)root=p;
    }
    // 插入一个值 v
    void insert(int v)
    {
        if(!root){root=new_node(v,0);return;}
        int p=root,f=0;
        while(p)
        {
            push_down(p);
            if(val[p]==v){num[p]++,push_up(p),splay(p);return;}
            f=p,p=ch[p][v>val[p]];
        }
        p=new_node(v,f),ch[f][v>val[f]]=p,splay(p);
    }
    // 删除一个值 v（删掉一次出现）
    void erase(int v)
    {
        int p=find_node(v);
        if(!p)return;                                // 不存在
        splay(p);
        if(num[p]>1){num[p]--,push_up(p);return;}
        int l=ch[p][0],r=ch[p][1];
        if(!l&&!r){root=0;return;}
        if(!l){fa[r]=0,ch[p][1]=0,root=r;return;}
        if(!r){fa[l]=0,ch[p][0]=0,root=l;return;}
        // 关键：左右子树的根要同时从 p 上摘下来（fa 和 ch 都清），
        // 否则 splay(q) 旋转时会顺着没摘干净的那条边把 r 又转进来，结点就丢了
        fa[l]=fa[r]=0,ch[p][0]=ch[p][1]=0;
        int q=l;
        while(ch[q][1])push_down(q),q=ch[q][1];
        splay(q);
        ch[q][1]=r,fa[r]=q,push_up(q);
    }
    // 找值为 v 的结点，没有返回 0
    int find_node(int v)
    {
        int p=root,last=0;
        while(p)
        {
            push_down(p); last=p;
            if(val[p]==v){splay(p); return p;}
            p=ch[p][v>val[p]];
        }
        if(last)splay(last);
        return 0;
    }
    // 严格小于 v 的个数（0 起）
    int get_rank(int v)
    {
        int p=root,last=0,ans=0;
        while(p)
        {
            push_down(p); last=p;
            if(v<=val[p])p=ch[p][0];
            else ans+=sz[ch[p][0]]+num[p],p=ch[p][1];
        }
        if(last)splay(last);
        return ans;
    }
    // 第 k 小（k 从 1 开始），不存在返回 0
    int kth(int k)
    {
        if(k<1||k>sz[root])return 0;
        int p=root;
        while(p)
        {
            push_down(p);
            if(k<=sz[ch[p][0]])p=ch[p][0];
            else if(k<=sz[ch[p][0]]+num[p]){int v=val[p]; splay(p); return v;}
            else k-=sz[ch[p][0]]+num[p],p=ch[p][1];
        }
        return 0;
    }
    int get_pre(int v)                               // 严格小于 v 的最大值
    {
        int p=root,last=0,ans=-INF;
        while(p)
        {
            push_down(p); last=p;
            if(val[p]<v)ans=val[p],p=ch[p][1];
            else p=ch[p][0];
        }
        if(last)splay(last);
        return ans;
    }
    int get_next(int v)                              // 严格大于 v 的最小值
    {
        int p=root,last=0,ans=INF;
        while(p)
        {
            push_down(p); last=p;
            if(val[p]>v)ans=val[p],p=ch[p][0];
            else p=ch[p][1];
        }
        if(last)splay(last);
        return ans;
    }
    // ---------- 以下是文艺平衡树 ----------
    // 直接按中序建一棵尽量平衡的树：区间 [l,r] 里放的是值 v..v+r-l
    int build(int l,int r,int v)
    {
        if(l>r)return 0;
        int mid=(l+r)>>1;
        int p=new_node(v+mid-l,0);
        ch[p][0]=build(l,mid-1,v);
        ch[p][1]=build(mid+1,r,v+mid-l+1);
        if(ch[p][0])fa[ch[p][0]]=p;
        if(ch[p][1])fa[ch[p][1]]=p;
        push_up(p);
        return p;
    }
    // 建出序列 1..n_，两个哨兵也占一个位置
    void init(int n_)
    {
        clear();
        root=build(0,n_+1,0);
        val[get_pos(1)]=-INF;
        val[get_pos(n_+2)]=INF;
    }
    // 中序第 k 个结点（1 起，含哨兵），一路 push_down
    int get_pos(int k)
    {
        int p=root;
        while(p)
        {
            push_down(p);
            if(k<=sz[ch[p][0]])p=ch[p][0];
            else if(k<=sz[ch[p][0]]+num[p])return p;
            else k-=sz[ch[p][0]]+num[p],p=ch[p][1];
        }
        return 0;
    }
    // 区间翻转 [l,r]（原序列下标，1 起）
    // 左哨兵占第 1 个中序位置，所以「l-1 位置」是第 l 个，「r+1 位置」是第 r+2 个；
    // 把这两个旋上来之后，夹在中间的整棵子树正好是 [l,r]
    void reverse(int l,int r)
    {
        int a=get_pos(l),b=get_pos(r+2);
        if(!a||!b)return;
        splay(a),splay(b,a);
        int t=ch[b][0];
        if(t)rev[t]^=1;
    }
    // 中序打印真实元素，调试用（结点记了出现次数，这里每个值只打一次）
    void print(int p)
    {
        if(!p)return;
        push_down(p);
        print(ch[p][0]);
        if(val[p]!=-INF&&val[p]!=INF)printf(" %d",val[p]);
        print(ch[p][1]);
    }
};
Splay sp;
\end{lstlisting}
% SOURCE: 02-数据结构/李超树.cpp
\subsection{李超树}

\note{Line\{a,\allowbreak{}b,\allowbreak{}c\}表示(\allowbreak{}a*x+\allowbreak{}b)\allowbreak{}/\allowbreak{}c，c>\allowbreak{}0，c=\allowbreak{}0为空直线；对象保存节点和查询域。better用交叉乘积比较；连续整数域和离散坐标按各入口选择，域外查询不可套用。}
\begin{lstlisting}
typedef __int128 lll;
struct Line
{
    ll a,b,c;
    Line(ll a_=0,ll b_=0,ll c_=0):a(a_),b(b_),c(c_){}
};
\end{lstlisting}
\note{分母为正；交叉乘积必须在 \_\_int128 范围内}
\begin{lstlisting}
bool better(const Line &a,const Line &b,ll x)
{
    if(!a.c)return false;
    if(!b.c)return true;
    return ((lll)a.a*x+a.b)*b.c>((lll)b.a*x+b.b)*a.c;
}
\end{lstlisting}
\note{动态开点，连续整数域与离散坐标共用；每次操作 O(\allowbreak{}log V)\allowbreak{}}
\begin{lstlisting}
struct LiChao
{
    struct Node
    {
        int lc=0,rc=0;
        Line s;
    };
    vector<Node> tr;
    vector<ll> xs;
    ll lo,hi;
    void init(ll l,ll r)
    {
        assert(l<=r);
        lo=l,hi=r,xs.clear(),tr.assign(2,Node());
    }
    void init(vector<ll> v)
    {
        assert(!v.empty());
        sort(v.begin(),v.end());
        v.erase(unique(v.begin(),v.end()),v.end());
        init(0,(ll)v.size()-1);
        xs=move(v);
    }
    ll coord(ll x)
    {
        return xs.empty()?x:xs[x];
    }
    void insert(Line s)
    {
        assert(s.c>0);
        insert(s,lo,hi,1);
    }
    void insert(Line s,ll l,ll r,int p)
    {
        if(!tr[p].s.c)
        {
            tr[p].s=s;
            return;
        }
        ll mid=(ll)((lll)l+((lll)r-l)/2);
        bool bl=better(s,tr[p].s,coord(l)),bm=better(s,tr[p].s,coord(mid));
        if(bm)swap(s,tr[p].s);
        if(l==r)return;
        bool left=bl!=bm;
        int q=left?tr[p].lc:tr[p].rc;
        if(!q)
        {
            q=(int)tr.size();
            tr.push_back(Node());
            if(left)tr[p].lc=q;
            else tr[p].rc=q;
        }
        if(left)insert(s,l,mid,q);
        else insert(s,mid+1,r,q);
    }
    // 返回最优直线，避免把分数截断成整数
    Line query(ll x)
    {
        ll id=x;
        if(!xs.empty())
        {
            id=lower_bound(xs.begin(),xs.end(),x)-xs.begin();
            assert(id<(ll)xs.size()&&xs[id]==x);
        }
        assert(lo<=id&&id<=hi);
        ll l=lo,r=hi;
        int p=1;
        Line ans;
        while(p)
        {
            if(better(tr[p].s,ans,x))ans=tr[p].s;
            if(l==r)break;
            ll mid=(ll)((lll)l+((lll)r-l)/2);
            if(id<=mid)p=tr[p].lc,r=mid;
            else p=tr[p].rc,l=mid+1;
        }
        return ans;
    }
};
\end{lstlisting}
% SOURCE: 02-数据结构/左偏树(可并堆).cpp
\subsection{左偏树(可并堆)}

\note{节点保存键、左右儿子和空路径长，0为空根；merge返回合并根并消耗两棵旧堆。pop返回删顶后的新根，top/\allowbreak{}pop要求非空；旧根不能作为独立堆继续操作。}
\note{小根可并堆：root 是节点编号，0 表示空；merge 后旧根不能再当独立堆使用。 两个待合并堆必须没有共享节点；相同权值允许，节点删除后不回收。 merge 沿右链递归，右子树空路径长不大于左边；合并/\allowbreak{}插入/\allowbreak{}删顶 O(\allowbreak{}log n)\allowbreak{}。}
\begin{lstlisting}
struct LeftistHeap
{
    struct Node{ll val=0;int l=0,r=0,d=0;};
    vector<Node> tr=vector<Node>(1);
    void clear(){tr.assign(1,Node{});}
    int merge(int x,int y)
    {
        if(!x||!y)return x|y;
        assert(x!=y);
        if(tr[x].val>tr[y].val)swap(x,y);
        tr[x].r=merge(tr[x].r,y);
        if(tr[tr[x].l].d<tr[tr[x].r].d)swap(tr[x].l,tr[x].r);
        tr[x].d=tr[tr[x].r].d+1;
        return x;
    }
    int insert(int root,ll x)
    {
        int p=tr.size();
        tr.push_back({x,0,0,1});
        return merge(root,p);
    }
    ll top(int root){assert(root);return tr[root].val;}
    int pop(int root)
    {
        assert(root);
        int p=merge(tr[root].l,tr[root].r);
        tr[root].l=tr[root].r=0,tr[root].d=1;
        return p;
    }
};
\end{lstlisting}
% SOURCE: 02-数据结构/SegmentTreeBeats.cpp
\subsection{SegmentTreeBeats}

\note{Beats对象保存t：sum、最大/\allowbreak{}次大值及次数、最小/\allowbreak{}次小值及次数、加法标记。输入0起，操作闭区间；update的op按实现选择加、chmin、chmax。只有阈值越过次极值的条件满足时才整段修改，否则递归。}
\note{区间 chmin/\allowbreak{}chmax/\allowbreak{}add/\allowbreak{}sum，0-\allowbreak{}indexed 闭区间；同时维护严格次大/\allowbreak{}次小、极值个数。 chmin 只有 max2<\allowbreak{}x<\allowbreak{}max1 才能整段更新；否则递归。chmax 对称；加法移动全部极值。 含 add 的通用 Beats 标准摊还界 O(\allowbreak{}(\allowbreak{}n+\allowbreak{}q)\allowbreak{}log\ensuremath{^2}n)\allowbreak{}，空间 O(\allowbreak{}n)\allowbreak{}。值/\allowbreak{}和远离 ±INF。}
\begin{lstlisting}
struct Beats
{
    static constexpr ll INF=LLONG_MAX/4;
    struct Node{ll sum=0,mx=-INF,mx2=-INF,mn=INF,mn2=INF,add=0;int cmx=0,cmn=0;};
    int n;vector<Node> t;
    Beats(const vector<ll> &a):n(a.size()),t(max(1,4*n)){if(n)build(1,0,n-1,a);}
    void pull(int p)
    {
        auto &a=t[p*2],&b=t[p*2+1],&c=t[p];c.sum=a.sum+b.sum;
        c.mx=max(a.mx,b.mx);c.cmx=(a.mx==c.mx?a.cmx:0)+(b.mx==c.mx?b.cmx:0);
        c.mx2=max(a.mx==c.mx?a.mx2:a.mx,b.mx==c.mx?b.mx2:b.mx);
        c.mn=min(a.mn,b.mn);c.cmn=(a.mn==c.mn?a.cmn:0)+(b.mn==c.mn?b.cmn:0);
        c.mn2=min(a.mn==c.mn?a.mn2:a.mn,b.mn==c.mn?b.mn2:b.mn);
    }
    void build(int p,int l,int r,const vector<ll> &a){if(l==r){t[p].sum=t[p].mn=t[p].mx=a[l];t[p].cmn=t[p].cmx=1;return;}int m=(l+r)/2;build(p*2,l,m,a);build(p*2+1,m+1,r,a);pull(p);}
    void plus(int p,int len,ll x){auto &a=t[p];a.sum+=len*x;a.mx+=x;a.mn+=x;if(a.mx2!=-INF)a.mx2+=x;if(a.mn2!=INF)a.mn2+=x;a.add+=x;}
    void lower(int p,ll x){auto &a=t[p];if(a.mx<=x)return;a.sum+=(x-a.mx)*a.cmx;if(a.mn==a.mx)a.mn=x;else if(a.mn2==a.mx)a.mn2=x;a.mx=x;}
    void upper(int p,ll x){auto &a=t[p];if(a.mn>=x)return;a.sum+=(x-a.mn)*a.cmn;if(a.mx==a.mn)a.mx=x;else if(a.mx2==a.mn)a.mx2=x;a.mn=x;}
    void push(int p,int l,int r)
    {
        int m=(l+r)/2;if(t[p].add){plus(p*2,m-l+1,t[p].add);plus(p*2+1,r-m,t[p].add);t[p].add=0;}
        for(int c:{p*2,p*2+1}){lower(c,t[p].mx);upper(c,t[p].mn);}
    }
    // op=0 加 x，op=1 chmin(x)，op=2 chmax(x)。调用区间必须在 [0,n)。
    void update(int p,int l,int r,int L,int R,ll x,int op)
    {
        if(R<l||r<L||(op==1&&t[p].mx<=x)||(op==2&&t[p].mn>=x))return;
        if(L<=l&&r<=R)
        {
            if(op==0){plus(p,r-l+1,x);return;}
            if(op==1&&t[p].mx2<x){lower(p,x);return;}
            if(op==2&&t[p].mn2>x){upper(p,x);return;}
        }
        push(p,l,r);int m=(l+r)/2;update(p*2,l,m,L,R,x,op);update(p*2+1,m+1,r,L,R,x,op);pull(p);
    }
    ll query(int p,int l,int r,int L,int R){if(R<l||r<L)return 0;if(L<=l&&r<=R)return t[p].sum;push(p,l,r);int m=(l+r)/2;return query(p*2,l,m,L,R)+query(p*2+1,m+1,r,L,R);}
    void update(int l,int r,ll x,int op){assert(0<=l&&l<=r&&r<n);update(1,0,n-1,l,r,x,op);}
    ll sum(int l,int r){assert(0<=l&&l<=r&&r<n);return query(1,0,n-1,l,r);}
};
\end{lstlisting}
% SOURCE: 02-数据结构/线段树(区间开方).cpp
\subsection{线段树(区间开方)}

\note{a[\allowbreak{}1.\allowbreak{}.\allowbreak{}n]\allowbreak{}须非负；seg维护区间和与最大值，最大值<\allowbreak{}=\allowbreak{}1的区间不再变化。br是对照缓冲，操作为整数平方根；大整数不能只依赖double开方后直接截断，应保留实现的校正。}
\begin{lstlisting}
const int N=100005;
\end{lstlisting}
\note{线段树：区间开方(\allowbreak{}向下取整)\allowbreak{} +\allowbreak{} 区间和 非负整数；64 位整数最多开方 7 次变成 0/\allowbreak{}1，用 flag 剪枝。区间和须不溢出。 均摊复杂度 O(\allowbreak{}(\allowbreak{}n+\allowbreak{}m)\allowbreak{} log n)\allowbreak{}}
\begin{lstlisting}
template<typename T,int NMAX=N>
struct SegmentTree
{
    T tr[NMAX*4];
    bool flag[NMAX*4];//该区间是否全 <=1，是就不用再递归
    T *src;//建树用的原数组(1-indexed)
    #define lp (p<<1)
    #define rp ((p<<1)|1)
    #define mid ((l+r)>>1)
    void build(int l,int r,int p)
    {
        if(l==r)
        {
            tr[p]=src[l],flag[p]=(tr[p]<=1);
            return;
        }
        build(l,mid,lp);
        build(mid+1,r,rp);
        tr[p]=tr[lp]+tr[rp],flag[p]=flag[lp]&&flag[rp];
    }
    void update(int L,int R,int l,int r,int p)//[L,R] 每个数开方
    {
        if(R<l||r<L)return;//无交集
        if(flag[p])return;//全 0/1，开方不变
        if(l==r)
        {
            T v=tr[p],x=(T)sqrtl((long double)v);
            while((__int128)x*x>v)--x;
            while((__int128)(x+1)*(x+1)<=v)++x;
            tr[p]=x,flag[p]=(x<=1);
            return;
        }
        update(L,R,l,mid,lp);
        update(L,R,mid+1,r,rp);
        tr[p]=tr[lp]+tr[rp],flag[p]=flag[lp]&&flag[rp];
    }
    T query(int L,int R,int l,int r,int p)
    {
        if(L<=l&&r<=R)return tr[p];
        T res=0;
        if(L<=mid)res+=query(L,R,l,mid,lp);
        if(R>mid)res+=query(L,R,mid+1,r,rp);
        return res;
    }
};
int n;
ll a[205],br[205];
SegmentTree<ll,205> seg;
\end{lstlisting}
% SOURCE: 02-数据结构/树状数组套权值线段树.cpp
\subsection{树状数组套权值线段树}

\note{对象的root[\allowbreak{}i]\allowbreak{}为第i个BIT节点的权值树根，节点池存子计数。位置0起、值已离散化到0.\allowbreak{}.\allowbreak{}sigma-\allowbreak{}1，kth(\allowbreak{}l,\allowbreak{}r,\allowbreak{}k)\allowbreak{}查询闭区间第k小；change先删旧值再加新值，调用者负责维护原数组。}
\note{在线点修改、区间第 k 小：位置 0.\allowbreak{}.\allowbreak{}n-\allowbreak{}1，值域离散化为 0.\allowbreak{}.\allowbreak{}\ensuremath{\sigma}-\allowbreak{}1，k 从 1 开始。 所有可能修改值预先加入压缩表；BIT 套动态权值线段树，O(\allowbreak{}log n log \ensuremath{\sigma})\allowbreak{} 每次。 roots 中存频数；change(\allowbreak{}pos,\allowbreak{}old,\allowbreak{}new)\allowbreak{} 前 old 必须确实是当前位置值。内存 O(\allowbreak{}(\allowbreak{}n+\allowbreak{}q)\allowbreak{}log n log \ensuremath{\sigma})\allowbreak{} 上界。}
\begin{lstlisting}
struct BITValueTree
{
    struct Node{int l=0,r=0,sum=0;};int n,sigma;vector<int> root;vector<Node> t={{}};
    BITValueTree(int n,int sigma):n(n),sigma(sigma),root(n+1){assert(sigma>0);}
    int add(int p,int l,int r,int x,int d)
    {
        if(!p){p=t.size();t.push_back({});}t[p].sum+=d;
        if(l<r){int m=(l+r)/2;if(x<=m){int z=add(t[p].l,l,m,x,d);t[p].l=z;}else{int z=add(t[p].r,m+1,r,x,d);t[p].r=z;}}
        return p;
    }
    void add(int pos,int value,int delta){assert(0<=pos&&pos<n&&0<=value&&value<sigma);for(int i=pos+1;i<=n;i+=i&-i)root[i]=add(root[i],0,sigma-1,value,delta);}
    void change(int pos,int old,int value){add(pos,old,-1);add(pos,value,1);}
    int kth(int l,int r,int k)const// 闭区间
    {
        assert(0<=l&&l<=r&&r<n&&1<=k&&k<=r-l+1);vector<int> a,b;
        for(int i=r+1;i;i-=i&-i)a.push_back(root[i]);for(int i=l;i;i-=i&-i)b.push_back(root[i]);
        int L=0,R=sigma-1;
        while(L<R)
        {int count=0,m=(L+R)/2;for(int p:a)count+=t[t[p].l].sum;for(int p:b)count-=t[t[p].l].sum;
            bool left=k<=count;if(!left)k-=count;for(int &p:a)p=left?t[p].l:t[p].r;for(int &p:b)p=left?t[p].l:t[p].r;if(left)R=m;else L=m+1;}
        return L;
    }
};
\end{lstlisting}
% SOURCE: 02-数据结构/可持久化Treap(fhq).cpp
\subsection{可持久化Treap(fhq)}

\note{节点池保存键、优先级、儿子和大小，根编号区分各版本；修改路径复制节点，旧版本不变。split/\allowbreak{}merge返回不同根，不能把非持久化的原地修改混入，容量按总操作次数估算。}
\note{隐式无旋 Treap，路径复制；版本根永不原地修改}
\begin{lstlisting}
struct Treap
{
    struct Node
    {
        int lc=0,rc=0,sz=0;
        unsigned key=0;
        ll val=0,sum=0,add=0;
        bool rev=false;
    };
    vector<Node> tr={Node()};
    mt19937 rng{19260817};
    int clone(int p)
    {
        if(!p)return 0;
        Node z=tr[p];
        tr.push_back(z);
        return (int)tr.size()-1;
    }
    int new_node(ll v)
    {
        Node z;
        z.sz=1,z.val=z.sum=v,z.key=rng();
        tr.push_back(z);
        return (int)tr.size()-1;
    }
    void push_up(int p)
    {
        int l=tr[p].lc,r=tr[p].rc;
        tr[p].sz=tr[l].sz+tr[r].sz+1;
        tr[p].sum=tr[l].sum+tr[r].sum+tr[p].val;
    }
    void set_add(int p,ll v)
    {
        if(p)tr[p].val+=v,tr[p].sum+=v*tr[p].sz,tr[p].add+=v;
    }
    void set_rev(int p)
    {
        if(p)swap(tr[p].lc,tr[p].rc),tr[p].rev=!tr[p].rev;
    }
    // p 已复制，带标记的儿子必须先复制再下推
    void push_down(int p)
    {
        if(!tr[p].add&&!tr[p].rev)return;
        int l=clone(tr[p].lc),r=clone(tr[p].rc);
        tr[p].lc=l,tr[p].rc=r;
        set_add(l,tr[p].add),set_add(r,tr[p].add);
        if(tr[p].rev)set_rev(l),set_rev(r);
        tr[p].add=0,tr[p].rev=false;
    }
    // 按前 k 个元素分裂，允许 k=0 或整段长度
    pair<int,int> split(int p,int k)
    {
        if(!p)return {0,0};
        if(k<=0)return {0,p};
        if(k>=tr[p].sz)return {p,0};
        p=clone(p),push_down(p);
        int s=tr[tr[p].lc].sz;
        if(k<=s)
        {
            pair<int,int> z=split(tr[p].lc,k);
            tr[p].lc=z.second,push_up(p);
            return {z.first,p};
        }
        pair<int,int> z=split(tr[p].rc,k-s-1);
        tr[p].rc=z.first,push_up(p);
        return {p,z.second};
    }
    int merge(int a,int b)
    {
        if(!a||!b)return a?a:b;
        if(tr[a].key<tr[b].key)
        {
            a=clone(a),push_down(a);
            tr[a].rc=merge(tr[a].rc,b),push_up(a);
            return a;
        }
        b=clone(b),push_down(b);
        tr[b].lc=merge(a,tr[b].lc),push_up(b);
        return b;
    }
    // 在前 k 个元素之后插入
    int insert(int p,int k,ll v)
    {
        assert(0<=k&&k<=tr[p].sz);
        pair<int,int> z=split(p,k);
        return merge(merge(z.first,new_node(v)),z.second);
    }
    int erase(int p,int k)
    {
        assert(1<=k&&k<=tr[p].sz);
        pair<int,int> z=split(p,k-1),w=split(z.second,1);
        return merge(z.first,w.second);
    }
    int update(int p,int l,int r,ll v,bool rev)
    {
        assert(1<=l&&l<=r&&r<=tr[p].sz);
        pair<int,int> z=split(p,l-1),w=split(z.second,r-l+1);
        int b=clone(w.first);
        set_add(b,v);
        if(rev)set_rev(b);
        return merge(merge(z.first,b),w.second);
    }
    // 只读查询，携带祖先标记，不产生新节点
    ll query(int p,int l,int r,ll add=0,bool rev=false)
    {
        if(!p||l>r)return 0;
        if(l==1&&r==tr[p].sz)return tr[p].sum+add*tr[p].sz;
        int a=rev?tr[p].rc:tr[p].lc,b=rev?tr[p].lc:tr[p].rc,s=tr[a].sz;
        ll ans=0,v=add+tr[p].add;
        bool f=rev^tr[p].rev;
        if(l<=s)ans+=query(a,l,min(r,s),v,f);
        if(l<=s+1&&s+1<=r)ans+=tr[p].val+add;
        if(r>s+1)ans+=query(b,max(1,l-s-1),r-s-1,v,f);
        return ans;
    }
    ll range_sum(int p,int l,int r)
    {
        assert(1<=l&&l<=r&&r<=tr[p].sz);
        return query(p,l,r);
    }
};
\end{lstlisting}
% SOURCE: 02-数据结构/链接剖分LCT.cpp
\subsection{链接剖分LCT}

\note{tr保存辅助Splay的儿子、父亲、翻转标记、点权与路径和；0为空点，下标1起。link/\allowbreak{}cut操作森林边，split暴露路径；访问前先下传翻转，不能把辅助父子关系当原树父子。}
\note{动态森林：维护点权路径和，操作均摊 O(\allowbreak{}log n)\allowbreak{}，下标从 1 开始}
\begin{lstlisting}
struct LCT
{
    struct Node { int ch[2]={0,0},fa=0; ll v=0,sum=0; bool rev=false; };
    vector<Node> tr;
    LCT(int n):tr(n+1) {}
    bool is_root(int x)
    {
        int f=tr[x].fa;
        return tr[f].ch[0]!=x&&tr[f].ch[1]!=x;
    }
    void pull(int x) { tr[x].sum=tr[tr[x].ch[0]].sum+tr[tr[x].ch[1]].sum+tr[x].v; }
    void reverse_node(int x)
    {
        if(x)swap(tr[x].ch[0],tr[x].ch[1]),tr[x].rev=!tr[x].rev;
    }
    void push(int x)
    {
        if(tr[x].rev)reverse_node(tr[x].ch[0]),reverse_node(tr[x].ch[1]),tr[x].rev=false;
    }
    void rotate(int x)
    {
        int y=tr[x].fa,z=tr[y].fa,k=tr[y].ch[1]==x,w=tr[x].ch[k^1];
        if(!is_root(y))tr[z].ch[tr[z].ch[1]==y]=x;
        tr[x].fa=z,tr[x].ch[k^1]=y,tr[y].fa=x,tr[y].ch[k]=w;
        if(w)tr[w].fa=y;
        pull(y),pull(x);
    }
    void splay(int x)
    {
        vector<int> st;
        for(int y=x;;y=tr[y].fa)
        {
            st.push_back(y);
            if(is_root(y))break;
        }
        for(int i=(int)st.size()-1;i>=0;i--)push(st[i]);
        while(!is_root(x))
        {
            int y=tr[x].fa,z=tr[y].fa;
            if(!is_root(y))rotate((tr[y].ch[0]==x)==(tr[z].ch[0]==y)?y:x);
            rotate(x);
        }
    }
    void access(int x)
    {
        for(int y=0;x;y=x,x=tr[x].fa)splay(x),tr[x].ch[1]=y,pull(x);
    }
    void makeroot(int x) { access(x),splay(x),reverse_node(x); }
    int findroot(int x)
    {
        access(x),splay(x),push(x);
        while(tr[x].ch[0])x=tr[x].ch[0],push(x);
        splay(x);
        return x;
    }
    bool connected(int x,int y) { return x==y||findroot(x)==findroot(y); }
    bool link(int x,int y)
    {
        makeroot(x);
        if(findroot(y)==x)return false;
        tr[x].fa=y;
        return true;
    }
    bool cut(int x,int y)
    {
        makeroot(x),access(y),splay(y);
        if(tr[y].ch[0]!=x||tr[x].ch[1])return false;
        tr[y].ch[0]=tr[x].fa=0,pull(y);
        return true;
    }
    void set_val(int x,ll v) { access(x),splay(x),tr[x].v=v,pull(x); }
    bool query(int x,int y,ll &sum)
    {
        if(!connected(x,y))return false;
        makeroot(x),access(y),splay(y),sum=tr[y].sum;
        return true;
    }
};
\end{lstlisting}
% SOURCE: 02-数据结构/珂朵莉树ODT.cpp
\subsection{珂朵莉树ODT}

\note{odt保存互不交叠的常值闭区间Node\{l,\allowbreak{}r,\allowbreak{}v\}，n为长度。split分裂出操作边界，赋值会合并覆盖范围；迭代器在erase后失效。随机数据适用性不代表任意序列的最坏复杂度保证。}
\begin{lstlisting}
const int N=100005;
\end{lstlisting}
\note{珂朵莉树(\allowbreak{}ODT /\allowbreak{} 颜色段均摊)\allowbreak{}：set 维护极长同色段 区间赋值遍历并删除的段可按生命周期均摊；每次 split 只新增 O(\allowbreak{}1)\allowbreak{} 段。 add/\allowbreak{}kth/\allowbreak{}求和等会遍历但不删除段，重复扫很多段可达 O(\allowbreak{}nq)\allowbreak{}，有赋值也不保证性能。 仅赋值型删段操作有上述确定性均摊；混合操作的随机期望依赖输入分布，不能当最坏保证。}
\begin{lstlisting}
struct Node
{
    int l,r;
    mutable ll v;
    Node(int l,int r,ll v):l(l),r(r),v(v){}
    bool operator<(const Node &o)const
    {
        return l<o.l;
    }
};
int n;
set<Node> odt;
\end{lstlisting}
\note{把位置 pos 所在的段拆成 [\allowbreak{}l,\allowbreak{}pos-\allowbreak{}1]\allowbreak{} 和 [\allowbreak{}pos,\allowbreak{}r]\allowbreak{}，返回起点为 pos 的迭代器}
\begin{lstlisting}
auto split(int pos)
{
    if(pos>n)return odt.end();//右端越界，直接返回尾迭代器
    auto it=odt.lower_bound(Node(pos,0,0));
    if(it!=odt.end()&&it->l==pos)return it;
    --it;
    int l=it->l,r=it->r;
    ll v=it->v;
    odt.erase(it);
    odt.insert(Node(l,pos-1,v));
    return odt.insert(Node(pos,r,v)).first;
}
\end{lstlisting}
\note{区间赋值，O(\allowbreak{}段数 log n)\allowbreak{}}
\begin{lstlisting}
void assign(int l,int r,ll v)
{
    auto itr=split(r+1),itl=split(l);
    odt.erase(itl,itr);
    odt.insert(Node(l,r,v));
}
\end{lstlisting}
\note{区间加，把 l.\allowbreak{}.\allowbreak{}r 覆盖到的每段整体加上 v}
\begin{lstlisting}
void add(int l,int r,ll v)
{
    auto itr=split(r+1),itl=split(l);
    for(auto it=itl;it!=itr;++it)it->v+=v;
}
\end{lstlisting}
\note{区间和}
\begin{lstlisting}
ll query_sum(int l,int r)
{
    auto itr=split(r+1),itl=split(l);
    ll s=0;
    for(auto it=itl;it!=itr;++it)s+=(ll)(it->r-it->l+1)*it->v;
    return s;
}
\end{lstlisting}
\note{区间内等于 v 的个数}
\begin{lstlisting}
int query_cnt(int l,int r,ll v)
{
    auto itr=split(r+1),itl=split(l);
    int c=0;
    for(auto it=itl;it!=itr;++it)
        if(it->v==v)c+=it->r-it->l+1;
    return c;
}
\end{lstlisting}
\note{区间第 k 小(\allowbreak{}1<\allowbreak{}=\allowbreak{}k<\allowbreak{}=\allowbreak{}r-\allowbreak{}l+\allowbreak{}1)\allowbreak{}，把覆盖到的段按值排序后累加长度}
\begin{lstlisting}
ll query_kth(int l,int r,int k)
{
    vector<pair<ll,int>> tmp;
    auto itr=split(r+1),itl=split(l);
    for(auto it=itl;it!=itr;++it)tmp.push_back({it->v,it->r-it->l+1});
    sort(tmp.begin(),tmp.end());
    for(auto &x:tmp)
    {
        if(k<=x.second)return x.first;
        k-=x.second;
    }
    return -1;//k 超过区间长度
}
\end{lstlisting}
% SOURCE: 02-数据结构/WaveletTree.cpp
\subsection{WaveletTree}

\note{对象的value是升序去重值，root为根；节点前缀表记前i项中有多少进入左子树。查询区间为0起半开[\allowbreak{}l,\allowbreak{}r)\allowbreak{}，k从0起；构造保存静态序列，修改须重建。}
\note{静态区间第 k 小与 <\allowbreak{}=\allowbreak{}x 计数，0-\allowbreak{}indexed 半开区间 [\allowbreak{}l,\allowbreak{}r)\allowbreak{}，k 从 0 开始。 离散化后稳定分流，构建 O(\allowbreak{}n log \ensuremath{\sigma})\allowbreak{}，查询 O(\allowbreak{}log \ensuremath{\sigma})\allowbreak{}，空间 O(\allowbreak{}n log \ensuremath{\sigma})\allowbreak{}。不支持点修改。}
\begin{lstlisting}
struct WaveletTree
{
    struct Node{int lo,hi,left=-1,right=-1;vector<int> pref;};
    vector<Node> t;vector<long long> value;int root=-1,n;
    int build(vector<int> a,int lo,int hi)
    {
        int p=t.size();t.push_back({lo,hi,-1,-1,{0}});if(lo==hi)return p;
        vector<int> l,r;int m=(lo+hi)/2;
        for(int x:a){if(x<=m)l.push_back(x);else r.push_back(x);t[p].pref.push_back(l.size());}
        if(!l.empty()){int c=build(l,lo,m);t[p].left=c;}
        if(!r.empty()){int c=build(r,m+1,hi);t[p].right=c;}
        return p;
    }
    WaveletTree(const vector<long long> &a):value(a),n(a.size())
    {
        sort(value.begin(),value.end());value.erase(unique(value.begin(),value.end()),value.end());
        vector<int> id;for(auto x:a)id.push_back(lower_bound(value.begin(),value.end(),x)-value.begin());
        if(n)root=build(id,0,value.size()-1);
    }
    long long kth(int l,int r,int k)const
    {
        assert(0<=l&&l<=r&&r<=n&&0<=k&&k<r-l);int p=root;
        while(t[p].lo<t[p].hi){int a=t[p].pref[l],b=t[p].pref[r];if(k<b-a)p=t[p].left,l=a,r=b;else k-=b-a,p=t[p].right,l-=a,r-=b;}
        return value[t[p].lo];
    }
    int count(int p,int l,int r,int x)const
    {if(p<0||l==r||x<t[p].lo)return 0;if(t[p].hi<=x)return r-l;int a=t[p].pref[l],b=t[p].pref[r];return count(t[p].left,a,b,x)+count(t[p].right,l-a,r-b,x);}
    int count_le(int l,int r,long long x)const{assert(0<=l&&l<=r&&r<=n);int id=upper_bound(value.begin(),value.end(),x)-value.begin()-1;return count(root,l,r,id);}
};
\end{lstlisting}
% SOURCE: 02-数据结构/KD-Tree.cpp
\subsection{KD-Tree}

\note{对象a为重排点集，tr存分割节点和包围盒，rt为根；查询传目标坐标等标量，返回最近距离平方。构造后为静态点集，修改须重建；空集返回LLONG\_MAX，坐标平方须在ll范围内。}
\note{静态二维 KD 树：中位数建树 O(\allowbreak{}n log n)\allowbreak{}，查询最坏 O(\allowbreak{}n)\allowbreak{} 坐标绝对值不超过 1e9，距离返回平方；空集返回 LLONG\_MAX}
\begin{lstlisting}
struct KDTree
{
    struct Point { ll x,y; };
    struct Node { Point p; ll lo[2],hi[2]; int l=0,r=0,sz=0; };
    vector<Point> a;
    vector<Node> tr;
    int rt=0;
    KDTree(vector<Point> b):a(b),tr(b.size()+1) { rt=build(0,(int)a.size()-1,0); }
    int build(int l,int r,int d)
    {
        if(l>r)return 0;
        int mid=(l+r)>>1,p=mid+1;
        nth_element(a.begin()+l,a.begin()+mid,a.begin()+r+1,[d](Point u,Point v) { return d?u.y<v.y:u.x<v.x; });
        tr[p].p=a[mid],tr[p].l=build(l,mid-1,d^1),tr[p].r=build(mid+1,r,d^1);
        tr[p].lo[0]=tr[p].hi[0]=a[mid].x,tr[p].lo[1]=tr[p].hi[1]=a[mid].y,tr[p].sz=1;
        for(int v:{tr[p].l,tr[p].r})if(v)
        {
            tr[p].sz+=tr[v].sz;
            for(int k=0;k<2;k++)tr[p].lo[k]=min(tr[p].lo[k],tr[v].lo[k]),tr[p].hi[k]=max(tr[p].hi[k],tr[v].hi[k]);
        }
        return p;
    }
    ll dist(Point u,Point v) { return (u.x-v.x)*(u.x-v.x)+(u.y-v.y)*(u.y-v.y); }
    ll bound(int p,Point q)
    {
        if(!p)return LLONG_MAX;
        ll ans=0,b[2]={q.x,q.y};
        for(int k=0;k<2;k++)
        {
            ll d=max({tr[p].lo[k]-b[k],b[k]-tr[p].hi[k],0LL});
            ans+=d*d;
        }
        return ans;
    }
    void nearest(int p,Point q,ll &ans)
    {
        if(!p)return;
        ans=min(ans,dist(q,tr[p].p));
        int l=tr[p].l,r=tr[p].r;
        if(bound(l,q)>bound(r,q))swap(l,r);
        if(bound(l,q)<ans)nearest(l,q,ans);
        if(bound(r,q)<ans)nearest(r,q,ans);
    }
    ll nearest(Point q)
    {
        ll ans=LLONG_MAX;
        nearest(rt,q,ans);
        return ans;
    }
    int rectangle(int p,ll x1,ll y1,ll x2,ll y2)
    {
        if(!p||tr[p].hi[0]<x1||tr[p].lo[0]>x2||tr[p].hi[1]<y1||tr[p].lo[1]>y2)return 0;
        if(x1<=tr[p].lo[0]&&tr[p].hi[0]<=x2&&y1<=tr[p].lo[1]&&tr[p].hi[1]<=y2)return tr[p].sz;
        Point q=tr[p].p;
        int ans=x1<=q.x&&q.x<=x2&&y1<=q.y&&q.y<=y2;
        return ans+rectangle(tr[p].l,x1,y1,x2,y2)+rectangle(tr[p].r,x1,y1,x2,y2);
    }
    int rectangle(ll x1,ll y1,ll x2,ll y2)
    {
        if(x1>x2||y1>y2)return 0;
        return rectangle(rt,x1,y1,x2,y2);
    }
};
\end{lstlisting}
% SOURCE: 02-数据结构/替罪羊树.cpp
\subsection{替罪羊树}

\note{节点记录键、有效计数与子树大小，删除标记区别实际节点和有效元素；root为树根。失衡后重建子树，查询排名按有效元素数；重建不能把已删节点当活元素。}
\note{替罪羊树：重复值计数，alpha 重构，单次操作均摊 O(\allowbreak{}log n)\allowbreak{}}
\begin{lstlisting}
struct Scapegoat
{
    struct Node { int l=0,r=0,v=0,cnt=0,sz=0,tot=0; };
    vector<Node> tr;
    int rt=0;
    Scapegoat() { tr.push_back(Node()); }
    void pull(int p)
    {
        tr[p].sz=tr[tr[p].l].sz+tr[tr[p].r].sz+tr[p].cnt;
        tr[p].tot=tr[tr[p].l].tot+tr[tr[p].r].tot+1;
    }
    void collect(int p,vector<int> &a)
    {
        if(!p)return;
        collect(tr[p].l,a);
        if(tr[p].cnt)a.push_back(p);
        collect(tr[p].r,a);
    }
    int build(vector<int> &a,int l,int r)
    {
        if(l>r)return 0;
        int mid=(l+r)>>1,p=a[mid];
        tr[p].l=build(a,l,mid-1),tr[p].r=build(a,mid+1,r);
        pull(p);
        return p;
    }
    int rebuild(int p)
    {
        vector<int> a;
        collect(p,a);
        return build(a,0,(int)a.size()-1);
    }
    int change(int p,int x,int d)
    {
        if(!p)
        {
            if(d<0)return 0;
            Node u;
            u.v=x,u.cnt=u.sz=u.tot=1;
            tr.push_back(u);
            return (int)tr.size()-1;
        }
        if(x==tr[p].v)tr[p].cnt=max(0,tr[p].cnt+d);
        else if(x<tr[p].v)tr[p].l=change(tr[p].l,x,d);
        else tr[p].r=change(tr[p].r,x,d);
        pull(p);
        // 用不同键数衡量结构重量，删除后清除过多墓碑
        int w=max(tr[tr[p].l].tot,tr[tr[p].r].tot);
        if(w*4>tr[p].tot*3||tr[p].sz*2<tr[p].tot)p=rebuild(p);
        return p;
    }
    void insert(int x) { rt=change(rt,x,1); }
    void erase(int x) { rt=change(rt,x,-1); }
    int rank(int x)
    {
        int p=rt,ans=1;
        while(p)
        {
            if(x<=tr[p].v)p=tr[p].l;
            else ans+=tr[tr[p].l].sz+tr[p].cnt,p=tr[p].r;
        }
        return ans;
    }
    bool kth(int k,int &x)
    {
        if(k<1||k>tr[rt].sz)return false;
        int p=rt;
        while(p)
        {
            int s=tr[tr[p].l].sz;
            if(k<=s)p=tr[p].l;
            else if(k<=s+tr[p].cnt) { x=tr[p].v; return true; }
            else k-=s+tr[p].cnt,p=tr[p].r;
        }
        return false;
    }
    bool prev(int x,int &v) { return kth(rank(x)-1,v); }
    bool next(int x,int &v)
    {
        int p=rt,k=0;
        while(p)
        {
            if(tr[p].v<=x)k+=tr[tr[p].l].sz+tr[p].cnt,p=tr[p].r;
            else p=tr[p].l;
        }
        return kth(k+1,v);
    }
};
\end{lstlisting}
% SOURCE: 02-数据结构/哈希表.cpp
\subsection{哈希表}

\note{hc/\allowbreak{}ho分别是链式/\allowbreak{}开放寻址哈希表，x是键；节点或槽位保存键、计数及访问标记。两种实现按需选一种，换数据集清空；开放寻址容量须留空槽，不能填满后无限探测。}
\note{哈希表（整数键）：拉链法 +\allowbreak{} 开放寻址法 拉链法 平均 O(\allowbreak{}1)\allowbreak{}，天然支持负数；开放寻址法 平均 O(\allowbreak{}1)\allowbreak{}，装太满会退化}
\begin{lstlisting}
const int N=100003;         // 容量；质数不能保证无冲突
\end{lstlisting}
\note{-\allowbreak{}-\allowbreak{}-\allowbreak{}-\allowbreak{}-\allowbreak{}-\allowbreak{}-\allowbreak{}-\allowbreak{}-\allowbreak{}-\allowbreak{}-\allowbreak{}-\allowbreak{}-\allowbreak{}-\allowbreak{}-\allowbreak{}-\allowbreak{} 拉链法 -\allowbreak{}-\allowbreak{}-\allowbreak{}-\allowbreak{}-\allowbreak{}-\allowbreak{}-\allowbreak{}-\allowbreak{}-\allowbreak{}-\allowbreak{}-\allowbreak{}-\allowbreak{}-\allowbreak{}-\allowbreak{}-\allowbreak{}-\allowbreak{} h[\allowbreak{}k]\allowbreak{} 存链表头（结点下标），e[\allowbreak{}]\allowbreak{} 存值，nxt[\allowbreak{}]\allowbreak{} 存后继，idx 从 0 开始}
\begin{lstlisting}
struct HashChain{
    int h[N],e[N],nxt[N],idx;
    void init()
    {
        memset(h,-1,sizeof(h));   // -1 表示空
        idx=0;
    }
    int get_key(int x)
    {
        return (x%N+N)%N;         // 负数取模也要落在 [0,N)
    }
    void insert(int x)            // O(1) 平均
    {
        int k=get_key(x);
        e[idx]=x,nxt[idx]=h[k],h[k]=idx++;
    }
    bool find(int x)              // O(1) 平均
    {
        int k=get_key(x);
        for(int i=h[k];i!=-1;i=nxt[i])
            if(e[i]==x)return true;
        return false;
    }
};
\end{lstlisting}
\note{-\allowbreak{}-\allowbreak{}-\allowbreak{}-\allowbreak{}-\allowbreak{}-\allowbreak{}-\allowbreak{}-\allowbreak{}-\allowbreak{}-\allowbreak{}-\allowbreak{}-\allowbreak{}-\allowbreak{}-\allowbreak{}-\allowbreak{}-\allowbreak{} 开放寻址法 -\allowbreak{}-\allowbreak{}-\allowbreak{}-\allowbreak{}-\allowbreak{}-\allowbreak{}-\allowbreak{}-\allowbreak{}-\allowbreak{}-\allowbreak{}-\allowbreak{}-\allowbreak{}-\allowbreak{}-\allowbreak{}-\allowbreak{}-\allowbreak{} h[\allowbreak{}k]\allowbreak{} 为 null 表示空位；线性探测，key 冲突就往后找 注意：表要开 2\textasciitilde{}3 倍，装载因子太高会退化成 O(\allowbreak{}n)\allowbreak{}}
\begin{lstlisting}
struct HashOpen{
    const int null=0x3f3f3f3f;    // 空标记，保证没有输入等于它
    int h[N];
    void init()
    {
        memset(h,0x3f,sizeof(h));
    }
    int get_key(int x)
    {
        return (x%N+N)%N;
    }
    // 已有位置或首个空位；表满且不存在时返回 -1。x 不能等于 null。
    int find(int x)
    {
        int k=get_key(x);
        int start=k;
        while(h[k]!=null&&h[k]!=x)
        {
            k++;
            if(k==N)k=0;          // 走到表尾绕回表头
            if(k==start)return -1;
        }
        return k;
    }
    bool insert(int x)            // 重复插入只留一份；满表返回 false
    {
        if(x==null)return false;
        int p=find(x);
        if(p<0)return false;
        h[p]=x;
        return true;
    }
    bool count(int x)
    {
        if(x==null)return false;
        int p=find(x);
        return p>=0&&h[p]==x;
    }
};
HashChain hc;
HashOpen ho;
int n,x;
\end{lstlisting}
% SOURCE: 02-数据结构/二叉搜索树.cpp
\subsection{二叉搜索树}

\note{t为树的节点池，root为根，0为空节点；x为单次操作的键，op为操作类型。节点维护子树信息，插入/\allowbreak{}删除后用返回值更新根；普通BST不保证平衡，不能把最坏复杂度当log n。}
\note{二叉搜索树 BST（数组版，无旋）：插入 /\allowbreak{} 删除 /\allowbreak{} 排名 /\allowbreak{} 第k小 /\allowbreak{} 前驱后继，维护 size 约定：不允许重复键（第二次数相同值直接忽略），随机数据下期望 O(\allowbreak{}log n)\allowbreak{}，退化会变 O(\allowbreak{}n)\allowbreak{} 排名约定：get\_rank(\allowbreak{}x)\allowbreak{} 返回「比 x 小的数个数」（0 起）}
\begin{lstlisting}
const int N=200005;
const int INF=0x3f3f3f3f;
int n,op,x;
struct BST{
    int val[N],sz[N],ls[N],rs[N];   // 0 号结点是空结点
    int tot;
    void clear(){tot=0;}
    int new_node(int x)
    {
        tot++;
        val[tot]=x,sz[tot]=1,ls[tot]=0,rs[tot]=0;
        return tot;
    }
    void push_up(int p){sz[p]=sz[ls[p]]+sz[rs[p]]+1;}
    // 插入 x，重复值忽略；返回新的根
    int insert(int p,int x)
    {
        if(p==0)return new_node(x);
        if(x<val[p])ls[p]=insert(ls[p],x);
        else if(x>val[p])rs[p]=insert(rs[p],x);
        else return p;              // 已存在，不动
        push_up(p);
        return p;
    }
    // 删除 x，返回新的根；用左子树最大值替换（前驱顶替）
    int erase(int p,int x)
    {
        if(p==0)return 0;
        if(x<val[p])ls[p]=erase(ls[p],x);
        else if(x>val[p])rs[p]=erase(rs[p],x);
        else
        {
            if(ls[p]==0)return rs[p];
            if(rs[p]==0)return ls[p];
            int q=ls[p];
            while(rs[q])q=rs[q];    // 找前驱，把它换到 p
            val[p]=val[q];
            ls[p]=erase(ls[p],val[q]);
        }
        push_up(p);
        return p;
    }
    bool find(int p,int x)
    {
        while(p)
        {
            if(x<val[p])p=ls[p];
            else if(x>val[p])p=rs[p];
            else return true;
        }
        return false;
    }
    // 比 x 小的数个数（0 起），O(树高)
    int get_rank(int p,int x)
    {
        int cnt=0;
        while(p)
        {
            if(x<=val[p])p=ls[p];
            else cnt+=sz[ls[p]]+1,p=rs[p];
        }
        return cnt;
    }
    // 第 k 小（k 从 1 开始），不存在返回 0
    int kth(int p,int k)
    {
        if(k<1||k>sz[p])return 0;
        while(p)
        {
            if(k<=sz[ls[p]])p=ls[p];
            else if(k==sz[ls[p]]+1)return val[p];
            else k-=sz[ls[p]]+1,p=rs[p];
        }
        return 0;
    }
    // 小于 x 的最大值，不存在返回 -INF
    int get_pre(int p,int x)
    {
        int ans=-INF;
        while(p)
        {
            if(val[p]<x)ans=val[p],p=rs[p];
            else p=ls[p];
        }
        return ans;
    }
    // 大于 x 的最小值，不存在返回 INF
    int get_next(int p,int x)
    {
        int ans=INF;
        while(p)
        {
            if(val[p]>x)ans=val[p],p=ls[p];
            else p=rs[p];
        }
        return ans;
    }
};
BST t;
int root;
\end{lstlisting}
\section{字符串}
% SOURCE: 03-字符串/字符串哈希.cpp
\subsection{字符串哈希}

\note{s为输入，H1、H2分别保存单/\allowbreak{}双模前缀哈希及幂；区间参数按各接口的闭区间约定填写。换串重新build，比较须长度相同、参数一致；双哈希仍有碰撞可能。}
\note{至多一次失配的最长匹配：先求 LCP，未到串尾则跳过一个字符，再求一次 LCP；两次查询。 LCP 上界截在两串剩余长度内；哈希有碰撞概率，需严格正确可用 SA+\allowbreak{}RMQ 的 LCP。}
\begin{lstlisting}
typedef unsigned long long ull;
const int N=1000005;
const ull base1=131;
const ull base2=13331;
int n,m;
string s;
\end{lstlisting}
\note{单哈希：自然溢出，预处理 O(\allowbreak{}n)\allowbreak{}，子串哈希 O(\allowbreak{}1)\allowbreak{}}
\begin{lstlisting}
struct Hash1
{
    ull h[N],p[N];
    void build(const string &str)
    {
        int len=str.length();
        p[0]=1;
        h[0]=0;
        for(int i=1;i<=len;i++)
        {
            p[i]=p[i-1]*base1;
            h[i]=h[i-1]*base1+(unsigned char)str[i-1];
        }
    }
    // 下标 1 开始，[l,r] 的哈希值
    ull get(int l,int r)
    {
        if(l>r)return 0;
        return h[r]-h[l-1]*p[r-l+1];
    }
};
\end{lstlisting}
\note{双哈希：取模 1e9+\allowbreak{}7 /\allowbreak{} 1e9+\allowbreak{}9，降低冲突概率}
\begin{lstlisting}
struct Hash2
{
    static const ull mod1=1000000007;
    static const ull mod2=1000000009;
    ull h1[N],h2[N],p1[N],p2[N];
    void build(const string &str)
    {
        int len=str.length();
        p1[0]=p2[0]=1;
        h1[0]=h2[0]=0;
        for(int i=1;i<=len;i++)
        {
            p1[i]=p1[i-1]*base1%mod1,p2[i]=p2[i-1]*base2%mod2;
            h1[i]=(h1[i-1]*base1+(unsigned char)str[i-1])%mod1;
            h2[i]=(h2[i-1]*base2+(unsigned char)str[i-1])%mod2;
        }
    }
    // 下标 1 开始，返回 [l,r] 的双哈希
    pair<ull,ull> get(int l,int r)
    {
        if(l>r)return {0,0};
        ull x=(h1[r]+mod1-h1[l-1]*p1[r-l+1]%mod1)%mod1;
        ull y=(h2[r]+mod2-h2[l-1]*p2[r-l+1]%mod2)%mod2;
        return {x,y};
    }
};
Hash1 H1;
Hash2 H2;
\end{lstlisting}
% SOURCE: 03-字符串/KMP.cpp
\subsection{KMP}

\note{s为文本，t为模式串，字符串从0起；nxt按匹配长度记录失配后的最长可复用前缀。模式变化后重建nxt，文本扫描保留重叠匹配；不要把匹配长度和字符下标混用。}
\note{长 n 串最长 border 长 b：最短周期 p=\allowbreak{}n-\allowbreak{}b；整个串完整重复还需 n\%p=\allowbreak{}=\allowbreak{}0。 p,\allowbreak{}q 都是周期且 n>\allowbreak{}=\allowbreak{}p+\allowbreak{}q-\allowbreak{}gcd(\allowbreak{}p,\allowbreak{}q)\allowbreak{} 时，gcd(\allowbreak{}p,\allowbreak{}q)\allowbreak{} 也是周期；周期不等于整串循环节。 border 链能变成树，前缀出现/\allowbreak{}祖先约束可结合 DFS 序与树状数组，不必重复匹配整串。 长度超过 n/\allowbreak{}2 的真 border 沿链按最短周期 p 等差递减，可整段跳；短 border 仍需另处理。}
\begin{lstlisting}
const int N=1000005;
int n,m;
string s,t;
int nxt[N];
\end{lstlisting}
\note{求模式串 p 的 nxt 数组，nxt[\allowbreak{}i]\allowbreak{} 表示 p[\allowbreak{}0.\allowbreak{}.\allowbreak{}i-\allowbreak{}1]\allowbreak{} 的最长公共前后缀长度 nxt[\allowbreak{}0]\allowbreak{}=\allowbreak{}-\allowbreak{}1 便于失配跳转；O(\allowbreak{}|p|)\allowbreak{}}
\begin{lstlisting}
void get_nxt(const string &p)
{
    int l=p.length();
    nxt[0]=-1;
    int i=0,j=-1;
    while(i<l)
    {
        if(j==-1||p[i]==p[j])//j=-1 表示从头重新比较
        {
            i++;
            j++;
            nxt[i]=j;
        }
        else j=nxt[j];
    }
}
\end{lstlisting}
\note{在 s 中找 p 的所有出现次数（允许重叠）；必须先 get\_nxt(\allowbreak{}p)\allowbreak{} 复杂度 O(\allowbreak{}|s|+\allowbreak{}|p|)\allowbreak{}}
\begin{lstlisting}
int kmp(const string &s,const string &p)
{
    int l1=s.length(),l2=p.length();
    int cnt=0,i=0,j=0;
    if(l2==0)return 0;
    while(i<l1)
    {
        if(j==l2)//匹配成功，继续找下一次
        {
            cnt++;
            j=nxt[j];
        }
        if(s[i]==p[j])
        {
            i++;
            j++;
        }
        else
        {
            if(j==0)i++;
            else j=nxt[j];
        }
    }
    if(j==l2)cnt++;
    return cnt;
}
\end{lstlisting}
\note{最小循环节长度：n-\allowbreak{}nxt[\allowbreak{}n]\allowbreak{} 若能整除 n 就是循环节，否则整串自己当循环节 例 "abcab" 长度 5、n-\allowbreak{}nxt[\allowbreak{}n]\allowbreak{}=\allowbreak{}3 不整除，返回 5}
\begin{lstlisting}
int min_cycle(const string &p)
{
    int l=p.length();
    if(!l)return 0;
    get_nxt(p);
    int c=l-nxt[l];
    if(l%c)return l;
    return c;
}
\end{lstlisting}
% SOURCE: 03-字符串/Z函数.cpp
\subsection{Z函数}

\note{s是计算Z值的串，t为扩展匹配示例的另一串；z[\allowbreak{}i]\allowbreak{}是所预处理串与其后缀的最长公共前缀长度，z[\allowbreak{}0]\allowbreak{}=\allowbreak{}0。字符串下标0起，换串重建Z数组。}
\begin{lstlisting}
const int N=1000005;
int n,m;
string s,t;
int z[N];// z[0]=0；z[i] = s 与 s[i..] 的最长公共前缀长度
\end{lstlisting}
\note{扩展 KMP /\allowbreak{} Z 数组，O(\allowbreak{}n)\allowbreak{} 维护 [\allowbreak{}l,\allowbreak{}r]\allowbreak{} 为当前已知最靠右的匹配段 s[\allowbreak{}l.\allowbreak{}.\allowbreak{}r]\allowbreak{}=\allowbreak{}=\allowbreak{}s[\allowbreak{}0.\allowbreak{}.\allowbreak{}r-\allowbreak{}l]\allowbreak{}}
\begin{lstlisting}
void get_z(const string &str)
{
    int len=str.length();
    z[0]=0;
    int l=0,r=0;
    for(int i=1;i<len;i++)
    {
        if(i<=r)z[i]=min(r-i+1,z[i-l]);//在匹配段内直接复用
        else z[i]=0;//必须在 while 前清零，否则会复用上一次调用的旧值
        while(i+z[i]<len&&str[z[i]]==str[i+z[i]])z[i]++;
        if(i+z[i]-1>r)l=i,r=i+z[i]-1;
    }
}
\end{lstlisting}
\note{扩展 KMP：ex[\allowbreak{}i]\allowbreak{} =\allowbreak{} t 与 s[\allowbreak{}i.\allowbreak{}.\allowbreak{}]\allowbreak{} 的最长公共前缀长度（s 是文本，t 是模式）}
\begin{lstlisting}
void get_ex(const string &s,const string &t,int *ex)
{
    int n=s.length(),m=t.length();
    get_z(t);
    int l=0,r=0;
    for(int i=0;i<n;i++)
    {
        if(i<=r)ex[i]=min(r-i+1,z[i-l]);
        else ex[i]=0;//同样先清零再用
        while(i+ex[i]<n&&ex[i]<m&&s[i+ex[i]]==t[ex[i]])ex[i]++;
        if(i+ex[i]-1>r)l=i,r=i+ex[i]-1;
    }
}
\end{lstlisting}
% SOURCE: 03-字符串/Manacher.cpp
\subsection{Manacher}

\note{s为原串；M操作全局变换串buf和半径rr，buf用负数哨兵区分任意字节，rr[\allowbreak{}i]\allowbreak{}为中心半径，数值等于相应最长回文长度。原串与变换串下标不同；初始化后通过对应映射查询，不直接把rr下标当原串位置。}
\begin{lstlisting}
const int N=1000005;
int n,m;
string s;
\end{lstlisting}
\note{预处理成 "\#a\#b\#a\#" 形式，两端加互不相同且不出现在原串的哨兵，省掉边界判断}
\begin{lstlisting}
int len_cur;
int buf[N<<1];// 负数哨兵与任意字节互不相同
int rr[N<<1];// rr[i]=以 i 为中心的回文半径，数值上等于该中心的最长回文长度
\end{lstlisting}
\note{求最长回文子串长度 /\allowbreak{} 回文子串个数，O(\allowbreak{}n)\allowbreak{}}
\begin{lstlisting}
struct Manacher
{
    // 传入原始串，返回转化后长度 len_cur，rr[] 为结果
    int build(const string &str)
    {
        int len=0,nn=str.length();
        buf[0]=-1;
        for(int j=0;j<nn;j++)
        {
            buf[++len]=-2;
            buf[++len]=(unsigned char)str[j];
        }
        buf[++len]=-2;
        buf[++len]=-3;
        len_cur=len;
        int r=0,c=0;
        rr[len]=0;
        for(int i=1;i<len;i++)
        {
            rr[i]=0;//同样先清零，避免复用上一次调用的旧半径
            if(i<=r)rr[i]=min(rr[(c<<1)-i],r-i);//镜像复制，右端不超过 r
            while(buf[i+rr[i]+1]==buf[i-rr[i]-1])rr[i]++;
            if(rr[i]+i>r)r=i+rr[i],c=i;
        }
        return len;
    }
    // 最长回文子串长度
    int longest()
    {
        int ans=0;
        for(int i=1;i<=len_cur;i++)if(rr[i]>ans)ans=rr[i];
        return ans;
    }
    // 回文子串个数（按出现位置计，可重复）
    // 每个中心贡献 (rr[i]+1)/2 个：偶数中心 rr 为偶数，奇数中心 rr 为奇数
    ll count_pal()
    {
        ll ans=0;
        for(int i=1;i<=len_cur;i++)ans+=(rr[i]+1)>>1;
        return ans;
    }
};
Manacher M;
\end{lstlisting}
% SOURCE: 03-字符串/AC自动机.cpp
\subsection{AC自动机}

\note{trie[\allowbreak{}u]\allowbreak{}[\allowbreak{}c]\allowbreak{}是转移，nxt[\allowbreak{}u]\allowbreak{}是fail，id[\allowbreak{}i]\allowbreak{}是第i个模式的终点；cnt记模式终止次数。root=\allowbreak{}1，tot是节点数，0为虚节点；容量按模式总长度设置。建立补全转移后不能继续直接插串，换模式集重置树和计数。}
\note{匹配文本前缀到状态 u：它的 fail 树祖先均为当前位置结束的匹配模式。 反过来模式结点 v 的出现次数/\allowbreak{}位置和，可统计 fail 子树内被访问状态；DFS 序转 BIT 区间。 多查询若对文本前缀长度单调，可共建一棵 AC，再整体二分按 mid 扫文本。}
\begin{lstlisting}
const int N=200005;
int n,m,root=1,tot=1;
string s;
int trie[N][26];// Trie 儿子
int nxt[N];// fail 指针
int cnt[N];// 该结点被多少个模式串覆盖
int id[N];// 每个模式串的终止结点
\end{lstlisting}
\note{插入模式串，返回终止结点编号}
\begin{lstlisting}
int insert(const string &str)
{
    int p=root;
    for(int i=0;i<(int)str.length();i++)
    {
        int k=str[i]-'a';
        if(!trie[p][k])trie[p][k]=++tot;
        p=trie[p][k];
    }
    cnt[p]++;//终止结点覆盖数 +1，重复模式串也累计
    return p;
}
\end{lstlisting}
\note{清空 Trie，多组数据用；O(\allowbreak{}tot)\allowbreak{}}
\begin{lstlisting}
void clear_trie()
{
    for(int i=0;i<=tot;i++)
    {
        cnt[i]=0,nxt[i]=0;
        for(int j=0;j<26;j++)trie[i][j]=0;
    }
    tot=1;
}
\end{lstlisting}
\note{建 fail 指针（BFS），根的儿子 fail 指向根；O(\allowbreak{}结点数*26)\allowbreak{}}
\begin{lstlisting}
void build_fail()
{
    queue<int> q;
    for(int j=0;j<26;j++)
    {
        if(trie[root][j])nxt[trie[root][j]]=root,q.push(trie[root][j]);
        else trie[root][j]=root;//根处失配回到根
    }
    while(!q.empty())
    {
        int u=q.front();
        q.pop();
        cnt[u]+=cnt[nxt[u]];//继承 fail 的覆盖数，统计时不必沿 fail 链跳
        for(int j=0;j<26;j++)
        {
            int v=trie[u][j];
            if(v)nxt[v]=trie[nxt[u]][j],q.push(v);
            else trie[u][j]=trie[nxt[u]][j];
        }
    }
}
\end{lstlisting}
\note{文本串在自动机上跑一遍，返回所有模式串出现次数之和（含重复计入） 复杂度 O(\allowbreak{}|文本|)\allowbreak{}}
\begin{lstlisting}
ll query(const string &txt)
{
    int p=root;
    ll ans=0;
    for(int i=0;i<(int)txt.length();i++)
    {
        p=trie[p][txt[i]-'a'];
        ans+=cnt[p];
    }
    return ans;
}
\end{lstlisting}
\note{普通 Trie 应用（独立于上面的 AC）：f[\allowbreak{}j]\allowbreak{} 为选恰 j 个已插入串的最大 LCP 长度。 g[\allowbreak{}h]\allowbreak{} 为深度 h 的节点最大经过次数，f[\allowbreak{}j]\allowbreak{}>\allowbreak{}=\allowbreak{}h 等价于 g[\allowbreak{}h]\allowbreak{}>\allowbreak{}=\allowbreak{}j，是同一单调阶梯的转置。 节点 cnt 从 c-\allowbreak{}1 变 c，只需令 f[\allowbreak{}c]\allowbreak{}=\allowbreak{}max(\allowbreak{}f[\allowbreak{}c]\allowbreak{},\allowbreak{}h)\allowbreak{}；g 用于解释，不必实际存储。 小写非空串，重复串保留重数；插入后返回 sum(\allowbreak{}f[\allowbreak{}j]\allowbreak{} xor j)\allowbreak{}，j=\allowbreak{}1.\allowbreak{}.\allowbreak{}已插入串数。}
\begin{lstlisting}
struct PrefixLCP
{
    vector<array<int,26>> tr{array<int,26>{}};
    vector<int> cnt{0},f{0};
    ll sum=0;
    ll insert(const string &s)
    {
        int r=f.size();f.push_back(0);sum+=r;// 新增 j=r 的基准项 0 xor r
        int p=0,h=0;
        for(char ch:s)
        {
            int x=ch-'a';++h;
            if(!tr[p][x])
            {
                tr[p][x]=tr.size();
                tr.push_back({});cnt.push_back(0);
            }
            p=tr[p][x];int c=++cnt[p];
            if(h>f[c]){sum-=f[c]^c;f[c]=h;sum+=f[c]^c;}
        }
        return sum;
    }
};
\end{lstlisting}
% SOURCE: 03-字符串/后缀数组.cpp
\subsection{后缀数组}

\note{s下标0起；sa按字典序存后缀起点，rk是逆排名，height存相邻后缀LCP。cnt、xs、yr、tmp是计数排序与倍增缓冲；建立后用排名查询LCP，容量按实现的哨兵和额外位置预留。}
\begin{lstlisting}
const int N=1000005;
int n,m;
string s;
\end{lstlisting}
\note{sa[\allowbreak{}i]\allowbreak{} 排名第 i 的后缀起点，rk[\allowbreak{}i]\allowbreak{} 后缀 i 的排名，height[\allowbreak{}i]\allowbreak{}=\allowbreak{}lcp(\allowbreak{}sa[\allowbreak{}i-\allowbreak{}1]\allowbreak{},\allowbreak{}sa[\allowbreak{}i]\allowbreak{})\allowbreak{}}
\begin{lstlisting}
int sa[N],rk[N],height[N];
int cnt[N],xs[N],yr[N],tmp[N];
\end{lstlisting}
\note{倍增 +\allowbreak{} 基数排序，O(\allowbreak{}n log n)\allowbreak{}；字符集按 ASCII 处理，数字/\allowbreak{}大小写字母都行 传入 0 下标字符串首地址，函数内只用局部长度 len，顺手回填全局 n 调用前保证 str 末尾有 '\textbackslash{}0'（空串请直接特判返回）}
\begin{lstlisting}
void build_sa(char *str)
{
    int len=strlen(str);
    n=len;
    if(len==0)return;
    int mx=max(len,300);
    for(int i=0;i<=mx;i++)cnt[i]=0;
    for(int i=0;i<len;i++)cnt[(unsigned char)str[i]]++;
    for(int i=1;i<=mx;i++)cnt[i]+=cnt[i-1];
    for(int i=len-1;i>=0;i--)sa[--cnt[(unsigned char)str[i]]]=i;
    int p=1;
    rk[sa[0]]=0;
    for(int i=1;i<len;i++)
    {
        if(str[sa[i]]!=str[sa[i-1]])p++;
        rk[sa[i]]=p-1;
    }
    for(int k=1;k<len&&p<len;k<<=1)
    {
        // 第二关键字排序：后半段为空的后缀排最前，其余按上一轮 rk 顺序
        int num=0;
        for(int i=len-k;i<len;i++)yr[num++]=i;
        for(int i=0;i<len;i++)if(sa[i]>=k)yr[num++]=sa[i]-k;
        for(int i=0;i<=p;i++)cnt[i]=0;
        for(int i=0;i<len;i++)cnt[rk[yr[i]]]++;
        for(int i=1;i<=p;i++)cnt[i]+=cnt[i-1];
        for(int i=len-1;i>=0;i--)sa[--cnt[rk[yr[i]]]]=yr[i];
        // 用 (第一关键字, 第二关键字) 重算 rk
        p=1;
        tmp[sa[0]]=0;
        for(int i=1;i<len;i++)
        {
            int a=sa[i-1],b=sa[i];
            int a2=(a+k<len)?rk[a+k]:-1;
            int b2=(b+k<len)?rk[b+k]:-1;
            if(rk[a]!=rk[b]||a2!=b2)p++;
            tmp[b]=p-1;
        }
        for(int i=0;i<len;i++)rk[i]=tmp[i];
    }
    // height 用结论 height[rk[i]]>=height[rk[i-1]]-1，O(n)
    int k=0;
    for(int i=0;i<len;i++)
    {
        if(rk[i]==0)
        {
            k=0;
            height[0]=0;
            continue;
        }
        if(k)k--;
        int j=sa[rk[i]-1];
        while(i+k<len&&j+k<len&&str[i+k]==str[j+k])k++;
        height[rk[i]]=k;
    }
}
\end{lstlisting}
% SOURCE: 03-字符串/后缀自动机SAM.cpp
\subsection{后缀自动机SAM}

\note{s是输入，last是当前整串所在状态，tot是状态数；len为状态最长长度，link为后缀链接，nxt为字符转移，siz是出现次数。cnt、id用于按len计数排序；克隆状态初始出现次数为0，统计时沿link从长到短汇总。}
\begin{lstlisting}
const int N=200005;
int n,m;
string s;
\end{lstlisting}
\note{后缀自动机：O(\allowbreak{}n)\allowbreak{} 建图，nxt 转移、link 后缀链接、len 该状态最长子串长度 siz[\allowbreak{}i]\allowbreak{} 为状态 i 的 endpos 集合大小（即该状态代表子串的出现次数）}
\begin{lstlisting}
int tot,last;
int len[N],link[N],siz[N],cnt[N+2],id[N+2];
int nxt[N][26];
\end{lstlisting}
\note{清空 SAM，多组数据用；根固定为 1}
\begin{lstlisting}
void sam_init()
{
    tot=1,last=1;
    len[1]=0,link[1]=0,siz[1]=0;
    for(int j=0;j<26;j++)nxt[1][j]=0;
}
\end{lstlisting}
\note{插入一个字符 c（'a'\textasciitilde{}'z'）}
\begin{lstlisting}
void sam_extend(int c)
{
    int cur=++tot;
    len[cur]=len[last]+1,link[cur]=0,siz[cur]=1;
    for(int j=0;j<26;j++)nxt[cur][j]=0;
    int p=last;
    while(p&&!nxt[p][c])nxt[p][c]=cur,p=link[p];
    if(!p)link[cur]=1;
    else
    {
        int q=nxt[p][c];
        if(len[p]+1==len[q])link[cur]=q;
        else
        {
            int clone=++tot;// 分裂出的克隆点，siz 记为 0，不算新出现
            len[clone]=len[p]+1,link[clone]=link[q],siz[clone]=0;
            for(int j=0;j<26;j++)nxt[clone][j]=nxt[q][j];
            while(p&&nxt[p][c]==q)nxt[p][c]=clone,p=link[p];
            link[q]=link[cur]=clone;
        }
    }
    last=cur;
}
\end{lstlisting}
\note{按 len 基数排序，得到 len 单调不减的拓扑序 id[\allowbreak{}0.\allowbreak{}.\allowbreak{}tot-\allowbreak{}1]\allowbreak{}；顺便求每个状态的出现次数 所有字符加入后调用一次；siz 累加后不要继续 extend/\allowbreak{}build。}
\begin{lstlisting}
void sam_build()
{
    int mx=tot+1;
    for(int i=0;i<=mx;i++)cnt[i]=0;
    for(int i=1;i<=tot;i++)cnt[len[i]+1]++;
    for(int i=1;i<=mx;i++)cnt[i]+=cnt[i-1];
    for(int i=tot;i>=1;i--)id[--cnt[len[i]+1]]=i;// 先减后放，id 下标落在 0..tot-1
    for(int i=tot-1;i>=0;i--)siz[link[id[i]]]+=siz[id[i]];// 沿 link 累加 endpos
}
\end{lstlisting}
\note{本质不同子串个数：每个状态贡献 len[\allowbreak{}i]\allowbreak{}-\allowbreak{}len[\allowbreak{}link[\allowbreak{}i]\allowbreak{}]\allowbreak{}}
\begin{lstlisting}
ll count_diff_substr()
{
    ll ans=0;
    for(int i=2;i<=tot;i++)ans+=len[i]-len[link[i]];
    return ans;
}
\end{lstlisting}
\note{求 s 与 t 的最长公共子串长度：t 在 s 的 SAM 上跑，失配时沿 link 跳}
\begin{lstlisting}
int longest_common(const string &t)
{
    int p=1,l=0,ans=0;
    for(int i=0;i<(int)t.length();i++)
    {
        int c=t[i]-'a';
        while(p&&!nxt[p][c])p=link[p],l=len[p];// 跳 link，当前匹配长度退到 len[p]
        if(nxt[p][c])p=nxt[p][c],l++;
        else p=1,l=0;
        ans=max(ans,l);
    }
    return ans;
}
\end{lstlisting}
% SOURCE: 03-字符串/回文自动机PAM.cpp
\subsection{回文自动机PAM}

\note{len、link、nxt分别是回文长度、最长真回文后缀和转移，last是当前最长回文后缀；siz为出现次数。str从1起，pos[\allowbreak{}u]\allowbreak{}记录一个出现位置的右端点。两个根长度为0、-\allowbreak{}1；每个新字符串须重新初始化根和状态。}
\begin{lstlisting}
const int N=200005;
int n,m;
string s;
\end{lstlisting}
\note{回文自动机：结点 0 是偶根 len=\allowbreak{}0，结点 1 是奇根 len=\allowbreak{}-\allowbreak{}1，奇根的 link 指向自己 nxt 转移、link 后缀回文链、len 该回文串长度、siz 该回文串出现次数}
\begin{lstlisting}
int tot,last;
int len[N],link[N],siz[N],cnt[N+2],id[N+2];
int nxt[N][26];
char str[N];// 1-indexed 原串，str[0] 放用不到的哨兵 '#'
int pos[N];// 第一次出现的右端点；内容是 str[pos[u]-len[u]+1..pos[u]]
\end{lstlisting}
\note{清空 PAM，多组数据用}
\begin{lstlisting}
void pam_init()
{
    tot=1,last=0;
    str[0]='#';
    len[0]=0,link[0]=1;
    len[1]=-1,link[1]=1;
    siz[0]=siz[1]=0;
    for(int j=0;j<26;j++)nxt[0][j]=nxt[1][j]=0;
}
\end{lstlisting}
\note{从 p 出发沿 link 跳，找到能接上位置 i 的最长回文后缀}
\begin{lstlisting}
int get_fail(int p,int i)
{
    while(str[i-len[p]-1]!=str[i])p=link[p];
    return p;
}
\end{lstlisting}
\note{在末尾加入第 i 个字符（1-\allowbreak{}indexed，str[\allowbreak{}i]\allowbreak{} 已赋值）}
\begin{lstlisting}
void pam_extend(int i)
{
    int c=str[i]-'a';
    int p=get_fail(last,i);
    if(!nxt[p][c])
    {
        int cur=++tot;
        len[cur]=len[p]+2,siz[cur]=1,pos[cur]=i;
        for(int j=0;j<26;j++)nxt[cur][j]=0;
        if(len[cur]==1)link[cur]=0;// 单字符回文的后缀链接是偶根
        else link[cur]=nxt[get_fail(link[p],i)][c];
        nxt[p][c]=cur;
    }
    else siz[nxt[p][c]]++;// 该回文串已经存在，出现次数 +1
    last=nxt[p][c];
}
\end{lstlisting}
\note{按 len 基数排序，得到 len 单调不减的拓扑序 id[\allowbreak{}0.\allowbreak{}.\allowbreak{}tot]\allowbreak{}，再沿 link 累加出现次数 len 最小是奇根的 -\allowbreak{}1，统一 +\allowbreak{}1 后落在 0.\allowbreak{}.\allowbreak{}tot+\allowbreak{}1，所以桶和 id 都开到 tot+\allowbreak{}1 所有字符加入后调用一次；累加后不要继续 extend/\allowbreak{}build。}
\begin{lstlisting}
void pam_build()
{
    int mx=tot+1;// 桶上界：len[i]+1 最大为 tot+1
    for(int i=0;i<=mx;i++)cnt[i]=0;
    for(int i=0;i<=tot;i++)cnt[len[i]+1]++;
    for(int i=1;i<=mx;i++)cnt[i]+=cnt[i-1];
    for(int i=tot;i>=0;i--)id[--cnt[len[i]+1]]=i;// 先减后放，保证 id[0] 也被填上
    for(int i=tot;i>=0;i--)
    {
        int u=id[i];
        if(u>=2)siz[link[u]]+=siz[u];// 只有真结点（非两个根）才沿 link 向下累加
    }
}
\end{lstlisting}
\note{本质不同回文子串个数：去掉两个根}
\begin{lstlisting}
int count_pal(){return tot-1;}
\end{lstlisting}
\note{最长回文子串长度：所有结点 len 的最大值}
\begin{lstlisting}
int longest_pal()
{
    int ans=0;
    for(int i=2;i<=tot;i++)ans=max(ans,len[i]);
    return ans;
}
\end{lstlisting}
% SOURCE: 03-字符串/最小表示法.cpp
\subsection{最小表示法}

\note{s为循环串；返回最小循环移位的起点，下标0起。比较通过取模访问，s本身不必复制成两倍；输出时从返回位置循环取字符。}
\begin{lstlisting}
const int N=100005;
int n,m;
string s;
\end{lstlisting}
\note{最小表示法：求长度为 n 的串循环同构串中字典序最小的那个的起始下标 返回原串下标（0 起），比较 s[\allowbreak{}i+\allowbreak{}k]\allowbreak{} 与 s[\allowbreak{}j+\allowbreak{}k]\allowbreak{}，O(\allowbreak{}n)\allowbreak{}}
\begin{lstlisting}
int min_show(const string &t)
{
    int len=t.length();
    if(len==0)return 0;
    int i=0,j=1,k=0;// i,j 为两个候选起点，k 为已匹配长度
    while(i<len&&j<len&&k<len)
    {
        int a=(unsigned char)t[(i+k)%len],b=(unsigned char)t[(j+k)%len];
        if(a==b)k++;
        else
        {
            if(a>b)i=i+k+1;// i 处更差，跳到 k+1 之后
            else j=j+k+1;
            if(i==j)j++;
            k=0;
        }
    }
    return min(i,j);
}
\end{lstlisting}
\note{返回最小表示的串本身}
\begin{lstlisting}
string min_str(const string &t)
{
    int p=min_show(t),len=t.length();
    string r="";
    for(int i=0;i<len;i++)r+=t[(p+i)%len];
    return r;
}
\end{lstlisting}
\note{最大表示法：把比较方向反过来即可}
\begin{lstlisting}
int max_show(const string &t)
{
    int len=t.length();
    if(len==0)return 0;
    int i=0,j=1,k=0;
    while(i<len&&j<len&&k<len)
    {
        int a=(unsigned char)t[(i+k)%len],b=(unsigned char)t[(j+k)%len];
        if(a==b)k++;
        else
        {
            if(a<b)i=i+k+1;
            else j=j+k+1;
            if(i==j)j++;
            k=0;
        }
    }
    return min(i,j);
}
\end{lstlisting}
% SOURCE: 03-字符串/序列自动机.cpp
\subsection{序列自动机}

\note{s为原串，nxt[\allowbreak{}pos]\allowbreak{}[\allowbreak{}c]\allowbreak{}表示从pos开始（含pos）的首个字符c位置，s.\allowbreak{}size(\allowbreak{})\allowbreak{}表示不存在。匹配后从找到位置+\allowbreak{}1继续跳转，换串重新建立，字符范围为小写字母。}
\begin{lstlisting}
const int N=100005;
int n,m;
string s;
\end{lstlisting}
\note{小写字母、0 下标：nxt[\allowbreak{}i]\allowbreak{}[\allowbreak{}c]\allowbreak{} 是从 i 开始（含 i）首个 c 的位置，s.\allowbreak{}size(\allowbreak{})\allowbreak{} 表示不存在。 从后往前递推，O(\allowbreak{}26n)\allowbreak{} 预处理}
\begin{lstlisting}
int nxt[N][26];
void build_seq(const string &t)
{
    s=t;
    int len=t.length();
    for(int c=0;c<26;c++)nxt[len][c]=len;// 末尾之后什么都没有，用 len 当哨兵
    for(int i=len-1;i>=0;i--)
    {
        for(int c=0;c<26;c++)nxt[i][c]=nxt[i+1][c];
        nxt[i][t[i]-'a']=i;
    }
}
\end{lstlisting}
\note{判断 t 是否为 s 的子序列，O(\allowbreak{}|t|)\allowbreak{}；匹配失败会停在下标 len 处}
\begin{lstlisting}
bool is_subseq(const string &t)
{
    int p=0,len=s.length();
    for(int i=0;i<(int)t.length();i++)
    {
        int c=t[i]-'a';
        if(p>len||nxt[p][c]>=len)return false;
        p=nxt[p][c]+1;// 跳到下一个位置之后继续找
    }
    return true;
}
\end{lstlisting}
\note{本质不同子序列个数（含空串），O(\allowbreak{}26n)\allowbreak{} dp[\allowbreak{}i]\allowbreak{} 表示前 i 个字符能形成的不同子序列个数，重复的用上一次出现位置减掉 注意：个数是指数级的，字符串稍长就会溢出，题目要求取模时给 dp 全程加 \%mod 即可}
\begin{lstlisting}
ll count_diff_subseq()
{
    int len=s.length();
    int last[26];
    for(int c=0;c<26;c++)last[c]=-1;
    ll dp[N];
    dp[0]=1;// 空串
    for(int i=1;i<=len;i++)
    {
        int c=s[i-1]-'a';
        dp[i]=dp[i-1]*2;
        if(last[c]!=-1)dp[i]-=dp[last[c]];// 减去上一次以 c 结尾造成的重复
        last[c]=i-1;
    }
    return dp[len];
}
\end{lstlisting}
% SOURCE: 03-字符串/Lyndon分解Duval.cpp
\subsection{Lyndon分解Duval}

\note{s是0起输入串，lyndon\_factorization(\allowbreak{})\allowbreak{}返回分解后的半开区间[\allowbreak{}l,\allowbreak{}r)\allowbreak{}，按顺序拼接还原s；输入不变。字符按unsigned char比较，空串返回空结果。}
\note{Duval：O(\allowbreak{}n)\allowbreak{} 把字符串拆成字典序非增的 Lyndon 串，返回 0-\allowbreak{}indexed 半开区间。 Lyndon 串严格小于自身任何非平凡循环移位；用于周期/\allowbreak{}最小表示、Lyndon 相关分段。 按 unsigned char 比较，空串返回空；最小表示法已有模板，此处不重复。 应用：各前缀最小后缀、周期分段；前缀最小后缀需跟踪 Duval 扫描状态，不能只取分解首块。}
\begin{lstlisting}
string s;
vector<pair<int,int>> lyndon_factorization()
{
    int n=s.size(),i=0;vector<pair<int,int>> ans;
    while(i<n){int j=i+1,k=i;while(j<n&&(unsigned char)s[k]<=(unsigned char)s[j]){if((unsigned char)s[k]<(unsigned char)s[j])k=i;else k++;j++;}int len=j-k;while(i<=k)ans.push_back({i,i+len}),i+=len;}
    return ans;
}
\end{lstlisting}
% SOURCE: 03-字符串/有限状态自动机DP.cpp
\subsection{有限状态自动机DP}

\note{ac存模式自动机及状态转移，cnt是计数示例的结果；状态代表已匹配的前缀，终止标记决定禁串/\allowbreak{}包含串条件。换模式集合重建ac；字符串长度与模数按对应计数入口设置。}
\begin{lstlisting}
const int N=26;
\end{lstlisting}
\note{字母表为连续小写字母 a.\allowbreak{}.\allowbreak{}a+\allowbreak{}sig-\allowbreak{}1；根为 1，0 不使用。 禁止串 AC 自动机 +\allowbreak{} 计数 DP；建图 O(\allowbreak{}总串长+\allowbreak{}状态数*sig)\allowbreak{}。 init -\allowbreak{}>\allowbreak{} insert 所有模式 -\allowbreak{}>\allowbreak{} build -\allowbreak{}>\allowbreak{} count；build 后不可再插入。}
\begin{lstlisting}
struct AutomatonDP
{
    struct Node
    {
        int go[N],nxt;
        bool bad;
        Node()
        {
            memset(go,0,sizeof(go));
            nxt=0,bad=false;
        }
    };
    vector<Node> tr;
    int sig;
    bool built;
    void init(int s)
    {
        if(s<1||s>N)throw invalid_argument("alphabet");
        sig=s,built=false;
        tr.assign(2,Node());
    }
    // 插入空模式会把根标坏，表示包括空串在内的所有串都被禁止。
    void insert(const string &s)
    {
        if(built)throw logic_error("insert after build");
        for(int i=0;i<(int)s.length();i++)
            if(s[i]<'a'||s[i]>='a'+sig)throw invalid_argument("character");
        int p=1;
        for(int i=0;i<(int)s.length();i++)
        {
            int c=s[i]-'a';
            if(!tr[p].go[c])
            {
                int v=tr.size();
                tr[p].go[c]=v;
                tr.push_back(Node());
            }
            p=tr[p].go[c];
        }
        tr[p].bad=true;//重复禁串不影响合法性
    }
    // BFS 补全转移，并继承 fail 链上的禁串信息；重复调用不重新建图。
    void build()
    {
        if(built)return;
        queue<int> q;
        tr[1].nxt=1;
        for(int c=0;c<sig;c++)
        {
            int v=tr[1].go[c];
            if(v)tr[v].nxt=1,q.push(v);
            else tr[1].go[c]=1;
        }
        while(!q.empty())
        {
            int u=q.front();
            q.pop();
            tr[u].bad=tr[u].bad||tr[tr[u].nxt].bad;
            for(int c=0;c<sig;c++)
            {
                int v=tr[u].go[c];
                if(v)tr[v].nxt=tr[tr[u].nxt].go[c],q.push(v);
                else tr[u].go[c]=tr[tr[u].nxt].go[c];
            }
        }
        built=true;
    }
    // O(n*状态数*sig)，滚动数组空间 O(状态数)，返回模 mod 的合法串数。
    // mod 必须为正；用 __int128 加法，支持正 ll 范围内的模数。
    ll count(int n,ll mod)
    {
        if(!built)throw logic_error("build first");
        if(n<0||mod<=0)throw invalid_argument("length or modulus");
        if(tr[1].bad)return 0;
        int tot=tr.size();
        vector<ll> f(tot),g(tot);
        f[1]=1%mod;
        for(int i=1;i<=n;i++)
        {
            fill(g.begin(),g.end(),0);
            for(int u=1;u<tot;u++)
            {
                if(!f[u]||tr[u].bad)continue;
                for(int c=0;c<sig;c++)
                {
                    int v=tr[u].go[c];
                    if(!tr[v].bad)g[v]=(ll)(((__int128)g[v]+f[u])%mod);
                }
            }
            f.swap(g);
        }
        ll ans=0;
        for(int u=1;u<tot;u++)ans=(ll)(((__int128)ans+f[u])%mod);
        return ans;
    }
};
AutomatonDP ac;
ll cnt=0;
mt19937 rnd(19260817);
\end{lstlisting}
% SOURCE: 03-字符串/广义后缀自动机.cpp
\subsection{广义后缀自动机}

\note{对象t存len、link及next，根为0；每次insert从根开始，区分多条原串。contains判断是否为任一原串的子串；克隆复制转移与链接，不将不同串直接拼成跨边界子串。}
\note{多串子串集合的 SAM，小写字母；每条串从根重新插入，不引入跨串子串。 已有转移也可能需要 clone，不能把普通 extend 原样 last=\allowbreak{}根 就用。 状态数 O(\allowbreak{}总长)\allowbreak{}，26 固定字母表；支持子串判定及不同子串数，不统计各串出现次数。}
\begin{lstlisting}
struct GeneralSAM
{
    struct Node{array<int,26> next{};int len=0,link=-1;};vector<Node> t={{}};
    int clone(int q,int len){Node v=t[q];v.len=len;t.push_back(v);return t.size()-1;}
    int extend(int last,int c)
    {
        int p=last,q=t[p].next[c];
        if(q)
        {
            if(t[p].len+1==t[q].len)return q;
            int z=clone(q,t[p].len+1);while(p>=0&&t[p].next[c]==q)t[p].next[c]=z,p=t[p].link;t[q].link=z;return z;
        }
        int cur=t.size();t.push_back({});t[cur].len=t[last].len+1;
        while(p>=0&&!t[p].next[c])t[p].next[c]=cur,p=t[p].link;
        if(p<0)t[cur].link=0;
        else
        {
            q=t[p].next[c];if(t[p].len+1==t[q].len)t[cur].link=q;
            else{int z=clone(q,t[p].len+1);while(p>=0&&t[p].next[c]==q)t[p].next[c]=z,p=t[p].link;t[q].link=t[cur].link=z;}
        }
        return cur;
    }
    void insert(const string &s){int last=0;for(char c:s){assert(c>='a'&&c<='z');last=extend(last,c-'a');}}
    bool contains(const string &s)const{int p=0;for(char c:s){if(c<'a'||c>'z')return false;p=t[p].next[c-'a'];if(!p)return false;}return true;}
    long long distinct()const{long long ans=0;for(int i=1;i<(int)t.size();i++)ans+=t[i].len-t[t[i].link].len;return ans;}
};
\end{lstlisting}
% SOURCE: 03-字符串/串联重复分组枚举.cpp
\subsection{串联重复分组枚举}

\note{s是0起输入串，tandem\_repeats(\allowbreak{})\allowbreak{}返回\{p,\allowbreak{}lo,\allowbreak{}hi\}：起点在闭区间lo.\allowbreak{}.\allowbreak{}hi时，连续两个长度p的块相同。p未必是最小周期；StringLCE分别维护正向和反向串，因此保留构造参数区分两个独立实例。}
\note{枚举 tandem repeat（相邻两个相同长度块），以 \{周期p,\allowbreak{}起点闭区间lo.\allowbreak{}.\allowbreak{}hi\} 压缩输出。 用 SA+\allowbreak{}LCP/\allowbreak{}RMQ 的 LCE 分组法覆盖 Main–Lorentz 的重复枚举用途，不枚举所有 O(\allowbreak{}n\ensuremath{^2})\allowbreak{} 个答案。 每个 (\allowbreak{}p,\allowbreak{}start)\allowbreak{} 恰出现一次，枚举 O(\allowbreak{}n log n)\allowbreak{}，本版 SA 排序构建 O(\allowbreak{}n log\ensuremath{^2}n)\allowbreak{}，空间 O(\allowbreak{}n log n)\allowbreak{}。 period 是这次两个块的长度，未必是字符串最小周期；输入任意 byte，不含隐式终止符。}
\begin{lstlisting}
struct StringLCE
{
    int n;vector<int>rank,lg;vector<vector<int>>st;
    StringLCE(const string&s):n(s.size()),rank(n),lg(n+1)
    {
        vector<int>sa(n),next(n),lcp(n);iota(sa.begin(),sa.end(),0);for(int i=0;i<n;i++)rank[i]=(unsigned char)s[i];
        for(int len=1;len<n;len*=2){auto key=[&](int i){return pair(rank[i],i+len<n?rank[i+len]:-1);};sort(sa.begin(),sa.end(),[&](int i,int j){return key(i)<key(j);});next[sa[0]]=0;for(int i=1;i<n;i++)next[sa[i]]=next[sa[i-1]]+(key(sa[i])!=key(sa[i-1]));rank=next;if(rank[sa.back()]==n-1)break;}
        if(n==1)rank[0]=0;
        int h=0;for(int i=0;i<n;i++){int r=rank[i];if(!r){h=0;continue;}int j=sa[r-1];while(i+h<n&&j+h<n&&s[i+h]==s[j+h])h++;lcp[r]=h;if(h)h--;}
        for(int i=2;i<=n;i++)lg[i]=lg[i/2]+1;st.push_back(lcp);for(int k=1;(1<<k)<=n;k++){st.push_back(vector<int>(n));for(int i=0;i+(1<<k)<=n;i++)st[k][i]=min(st[k-1][i],st[k-1][i+(1<<(k-1))]);}
    }
    int query(int i,int j)const{if(i<0||j<0||i>=n||j>=n)return 0;if(i==j)return n-i;int a=rank[i],b=rank[j];if(a>b)swap(a,b);int k=lg[b-a];return min(st[k][a+1],st[k][b-(1<<k)+1]);}
};
string s;
vector<array<int,3>> tandem_repeats()
{
    int n=s.size();StringLCE forward(s);string rev=s;reverse(rev.begin(),rev.end());StringLCE backward(rev);vector<array<int,3>>ans;
    for(int p=1;p*2<=n;p++)for(int i=0;i+p<n;i+=p)
    {int left=i?backward.query(n-i,n-i-p):0,right=forward.query(i,i+p);int l=max({0,i-left,i-p+1}),r=min({i,i+right-p,n-2*p});if(l<=r)ans.push_back({p,l,r});}
    return ans;
}
\end{lstlisting}
% SOURCE: 03-字符串/后缀树Ukkonen.cpp
\subsection{后缀树Ukkonen}

\note{对象保存原串、节点边区间、后缀链接及活动点；边以原串半开区间引用，避免复制子串。输入小写字母，内部终止字符为'\{'；contains查询模式，换串重新构造。}
\note{Ukkonen 后缀树，完整字符串构建 O(\allowbreak{}27n)\allowbreak{}，小写字母，内部追加唯一终止符 '\{'。 边为 text[\allowbreak{}l[\allowbreak{}v]\allowbreak{}.\allowbreak{}.\allowbreak{}r[\allowbreak{}v]\allowbreak{})\allowbreak{}，root=\allowbreak{}0，节点 1 为辅助根；parent/\allowbreak{}link 是父亲/\allowbreak{}后缀链接。 比 SAM/\allowbreak{}SA 更适合需要压缩后缀 trie 边的题。动态 vector，最多约 2n+\allowbreak{}3 节点。 构建核心参考 KACTL SuffixTree.\allowbreak{}h /\allowbreak{} e-\allowbreak{}maxx Ukkonen；原始接口已改成动态存储。}
\begin{lstlisting}
struct SuffixTree
{
    string text;int original,v=0,q=0,nodes=2;
    vector<array<int,27>> next;vector<int> l,r,parent,link;
    int code(char c)const{return c=='{'?26:c-'a';}
    void append(int i,int c)
    {
        again:
        if(r[v]<=q){if(next[v][c]<0){next[v][c]=nodes;l[nodes]=i;parent[nodes++]=v;v=link[v];q=r[v];goto again;}v=next[v][c];q=l[v];}
        if(q==-1||c==code(text[q]))q++;
        else
        {
            l[nodes+1]=i;parent[nodes+1]=nodes;l[nodes]=l[v];r[nodes]=q;parent[nodes]=parent[v];next[nodes][c]=nodes+1;next[nodes][code(text[q])]=v;
            l[v]=q;parent[v]=nodes;next[parent[nodes]][code(text[l[nodes]])]=nodes;
            v=link[parent[nodes]];q=l[nodes];while(q<r[nodes]){v=next[v][code(text[q])];q+=r[v]-l[v];}
            link[nodes]=q==r[nodes]?v:nodes+2;q=r[v]-(q-r[nodes]);nodes+=2;goto again;
        }
    }
    SuffixTree(string s):text(s+"{"),original(s.size()),next(2*text.size()+3),l(next.size()),r(next.size(),text.size()),parent(next.size()),link(next.size())
    {
        for(char c:s)assert(c>='a'&&c<='z');for(auto &a:next)a.fill(-1);next[1].fill(0);
        link[0]=1;l[0]=l[1]=-1;r[0]=r[1]=0;for(int i=0;i<(int)text.size();i++)append(i,code(text[i]));
    }
    bool contains(const string &s)const
    {
        int u=0,i=0;while(i<(int)s.size()){if(s[i]<'a'||s[i]>'z')return false;u=next[u][code(s[i])];if(u<0)return false;for(int j=l[u];j<r[u]&&i<(int)s.size();j++,i++)if(text[j]!=s[i])return false;}return true;
    }
};
\end{lstlisting}
\section{图论}
% SOURCE: 04-图论/存图(vector邻接表).cpp
\subsection{存图(vector邻接表)}

\note{adj[\allowbreak{}u]\allowbreak{}存Edge\{to,\allowbreak{}w\}，n/\allowbreak{}m为点边数，点1.\allowbreak{}.\allowbreak{}n；有向边存一次，无向边存两次。重建时清空各邻接向量；需要区分重边、反向边时另加边编号，不能只按端点判断。}
\begin{lstlisting}
const int N=100005;
int n,m;
struct Edge{int to,w;};
vector<Edge> adj[N];
void add_edge(int u,int v,int w){adj[u].push_back({v,w});}
void add_undirected(int u,int v,int w){add_edge(u,v,w);add_edge(v,u,w);}
\end{lstlisting}
\note{遍历：for(\allowbreak{}auto [\allowbreak{}v,\allowbreak{}w]\allowbreak{}:adj[\allowbreak{}u]\allowbreak{})\allowbreak{}；重建：for(\allowbreak{}int u=\allowbreak{}1;u<\allowbreak{}=\allowbreak{}n;u+\allowbreak{}+\allowbreak{})\allowbreak{}adj[\allowbreak{}u]\allowbreak{}.\allowbreak{}clear(\allowbreak{})\allowbreak{}; 无向边存两次；要区分重边/\allowbreak{}反向边时在 Edge 里增加 id。}
% SOURCE: 04-图论/Dijkstra.cpp
\subsection{Dijkstra}

\note{adj[\allowbreak{}u]\allowbreak{}存出边，dis为源点s到各点的距离，vis为最短距离已确定标记，n/\allowbreak{}m为规模。点1.\allowbreak{}.\allowbreak{}n，边权非负；每次运行重置距离和标记，换图清空adj。}
\begin{lstlisting}
const int N=100005;
const ll INF=1e18;// 距离上界，加法前判 INF 防溢出
int n,m,s;
struct Edge
{
    int to,w;
};
vector<Edge> adj[N];
ll dis[N];
bool vis[N];
void add_edge(int u,int v,int c)
{
    adj[u].push_back({v,c});
}
void add_undirected(int u,int v,int c)
{
    add_edge(u,v,c);
    add_edge(v,u,c);
}
struct Node// 堆里的小根堆结点
{
    int u;
    ll d;
    bool operator>(const Node &b)const
    {
        return d>b.d;
    }
};
void dijkstra_naive(int s)// 朴素 O(n^2+m)，稠密图更稳
{
    for(int i=1;i<=n;i++)dis[i]=INF,vis[i]=0;
    dis[s]=0;
    for(int i=1;i<=n;i++)
    {
        int u=0;
        for(int j=1;j<=n;j++)
            if(!vis[j]&&(u==0||dis[j]<dis[u]))u=j;
        if(u==0||dis[u]==INF)break;// 剩下的点都不可达
        vis[u]=1;
        for(int k=0;k<(int)adj[u].size();k++)
        {
            int v=adj[u][k].to,c=adj[u][k].w;
            if(dis[u]+c<dis[v])dis[v]=dis[u]+c;
        }
    }
}
void dijkstra_heap(int s)// 堆优化 O((n+m) log n)，只适合非负权
{
    for(int i=1;i<=n;i++)dis[i]=INF,vis[i]=0;
    dis[s]=0;
    priority_queue<Node,vector<Node>,greater<Node> > pq;
    Node st;
    st.u=s,st.d=0;
    pq.push(st);
    while(!pq.empty())
    {
        int u=pq.top().u;
        pq.pop();
        if(vis[u])continue;// vis 判重，每个点只出堆一次
        vis[u]=1;
        for(int k=0;k<(int)adj[u].size();k++)
        {
            int v=adj[u][k].to,c=adj[u][k].w;
            if(dis[u]+c<dis[v])
            {
                dis[v]=dis[u]+c;
                Node tmp;
                tmp.u=v,tmp.d=dis[v];
                pq.push(tmp);
            }
        }
    }
}
\end{lstlisting}
% SOURCE: 04-图论/01BFS.cpp
\subsection{01BFS}

\note{adj[\allowbreak{}u]\allowbreak{}存\{v,\allowbreak{}w\}，w只能为0/\allowbreak{}1；d为最短距离，q保存距离与节点，入口bfs01(\allowbreak{}s)\allowbreak{}重置二者。点编号为adj的合法下标，支持0起；换图重新设置adj，不可达为INF。}
\begin{lstlisting}
const int INF=INT_MAX/4;
vector<vector<pair<int,int>>> adj;
vector<int> d;
deque<pair<int,int>> q;
\end{lstlisting}
\note{adj[\allowbreak{}u]\allowbreak{} 存 \{v,\allowbreak{}w\}，只允许 w=\allowbreak{}0/\allowbreak{}1；点编号取 adj 的合法下标，支持 0 起。 松弛成功时：0 边入队首，1 边入队尾；不能用普通 BFS 首次访问就锁定。 O(\allowbreak{}n+\allowbreak{}m)\allowbreak{}，不可达为 INF；队列带入队距离，跳过过期项。}
\begin{lstlisting}
void bfs01(int s)
{
    int n=adj.size();
    assert(0<=s&&s<n);
    d.assign(n,INF);q.clear();
    d[s]=0,q.push_front({0,s});
    while(!q.empty())
    {
        auto [dist,u]=q.front();q.pop_front();
        if(dist!=d[u])continue;
        for(auto [v,w]:adj[u])
        {
            assert(0<=v&&v<n&&(w==0||w==1));
            if(d[v]<=dist+w)continue;
            d[v]=dist+w;
            if(w==0)q.push_front({d[v],v});
            else q.push_back({d[v],v});
        }
    }
}
\end{lstlisting}
% SOURCE: 04-图论/Floyd.cpp
\subsection{Floyd}

\note{d[\allowbreak{}i]\allowbreak{}[\allowbreak{}j]\allowbreak{}初始为边权，无边INF、对角0，重边取最小；运行后变为全源最短距离。reach是布尔可达矩阵，供闭包例子使用。中间点k必须在最外层，负环存在时普通最短距离不再有定义。}
\begin{lstlisting}
const int N=505;
const ll INF=1e18;
int n,m;
ll d[N][N];// 邻接矩阵，无边为 INF，d[i][i]=0
bool reach[N][N];// 传递闭包：i 能否到 j
void floyd()// 全源最短路 O(n^3)，允许负权边，不能有负环
{
    for(int k=1;k<=n;k++)
        for(int i=1;i<=n;i++)
            for(int j=1;j<=n;j++)
                if(d[i][k]<INF&&d[k][j]<INF&&d[i][k]+d[k][j]<d[i][j])d[i][j]=d[i][k]+d[k][j];
}
void closure()// 传递闭包 O(n^3)
{
    for(int k=1;k<=n;k++)
        for(int i=1;i<=n;i++)
            for(int j=1;j<=n;j++)
                if(reach[i][k]&&reach[k][j])reach[i][j]=1;
}
\end{lstlisting}
% SOURCE: 04-图论/Bellman-Ford.cpp
\subsection{Bellman-Ford}

\note{e[\allowbreak{}1.\allowbreak{}.\allowbreak{}m]\allowbreak{}为边表，dis是源点s的距离，n为点数。遍历整张边表松弛，不可达为INF；负环判定只覆盖源点可达部分，检测任意负环需让所有点初始可达。}
\begin{lstlisting}
const int N=100005;
const ll INF=1e18;
int n,m,s;
struct Edge
{
    int u,v,w;
};
Edge e[N];// 边表，Bellman-Ford 只需要能枚举所有边
ll dis[N];
int bellman_ford(int s)// O(n*m)：n-1 轮松弛 + 第 n 轮判负环
{
    for(int i=1;i<=n;i++)dis[i]=INF;
    dis[s]=0;
    int flag=0;
    for(int i=1;i<=n;i++)
    {
        flag=0;
        for(int j=1;j<=m;j++)
            if(dis[e[j].u]<INF&&dis[e[j].u]+e[j].w<dis[e[j].v])
            {
                dis[e[j].v]=dis[e[j].u]+e[j].w;
                flag=1;
            }
        if(!flag)break;// 一整轮没松弛，提前结束
        if(i==n&&flag)return 1;// 第 n 轮还能松弛 -> 有可达负环
    }
    return 0;
}
\end{lstlisting}
% SOURCE: 04-图论/SPFA.cpp
\subsection{SPFA}

\note{adj为邻接表，dis为距离，inq表示是否已在队列，cnt为当前松弛路径边数。cnt>\allowbreak{}=\allowbreak{}n判源点可达负环；换图清邻接，运行前重置队列状态，最坏时间仍为O(\allowbreak{}nm)\allowbreak{}。}
\begin{lstlisting}
const int N=100005;
const ll INF=1e18;
int n,m,s;
struct Edge
{
    int to,w;
};
vector<Edge> adj[N];
ll dis[N];
int cnt[N];// cnt[u]：当前松弛路径的边数，>=n 说明有可达负环
bool inq[N];
int spfa(int s)// 最坏 O(n*m)；返回 1 表示存在可达负环
{
    for(int i=1;i<=n;i++)dis[i]=INF,cnt[i]=0,inq[i]=0;
    queue<int> q;
    dis[s]=0,inq[s]=1,cnt[s]=0;
    q.push(s);
    while(!q.empty())
    {
        int u=q.front();
        q.pop();
        inq[u]=0;
        for(int i=0;i<(int)adj[u].size();i++)
        {
            int v=adj[u][i].to,c=adj[u][i].w;
            if(dis[u]+c<dis[v])
            {
                dis[v]=dis[u]+c;
                cnt[v]=cnt[u]+1;
                if(cnt[v]>=n)return 1;
                if(!inq[v])
                {
                    q.push(v);
                    inq[v]=1;
                }
            }
        }
    }
    return 0;
}
\end{lstlisting}
% SOURCE: 04-图论/拓扑排序.cpp
\subsection{拓扑排序}

\note{adj为有向邻接表，in/\allowbreak{}out为入度/\allowbreak{}出度，q为拓扑队列，dp按示例存路径计数或最长路。点1.\allowbreak{}.\allowbreak{}n；排序使用工作入度，换图重新统计，排出点数不足n表示有环。}
\note{最小堆给字典序最小拓扑序，不能保证任意 DAG 的逆序数最少。 反例边 5-\allowbreak{}>\allowbreak{}1、5-\allowbreak{}>\allowbreak{}2：堆顺序 3,\allowbreak{}4,\allowbreak{}5,\allowbreak{}1,\allowbreak{}2 有 6 个逆序，5,\allowbreak{}1,\allowbreak{}2,\allowbreak{}3,\allowbreak{}4 只有 4 个。 “按约束排列值”的构造可能需要 p[\allowbreak{}topo[\allowbreak{}i]\allowbreak{}]\allowbreak{}=\allowbreak{}i，输出排名数组而非 topo 本身。}
\begin{lstlisting}
const int N=100005;
int n,m;
vector<int> adj[N];
int in[N],out[N];
int dp[N];// dp[u]：从所有入度为 0 的点走到 u 的方案数 / 最长路长度
int q[N];// 手写队列，避免 STL queue 在拓扑里反复 push
int topo(int mode)// mode=0 统计方案数(取模)，mode=1 最长路；返回拓扑点数
{
    int head=0,tail=0;
    for(int i=1;i<=n;i++)if(in[i]==0)q[tail++]=i;
    int cnt=0;
    while(head<tail)
    {
        int u=q[head++];
        cnt++;
        for(int i=0;i<(int)adj[u].size();i++)
        {
            int v=adj[u][i];
            if(mode==0)dp[v]=(dp[v]+dp[u])%80112002;
            else dp[v]=max(dp[v],dp[u]+1);
            if(--in[v]==0)q[tail++]=v;// 入度减到 0 立刻能入队
        }
    }
    if(cnt<n)return -1;// 有点没进队 -> 有环
    return cnt;
}
\end{lstlisting}
% SOURCE: 04-图论/最小生成树Kruskal.cpp
\subsection{最小生成树Kruskal}

\note{e为边表，fa/\allowbreak{}sz为DSU；adj、dis、vis属于Prim替代例。Kruskal排序会改变边表顺序，选不足n-\allowbreak{}1条表示断连；若题目要森林可保留已选边，不能把返回值直接当连通图MST。}
\begin{lstlisting}
const int N=200005;
const ll INF=1e18;
int n,m;
struct Edge
{
    int u,v,w;
};
Edge e[N];
int fa[N],sz[N];
bool cmp(Edge x,Edge y)
{
    return x.w<y.w;
}
int findd(int x)
{
    if(fa[x]!=x)fa[x]=findd(fa[x]);
    return fa[x];
}
void unionn(int x,int y)
{
    x=findd(x),y=findd(y);
    if(x==y)return;
    if(sz[x]<sz[y])swap(x,y);
    fa[y]=x,sz[x]+=sz[y];
}
ll kruskal()// O(m log m)，返回最小生成树边权和；不连通返回 -1
{
    for(int i=1;i<=n;i++)fa[i]=i,sz[i]=1;
    sort(e+1,e+m+1,cmp);
    ll sum=0;
    int cnt=0;
    for(int i=1;i<=m;i++)
    {
        int u=findd(e[i].u),v=findd(e[i].v);
        if(u==v)continue;// 两端已连通，这条边舍掉
        unionn(u,v);
        sum+=e[i].w;
        cnt++;
        if(cnt==n-1)break;// 已经选了 n-1 条边，提前结束
    }
    if(cnt<n-1)return -1;
    return sum;
}
\end{lstlisting}
\note{-\allowbreak{}-\allowbreak{}-\allowbreak{}-\allowbreak{}-\allowbreak{}-\allowbreak{}-\allowbreak{}-\allowbreak{}-\allowbreak{}-\allowbreak{} 附：Prim 堆优化版 O(\allowbreak{}(\allowbreak{}n+\allowbreak{}m)\allowbreak{} log n)\allowbreak{}，稠密图也可用朴素版 -\allowbreak{}-\allowbreak{}-\allowbreak{}-\allowbreak{}-\allowbreak{}-\allowbreak{}-\allowbreak{}-\allowbreak{}-\allowbreak{}-\allowbreak{}}
\begin{lstlisting}
vector<pair<int,int>> adj[N];
void add_edge(int u,int v,int c)
{
    adj[u].push_back({v,c});
}
bool vis[N];
ll dis[N];
struct Node
{
    int u;
    ll d;
    bool operator>(const Node &b)const
    {
        return d>b.d;
    }
};
ll prim(int s)// 从 s 出发，返回 MST 权值和；不连通返回 -1
{
    for(int i=1;i<=n;i++)dis[i]=INF,vis[i]=0;
    dis[s]=0;
    priority_queue<Node,vector<Node>,greater<Node> > pq;
    Node st;
    st.u=s,st.d=0;
    pq.push(st);
    ll sum=0;
    int cnt=0;
    while(!pq.empty())
    {
        int u=pq.top().u;
        ll d=pq.top().d;
        pq.pop();
        if(vis[u])continue;
        vis[u]=1;
        sum+=d;
        cnt++;
        for(auto [v,w]:adj[u])
            if(w<dis[v])
            {
                dis[v]=w;
                Node tmp;
                tmp.u=v,tmp.d=w;
                pq.push(tmp);
            }
    }
    if(cnt<n)return -1;
    return sum;
}
\end{lstlisting}
% SOURCE: 04-图论/树的遍历与dfs序.cpp
\subsection{树的遍历与dfs序}

\note{adj存邻点，多测清空后重建；in/\allowbreak{}out为子树DFS区间，rnk为DFS序到点的逆映射，par/\allowbreak{}sz为父亲和大小；euler/\allowbreak{}first/\allowbreak{}edep供欧拉序LCA，st/\allowbreak{}elog为RMQ。stk/\allowbreak{}it模拟递归，it是邻接表位置，从0起；不同序列不能混用下标，换树重置timer及首次位置。}
\note{树的遍历与 dfs 序：dfs 序 /\allowbreak{} 欧拉序 /\allowbreak{} 时间戳 三个数组：in[\allowbreak{}u]\allowbreak{},\allowbreak{}out[\allowbreak{}u]\allowbreak{} 是子树对应的时间戳区间；rnk[\allowbreak{}i]\allowbreak{} 是第 i 个时间戳上的点 子树 → 区间：子树 u 恰好是 [\allowbreak{}in[\allowbreak{}u]\allowbreak{},\allowbreak{} out[\allowbreak{}u]\allowbreak{}]\allowbreak{}（in 就是 dfn） euler 数组是长度 2n-\allowbreak{}1 的欧拉序（每步都记，配合 ST 表可 O(\allowbreak{}1)\allowbreak{} 求 LCA） dfs 全部写成迭代版（显式栈），n=\allowbreak{}1e5 的链也不会爆栈}
\begin{lstlisting}
const int N=100005;
int n,root;
vector<int> adj[N];
int val[N];
int in[N],out[N],rnk[N],par[N],sz[N],timer_;
int euler[N<<1],first[N],edep[N<<1],elog[N<<1],st[20][N<<1];// 欧拉序 + ST 表
int stk[N],it[N];// 显式栈：stk 存点，it 存下一邻点的位置，从 0 起（模拟递归）
void add_edge(int u,int v)
{
    adj[u].push_back(v);
    adj[v].push_back(u);
}
\end{lstlisting}
\note{迭代 dfs 求 dfs 序 +\allowbreak{} 时间戳 +\allowbreak{} 子树大小 +\allowbreak{} 欧拉序，O(\allowbreak{}n)\allowbreak{}}
\begin{lstlisting}
void get_dfn(int rt)
{
    timer_=0;
    int tp=0,cnt=0;
    stk[++tp]=rt,it[tp]=0,par[rt]=0,sz[rt]=1;
    in[rt]=++timer_,rnk[timer_]=rt;
    euler[++cnt]=rt,first[rt]=cnt,edep[cnt]=0;
    while(tp)
    {
        int u=stk[tp];
        if(it[tp]<(int)adj[u].size())
        {
            int v=adj[u][it[tp]++];// 回溯时接着走下一个邻点
            if(v==par[u])continue;
            par[v]=u,sz[v]=1;
            in[v]=++timer_,rnk[timer_]=v;
            euler[++cnt]=v,first[v]=cnt,edep[cnt]=edep[first[u]]+1;
            stk[++tp]=v,it[tp]=0;
        }
        else
        {
            out[u]=timer_;
            tp--;
            if(tp)
            {
                sz[stk[tp]]+=sz[u];// 回到父亲，把大小并上去
                euler[++cnt]=stk[tp],edep[cnt]=edep[first[stk[tp]]];// 欧拉序回退也记一次
            }
        }
    }
    for(int i=2;i<=cnt;i++)elog[i]=elog[i>>1]+1;// ST 表预处理的 log
    for(int i=1;i<=cnt;i++)st[0][i]=i;
    for(int k=1;k<20;k++)
        for(int i=1;i+(1<<k)-1<=cnt;i++)
        {
            int a=st[k-1][i],b=st[k-1][i+(1<<(k-1))];
            st[k][i]=edep[a]<edep[b]?a:b;
        }
}
\end{lstlisting}
\note{欧拉序 +\allowbreak{} ST 表求 LCA，O(\allowbreak{}1)\allowbreak{}（建表 O(\allowbreak{}n log n)\allowbreak{}）}
\begin{lstlisting}
int lca(int x,int y)
{
    if(first[x]>first[y])swap(x,y);
    int l=first[x],r=first[y],k=elog[r-l+1];
    int a=st[k][l],b=st[k][r-(1<<k)+1];
    return edep[a]<edep[b]?euler[a]:euler[b];
}
\end{lstlisting}
\note{dfs 序上的树状数组：单点加 +\allowbreak{} 区间和，O(\allowbreak{}log n)\allowbreak{}，用来做「子树求和 /\allowbreak{} 子树加」}
\begin{lstlisting}
struct BIT{
    ll tr[N];
    void add(int x,ll k)
    {
        for(;x<=n;x+=x&(-x))tr[x]+=k;
    }
    ll query(int x)
    {
        ll s=0;
        for(;x>0;x-=x&(-x))s+=tr[x];
        return s;
    }
    ll range(int l,int r){return query(r)-query(l-1);}
}bit;
\end{lstlisting}
% SOURCE: 04-图论/LCA(倍增).cpp
\subsection{LCA(倍增)}

\note{adj存邻点，多测清空后重建；dep为深度，fa[\allowbreak{}u]\allowbreak{}[\allowbreak{}k]\allowbreak{}为2\textasciicircum{}k祖先，root为根，0为父亲哨兵；diff用于边路径差分，汇总后diff[\allowbreak{}u]\allowbreak{}是父边覆盖次数。点1.\allowbreak{}.\allowbreak{}n，倍增层数按最大深度设置；换树重建父亲和深度，根没有父边。}
\begin{lstlisting}
const int N=100005;
int n,m,root;
vector<int> adj[N];
int dep[N],fa[N][20];// fa[u][k]：u 往上跳 2^k 步的祖先
int diff[N];// 树上差分数组
void add_edge(int u,int v)
{
    adj[u].push_back(v);
}
void add_undirected(int u,int v)
{
    add_edge(u,v);
    add_edge(v,u);
}
void dfs(int u,int f)// 预处理深度和倍增表，O(n log n)
{
    dep[u]=dep[f]+1;
    fa[u][0]=f;
    for(int k=1;k<20;k++)fa[u][k]=fa[fa[u][k-1]][k-1];
    for(int v:adj[u])
        if(v!=f)dfs(v,u);
}
int lca(int x,int y)// 倍增求 LCA，O(log n)
{
    if(dep[x]<dep[y])swap(x,y);
    int d=dep[x]-dep[y];
    for(int k=0;k<20;k++)if(d>>k&1)x=fa[x][k];// 先把深的提到同一层
    if(x==y)return x;
    for(int k=19;k>=0;k--)
        if(fa[x][k]!=fa[y][k])x=fa[x][k],y=fa[y][k];
    return fa[x][0];
}
int dist(int x,int y)// 树上两点距离（边权为 1，带权就把 dep 换成前缀和）
{
    int f=lca(x,y);
    return dep[x]+dep[y]-2*dep[f];
}
void add_path(int x,int y)// 路径 (x,y) 差分打标记
{
    int f=lca(x,y);
    diff[x]++,diff[y]++,diff[f]-=2;// 边差分：点上减 2 次
}
void collect(int u,int f)// 自底向上合并差分，得到每条边被覆盖的次数
{
    for(int v:adj[u])
        if(v!=f)
        {
            collect(v,u);
            diff[u]+=diff[v];
        }
}
\end{lstlisting}
% SOURCE: 04-图论/树的重心与直径.cpp
\subsection{树的重心与直径}

\note{adj存邻点，多测清空后重建；que/\allowbreak{}order为BFS队列与顺序，fa/\allowbreak{}sz求子树大小，dis为当前BFS距离；da/\allowbreak{}db为直径端点，dia为边数长度，dist\_db保留从db的距离。无权树可用两次BFS，带权或负权版本不能直接套这份距离递推。}
\note{树的重心、直径、树上最远点 直径两种写法：两遍 bfs /\allowbreak{} 树形 DP；求法都是 O(\allowbreak{}n)\allowbreak{} 重心：删掉它之后最大连通块最小（重心最多两个，这里求编号最小的那个） 最远点：任意点 x 的某个最远点可取直径端点，O(\allowbreak{}1)\allowbreak{}（见 build\_far/\allowbreak{}far\_node/\allowbreak{}eccentricity） 递归写法给 n<\allowbreak{}=\allowbreak{}3e4 用；n=\allowbreak{}1e5 的链请用下面的迭代版（显式栈），否则可能爆栈}
\begin{lstlisting}
const int N=100005;
int n;
vector<int> adj[N];
int dis[N],que[N],order[N],fa[N],sz[N];// que 是 bfs 队列，order 是 bfs 序
int da,db,dia;// 直径两端点与长度
\end{lstlisting}
\note{带权图把 adj 的元素换成 \{to,\allowbreak{}w\}，dis 累加边权即可}
\begin{lstlisting}
void add_edge(int u,int v)
{
    adj[u].push_back(v);
    adj[v].push_back(u);
}
\end{lstlisting}
\note{从 s 出发 bfs，返回最远点的编号；dis 顺手算好，O(\allowbreak{}n)\allowbreak{}，迭代不爆栈}
\begin{lstlisting}
int bfs(int s)
{
    for(int i=1;i<=n;i++)dis[i]=-1;
    int hd=0,tl=0,far=s;
    que[tl++]=s,dis[s]=0;
    while(hd<tl)
    {
        int u=que[hd++];
        if(dis[u]>dis[far])far=u;
        for(int v:adj[u])
        {
            if(dis[v]==-1)dis[v]=dis[u]+1,que[tl++]=v;
        }
    }
    return far;
}
\end{lstlisting}
\note{两遍 bfs：第一次找端点，第二次定长度，O(\allowbreak{}n)\allowbreak{}}
\begin{lstlisting}
int get_diameter()
{
    da=bfs(1);
    db=bfs(da);
    dia=dis[db];
    return dia;
}
\end{lstlisting}
\note{树形 DP 求直径：dp1/\allowbreak{}dp2 记向下最长、次长，O(\allowbreak{}n)\allowbreak{}，迭代不爆栈}
\begin{lstlisting}
int get_diameter_dp()
{
    int hd=0,tl=0;
    for(int i=1;i<=n;i++)fa[i]=0;
    que[tl++]=1,fa[1]=-1,dis[1]=0;
    while(hd<tl)
    {
        int u=que[hd++];
        order[hd]=u;// order[1..tl] 就是 bfs 序
        for(int v:adj[u])
        {
            if(v==fa[u])continue;
            fa[v]=u,dis[v]=dis[u]+1,que[tl++]=v;
        }
    }
    int dp1[N],dp2[N],ans=0;
    for(int i=tl;i>=1;i--)// 逆序保证儿子先算完
    {
        int u=order[i];
        dp1[u]=dp2[u]=0;
        for(int v:adj[u])
        {
            if(v==fa[u])continue;
            int d=dp1[v]+1;
            if(d>dp1[u])dp2[u]=dp1[u],dp1[u]=d;
            else if(d>dp2[u])dp2[u]=d;
        }
        ans=max(ans,dp1[u]+dp2[u]);
    }
    return ans;
}
\end{lstlisting}
\note{求重心：最大子树最小的点，O(\allowbreak{}n)\allowbreak{}，迭代不爆栈}
\begin{lstlisting}
int get_centroid()
{
    int hd=0,tl=0;
    for(int i=1;i<=n;i++)fa[i]=0;
    que[tl++]=1,fa[1]=-1;
    while(hd<tl)
    {
        int u=que[hd++];
        order[hd]=u;
        for(int v:adj[u])
        {
            if(v==fa[u])continue;
            fa[v]=u,que[tl++]=v;
        }
    }
    for(int i=tl;i>=1;i--)
    {
        int u=order[i];
        sz[u]=1;
        for(int v:adj[u])
            if(v!=fa[u])sz[u]+=sz[v];
    }
    int best=1,mx=n+1;
    for(int u=1;u<=n;u++)
    {
        int m=n-sz[u];// 父亲那一边的块
        for(int v:adj[u])
            if(v!=fa[u])m=max(m,sz[v]);
        if(m<mx)mx=m,best=u;
    }
    return best;
}
\end{lstlisting}
\note{树上最远点 /\allowbreak{} 偏心距：max(\allowbreak{}dist(\allowbreak{}x,\allowbreak{}da)\allowbreak{},\allowbreak{}dist(\allowbreak{}x,\allowbreak{}db)\allowbreak{})\allowbreak{}，并列时还可能有其他最远点。 前提：dis 是 bfs(\allowbreak{}da)\allowbreak{} 的结果（dis[\allowbreak{}da]\allowbreak{}=\allowbreak{}0,\allowbreak{} dis[\allowbreak{}db]\allowbreak{}=\allowbreak{}dia） 注意：不能写成 dia-\allowbreak{}dis[\allowbreak{}x]\allowbreak{}！那只在 x 落在直径路径上时才对， 一般点要分别算 d(\allowbreak{}x,\allowbreak{}da)\allowbreak{}=\allowbreak{}dis[\allowbreak{}x]\allowbreak{} 和 d(\allowbreak{}x,\allowbreak{}db)\allowbreak{}，其中 d(\allowbreak{}x,\allowbreak{}db)\allowbreak{} 需要 dist\_db[\allowbreak{}]\allowbreak{}}
\begin{lstlisting}
int dist_db[N];// bfs(db) 的预处理结果，供 d_to_db 使用
void build_far()// O(n)，先调 get_diameter() 再用
{
    bfs(db);
    for(int i=1;i<=n;i++)dist_db[i]=dis[i];
    bfs(da);
}
int far_node(int x)
{
    return dist_db[x]>dis[x]?db:da;// 到 db 更远就选 db
}
int eccentricity(int x)
{
    return max(dis[x],dist_db[x]);// max(d(x,da), d(x,db))
}
\end{lstlisting}
\note{朴素递归版求直径（n 小的时候用，链会爆栈）}
\begin{lstlisting}
int dfs_dia(int u,int f,int &res)
{
    int a=0;// 往下的最长链，次长的不单独存
    for(int v:adj[u])
    {
        if(v==f)continue;
        int d=dfs_dia(v,u,res)+1;
        res=max(res,a+d),a=max(a,d);
    }
    return a;
}
\end{lstlisting}
% SOURCE: 04-图论/Tarjan强连通分量.cpp
\subsection{Tarjan强连通分量}

\note{adj存邻点，多测清空后重建；dfn/\allowbreak{}low为访问序与可回溯序，stk/\allowbreak{}in\_stk维护未出栈点，belong[\allowbreak{}u]\allowbreak{}为分量号，sz为分量大小。dag/\allowbreak{}val/\allowbreak{}dp是后续缩点应用的预留状态，模板不会自动填好；换图重置计数、数组与邻接。}
\note{适用：有向图缩点后做 DAG DP、判互相可达；单点也属于 SCC。 参数：u 为 1.\allowbreak{}.\allowbreak{}n 的顶点；图不连通时对每个未访问点调用 tarjan。 关键：low 是可追溯到的栈内最早时间；已出栈分量不能参与回退。 结论：belong[\allowbreak{}u]\allowbreak{} 是分量编号，sz[\allowbreak{}id]\allowbreak{} 是大小；跨分量边构成 DAG。 易错：多组清空 dfn/\allowbreak{}low/\allowbreak{}in\_stk/\allowbreak{}belong/\allowbreak{}sz、adj/\allowbreak{}dag，timer/\allowbreak{}top/\allowbreak{}scc\_cnt 归零。 复杂度：全图 O(\allowbreak{}n+\allowbreak{}m)\allowbreak{}，空间 O(\allowbreak{}n+\allowbreak{}m)\allowbreak{}；递归长链可能爆栈，val/\allowbreak{}dp 的和注意 int 范围。}
\begin{lstlisting}
const int N=100005;
int n,m;
vector<int> adj[N];
int dfn[N],low[N],stk[N],in_stk[N],belong[N];
int timer=0,top=0,scc_cnt=0;
vector<int> dag[N];// 缩点后的 DAG
int sz[N];// 每个 SCC 的点数
int val[N];// 每个 SCC 的点权和（DP 用）
int dp[N];// 缩点 DAG 上 DP 的值
\end{lstlisting}
\note{O(\allowbreak{}1)\allowbreak{}，加有向 u-\allowbreak{}>\allowbreak{}v 边。}
\begin{lstlisting}
void add_edge(int u,int v)
{
    adj[u].push_back(v);
}
\end{lstlisting}
\note{O(\allowbreak{}1)\allowbreak{}，连加双向边；SCC 问题通常只需 add\_edge，勿误用双向图。}
\begin{lstlisting}
void add_undirected(int u,int v)
{
    add_edge(u,v);
    add_edge(v,u);
}
\end{lstlisting}
\note{全图 O(\allowbreak{}n+\allowbreak{}m)\allowbreak{}，u 是当前入口；low[\allowbreak{}u]\allowbreak{}=\allowbreak{}=\allowbreak{}dfn[\allowbreak{}u]\allowbreak{} 时弹栈直到 u，恰好一个 SCC。}
\begin{lstlisting}
void tarjan(int u)// Tarjan 求 SCC，O(n+m)
{
    dfn[u]=low[u]=++timer;
    stk[++top]=u,in_stk[u]=1;
    for(int v:adj[u])
    {
        if(!dfn[v])
        {
            tarjan(v);
            low[u]=min(low[u],low[v]);
        }
        else if(in_stk[v])low[u]=min(low[u],dfn[v]);// 只回退栈内的点
    }
    if(dfn[u]==low[u])// u 是这个 SCC 的根
    {
        scc_cnt++;
        int x;
        do
        {
            x=stk[top--];
            in_stk[x]=0;
            belong[x]=scc_cnt;
            sz[scc_cnt]++;
        }while(x!=u);
    }
}
\end{lstlisting}
% SOURCE: 04-图论/Tarjan割点与桥.cpp
\subsection{Tarjan割点与桥}

\note{g[\allowbreak{}u]\allowbreak{}存\{邻点,\allowbreak{}有向边编号\}；dfn/\allowbreak{}low为Tarjan序，is\_cut标割点，ea/\allowbreak{}eb列桥，cut\_edge标桥边。num从1起使成对边能用异或1取反向；跳过父边反向编号，不能跳过所有通向父点的边，多测清空边标记和计数。}
\note{适用：无向图删除点或边后的连通性，割点与桥必须分别判断。 参数：顶点 1.\allowbreak{}.\allowbreak{}n；根调用 tarjan(\allowbreak{}u,\allowbreak{}0)\allowbreak{}，fa 是入边编号而非父顶点。 关键：桥要求 low[\allowbreak{}v]\allowbreak{}>\allowbreak{}dfn[\allowbreak{}u]\allowbreak{}，割点允许相等；根必须有至少两棵 DFS 子树。 从 2 开始成对加边，i\textasciicircum{}1 为反向边；重边仅跳过真正的父边。 vector 邻接表保存 \{目标,\allowbreak{}有向边编号\}；重建时 num=\allowbreak{}1，清空 g 与各标记。 DFS 树边用 low[\allowbreak{}v]\allowbreak{}，已访问邻点用 dfn[\allowbreak{}v]\allowbreak{}；否则会把跨割点的返祖信息错误传播。 复杂度：全图 O(\allowbreak{}n+\allowbreak{}m)\allowbreak{}，空间 O(\allowbreak{}n+\allowbreak{}m)\allowbreak{}，递归深度 O(\allowbreak{}n)\allowbreak{}；多组须清空标记与计数。}
\begin{lstlisting}
const int N=100005;
int n,m;
vector<pair<int,int>> g[N];
int num=1;
int dfn[N],low[N];
int is_cut[N];// 是否为割点
int timer=0;
int ea[N],eb[N],ecnt=0;// 桥的列表
vector<int> cut_edge(2);// 按有向边编号标记桥
\end{lstlisting}
\note{O(\allowbreak{}1)\allowbreak{}，加一个方向的邻接边；编号由 num 决定，不能单独打乱配对。}
\begin{lstlisting}
void add_edge(int u,int v)
{
    g[u].push_back({v,++num});cut_edge.push_back(0);
}
\end{lstlisting}
\note{O(\allowbreak{}1)\allowbreak{}，按顺序加两个方向；注意现有边编号与异或反边约定的风险。}
\begin{lstlisting}
void add_undirected(int u,int v)
{
    add_edge(u,v);
    add_edge(v,u);
}
\end{lstlisting}
\note{全图 O(\allowbreak{}n+\allowbreak{}m)\allowbreak{}，跳过入边的反边以保留重边回边；结果在 is\_cut、ea/\allowbreak{}eb、cut\_edge。}
\begin{lstlisting}
void tarjan(int u,int fa)// 无向图求割点与桥，O(n+m)；fa 是边的入边编号，避免走回父亲
{
    dfn[u]=low[u]=++timer;
    int child=0;
    for(auto [v,i]:g[u])
    {
        if(!dfn[v])
        {
            child++;
            tarjan(v,i);
            low[u]=min(low[u],low[v]);
            if(low[v]>dfn[u])// 子树回不到 u 及更早 -> 这条边是桥
            {
                ecnt++;
                ea[ecnt]=u,eb[ecnt]=v;
                cut_edge[i]=cut_edge[i^1]=1;
            }
            if(fa&&low[v]>=dfn[u])is_cut[u]=1;// 非根结点有子树回不到上面 -> 割点
        }
        else if(i!=(fa^1))low[u]=min(low[u],dfn[v]);// 反向边走一次即可
    }
    if(!fa&&child>=2)is_cut[u]=1;// 根结点有两棵以上子树才是割点
}
\end{lstlisting}
% SOURCE: 04-图论/2-SAT.cpp
\subsection{2-SAT}

\note{n为布尔变量数，adj保存蕴含边；真/\allowbreak{}假节点通过id映射到1.\allowbreak{}.\allowbreak{}2n。dfn/\allowbreak{}low/\allowbreak{}scc/\allowbreak{}stk/\allowbreak{}ins是Tarjan状态，val[\allowbreak{}1.\allowbreak{}.\allowbreak{}n]\allowbreak{}为解；solve重置SCC状态而保留约束图，solve\_lex还会追加固定条件，换约束集调用clear\_all。}
\begin{lstlisting}
const int N=200005;
int n,m;
vector<int> adj[N];// 蕴含边：u -> v 表示 u 真则 v 真
int dfn[N],low[N],scc[N],stk[N],ins[N],tim,top,cnt;
int val[N];// val[i]：变量 i 的取值（0/1）
int id(int x,int v)// 变量 x 取值为 v 时对应的结点：1..n 为真，n+1..2n 为假
{
    return v?x:x+n;
}
void add_clause(int x,int vx,int y,int vy)// 加条件 (x==vx) 或 (y==vy)
{
    adj[id(x,!vx)].push_back(id(y,vy));// x 不取 vx，就只能 y 取 vy
    adj[id(y,!vy)].push_back(id(x,vx));
}
void tarjan(int u)// 缩点，O(n+m)
{
    dfn[u]=low[u]=++tim;
    stk[++top]=u,ins[u]=1;
    for(int i=0;i<(int)adj[u].size();i++)
    {
        int v=adj[u][i];
        if(!dfn[v])tarjan(v),low[u]=min(low[u],low[v]);
        else if(ins[v])low[u]=min(low[u],dfn[v]);
    }
    if(dfn[u]==low[u])
    {
        cnt++;
        int v;
        do
        {
            v=stk[top--];
            ins[v]=0;
            scc[v]=cnt;
        }while(v!=u);
    }
}
int solve()// 返回是否有解，有解时 val[] 是一组可行赋值
{
    // 重新判定只清 SCC 工作数组，保留图与此前固定条件。
    for(int i=1;i<=2*n;i++)dfn[i]=low[i]=scc[i]=ins[i]=0;
    tim=top=cnt=0;
    for(int i=1;i<=2*n;i++)if(!dfn[i])tarjan(i);
    for(int i=1;i<=n;i++)
    {
        if(scc[i]==scc[i+n])return 0;// 真和假互相可达 -> 矛盾
        val[i]=scc[i]<scc[i+n];// Tarjan 出栈序是反拓扑序，编号小的在拓扑序后，取它为真
    }
    return 1;
}
\end{lstlisting}
\note{小规模字典序最小赋值：按变量 1.\allowbreak{}.\allowbreak{}n 优先取 0；O(\allowbreak{}n*(\allowbreak{}n+\allowbreak{}m)\allowbreak{})\allowbreak{}，图中额外保留 n 条强制边。 SCC 编号给出的任意解不保证字典序；单位子句 x=\allowbreak{}v 只需连 !(\allowbreak{}x=\allowbreak{}v)\allowbreak{}-\allowbreak{}>\allowbreak{}(\allowbreak{}x=\allowbreak{}v)\allowbreak{}。}
\begin{lstlisting}
int solve_lex()
{
    if(!solve())return 0;
    for(int x=1;x<=n;x++)
    {
        adj[id(x,1)].push_back(id(x,0)); // 试 x=0
        if(!solve())
        {
            adj[id(x,1)].pop_back();
            adj[id(x,0)].push_back(id(x,1)); // 前面固定的条件保留，改成 x=1
        }
    }
    return solve();
}
void clear_all()// 多组数据清空
{
    for(int i=1;i<=2*n;i++)adj[i].clear(),dfn[i]=low[i]=scc[i]=ins[i]=0;
    tim=top=cnt=0;
}
\end{lstlisting}
% SOURCE: 04-图论/差分约束.cpp
\subsection{差分约束}

\note{adj存\{邻点,\allowbreak{}权值\}，清空范围含超级源0；dis[\allowbreak{}i]\allowbreak{}对应变量x\_i，边u到v权c表示x\_v<\allowbreak{}=\allowbreak{}x\_u+\allowbreak{}c；0为超级源，cnt/\allowbreak{}inq为负环检测和队列状态。solve会追加超级源边，多次调用前重建图，结果为一组可行势能，未必满足题目额外的最优目标。}
\begin{lstlisting}
const int N=100005;
const ll INF=1e18;
int n,m;
struct Edge{int to;ll w;};
vector<Edge> adj[N];
ll dis[N];// dis[i] 就是变量 x_i 的一组可行解
int cnt[N],inq[N];
void add_edge(int u,int v,ll c)// 有向边 u->v 权 c，表示 x_v <= x_u + c
{
    adj[u].push_back({v,c});
}
void add_leq(int u,int v,ll c)// 条件 x_v - x_u <= c，直接就是一条 u->v 边权 c
{
    add_edge(u,v,c);
}
void add_geq(int u,int v,ll c)// 条件 x_v - x_u >= c，转成 x_u - x_v <= -c，即 v->u 边权 -c
{
    add_edge(v,u,-c);
}
int spfa(int s)// 从超级源点 0 跑最短路，返回 1 表示有负环（约束矛盾，无解）
{
    for(int i=0;i<=n;i++)dis[i]=INF,cnt[i]=0,inq[i]=0;
    queue<int> q;
    dis[s]=0,inq[s]=1,q.push(s);
    while(!q.empty())
    {
        int u=q.front();
        q.pop();
        inq[u]=0;
        for(auto [v,w]:adj[u])
        {
            if(dis[u]+w<dis[v])
            {
                dis[v]=dis[u]+w;
                cnt[v]=cnt[u]+1;
                if(cnt[v]>n)return 1;// 最短路用到的边数超过 n，必有负环
                if(!inq[v])inq[v]=1,q.push(v);
            }
        }
    }
    return 0;
}
int solve()// 建超级源点 0 连向所有变量，返回是否有解
{
    for(int i=1;i<=n;i++)add_edge(0,i,0);// 每个变量都可以取 <= 0 的值，等价于固定一组参照
    return !spfa(0);
}
void clear_all()
{
    for(int i=0;i<=n;i++)adj[i].clear();
}
\end{lstlisting}
% SOURCE: 04-图论/二分图判定.cpp
\subsection{二分图判定}

\note{adj为无向图，col为0未染色、1/\allowbreak{}2两侧，c1/\allowbreak{}c2计当前连通块两侧点数。断连图逐块启动染色，遇同色边即失败；换图清空颜色，孤点也是一个连通块。}
\begin{lstlisting}
const int N=100005;
int n,m;
vector<int> adj[N];
int col[N];// 0 未染色，1 / 2 两种颜色（黑白）
int c1,c2;// 当前连通块两种颜色的点数
int dfs_color(int u,int c)// 深搜染色，返回 0 表示出现奇环（不是二分图）
{
    col[u]=c;
    if(c==1)c1++;
    else c2++;
    for(int i=0;i<(int)adj[u].size();i++)
    {
        int v=adj[u][i];
        if(!col[v])
        {
            if(!dfs_color(v,3-c))return 0;// 3-c 就是换成另一种颜色
        }
        else if(col[v]==c)return 0;// 相邻同色 -> 矛盾
    }
    return 1;
}
int bfs_color(int s)// 宽搜染色，写法与 dfs 等价
{
    queue<int> q;
    col[s]=1,c1=1,c2=0;
    q.push(s);
    while(!q.empty())
    {
        int u=q.front();
        q.pop();
        for(int i=0;i<(int)adj[u].size();i++)
        {
            int v=adj[u][i];
            if(!col[v])
            {
                col[v]=3-col[u];
                if(col[v]==1)c1++;
                else c2++;
                q.push(v);
            }
            else if(col[v]==col[u])return 0;
        }
    }
    return 1;
}
\end{lstlisting}
% SOURCE: 04-图论/二分图最大匹配(匈牙利).cpp
\subsection{二分图最大匹配(匈牙利)}

\note{adj[\allowbreak{}u]\allowbreak{}存左点的右邻居，match[\allowbreak{}v]\allowbreak{}是右点对应的左点，0未匹配；vis为本轮增广访问的右点。两侧编号分别从1起，不能混成同一侧；每次新增广清vis，换图清match与邻接。}
\begin{lstlisting}
const int N=1005;
int n,m;// 左右各 n 个点（左右点数不同就取较大的上界，单独记 nl,nr）
vector<int> adj[N];// adj[u] 存左部点 u 能匹配的右部点
int match[N];// match[v]：右部点 v 匹配的左部点
int vis[N];// 本轮增广中右部点是否访问过
int dfs(int u)// 从左部点 u 出发找增广路，O(m)
{
    for(int i=0;i<(int)adj[u].size();i++)
    {
        int v=adj[u][i];
        if(vis[v])continue;// 本轮已经试过，跳过
        vis[v]=1;
        if(!match[v]||dfs(match[v]))// v 没匹配，或者能让它原来的搭档腾位置
        {
            match[v]=u;
            return 1;
        }
    }
    return 0;
}
int hungarian(int nl)// 匈牙利求最大匹配，O(n*m)
{
    for(int i=1;i<=n;i++)match[i]=0;
    int ans=0;
    for(int u=1;u<=nl;u++)
    {
        for(int i=1;i<=n;i++)vis[i]=0;// 每个左部点重新清空标记
        if(dfs(u))ans++;
    }
    return ans;
}
\end{lstlisting}
% SOURCE: 04-图论/二分图最小点覆盖.cpp
\subsection{二分图最小点覆盖}

\note{adj[\allowbreak{}1.\allowbreak{}.\allowbreak{}nl]\allowbreak{}存左点的右邻居；match[\allowbreak{}1.\allowbreak{}.\allowbreak{}nr]\allowbreak{}存右点在最大匹配中的左伴侣，0为未匹配。先求最大匹配，再vertex\_cover(\allowbreak{})\allowbreak{}，返回\{覆盖的左点,\allowbreak{}覆盖的右点\}。每次重建交错路访问状态，原邻接表和匹配不变。}
\note{必须先求最大匹配（仅“无法再直接配对”的极大匹配不够）。 adj[\allowbreak{}1.\allowbreak{}.\allowbreak{}nl]\allowbreak{} 是左点的右邻居；match[\allowbreak{}1.\allowbreak{}.\allowbreak{}nr]\allowbreak{} 是右点匹配的左点，0 表示未匹配。 从未匹配左点沿交错路：左-\allowbreak{}>\allowbreak{}右走非匹配边，右-\allowbreak{}>\allowbreak{}左走匹配边。 最小点覆盖 =\allowbreak{} 未到达左点 +\allowbreak{} 到达右点；返回左右编号，O(\allowbreak{}n+\allowbreak{}m)\allowbreak{}。 最大独立集取该覆盖的补集；本结论仅适用于二分图。}
\begin{lstlisting}
vector<vector<int>> adj;
vector<int> match;
pair<vector<int>,vector<int>> vertex_cover()
{
    int nl=(int)adj.size()-1,nr=(int)match.size()-1;
    vector<int> ml(nl+1),vl(nl+1),vr(nr+1);
    for(int v=1;v<=nr;v++)if(match[v])ml[match[v]]=v;
    queue<int> q;
    for(int u=1;u<=nl;u++)if(!ml[u])vl[u]=1,q.push(u);
    while(!q.empty())
    {
        int u=q.front();q.pop();
        for(int v:adj[u])if(v!=ml[u]&&!vr[v])
        {
            vr[v]=1;
            int w=match[v];
            if(w&&!vl[w])vl[w]=1,q.push(w);
        }
    }
    pair<vector<int>,vector<int>> ans;
    for(int u=1;u<=nl;u++)if(!vl[u])ans.first.push_back(u);
    for(int v=1;v<=nr;v++)if(vr[v])ans.second.push_back(v);
    return ans;
}
\end{lstlisting}
% SOURCE: 04-图论/网络流Dinic.cpp
\subsection{网络流Dinic}

\note{adj[\allowbreak{}u]\allowbreak{}存残量边在vector<\allowbreak{}Edge>\allowbreak{} e中的下标；e[\allowbreak{}id]\allowbreak{}.\allowbreak{}to/\allowbreak{}cap为终点和剩余容量，编号从0起，id\textasciicircum{}1是反向边。s/\allowbreak{}t为源汇，dep为层次，cur[\allowbreak{}u]\allowbreak{}是adj[\allowbreak{}u]\allowbreak{}中的位置，从0起；dinic修改容量，再调用只求新增流。重新求原网络须清空adj、e后建图，或恢复初始容量。}
\note{适用：容量限制、二分图匹配、最小割；最大流等于最小割容量。 编号：顶点 1.\allowbreak{}.\allowbreak{}n；s,\allowbreak{}t 必须不同，容量 c>\allowbreak{}=\allowbreak{}0。 存图：adj 存边编号，e 动态保存成对残量边；两条边占两个元素。 关键：只走层数加一的边；当前弧按邻接表位置推进，失败边不反复扫描。 易错：调用会修改 e[\allowbreak{}id]\allowbreak{}.\allowbreak{}cap；重新求原图最大流须重建，清空 adj 与 e。 复杂度：一般图 O(\allowbreak{}n\textasciicircum{}2*m)\allowbreak{}，空间 O(\allowbreak{}n+\allowbreak{}m)\allowbreak{}；递归深度最坏 O(\allowbreak{}n)\allowbreak{}。 最大权闭合图：选 u 必须选 v 就连 u-\allowbreak{}>\allowbreak{}v(\allowbreak{}INF)\allowbreak{}；正权连 S-\allowbreak{}>\allowbreak{}u(\allowbreak{}w)\allowbreak{}，负权连 u-\allowbreak{}>\allowbreak{}T(\allowbreak{}-\allowbreak{}w)\allowbreak{}。 最优值=\allowbreak{}正权总和-\allowbreak{}最小割；最后残量图中 S 可达点即选集。INF 大于所有有限容量总和且不溢出。}
\begin{lstlisting}
const int N=100005;
const ll INF=1e18;
int n,m,s,t;
struct Edge{int to;ll cap;};
vector<Edge> e;
vector<int> adj[N];// 存 e 的下标，从 0 起；id^1 是反向边
int dep[N],cur[N];
\end{lstlisting}
\note{O(\allowbreak{}1)\allowbreak{}，加容量为 c 的 u-\allowbreak{}>\allowbreak{}v 边；反向边初始为 0，后续用于撤销旧流。}
\begin{lstlisting}
void add_edge(int u,int v,ll c)
{
    int id=e.size();
    e.push_back({v,c});
    e.push_back({u,0});
    adj[u].push_back(id);
    adj[v].push_back(id^1);
}
\end{lstlisting}
\note{O(\allowbreak{}n+\allowbreak{}m)\allowbreak{}，按正残量边分层；每轮重置 dep，不能沿已满的边走。}
\begin{lstlisting}
int bfs()// 分层，O(m)；返回能否到达汇点
{
    for(int i=1;i<=n;i++)dep[i]=-1;
    queue<int> q;
    dep[s]=0;
    q.push(s);
    while(!q.empty())
    {
        int u=q.front();
        q.pop();
        for(int i:adj[u])
            if(e[i].cap>0&&dep[e[i].to]<0)
            {
                dep[e[i].to]=dep[u]+1;
                q.push(e[i].to);
            }
    }
    return dep[t]>=0;
}
\end{lstlisting}
\note{沿分层图从 u 推至 t，flow 为本次上限；单次最坏 O(\allowbreak{}n*m)\allowbreak{}，一轮阻塞流 O(\allowbreak{}n*m)\allowbreak{}。}
\begin{lstlisting}
ll dfs(int u,ll flow)// 沿分层图推流，当前弧优化保证每条边只扫一次
{
    if(u==t)return flow;
    for(int &p=cur[u];p<(int)adj[u].size();p++)
    {
        int i=adj[u][p],v=e[i].to;
        if(e[i].cap>0&&dep[v]==dep[u]+1)
        {
            ll f=dfs(v,min(flow,e[i].cap));
            if(f>0)
            {
                e[i].cap-=f;
                e[i^1].cap+=f;// 反向边加回去，支持退流
                return f;
            }
        }
    }
    return 0;
}
\end{lstlisting}
\note{O(\allowbreak{}n\textasciicircum{}2*m)\allowbreak{}，返回当前残量图还能增广的流量；每轮 BFS 后必须重置 cur。}
\begin{lstlisting}
ll dinic()// 最大流，O(n^2 m)，随机图/二分图很快
{
    ll ans=0;
    while(bfs())
    {
        for(int i=1;i<=n;i++)cur[i]=0;
        ll f;
        while((f=dfs(s,INF))>0)ans+=f;
    }
    return ans;
}
\end{lstlisting}
% SOURCE: 04-图论/费用流MCMF.cpp
\subsection{费用流MCMF}

\note{adj[\allowbreak{}u]\allowbreak{}存vector<\allowbreak{}Edge>\allowbreak{} e的下标；e[\allowbreak{}id]\allowbreak{}.\allowbreak{}to/\allowbreak{}cap/\allowbreak{}fee为终点、剩余容量和单位费用，编号从0起，id\textasciicircum{}1是反向边，反向费用取负。h为势能，dis为约化费用距离，pre为入边编号；ans\_flow/\allowbreak{}ans\_cost是本次调用的结果。多测清空adj、e后重建，费用乘流量须装入ll。}
\note{适用：运输、分配、带费用匹配；先最大化流量，再最小化总费用。 参数：顶点 1.\allowbreak{}.\allowbreak{}n，s!=\allowbreak{}t；c 是非负容量，f 是单位费用，可为负。 前提：源点可达残量图不能有负费用环；初始化 SPFA 没有负环检测。 关键：fee+\allowbreak{}h[\allowbreak{}u]\allowbreak{}-\allowbreak{}h[\allowbreak{}v]\allowbreak{} 为约化费用；势能保持最短路选择与真实费用一致。 易错：e[\allowbreak{}id]\allowbreak{}.\allowbreak{}cap 会被原地修改；重建时 清空 adj 与 e；反向边也占一个元素。 复杂度：初始化最坏 O(\allowbreak{}n*m)\allowbreak{}，A 次增广 O(\allowbreak{}A*(\allowbreak{}n+\allowbreak{}m)\allowbreak{}*log(\allowbreak{}n+\allowbreak{}m)\allowbreak{})\allowbreak{}；f*fee 和总费用须能放入 ll。 有限容量负费用边可先预置到上界，计入费用与点供需，再用其反向边撤销多余预流。 预置后必须平衡每个点的流量；这不是“直接灌满负边然后跑普通 s-\allowbreak{}t 流”，更不能用于无界负环。}
\begin{lstlisting}
const int N=100005;
const ll INF=1e18;
int n,m,s,t;
struct Edge{int to;ll cap,fee;};
vector<Edge> e;
vector<int> adj[N];// 存 e 的下标，从 0 起；id^1 是反向边
ll dis[N],h[N];// dis：本次的约化费用最短路；h：势能，保证约化费用非负
int vis[N],pre[N];// pre[v]：最短路上到 v 的那条入边编号
ll ans_flow,ans_cost;
\end{lstlisting}
\note{O(\allowbreak{}1)\allowbreak{}，添加 u-\allowbreak{}>\allowbreak{}v；c 为容量，f 为单价，反向边单价取负以抵消旧费用。}
\begin{lstlisting}
void add_edge(int u,int v,ll c,ll f)
{
    int id=e.size();
    e.push_back({v,c,f});
    e.push_back({u,0,-f});
    adj[u].push_back(id);
    adj[v].push_back(id^1);
}
\end{lstlisting}
\note{最坏 O(\allowbreak{}n*m)\allowbreak{}，从全局 s 求初始势能；先处理负边才能使用 Dijkstra。}
\begin{lstlisting}
void spfa_init()// 先用 SPFA 求一遍初始势能，这样有负费用边也能跑 Dijkstra
{
    for(int i=1;i<=n;i++)h[i]=INF,vis[i]=0;
    queue<int> q;
    h[s]=0,vis[s]=1,q.push(s);
    while(!q.empty())
    {
        int u=q.front();
        q.pop();
        vis[u]=0;
        for(int i:adj[u])
        {
            int v=e[i].to;
            if(e[i].cap>0&&h[u]+e[i].fee<h[v])
            {
                h[v]=h[u]+e[i].fee;
                if(!vis[v])vis[v]=1,q.push(v);
            }
        }
    }
    for(int i=1;i<=n;i++)if(h[i]==INF)h[i]=0;// 不可达点势能置 0
}
\end{lstlisting}
\note{O(\allowbreak{}(\allowbreak{}n+\allowbreak{}m)\allowbreak{}*log(\allowbreak{}n+\allowbreak{}m)\allowbreak{})\allowbreak{}，仅走正容量边；pre[\allowbreak{}v]\allowbreak{} 保存入边，返回 t 是否可达。}
\begin{lstlisting}
int dijkstra()// 沿约化费用最短路增广，O(m log n)；返回能否找到增广路
{
    for(int i=1;i<=n;i++)dis[i]=INF,vis[i]=0,pre[i]=0;
    priority_queue<pair<ll,int>,vector<pair<ll,int>>,greater<pair<ll,int>>> q;
    dis[s]=0,q.push(make_pair(0LL,s));
    while(!q.empty())
    {
        pair<ll,int> p=q.top();
        q.pop();
        int u=p.second;
        if(vis[u])continue;// 过期的旧距离，跳过
        vis[u]=1;
        for(int i:adj[u])
        {
            int v=e[i].to;
            if(e[i].cap>0&&dis[u]+e[i].fee+h[u]-h[v]<dis[v])
            {
                dis[v]=dis[u]+e[i].fee+h[u]-h[v];
                pre[v]=i;
                q.push(make_pair(dis[v],v));
            }
        }
    }
    return dis[t]<INF;
}
\end{lstlisting}
\note{初始化加 A 次最短路增广；ans\_flow/\allowbreak{}ans\_cost 为输出，沿 pre 的反边终点回溯父点。}
\begin{lstlisting}
void mcmf()// 最小费用最大流，结果放 ans_flow / ans_cost
{
    ans_flow=0,ans_cost=0;
    // 最坏 O(n*m)，从全局 s 求初始势能；先处理负边才能使用 Dijkstra。
    spfa_init();
    while(dijkstra())
    {
        for(int i=1;i<=n;i++)if(dis[i]<INF)h[i]+=dis[i];// 累加势能
        ll f=INF;
        for(int v=t;v!=s;v=e[pre[v]^1].to)f=min(f,e[pre[v]].cap);// 沿路径找瓶颈
        for(int v=t;v!=s;v=e[pre[v]^1].to)
        {
            e[pre[v]].cap-=f;
            e[pre[v]^1].cap+=f;
            ans_cost+=f*e[pre[v]].fee;// 按真实费用累加
        }
        ans_flow+=f;
    }
}
\end{lstlisting}
% SOURCE: 04-图论/上下界网络流.cpp
\subsection{上下界网络流}

\note{eu/\allowbreak{}ev与elow/\allowbreak{}eup为原边端点和上下界；adj存vector<\allowbreak{}Edge>\allowbreak{} e的下标，编号从0起，id\textasciicircum{}1是反向边，cur是邻接表中的位置。eidx定位原边的残量正向边，实际流量为elow[\allowbreak{}i]\allowbreak{}+\allowbreak{}e[\allowbreak{}eidx[\allowbreak{}i]\allowbreak{}\textasciicircum{}1]\allowbreak{}.\allowbreak{}cap；d为下界产生的净流入需求，S/\allowbreak{}T是超级源汇。build清空adj、e并重建容量和需求，solve\_lr先满足需求再移除辅助边。}
\note{适用：每条边必须达到最低流量的可行流、最大流、最小流。 编号：原点 1.\allowbreak{}.\allowbreak{}n，边 1.\allowbreak{}.\allowbreak{}m；超级源汇 n+\allowbreak{}1,\allowbreak{}n+\allowbreak{}2，必须均小于 N。 参数：eu/\allowbreak{}ev 是端点，0<\allowbreak{}=\allowbreak{}elow<\allowbreak{}=\allowbreak{}eup；有源汇时 s!=\allowbreak{}t。 关键：先固定下界，再用 eup-\allowbreak{}elow 建残量边；d 记录下界造成的收支差。 结论：超级源总流量等于 need 才可行；原边实际流量为 elow[\allowbreak{}i]\allowbreak{}+\allowbreak{}e[\allowbreak{}eidx[\allowbreak{}i]\allowbreak{}\textasciicircum{}1]\allowbreak{}.\allowbreak{}cap。 易错：e 动态容纳原边、平衡边和 t-\allowbreak{}>\allowbreak{}s 边及反边；INF 须大于所需总流量。 边界：最小流写法允许 base-\allowbreak{}dinic(\allowbreak{}t,\allowbreak{}s)\allowbreak{} 为负；题目若只接受非负流需另核对约定。}
\begin{lstlisting}
const int N=1005;
const ll INF=1e18;
int n,m,s,t,S,T;// S,T 是超级源汇
struct Edge{int to;ll cap;};
vector<Edge> e;
vector<int> adj[N];// 存 e 的下标，从 0 起；id^1 是反向边
int dep[N],cur[N];
ll d[N];// d[i]>0：i 还需要流入这么多；d[i]<0：i 需要流出这么多
int eu[N],ev[N],eidx[N];// eidx[i]：第 i 条上下界边对应的正向边编号
ll elow[N],eup[N];
\end{lstlisting}
\note{O(\allowbreak{}1)\allowbreak{}，加入 u-\allowbreak{}>\allowbreak{}v 的剩余容量 c；端点可含超级源汇。}
\begin{lstlisting}
void add_edge(int u,int v,ll c)
{
    int id=e.size();
    e.push_back({v,c});
    e.push_back({u,0});
    adj[u].push_back(id);
    adj[v].push_back(id^1);
}
\end{lstlisting}
\note{O(\allowbreak{}n+\allowbreak{}m)\allowbreak{}，ss/\allowbreak{}tt 是本轮源汇；只给正残量边分层。}
\begin{lstlisting}
int bfs(int ss,int tt)// 分层，O(m)
{
    for(int i=1;i<=n+2;i++)dep[i]=-1;
    queue<int> q;
    dep[ss]=0;
    q.push(ss);
    while(!q.empty())
    {
        int u=q.front();
        q.pop();
        for(int i:adj[u])
            if(e[i].cap>0&&dep[e[i].to]<0)
            {
                dep[e[i].to]=dep[u]+1;
                q.push(e[i].to);
            }
    }
    return dep[tt]>=0;
}
\end{lstlisting}
\note{从 u 到 tt 推至多 flow；单次最坏 O(\allowbreak{}n*m)\allowbreak{}，当前弧避免同轮反复试死边。}
\begin{lstlisting}
ll dfs(int u,int tt,ll flow)// 沿分层图推流，当前弧优化
{
    if(u==tt)return flow;
    for(int &p=cur[u];p<(int)adj[u].size();p++)
    {
        int i=adj[u][p],v=e[i].to;
        if(e[i].cap>0&&dep[v]==dep[u]+1)
        {
            ll f=dfs(v,tt,min(flow,e[i].cap));
            if(f>0)
            {
                e[i].cap-=f;
                e[i^1].cap+=f;
                return f;
            }
        }
    }
    return 0;
}
\end{lstlisting}
\note{O(\allowbreak{}n\textasciicircum{}2*m)\allowbreak{}，返回 ss-\allowbreak{}>\allowbreak{}tt 的新增流量，并修改残量容量。}
\begin{lstlisting}
ll dinic(int ss,int tt)// 最大流，O(n^2 m)
{
    ll ans=0;
    while(bfs(ss,tt))
    {
        for(int i=1;i<=n+2;i++)cur[i]=0;
        ll f;
        while((f=dfs(ss,tt,INF))>0)ans+=f;
    }
    return ans;
}
\end{lstlisting}
\note{O(\allowbreak{}n+\allowbreak{}m)\allowbreak{}，从 1.\allowbreak{}.\allowbreak{}m 的输入边重建；eidx 指向正边，清空旧平衡量。}
\begin{lstlisting}
void build()// 按 eu/ev/elow/eup 重新建图：上下界边先默认流下界，剩下的容量建成普通边
{
    for(int i=1;i<=n+2;i++)adj[i].clear(),d[i]=0;
    e.clear();
    for(int i=1;i<=m;i++)
    {
        add_edge(eu[i],ev[i],eup[i]-elow[i]);
        eidx[i]=(int)e.size()-2;
        d[eu[i]]-=elow[i],d[ev[i]]+=elow[i];// 下界先当满流记账
    }
}
\end{lstlisting}
\note{type=\allowbreak{}0：无源汇可行流；type=\allowbreak{}1：有源汇上下界最大流；type=\allowbreak{}2：有源汇上下界最小流 返回流量大小，-\allowbreak{}1 表示无解 O(\allowbreak{}n\textasciicircum{}2*m)\allowbreak{}，type=\allowbreak{}0/\allowbreak{}1/\allowbreak{}2 对应可行/\allowbreak{}最大/\allowbreak{}最小；返回 -\allowbreak{}1 也可能与负最小流混淆。}
\begin{lstlisting}
ll solve_lr(int type)
{
    build();
    int e_ts=0;
    if(type)add_edge(t,s,INF),e_ts=(int)e.size()-2;// 人为加 t->s 的无穷边，把有源汇变成无源汇
    S=n+1,T=n+2;
    ll need=0;
    for(int i=1;i<=n;i++)
    {
        if(d[i]>0)add_edge(S,i,d[i]),need+=d[i];// 缺流入的点由超级源补
        else if(d[i]<0)add_edge(i,T,-d[i]);// 多出来的流出丢给超级汇
    }
    if(dinic(S,T)!=need)return -1;// 超级源没流满 -> 下界无法同时满足
    if(type==0)return 0;
    ll base=e[e_ts^1].cap;// t->s 边上流过的量，就是当前 s->t 的流量
    for(int i=0;i<(int)e.size();i++)
        if(e[i].to==S||e[i].to==T)e[i].cap=0,e[i^1].cap=0;// 拆掉超级源汇的边
    e[e_ts].cap=0,e[e_ts^1].cap=0;// 拆掉人为边
    if(type==1)return base+dinic(s,t);// 还能继续增广就是最大流
    return base-dinic(t,s);// 能退掉的流量越多，剩下的就是最小流
}
\end{lstlisting}
% SOURCE: 04-图论/欧拉路径与欧拉回路.cpp
\subsection{欧拉路径与欧拉回路}

\note{adj/\allowbreak{}dadj分别存无向/\allowbreak{}有向\{邻点,\allowbreak{}id\}，两种id各从1起，ecnt/\allowbreak{}decnt为边数；无向边两端共用id。deg与din/\allowbreak{}dout为度数，used/\allowbreak{}dvis按边编号标记；it为邻接表位置，从0起，stk/\allowbreak{}path\_按边数分配，path\_[\allowbreak{}1.\allowbreak{}.\allowbreak{}pcnt]\allowbreak{}为结果。init清空对应邻接表，m须等于边数；求解时重置访问标记，可重复求解。先检查度与有边点连通，孤点不影响。}
\note{欧拉路径 /\allowbreak{} 欧拉回路（Hierholzer 算法，用显式栈写，n=\allowbreak{}1e5 不爆栈） 无向图：所有点度数为偶 → 欧拉回路；恰有 0/\allowbreak{}2 个奇度点 → 有欧拉路径 有向图：所有点入度=\allowbreak{}出度 → 欧拉回路；恰有 1 个出度比入度大 1（起点）、 1 个入度比出度大 1（终点）、其余相等 → 有欧拉路径 前置条件：忽略孤立点后图连通}
\begin{lstlisting}
const int N=100005;
int n,m;
struct Edge{int to,id;};
vector<Edge> adj[N],dadj[N];// 无向边两端共用 id；有向边各有一个 id
int ecnt,decnt,deg[N],din[N],dout[N],it[N],pcnt;
vector<int> used,dvis,stk,path_;// 路径和栈按边数扩展，path_[1..pcnt] 为结果
void init_undirected(int n_)
{
    n=n_,ecnt=0,pcnt=0;
    for(int i=1;i<=n;i++)adj[i].clear(),deg[i]=0;
    used.assign(1,0);
}
void add_uedge(int u,int v)// 无向边，一次加一对
{
    ++ecnt;
    adj[u].push_back({v,ecnt});
    adj[v].push_back({u,ecnt});
    used.push_back(0);
    deg[u]++,deg[v]++;
}
void init_directed(int n_)
{
    n=n_,decnt=0,pcnt=0;
    for(int i=1;i<=n;i++)dadj[i].clear(),din[i]=0,dout[i]=0;
    dvis.assign(1,0);
}
void add_dedge(int u,int v)// 有向边
{
    dadj[u].push_back({v,++decnt});
    dvis.push_back(0);
    dout[u]++,din[v]++;
}
\end{lstlisting}
\note{忽略孤立点后是否连通（无向）}
\begin{lstlisting}
int connected_undirected()
{
    int s=0;
    for(int i=1;i<=n;i++)if(deg[i]){s=i;break;}
    if(!s)return 1;// 没有边
    int vis2[N]={0},q[N],hd=0,tl=0,cnt=0,all=0;
    for(int i=1;i<=n;i++)if(deg[i])all++;
    vis2[s]=1,q[tl++]=s;
    while(hd<tl)
    {
        int u=q[hd++];
        cnt++;
        for(auto [v,id]:adj[u])
            if(!vis2[v])vis2[v]=1,q[tl++]=v;
    }
    return cnt==all;
}
\end{lstlisting}
\note{有向图弱连通（按无向边判连通，需要另建一份邻接，这里偷懒用并查集）}
\begin{lstlisting}
int ufa[N];
int findd(int x){return ufa[x]==x?x:ufa[x]=findd(ufa[x]);}
int connected_directed()
{
    for(int i=1;i<=n;i++)ufa[i]=i;
    int all=0;
    for(int i=1;i<=n;i++)if(din[i]||dout[i])all++;
    for(int u=1;u<=n;u++)
        for(auto [v,id]:dadj[u])
            ufa[findd(u)]=findd(v);
    if(!all)return 1;
    int r=-1;
    for(int i=1;i<=n;i++)if(din[i]||dout[i]){r=findd(i);break;}
    for(int i=1;i<=n;i++)if((din[i]||dout[i])&&findd(i)!=r)return 0;
    return 1;
}
\end{lstlisting}
\note{Hierholzer：无向图。返回 0 表示不存在，1 表示路径，2 表示回路；路径存在 path\_[\allowbreak{}1.\allowbreak{}.\allowbreak{}pcnt]\allowbreak{} 里}
\begin{lstlisting}
int euler_undirected(int &s)
{
    pcnt=0;
    int odd=0;
    for(int i=1;i<=n;i++)
        if(deg[i]&1)odd++,s=i;
    if(odd!=0&&odd!=2)return 0;
    if(odd==2)
    {
        int c=0;
        for(int i=1;i<=n;i++)if(deg[i]&1){c++;if(c==1)s=i;}
    }
    else
    {
        s=0;
        for(int i=1;i<=n;i++)if(deg[i]){s=i;break;}
    }
    if(!s)return 2;
    if(!connected_undirected())return 0;
    for(int i=1;i<=n;i++)it[i]=0;
    used.assign(ecnt+1,0);
    stk.resize(ecnt+2),path_.resize(ecnt+2);
    int tp=0;
    stk[++tp]=s,pcnt=0;
    while(tp)
    {
        int u=stk[tp];
        int &p=it[u];
        while(p<(int)adj[u].size()&&used[adj[u][p].id])p++;
        if(p<(int)adj[u].size())
        {
            auto [v,id]=adj[u][p++];
            used[id]=1;
            stk[++tp]=v;
        }
        else
        {
            path_[++pcnt]=u;// 出栈顺序就是欧拉路的逆序
            tp--;
        }
    }
    reverse(path_.begin()+1,path_.begin()+pcnt+1);
    if(pcnt!=m+1)return 0;// 边没走完说明图不连通
    return odd==2?1:2;
}
\end{lstlisting}
\note{Hierholzer：有向图。返回 0/\allowbreak{}1/\allowbreak{}2 同上}
\begin{lstlisting}
int euler_directed(int &s)
{
    pcnt=0;
    int c1=0,c2=0;
    s=0;
    for(int i=1;i<=n;i++)
    {
        if(dout[i]-din[i]==1)c1++,s=i;
        else if(din[i]-dout[i]==1)c2++;
        else if(din[i]!=dout[i])return 0;
    }
    if(!((c1==1&&c2==1)||(c1==0&&c2==0)))return 0;
    if(c1==0)
    {
        for(int i=1;i<=n;i++)if(dout[i]){s=i;break;}
    }
    if(!s)return 2;
    if(!connected_directed())return 0;
    for(int i=1;i<=n;i++)it[i]=0;
    dvis.assign(decnt+1,0);
    stk.resize(decnt+2),path_.resize(decnt+2);
    int tp=0;
    stk[++tp]=s,pcnt=0;
    while(tp)
    {
        int u=stk[tp];
        int &p=it[u];
        while(p<(int)dadj[u].size()&&dvis[dadj[u][p].id])p++;
        if(p<(int)dadj[u].size())
        {
            auto [v,id]=dadj[u][p++];
            dvis[id]=1;
            stk[++tp]=v;
        }
        else
        {
            path_[++pcnt]=u;
            tp--;
        }
    }
    reverse(path_.begin()+1,path_.begin()+pcnt+1);
    if(pcnt!=m+1)return 0;
    return c1==1?1:2;
}
\end{lstlisting}
% SOURCE: 04-图论/基环树.cpp
\subsection{基环树}

\note{adj[\allowbreak{}u]\allowbreak{}存\{邻点,\allowbreak{}id\}，一条无向边两端共用从1起的id，ecnt为无向边数；pe只跳过同一父边，支持重边和自环。val为点权，ring[\allowbreak{}1.\allowbreak{}.\allowbreak{}rn]\allowbreak{}按环顺序存点，on\_ring标环点；par/\allowbreak{}vis与fstk/\allowbreak{}fit找环，fit为邻接表位置，从0起。dp0/\allowbreak{}dp1为树枝不选/\allowbreak{}选点最优值；环上固定首点状态做两遍DP。换图清空adj、ecnt；有向找环须单独建每点一条出边。}
\note{基环树：n 个点 n 条边的连通图，有且只有一个环 内容：找环（无向/\allowbreak{}有向两种） +\allowbreak{} 环上 DP（例题：基环树最大独立集，即「没有上司的舞会」带环版） 最大独立集做法：先对每个环点挂的树做树形 DP，再把环拆成链做两遍线性 DP 找环与 DP 全部迭代实现，n=\allowbreak{}1e5 的链+\allowbreak{}环不会爆栈}
\begin{lstlisting}
const int N=100005;
int n,m,ecnt;
struct Edge{int to,id;};
vector<Edge> adj[N];
int val[N];
int par[N],pe[N],vis[N];        // par：父亲，pe：从父亲过来的无向边编号（两向共用）
int ring[N],rn;                 // 环上的点，按环顺序（按环边相邻，编号不必单调）
int fstk[N],fit[N];             // 找环用的显式栈（放全局，避免递归/大数组爆栈）
int mark[N],tmp[N];             // mark：祖先标记；tmp：暂存另一支路径
ll dp0[N],dp1[N];              // dp0：不选 u，dp1：选 u（子树部分）
int f0[N],f1[N];                // 环上线性 DP
void add_edge(int u,int v)
{
    ++ecnt;
    adj[u].push_back({v,ecnt});
    adj[v].push_back({u,ecnt});
}
\end{lstlisting}
\note{迭代 dfs 找无向图（n 点 n 边）里的那个环，O(\allowbreak{}n)\allowbreak{} 做法：dfs 时碰到已访问的点 (\allowbreak{}u,\allowbreak{}v)\allowbreak{}，就说明树边链 u→根 和 v→根 在某个点汇合，环 =\allowbreak{} 两支接起来}
\begin{lstlisting}
void find_ring_undirected(int rt)
{
    for(int i=1;i<=n;i++)vis[i]=0,par[i]=0,pe[i]=0,mark[i]=0;
    rn=0;
    int tp=0,cu=0,cv=0;
    vis[rt]=1,fstk[++tp]=rt,fit[tp]=0;
    while(tp)
    {
        int u=fstk[tp];
        if(fit[tp]<(int)adj[u].size())
        {
            auto [v,id]=adj[u][fit[tp]++];
            if(id==pe[u])continue;// 只跳同一条父边，保留重边
            if(vis[v]){cu=u,cv=v;break;}// 碰到走过的点 → 找到环
            vis[v]=1,par[v]=u,pe[v]=id;
            fstk[++tp]=v,fit[tp]=0;
        }
        else tp--;
    }
    if(!cu){rn=0;return;}
    // 先打出 cu 到根的链，再从 cv 往上走到第一个落在链上的点 l（用时间戳判，不用清数组）
    // 环 = cu→l 这一支 + cv→l 这一支接起来，正好是「树边 + 那条回边」构成的唯一环
    // 注意 l 可能就是 cv 自己（cv 是 cu 的祖先），所以要从 cv 自身开始判
    for(int i=1;i<=n;i++)mark[i]=0;
    for(int p=cu;p;p=par[p])mark[p]=1;
    int l=0;
    for(int p=cv;p;p=par[p])if(mark[p]){l=p;break;}
    if(!l){rn=0;return;}
    rn=0;
    for(int p=cu;p!=l;p=par[p])ring[++rn]=p;// cu 这一支
    ring[++rn]=l;
    int tc=0;
    for(int p=cv;p!=l;p=par[p])tmp[++tc]=p;// cv 这一支，反着接上去
    for(int i=tc;i>=1;i--)ring[++rn]=tmp[i];
}
\end{lstlisting}
\note{有向图单独建 adj[\allowbreak{}u]\allowbreak{}.\allowbreak{}push\_back(\allowbreak{}\{v,\allowbreak{}+\allowbreak{}+\allowbreak{}ecnt\})\allowbreak{}，不能用双向 add\_edge。 有向图找环（每个点出度为 1 /\allowbreak{} 内向基环树都适用）：沿出边走，走过就说明有环}
\begin{lstlisting}
void find_ring_directed(int rt)
{
    for(int i=1;i<=n;i++)vis[i]=0;
    rn=0;
    int p=rt;
    while(!vis[p])vis[p]=1,p=adj[p][0].to;
    int q=adj[p][0].to;
    while(q!=p)ring[++rn]=q,q=adj[q][0].to;
    ring[++rn]=p;
}
\end{lstlisting}
\note{树形 DP：以环点 rt 为根，只走 rt 挂的那棵树（环上的其他点都不进去），迭代逆序，O(\allowbreak{}n)\allowbreak{} 结果在 dp0/\allowbreak{}dp1 里：dp0[\allowbreak{}u]\allowbreak{} 不选 u，dp1[\allowbreak{}u]\allowbreak{} 选 u（只算这块子树）}
\begin{lstlisting}
int on_ring[N];
int torder[N],tfa[N],tq[N];// tree_dp 用的全局数组（放栈上会爆栈）
void tree_dp(int rt)
{
    int hd=0,tl=0,tot=0;
    tq[tl++]=rt,tfa[rt]=0,dp0[rt]=dp1[rt]=0;
    while(hd<tl)
    {
        int u=tq[hd++];
        torder[++tot]=u;
        for(auto [v,id]:adj[u])
        {
            if(v==tfa[u]||on_ring[v])continue;// 环上别的点不走，只走 rt 挂的树
            tfa[v]=u,tq[tl++]=v;
        }
    }
    for(int i=tot;i>=1;i--)
    {
        int u=torder[i];
        dp0[u]=0,dp1[u]=val[u];
        for(auto [v,id]:adj[u])
        {
            if(v==tfa[u]||on_ring[v])continue;
            dp0[u]+=max(dp0[v],dp1[v]);
            dp1[u]+=dp0[v];
        }
    }
}
\end{lstlisting}
\note{基环树最大独立集：先各环点挂树 DP，再在两段链上 DP，O(\allowbreak{}n)\allowbreak{}}
\begin{lstlisting}
ll max_independent_set(int rt)
{
    find_ring_undirected(rt);
    for(int i=1;i<=n;i++)on_ring[i]=0;
    for(int i=1;i<=rn;i++)on_ring[ring[i]]=1;
    for(int i=1;i<=rn;i++)tree_dp(ring[i]);
    // 首点不选、首点必选；末点不能与首点同时选
    const ll neg=-4e18;
    ll ans=0;
    for(int take=0;take<=(rn>1);take++)
    {
        ll f0=take?neg:dp0[ring[1]],f1=take?dp1[ring[1]]:neg;
        for(int i=2;i<=rn;i++)
        {
            int u=ring[i];
            ll g0=max(f0,f1)+dp0[u],g1=f0+dp1[u];
            f0=g0,f1=g1;
        }
        ans=max(ans,take?f0:max(f0,f1));
    }
    return ans;
}
\end{lstlisting}
% SOURCE: 04-图论/树上启发式合并DSUonTree.cpp
\subsection{树上启发式合并DSUonTree}

\note{adj存邻点，多测清空后重建；col为颜色，siz/\allowbreak{}heavy为大小与重儿子，in\_/\allowbreak{}out\_/\allowbreak{}rnk\_定位子树；cnt为颜色频数，color\_cnt/\allowbreak{}color\_mx及cur\_sum/\allowbreak{}cur\_mx为当前统计。added记录碰过的颜色以便清理；保留重子树再加入轻子树，换树清除残留计数。}
\note{树上启发式合并（DSU on tree /\allowbreak{} 小并大） 例题：每棵子树内不同颜色数 color\_cnt[\allowbreak{}u]\allowbreak{}、出现次数最多的颜色的出现次数 color\_mx[\allowbreak{}u]\allowbreak{} 复杂度 O(\allowbreak{}n log n)\allowbreak{}（重儿子保留贡献，轻儿子暴力重算） 预处理 dfs 序、子树大小、重儿子都是迭代的；get\_cnt 也用显式栈写，链状数据不爆栈}
\begin{lstlisting}
const int N=100005;
int n;
vector<int> adj[N];
int col[N],fa_[N],siz[N],heavy[N],in_[N],out_[N],rnk_[N],timer_;
int cnt[N],color_cnt[N],color_mx[N];// cnt：颜色当前出现次数
int cur_sum,cur_mx,touched,added[N];// touched/added：记录碰过的颜色，清空时只清这些
void add_edge(int u,int v)
{
    adj[u].push_back(v);
    adj[v].push_back(u);
}
\end{lstlisting}
\note{迭代求 dfs 序 /\allowbreak{} 子树大小 /\allowbreak{} 重儿子，O(\allowbreak{}n)\allowbreak{}}
\begin{lstlisting}
void get_order(int rt)
{
    int tl=0,order[N],stk[N];
    fa_[rt]=0,timer_=0;
    int tp=0;
    stk[++tp]=rt;
    while(tp)// 第一次：只求父亲和顺序
    {
        int u=stk[tp--];
        in_[u]=++timer_,rnk_[timer_]=u,order[++tl]=u;
        for(int v:adj[u])
        {
            if(v==fa_[u])continue;
            fa_[v]=u,stk[++tp]=v;
        }
    }
    for(int i=tl;i>=1;i--)// 逆序：子树大小与重儿子
    {
        int u=order[i];
        siz[u]=1,heavy[u]=0;
        for(int v:adj[u])
        {
            if(v==fa_[u])continue;
            siz[u]+=siz[v];
            if(siz[v]>siz[heavy[u]])heavy[u]=v;
        }
        out_[u]=in_[u]+siz[u]-1;// 子树是 dfs 序上一段区间
    }
}
void update(int u)// 把 u 这个点的颜色统计进来，O(1)
{
    int c=col[u];
    if(cnt[c]==0)cur_sum++;
    cnt[c]++;
    if(cnt[c]>cur_mx)cur_mx=cnt[c];
    if(cnt[c]==1)added[touched++]=c;
}
\end{lstlisting}
\note{迭代版 DSU on tree（用显式栈模拟上面那个递归，长链也不爆栈） 状态机：0 先处理轻儿子(\allowbreak{}keep=\allowbreak{}0)\allowbreak{} → 1 处理重儿子(\allowbreak{}keep=\allowbreak{}1)\allowbreak{} → 2 暴力加轻儿子+\allowbreak{}自己 → 3 收尾 递归版见注释里的 get\_cnt\_rec，逻辑完全一样}
\begin{lstlisting}
void get_cnt(int root,int keep)
{
    int st[N],state[N],saved[N],tp=0;
    for(int i=0;i<touched;i++)cnt[added[i]]=0;
    touched=cur_sum=cur_mx=0;
    st[++tp]=root,state[tp]=0,saved[tp]=keep;
    while(tp)
    {
        int u=st[tp],s=state[tp];
        if(s==0)
        {
            state[tp]=1;
            for(int v:adj[u])
            {
                if(v==fa_[u]||v==heavy[u])continue;
                st[++tp]=v,state[tp]=0,saved[tp]=0;// 轻儿子：算完丢掉
            }
        }
        else if(s==1)
        {
            state[tp]=2;
            if(heavy[u])st[++tp]=heavy[u],state[tp]=0,saved[tp]=1;// 重儿子：贡献保留
        }
        else if(s==2)
        {
            state[tp]=3;
            for(int v:adj[u])
            {
                if(v==fa_[u]||v==heavy[u])continue;
                for(int j=in_[v];j<=out_[v];j++)update(rnk_[j]);// 轻儿子暴力加
            }
            update(u);
            color_cnt[u]=cur_sum,color_mx[u]=cur_mx;// 此刻桶里正好是子树 u
        }
        else
        {
            if(!saved[tp])
            {
                for(int i=0;i<touched;i++)cnt[added[i]]=0;// 只清碰过的颜色
                touched=0,cur_sum=0,cur_mx=0;
            }
            tp--;
        }
    }
}
/* 递归版（n 小时更好看懂，链状数据会爆栈，正式用上面的迭代版）
void get_cnt_rec(int u,int keep)
{
    for(int v:adj[u])
        if(v!=fa_[u]&&v!=heavy[u])get_cnt_rec(v,0);
    if(heavy[u])get_cnt_rec(heavy[u],1);
    for(int v:adj[u])
        if(v!=fa_[u]&&v!=heavy[u])
            for(int j=in_[v];j<=out_[v];j++)update(rnk_[j]);
    update(u);
    color_cnt[u]=cur_sum,color_mx[u]=cur_mx;
    if(!keep)
    {
        for(int i=0;i<touched;i++)cnt[added[i]]=0;
        touched=0,cur_sum=0,cur_mx=0;
    }
}
*/
\end{lstlisting}
% SOURCE: 04-图论/点分治.cpp
\subsection{点分治}

\note{对象adj为无权树，k是距离，vis标已删除重心，cnt[\allowbreak{}d]\allowbreak{}计已处理子树距离桶；答案是无序异点对数。每棵子树先查询后入桶，used只清本轮触及项，solve后不能不重置vis就重复求解。}
\note{适用：无权树精确距离点对计数；统计无序且两端不同的点对。 参数：顶点 1.\allowbreak{}.\allowbreak{}n，n>\allowbreak{}=\allowbreak{}1，k 是边数距离；输入必须为连通树。 状态：vis 表示已删除重心，cnt[\allowbreak{}d]\allowbreak{} 是此前子树到当前重心距离为 d 的点数。 关键：每棵子树先查询、后入桶，避免把同一子树内部点对错误算作过重心路径。 易错：used 仅清本轮触及的桶，递归前必须清净；答案最多 n*(\allowbreak{}n-\allowbreak{}1)\allowbreak{}/\allowbreak{}2，用 ll。 复杂度：总 O(\allowbreak{}n log n)\allowbreak{}、空间 O(\allowbreak{}n)\allowbreak{}；重心递归深度 O(\allowbreak{}log n)\allowbreak{}，子树遍历用显式队列。}
\note{O(\allowbreak{}n log n)\allowbreak{}，统计距离为 k 的无序异点对；遍历子树不用递归}
\begin{lstlisting}
struct Centroid
{
    int n,k;
    vector<vector<int> > adj;
    vector<int> fa,sz,vis,cnt;
    Centroid(int n,int k):n(n),k(k),adj(n+1),fa(n+1),sz(n+1),vis(n+1),cnt(n+1){}
    // 摊还 O(1)，加无向树边 u-v；不要加入自环或额外环。
    void add_edge(int u,int v)
    {
        adj[u].push_back(v),adj[v].push_back(u);
    }
    int center(int rt)
    {
        vector<int> q(1,rt);
        fa[rt]=0;
        for(int i=0;i<(int)q.size();i++)
            for(int v:adj[q[i]])if(v!=fa[q[i]]&&!vis[v])fa[v]=q[i],q.push_back(v);
        int tot=q.size(),c=rt,best=tot;
        for(int i=tot-1;i>=0;i--)
        {
            int u=q[i],mx=0;
            sz[u]=1;
            for(int v:adj[u])if(fa[v]==u&&!vis[v])sz[u]+=sz[v],mx=max(mx,sz[v]);
            mx=max(mx,tot-sz[u]);
            if(mx<best)best=mx,c=u;
        }
        return c;
    }
    // 本块 O(块大小)，连同递归 O(块大小*log(块大小))；rt 是当前块入口，返回其点对数。
    ll solve(int rt)
    {
        int c=center(rt);
        vis[c]=1;
        ll ans=0;
        vector<int> used(1,0);
        cnt[0]=1;
        for(int v:adj[c])if(!vis[v])
        {
            vector<array<int,3> > q(1,{v,c,1});
            vector<int> d;
            for(int i=0;i<(int)q.size();i++)
            {
                int u=q[i][0],p=q[i][1],dep=q[i][2];
                if(dep>k)continue;
                d.push_back(dep),ans+=cnt[k-dep];
                for(int w:adj[u])if(w!=p&&!vis[w])q.push_back({w,u,dep+1});
            }
            for(int dep:d)
            {
                if(!cnt[dep])used.push_back(dep);
                cnt[dep]++;
            }
        }
        for(int d:used)cnt[d]=0;
        for(int v:adj[c])if(!vis[v])ans+=solve(v);
        return ans;
    }
    ll run()
    {
        if(k<=0||k>=n)return 0;
        fill(vis.begin(),vis.end(),0);
        return solve(1);
    }
};
\end{lstlisting}
% SOURCE: 04-图论/虚树.cpp
\subsection{虚树}

\note{对象adj为原树，dfn/\allowbreak{}ed/\allowbreak{}dep及祖先表支持LCA，tr为当前虚树，mark标关键点，cnt供子树汇总。每次给关键点列表，按DFS序加相邻LCA建树；清理上次触及的节点，不能遍历旧虚边。}
\begin{lstlisting}
struct VirtualTree
{
    int n,tim,lg;
    vector<vector<int>> adj,up,tr;
    vector<int> dep,dfn,ed,mark,cnt;
    VirtualTree(int n):n(n),tim(0),lg(1),adj(n+1),tr(n+1),dep(n+1),dfn(n+1),ed(n+1),mark(n+1),cnt(n+1)
    {
        while((1<<lg)<=n)lg++;
        up.assign(lg,vector<int>(n+1));
    }
    void add_edge(int u,int v)
    {
        adj[u].push_back(v),adj[v].push_back(u);
    }
    // 迭代预处理，O(n log n)，避免链上递归爆栈
    void build()
    {
        vector<pair<int,int>> st;
        st.push_back({1,0});
        while(!st.empty())
        {
            int u=st.back().first,&i=st.back().second;
            if(i==0)
            {
                dfn[u]=++tim;
                for(int j=1;j<lg;j++)up[j][u]=up[j-1][up[j-1][u]];
            }
            if(i==(int)adj[u].size())
            {
                ed[u]=tim,st.pop_back();
                continue;
            }
            int v=adj[u][i++];
            if(v==up[0][u])continue;
            up[0][v]=u,dep[v]=dep[u]+1;
            st.push_back({v,0});
        }
    }
    bool anc(int u,int v)
    {
        return dfn[u]<=dfn[v]&&ed[v]<=ed[u];
    }
    int lca(int u,int v)
    {
        if(anc(u,v))return u;
        for(int j=lg-1;j>=0;j--)if(up[j][u]&&!anc(up[j][u],v))u=up[j][u];
        return up[0][u];
    }
    // 栈构建虚树，DP 求所有关键点对距离和；O(k log k+k log n)
    ll query(vector<int> a)
    {
        if(a.empty())return 0;
        sort(a.begin(),a.end(),[&](int u,int v){return dfn[u]<dfn[v];});
        a.erase(unique(a.begin(),a.end()),a.end());
        int k=a.size();
        for(int u:a)mark[u]=1;
        for(int i=1;i<k;i++)a.push_back(lca(a[i-1],a[i]));
        sort(a.begin(),a.end(),[&](int u,int v){return dfn[u]<dfn[v];});
        a.erase(unique(a.begin(),a.end()),a.end());
        vector<int> st;
        for(int u:a)tr[u].clear(),cnt[u]=0;
        for(int u:a)
        {
            while(!st.empty()&&!anc(st.back(),u))st.pop_back();
            if(!st.empty())tr[st.back()].push_back(u);
            st.push_back(u);
        }
        ll ans=0;
        for(int i=(int)a.size()-1;i>=0;i--)
        {
            int u=a[i];
            cnt[u]+=mark[u];
            for(int v:tr[u])
            {
                ans+=1LL*(dep[v]-dep[u])*cnt[v]*(k-cnt[v]);
                cnt[u]+=cnt[v];
            }
            mark[u]=0;
        }
        return ans;
    }
};
\end{lstlisting}
% SOURCE: 04-图论/Kruskal重构树.cpp
\subsection{Kruskal重构树}

\note{对象edges为带权原边，ch为重构树儿子，val为合并阈值，cnt为原点数量；点1.\allowbreak{}.\allowbreak{}n，内部节点另外编号。lca跨连通块为0，同点叶子阈值不是实际路径边权；换边权重新构建。}
\note{无向图按边权升序合并：原点为叶子，新结点权为合并边权；支持重边、负权、不连通。 lca(\allowbreak{}u,\allowbreak{}v)\allowbreak{} 的权 =\allowbreak{} 两点路径中最大边权的最小值；异连通块返回 0，同点无边须另处理。 component(\allowbreak{}u,\allowbreak{}w)\allowbreak{} 返回只保留边权<\allowbreak{}=\allowbreak{}w 时 u 所在连通块的树根，cnt[\allowbreak{}root]\allowbreak{} 为原点数。 构建 O(\allowbreak{}m log m+\allowbreak{}n log n)\allowbreak{}，查询 O(\allowbreak{}log n)\allowbreak{}，空间 O(\allowbreak{}m+\allowbreak{}n log n)\allowbreak{}；静态图，修改边后重建。}
\begin{lstlisting}
struct KruskalTree
{
    struct Edge{int u,v;ll w;};
    int n,tot,lg;
    vector<vector<int>> ch,up;
    vector<int> dep,root,cnt;
    vector<ll> val;
    // 原点 1..n，n>=1；传入边副本，不修改调用者。
    KruskalTree(int n,vector<Edge> e):n(n),tot(n),lg(1),ch(2*n),dep(2*n),root(2*n),cnt(2*n,1),val(2*n,LLONG_MIN)
    {
        assert(n>=1);
        vector<int> fa(n+1),sz(n+1,1),node(n+1),par(2*n);
        iota(fa.begin(),fa.end(),0),iota(node.begin(),node.end(),0);
        auto find=[&](int x){while(x!=fa[x])x=fa[x]=fa[fa[x]];return x;};
        sort(e.begin(),e.end(),[](Edge a,Edge b){return a.w<b.w;});
        for(auto [u,v,w]:e)
        {
            int a=find(u),b=find(v);if(a==b)continue;
            int x=node[a],y=node[b];++tot;
            ch[tot]={x,y},par[x]=par[y]=tot,val[tot]=w,cnt[tot]=cnt[x]+cnt[y];
            if(sz[a]<sz[b])swap(a,b);
            fa[b]=a,sz[a]+=sz[b],node[a]=tot;
        }
        while((1LL<<lg)<=tot)++lg;
        up.assign(lg,vector<int>(tot+1));
        // 父编号总大于儿子，倒序即可处理整片森林，避免深链递归。
        for(int u=tot;u>=1;--u)
        {
            up[0][u]=par[u],dep[u]=dep[par[u]]+1;
            root[u]=par[u]?root[par[u]]:u;
            for(int j=1;j<lg;++j)up[j][u]=up[j-1][up[j-1][u]];
        }
    }
    int lca(int u,int v)const
    {
        if(root[u]!=root[v])return 0;
        if(dep[u]<dep[v])swap(u,v);
        for(int j=lg-1;j>=0;--j)if((dep[u]-dep[v])>>j&1)u=up[j][u];
        if(u==v)return u;
        for(int j=lg-1;j>=0;--j)if(up[j][u]!=up[j][v])u=up[j][u],v=up[j][v];
        return up[0][u];
    }
    int component(int u,ll w)const
    {
        for(int j=lg-1;j>=0;--j)if(up[j][u]&&val[up[j][u]]<=w)u=up[j][u];
        return u;
    }
};
\end{lstlisting}
\note{例：KruskalTree t(\allowbreak{}3,\allowbreak{}\{\{1,\allowbreak{}2,\allowbreak{}5\},\allowbreak{}\{2,\allowbreak{}3,\allowbreak{}8\}\})\allowbreak{}; t.\allowbreak{}val[\allowbreak{}t.\allowbreak{}lca(\allowbreak{}1,\allowbreak{}3)\allowbreak{}]\allowbreak{}=\allowbreak{}=\allowbreak{}8。 同权边会形成多个同权祖先；阈值查询要爬到最高合法祖先，不能停在第一次合并。}
% SOURCE: 04-图论/线段树优化建图.cpp
\subsection{线段树优化建图}

\topic{子树同深度约束压成区间}
\note{DFS 序用半开子树区间 $[tin_u,tout_u)$；按绝对深度分桶，桶内按 $tin$ 排序。目标深度桶中两次 \texttt{lower\_bound} 定位子树集合，用桶内编号作线段树区间坐标。各桶总长度 $O(n)$，不需开“深度乘点数”个节点。}
\note{禁止 $u$ 与区间内任一 $x$ 同选时，蕴含为 $u\to\neg x$ 和 $x\to\neg u$，分别压成点到区间、区间到点。辅助节点仅压缩可达性，不代表布尔变量，不可随意用 xor 1 取反；需排除 $u$ 时拆开它所在位置（2025 武汉 B）。}

\note{对象g存\{目标,\allowbreak{}权\}，原点1.\allowbreak{}.\allowbreak{}n，辅助点追加；两棵树分别处理点到区间、区间到点。范围为闭区间，最短路按g.\allowbreak{}size(\allowbreak{})\allowbreak{}分配状态；流量、蕴含等应用还需遵守各模型的边语义。}
\note{原点 1.\allowbreak{}.\allowbreak{}n；闭区间连边。出树父-\allowbreak{}>\allowbreak{}子，入树子-\allowbreak{}>\allowbreak{}父；两树在叶子共用原点。 建图 O(\allowbreak{}n)\allowbreak{}，一次点/\allowbreak{}区间或区间/\allowbreak{}区间连边 O(\allowbreak{}log n)\allowbreak{}；共 O(\allowbreak{}n+\allowbreak{}q)\allowbreak{} 点、O(\allowbreak{}n+\allowbreak{}q log n)\allowbreak{} 边。 g 是 vector 邻接表，边为 \{目标,\allowbreak{}权重\}；Dijkstra 时所有权重非负，距离数组按 g.\allowbreak{}size(\allowbreak{})\allowbreak{} 分配。}
\begin{lstlisting}
struct RangeGraph
{
    int n;
    vector<int> in,out;
    vector<vector<pair<int,ll>>> g;
    RangeGraph(int n):n(n),in(4*n+4),out(4*n+4),g(n+1){assert(n>0);build(1,1,n);}
    int new_node(){g.emplace_back();return (int)g.size()-1;}
    void add_edge(int u,int v,ll w){g[u].push_back({v,w});}
    void build(int p,int l,int r)
    {
        if(l==r){in[p]=out[p]=l;return;}
        in[p]=new_node(),out[p]=new_node();int m=(l+r)/2;
        build(p*2,l,m),build(p*2+1,m+1,r);
        for(int c:{p*2,p*2+1})add_edge(out[p],out[c],0),add_edge(in[c],in[p],0);
    }
    // cover(u,L,R,w,type)：type=0 为 u->区间，1 为区间->u；u 可是辅助点。
    void cover(int p,int l,int r,int u,int L,int R,ll w,int type)
    {
        if(L<=l&&r<=R){if(type)add_edge(in[p],u,w);else add_edge(u,out[p],w);return;}
        int m=(l+r)/2;
        if(L<=m)cover(p*2,l,m,u,L,R,w,type);
        if(R>m)cover(p*2+1,m+1,r,u,L,R,w,type);
    }
    void point_to_range(int u,int l,int r,ll w)
    {
        if(l>r)return;
        assert(1<=l&&r<=n);cover(1,1,n,u,l,r,w,0);
    }
    void range_to_point(int l,int r,int u,ll w)
    {
        if(l>r)return;
        assert(1<=l&&r<=n);cover(1,1,n,u,l,r,w,1);
    }
    void range_to_range(int l,int r,int L,int R,ll w)
    {
        if(l>r||L>R)return;
        int u=new_node();range_to_point(l,r,u,0),point_to_range(u,L,R,w);
    }
};
\end{lstlisting}
\note{区间-\allowbreak{}>\allowbreak{}区间经过一个新中转点，整条路径只计一次 w，避免 O(\allowbreak{}log\ensuremath{^2} n)\allowbreak{} 两两连边。 用于流：辅助树边须改为 INF 容量，不能把这里的 0 权当作 0 容量。 用于 2-\allowbreak{}SAT：辅助点只压缩可达性，仍要补齐原变量约束的逆否蕴含，不能给辅助点随意配否定。 子树同深度：tin/\allowbreak{}tout 用半开区间 [\allowbreak{}tin[\allowbreak{}u]\allowbreak{},\allowbreak{}tout[\allowbreak{}u]\allowbreak{})\allowbreak{}，按绝对深度把 tin 排进各桶。 在目标深度桶中 lower\_bound(\allowbreak{}tin[\allowbreak{}u]\allowbreak{})\allowbreak{}、lower\_bound(\allowbreak{}tout[\allowbreak{}u]\allowbreak{})\allowbreak{} 即所求连续段，O(\allowbreak{}log n)\allowbreak{}。 桶内顺序作为区间坐标；若按目标深度分组建图，各桶总长度 O(\allowbreak{}n)\allowbreak{}，无需开 深度*n 个点。 禁止 u 与区间内 x 同选：同时连 u-\allowbreak{}>\allowbreak{}非x、x-\allowbreak{}>\allowbreak{}非u 两个方向，可分别压成点到区间和区间到点。 图上的辅助点不代表布尔变量，不可对它按 xor 1 取反；具体题目排除 u 时拆开对应位置。}
% SOURCE: 04-图论/圆方树.cpp
\subsection{圆方树}

\note{原点与点双分量节点共同组成圆方森林；访问序和low用于提取点双，栈记录尚未归属的边/\allowbreak{}点。原图编号与新增方点编号必须区分，断连图逐块建立，割点可连接多个方点。}
\note{无向图点双圆方森林，O(\allowbreak{}n+\allowbreak{}m)\allowbreak{}；忽略自环，支持重边 Tarjan 使用递归，大链应增大栈或改成显式栈}
\begin{lstlisting}
struct BlockTree
{
    int n,num,tim;
    vector<vector<pair<int,int>>> adj;
    vector<vector<int>> tr;
    vector<int> dfn,low,st;
    BlockTree(int n):n(n),num(n),tim(0),adj(n+1),tr(2*n+1),dfn(n+1),low(n+1){}
    void add_edge(int u,int v,int id)
    {
        if(u==v)return;
        adj[u].push_back({v,id}),adj[v].push_back({u,id});
    }
    void dfs(int u,int pe)
    {
        dfn[u]=low[u]=++tim,st.push_back(u);
        for(pair<int,int> e:adj[u])
        {
            int v=e.first,id=e.second;
            if(id==pe)continue;
            if(!dfn[v])
            {
                dfs(v,id),low[u]=min(low[u],low[v]);
                if(low[v]>=dfn[u])
                {
                    ++num;
                    tr[num].push_back(u),tr[u].push_back(num);
                    int x;
                    do
                    {
                        x=st.back(),st.pop_back();
                        tr[num].push_back(x),tr[x].push_back(num);
                    }while(x!=v);
                }
            }
            else low[u]=min(low[u],dfn[v]);
        }
    }
    void build()
    {
        for(int u=1;u<=n;u++)if(!dfn[u])dfs(u,-1),st.pop_back();
    }
    // 圆方树路径上的原点即必经点，包含两个端点；不连通返回 -1
    // 单次 O(n)，多询问可预处理 LCA 与路径前缀和至 O(log n)
    int query(int s,int t)
    {
        vector<int> fa(num+1,-1);
        queue<int> q;
        fa[s]=0,q.push(s);
        while(!q.empty())
        {
            int u=q.front();
            q.pop();
            for(int v:tr[u])if(fa[v]<0)fa[v]=u,q.push(v);
        }
        if(fa[t]<0)return -1;
        int ans=0;
        for(int u=t;u;u=fa[u])ans+=u<=n;
        return ans;
    }
};
\end{lstlisting}
% SOURCE: 04-图论/长链剖分.cpp
\subsection{长链剖分}

\note{对象adj为树，len/\allowbreak{}son选深度最长儿子，off定位共享深度桶，buf存各相对深度数量，ans存最优深度。并列取更小深度；长儿子继承桶，轻儿子另合并，重建时清缓冲与偏移。}
\note{O(\allowbreak{}n)\allowbreak{}，长儿子继承深度桶，求子树中结点最多的相对深度；并列取小}
\begin{lstlisting}
struct LongChain
{
    int n,tot;
    vector<vector<int> > adj;
    vector<int> fa,len,son,off,ans,buf,ord;
    LongChain(int n):n(n),tot(0),adj(n+1),fa(n+1),len(n+1),son(n+1),off(n+1),ans(n+1),buf(n+1){}
    void add_edge(int u,int v)
    {
        adj[u].push_back(v),adj[v].push_back(u);
    }
    vector<int> run()
    {
        ord.assign(1,1),fa[1]=0,tot=0;
        fill(buf.begin(),buf.end(),0);
        for(int i=0;i<(int)ord.size();i++)for(int v:adj[ord[i]])if(v!=fa[ord[i]])fa[v]=ord[i],ord.push_back(v);
        for(int i=n-1;i>=0;i--)
        {
            int u=ord[i];
            len[u]=1,son[u]=0;
            for(int v:adj[u])if(fa[v]==u&&len[v]+1>len[u])len[u]=len[v]+1,son[u]=v;
        }
        off[1]=tot,tot+=len[1];
        for(int u:ord)for(int v:adj[u])if(fa[v]==u)
        {
            if(v==son[u])off[v]=off[u]+1;
            else off[v]=tot,tot+=len[v];
        }
        for(int i=n-1;i>=0;i--)
        {
            int u=ord[i],p=off[u];
            buf[p]=1,ans[u]=son[u]?ans[son[u]]+1:0;
            for(int v:adj[u])if(fa[v]==u&&v!=son[u])
                for(int j=0;j<len[v];j++)
                {
                    buf[p+j+1]+=buf[off[v]+j];
                    int d=j+1;
                    if(buf[p+d]>buf[p+ans[u]]||(buf[p+d]==buf[p+ans[u]]&&d<ans[u]))ans[u]=d;
                }
            if(buf[p]==buf[p+ans[u]])ans[u]=0;
        }
        return ans;
    }
};
\end{lstlisting}
% SOURCE: 04-图论/KM算法(最大权匹配).cpp
\subsection{KM算法(最大权匹配)}

\note{w[\allowbreak{}i]\allowbreak{}[\allowbreak{}j]\allowbreak{}为左右各n点的收益，缺边设足够小值且须存在真实完美匹配；lx/\allowbreak{}ly为顶标，slack为右侧松弛，match[\allowbreak{}v]\allowbreak{}为右点伴侣。visx/\allowbreak{}visy为交错树访问状态，结果仍需检查是否选中缺边。}
\begin{lstlisting}
const int N=505;
const ll INF=1e18;
int n;// 左右各 n 个点
ll w[N][N];// 边权；不成边的位置先设成很小的数（如 -1e12），要求存在完美匹配
ll lx[N],ly[N],slack[N];// 左右顶标、右部点的松弛量
int match[N];// match[v]：右部点 v 匹配的左部点
int visx[N],visy[N];// 本次增广中，左右部点是否在交错树上
int dfs(int u)// 从左部点 u 出发增广，只走 lx[u]+ly[v]==w[u][v] 的相等边
{
    visx[u]=1;
    for(int v=1;v<=n;v++)
    {
        if(visy[v])continue;
        ll d=lx[u]+ly[v]-w[u][v];
        if(d==0)
        {
            visy[v]=1;
            if(!match[v]||dfs(match[v]))// v 空着，或者能让它的搭档换一条路
            {
                match[v]=u;
                return 1;
            }
        }
        else if(d<slack[v])slack[v]=d;// 记下最小差值，供顶标调整用
    }
    return 0;
}
ll km()// 二分图最大权完美匹配，O(n^4) 最坏；本版反复 DFS，需 O(n^3) 时用保留交错树的增广写法
{
    for(int i=1;i<=n;i++)
    {
        lx[i]=-INF,ly[i]=0,match[i]=0;
        for(int j=1;j<=n;j++)lx[i]=max(lx[i],w[i][j]);// 左顶标取该行最大边权
    }
    for(int u=1;u<=n;u++)
    {
        for(int i=1;i<=n;i++)slack[i]=INF,visx[i]=0,visy[i]=0;
        while(!dfs(u))// 增广失败就调顶标，直到能找到相等边
        {
            ll d=INF;
            for(int i=1;i<=n;i++)if(!visy[i])d=min(d,slack[i]);
            for(int i=1;i<=n;i++)
            {
                if(visx[i])lx[i]-=d;
                if(visy[i])ly[i]+=d;
                else slack[i]-=d;
            }
            for(int i=1;i<=n;i++)visx[i]=0,visy[i]=0;
        }
    }
    ll ans=0;
    for(int v=1;v<=n;v++)ans+=w[match[v]][v];
    return ans;
}
\end{lstlisting}
% SOURCE: 04-图论/带权二分图匹配Hungarian.cpp
\subsection{带权二分图匹配Hungarian}

\note{c[\allowbreak{}i]\allowbreak{}[\allowbreak{}j]\allowbreak{}为左i分配给右j的费用，行n、列m且n<\allowbreak{}=\allowbreak{}m，均0起。返回\{最小总费用,\allowbreak{}每行匹配列\}；u,\allowbreak{}v为势能，p为列匹配，way记增广前驱，每次重建。缺边设大费用，最后检查是否选中。}
\note{最小费用匹配：n 行、m 列，n<\allowbreak{}=\allowbreak{}m，所有行都要分配不同列；O(\allowbreak{}n\ensuremath{^2}m)\allowbreak{}，0-\allowbreak{}indexed。 完整费用矩阵，负费用允许；缺边设大数并在结果中检查是否选中。最大权取负费用。 返回 \{费用,\allowbreak{}每行匹配列\}；费用及势能/\allowbreak{}差值须放入 ll，建议 |cost|*n <\allowbreak{}<\allowbreak{} 1e18。}
\begin{lstlisting}
vector<vector<ll>> c;
pair<ll,vector<int>> hungarian()
{
    int n=c.size(),m=n?c[0].size():0;assert(n<=m);
    vector<ll> u(n+1),v(m+1);vector<int> p(m+1),way(m+1);
    for(int i=1;i<=n;i++)
    {
        p[0]=i;int j0=0;vector<ll> d(m+1,LLONG_MAX/4);vector<bool> used(m+1);
        do
        {
            used[j0]=true;int i0=p[j0],j1=0;ll delta=LLONG_MAX/4;
            for(int j=1;j<=m;j++)if(!used[j])
            {ll x=c[i0-1][j-1]-u[i0]-v[j];if(x<d[j])d[j]=x,way[j]=j0;if(d[j]<delta)delta=d[j],j1=j;}
            for(int j=0;j<=m;j++)if(used[j])u[p[j]]+=delta,v[j]-=delta;else d[j]-=delta;
            j0=j1;
        }while(p[j0]);
        do{int j1=way[j0];p[j0]=p[j1];j0=j1;}while(j0);
    }
    vector<int> match(n);for(int j=1;j<=m;j++)if(p[j])match[p[j]-1]=j-1;
    return {-v[0],match};
}
\end{lstlisting}
% SOURCE: 04-图论/二分图匹配HopcroftKarp.cpp
\subsection{二分图匹配HopcroftKarp}

\note{对象n/\allowbreak{}m分别为左右点数，两侧0起；g存左侧出边，a/\allowbreak{}b分别为左右伴侣，-\allowbreak{}1未匹配。距离层和当前弧每轮重建，solve在当前匹配上继续增广，换图重新构造。}
\note{左侧 n、右侧 m，均从 0 编号；O(\allowbreak{}E sqrt(\allowbreak{}V)\allowbreak{})\allowbreak{}，vector 邻接表只存左到右。 bfs 建最短增广层次，dfs 每轮只找最短路径；返回匹配数，a/\allowbreak{}b 保存双方匹配，-\allowbreak{}1 未匹配。}
\begin{lstlisting}
struct HopcroftKarp
{
    int n,m,limit;vector<vector<int>> g;vector<int> a,b,d,it;
    HopcroftKarp(int n,int m):n(n),m(m),g(n),a(n,-1),b(m,-1),d(n),it(n){}
    bool bfs()
    {
        queue<int> q;fill(d.begin(),d.end(),-1);limit=INT_MAX;
        for(int u=0;u<n;u++)if(a[u]<0)d[u]=0,q.push(u);
        while(!q.empty())
        {
            int u=q.front();q.pop();if(d[u]>=limit)continue;
            for(int v:g[u])if(b[v]<0)limit=d[u]+1;else if(d[b[v]]<0)d[b[v]]=d[u]+1,q.push(b[v]);
        }
        return limit<INT_MAX;
    }
    bool dfs(int u)
    {
        for(int &i=it[u];i<(int)g[u].size();i++)
        {
            int v=g[u][i];
            if((b[v]<0&&d[u]+1==limit)||(b[v]>=0&&d[b[v]]==d[u]+1&&dfs(b[v])))
            {a[u]=v;b[v]=u;return true;}
        }
        d[u]=-1;return false;
    }
    int solve(){int ans=count_if(a.begin(),a.end(),[](int x){return x>=0;});while(bfs()){fill(it.begin(),it.end(),0);for(int u=0;u<n;u++)if(a[u]<0&&dfs(u))ans++;}return ans;}
};
\end{lstlisting}
% SOURCE: 04-图论/树同构(AHU).cpp
\subsection{树同构(AHU)}

\note{a/\allowbreak{}b是两棵独立无向树的邻接表，点1起；ids把排序后的子树类型精确映射成编号。比较两树必须共享ids，不能各自编号后直接比较整数；因此保留两个树参数。}
\note{无向无标号树，vector 邻接表，点 1.\allowbreak{}.\allowbreak{}n；输入必须为合法树，空树用长度 1 的 adj。 根的类型=\allowbreak{}排序后的子树类型列表；共享字典精确编号，没有随机哈希碰撞。 无根树从重心取根，最多两个；O(\allowbreak{}n log n)\allowbreak{}，空间 O(\allowbreak{}n)\allowbreak{}，遍历用栈避免长链递归。}
\begin{lstlisting}
vector<int> tree_centroids(const vector<vector<int>> &adj)
{
    int n=(int)adj.size()-1;
    if(!n)return {};
    vector<int> fa(n+1),sz(n+1,1),order{1},ans;
    for(int i=0;i<(int)order.size();i++)
    {
        int u=order[i];
        for(int v:adj[u])if(v!=fa[u])fa[v]=u,order.push_back(v);
    }
    for(int i=n-1;i>=0;i--)
    {
        int u=order[i],mx=0;
        for(int v:adj[u])if(v!=fa[u])sz[u]+=sz[v],mx=max(mx,sz[v]);
        if(max(mx,n-sz[u])<=n/2)ans.push_back(u);
    }
    return ans;
}
int tree_code(const vector<vector<int>> &adj,int root,map<vector<int>,int> &ids)
{
    int n=(int)adj.size()-1;
    vector<int> fa(n+1),code(n+1),order{root};
    for(int i=0;i<n;i++)
    {
        int u=order[i];
        for(int v:adj[u])if(v!=fa[u])fa[v]=u,order.push_back(v);
    }
    for(int i=n-1;i>=0;i--)
    {
        int u=order[i];
        vector<int> children;
        for(int v:adj[u])if(v!=fa[u])children.push_back(code[v]);
        sort(children.begin(),children.end());
        auto [it,inserted]=ids.emplace(move(children),ids.size()+1);
        code[u]=it->second;
    }
    return code[root];
}
bool same_tree(const vector<vector<int>> &a,const vector<vector<int>> &b)
{
    assert(!a.empty()&&!b.empty());
    if(a.size()!=b.size())return false;
    if(a.size()==1)return true;
    auto ca=tree_centroids(a),cb=tree_centroids(b);
    if(ca.size()!=cb.size())return false;
    map<vector<int>,int> ids;
    int x=tree_code(a,ca[0],ids);
    for(int root:cb)if(tree_code(b,root,ids)==x)return true;
    return false;
}
\end{lstlisting}
% SOURCE: 04-图论/Pruefer序列.cpp
\subsection{Pruefer序列}

\note{g是0起无向标号树，n>\allowbreak{}=\allowbreak{}2；prufer\_encode(\allowbreak{})\allowbreak{}生成全局code并返回，长度n-\allowbreak{}2。prufer\_decode(\allowbreak{})\allowbreak{}读code，重建g并返回；度数和叶子堆每次初始化。编码不改原图，解码覆盖g。}
\note{标号树与 Prüfer 序列的双射：n>\allowbreak{}=\allowbreak{}2，编号 0.\allowbreak{}.\allowbreak{}n-\allowbreak{}1，序列长度 n-\allowbreak{}2，O(\allowbreak{}n log n)\allowbreak{}。 每次删编号最小叶子、输出其邻居；点出现次数=\allowbreak{}度数-\allowbreak{}1。Cayley：n\textasciicircum{}(\allowbreak{}n-\allowbreak{}2)\allowbreak{} 棵标号树。 给定度数 d\_i>\allowbreak{}=\allowbreak{}1 且 \ensuremath{\Sigma}d\_i=\allowbreak{}2n-\allowbreak{}2，树数量为 (\allowbreak{}n-\allowbreak{}2)\allowbreak{}!/\allowbreak{}\ensuremath{\Pi}(\allowbreak{}d\_i-\allowbreak{}1)\allowbreak{}!。}
\begin{lstlisting}
vector<vector<int>> g;
vector<int> code;
vector<int> prufer_encode()
{
    int n=g.size();assert(n>=2);vector<int>d(n);code.clear();priority_queue<int,vector<int>,greater<int>> q;
    for(int i=0;i<n;i++){d[i]=g[i].size();if(d[i]==1)q.push(i);}
    for(int k=0;k<n-2;k++){int u=q.top();q.pop();d[u]=0;int v=-1;for(int x:g[u])if(d[x]){v=x;break;}assert(v>=0);code.push_back(v);if(--d[v]==1)q.push(v);}
    return code;
}
vector<vector<int>> prufer_decode()
{
    int n=code.size()+2;vector<int>d(n,1);for(int x:code){assert(x>=0&&x<n);d[x]++;}
    priority_queue<int,vector<int>,greater<int>> q;for(int i=0;i<n;i++)if(d[i]==1)q.push(i);g.assign(n,{});
    auto edge=[&](int u,int v){g[u].push_back(v);g[v].push_back(u);};
    for(int v:code){int u=q.top();q.pop();edge(u,v);if(--d[v]==1)q.push(v);}
    int u=q.top();q.pop();edge(u,q.top());return g;
}
\end{lstlisting}
% SOURCE: 04-图论/动态点分治.cpp
\subsection{动态点分治}

\note{对象path[\allowbreak{}u]\allowbreak{}存u到各层重心的距离，on为是否激活，bag[\allowbreak{}c]\allowbreak{}保存激活点到重心c的距离多重集。构造时自动预处理，再set\_active/\allowbreak{}query；重复激活不会叠加，取消时仅删一个副本，无激活点返回INF。}
\note{点分树：带非负边权树，动态点亮/\allowbreak{}熄灭，查询距最近亮点距离；空集合返回 INF。 0-\allowbreak{}indexed；预存每点到重心祖先的距离，构建 O(\allowbreak{}n log n)\allowbreak{}，切换 O(\allowbreak{}log\ensuremath{^2}n)\allowbreak{}，查询 O(\allowbreak{}log n)\allowbreak{}。 每个重心用 multiset 保存亮点距离；删除一个副本，重复距离不能全删。建树遍历迭代。}
\begin{lstlisting}
struct DynamicCentroid
{
    using ll=long long;static constexpr ll INF=LLONG_MAX/4;
    vector<vector<pair<int,ll>>> g,path;vector<bool> removed,on;vector<int> parent,size;vector<multiset<ll>> bag;
    DynamicCentroid(vector<vector<pair<int,ll>>> g):g(g),path(g.size()),removed(g.size()),on(g.size()),parent(g.size()),size(g.size()),bag(g.size()){if(!g.empty())build(0);}
    void build(int entry)
    {
        vector<int> order={entry};parent[entry]=-1;
        for(int i=0;i<(int)order.size();i++){int u=order[i];for(auto [v,w]:g[u])if(!removed[v]&&v!=parent[u])parent[v]=u,order.push_back(v);}
        int c=entry,total=order.size(),best=total;
        for(auto it=order.rbegin();it!=order.rend();it++){int u=*it;size[u]=1;int largest=0;for(auto [v,w]:g[u])if(!removed[v]&&parent[v]==u)size[u]+=size[v],largest=max(largest,size[v]);largest=max(largest,total-size[u]);if(largest<best)best=largest,c=u;}
        vector<tuple<int,int,ll>> q={{c,-1,0}};
        for(int i=0;i<(int)q.size();i++){auto [u,p,d]=q[i];path[u].push_back({c,d});for(auto [v,w]:g[u])if(!removed[v]&&v!=p)q.push_back({v,u,d+w});}
        removed[c]=true;for(auto [v,w]:g[c])if(!removed[v])build(v);
    }
    void set_active(int u,bool active){if(on[u]==active)return;on[u]=active;for(auto [c,d]:path[u])if(active)bag[c].insert(d);else bag[c].erase(bag[c].find(d));}
    ll query(int u)const{ll ans=INF;for(auto [c,d]:path[u])if(!bag[c].empty())ans=min(ans,d+*bag[c].begin());return ans;}
};
\end{lstlisting}
% SOURCE: 04-图论/曼哈顿最小生成树.cpp
\subsection{曼哈顿最小生成树}

\note{p是0起整数坐标点集，返回\{权,\allowbreak{}u,\allowbreak{}v\}候选边；只生成足够求MST的稀疏边，最终还需Kruskal。四方向变换在副本上进行，原点编号保留，坐标差与距离用ll。}
\note{二维点的曼哈顿距离完全图 MST：O(\allowbreak{}n log n)\allowbreak{} 生成候选边，再 Kruskal。 顶点 0.\allowbreak{}.\allowbreak{}n-\allowbreak{}1；返回 \{距离,\allowbreak{}u,\allowbreak{}v\}，允许重合点。坐标绝对值及距离须远离 ll 上界。 方向扫描参考 KACTL ManhattanMST.\allowbreak{}h（CC0）；点按值传入，不改变原坐标。}
\begin{lstlisting}
vector<array<ll,3>> manhattan_edges(vector<pair<ll,ll>> p)
{
    int n=p.size();vector<int> id(n);iota(id.begin(),id.end(),0);
    vector<array<ll,3>> e;
    for(int k=0;k<4;k++)
    {
        sort(id.begin(),id.end(),[&](int a,int b){return (__int128)p[a].first+p[a].second<(__int128)p[b].first+p[b].second;});
        map<ll,int> sweep;
        for(int i:id)
        {
            for(auto it=sweep.lower_bound(-p[i].second);it!=sweep.end();)
            {
                int j=it->second;ll dx=p[i].first-p[j].first,dy=p[i].second-p[j].second;
                if(dy>dx)break;
                e.push_back({dx+dy,i,j});it=sweep.erase(it);
            }
            sweep[-p[i].second]=i;
        }
        for(auto &[x,y]:p)if(k&1)x=-x;else swap(x,y);
    }
    return e;
}
\end{lstlisting}
% SOURCE: 04-图论/全局最小割StoerWagner.cpp
\subsection{全局最小割StoerWagner}

\note{w是对称非负权邻接矩阵，重边权相加，对角置0；返回\{割权,\allowbreak{}一侧原点编号\}。工作副本在收缩时改写，原矩阵不变；少于2点返回零和空侧。}
\note{无向非负权图全局最小割：O(\allowbreak{}n\ensuremath{^3})\allowbreak{} 时间，O(\allowbreak{}n\ensuremath{^2})\allowbreak{} 空间；允许重边合并。 输入对称邻接矩阵，0.\allowbreak{}.\allowbreak{}n-\allowbreak{}1，自环忽略；n<\allowbreak{}2 返回 0。返回割值及割的一侧。 每轮最大邻接序末尾两个点合并；不是只求固定 s-\allowbreak{}t 的割。}
\begin{lstlisting}
pair<ll,vector<int>> global_min_cut(vector<vector<ll>> w)
{
    int n=w.size();if(n<2)return {0,{}};
    vector<int> v(n);iota(v.begin(),v.end(),0);
    vector<vector<int>> group(n);for(int i=0;i<n;i++)group[i]={i};
    ll ans=LLONG_MAX;vector<int> cut;
    while(v.size()>1)
    {
        vector<ll> d(n);vector<bool> used(n);int prev=-1;
        for(int k=0;k<(int)v.size();k++)
        {
            int t=-1;for(int u:v)if(!used[u]&&(t<0||d[u]>d[t]))t=u;
            if(k+1==(int)v.size())
            {
                if(d[t]<ans)ans=d[t],cut=group[t];
                for(int u:v)if(u!=prev&&u!=t)w[prev][u]+=w[t][u],w[u][prev]=w[prev][u];
                group[prev].insert(group[prev].end(),group[t].begin(),group[t].end());
                v.erase(find(v.begin(),v.end(),t));break;
            }
            used[t]=true;prev=t;for(int u:v)if(!used[u])d[u]+=w[t][u];
        }
    }
    return {ans,cut};
}
\end{lstlisting}
% SOURCE: 04-图论/割树GomoryHu.cpp
\subsection{割树GomoryHu}

\note{点0.\allowbreak{}.\allowbreak{}n-\allowbreak{}1，edges每项为\{u,\allowbreak{}v,\allowbreak{}非负容量\}，表示无向边。返回n-\allowbreak{}1条\{u,\allowbreak{}v,\allowbreak{}割值\}；树路上的最小边权等于原图两点最小割。每次求割重建CutDinic实例和残量图，不复用上轮容量。}
\note{无向非负容量图的割树：n-\allowbreak{}1 次最大流，编号 0.\allowbreak{}.\allowbreak{}n-\allowbreak{}1；每对最小割是树路径最小边。 返回 \{u,\allowbreak{}v,\allowbreak{}割值\} 的 n-\allowbreak{}1 条边；不连通时允许 0 权边。每轮重建残量图。 内置 vector Dinic；流量总和须放入 ll，递归推流深度 O(\allowbreak{}n)\allowbreak{}。}
\begin{lstlisting}
struct CutDinic
{
    struct E{int v,rev;ll c;};vector<vector<E>> g;vector<int> d,it;
    CutDinic(int n):g(n),d(n),it(n){}
    void add(int u,int v,ll c){if(u==v)return;int a=g[u].size(),b=g[v].size();g[u].push_back({v,b,c});g[v].push_back({u,a,c});}
    bool bfs(int s,int t){fill(d.begin(),d.end(),-1);queue<int> q;d[s]=0;q.push(s);while(!q.empty()){int u=q.front();q.pop();for(auto e:g[u])if(e.c&&d[e.v]<0)d[e.v]=d[u]+1,q.push(e.v);}return d[t]>=0;}
    ll dfs(int u,int t,ll f){if(u==t)return f;for(int &i=it[u];i<(int)g[u].size();i++){auto &e=g[u][i];if(e.c&&d[e.v]==d[u]+1){ll z=dfs(e.v,t,min(f,e.c));if(z){e.c-=z;g[e.v][e.rev].c+=z;return z;}}}return 0;}
    ll flow(int s,int t){ll ans=0;while(bfs(s,t)){fill(it.begin(),it.end(),0);while(ll z=dfs(s,t,LLONG_MAX/4))ans+=z;}return ans;}
};
int n;
vector<array<ll,3>> edges;
vector<array<ll,3>> gomory_hu()
{
    vector<int> p(n);vector<array<ll,3>> tree;
    for(int s=1;s<n;s++)
    {
        CutDinic f(n);for(auto [u,v,c]:edges)f.add(u,v,c);
        ll cut=f.flow(s,p[s]);tree.push_back({s,p[s],cut});
        for(int j=s+1;j<n;j++)if(p[j]==p[s]&&f.d[j]>=0)p[j]=s;
    }
    return tree;
}
\end{lstlisting}
% SOURCE: 04-图论/斯坦纳树.cpp
\subsection{斯坦纳树}

\note{g[\allowbreak{}u]\allowbreak{}存\{v,\allowbreak{}非负边权\}，点0起；terminal列出要连通的终端。dp[\allowbreak{}S]\allowbreak{}[\allowbreak{}v]\allowbreak{}是连接终端集合S且包含v的最小代价，每次按终端数重建；返回全部终端的最小代价，不连通为INF，空终端返回0。}
\note{无向非负权图，连接 k 个指定终端的最小代价子图；k 小，顶点 0.\allowbreak{}.\allowbreak{}n-\allowbreak{}1。 dp[\allowbreak{}S]\allowbreak{}[\allowbreak{}v]\allowbreak{}：连接终端集合 S、包含 v；子集合并后做多源 Dijkstra。 O(\allowbreak{}3\textasciicircum{}k*n +\allowbreak{} 2\textasciicircum{}k*(\allowbreak{}n+\allowbreak{}m)\allowbreak{}log n)\allowbreak{}，空间 O(\allowbreak{}2\textasciicircum{}k*n)\allowbreak{}；不连通返回 INF，k=\allowbreak{}0 返回 0。}
\begin{lstlisting}
vector<vector<pair<int,ll>>> g;
vector<int> terminal;
vector<vector<ll>> dp;
ll steiner_tree()
{
    const ll INF=LLONG_MAX/4;int n=g.size(),k=terminal.size();
    dp.clear();if(!k)return 0;assert(k<25);int lim=1<<k;
    dp.assign(lim,vector<ll>(n,INF));
    for(int i=0;i<k;i++)dp[1<<i][terminal[i]]=0;
    for(int s=1;s<lim;s++)
    {
        for(int a=(s-1)&s;a;a=(a-1)&s)if(a<(s^a))
            for(int v=0;v<n;v++)dp[s][v]=min(dp[s][v],dp[a][v]+dp[s^a][v]);
        priority_queue<pair<ll,int>,vector<pair<ll,int>>,greater<pair<ll,int>>> q;
        for(int v=0;v<n;v++)if(dp[s][v]<INF)q.push({dp[s][v],v});
        while(!q.empty())
        {
            auto [d,u]=q.top();q.pop();if(d!=dp[s][u])continue;
            for(auto [v,w]:g[u])if(w<=INF-d&&d+w<dp[s][v])dp[s][v]=d+w,q.push({d+w,v});
        }
    }
    return dp.back()[terminal[0]];
}
\end{lstlisting}
% SOURCE: 04-图论/最小树形图(朱刘算法).cpp
\subsection{最小树形图(朱刘算法)}

\note{e为有向边，root为根，in/\allowbreak{}pre为每点最小入边及前驱，idd/\allowbreak{}vis用于找环缩点，ok独立表示是否存在解。点0起；算法原地改写e及n，负权允许，不用答案-\allowbreak{}1判断无解，再求解须恢复原图。}
\begin{lstlisting}
const int N=105;
const ll INF=1e15;
int n,m,root;// 模板内部用 0..n-1 编号，缩点后 n 会变小
int ok;// 1 表示最小树形图存在（边权可以为负，所以别拿 -1 当无解标志）
struct Edge
{
    int u,v;
    ll w;
};
Edge e[N*N];// 边表，每轮缩点后原地修改
ll in[N];// 每个点当前选中的最小入边权
int pre[N],idd[N],vis[N];
ll zhuliu()// 有向图最小树形图（朱刘/Edmonds），O(n*m)；结果存返回值，无解时 ok=0
{
    ll ans=0;
    ok=1;
    while(1)
    {
        for(int i=0;i<n;i++)in[i]=INF;
        for(int i=1;i<=m;i++)
            if(e[i].u!=e[i].v&&e[i].w<in[e[i].v])
                in[e[i].v]=e[i].w,pre[e[i].v]=e[i].u;// 自环不考虑
        for(int i=0;i<n;i++)
            if(i!=root&&in[i]==INF)
            {
                ok=0;// 有人没有入边 -> 根到不了它，无解
                return 0;
            }
        int cnt=0;
        for(int i=0;i<n;i++)idd[i]=-1,vis[i]=-1;
        in[root]=0;
        for(int i=0;i<n;i++)
        {
            ans+=in[i];
            int v=i;
            while(vis[v]!=i&&idd[v]==-1&&v!=root)
                vis[v]=i,v=pre[v];// 顺着入边往上爬，找环
            if(v!=root&&idd[v]==-1)// 找到一个没见过的新环，缩成一点
            {
                for(int u=pre[v];u!=v;u=pre[u])idd[u]=cnt;
                idd[v]=cnt++;
            }
        }
        if(!cnt)break;// 一个环都没有，已经是最小树形图
        for(int i=0;i<n;i++)if(idd[i]==-1)idd[i]=cnt++;// 不在环里的点各成一点
        for(int i=1;i<=m;i++)
        {
            int v=e[i].v;
            e[i].u=idd[e[i].u],e[i].v=idd[v];
            if(e[i].u!=e[i].v)e[i].w-=in[v];// 环内的边扣掉原本选的入边，抵消重复计算
        }
        n=cnt;
        root=idd[root];
    }
    return ans;
}
\end{lstlisting}
% SOURCE: 04-图论/最大团与极大团.cpp
\subsection{最大团与极大团}

\note{MaximumClique维护邻接关系、当前团和最佳团；最大团返回点集。maximal\_cliques的g是每点邻居的64位掩码，callback接收每个极大团；极大并非最大，点数和移位须符合位宽。}
\note{最大团：着色数是可扩展规模上界，分支界限搜索；最坏指数级。n<\allowbreak{}=\allowbreak{}200，0-\allowbreak{}indexed。 输入无自环对称 bitset 邻接表；一般图最大独立集等于补图最大团。 极大团枚举：Bron–Kerbosch +\allowbreak{} 枢轴，n<\allowbreak{}=\allowbreak{}63；极大是不再能扩展，未必规模最大。}
\begin{lstlisting}
struct MaximumClique
{
    vector<bitset<200>> g;vector<int> cur,best;
    MaximumClique(vector<bitset<200>> g):g(g){assert(g.size()<=200);}
    void search(vector<int> p)
    {
        vector<int> order,bound;int color=0;
        while(!p.empty())
        {
            ++color;vector<int> rest;bitset<200> used;
            for(int v:p)if((g[v]&used).none())used[v]=1,order.push_back(v),bound.push_back(color);else rest.push_back(v);
            p.swap(rest);
        }
        for(int i=(int)order.size()-1;i>=0;i--)
        {
            if(cur.size()+bound[i]<=best.size())return;
            int v=order[i];cur.push_back(v);vector<int> next;
            for(int j=0;j<i;j++)if(g[v][order[j]])next.push_back(order[j]);
            if(next.empty()){if(cur.size()>best.size())best=cur;}else search(next);
            cur.pop_back();
        }
    }
    vector<int> solve(){cur.clear();best.clear();vector<int> p(g.size());iota(p.begin(),p.end(),0);search(p);return best;}
};
\end{lstlisting}
\note{callback 接收一个团的顶点掩码；空图约定枚举空团。输出数量本身可能指数级。}
\begin{lstlisting}
template<class Callback>
void maximal_cliques(const vector<unsigned long long> &g,Callback callback)
{
    using U=unsigned long long;int n=g.size();assert(n<=63);
    auto dfs=[&](auto&&dfs,U r,U p,U x)->void
    {
        if(!p&&!x){callback(r);return;}
        U union_set=p|x;int pivot=-1,degree=-1;
        while(union_set){int u=__builtin_ctzll(union_set);union_set&=union_set-1;int d=__builtin_popcountll(p&g[u]);if(d>degree)degree=d,pivot=u;}
        U cand=p&(pivot<0?~U(0):~g[pivot]);
        while(cand){int v=__builtin_ctzll(cand);U bit=U(1)<<v;cand^=bit;dfs(dfs,r|bit,p&g[v],x&g[v]);p^=bit;x|=bit;}
    };
    dfs(dfs,0,(1ULL<<n)-1,0);
}
\end{lstlisting}
% SOURCE: 04-图论/网络流PushRelabel.cpp
\subsection{网络流PushRelabel}

\note{对象g保存残余边及反向下标，ec为超额流，H为高度；calc(\allowbreak{}s,\allowbreak{}t)\allowbreak{}消耗网络并返回最大流。leftOfMinCut读取求解后的高度侧别，重新求原网络须重建容量。}
\note{最高标号 Push–Relabel +\allowbreak{} gap，vector 存图，0-\allowbreak{}indexed；O(\allowbreak{}V\ensuremath{^2} sqrt(\allowbreak{}E)\allowbreak{})\allowbreak{} 标准界。 addEdge(\allowbreak{}u,\allowbreak{}v,\allowbreak{}c,\allowbreak{}rcap=\allowbreak{}0)\allowbreak{}，有向边反向初始 0，无向容量可 rcap=\allowbreak{}c。容量及总流须放入 ll。 calc(\allowbreak{}s,\allowbreak{}t)\allowbreak{} 仅调用一次，s!=\allowbreak{}t；边加完后再调用，运行期间 vector 不得扩容。 leftOfMinCut(\allowbreak{}a)\allowbreak{} 表示最后残量图源侧；整理自 KACTL PushRelabel.\allowbreak{}h，Simon Lindholm，CC0。}
\begin{lstlisting}
struct PushRelabel {
    struct Edge {
        int dest, back;
        ll f, c;
    };
    vector<vector<Edge>> g;
    vector<ll> ec;
    vector<Edge*> cur;
    vector<vector<int>> hs; vector<int> H;
    PushRelabel(int n) : g(n), ec(n), cur(n), hs(2*n), H(n) {}
    void addEdge(int s, int t, ll cap, ll rcap=0) {
        if (s == t) return;
        g[s].push_back({t, (int)g[t].size(), 0, cap});
        g[t].push_back({s, (int)g[s].size()-1, 0, rcap});
    }
    void addFlow(Edge& e, ll f) {
        Edge &back = g[e.dest][e.back];
        if (!ec[e.dest] && f) hs[H[e.dest]].push_back(e.dest);
        e.f += f; e.c -= f; ec[e.dest] += f;
        back.f -= f; back.c += f; ec[back.dest] -= f;
    }
    ll calc(int s, int t) {
        int v = (int)g.size(); H[s] = v; ec[t] = 1;
        vector<int> co(2*v); co[0] = v-1;
        for(int i=0;i<v;i++) cur[i] = g[i].data();
        for (Edge& e : g[s]) addFlow(e, e.c);
        for (int hi = 0;;) {
            while (hs[hi].empty()) if (!hi--) return -ec[s];
            int u = hs[hi].back(); hs[hi].pop_back();
            while (ec[u] > 0)  // discharge u
                if (cur[u] == g[u].data() + (int)g[u].size()) {
                    H[u] = 1e9;
                    for (Edge& e : g[u]) if (e.c && H[u] > H[e.dest]+1)
                        H[u] = H[e.dest]+1, cur[u] = &e;
                    if (++co[H[u]], !--co[hi] && hi < v)
                        for(int i=0;i<v;i++) if (hi < H[i] && H[i] < v)
                            --co[H[i]], H[i] = v + 1;
                    hi = H[u];
                } else if (cur[u]->c && H[u] == H[cur[u]->dest]+1)
                    addFlow(*cur[u], min(ec[u], cur[u]->c));
                else ++cur[u];
        }
    }
    bool leftOfMinCut(int a) { return H[a] >= (int)g.size(); }
};
\end{lstlisting}
% SOURCE: 04-图论/一般图匹配(带花树).cpp
\subsection{一般图匹配(带花树)}

\note{对象邻接表为无向图，match保存伴侣，未匹配哨兵按接口设置；base为花基，parent为交错树前驱。缩花只改变搜索视图，增广后更新真实匹配；换图重新初始化对象。}
\note{O(\allowbreak{}n\textasciicircum{}3)\allowbreak{}，Edmonds 缩花，顶点从 1 开始；允许奇环}
\begin{lstlisting}
struct Blossom
{
    int n;
    vector<vector<int> > adj;
    vector<int> match,fa,base,vis,flower;
    Blossom(int n):n(n),adj(n+1),match(n+1),fa(n+1),base(n+1),vis(n+1),flower(n+1){}
    void add_edge(int u,int v)
    {
        if(u!=v)adj[u].push_back(v),adj[v].push_back(u);
    }
    int lca(int u,int v)
    {
        vector<int> mark(n+1);
        while(true)
        {
            u=base[u],mark[u]=1;
            if(!match[u])break;
            u=fa[match[u]];
        }
        while(true)
        {
            v=base[v];
            if(mark[v])return v;
            v=fa[match[v]];
        }
    }
    void mark_path(int u,int b,int v)
    {
        while(base[u]!=b)
        {
            flower[base[u]]=flower[base[match[u]]]=1;
            fa[u]=v,v=match[u],u=fa[match[u]];
        }
    }
    int augment(int s)
    {
        fill(vis.begin(),vis.end(),0),fill(fa.begin(),fa.end(),0);
        iota(base.begin(),base.end(),0);
        vector<int> q(1,s);
        vis[s]=1;
        for(int i=0;i<(int)q.size();i++)
        {
            int u=q[i];
            for(int v:adj[u])
            {
                if(base[u]==base[v]||match[u]==v)continue;
                if(v==s||(match[v]&&fa[match[v]]))
                {
                    int b=lca(u,v);
                    fill(flower.begin(),flower.end(),0);
                    mark_path(u,b,v),mark_path(v,b,u);
                    for(int w=1;w<=n;w++)if(flower[base[w]])
                    {
                        base[w]=b;
                        if(!vis[w])vis[w]=1,q.push_back(w);
                    }
                }
                else if(!fa[v])
                {
                    fa[v]=u;
                    if(!match[v])
                    {
                        while(v)
                        {
                            int p=fa[v],w=match[p];
                            match[v]=p,match[p]=v,v=w;
                        }
                        return 1;
                    }
                    vis[match[v]]=1,q.push_back(match[v]);
                }
            }
        }
        return 0;
    }
    int run()
    {
        fill(match.begin(),match.end(),0);
        int ans=0;
        for(int i=1;i<=n;i++)if(!match[i])ans+=augment(i);
        return ans;
    }
};
\end{lstlisting}
% SOURCE: 04-图论/一般图带权匹配.cpp
\subsection{一般图带权匹配}

\note{Blossom的权矩阵存真实点及收缩花，match为伴侣，lab为对偶标号，flower记录花的内部顺序。容量需覆盖收缩节点，内部编号按源实现约定；重复边取更优权，负边及缺边的处理见入口前提。}
\note{一般图最大总权匹配，O(\allowbreak{}n\ensuremath{^3})\allowbreak{} 时间、O(\allowbreak{}n\ensuremath{^2})\allowbreak{} 空间；允许不匹配，并非先最大匹配数。 顶点 1.\allowbreak{}.\allowbreak{}n，match[\allowbreak{}u]\allowbreak{}=\allowbreak{}0 未匹配；只保留正权，非正权可舍弃。重边取最大值，自环忽略。 solve 返回 \{总权,\allowbreak{}匹配边数\}，对象一次构建后只 solve 一次；边权/\allowbreak{}双倍顶标远离 1e18。 整理自 ShahjalalShohag/\allowbreak{}code-\allowbreak{}library/\allowbreak{}Graph Theory/\allowbreak{}Blossom Algorithm Weighted.\allowbreak{}cpp；动态 vector 存储。 花收缩维护奇环、花展开恢复交错树；辅助函数承担算法不变量，不能任意合并删除。}
\note{Complexity: O(\allowbreak{}n\textasciicircum{}3)\allowbreak{} It finds maximum cost matching on a general graph,\allowbreak{} not necessarily with maximum matching 1-\allowbreak{}indexed}
\begin{lstlisting}
struct Blossom {
  const long long inf = 1e18;
  struct edge {
    int u, v;
    long long w;
  };
  vector<vector<edge>> g;
  long long reduced(const edge &e) const {return lab[e.u]+lab[e.v]-2*g[e.u][e.v].w;}
  int n,n_x;
  vector<int> match,slack,st,pa,S,vis;vector<vector<int>> flower_from;
  vector<long long> lab;
  vector<vector<int>> flower;
  deque<int> q;
  explicit Blossom(int _n):g(2*_n+1,vector<edge>(2*_n+1)),n(_n),n_x(_n),
      match(2*n+1),slack(2*n+1),st(2*n+1),pa(2*n+1),S(2*n+1),vis(2*n+1),
      flower_from(2*n+1,vector<int>(2*n+1)),lab(2*n+1),flower(2*n+1) {
    q = deque<int>();
    for (int u = 1; u <= n * 2; ++u) {
      match[u] = slack[u] = st[u] = pa[u] = S[u] = vis[u] = lab[u] = 0;
      for (int v = 1; v <= n * 2; ++v) {
        g[u][v] = edge{u, v, 0};
        flower_from[u][v] = 0;
      }
      flower[u].clear();
    }
  }
  void add_edge(int u, int v, long long w) {
    assert(1<=u&&u<=n&&1<=v&&v<=n);if(u==v)return;
    g[u][v].w = max(g[u][v].w, w);
    g[v][u].w = max(g[v][u].w, w);
  }
  inline void update_slack(int u, int x) {
    if (!slack[x] || reduced(g[u][x]) < reduced(g[slack[x]][x])) slack[x] = u;
  }
  inline void set_slack(int x) {
    slack[x] = 0;
    for (int u = 1; u <= n; ++u) {
      if (g[u][x].w > 0 && st[u] != x && S[st[u]] == 0) update_slack(u, x);
    }
  }
  inline void q_push(int x) {
    if (x <= n) return q.push_back(x);
    for (int i = 0; i < (int)flower[x].size(); i++) q_push(flower[x][i]);
  }
  inline void set_st(int x, int b) {
    st[x] = b;
    if (x <= n) return;
    for (int i = 0; i < (int)flower[x].size(); ++i) set_st(flower[x][i], b);
  }
  inline int get_pr(int b, int xr) {
    int pr = find(flower[b].begin(), flower[b].end(), xr) - flower[b].begin();
    if (pr % 2 == 1) {
      reverse(flower[b].begin() + 1, flower[b].end());
      return (int)flower[b].size() - pr;
    } else return pr;
  }
  inline void set_match(int u, int v) {
    match[u] = g[u][v].v;
    if (u <= n) return;
    edge e = g[u][v];
    int xr = flower_from[u][e.u], pr = get_pr(u, xr);
    for (int i = 0; i < pr; ++i) set_match(flower[u][i], flower[u][i ^ 1]);
    set_match(xr, v);
    rotate(flower[u].begin(), flower[u].begin() + pr, flower[u].end());
  }
  inline void augment(int u, int v) {
    int xnv = st[match[u]];
    set_match(u, v);
    if (!xnv) return;
    set_match(xnv, st[pa[xnv]]);
    augment(st[pa[xnv]], xnv);
  }
  inline int get_lca(int u, int v) {
    static int t = 0;
    for (++t; u || v; swap(u, v)) {
      if (u == 0) continue;
      if (vis[u] == t) return u;
      vis[u] = t;
      u = st[match[u]];
      if (u) u = st[pa[u]];
    }
    return 0;
  }
  inline void add_blossom(int u, int lca, int v) {
    int b = n + 1;
    while(b <= n_x && st[b]) ++b;
    if (b > n_x) ++n_x;
    lab[b] = 0, S[b] = 0;
    match[b] = match[lca];
    flower[b].clear();
    flower[b].push_back(lca);
    for (int x = u, y; x != lca; x = st[pa[y]]) {
      flower[b].push_back(x), flower[b].push_back(y = st[match[x]]), q_push(y);
    }
    reverse(flower[b].begin() + 1, flower[b].end());
    for (int x = v, y; x != lca; x = st[pa[y]]) {
      flower[b].push_back(x), flower[b].push_back(y = st[match[x]]), q_push(y);
    }
    set_st(b, b);
    for (int x = 1; x <= n_x; ++x) g[b][x].w = g[x][b].w = 0;
    for (int x = 1; x <= n; ++x) flower_from[b][x] = 0;
    for (int i = 0; i < (int)flower[b].size(); ++i) {
      int xs = flower[b][i];
      for (int x = 1; x <= n_x; ++x) {
        if (g[b][x].w == 0 || reduced(g[xs][x]) < reduced(g[b][x]))
          g[b][x] = g[xs][x], g[x][b] = g[x][xs];
      }
      for (int x = 1; x <= n; ++x) {
        if (flower_from[xs][x]) flower_from[b][x] = xs;
      }
    }
    set_slack(b);
  }
  inline void expand_blossom(int b) { // S[b] == 1
    for (int i = 0; i < (int)flower[b].size(); ++i) set_st(flower[b][i], flower[b][i]);
    int xr = flower_from[b][g[b][pa[b]].u], pr = get_pr(b, xr);
    for (int i = 0; i < pr; i += 2) {
      int xs = flower[b][i], xns = flower[b][i + 1];
      pa[xs] = g[xns][xs].u;
      S[xs] = 1, S[xns] = 0;
      slack[xs] = 0, set_slack(xns);
      q_push(xns);
    }
    S[xr] = 1, pa[xr] = pa[b];
    for (int i = pr + 1; i < (int)flower[b].size(); ++i) {
      int xs = flower[b][i];
      S[xs] = -1, set_slack(xs);
    }
    st[b] = 0;
  }
  inline bool on_found_edge(const edge &e) {
    int u = st[e.u], v = st[e.v];
    if (S[v] == -1) {
      pa[v] = e.u, S[v] = 1;
      int nu = st[match[v]];
      slack[v] = slack[nu] = 0;
      S[nu] = 0, q_push(nu);
    } else if (S[v] == 0) {
      int lca = get_lca(u, v);
      if (!lca) return augment(u, v), augment(v, u), 1;
      else add_blossom(u, lca, v);
    }
    return 0;
  }
  inline bool matching() {
    fill(S.begin(),S.end(),-1), fill(slack.begin(),slack.end(),0);
    q.clear();
    for (int x = 1; x <= n_x; ++x) {
      if (st[x] == x && !match[x]) pa[x] = 0, S[x] = 0, q_push(x);
    }
    if (q.empty()) return 0;
    for (;;) {
      while((int)q.size()) {
        int u = q.front();
        q.pop_front();
        if (S[st[u]] == 1)continue;
        for (int v = 1; v <= n; ++v) {
          if (g[u][v].w > 0 && st[u] != st[v]) {
            if (reduced(g[u][v]) == 0) {
              if (on_found_edge(g[u][v])) return 1;
            } else update_slack(u, st[v]);
          }
        }
      }
      long long d = inf;
      for (int b = n + 1; b <= n_x; ++b) {
        if (st[b] == b && S[b] == 1) d = min(d, lab[b] / 2);
      }
      for (int x = 1; x <= n_x; ++x) {
        if (st[x] == x && slack[x]) {
          if (S[x] == -1) d = min(d, reduced(g[slack[x]][x]));
          else if (S[x] == 0) d = min(d, reduced(g[slack[x]][x]) / 2);
        }
      }
      for (int u = 1; u <= n; ++u) {
        if (S[st[u]] == 0) {
          if (lab[u] <= d)return 0;
          lab[u] -= d;
        } else if (S[st[u]] == 1) lab[u] += d;
      }
      for (int b = n + 1; b <= n_x; ++b) {
        if (st[b] == b) {
          if (S[st[b]] == 0) lab[b] += d * 2;
          else if (S[st[b]] == 1) lab[b] -= d * 2;
        }
      }
      q.clear();
      for (int x = 1; x <= n_x; ++x) {
        if (st[x] == x && slack[x] && st[slack[x]] != x && reduced(g[slack[x]][x]) == 0)
          if (on_found_edge(g[slack[x]][x])) return 1;
      }
      for (int b = n + 1; b <= n_x; ++b) {
        if (st[b] == b && S[b] == 1 && lab[b] == 0) expand_blossom(b);
      }
    }
    return 0;
  }
  pair<long long, int> solve() {
    fill(match.begin(),match.end(),0);
    n_x = n;
    int cnt = 0;
    long long ans = 0;
    for (int u = 0; u <= n; ++u) st[u] = u, flower[u].clear();
    long long w_max = 0;
    for (int u = 1; u <= n; ++u) {
      for (int v = 1; v <= n; ++v) {
        flower_from[u][v] = (u == v ? u : 0);
        w_max = max(w_max, g[u][v].w);
      }
    }
    for (int u = 1; u <= n; ++u) lab[u] = w_max;
    while (matching()) ++cnt;
    for (int u = 1; u <= n; ++u) {
      if (match[u] && match[u] < u) ans += g[u][match[u]].w;
    }
    return make_pair(ans, cnt);
  }
};
\end{lstlisting}
% SOURCE: 04-图论/支配树.cpp
\subsection{支配树}

\note{g为0起有向邻接表，root是入口。返回每点的直接支配点，入口及不可达点为-\allowbreak{}1。dfn和ord互相映射原点与DFS编号，parent为DFS父亲，semi为半支配点，ancestor与label用于并查集求代表；每次重建全部状态。}
\note{Lengauer–Tarjan 支配树：固定入口 root 到 v 的所有路径都经过 idom[\allowbreak{}v]\allowbreak{}。 有向图 0-\allowbreak{}indexed，返回直接支配点；入口和不可达点为 -\allowbreak{}1。O(\allowbreak{}(\allowbreak{}n+\allowbreak{}m)\allowbreak{}log n)\allowbreak{} 标准界。 DFS 与并查集压缩均迭代，避免长链爆栈；割点算法不能替代入口相关的有向支配。}
\begin{lstlisting}
vector<vector<int>> g;
vector<int> dominator_tree(int root)
{
    int n=g.size();assert(0<=root&&root<n);vector<int> dfn(n),ord={-1},parent(n+1),idx(n),stack={root};
    dfn[root]=1;ord.push_back(root);
    while(!stack.empty()){int u=stack.back();if(idx[u]==(int)g[u].size()){stack.pop_back();continue;}int v=g[u][idx[u]++];if(!dfn[v]){dfn[v]=ord.size();parent[dfn[v]]=dfn[u];ord.push_back(v);stack.push_back(v);}}
    int count=ord.size()-1;vector<vector<int>> pred(count+1),bucket(count+1);
    for(int u=0;u<n;u++)if(dfn[u])for(int v:g[u])if(dfn[v])pred[dfn[v]].push_back(dfn[u]);
    vector<int> semi(count+1),label(count+1),ancestor(count+1),dom(count+1);iota(semi.begin(),semi.end(),0);label=semi;
    auto representative=[&](int v){vector<int> path;int x=v;while(ancestor[x]&&ancestor[ancestor[x]])path.push_back(x),x=ancestor[x];for(auto it=path.rbegin();it!=path.rend();it++){int u=*it,a=ancestor[u];if(semi[label[a]]<semi[label[u]])label[u]=label[a];ancestor[u]=ancestor[a];}return label[v];};
    for(int i=count;i>1;i--)
    {
        for(int v:pred[i])semi[i]=min(semi[i],semi[representative(v)]);
        bucket[semi[i]].push_back(i);ancestor[i]=parent[i];
        for(int v:bucket[parent[i]]){int u=representative(v);dom[v]=semi[u]<semi[v]?u:parent[i];}bucket[parent[i]].clear();
    }
    vector<int> answer(n,-1);
    for(int i=2;i<=count;i++){if(dom[i]!=semi[i])dom[i]=dom[dom[i]];answer[ord[i]]=ord[dom[i]];}
    return answer;
}
\end{lstlisting}
% SOURCE: 04-图论/稳定匹配.cpp
\subsection{稳定匹配}

\note{两侧各n人，编号0起；a[\allowbreak{}i]\allowbreak{}为左人i对右人的完整偏好顺序，b[\allowbreak{}j]\allowbreak{}为右人j对左人的偏好顺序，越前越喜欢。返回每个左人的伴侣；rank、提议进度、左右匹配和队列每次重建。}
\note{Gale–Shapley：两侧各 n 人，严格、完整偏好列表，编号 0.\allowbreak{}.\allowbreak{}n-\allowbreak{}1，越前越喜欢。 返回左侧的伴侣；O(\allowbreak{}n\ensuremath{^2})\allowbreak{}，提议的左侧在所有稳定匹配中最优。目标不是最大总权。}
\begin{lstlisting}
vector<vector<int>> a,b;
vector<int> stable_matching()
{
    int n=a.size();assert((int)b.size()==n);
    vector<vector<int>> rank(n,vector<int>(n));
    for(int j=0;j<n;j++)for(int k=0;k<n;k++)rank[j][b[j][k]]=k;
    vector<int> next(n),left(n,-1),right(n,-1);queue<int> q;
    for(int i=0;i<n;i++)q.push(i);
    while(!q.empty())
    {
        int i=q.front();q.pop();int j=a[i][next[i]++],old=right[j];
        if(old<0||rank[j][i]<rank[j][old])
        {left[i]=j;right[j]=i;if(old>=0)left[old]=-1,q.push(old);}
        else q.push(i);
    }
    return left;
}
\end{lstlisting}
% SOURCE: 04-图论/一般图边染色.cpp
\subsection{一般图边染色}

\note{N为点数，eds为0起无向简单图边表，返回同顺序的颜色编号；cc记录点度，颜色表维护每点缺失颜色。重边不能直接使用这份简单图算法，空边集按实现边界处理。}
\note{一般无向简单图边染色，O(\allowbreak{}nm)\allowbreak{}，最多 \ensuremath{\Delta}+\allowbreak{}1 色，颜色从 0 开始；不允许重边、自环。 参考 KACTL EdgeColoring.\allowbreak{}h，Simon Lindholm，CC0；Misra–Gries 扇形旋转与双色链交换。 二分图可做到 \ensuremath{\Delta} 色，本函数不保证；编号 0.\allowbreak{}.\allowbreak{}n-\allowbreak{}1，返回与输入边一一对应的颜色。 二分图（含重边）只需 \ensuremath{\Delta} 色；本一般简单图算法只保证 \ensuremath{\Delta}+\allowbreak{}1，不能替代要求 \ensuremath{\Delta} 色的构造。 二分图构造可沿两端缺色组成的双色交替链换色；判定/\allowbreak{}构造模型要分清点染色与边染色。}
\begin{lstlisting}
vector<int> edgeColoring(int N, vector<pair<int,int>> eds) {
    vector<int> cc(N + 1), ret((int)eds.size()), fan(N), free(N), loc;
    for (pair<int,int> e : eds) ++cc[e.first], ++cc[e.second];
    int u, v, ncols = *max_element(cc.begin(),cc.end()) + 1;
    vector<vector<int>> adj(N, vector<int>(ncols, -1));
    for (pair<int,int> e : eds) {
        tie(u, v) = e;
        fan[0] = v;
        loc.assign(ncols, 0);
        int at = u, end = u, d, c = free[u], ind = 0, i = 0;
        while (d = free[v], !loc[d] && (v = adj[u][d]) != -1)
            loc[d] = ++ind, cc[ind] = d, fan[ind] = v;
        cc[loc[d]] = c;
        for (int cd = d; at != -1; cd ^= c ^ d, at = adj[at][cd])
            swap(adj[at][cd], adj[end = at][cd ^ c ^ d]);
        while (adj[fan[i]][d] != -1) {
            int left = fan[i], right = fan[++i], e = cc[i];
            adj[u][e] = left;
            adj[left][e] = u;
            adj[right][e] = -1;
            free[right] = e;
        }
        adj[u][d] = fan[i];
        adj[fan[i]][d] = u;
        for (int y : {fan[0], u, end})
            for (int& z = free[y] = 0; adj[y][z] != -1; z++);
    }
    for(int i=0;i<(int)eds.size();i++)
        for (tie(u, v) = eds[i]; adj[u][ret[i]] != v;) ++ret[i];
    return ret;
}
\end{lstlisting}
\section{数学}
% SOURCE: 05-数学/快速幂.cpp
\subsection{快速幂}

\note{a为底数，n为非负指数，mod为模数；返回a\textasciicircum{}n的对应形式。乘法单位元为1，模版需归一负底数并防模乘溢出；负指数必须另求逆，不能直接右移循环。}
\note{适用：模指数、模乘；指数按二进制拆开，幂底每轮平方。 参数：n>\allowbreak{}=\allowbreak{}0 是指数，mod>\allowbreak{}0；qmul 的 b 是乘数，qpow 的 n 才是指数。 乘法用 \_\_int128；倍增模乘用比较减法，mod 接近 LLONG\_MAX 也不溢出。 底数、乘数可负，返回 [\allowbreak{}0,\allowbreak{}mod)\allowbreak{}；零次幂是 1\%mod。}
\note{结论：大模数优先选 qpow128/\allowbreak{}qmul128；非负指数范围内时间 O(\allowbreak{}log n)\allowbreak{}，辅助空间 O(\allowbreak{}1)\allowbreak{}。}
\begin{lstlisting}
typedef unsigned long long ull;
typedef __int128 lll;
const int N=100005;
\end{lstlisting}
\note{O(\allowbreak{}log n)\allowbreak{}，快速幂 (\allowbreak{}a\textasciicircum{}n)\allowbreak{} \% mod}
\begin{lstlisting}
ll qpow(ll a,ll n,ll mod)
{
    ll res=1%mod;
    a%=mod;
    if(a<0)a+=mod;
    while(n)
    {
        if(n&1)res=(__int128)res*a%mod;
        a=(__int128)a*a%mod;
        n>>=1;
    }
    return res;
}
\end{lstlisting}
\note{O(\allowbreak{}log n)\allowbreak{}，快速乘 (\allowbreak{}a*b)\allowbreak{} \% mod，防 long long 溢出（mod <\allowbreak{} 2\textasciicircum{}63）}
\begin{lstlisting}
ll qmul(ll a,ll b,ll mod)
{
    ll res=0;
    a%=mod;
    if(a<0)a+=mod;
    b%=mod;
    if(b<0)b+=mod;
    while(b)
    {
        if(b&1)res=res>=mod-a?res-(mod-a):res+a;
        a=a>=mod-a?a-(mod-a):a+a;
        b>>=1;
    }
    return res;
}
ll qmul128(ll a,ll b,ll mod)
{
    lll t=(lll)a*b;
    t%=mod; return (ll)(t<0?t+mod:t);
}
\end{lstlisting}
\note{O(\allowbreak{}log n)\allowbreak{}，快速幂 +\allowbreak{} 快速乘版，要求 mod*mod 会爆 long long 时用 O(\allowbreak{}log n*log mod)\allowbreak{}，仅用无溢出的倍增模乘。}
\begin{lstlisting}
ll qpow_safe(ll a,ll n,ll mod)
{
    ll res=1%mod;
    a%=mod;
    if(a<0)a+=mod;
    while(n)
    {
        if(n&1)res=qmul(res,a,mod);
        a=qmul(a,a,mod);
        n>>=1;
    }
    return res;
}
\end{lstlisting}
\note{O(\allowbreak{}log n)\allowbreak{}，用 \_\_int128 做乘法再取模，跑得比 qmul 快，与 qpow\_safe 互为对照 O(\allowbreak{}log n)\allowbreak{}，使用宽整数模乘并归一底数；n=\allowbreak{}0 正确返回 1\%mod。}
\begin{lstlisting}
ll qpow128(ll a,ll n,ll mod)
{
    ll res=1%mod;
    a%=mod;
    if(a<0)a+=mod;
    while(n)
    {
        if(n&1)res=qmul128(res,a,mod);
        a=qmul128(a,a,mod);
        n>>=1;
    }
    return res;
}
\end{lstlisting}
% SOURCE: 05-数学/乘法逆元.cpp
\subsection{乘法逆元}

\note{n为预处理上界，p为模数，inv是1.\allowbreak{}.\allowbreak{}n的逆元，fact/\allowbreak{}inv\_fact是0.\allowbreak{}.\allowbreak{}n的阶乘及逆阶乘。线性递推和费马法要求素数且n<\allowbreak{}p；换模数必须重新建表，0没有乘法逆元。}
\note{适用：模意义下除法、组合数；除以 a 等价于乘逆元，但逆元存在需 gcd(\allowbreak{}a,\allowbreak{}p)\allowbreak{}=\allowbreak{}1。 参数：模数 p>\allowbreak{}1；费马及线性递推版要求素数 p，指数非负，输入余数先归一。 下标：inv 为 1.\allowbreak{}.\allowbreak{}n，fact/\allowbreak{}inv\_fact 为 0.\allowbreak{}.\allowbreak{}n，0<\allowbreak{}=\allowbreak{}n<\allowbreak{}N 且 n<\allowbreak{}p。 关键：阶乘逆元从最高下标逆推，只需一次快速幂，查询 O(\allowbreak{}1)\allowbreak{}。 易错：inv\_exgcd 未检查 gcd；模数变化后阶乘表必须重建；ll 乘积需保证不溢出。 复杂度：单个逆元 O(\allowbreak{}log p)\allowbreak{}，批量 O(\allowbreak{}n)\allowbreak{}，阶乘预处理 O(\allowbreak{}n+\allowbreak{}log p)\allowbreak{}，空间 O(\allowbreak{}n)\allowbreak{}。}
\begin{lstlisting}
const int N=3000006;
ll n,p;
ll inv[N],fact[N],inv_fact[N];
\end{lstlisting}
\note{O(\allowbreak{}log p)\allowbreak{}，费马小定理求逆元，要求 p 为素数且 gcd(\allowbreak{}a,\allowbreak{}p)\allowbreak{}=\allowbreak{}1 O(\allowbreak{}log n)\allowbreak{}，计算 a\textasciicircum{}n mod mod；函数本身是幂，求逆需指数 p-\allowbreak{}2。}
\begin{lstlisting}
ll qpow(ll a,ll n,ll mod)
{
    ll res=1;
    a%=mod;
    while(n)
    {
        if(n&1)(res*=a)%=mod;
        (a*=a)%=mod;
        n>>=1;
    }
    return res;
}
\end{lstlisting}
\note{O(\allowbreak{}log p)\allowbreak{}，返回 a 的逆元；p 为素数且 a 不是 p 的倍数。}
\begin{lstlisting}
ll inv_fermat(ll a,ll p)
{
    return qpow(a,p-2,p);
}
\end{lstlisting}
\note{O(\allowbreak{}log p)\allowbreak{}，exgcd 求逆元，只需 gcd(\allowbreak{}a,\allowbreak{}p)\allowbreak{}=\allowbreak{}1 O(\allowbreak{}log min(\allowbreak{}a,\allowbreak{}b)\allowbreak{})\allowbreak{}，返回 gcd，引用 x/\allowbreak{}y 输出 ax+\allowbreak{}by=\allowbreak{}gcd。}
\begin{lstlisting}
ll exgcd(ll a,ll b,ll &x,ll &y)
{
    if(!b){x=1,y=0;return a;}
    ll xx,yy;
    ll g=exgcd(b,a%b,xx,yy);
    x=yy,y=xx-(a/b)*yy;
    return g;
}
\end{lstlisting}
\note{O(\allowbreak{}log p)\allowbreak{}，返回最小非负逆元；a 建议传非负余数，互素前提由调用者保证。}
\begin{lstlisting}
ll inv_exgcd(ll a,ll p)
{
    ll x,y;
    exgcd(a,p,x,y);
    return (x%p<0?x%p+p:x%p);
}
\end{lstlisting}
\note{O(\allowbreak{}n)\allowbreak{}，线性递推求 1.\allowbreak{}.\allowbreak{}n 的全部逆元（p 为素数且 n<\allowbreak{}p） inv[\allowbreak{}i]\allowbreak{} =\allowbreak{} -\allowbreak{}(\allowbreak{}p/\allowbreak{}i)\allowbreak{} * inv[\allowbreak{}p\%i]\allowbreak{} (\allowbreak{}mod p)\allowbreak{} O(\allowbreak{}n)\allowbreak{}，求 inv[\allowbreak{}1.\allowbreak{}.\allowbreak{}n]\allowbreak{}；至少 n>\allowbreak{}=\allowbreak{}1，递推访问 p\%i<\allowbreak{}i 的已求逆元。}
\begin{lstlisting}
void inv_init(int n,ll p)
{
    inv[1]=1;
    for(int i=2;i<=n;i++)
        inv[i]=((p-p/i*inv[p%i])%p+p)%p;
}
\end{lstlisting}
\note{O(\allowbreak{}n)\allowbreak{}，阶乘与阶乘逆元预处理，配合组合数用 O(\allowbreak{}n+\allowbreak{}log p)\allowbreak{}，n 是预处理上界，p 为素数；fact[\allowbreak{}n]\allowbreak{} 不得为 0。}
\begin{lstlisting}
void fact_init(int n,ll p)
{
    fact[0]=1;
    for(int i=1;i<=n;i++)fact[i]=fact[i-1]*i%p;
    inv_fact[n]=qpow(fact[n],p-2,p);
    for(int i=n;i>=1;i--)inv_fact[i-1]=inv_fact[i]*i%p;
}
\end{lstlisting}
\note{O(\allowbreak{}1)\allowbreak{}，要求 n<\allowbreak{}p 且 p 为素数 O(\allowbreak{}1)\allowbreak{}，返回 C(\allowbreak{}n,\allowbreak{}m)\allowbreak{} mod p；n 须在同模数预处理范围内，非法 m 返回 0。}
\begin{lstlisting}
ll C_small(ll n,ll m,ll p)
{
    if(m<0||m>n||n<0)return 0;
    return fact[n]*inv_fact[m]%p*inv_fact[n-m]%p;
}
\end{lstlisting}
% SOURCE: 05-数学/扩展欧几里得.cpp
\subsection{扩展欧几里得}

\note{返回gcd，x/\allowbreak{}y通过引用输出ax+\allowbreak{}by=\allowbreak{}gcd的系数；这些标量输出保留参数。线性丢番图需常数项整除gcd，再沿通解参数取范围；负数余数归一及零系数边界按接口处理。}
\note{适用：线性丢番图方程、逆元、同余方程；有解条件是 gcd 整除常数项。 参数：a/\allowbreak{}b 为整数系数，x/\allowbreak{}y 是引用输出；solve\_cong 的 m 必须为正。 关键：由 b*x1+\allowbreak{}(\allowbreak{}a\%b)\allowbreak{}*y1=\allowbreak{}g 回代成 a*y1+\allowbreak{}b*(\allowbreak{}x1-\allowbreak{}(\allowbreak{}a/\allowbreak{}b)\allowbreak{}*y1)\allowbreak{}=\allowbreak{}g。 结论：ax=\allowbreak{}b (\allowbreak{}mod m)\allowbreak{} 的解间隔是 m/\allowbreak{}g；返回一个最小非负解，不是全部解。 exgcd 返回非负 gcd；系数中间值须在 ll 范围，不能用 (\allowbreak{}0,\allowbreak{}0)\allowbreak{} 的 gcd 作除数。 复杂度：欧几里得 O(\allowbreak{}log min(\allowbreak{} a ,\allowbreak{} b )\allowbreak{})\allowbreak{}，递归空间同量级；逆元函数不会检查互素。}
\begin{lstlisting}
const int N=100005;
\end{lstlisting}
\note{O(\allowbreak{}log min(\allowbreak{}|a|,\allowbreak{}|b|)\allowbreak{})\allowbreak{}，求 ax+\allowbreak{}by=\allowbreak{}gcd(\allowbreak{}a,\allowbreak{}b)\allowbreak{} 的一组解，返回 gcd |a|、|b|<\allowbreak{}=\allowbreak{}LLONG\_MAX；返回非负 gcd，含负系数时仍保持等式。 O(\allowbreak{}log min(\allowbreak{}|a|,\allowbreak{}|b|)\allowbreak{})\allowbreak{}，输出系数 x/\allowbreak{}y；a=\allowbreak{}b=\allowbreak{}0 时返回 0，不能随后拿来做除数。}
\begin{lstlisting}
ll exgcd(ll a,ll b,ll &x,ll &y)
{
    if(!b)
    {
        x=a<0?-1:1,y=0;
        return a<0?-a:a;
    }
    ll xx,yy;
    ll g=exgcd(b,a%b,xx,yy);
    x=yy;
    y=xx-(a/b)*yy;
    return g;
}
\end{lstlisting}
\note{同余方程 ax =\allowbreak{}=\allowbreak{}=\allowbreak{} b (\allowbreak{}mod m)\allowbreak{}，返回最小非负解，无解返回 -\allowbreak{}1 解集：x =\allowbreak{} x0 +\allowbreak{} k*(\allowbreak{}m/\allowbreak{}g)\allowbreak{} O(\allowbreak{}log m)\allowbreak{}，a 是系数、b 是目标余数，m 是正模数；无解返回 -\allowbreak{}1。}
\begin{lstlisting}
ll solve_cong(ll a,ll b,ll m)
{
    a%=m;if(a<0)a+=m;
    b%=m;if(b<0)b+=m;
    ll x,y;
    ll g=exgcd(a,m,x,y);
    if(b%g)return -1;// 裴蜀定理：g|b 才有解
    ll t=m/g;
    x=(ll)((__int128)x*(b/g)%t);
    x=(x%t<0?x%t+t:x%t);
    return x;
}
\end{lstlisting}
\note{O(\allowbreak{}log p)\allowbreak{}，模 p 下的逆元，要求 gcd(\allowbreak{}a,\allowbreak{}p)\allowbreak{}=\allowbreak{}1，返回最小非负解 O(\allowbreak{}log p)\allowbreak{}，gcd(\allowbreak{}a,\allowbreak{}p)\allowbreak{}=\allowbreak{}1 且 a>\allowbreak{}=\allowbreak{}0、p>\allowbreak{}1 时返回逆元代表。}
\begin{lstlisting}
ll inv_exgcd(ll a,ll p)
{
    ll x,y;
    exgcd(a,p,x,y);
    return (x%p<0?x%p+p:x%p);
}
\end{lstlisting}
\note{裴蜀定理：a,\allowbreak{}b 能线性组合出的最小正整数为 gcd(\allowbreak{}a,\allowbreak{}b)\allowbreak{} O(\allowbreak{}log min(\allowbreak{}|a|,\allowbreak{}|b|)\allowbreak{})\allowbreak{}，当前直接返回 exgcd 的值；负输入不能保证是最小正整数。}
\begin{lstlisting}
ll bezout_min(ll a,ll b)
{
    ll x,y;
    return exgcd(a,b,x,y);
}
\end{lstlisting}
% SOURCE: 05-数学/质因数分解与约数.cpp
\subsection{质因数分解与约数}

\note{n为正整数，因子表存\{素数,\allowbreak{}指数\}，约数由各素数幂组合生成。结果是否排序按入口确定，1的因子表为空但约数有1；试除范围必须覆盖剩余合数的可能小因子。}
\note{适用：单个正整数的因子、约数、phi；大批查询可换线性筛，大整数分解可换 Rho。 参数：所有函数的 n>\allowbreak{}=\allowbreak{}1；返回向量下标 0 起，factor 的第二项为因子指数。 关键：试除到剩余 n 的平方根；循环结束剩余数若大于 1 必为素数。 结论：约数个数为指数加一的乘积；phi(\allowbreak{}n)\allowbreak{}=\allowbreak{}n*各不同素因子的 (\allowbreak{}1-\allowbreak{}1/\allowbreak{}p)\allowbreak{} 乘积。 试除用 i<\allowbreak{}=\allowbreak{}n/\allowbreak{}i 避免平方溢出；div\_sum 的最终答案仍须能存入 ll。 复杂度：试除与枚举均 O(\allowbreak{}sqrt n)\allowbreak{} 上界，约数向量空间 O(\allowbreak{}约数个数)\allowbreak{}。}
\note{O(\allowbreak{}sqrt n)\allowbreak{}，试除法分解质因数，返回 (\allowbreak{}素因子,\allowbreak{}指数)\allowbreak{} 对，按因子升序 O(\allowbreak{}sqrt n)\allowbreak{}，返回升序 (\allowbreak{}素因子,\allowbreak{}重数)\allowbreak{}；输入 n=\allowbreak{}1 返回空因子表。}
\begin{lstlisting}
vector<pair<ll,int>> factor(ll n)
{
    vector<pair<ll,int>> f;
    for(ll i=2;i<=n/i;i++)
    {
        if(n%i)continue;
        int c=0;
        while(n%i==0)n/=i,c++;
        f.push_back(make_pair(i,c));// 注意这里保序
    }
    if(n>1)f.push_back(make_pair(n,1));
    return f;
}
\end{lstlisting}
\note{O(\allowbreak{}sqrt n)\allowbreak{}，枚举所有正约数，输出无序 O(\allowbreak{}sqrt n)\allowbreak{}，返回无序正约数；i=\allowbreak{}=\allowbreak{}n/\allowbreak{}i 仅加入一次，自行排序以获得升序。}
\begin{lstlisting}
vector<ll> get_div(ll n)
{
    vector<ll> v;
    for(ll i=1;i<=n/i;i++)
        if(n%i==0)
        {
            v.push_back(i);
            if(i!=n/i)v.push_back(n/i);
        }
    return v;
}
\end{lstlisting}
\note{O(\allowbreak{}sqrt n)\allowbreak{}，约数个数（配合 factor 可做到 O(\allowbreak{}因子个数)\allowbreak{}） O(\allowbreak{}sqrt n)\allowbreak{}，返回正约数个数；完全平方数的中间约数不重复计。}
\begin{lstlisting}
int div_cnt(ll n)
{
    int r=0;
    for(ll i=1;i<=n/i;i++)
        if(n%i==0)r+=(i==n/i?1:2);
    return r;
}
\end{lstlisting}
\note{O(\allowbreak{}sqrt n)\allowbreak{}，约数之和，防溢出用 \_\_int128 累加 O(\allowbreak{}sqrt n)\allowbreak{}，对互补约数求和；\_\_int128 中间值最后转回 ll。}
\begin{lstlisting}
ll div_sum(ll n)
{
    __int128 s=0;
    for(ll i=1;i<=n/i;i++)
        if(n%i==0)
        {
            s+=i;
            if(i!=n/i)s+=n/i;
        }
    return (ll)s;
}
\end{lstlisting}
\note{O(\allowbreak{}sqrt n)\allowbreak{}，单个数的欧拉函数 O(\allowbreak{}sqrt n)\allowbreak{}，每个不同素因子只应用一次 r=\allowbreak{}r/\allowbreak{}p*(\allowbreak{}p-\allowbreak{}1)\allowbreak{}，phi(\allowbreak{}1)\allowbreak{}=\allowbreak{}1。}
\begin{lstlisting}
ll phi_one(ll n)
{
    ll r=n;
    for(ll i=2;i<=n/i;i++)
        if(n%i==0)
        {
            r=r/i*(i-1);
            while(n%i==0)n/=i;
        }
    if(n>1)r=r/n*(n-1);
    return r;
}
\end{lstlisting}
% SOURCE: 05-数学/线性筛与欧拉函数.cpp
\subsection{线性筛与欧拉函数}

\note{n为筛上界，prime[\allowbreak{}1.\allowbreak{}.\allowbreak{}cnt]\allowbreak{}为素数，vis为合数标记；phi/\allowbreak{}mu为欧拉函数/\allowbreak{}莫比乌斯函数，d为约数数目，dmin为最小质因子指数。1的函数值按实现初始化，多次筛选前重置计数及所用表。}
\note{适用：批量求 phi/\allowbreak{}mu/\allowbreak{}约数个数，数据上界能放入数组时一次 O(\allowbreak{}n)\allowbreak{} 预处理。 下标：数值范围 1.\allowbreak{}.\allowbreak{}n，1<\allowbreak{}=\allowbreak{}n<\allowbreak{}N；prime[\allowbreak{}1.\allowbreak{}.\allowbreak{}cnt]\allowbreak{} 为升序素数，1 不是素数。 关键：合数只由最小质因子筛一次，遇 i\%prime[\allowbreak{}j]\allowbreak{}=\allowbreak{}=\allowbreak{}0 后必须 break。 状态：dmin[\allowbreak{}i]\allowbreak{} 为最小质因子指数；加同一因子时约数个数仅替换该指数贡献。 易错：重复 get\_prime 前清 vis、prime 与相关表，单改 cnt 不足以重新筛。 结论：phi(\allowbreak{}1)\allowbreak{}=\allowbreak{}mu(\allowbreak{}1)\allowbreak{}=\allowbreak{}d(\allowbreak{}1)\allowbreak{}=\allowbreak{}1；mu=\allowbreak{}0 表示有平方因子，试除参考版是 O(\allowbreak{}sqrt x)\allowbreak{}。}
\begin{lstlisting}
const int N=1000006;
int n;
int prime[N],cnt;// prime[1..cnt] 存的素数
bool vis[N];// 合数标记
int phi[N];// 欧拉函数
int mu[N];// 莫比乌斯函数
int d[N];// 约数个数
int dmin[N];// 最小质因子在该数中的指数
\end{lstlisting}
\note{O(\allowbreak{}n)\allowbreak{}，线性筛同时求 phi /\allowbreak{} mu /\allowbreak{} d O(\allowbreak{}n)\allowbreak{}，n 为最大查询值；同因子与新因子两种递推不能混用。}
\begin{lstlisting}
void get_prime(int n)
{
    cnt=0;
    fill(vis,vis+n+1,false);
    vis[1]=true;
    phi[1]=1,mu[1]=1,d[1]=1,dmin[1]=0;
    for(int i=2;i<=n;i++)
    {
        if(!vis[i])
        {
            prime[++cnt]=i;
            phi[i]=i-1;// 素数的 phi
            mu[i]=-1;// 素数有 1 个质因子
            d[i]=2;// 1 和自身
            dmin[i]=1;
        }
        for(int j=1;j<=cnt&&(ll)i*prime[j]<=n;j++)
        {
            int t=i*prime[j];
            vis[t]=true;
            if(i%prime[j]==0)
            {
                // prime[j] 已在 i 里出现过
                phi[t]=phi[i]*prime[j];
                mu[t]=0;
                dmin[t]=dmin[i]+1;
                d[t]=d[i]/(dmin[i]+1)*(dmin[i]+2);
                break;// 每个合数只被最小质因子筛一次
            }
            phi[t]=phi[i]*(prime[j]-1);
            mu[t]=-mu[i];
            dmin[t]=1;
            d[t]=d[i]*2;
        }
    }
}
\end{lstlisting}
% SOURCE: 05-数学/组合数与Lucas定理.cpp
\subsection{组合数与Lucas定理}

\note{fact/\allowbreak{}inv\_fact按当前素数模建立，普通组合查询下标小于模数；Lucas把n、m逐位拆成该素数进制后相乘。模数改变须重建表，大素数不能默认枚举整个模数范围。}
\note{适用：素数模数的组合计数；n 很大而 p 较小时用 Lucas 按 p 进制拆位。 参数：C 的 n 为总数、m 为选取数；阶乘表下标 0.\allowbreak{}.\allowbreak{}预处理界，查询必须同模数。 前提：p 为素数，预处理界<\allowbreak{}N 且 <\allowbreak{}p；Lucas 需预处理到 p-\allowbreak{}1，因此 p 不能超过容量。 关键：阶乘逆元从 n 逆推；Lucas 每一位组合数为 0 则整体为 0。 易错：不能把固定 mod=\allowbreak{}1e9+\allowbreak{}7 直接拿来预处理 p-\allowbreak{}1；普通 ll 模乘要求 p\textasciicircum{}2 可存。 复杂度：阶乘 O(\allowbreak{}预处理界+\allowbreak{}log p)\allowbreak{}，C 为 O(\allowbreak{}1)\allowbreak{}，Lucas O(\allowbreak{}log\_p n)\allowbreak{}；pf 每层还要枚举余数。}
\begin{lstlisting}
const int N=200005;
const ll mod=1000000007;
ll fact[N],inv_fact[N];
\end{lstlisting}
\note{O(\allowbreak{}log n)\allowbreak{}，快速幂 O(\allowbreak{}log n)\allowbreak{}，a 是底数、n 是非负指数、mod 是正模数。}
\begin{lstlisting}
ll qpow(ll a,ll n,ll mod)
{
    ll res=1;
    a%=mod;
    while(n)
    {
        if(n&1)(res*=a)%=mod;
        (a*=a)%=mod;
        n>>=1;
    }
    return res;
}
\end{lstlisting}
\note{O(\allowbreak{}N)\allowbreak{}，阶乘与阶乘逆元预处理 注意：必须保证 N <\allowbreak{} p，否则 fact[\allowbreak{}N]\allowbreak{} =\allowbreak{}=\allowbreak{}=\allowbreak{} 0 (\allowbreak{}mod p)\allowbreak{} 会让逆元链整体失效 O(\allowbreak{}n+\allowbreak{}log p)\allowbreak{}，预处理 0.\allowbreak{}.\allowbreak{}n 的 fact/\allowbreak{}inv\_fact；p 是素数且 n<\allowbreak{}p。}
\begin{lstlisting}
void C_init(int n,ll p)
{
    fact[0]=1;
    for(int i=1;i<=n;i++)fact[i]=fact[i-1]*i%p;
    inv_fact[n]=qpow(fact[n],p-2,p);
    for(int i=n;i>=1;i--)inv_fact[i-1]=inv_fact[i]*i%p;
}
\end{lstlisting}
\note{O(\allowbreak{}1)\allowbreak{}，C(\allowbreak{}n,\allowbreak{}m)\allowbreak{} mod p，要求 0<\allowbreak{}=\allowbreak{}n<\allowbreak{}p 且 p 为素数 O(\allowbreak{}1)\allowbreak{}，求 C(\allowbreak{}n,\allowbreak{}m)\allowbreak{} mod p；n 在预处理范围，m 不合法返回 0。}
\begin{lstlisting}
ll C(ll n,ll m,ll p)
{
    if(m<0||m>n||n<0)return 0;
    return fact[n]*inv_fact[m]%p*inv_fact[n-m]%p;
}
const int C_MAX=200000;// 阶乘预处理的数组上限；调用时传的 n 要同时满足 n < p
\end{lstlisting}
\note{O(\allowbreak{}log\_p n)\allowbreak{}，Lucas 定理：C(\allowbreak{}n,\allowbreak{}m)\allowbreak{} mod p =\allowbreak{} C(\allowbreak{}n\%p,\allowbreak{}m\%p)\allowbreak{}*C(\allowbreak{}n/\allowbreak{}p,\allowbreak{}m/\allowbreak{}p)\allowbreak{}，p 为素数 要求已 C\_init(\allowbreak{}p-\allowbreak{}1,\allowbreak{}p)\allowbreak{}（即预处理到 p-\allowbreak{}1）；p 很大且 n<\allowbreak{}p 时直接算 C(\allowbreak{}n,\allowbreak{}m)\allowbreak{} O(\allowbreak{}log\_p n)\allowbreak{}，n/\allowbreak{}m 非负，p 为小素数；取 n\%p/\allowbreak{}m\%p 为本位组合数再递归。}
\begin{lstlisting}
ll Lucas(ll n,ll m,ll p)
{
    if(m<0||m>n)return 0;
    if(m==0)return 1;
    return C(n%p,m%p,p)*Lucas(n/p,m/p,p)%p;
}
\end{lstlisting}
% SOURCE: 05-数学/容斥原理与排列组合.cpp
\subsection{容斥原理与排列组合}

\textbf{范德蒙德卷积}
\[
\sum_j\binom rj\binom s{k-j}=\binom{r+s}k.
\]
\note{对 $a_1\le\cdots\le a_n$，枚举所有子集，$a_i$ 作为子集第 $j$ 小元素的次数为
$\binom{i-1}{j-1}2^{n-i}$。若该位置权重依次为 $1,1,2,4,\ldots$，总系数为}
\[
\frac{3^{i-1}+1}{2}\,2^{n-i}.
\]
\note{两个不同指定对象分别禁放两个指定端点，$k\ge2$：
$k!-2(k-1)!+(k-2)!$。模意义除法须可逆。}

\note{fac为阶乘表，inv按本文件预处理保存逆阶乘；容斥项使用ll计数或按固定模数计算。组合对象的计数向量表示各类重复数量；空集项是否包含由求和范围决定，不能遗漏零项。}
\begin{lstlisting}
const ll MOD=1000000007;
const int N=100005;
ll fac[N],inv[N];
ll qpow(ll a,ll n) // n>=0，0<=a<MOD
{
    ll r=1;
    while(n)
    {
        if(n&1)r=r*a%MOD;
        a=a*a%MOD,n>>=1;
    }
    return r;
}
\end{lstlisting}
\note{n<\allowbreak{}MOD，O(\allowbreak{}n)\allowbreak{} 预处理，O(\allowbreak{}1)\allowbreak{} 组合数}
\begin{lstlisting}
void init(int n)
{
    assert(n>=0&&n<N);
    fac[0]=1;
    for(int i=1;i<=n;i++)fac[i]=fac[i-1]*i%MOD;
    inv[n]=qpow(fac[n],MOD-2);
    for(int i=n;i>=1;i--)inv[i-1]=inv[i]*i%MOD;
}
ll C(int n,int m) // n 在预处理范围内；m 不合法时返回 0
{
    if(m<0||m>n)return 0;
    return fac[n]*inv[m]%MOD*inv[n-m]%MOD;
}
\end{lstlisting}
\note{O(\allowbreak{}n)\allowbreak{}，容斥排除固定点；n=\allowbreak{}0 时空排列有一种}
\begin{lstlisting}
ll derange(int n)
{
    ll ans=0;
    for(int i=0;i<=n;i++)
    {
        ll x=C(n,i)*fac[n-i]%MOD;
        ans=(ans+(i&1?MOD-x:x))%MOD;
    }
    return (ll)ans;
}
ll multiset_count(const vector<int> &cnt) // cnt 非负，总和在预处理范围内
{
    int n=accumulate(cnt.begin(),cnt.end(),0);
    ll ans=fac[n];
    for(int x:cnt)ans=ans*inv[x]%MOD;
    return (ll)ans;
}
\end{lstlisting}
\note{O(\allowbreak{}m*2\textasciicircum{}m)\allowbreak{}，[\allowbreak{}1,\allowbreak{}n]\allowbreak{} 中不被任一 d 整除的数；d 必须为正}
\begin{lstlisting}
ll avoid(ll n,const vector<ll> &d)
{
    int m=d.size();
    assert(n>=0&&m<25);
    __int128 ans=n;
    for(int s=1;s<(1<<m);s++)
    {
        ll l=1;
        bool over=false;
        for(int i=0;i<m;i++)if(s>>i&1)
        {
            assert(d[i]>0);
            ll x=d[i]/gcd(l,d[i]);
            if(l>n/x){over=true;break;}
            l*=x;
        }
        if(over)continue; // 最小公倍数已超过 n，该交集贡献为 0；不构造 n+1
        ll x=n/l;
        if(__builtin_popcount((unsigned)s)&1)ans-=x;
        else ans+=x;
    }
    return (ll)ans;
}
\end{lstlisting}
\note{范德蒙德：sum(\allowbreak{}k)\allowbreak{} C(\allowbreak{}a,\allowbreak{}k)\allowbreak{}C(\allowbreak{}b,\allowbreak{}r-\allowbreak{}k)\allowbreak{}=\allowbreak{}C(\allowbreak{}a+\allowbreak{}b,\allowbreak{}r)\allowbreak{}，越界项为 0。 例 sum(\allowbreak{}k)\allowbreak{} C(\allowbreak{}i-\allowbreak{}1,\allowbreak{}k)\allowbreak{}C(\allowbreak{}n-\allowbreak{}i,\allowbreak{}k)\allowbreak{}=\allowbreak{}C(\allowbreak{}n-\allowbreak{}1,\allowbreak{}n-\allowbreak{}i)\allowbreak{}，常能消去逐中位数枚举。 排序后算子集贡献：第 i 个元素作为子集第 j 个数，系数为 C(\allowbreak{}i-\allowbreak{}1,\allowbreak{}j-\allowbreak{}1)\allowbreak{}*2\textasciicircum{}(\allowbreak{}n-\allowbreak{}i)\allowbreak{}。 若该位置权为 F[\allowbreak{}1]\allowbreak{}=\allowbreak{}1、F[\allowbreak{}j]\allowbreak{}=\allowbreak{}2\textasciicircum{}(\allowbreak{}j-\allowbreak{}2)\allowbreak{}(\allowbreak{}j>\allowbreak{}=\allowbreak{}2)\allowbreak{}，总系数=\allowbreak{}(\allowbreak{}3\textasciicircum{}(\allowbreak{}i-\allowbreak{}1)\allowbreak{}+\allowbreak{}1)\allowbreak{}/\allowbreak{}2*2\textasciicircum{}(\allowbreak{}n-\allowbreak{}i)\allowbreak{}。 “只在最短可行链计数”且两个端点都必须移动：k>\allowbreak{}=\allowbreak{}2 时 k!-\allowbreak{}2(\allowbreak{}k-\allowbreak{}1)\allowbreak{}!+\allowbreak{}(\allowbreak{}k-\allowbreak{}2)\allowbreak{}!；单点另计。 Hall：二分图左侧全匹配 iff 每个左点子集 S 都满足 |N(\allowbreak{}S)\allowbreak{}|>\allowbreak{}=\allowbreak{}|S|；小组数可枚举组子集。 带需求量的组展开成单位点或用流判定，不能只检查单组邻居够不够。}
% SOURCE: 05-数学/中国剩余定理.cpp
\subsection{中国剩余定理}

\note{a[\allowbreak{}1.\allowbreak{}.\allowbreak{}n]\allowbreak{}存余数，m[\allowbreak{}1.\allowbreak{}.\allowbreak{}n]\allowbreak{}存正模数；crt(\allowbreak{})\allowbreak{}要求两两互素，crt\_ex(\allowbreak{})\allowbreak{}允许不互素，无解为-\allowbreak{}1。p,\allowbreak{}s为模数前后缀积，每次重算，N须预留n+\allowbreak{}1槽；返回最小非负解，最终总模数须可存于ll。}
\note{适用：多个余数条件合并，普通 CRT 要模数两两互素，扩展版无需互素。 下标：a/\allowbreak{}m 的有效范围 1.\allowbreak{}.\allowbreak{}n，n>\allowbreak{}=\allowbreak{}1；所有模数为正，a[\allowbreak{}i]\allowbreak{} 为对应余数。 结论：可行解按总模数或 lcm 周期重复，返回最小非负代表。 关键：扩展合并要求 gcd(\allowbreak{}m1,\allowbreak{}m2)\allowbreak{} 整除 a2-\allowbreak{}a1，再解出模 m2/\allowbreak{}g 的增量。 易错：总乘积、lcm、a1+\allowbreak{}m1*x 的存储都必须在 ll 范围内；\_\_int128 中间值不代表最终赋值安全。 复杂度：n 次欧几里得合并约 O(\allowbreak{}n log M)\allowbreak{}，数组版空间 O(\allowbreak{}n)\allowbreak{}。}
\begin{lstlisting}
const int N=100005;
\end{lstlisting}
\note{O(\allowbreak{}log)\allowbreak{}，扩展欧几里得 O(\allowbreak{}log min(\allowbreak{}a,\allowbreak{}b)\allowbreak{})\allowbreak{}，x/\allowbreak{}y 引用返回裴蜀系数；模数使用正数。}
\begin{lstlisting}
ll exgcd(ll a,ll b,ll &x,ll &y)
{
    if(!b){x=1,y=0;return a;}
    ll xx,yy;
    ll g=exgcd(b,a%b,xx,yy);
    x=yy,y=xx-(a/b)*yy;
    return g;
}
\end{lstlisting}
\note{模数两两互质版 CRT：x =\allowbreak{}=\allowbreak{}=\allowbreak{} a[\allowbreak{}i]\allowbreak{} (\allowbreak{}mod m[\allowbreak{}i]\allowbreak{})\allowbreak{}，返回最小非负解 用前缀积 p[\allowbreak{}i]\allowbreak{}=\allowbreak{}m[\allowbreak{}1.\allowbreak{}.\allowbreak{}i]\allowbreak{}、后缀积 s[\allowbreak{}i]\allowbreak{}=\allowbreak{}m[\allowbreak{}i.\allowbreak{}.\allowbreak{}n]\allowbreak{} 把 M/\allowbreak{}m[\allowbreak{}i]\allowbreak{} 控制在 1e18 内 总模数乘积应 <\allowbreak{}=\allowbreak{} 1e18，否则解可能溢出 O(\allowbreak{}n log M)\allowbreak{}，合并 1.\allowbreak{}.\allowbreak{}n 条互素同余；前后缀需 n+\allowbreak{}1 槽，输入乘积须可存。}
\begin{lstlisting}
int n;
ll a[N],m[N],p[N],s[N];
ll crt()
{
    p[0]=1;
    for(int i=1;i<=n;i++)p[i]=p[i-1]*m[i];
    s[n+1]=1;
    for(int i=n;i>=1;i--)s[i]=s[i+1]*m[i];
    ll M=p[n],ans=0;
    for(int i=1;i<=n;i++)
    {
        ll Mi=p[i-1]*s[i+1];// M / m[i]
        ll x,y;
        exgcd(Mi%m[i],m[i],x,y);// Mi 在模 m[i] 下的逆元
        x=(x%m[i]+m[i])%m[i];
        ans=(ans+(__int128)a[i]*Mi%M*x)%M;
    }
    return (ans%M<0?ans%M+M:ans%M);
}
\end{lstlisting}
\note{扩展 CRT（模数不必互质）：合并 x =\allowbreak{}=\allowbreak{}=\allowbreak{} a1 (\allowbreak{}mod m1)\allowbreak{} 与 x =\allowbreak{}=\allowbreak{}=\allowbreak{} a2 (\allowbreak{}mod m2)\allowbreak{} 返回是否可合并；可合并时 a1 为新余数、m1 为新模数，均取最小非负 O(\allowbreak{}log)\allowbreak{} O(\allowbreak{}log min(\allowbreak{}m1,\allowbreak{}m2)\allowbreak{})\allowbreak{}，成功时原地更新 a1/\allowbreak{}m1；失败返回 false，m1 是 lcm 周期。}
\begin{lstlisting}
bool crt_merge(ll &a1,ll &m1,ll a2,ll m2)
{
    a1%=m1;if(a1<0)a1+=m1;
    a2%=m2;if(a2<0)a2+=m2;
    ll x,y;
    ll g=exgcd(m1,m2,x,y);
    ll d=a2-a1;
    if(d%g)return false;// 无解
    ll t=m2/g;
    x=(ll)((__int128)(d/g)%t*x%t);
    x=(x%t<0?x%t+t:x%t);
    __int128 next_mod=(__int128)(m1/g)*m2;
    assert(next_mod<=LLONG_MAX);
    a1=(a1+(__int128)m1*x)%next_mod;
    m1=(ll)next_mod;// lcm
    a1=(a1%m1<0?a1%m1+m1:a1%m1);
    return true;
}
\end{lstlisting}
\note{O(\allowbreak{}n log)\allowbreak{}，扩展 CRT 数组版，无解返回 -\allowbreak{}1（要求 lcm 不超过 9e18） O(\allowbreak{}n log M)\allowbreak{}，输入模数可不互素，返回最小非负解或 -\allowbreak{}1；n 必须非零。}
\begin{lstlisting}
ll crt_ex()
{
    ll ra=a[1],rm=m[1];
    for(int i=2;i<=n;i++)
        if(!crt_merge(ra,rm,a[i],m[i]))return -1;
    return (ra%rm<0?ra%rm+rm:ra%rm);
}
\end{lstlisting}
% SOURCE: 05-数学/矩阵快速幂.cpp
\subsection{矩阵快速幂}

\note{M为矩阵阶数，Matrix对象存元素；乘法的左右矩阵是独立运算对象，因此保留参数。幂的单位元为M阶单位矩阵，转移方向取决于状态是列向量还是行向量，模数及乘法宽度按实现设置。}
\note{适用：固定阶线性递推与状态转移重复执行，指数很大时用平方加速。 下标：Matrix.\allowbreak{}a[\allowbreak{}1.\allowbreak{}.\allowbreak{}M]\allowbreak{}[\allowbreak{}1.\allowbreak{}.\allowbreak{}M]\allowbreak{}，数组大小 10，故 1<\allowbreak{}=\allowbreak{}M<\allowbreak{}=\allowbreak{}9；fib 必须使用 M=\allowbreak{}2。 参数：qpow 的 n 是非负指数，矩阵元素先归一到 [\allowbreak{}0,\allowbreak{}mod)\allowbreak{}。 关键：乘法不交换，res*base 的次序必须符合既定状态列向量/\allowbreak{}行向量约定。 结论：零次幂是单位矩阵；斐波那契初值 F(\allowbreak{}1)\allowbreak{}=\allowbreak{}F(\allowbreak{}2)\allowbreak{}=\allowbreak{}1，n<\allowbreak{}=\allowbreak{}0 返回 0。 复杂度：矩阵乘 O(\allowbreak{}M\textasciicircum{}3)\allowbreak{}，幂 O(\allowbreak{}M\textasciicircum{}3 log n)\allowbreak{}，空间 O(\allowbreak{}M\textasciicircum{}2)\allowbreak{}；不要把 M 开到数组容量之外。}
\begin{lstlisting}
const int N=100005;
const ll mod=1000000007;
int M=2;// 矩阵阶数，用之前改成实际大小
\end{lstlisting}
\note{矩阵乘法，O(\allowbreak{}M\textasciicircum{}3)\allowbreak{}，取模版}
\begin{lstlisting}
struct Matrix{
    ll a[10][10];
    Matrix(){memset(a,0,sizeof(a));}
};
\end{lstlisting}
\note{单位矩阵 O(\allowbreak{}M)\allowbreak{}，构造对角为 1 的单位矩阵，配合零次幂。}
\begin{lstlisting}
Matrix I()
{
    Matrix r;
    for(int i=1;i<=M;i++)r.a[i][i]=1;
    return r;
}
\end{lstlisting}
\note{O(\allowbreak{}M\textasciicircum{}3)\allowbreak{}，x/\allowbreak{}y 同为 M 阶，返回 x*y；跳过零项节省无效乘加。}
\begin{lstlisting}
Matrix mul(Matrix x,Matrix y)
{
    Matrix r;
    for(int i=1;i<=M;i++)
        for(int k=1;k<=M;k++)
        {
            if(!x.a[i][k])continue;
            for(int j=1;j<=M;j++)
                r.a[i][j]=(r.a[i][j]+x.a[i][k]*y.a[k][j])%mod;
        }
    return r;
}
\end{lstlisting}
\note{矩阵快速幂，O(\allowbreak{}M\textasciicircum{}3 log n)\allowbreak{} O(\allowbreak{}M\textasciicircum{}3 log n)\allowbreak{}，base 为转移矩阵，n>\allowbreak{}=\allowbreak{}0；返回 base 的 n 次幂。}
\begin{lstlisting}
Matrix qpow(Matrix base,ll n)
{
    Matrix res=I();
    while(n)
    {
        if(n&1)res=mul(res,base);
        base=mul(base,base);
        n>>=1;
    }
    return res;
}
\end{lstlisting}
\note{斐波那契：F(\allowbreak{}1)\allowbreak{}=\allowbreak{}F(\allowbreak{}2)\allowbreak{}=\allowbreak{}1 [\allowbreak{}F(\allowbreak{}n+\allowbreak{}1)\allowbreak{} F(\allowbreak{}n)\allowbreak{}]\allowbreak{}\textasciicircum{}T =\allowbreak{} [\allowbreak{}[\allowbreak{}1,\allowbreak{}1]\allowbreak{},\allowbreak{}[\allowbreak{}1,\allowbreak{}0]\allowbreak{}]\allowbreak{}\textasciicircum{}(\allowbreak{}n-\allowbreak{}1)\allowbreak{} * [\allowbreak{}F(\allowbreak{}2)\allowbreak{} F(\allowbreak{}1)\allowbreak{}]\allowbreak{}\textasciicircum{}T O(\allowbreak{}log n)\allowbreak{}，需 M=\allowbreak{}2；使用 [\allowbreak{}[\allowbreak{}1,\allowbreak{}1]\allowbreak{},\allowbreak{}[\allowbreak{}1,\allowbreak{}0]\allowbreak{}]\allowbreak{}，返回第 n 项 mod mod。}
\begin{lstlisting}
ll fib(ll n)
{
    if(n<=0)return 0;
    if(n<=2)return 1;
    Matrix A;
    A.a[1][1]=1,A.a[1][2]=1,A.a[2][1]=1,A.a[2][2]=0;
    Matrix r=qpow(A,n-1);
    return r.a[1][1];
}
\end{lstlisting}
% SOURCE: 05-数学/高斯消元.cpp
\subsection{高斯消元}

\note{a[\allowbreak{}1.\allowbreak{}.\allowbreak{}n]\allowbreak{}[\allowbreak{}1.\allowbreak{}.\allowbreak{}n]\allowbreak{}为系数，a[\allowbreak{}i]\allowbreak{}[\allowbreak{}n+\allowbreak{}1]\allowbreak{}为常数项，消元原地修改增广矩阵。返回状态区分唯一解、无解、多解，不能无条件读取对角结果；再次求解前恢复原方程，EPS按数值尺度调整。}
\note{适用：实数线性方程组，判断唯一解、无穷多解、无解。 下标：a[\allowbreak{}1.\allowbreak{}.\allowbreak{}n]\allowbreak{}[\allowbreak{}1.\allowbreak{}.\allowbreak{}n+\allowbreak{}1]\allowbreak{} 为增广矩阵，n+\allowbreak{}1<\allowbreak{}N；最后一列是常数项。 关键：按列选绝对值最大主元减小误差；归一主元后消去该列其余行。 结论：0 表示唯一解且 a[\allowbreak{}i]\allowbreak{}[\allowbreak{}n+\allowbreak{}1]\allowbreak{} 为答案，1 表示无穷多解，-\allowbreak{}1 表示矛盾。 易错：原地覆盖 a；EPS 是绝对阈值，尺度悬殊或病态矩阵需额外误差分析。 复杂度：O(\allowbreak{}n\textasciicircum{}3)\allowbreak{} 时间、O(\allowbreak{}n\textasciicircum{}2)\allowbreak{} 存储；适用于实数，不直接用于模数方程。}
\begin{lstlisting}
const int N=105;
const double EPS=1e-8;
int n;
double a[N][N];// a[i][1..n] 系数，a[i][n+1] 常数项
\end{lstlisting}
\note{返回值：0 唯一解（解在 a[\allowbreak{}i]\allowbreak{}[\allowbreak{}n+\allowbreak{}1]\allowbreak{}），1 无穷多解，-\allowbreak{}1 无解 O(\allowbreak{}n\textasciicircum{}3)\allowbreak{}，列主元高斯-\allowbreak{}约当消元 O(\allowbreak{}n\textasciicircum{}3)\allowbreak{}，n 是方程数与未知数数；主元行 r 与列 c 不一定同步，零列是自由变量。}
\begin{lstlisting}
int gauss(int n)
{
    int r=1;// 当前主元行
    for(int c=1;c<=n;c++)
    {
        int pivot=r;
        for(int i=r+1;i<=n;i++)
            if(fabs(a[i][c])>fabs(a[pivot][c]))pivot=i;
        if(fabs(a[pivot][c])<EPS)continue;// 这一列全 0，跳过
        for(int j=1;j<=n+1;j++)swap(a[r][j],a[pivot][j]);
        double div=a[r][c];
        for(int j=1;j<=n+1;j++)a[r][j]/=div;
        for(int i=1;i<=n;i++)
        {
            if(i==r)continue;
            double mul=a[i][c];
            if(fabs(mul)<EPS)continue;
            for(int j=1;j<=n+1;j++)a[i][j]-=mul*a[r][j];
        }
        r++;
    }
    if(r<=n)
    {
        for(int i=r;i<=n;i++)
            if(fabs(a[i][n+1])>EPS)return -1;// 0=非0，无解
        return 1;// 有效方程数小于未知数个数，无穷多解
    }
    return 0;
}
\end{lstlisting}
% SOURCE: 05-数学/模意义与异或高斯消元.cpp
\subsection{模意义与异或高斯消元}

\note{模版系数矩阵与常数项组成线性系统，p为素数模；异或版每行是二进制系数位及常数位。x给自由变量取0的一组解，rank给异或版的秩，返回值区分无解/\allowbreak{}唯一/\allowbreak{}多解；未直接输出全部解的基，矩阵宽度由不同系统决定，保留对象参数。}
\note{增广矩阵 a：n 行、m+\allowbreak{}1 列，最后一列常数；p 为素数，允许负系数。 返回 -\allowbreak{}1 无解、0 唯一解、1 多解；x 返回一组解（自由变量取 0），原 a 不变。 不能套到合数模数！消元 O(\allowbreak{}n*m*min(\allowbreak{}n,\allowbreak{}m)\allowbreak{})\allowbreak{}，每个主元另需 O(\allowbreak{}log p)\allowbreak{} 求逆。}
\begin{lstlisting}
int gauss_mod(vector<vector<ll>> a,int m,ll p,vector<ll> &x)
{
    assert(m>=0&&p>=2);
    int n=a.size(),r=0;
    vector<int> where(m,-1);
    for(auto &row:a)
    {
        assert((int)row.size()==m+1);
        for(ll &v:row){v%=p;if(v<0)v+=p;}
    }
    auto power=[&](ll v,ll e){ll z=1;for(;e;e>>=1,v=(__int128)v*v%p)if(e&1)z=(__int128)z*v%p;return z;};
    for(int c=0;c<m&&r<n;c++)
    {
        int k=r;
        while(k<n&&!a[k][c])++k;
        if(k==n)continue;
        swap(a[k],a[r]),where[c]=r;
        ll inv=power(a[r][c],p-2);
        for(int j=c;j<=m;j++)a[r][j]=(__int128)a[r][j]*inv%p;
        for(int i=0;i<n;i++)if(i!=r&&a[i][c])
        {
            ll t=a[i][c];
            for(int j=c;j<=m;j++)
            {
                ll v=(a[i][j]-(__int128)t*a[r][j])%p;
                a[i][j]=v<0?v+p:v;
            }
        }
        ++r;
    }
    x.assign(m,0);
    for(int i=r;i<n;i++)if(a[i][m])return -1;
    for(int c=0;c<m;c++)if(where[c]>=0)x[c]=a[where[c]][m];
    return r==m?0:1;
}
\end{lstlisting}
\note{GF(\allowbreak{}2)\allowbreak{}：第 0.\allowbreak{}.\allowbreak{}m-\allowbreak{}1 位是系数，第 m 位是常数；m<\allowbreak{}B。 返回约定同上；原矩阵会被修改，rank 给出秩，多解时共有 2\textasciicircum{}(\allowbreak{}m-\allowbreak{}rank)\allowbreak{} 组。 bitset 行消元，O(\allowbreak{}n*m*ceil(\allowbreak{}B/\allowbreak{}机器字长)\allowbreak{})\allowbreak{}；x 的自由变量取 0。}
\begin{lstlisting}
template<size_t B>
int gauss_xor(vector<bitset<B>> &a,int m,bitset<B> &x,int &rank)
{
    assert(m>=0&&(size_t)m<B);
    int n=a.size();rank=0;x.reset();
    vector<int> where(m,-1);
    for(int c=0;c<m&&rank<n;c++)
    {
        int k=rank;
        while(k<n&&!a[k][c])++k;
        if(k==n)continue;
        swap(a[k],a[rank]),where[c]=rank;
        for(int i=0;i<n;i++)if(i!=rank&&a[i][c])a[i]^=a[rank];
        ++rank;
    }
    for(int i=rank;i<n;i++)if(a[i][m])return -1;
    for(int c=0;c<m;c++)if(where[c]>=0)x[c]=a[where[c]][m];
    return rank==m?0:1;
}
\end{lstlisting}
% SOURCE: 05-数学/线性基.cpp
\subsection{线性基}

\note{普通基按最高位存主元；规范形用于空间判等，区间基另记录pos为主元可用的最右位置。输入按无符号位宽，空集异或为0；位置阈值控制可用主元，不能用普通基直接删除过期向量。}
\note{适用：集合子集异或的最大值、可表示性与不同异或值的第 k 小。 位数：只处理 ll 的 0.\allowbreak{}.\allowbreak{}62 位，输入须非负；b[\allowbreak{}i]\allowbreak{} 最高位为 i，秩 cnt<\allowbreak{}=\allowbreak{}63。 关键：高位消元保证每个最高位唯一，相关向量不增加秩；insert 返回是否独立。 易错：query\_kth 包含空集 0，k 是 1-\allowbreak{}indexed；先 build，后插入新向量要重新 build。 k 用 unsigned long long，秩 63 时最后一项的序号是 2\textasciicircum{}63；zero 不改变不同值的计数。 空间判等：普通基依赖插入顺序；各自 build 后逐项比较 b[\allowbreak{}0.\allowbreak{}.\allowbreak{}62]\allowbreak{} 才是唯一规范形。 如 \{3,\allowbreak{}5\} 与 \{6,\allowbreak{}3\} 张成同一空间。判等不比较 zero：它记录输入相关性，不是空间性质。 哈希只能加速筛选，须考虑碰撞；全零非空区间也会生成零空间，不能直接丢弃。}
\note{异或线性基：支持插入、最大异或和、第 k 小异或和、判断能否异或得到 b[\allowbreak{}i]\allowbreak{} 存最高位为 i 的基向量；只有 build(\allowbreak{})\allowbreak{} 之后才能保证 b 是简化阶梯形}
\begin{lstlisting}
struct LinearBasis{
    ll b[64];
    int cnt;// 基的大小（秩）
    bool zero;// 是否存在异或为 0 的非空子集
    LinearBasis(){memset(b,0,sizeof(b));cnt=0;zero=false;}
    // O(63)，插入 x，返回 x 是否独立
    bool insert(ll x)
    {
        for(int i=62;i>=0;i--)
        {
            if(!(x>>i&1))continue;
            if(!b[i]){b[i]=x;cnt++;return true;}
            x^=b[i];
        }
        zero=true;// 异或成 0
        return false;
    }
    // O(63)，最大异或和
    ll query_max()
    {
        ll res=0;
        for(int i=62;i>=0;i--)
            if((res^b[i])>res)res^=b[i];
        return res;
    }
    // O(63)，能否把 x 异或出来（含 0，0 恒可）
    bool check(ll x)
    {
        if(!x)return true;
        for(int i=62;i>=0;i--)
        {
            if(!(x>>i&1))continue;
            if(!b[i])return false;
            x^=b[i];
        }
        return true;
    }
    // O(63^2)，化成唯一约化形；所有主元位在其他向量里均为 0
    // 第 k 小查询前必须先调用
    void build()
    {
        for(int i=62;i>=0;i--)
        {
            if(!b[i])continue;
            for(int j=i-1;j>=0;j--)
                if(b[i]>>j&1)b[i]^=b[j];
        }
    }
    // O(63)，枚举张成里的数，k 从 1 开始且 1<=k<=2^cnt
    // 第 1 小一定是 0（空集或某个异或为 0 的组合），之后是非零值升序
    ll query_kth(unsigned long long k)
    {
        if(k<1||k>(1ULL<<cnt))return -1;
        if(k==1)return 0;
        k--;// 去掉 0，剩下按秩枚举
        ll res=0;
        for(int i=0,j=0;i<=62;i++)
            if(b[i])
            {
                if(k>>j&1)res^=b[i];
                j++;
            }
        return res;
    }
};
\end{lstlisting}
\note{区间异或：按位置 1.\allowbreak{}.\allowbreak{}n 插入 a[\allowbreak{}r]\allowbreak{}，此时 query\_max(\allowbreak{}l)\allowbreak{} 查 [\allowbreak{}l,\allowbreak{}r]\allowbreak{} 的最大子集异或。 同主元优先保留靠右的向量，被换出的旧向量继续消元；不能只覆盖而丢掉它。 每个 b[\allowbreak{}i]\allowbreak{} 能由 pos[\allowbreak{}i]\allowbreak{}.\allowbreak{}.\allowbreak{}r 的元素表示，筛选 pos[\allowbreak{}i]\allowbreak{}>\allowbreak{}=\allowbreak{}l 后恰好张成 [\allowbreak{}l,\allowbreak{}r]\allowbreak{}。 插入、查询 O(\allowbreak{}63)\allowbreak{}，空间 O(\allowbreak{}63)\allowbreak{}。离线询问按右端点排序，不必保存所有前缀。}
\begin{lstlisting}
struct RangeLinearBasis
{
    ll b[63]{};
    int pos[63]{};
    void insert(ll x,int p)
    {
        for(int i=62;i>=0;i--)
        {
            if(!(x>>i&1))continue;
            if(!b[i]){b[i]=x;pos[i]=p;return;}
            if(pos[i]<p){swap(b[i],x);swap(pos[i],p);}
            x^=b[i];
        }
    }
    ll query_max(int l)
    {
        ll res=0;
        for(int i=62;i>=0;i--)
            if(pos[i]>=l&&(res^b[i])>res)res^=b[i];
        return res;
    }
};
\end{lstlisting}
\note{求区间空间的规范形：把满足 pos[\allowbreak{}i]\allowbreak{}>\allowbreak{}=\allowbreak{}l 的向量插入普通基，再 build，O(\allowbreak{}63\textasciicircum{}2)\allowbreak{}。 固定 r 移动 l，仅跨过至多 63 个有效位置时空间会改变，可跳过其余左端点。 位置基不能直接 build：反消可能混入更早的位置，破坏区间筛选的不变量。}
% SOURCE: 05-数学/博弈论(Nim与SG).cpp
\subsection{博弈论(Nim与SG)}

\note{n是堆数，a[\allowbreak{}1.\allowbreak{}.\allowbreak{}n]\allowbreak{}存非负石子数；nim\_win(\allowbreak{})\allowbreak{}判断胜负，nim\_first(\allowbreak{})\allowbreak{}返回\{堆号,\allowbreak{}操作后剩余数\}，不是取走数量。sg[\allowbreak{}x]\allowbreak{}存限取子游戏的SG值，改最大取数k后重新sg\_init。}
\note{适用：公平、有限、无平局的正常玩法，不能行动者输；独立子游戏 SG 值异或。 下标：堆 a[\allowbreak{}1.\allowbreak{}.\allowbreak{}n]\allowbreak{}，石子数非负；SG 下标 x 是剩余石子数，0<\allowbreak{}=\allowbreak{}x<\allowbreak{}N。 参数：get\_sg 的 k 是单次最多取数，0<\allowbreak{}=\allowbreak{}k<\allowbreak{}N；改变 k 前必须 sg\_init。 关键：mex 取所有后继 SG 中未出现的最小非负整数，终态 sg[\allowbreak{}0]\allowbreak{}=\allowbreak{}0。 易错：Nim 任意取数结论不能直接套到限取游戏或最后取者输；\_\_builtin\_clz 不接受 0。 结论：Nim 异或和为 0 为必败态；首次求 SG 至多 O(\allowbreak{}x*k)\allowbreak{}，递归深度 O(\allowbreak{}x)\allowbreak{}。}
\begin{lstlisting}
const int N=1005;
int n;
int a[N];// 每堆石子数
\end{lstlisting}
\note{Nim 博弈：异或和非 0 先手必胜，O(\allowbreak{}n)\allowbreak{} 取石子游戏，每步可从任意一堆取任意多个 O(\allowbreak{}n)\allowbreak{}，判断普通 Nim 先手胜负；a 的有效下标是 1.\allowbreak{}.\allowbreak{}n。}
\begin{lstlisting}
bool nim_win()
{
    int s=0;
    for(int i=1;i<=n;i++)s^=a[i];
    return s!=0;
}
\end{lstlisting}
\note{O(\allowbreak{}n)\allowbreak{}，Nim 的必胜第一步：返回 (\allowbreak{}堆编号,\allowbreak{} 取后剩余)\allowbreak{}，无必胜步返回 0 原理：找最高位，把某堆改成 a[\allowbreak{}i]\allowbreak{}\textasciicircum{}(\allowbreak{}s)\allowbreak{} 使异或和为 0 O(\allowbreak{}n)\allowbreak{}，返回堆编号及取后剩余数；返回 (\allowbreak{}0,\allowbreak{}0)\allowbreak{} 表示没有必胜步，不是取走数量。}
\begin{lstlisting}
pair<int,int> nim_first()
{
    int s=0;
    for(int i=1;i<=n;i++)s^=a[i];
    if(!s)return make_pair(0,0);
    int hb=31-__builtin_clz(s);// s 的最高位
    for(int i=1;i<=n;i++)
        if(a[i]>>hb&1)
        {
            int t=a[i]^s;
            if(t<a[i])return make_pair(i,t);
        }
    return make_pair(0,0);
}
\end{lstlisting}
\note{SG 函数（记忆化），位置为「剩余石子数」，可取 1.\allowbreak{}.\allowbreak{}k 个 单堆的 SG 值，O(\allowbreak{}n*k)\allowbreak{}；vis 必须是每次调用独立的，不能共用 static 数组}
\begin{lstlisting}
int sg[N];
\end{lstlisting}
\note{首次 O(\allowbreak{}x*k)\allowbreak{}，命中缓存 O(\allowbreak{}1)\allowbreak{}；x 是石子数，k 固定，局部 vis 避免递归互相污染。}
\begin{lstlisting}
int get_sg(int x,int k)
{
    if(sg[x]>=0)return sg[x];
    bool vis[N];
    for(int i=0;i<=k;i++)vis[i]=false;
    for(int i=1;i<=k&&i<=x;i++)vis[get_sg(x-i,k)]=true;
    int mex=0;
    while(vis[mex])mex++;
    return sg[x]=mex;
}
\end{lstlisting}
\note{清空 SG 表，调用前初始化（sg[\allowbreak{}0]\allowbreak{}=\allowbreak{}0，其余 -\allowbreak{}1） O(\allowbreak{}N)\allowbreak{}，整张 SG 表置 -\allowbreak{}1 再设终态；更换游戏取数规则必须重置。}
\begin{lstlisting}
void sg_init()
{
    memset(sg,-1,sizeof(sg));
    sg[0]=0;
}
\end{lstlisting}
% SOURCE: 05-数学/莫比乌斯反演与整除分块.cpp
\subsection{莫比乌斯反演与整除分块}

\note{mu与sum\_mu是莫比乌斯值和前缀和，prime/\allowbreak{}cnt/\allowbreak{}vis供筛；n为预处理范围。按商相同的闭区间[\allowbreak{}l,\allowbreak{}r]\allowbreak{}聚合前缀差，带权数组g的含义由具体反演式确定，不能只按变量名代入。}
\note{适用：gcd 计数、互素点对、按约数求和；mu 把 gcd=\allowbreak{}1 条件转为倍数计数。 参数：筛界 1<\allowbreak{}=\allowbreak{}n<\allowbreak{}N；查询须已筛到所用的 mu 前缀上界。 下标：mu/\allowbreak{}sum\_mu 从 1 开始，sum\_mu[\allowbreak{}0]\allowbreak{}=\allowbreak{}0；g 的有效定义域是 1.\allowbreak{}.\allowbreak{}gn。 关键：单商块末端 n/\allowbreak{}(\allowbreak{}n/\allowbreak{}l)\allowbreak{}，双商取两个末端的较小值，区间均为闭区间。 结论：sum\_coprime 数 1<\allowbreak{}=\allowbreak{}i<\allowbreak{}=\allowbreak{}n、1<\allowbreak{}=\allowbreak{}j<\allowbreak{}=\allowbreak{}m 的有序点对，复杂度 O(\allowbreak{}sqrt(\allowbreak{}n)\allowbreak{}+\allowbreak{}sqrt(\allowbreak{}m)\allowbreak{})\allowbreak{} 上界。 易错：多次筛前清 vis；前缀差先转 ll 再乘商，答案及中间乘积仍需满足 ll 范围。}
\begin{lstlisting}
const int N=1000006;
int n;
int prime[N],cnt;// prime[1..cnt] 存的素数
bool vis[N];// 合数标记
int mu[N];// 莫比乌斯函数
ll sum_mu[N];// mu 的前缀和（杜教筛 / 整除分块都要用）
\end{lstlisting}
\note{O(\allowbreak{}n)\allowbreak{}，线性筛出 mu 并求前缀和 O(\allowbreak{}n)\allowbreak{}，线性筛并求 1.\allowbreak{}.\allowbreak{}n 前缀；含平方因子的 mu 为 0，其余按素因子个数变号。}
\begin{lstlisting}
void get_mu(int n)
{
    cnt=0;
    fill(vis,vis+n+1,false);
    sum_mu[0]=0;
    mu[1]=1;
    for(int i=2;i<=n;i++)
    {
        if(!vis[i])
        {
            prime[++cnt]=i;
            mu[i]=-1;// 素数有 1 个质因子
        }
        for(int j=1;j<=cnt&&(ll)i*prime[j]<=n;j++)
        {
            int t=i*prime[j];
            vis[t]=true;
            if(i%prime[j]==0)
            {
                mu[t]=0;// 出现平方因子
                break;// 每个合数只被最小质因子筛一次
            }
            mu[t]=-mu[i];
        }
    }
    for(int i=1;i<=n;i++)sum_mu[i]=sum_mu[i-1]+mu[i];
}
\end{lstlisting}
\note{O(\allowbreak{}sqrt n)\allowbreak{}，整除分块：求 \ensuremath{\Sigma}\_\{i=\allowbreak{}1\}\textasciicircum{}\{n\} mu[\allowbreak{}i]\allowbreak{}*(\allowbreak{}n/\allowbreak{}i)\allowbreak{} O(\allowbreak{}sqrt n)\allowbreak{}，求 mu[\allowbreak{}i]\allowbreak{}*floor(\allowbreak{}n/\allowbreak{}i)\allowbreak{} 的和；n=\allowbreak{}0 时空和为 0。}
\begin{lstlisting}
ll sum_mu_div(int n)
{
    ll res=0;
    for(int l=1,r;l<=n;l=r+1)
    {
        r=n/(n/l);// [l,r] 内 n/i 相同
        res+=(ll)(sum_mu[r]-sum_mu[l-1])*(n/l);
    }
    return res;
}
\end{lstlisting}
\note{O(\allowbreak{}sqrt n)\allowbreak{}，整除分块：求 \ensuremath{\Sigma}\_\{i=\allowbreak{}1\}\textasciicircum{}\{n\} mu[\allowbreak{}i]\allowbreak{}*g(\allowbreak{}n/\allowbreak{}i)\allowbreak{}，g 是任意整数函数 要求 g 可 O(\allowbreak{}1)\allowbreak{} 求值；分块时 g 的参数 n/\allowbreak{}i 只会取到 O(\allowbreak{}sqrt n)\allowbreak{} 个不同值 O(\allowbreak{}sqrt n)\allowbreak{}，g[\allowbreak{}t]\allowbreak{} 是商 t 的函数值，gn 是最大合法下标；n>\allowbreak{}gn 返回 -\allowbreak{}1。}
\begin{lstlisting}
ll sum_mu_mul_g(int n,ll g[],int gn)
{
    if(n>gn)return -1;// g 的定义域不够，调用者自己保证 n<=gn
    ll res=0;
    for(int l=1,r;l<=n;l=r+1)
    {
        r=n/(n/l);// [l,r] 内 n/i 相同
        res+=(ll)(sum_mu[r]-sum_mu[l-1])*g[n/l];
    }
    return res;
}
\end{lstlisting}
\note{O(\allowbreak{}sqrt n)\allowbreak{}，合并等商区间求 floor(\allowbreak{}n/\allowbreak{}i)\allowbreak{} 总和；原 O(\allowbreak{}log n)\allowbreak{} 注释不适用于此循环。}
\begin{lstlisting}
ll sum_floor(int n)
{
    ll res=0;
    for(int l=1,r;l<=n;l=r+1)
    {
        r=n/(n/l);
        res+=(ll)(r-l+1)*(n/l);
    }
    return res;
}
\end{lstlisting}
\note{O(\allowbreak{}sqrt n)\allowbreak{}，模板核心：\ensuremath{\Sigma}\ensuremath{\Sigma}[\allowbreak{}gcd(\allowbreak{}i,\allowbreak{}j)\allowbreak{}=\allowbreak{}=\allowbreak{}1]\allowbreak{}，i<\allowbreak{}=\allowbreak{}n,\allowbreak{}j<\allowbreak{}=\allowbreak{}m 反演后 =\allowbreak{} \ensuremath{\Sigma}\_\{d=\allowbreak{}1\}\textasciicircum{}\{n\} mu[\allowbreak{}d]\allowbreak{}*(\allowbreak{}n/\allowbreak{}d)\allowbreak{}*(\allowbreak{}m/\allowbreak{}d)\allowbreak{} O(\allowbreak{}sqrt(\allowbreak{}n)\allowbreak{}+\allowbreak{}sqrt(\allowbreak{}m)\allowbreak{})\allowbreak{} 上界，利用 mu 前缀计互素有序对；筛界至少 min(\allowbreak{}n,\allowbreak{}m)\allowbreak{}。}
\begin{lstlisting}
ll sum_coprime(int n,int m)
{
    if(n>m)swap(n,m);
    ll res=0;
    for(int l=1,r;l<=n;l=r+1)
    {
        r=min(n/(n/l),m/(m/l));// 两个商同时不变的最远位置
        res+=(ll)(sum_mu[r]-sum_mu[l-1])*(n/l)*(m/l);
    }
    return res;
}
\end{lstlisting}
% SOURCE: 05-数学/类欧几里得(floor_sum).cpp
\subsection{类欧几里得(floor\_sum)}

\note{n是项数，i遍历0.\allowbreak{}.\allowbreak{}n-\allowbreak{}1，求floor(\allowbreak{}(\allowbreak{}a*i+\allowbreak{}b)\allowbreak{}/\allowbreak{}m)\allowbreak{}之和，m>\allowbreak{}0；a/\allowbreak{}b可为有符号参数。先拆整数部分，再交换余数问题，乘积与总和采用实现的宽整数范围。}
\note{\ensuremath{\Sigma}\_\{i=\allowbreak{}0\}\textasciicircum{}\{n-\allowbreak{}1\} floor(\allowbreak{}(\allowbreak{}a*i+\allowbreak{}b)\allowbreak{}/\allowbreak{}m)\allowbreak{}，n>\allowbreak{}=\allowbreak{}0、m>\allowbreak{}0，a/\allowbreak{}b 可负，返回值须放入 ll。 C+\allowbreak{}+\allowbreak{} 负数除法向 0 截断，先改为向下取整；a*n+\allowbreak{}b 与累计用 \_\_int128。 去掉整商后交换横纵轴，类似欧几里得；O(\allowbreak{}log m)\allowbreak{}，空间 O(\allowbreak{}1)\allowbreak{}。 FS 记本函数，u=\allowbreak{}(\allowbreak{}f*t+\allowbreak{}x)\allowbreak{} mod m、v=\allowbreak{}(\allowbreak{}g*t+\allowbreak{}y)\allowbreak{} mod m，余数规范到 [\allowbreak{}0,\allowbreak{}m)\allowbreak{}。 \#满足 u<\allowbreak{}v (\allowbreak{}0<\allowbreak{}=\allowbreak{}t<\allowbreak{}n)\allowbreak{}=\allowbreak{}n-\allowbreak{}FS(\allowbreak{}n,\allowbreak{}m,\allowbreak{}g,\allowbreak{}y)\allowbreak{}+\allowbreak{}FS(\allowbreak{}n,\allowbreak{}m,\allowbreak{}f,\allowbreak{}x)\allowbreak{}+\allowbreak{}FS(\allowbreak{}n,\allowbreak{}m,\allowbreak{}g-\allowbreak{}f,\allowbreak{}y-\allowbreak{}x-\allowbreak{}1)\allowbreak{}。 因为 [\allowbreak{}u<\allowbreak{}v]\allowbreak{}=\allowbreak{}1+\allowbreak{}floor(\allowbreak{}(\allowbreak{}v-\allowbreak{}u-\allowbreak{}1)\allowbreak{}/\allowbreak{}m)\allowbreak{}；第三项可能有负参数，f/\allowbreak{}g/\allowbreak{}x/\allowbreak{}y 的差也须防溢出。}
\begin{lstlisting}
ll floor_sum(ll n,ll m,ll a,ll b)
{
    assert(n>=0&&m>0);
    ll qa=a/m,qb=b/m;
    a%=m,b%=m;
    if(a<0)a+=m,--qa;
    if(b<0)b+=m,--qb;
    __int128 ans=(__int128)n*(n-1)/2*qa+(__int128)n*qb;
    while(1)
    {
        if(a>=m)ans+=(__int128)n*(n-1)/2*(a/m),a%=m;
        if(b>=m)ans+=(__int128)n*(b/m),b%=m;
        __int128 y=(__int128)a*n+b;
        if(y<m)break;
        n=y/m,b=y%m,swap(a,m);
    }
    assert(LLONG_MIN<=ans&&ans<=LLONG_MAX);
    return (ll)ans;
}
\end{lstlisting}
% SOURCE: 05-数学/MillerRabin与PollardRho.cpp
\subsection{MillerRabin与PollardRho}

\note{n为待判定/\allowbreak{}分解整数，因子向量允许重复，整理后得到质因子及指数。Miller–Rabin基集针对实现声明的位宽，Rho使用随机多项式；1无质因子，模乘保持宽类型。}
\note{适用：大整数素性判断和质因数分解；实际接口是正的有符号 ll。 参数：n<\allowbreak{}=\allowbreak{}LLONG\_MAX；分解入口 factor 要 n>\allowbreak{}=\allowbreak{}1，pollard\_rho 只接受合数。 关键：n-\allowbreak{}1=\allowbreak{}d*2\textasciicircum{}s，强伪素数测试不断平方；固定七个底数覆盖接口可表示范围。 随机：Rho 返回的是任意非平凡因子，不保证素数，继续递归分解。 模乘和 Rho 的平方加常数均用 \_\_int128；仅支持正的 signed ll。 复杂度：七底数素性判定 O(\allowbreak{}log n)\allowbreak{} 次模乘；Rho 期望约 O(\allowbreak{}n\textasciicircum{}(\allowbreak{}1/\allowbreak{}4)\allowbreak{})\allowbreak{}，无确定时间上界。}
\begin{lstlisting}
typedef unsigned long long ull;
typedef __int128 lll;
\end{lstlisting}
\note{O(\allowbreak{}log\textasciicircum{}3 n)\allowbreak{}，Miller-\allowbreak{}Rabin 素性判定，确定性基组覆盖 64 位 O(\allowbreak{}1)\allowbreak{}，计算非负 a*b mod mod；mod>\allowbreak{}0，先扩为 \_\_int128。}
\begin{lstlisting}
ll qmul(ll a,ll b,ll mod)
{
    return (ll)((lll)a*b%mod);
}
\end{lstlisting}
\note{O(\allowbreak{}log n)\allowbreak{}，模 mod 快速幂；n 是非负指数。}
\begin{lstlisting}
ll qpow(ll a,ll n,ll mod)
{
    ll res=1%mod;
    a%=mod;
    while(n)
    {
        if(n&1)res=qmul(res,a,mod);
        a=qmul(a,a,mod);
        n>>=1;
    }
    return res;
}
\end{lstlisting}
\note{O(\allowbreak{}log n)\allowbreak{} 次宽整数模乘，返回 n 是否为素数；底数为 n 的倍数时跳过。}
\begin{lstlisting}
bool miller_rabin(ll n)
{
    if(n<2)return false;
    if(n==2||n==3)return true;
    if(n%2==0)return false;
    // 写成 n-1 = d*2^s
    ll d=n-1;
    int s=0;
    while(!(d&1))d>>=1,s++;
    ll base[7]={2,325,9375,28178,450775,9780504,1795265022};
    for(int i=0;i<7;i++)
    {
        ll a=base[i]%n;
        if(a==0)continue;
        ll x=qpow(a,d,n);
        if(x==1||x==n-1)continue;
        bool ok=false;
        for(int j=1;j<s;j++)
        {
            x=qmul(x,x,n);
            if(x==n-1){ok=true;break;}
        }
        if(!ok)return false;// 一定是合数
    }
    return true;
}
\end{lstlisting}
\note{O(\allowbreak{}n\textasciicircum{}(\allowbreak{}1/\allowbreak{}4)\allowbreak{})\allowbreak{}，Pollard-\allowbreak{}Rho 找 n 的一个非平凡因子，n 必须是合数 期望约 O(\allowbreak{}n\textasciicircum{}(\allowbreak{}1/\allowbreak{}4)\allowbreak{})\allowbreak{}，找合数 n 的因子；批量积减少 gcd 次数，d=\allowbreak{}=\allowbreak{}n 表示本轮失败。}
\begin{lstlisting}
ll pollard_rho(ll n)
{
    if(n%2==0)return 2;
    if(n%3==0)return 3;
    while(true)
    {
        static mt19937_64 rnd(chrono::steady_clock::now().time_since_epoch().count());
        ll c=uniform_int_distribution<ll>(1,n-1)(rnd);
        ll x=uniform_int_distribution<ll>(1,n-1)(rnd),y=x,d=1;
        // 倍增步长 + gcd 批量化
        ll q=1;
        for(ll len=1;d==1;len<<=1)
        {
            ll tx=x;
            for(ll i=1;i<=len;i++)
            {
                x=((lll)qmul(x,x,n)+c)%n;
                q=qmul(q,abs(x-y),n);
                if(i%127==0)
                {
                    d=__gcd(q,n);
                    if(d>1)break;
                }
            }
            if(d==1)d=__gcd(q,n);
            y=x;
            if(d==n){x=tx;break;}// 失败重来
        }
        if(d>1&&d<n)return d;
    }
}
\end{lstlisting}
\note{O(\allowbreak{}n\textasciicircum{}(\allowbreak{}1/\allowbreak{}4)\allowbreak{} log n)\allowbreak{}，递归分解出全部素因子（不排序、含重数） 随机递归分解，v 为追加输出而非覆盖；n=\allowbreak{}1 不追加，重复因子保留，结果需自行排序。}
\begin{lstlisting}
void factor(ll n,vector<ll> &v)
{
    if(n==1)return;
    if(miller_rabin(n)){v.push_back(n);return;}
    ll d=pollard_rho(n);
    // 随机递归分解，v 为追加输出而非覆盖；n=1 不追加，重复因子保留，结果需自行排序。
    factor(d,v);
    // 随机递归分解，v 为追加输出而非覆盖；n=1 不追加，重复因子保留，结果需自行排序。
    factor(n/d,v);
}
\end{lstlisting}
% SOURCE: 05-数学/BSGS与扩展BSGS.cpp
\subsection{BSGS与扩展BSGS}

\note{a、b、mod分别是底数、目标和模数，求a\textasciicircum{}x\ensuremath{\equiv}b；扩展版先剥离与模数的公因子。指数从0起，无解使用接口哨兵，模1、0底数及溢出边界见实现。}
\note{适用：求最小 x>\allowbreak{}=\allowbreak{}0 使 a\textasciicircum{}x=\allowbreak{}b (\allowbreak{}mod p)\allowbreak{}，无解返回 -\allowbreak{}1。 参数：a,\allowbreak{}b 是底数与目标余数，p>\allowbreak{}=\allowbreak{}1；p=\allowbreak{}1 返回 0，约定 a\textasciicircum{}0=\allowbreak{}1。 关键：指数写成 i*m+\allowbreak{}j，baby 表保留相同余数的最小 j，按 i 递增取最小解。 扩展：逐次消去 gcd(\allowbreak{}a,\allowbreak{}p)\allowbreak{}，记录已消去的指数 cnt；b 不能被公因子整除即无解。 易错：模乘用 \_\_int128，但 (\allowbreak{}a\%p+\allowbreak{}p)\allowbreak{} 等归一化加法仍须不溢出 ll。 复杂度：排序版 O(\allowbreak{}sqrt(\allowbreak{}p)\allowbreak{}*log p)\allowbreak{} 时间、O(\allowbreak{}sqrt(\allowbreak{}p)\allowbreak{})\allowbreak{} 空间；大模数内存是瓶颈。}
\begin{lstlisting}
typedef __int128 lll;
\end{lstlisting}
\note{O(\allowbreak{}log n)\allowbreak{}，模 p 快速幂；n>\allowbreak{}=\allowbreak{}0，乘法先扩为 \_\_int128。}
\begin{lstlisting}
ll qpow(ll a,ll n,ll p)
{
    ll ans=1%p;
    for(a%=p;n;n>>=1,a=(lll)a*a%p)if(n&1)ans=(lll)ans*a%p;
    return ans;
}
\end{lstlisting}
\note{O(\allowbreak{}log min(\allowbreak{}a,\allowbreak{}b)\allowbreak{})\allowbreak{}，返回 gcd，引用 x/\allowbreak{}y 输出 ax+\allowbreak{}by=\allowbreak{}gcd 的系数。}
\begin{lstlisting}
ll exgcd(ll a,ll b,ll &x,ll &y)
{
    if(!b){x=1,y=0;return a;}
    ll u,v,g=exgcd(b,a%b,u,v);
    x=v,y=u-a/b*v;
    return g;
}
\end{lstlisting}
\note{O(\allowbreak{}log p)\allowbreak{}，返回 a 的最小非负逆元；调用者须保证 gcd(\allowbreak{}a,\allowbreak{}p)\allowbreak{}=\allowbreak{}1。}
\begin{lstlisting}
ll inv_mod(ll a,ll p)
{
    ll x,y;
    exgcd(a,p,x,y);
    return (x%p<0?x%p+p:x%p);
}
\end{lstlisting}
\note{O(\allowbreak{}sqrt(\allowbreak{}p)\allowbreak{} log p)\allowbreak{}，互素离散对数，返回最小非负指数 O(\allowbreak{}sqrt(\allowbreak{}p)\allowbreak{}*log p)\allowbreak{}，仅处理底数与模数互素的离散对数；baby 表排序去重保留最小指数。}
\begin{lstlisting}
ll bsgs(ll a,ll b,ll p)
{
    assert(p>=1);
    if(p==1)return 0;
    a=(a%p<0?a%p+p:a%p),b=(b%p<0?b%p+p:b%p);
    if(b==1)return 0;
    if(gcd(a,p)!=1)return -1;// 非互素请用扩展版本
    ll m=sqrtl(p)+1,cur=1;
    vector<pair<ll,ll> > v;
    for(ll j=0;j<m;j++)v.push_back({cur,j}),cur=(lll)cur*a%p;
    sort(v.begin(),v.end());
    v.erase(unique(v.begin(),v.end(),[](pair<ll,ll> x,pair<ll,ll> y){return x.first==y.first;}),v.end());
    ll step=inv_mod(qpow(a,m,p),p);
    cur=b;
    for(ll i=0;i<=m;i++,cur=(lll)cur*step%p)
    {
        int j=lower_bound(v.begin(),v.end(),make_pair(cur,-1LL))-v.begin();
        if(j<(int)v.size()&&v[j].first==cur)return i*m+v[j].second;
    }
    return -1;
}
\end{lstlisting}
\note{消去公因子后做 BSGS，p>\allowbreak{}=\allowbreak{}1，约定 a\textasciicircum{}0=\allowbreak{}1 消公因子 O(\allowbreak{}log p)\allowbreak{} 次后做 BSGS；mul 保存被消部分，最终指数加 cnt。}
\begin{lstlisting}
ll exbsgs(ll a,ll b,ll p)
{
    assert(p>=1);
    if(p==1)return 0;
    a=(a%p<0?a%p+p:a%p),b=(b%p<0?b%p+p:b%p);
    if(b==1)return 0;
    ll cnt=0,mul=1,g;
    while((g=gcd(a,p))>1)
    {
        if(b%g)return -1;
        b/=g,p/=g,cnt++;
        mul=(lll)mul*(a/g)%p;
        if(mul==b)return cnt;
    }
    ll ans=bsgs(a,(lll)b*inv_mod(mul,p)%p,p);
    return ans<0?-1:ans+cnt;
}
\end{lstlisting}
% SOURCE: 05-数学/NTT.cpp
\subsection{NTT}

\note{系数从常数项起，变换长度是2的幂且整除模数-\allowbreak{}1；rev及单位根缓存对应当前长度。原地变换改系数，卷积保留两个独立操作数；逆变换乘长度逆元，不能任意更换模数和原根。}
\note{适用：模 998244353 的整数卷积；没有 FFT 的浮点舍入误差。 下标：多项式系数 0.\allowbreak{}.\allowbreak{}n-\allowbreak{}1/\allowbreak{}0.\allowbreak{}.\allowbreak{}m-\allowbreak{}1；数组长度是项数而非次数。 长度：n,\allowbreak{}m>\allowbreak{}=\allowbreak{}1；补零长度 len 为 2 的幂，必须 <\allowbreak{}=\allowbreak{}N 且整除 mod-\allowbreak{}1。 关键：蝶形把两半的偶奇贡献合并；逆变换使用逆单位根并乘长度逆元。 易错：输入须先归一到 [\allowbreak{}0,\allowbreak{}mod)\allowbreak{}；负系数直接传入会留下负余数。 复杂度：卷积 O(\allowbreak{}L log L)\allowbreak{}、空间 O(\allowbreak{}L)\allowbreak{}，L 是补零长度；静态缓冲不能并发共享。}
\begin{lstlisting}
const int N=1<<19;
const ll mod=998244353;// 常用 NTT 模数，原根 3，mod-1=119*2^23
const ll g=3;
\end{lstlisting}
\note{O(\allowbreak{}log n)\allowbreak{}，快速幂 O(\allowbreak{}log n)\allowbreak{}，a 为底数、n 为非负指数；此处模数约 1e9，ll 模乘安全。}
\begin{lstlisting}
ll qpow(ll a,ll n,ll mod)
{
    ll res=1;
    a%=mod;
    while(n)
    {
        if(n&1)(res*=a)%=mod;
        (a*=a)%=mod;
        n>>=1;
    }
    return res;
}
\end{lstlisting}
\note{O(\allowbreak{}n log n)\allowbreak{}，NTT 迭代版 in-\allowbreak{}place；inv=\allowbreak{}0 正变换，inv=\allowbreak{}1 逆变换 n 必须是 2 的幂且 n | (\allowbreak{}mod-\allowbreak{}1)\allowbreak{}，len 传实际长度 O(\allowbreak{}n log n)\allowbreak{}，a[\allowbreak{}0.\allowbreak{}.\allowbreak{}n-\allowbreak{}1]\allowbreak{} 原地变换，inv=\allowbreak{}0/\allowbreak{}1 表示正/\allowbreak{}逆；位反转保证迭代合并顺序。}
\begin{lstlisting}
void ntt(ll a[],int n,int inv)
{
    // 位反转置换
    for(int i=1,j=0;i<n;i++)
    {
        int bit=n>>1;
        for(;j&bit;bit>>=1)j^=bit;
        j^=bit;
        if(i<j)swap(a[i],a[j]);
    }
    for(int len=2;len<=n;len<<=1)
    {
        ll w=qpow(g,(mod-1)/len,mod);
        if(inv)w=qpow(w,mod-2,mod);
        for(int i=0;i<n;i+=len)
        {
            ll wn=1;
            for(int k=0;k<len/2;k++)
            {
                ll u=a[i+k],v=a[i+k+len/2]*wn%mod;
                a[i+k]=(u+v)%mod;
                a[i+k+len/2]=(u-v+mod)%mod;
                wn=wn*w%mod;
            }
        }
    }
    if(inv)
    {
        ll ninv=qpow(n,mod-2,mod);
        for(int i=0;i<n;i++)a[i]=a[i]*ninv%mod;
    }
}
\end{lstlisting}
\note{O(\allowbreak{}n log n)\allowbreak{}，多项式乘法，结果为 a[\allowbreak{}0.\allowbreak{}.\allowbreak{}n-\allowbreak{}1]\allowbreak{} * b[\allowbreak{}0.\allowbreak{}.\allowbreak{}m-\allowbreak{}1]\allowbreak{} 的系数 O(\allowbreak{}L log L)\allowbreak{}，输出 c[\allowbreak{}0.\allowbreak{}.\allowbreak{}clen-\allowbreak{}1]\allowbreak{}；clen 引用返回 n+\allowbreak{}m-\allowbreak{}1，c 至少有这些槽。}
\begin{lstlisting}
void poly_mul(ll a[],int n,ll b[],int m,ll c[],int &clen)
{
    if(n<=0||m<=0){clen=0;return;}
    int len=1;
    while(len<n+m-1)len<<=1;
    assert(len<=N);
    static ll x[N],y[N];
    for(int i=0;i<len;i++)x[i]=(i<n?(a[i]%mod+mod)%mod:0),y[i]=(i<m?(b[i]%mod+mod)%mod:0);
    ntt(x,len,0);
    ntt(y,len,0);
    for(int i=0;i<len;i++)x[i]=x[i]*y[i]%mod;
    ntt(x,len,1);
    clen=n+m-1;
    for(int i=0;i<clen;i++)c[i]=x[i];
}
\end{lstlisting}
% SOURCE: 05-数学/FFT.cpp
\subsection{FFT}

\note{输入系数向量按常数项到高次项排列，变换缓冲为复数且长度为2的幂；两个多项式是独立操作数，保留参数。整数卷积由逆变换四舍五入，浮点精度与范围限制不能省略。}
\note{适用：整数多项式卷积；结果以 llround 还原整数，系数太大会有精度误差。 下标：a[\allowbreak{}i]\allowbreak{}/\allowbreak{}b[\allowbreak{}i]\allowbreak{} 对应 x\textasciicircum{}i，0-\allowbreak{}indexed；空多项式相乘返回空。 参数：fft 的 inv=\allowbreak{}1 正变换、inv=\allowbreak{}-\allowbreak{}1 逆变换，不是 0/\allowbreak{}1 约定。 关键：长度补到 >\allowbreak{}=\allowbreak{}a.\allowbreak{}size(\allowbreak{})\allowbreak{}+\allowbreak{}b.\allowbreak{}size(\allowbreak{})\allowbreak{}-\allowbreak{}1 的 2 的幂，避免循环卷积混叠。 易错：double 约 53 位有效二进制位，负系数也要舍入；高精度或大系数考虑拆分或 NTT。 复杂度：O(\allowbreak{}L log L)\allowbreak{} 时间、O(\allowbreak{}L)\allowbreak{} 空间；输入按值传递，不修改调用者数组。}
\begin{lstlisting}
const double PI=acos(-1.0);
\end{lstlisting}
\note{O(\allowbreak{}n log n)\allowbreak{}，复数 FFT，系数规模需保证浮点舍入可靠 O(\allowbreak{}L log L)\allowbreak{}，原地变换长度 L 的复数向量，L>\allowbreak{}=\allowbreak{}1 且为 2 的幂；逆变换最后除 L。}
\begin{lstlisting}
void fft(vector<complex<double> > &a,int inv)
{
    int n=a.size();
    for(int i=1,j=0;i<n;i++)
    {
        int bit=n>>1;
        for(;j&bit;bit>>=1)j^=bit;
        j^=bit;
        if(i<j)swap(a[i],a[j]);
    }
    for(int len=2;len<=n;len<<=1)
    {
        complex<double> wlen(cos(2*PI/len),inv*sin(2*PI/len));
        for(int i=0;i<n;i+=len)
        {
            complex<double> w(1,0);
            for(int j=0;j<len/2;j++,w*=wlen)
            {
                complex<double> u=a[i+j],v=a[i+j+len/2]*w;
                a[i+j]=u+v,a[i+j+len/2]=u-v;
            }
        }
    }
    if(inv==-1)for(complex<double> &x:a)x/=n;
}
\end{lstlisting}
\note{O(\allowbreak{}L log L)\allowbreak{}，输出 a.\allowbreak{}size(\allowbreak{})\allowbreak{}+\allowbreak{}b.\allowbreak{}size(\allowbreak{})\allowbreak{}-\allowbreak{}1 个整数系数；点乘频域等价于时域卷积。}
\begin{lstlisting}
vector<ll> multiply(vector<ll> a,vector<ll> b)
{
    if(a.empty()||b.empty())return {};
    int sz=a.size()+b.size()-1,n=1;
    while(n<sz)n<<=1;
    vector<complex<double> > x(n),y(n);
    for(int i=0;i<(int)a.size();i++)x[i]=a[i];
    for(int i=0;i<(int)b.size();i++)y[i]=b[i];
    fft(x,1),fft(y,1);
    for(int i=0;i<n;i++)x[i]*=y[i];
    fft(x,-1);
    vector<ll> ans(sz);
    for(int i=0;i<sz;i++)ans[i]=llround(x[i].real());
    return ans;
}
\end{lstlisting}
% SOURCE: 05-数学/FWT(按位卷积).cpp
\subsection{FWT(按位卷积)}

\note{a/\allowbreak{}b按整数掩码为下标，长度必须相同且为2的幂；type选择OR/\allowbreak{}AND/\allowbreak{}XOR，inverse选择逆变换。fwt原地改缓冲，卷积参数区分两份输入；XOR逆变换需2可逆。}
\begin{lstlisting}
const ll MOD=998244353,INV2=(MOD+1)/2;
\end{lstlisting}
\note{type：0 XOR，1 AND，2 OR；长度必须为 2 的幂，系数先归一到 [\allowbreak{}0,\allowbreak{}MOD)\allowbreak{}。 XOR：(\allowbreak{}u+\allowbreak{}v,\allowbreak{}u-\allowbreak{}v)\allowbreak{}，逆变换每层乘 1/\allowbreak{}2；AND：(\allowbreak{}u+\allowbreak{}v,\allowbreak{}v)\allowbreak{}；OR：(\allowbreak{}u,\allowbreak{}u+\allowbreak{}v)\allowbreak{}。 与 FFT 不同，结果下标是 i\textasciicircum{}j /\allowbreak{} i\&j /\allowbreak{} i|j，长度取 >\allowbreak{}=\allowbreak{}max(\allowbreak{}输入项数)\allowbreak{}。}
\begin{lstlisting}
void fwt(vector<ll> &a,int type,bool inverse=false)
{
    int n=a.size();
    assert(n>0&&(n&(n-1))==0&&0<=type&&type<=2);
    for(int len=1;len<n;len*=2)
        for(int i=0;i<n;i+=len*2)
            for(int j=0;j<len;j++)
            {
                ll u=a[i+j],v=a[i+j+len];
                if(type==0)
                {
                    a[i+j]=(u+v)%MOD,a[i+j+len]=(u-v+MOD)%MOD;
                    if(inverse)a[i+j]=a[i+j]*INV2%MOD,a[i+j+len]=a[i+j+len]*INV2%MOD;
                }
                else if(type==1)a[i+j]=(u+(inverse?MOD-v:v))%MOD;
                else a[i+j+len]=(v+(inverse?MOD-u:u))%MOD;
            }
}
\end{lstlisting}
\note{O(\allowbreak{}L log L)\allowbreak{}，返回补到 L 项后的卷积；XOR 逆变换要求 2 在模数下可逆。}
\begin{lstlisting}
vector<ll> bit_convolution(vector<ll> a,vector<ll> b,int type)
{
    if(a.empty()||b.empty())return {};
    int n=1;
    while(n<(int)max(a.size(),b.size()))n*=2;
    a.resize(n),b.resize(n);
    for(ll &x:a){x%=MOD;if(x<0)x+=MOD;}
    for(ll &x:b){x%=MOD;if(x<0)x+=MOD;}
    fwt(a,type),fwt(b,type);
    for(int i=0;i<n;i++)a[i]=a[i]*b[i]%MOD;
    fwt(a,type,true);
    return a;
}
\end{lstlisting}
% SOURCE: 05-数学/多项式求逆与ln_exp.cpp
\subsection{多项式求逆与ln\_exp}

\note{Bell 数的指数生成函数；最终系数乘 $n!$，模数须允许所需分母求逆。}
\[
\sum_{n\ge0}B_n\frac{x^n}{n!}=\exp(e^x-1).
\]

\note{向量下标为次数，n为截断项数；Newton迭代同时持有原多项式和多个中间结果，保留独立对象参数。求逆常数项非零，ln常数项为1，exp常数项为0，模数及可逆范围按实现设置。}
\note{适用：生成函数运算；均在模 P 的形式幂级数中计算。 参数：a[\allowbreak{}i]\allowbreak{} 为 x\textasciicircum{}i 的系数，0-\allowbreak{}indexed；n>\allowbreak{}=\allowbreak{}1 为保留项数，结果取模 x\textasciicircum{}n。 前提：输入非空且系数在 [\allowbreak{}0,\allowbreak{}P)\allowbreak{}，求逆 a[\allowbreak{}0]\allowbreak{}!=\allowbreak{}0，ln 要 a[\allowbreak{}0]\allowbreak{}=\allowbreak{}1，exp 要 a[\allowbreak{}0]\allowbreak{}=\allowbreak{}0。 关键：牛顿迭代每次把正确项数翻倍；低次项截断不影响后续模 x\textasciicircum{}n 的结果。 易错：积分除以次数，需 n<\allowbreak{}P；完整卷积长度补成 2 的幂且不超过 2\textasciicircum{}23。 复杂度：求逆/\allowbreak{}ln/\allowbreak{}exp 按输入截至 O(\allowbreak{}n)\allowbreak{} 时 O(\allowbreak{}n log n)\allowbreak{}，额外空间 O(\allowbreak{}n)\allowbreak{}；超长输入会增加卷积成本。 Bell 数 EGF：sum(\allowbreak{}n>\allowbreak{}=\allowbreak{}0)\allowbreak{} B[\allowbreak{}n]\allowbreak{}*x\textasciicircum{}n/\allowbreak{}n! =\allowbreak{} exp(\allowbreak{}exp(\allowbreak{}x)\allowbreak{}-\allowbreak{}1)\allowbreak{}。 先构造 exp(\allowbreak{}x)\allowbreak{}-\allowbreak{}1 的系数（常数 0，其余 1/\allowbreak{}n!），求一次多项式 exp，再乘 n! 取 B[\allowbreak{}n]\allowbreak{}。 需 n<\allowbreak{}模数、阶乘可逆；普通生成函数的系数不能直接当 B[\allowbreak{}n]\allowbreak{}。}
\begin{lstlisting}
const int P=998244353;
\end{lstlisting}
\note{O(\allowbreak{}log b)\allowbreak{}，计算 a\textasciicircum{}b mod P；b 是非负指数，P 为固定素数。}
\begin{lstlisting}
int qpow(int a,int b)
{
    int ans=1;
    for(;b;b>>=1,a=(ll)a*a%P)if(b&1)ans=(ll)ans*a%P;
    return ans;
}
\end{lstlisting}
\note{O(\allowbreak{}L log L)\allowbreak{}，原地变换 a[\allowbreak{}0.\allowbreak{}.\allowbreak{}L-\allowbreak{}1]\allowbreak{}；L 是 2 的幂，inv=\allowbreak{}0 正变换，非零为逆变换。}
\begin{lstlisting}
void ntt(vector<int> &a,int inv)
{
    int n=a.size();
    for(int i=1,j=0;i<n;i++)
    {
        int bit=n>>1;
        for(;j&bit;bit>>=1)j^=bit;
        j^=bit;
        if(i<j)swap(a[i],a[j]);
    }
    for(int len=2;len<=n;len<<=1)
    {
        int wlen=qpow(3,(P-1)/len);
        if(inv)wlen=qpow(wlen,P-2);
        for(int i=0;i<n;i+=len)
        {
            int w=1;
            for(int j=0;j<len/2;j++,w=(ll)w*wlen%P)
            {
                int u=a[i+j],v=(ll)a[i+j+len/2]*w%P;
                a[i+j]=(u+v)%P,a[i+j+len/2]=(u-v+P)%P;
            }
        }
    }
    if(inv)
    {
        int v=qpow(n,P-2);
        for(int &x:a)x=(ll)x*v%P;
    }
}
\end{lstlisting}
\note{NTT 长度不超过 2\textasciicircum{}23，输入系数在 [\allowbreak{}0,\allowbreak{}P)\allowbreak{} O(\allowbreak{}L log L)\allowbreak{}，a/\allowbreak{}b 为系数副本，k>\allowbreak{}=\allowbreak{}0 为结果项数；补零后点乘，再截断或补到 k 项。}
\begin{lstlisting}
vector<int> mul(vector<int> a,vector<int> b,int k)
{
    if(a.empty()||b.empty())return vector<int>(k);
    int sz=a.size()+b.size()-1,n=1;
    while(n<sz)n<<=1;
    assert(n<=(1<<23));
    a.resize(n),b.resize(n);
    ntt(a,0),ntt(b,0);
    for(int i=0;i<n;i++)a[i]=(ll)a[i]*b[i]%P;
    ntt(a,1),a.resize(k);
    return a;
}
\end{lstlisting}
\note{牛顿迭代 O(\allowbreak{}n log n)\allowbreak{}，a[\allowbreak{}0]\allowbreak{}!=\allowbreak{}0 O(\allowbreak{}n log n)\allowbreak{}，返回逆级数；b*(\allowbreak{}2-\allowbreak{}a*b)\allowbreak{} 使误差平方，缺失的输入项视为 0。}
\begin{lstlisting}
vector<int> inverse(const vector<int> &a,int n)
{
    assert(!a.empty()&&a[0]);
    vector<int> b(1,qpow(a[0],P-2));
    for(int len=2;(int)b.size()<n;len<<=1)
    {
        int k=min(len,n);
        vector<int> c(a.begin(),a.begin()+min(k,(int)a.size()));
        c=mul(c,b,k);
        for(int &x:c)x=(P-x)%P;
        c[0]=(c[0]+2)%P;
        b=mul(b,c,k);
    }
    return b;
}
\end{lstlisting}
\note{ln(\allowbreak{}a)\allowbreak{}=\allowbreak{}积分(\allowbreak{}a'/\allowbreak{}a)\allowbreak{}，要求 a[\allowbreak{}0]\allowbreak{}=\allowbreak{}1，O(\allowbreak{}n log n)\allowbreak{} O(\allowbreak{}n log n)\allowbreak{}，导数乘逆再积分；常数项固定为 0，iv[\allowbreak{}i]\allowbreak{} 是整数 i 的模逆。}
\begin{lstlisting}
vector<int> logarithm(const vector<int> &a,int n)
{
    assert(a[0]==1);
    vector<int> d(max(0,(int)a.size()-1));
    for(int i=1;i<(int)a.size();i++)d[i-1]=(ll)a[i]*i%P;
    vector<int> c=mul(d,inverse(a,n),n-1),ans(n),iv(n,1);
    for(int i=2;i<n;i++)iv[i]=P-(ll)(P/i)*iv[P%i]%P;
    for(int i=1;i<n;i++)ans[i]=(ll)c[i-1]*iv[i]%P;
    return ans;
}
\end{lstlisting}
\note{牛顿迭代 b*=\allowbreak{}1+\allowbreak{}a-\allowbreak{}ln(\allowbreak{}b)\allowbreak{}，要求 a[\allowbreak{}0]\allowbreak{}=\allowbreak{}0，O(\allowbreak{}n log n)\allowbreak{} O(\allowbreak{}n log n)\allowbreak{}，返回 exp(\allowbreak{}a)\allowbreak{}；b*(\allowbreak{}1+\allowbreak{}a-\allowbreak{}ln(\allowbreak{}b)\allowbreak{})\allowbreak{} 修正高次误差，b[\allowbreak{}0]\allowbreak{} 始终为 1。}
\begin{lstlisting}
vector<int> exponential(const vector<int> &a,int n)
{
    assert(a[0]==0);
    vector<int> b(1,1);
    for(int len=2;(int)b.size()<n;len<<=1)
    {
        int k=min(len,n);
        vector<int> c=logarithm(b,k);
        for(int i=0;i<k;i++)c[i]=((i<(int)a.size()?a[i]:0)-c[i]+P)%P;
        c[0]=(c[0]+1)%P;
        b=mul(b,c,k);
    }
    return b;
}
\end{lstlisting}
\note{以下扩展共用本文件 mul/\allowbreak{}inverse/\allowbreak{}ln/\allowbreak{}exp：固定模 998244353，系数已归一化，n<\allowbreak{}P。}
\begin{lstlisting}
int power_u64(int a,unsigned long long e){int z=1;for(;e;e>>=1,a=(ll)a*a%P)if(e&1)z=(ll)z*a%P;return z;}
\end{lstlisting}
\note{a\textasciicircum{}k mod x\textasciicircum{}n：分离零前缀、首项常数，再 exp(\allowbreak{}k*ln(\allowbreak{}a)\allowbreak{})\allowbreak{}；k=\allowbreak{}0 返回 1，O(\allowbreak{}n log n)\allowbreak{}。}
\begin{lstlisting}
vector<int> polynomial_power(const vector<int> &a,unsigned long long k,int n)
{
    if(n<=0)return {};vector<int> ans(n);if(!k){ans[0]=1;return ans;}
    int lead=0;while(lead<(int)a.size()&&!a[lead])lead++;
    if(lead==(int)a.size()||(lead&&k>=(unsigned long long)(n+lead-1)/lead))return ans;
    int shift=lead*k,len=n-shift;vector<int> b(len);int inv=qpow(a[lead],P-2);
    for(int i=0;i<len&&lead+i<(int)a.size();i++)b[i]=(ll)a[lead+i]*inv%P;
    auto log=logarithm(b,len);for(int &x:log)x=(ll)x*(k%P)%P;auto c=exponential(log,len);int scale=power_u64(a[lead],k);
    for(int i=0;i<len;i++)ans[shift+i]=(ll)c[i]*scale%P;return ans;
}
\end{lstlisting}
\note{Tonelli–Shanks，固定奇素数 P；非剩余返回 -\allowbreak{}1，选择两个根中较小者。}
\begin{lstlisting}
int scalar_sqrt(int a)
{
    if(!a)return 0;if(qpow(a,(P-1)/2)!=1)return -1;
    int q=P-1,m=0;while(!(q&1))q>>=1,m++;int z=2;while(qpow(z,(P-1)/2)==1)z++;
    int c=qpow(z,q),x=qpow(a,(q+1)/2),t=qpow(a,q);
    while(t!=1){int i=0,u=t;while(u!=1)u=(ll)u*u%P,i++;int b=qpow(c,1<<(m-i-1));x=(ll)x*b%P;c=(ll)b*b%P;t=(ll)t*c%P;m=i;}
    return min(x,P-x);
}
\end{lstlisting}
\note{b\ensuremath{^2}=\allowbreak{}a mod x\textasciicircum{}n；非零首项次数必须为偶数且系数为二次剩余。false 无解，O(\allowbreak{}n log n)\allowbreak{}。}
\begin{lstlisting}
bool polynomial_sqrt(const vector<int> &a,int n,vector<int> &ans)
{
    ans.assign(n,0);if(!n)return true;int lead=0;while(lead<min(n,(int)a.size())&&!a[lead])lead++;
    if(lead==min(n,(int)a.size()))return true;if(lead&1)return false;int x=scalar_sqrt(a[lead]);if(x<0)return false;
    int len=n-lead;vector<int> f(len);for(int i=0;i<len&&lead+i<(int)a.size();i++)f[i]=a[lead+i];vector<int>b={x};
    while((int)b.size()<len){int k=min(len,2*(int)b.size());vector<int> c(f.begin(),f.begin()+k);c=mul(c,inverse(b,k),k);b.resize(k);for(int i=0;i<k;i++)b[i]=(ll)(b[i]+c[i])*((P+1)/2)%P;}
    for(int i=0;i<len;i++)ans[lead/2+i]=b[i];return true;
}
void trim_poly(vector<int> &a){while(!a.empty()&&!a.back())a.pop_back();}
\end{lstlisting}
\note{商和余数：a=\allowbreak{}b*q+\allowbreak{}r，deg r<\allowbreak{}deg b；零多项式用空 vector，b 不能为零。O(\allowbreak{}n log n)\allowbreak{}。}
\begin{lstlisting}
pair<vector<int>,vector<int>> polynomial_divmod(vector<int> a,vector<int> b)
{
    trim_poly(a);trim_poly(b);assert(!b.empty());if(a.size()<b.size())return {{},a};
    int k=a.size()-b.size()+1;vector<int> ra=a,rb=b;reverse(ra.begin(),ra.end());reverse(rb.begin(),rb.end());ra.resize(k);rb.resize(min(k,(int)rb.size()));
    auto q=mul(ra,inverse(rb,k),k);reverse(q.begin(),q.end());auto c=mul(b,q,a.size());for(int i=0;i<(int)a.size();i++)a[i]=(a[i]-c[i]+P)%P;
    a.resize(b.size()-1);trim_poly(a);trim_poly(q);return {q,a};
}
\end{lstlisting}
\note{乘积树：多点求值与快速插值 O(\allowbreak{}(\allowbreak{}n+\allowbreak{}m)\allowbreak{}log\ensuremath{^2}(\allowbreak{}n+\allowbreak{}m)\allowbreak{})\allowbreak{}；点数 m，所有点在 [\allowbreak{}0,\allowbreak{}P)\allowbreak{}。 插值要求点两两不同；各叶节点为 (\allowbreak{}x-\allowbreak{}x\_i)\allowbreak{}，用导数求分母，并自底向上合并。}
\begin{lstlisting}
struct PolynomialPoints
{
    int n;vector<int> x;vector<vector<int>> tree;
    PolynomialPoints(vector<int> x):n(x.size()),x(x),tree(max(1,4*n)){if(n)build(1,0,n);}
    void build(int p,int l,int r){if(r-l==1){tree[p]={(P-x[l])%P,1};return;}int m=(l+r)/2;build(p*2,l,m);build(p*2+1,m,r);tree[p]=mul(tree[p*2],tree[p*2+1],r-l+1);}
    void evaluate(int p,int l,int r,const vector<int> &a,vector<int> &ans)const
    {auto rem=polynomial_divmod(a,tree[p]).second;if(r-l==1){ans[l]=rem.empty()?0:rem[0];return;}int m=(l+r)/2;evaluate(p*2,l,m,rem,ans);evaluate(p*2+1,m,r,rem,ans);}
    vector<int> evaluate(const vector<int> &a)const{vector<int>ans(n);if(n)evaluate(1,0,n,a,ans);return ans;}
    vector<int> combine(int p,int l,int r,const vector<int>&w)const
    {if(r-l==1)return {w[l]};int m=(l+r)/2;auto a=mul(combine(p*2,l,m,w),tree[p*2+1],r-l),b=mul(combine(p*2+1,m,r,w),tree[p*2],r-l);for(int i=0;i<r-l;i++)a[i]=(a[i]+b[i])%P;return a;}
    vector<int> interpolate(const vector<int> &y)const
    {assert(y.size()==x.size());if(!n)return {};vector<int>d(n);for(int i=1;i<=n;i++)d[i-1]=(ll)tree[1][i]*i%P;auto w=evaluate(d);for(int i=0;i<n;i++){assert(w[i]);w[i]=(ll)y[i]*qpow(w[i],P-2)%P;}return combine(1,0,n,w);}
};
\end{lstlisting}
\note{a(\allowbreak{}b(\allowbreak{}x)\allowbreak{})\allowbreak{} mod x\textasciicircum{}n，b[\allowbreak{}0]\allowbreak{}=\allowbreak{}0。分块复合（Brent–Kung 的朴素线性组合版）。 O(\allowbreak{}n\ensuremath{^2} +\allowbreak{} sqrt(\allowbreak{}n)\allowbreak{}*M(\allowbreak{}n)\allowbreak{})\allowbreak{}，M 为卷积成本；比 Horner 的 n 次 NTT 少，但不是最新准线性算法。}
\begin{lstlisting}
vector<int> polynomial_compose(vector<int> a,vector<int> b,int n)
{
    if(!n)return {};assert(b.empty()||b[0]==0);a.resize(n);b.resize(n);int block=max(1,(int)sqrt(n));
    vector<vector<int>> baby(block+1,vector<int>(n));baby[0][0]=1;for(int j=1;j<=block;j++)baby[j]=mul(baby[j-1],b,n);
    vector<int> ans(n);
    for(int l=(n-1)/block*block;l>=0;l-=block){ans=mul(ans,baby[block],n);for(int j=0;j<block&&l+j<n;j++)if(a[l+j])for(int i=0;i<n;i++)ans[i]=(ans[i]+(ll)a[l+j]*baby[j][i])%P;}
    return ans;
}
\end{lstlisting}
% SOURCE: 05-数学/任意模卷积.cpp
\subsection{任意模卷积}

\note{a/\allowbreak{}b为两个多项式系数向量，mod为目标模数；多个NTT素数分别求卷积，再CRT还原。可恢复范围取决于辅助模数乘积与系数和界，任意目标模数不等于无限数值范围。}
\note{任意正模 mod<\allowbreak{}=\allowbreak{}2\textasciicircum{}31-\allowbreak{}1 的整数卷积；负系数先归一化，多模 NTT +\allowbreak{} CRT 精确重构。 O(\allowbreak{}L log L)\allowbreak{}，L 为补齐的 2 次幂且 <\allowbreak{}=\allowbreak{}2\textasciicircum{}24；三模积超过本范围内整数卷积系数上界。 返回 mod 下的系数，mod=\allowbreak{}1 全为 0，空输入返回空；内存为 O(\allowbreak{}L)\allowbreak{}。}
\begin{lstlisting}
ll pow_mod(ll a,ll b,ll p){ll z=1;for(;b;b>>=1,a=a*a%p)if(b&1)z=z*a%p;return z;}
vector<int> convolution_prime(const vector<ll> &a,const vector<ll> &b,int p,int root,int len)
{
    vector<ll> x(len),y(len);for(int i=0;i<(int)a.size();i++)x[i]=a[i]%p;for(int i=0;i<(int)b.size();i++)y[i]=b[i]%p;
    auto ntt=[&](vector<ll> &f,bool inv)
    {
        for(int i=1,j=0;i<len;i++){int bit=len>>1;for(;j&bit;bit>>=1)j^=bit;j^=bit;if(i<j)swap(f[i],f[j]);}
        for(int l=2;l<=len;l*=2){ll z=pow_mod(root,(p-1)/l,p);if(inv)z=pow_mod(z,p-2,p);for(int i=0;i<len;i+=l){ll w=1;for(int j=0;j<l/2;j++,w=w*z%p){ll u=f[i+j],v=f[i+j+l/2]*w%p;f[i+j]=(u+v)%p;f[i+j+l/2]=(u-v+p)%p;}}}
        if(inv){ll z=pow_mod(len,p-2,p);for(ll &v:f)v=v*z%p;}
    };
    ntt(x,false);ntt(y,false);for(int i=0;i<len;i++)x[i]=x[i]*y[i]%p;ntt(x,true);
    return vector<int>(x.begin(),x.begin()+a.size()+b.size()-1);
}
vector<int> convolution_mod(vector<ll> a,vector<ll> b,int mod)
{
    assert(mod>0);if(a.empty()||b.empty())return {};
    int size=a.size()+b.size()-1,len=1;while(len<size)len*=2;assert(len<=(1<<24));
    for(ll &x:a)x=(x%mod+mod)%mod;for(ll &x:b)x=(x%mod+mod)%mod;
    const ll p=167772161,q=469762049,r=754974721;
    auto x=convolution_prime(a,b,p,3,len),y=convolution_prime(a,b,q,3,len),z=convolution_prime(a,b,r,11,len);
    ll ip=pow_mod(p,q-2,q),ipq=pow_mod(p*q%r,r-2,r);vector<int> ans(size);
    for(int i=0;i<size;i++)
    {
        ll t=(y[i]-x[i]+q)%q*ip%q;ll low=x[i]+p*t;
        ll u=(z[i]-low%r+r)%r*ipq%r;
        ans[i]=(low+(__int128)p*q*u)%mod;
    }
    return ans;
}
\end{lstlisting}
% SOURCE: 05-数学/拉格朗日插值.cpp
\subsection{拉格朗日插值}

\note{采样值向量对应声明的横坐标，查询x为目标点；前后缀积省去逐项乘法。等距插值不能替代任意点插值，横坐标差在模意义下必须可逆，采样点上的查询直接返回原值。}
\note{适用：已知连续整数点值，求低次数多项式在新点的值，如幂和求和。 下标：y[\allowbreak{}0.\allowbreak{}.\allowbreak{}n]\allowbreak{}=\allowbreak{}f(\allowbreak{}0.\allowbreak{}.\allowbreak{}n)\allowbreak{}，非空，次数<\allowbreak{}=\allowbreak{}n<\allowbreak{}P；eval 的 a[\allowbreak{}i]\allowbreak{} 是 x\textasciicircum{}i 系数。 参数：x 是求值点，入口按模 P 归一，y 和 a 建议预先归一到 [\allowbreak{}0,\allowbreak{}P)\allowbreak{}。 关键：前后缀积算去掉 i 后的所有 (\allowbreak{}x-\allowbreak{}j)\allowbreak{}，分母为 i!*(\allowbreak{}n-\allowbreak{}i)\allowbreak{}!*(\allowbreak{}-\allowbreak{}1)\allowbreak{}\textasciicircum{}(\allowbreak{}n-\allowbreak{}i)\allowbreak{}。 易错：x 对应已知节点时直接返回 y[\allowbreak{}x]\allowbreak{}，不做除以零；节点必须不同模 P。 复杂度：插值 O(\allowbreak{}n+\allowbreak{}log P)\allowbreak{}，空间 O(\allowbreak{}n)\allowbreak{}，Horner 系数求值 O(\allowbreak{}次数)\allowbreak{}。 参数计数先证明次数：DAG 上每条路径至多 n-\allowbreak{}1 条边，每边至多乘一次参数，则次数<\allowbreak{}=\allowbreak{}n-\allowbreak{}1。 取 n 个不同节点值即可；含参数的分母、任意函数或取模分支不能直接宣称低次多项式。}
\begin{lstlisting}
const int P=998244353;
\end{lstlisting}
\note{O(\allowbreak{}log b)\allowbreak{}，固定模 P 快速幂；a 为归一底数，b>\allowbreak{}=\allowbreak{}0。}
\begin{lstlisting}
ll qpow(ll a,ll b)
{
    ll ans=1;
    for(;b;b>>=1,a=a*a%P)if(b&1)ans=ans*a%P;
    return ans;
}
\end{lstlisting}
\note{已知 y[\allowbreak{}i]\allowbreak{}=\allowbreak{}f(\allowbreak{}i)\allowbreak{}，次数不超过 n；n<\allowbreak{}P，O(\allowbreak{}n)\allowbreak{} 求 f(\allowbreak{}x)\allowbreak{} O(\allowbreak{}n+\allowbreak{}log P)\allowbreak{}，y 是连续节点值的副本、x 是查询点；预处理逆阶乘只求一次幂。}
\begin{lstlisting}
int interpolate(vector<int> y,ll x)
{
    int n=(int)y.size()-1;
    x=(x%P+P)%P;
    if(x<=n)return y[x];
    vector<ll> pre(n+2,1),suf(n+2,1),fac(n+1,1),inv(n+1);
    for(int i=0;i<=n;i++)pre[i+1]=pre[i]*(x-i+P)%P;
    for(int i=n;i>=0;i--)suf[i]=suf[i+1]*(x-i+P)%P;
    for(int i=1;i<=n;i++)fac[i]=fac[i-1]*i%P;
    inv[n]=qpow(fac[n],P-2);
    for(int i=n;i>=1;i--)inv[i-1]=inv[i]*i%P;
    ll ans=0;
    for(int i=0;i<=n;i++)
    {
        ll v=pre[i]*suf[i+1]%P*inv[i]%P*inv[n-i]%P*y[i]%P;
        if((n-i)&1)ans=(ans-v+P)%P;
        else ans=(ans+v)%P;
    }
    return ans;
}
\end{lstlisting}
\note{O(\allowbreak{}a.\allowbreak{}size(\allowbreak{})\allowbreak{})\allowbreak{}，按系数向量求 f(\allowbreak{}x)\allowbreak{}，倒序 Horner 避免逐项求幂；空向量返回 0。}
\begin{lstlisting}
int eval(vector<int> a,ll x)
{
    x=(x%P+P)%P;
    ll ans=0;
    for(int i=(int)a.size()-1;i>=0;i--)ans=(ans*x+a[i])%P;
    return ans;
}
\end{lstlisting}
% SOURCE: 05-数学/BM与快速线性递推.cpp
\subsection{BM与快速线性递推}

\note{输入向量是模意义序列或递推系数，系数第i项乘前i+\allowbreak{}1项；BM求最短递推，查询另需足够初值。不同多项式/\allowbreak{}初值是独立对象，保留运算参数；模数须满足逆元条件。}
\begin{lstlisting}
const ll MOD=998244353;
ll mod_norm(ll x){x%=MOD;return x<0?x+MOD:x;}
ll qpow(ll a,ll b)
{
    ll r=1;
    for(;b;b>>=1,a=a*a%MOD)if(b&1)r=r*a%MOD;
    return r;
}
\end{lstlisting}
\note{一、BM：返回 c，使 f[\allowbreak{}n]\allowbreak{}=\allowbreak{}c[\allowbreak{}0]\allowbreak{}f[\allowbreak{}n-\allowbreak{}1]\allowbreak{}+\allowbreak{}.\allowbreak{}.\allowbreak{}.\allowbreak{}+\allowbreak{}c[\allowbreak{}k-\allowbreak{}1]\allowbreak{}f[\allowbreak{}n-\allowbreak{}k]\allowbreak{}，O(\allowbreak{}s*k)\allowbreak{}，最坏 O(\allowbreak{}s\ensuremath{^2})\allowbreak{}。 输入为连续前 s 项，固定素数模数；全零返回空。仅证明有限前缀的递推关系。 若无限序列最短阶数<\allowbreak{}=\allowbreak{}K，前 2K 项足以恢复；不能凭任意短前缀保证后续正确。}
\begin{lstlisting}
vector<ll> berlekamp_massey(vector<ll> a)
{
    for(ll &x:a)x=mod_norm(x);
    vector<ll> c{1},b{1};
    int k=0,m=1;
    ll last=1;
    for(int n=0;n<(int)a.size();n++)
    {
        ll d=a[n];
        for(int i=1;i<=k;i++)d=(d+c[i]*a[n-i])%MOD;
        if(!d){++m;continue;}
        auto old=c;
        ll coef=d*qpow(last,MOD-2)%MOD;
        c.resize(max(c.size(),b.size()+m));
        for(int i=0;i<(int)b.size();i++)c[i+m]=mod_norm(c[i+m]-coef*b[i]%MOD);
        if(2*k<=n)k=n+1-k,b=move(old),last=d,m=1;
        else ++m;
    }
    c.resize(k+1),c.erase(c.begin());
    for(ll &x:c)x=mod_norm(-x);
    return c;
}
\end{lstlisting}
\note{二、已知递推求第 n 项：x\textasciicircum{}n mod (\allowbreak{}x\textasciicircum{}k-\allowbreak{}c[\allowbreak{}0]\allowbreak{}x\textasciicircum{}(\allowbreak{}k-\allowbreak{}1)\allowbreak{}-\allowbreak{}.\allowbreak{}.\allowbreak{}.\allowbreak{}-\allowbreak{}c[\allowbreak{}k-\allowbreak{}1]\allowbreak{})\allowbreak{}。 O(\allowbreak{}k\ensuremath{^2} log n)\allowbreak{}，空间 O(\allowbreak{}k)\allowbreak{}；init 至少含 f[\allowbreak{}0.\allowbreak{}.\allowbreak{}k-\allowbreak{}1]\allowbreak{}，n 从 0 开始。 c 为空约定为零序列；输入系数/\allowbreak{}初值允许负数。只需本区与 mod\_norm/\allowbreak{}MOD 即可使用。}
\begin{lstlisting}
ll linear_nth(vector<ll> init,vector<ll> c,unsigned long long n)
{
    int k=c.size();
    assert(init.size()>=c.size());
    if(!k)return 0;
    for(ll &x:init)x=mod_norm(x);
    for(ll &x:c)x=mod_norm(x);
    if(n<(unsigned)k)return init[n];
    auto mul=[&](const vector<ll> &a,const vector<ll> &b)
    {
        vector<ll> t(2*k-1);
        for(int i=0;i<k;i++)for(int j=0;j<k;j++)t[i+j]=(t[i+j]+a[i]*b[j])%MOD;
        for(int i=2*k-2;i>=k;i--)
            for(int j=1;j<=k;j++)t[i-j]=(t[i-j]+t[i]*c[j-1])%MOD;
        t.resize(k);return t;
    };
    vector<ll> a(k),b(k);
    a[0]=1;
    if(k==1)b[0]=c[0];else b[1]=1;
    for(;n;n>>=1,b=mul(b,b))if(n&1)a=mul(a,b);
    ll ans=0;
    for(int i=0;i<k;i++)ans=(ans+a[i]*init[i])%MOD;
    return ans;
}
\end{lstlisting}
% SOURCE: 05-数学/卡特兰数.cpp
\subsection{卡特兰数}

\note{fact/\allowbreak{}inv\_fact是阶乘与逆阶乘表，mod为模数；组合公式访问到2n，容量须覆盖它。按模运算公式的素数与可逆条件预处理，不能直接把素数模版本用于任意合数。}
\note{适用：合法括号序列、出栈顺序、不同二叉树形态、不交叉匹配。 参数：n 是成对元素数，Cat[\allowbreak{}0]\allowbreak{}=\allowbreak{}1；阶乘需先预处理到 2*n，且 2*n<\allowbreak{}mod、<\allowbreak{}N。 关键：Cat[\allowbreak{}n]\allowbreak{}=\allowbreak{}C(\allowbreak{}2n,\allowbreak{}n)\allowbreak{}/\allowbreak{}(\allowbreak{}n+\allowbreak{}1)\allowbreak{}；首个匹配分割左右子结构得到递推。 下标：cat\_dp 的 g[\allowbreak{}i]\allowbreak{}[\allowbreak{}j]\allowbreak{} 是剩余未入栈数与栈内数；两者相加不超过初始 n。 易错：mod 必须为素数以求逆；DP 与朴素数组要求 n<\allowbreak{}1005，乘积应在 ll 范围。 复杂度：公式查询 O(\allowbreak{}log mod)\allowbreak{}，双状态 DP 和卷积递推 O(\allowbreak{}n\textasciicircum{}2)\allowbreak{}，空间分别 O(\allowbreak{}n\textasciicircum{}2)\allowbreak{}/\allowbreak{}O(\allowbreak{}n)\allowbreak{}。}
\begin{lstlisting}
const int N=1000006;
ll fact[N],inv_fact[N];
ll mod=1000000007;
\end{lstlisting}
\note{O(\allowbreak{}log n)\allowbreak{}，快速幂 O(\allowbreak{}log n)\allowbreak{}，计算 a 的非负 n 次幂；mod 为正模数。}
\begin{lstlisting}
ll qpow(ll a,ll n,ll mod)
{
    ll res=1;
    a%=mod;
    while(n)
    {
        if(n&1)(res*=a)%=mod;
        (a*=a)%=mod;
        n>>=1;
    }
    return res;
}
\end{lstlisting}
\note{O(\allowbreak{}N)\allowbreak{}，阶乘与阶乘逆元，卡特兰数组合公式要用 O(\allowbreak{}n+\allowbreak{}log p)\allowbreak{}，设置全局 mod=\allowbreak{}p，预处理阶乘 0.\allowbreak{}.\allowbreak{}n；p 为素数且 n<\allowbreak{}p。}
\begin{lstlisting}
void C_init(int n,ll p)
{
    mod=p;
    fact[0]=1;
    for(int i=1;i<=n;i++)fact[i]=fact[i-1]*i%p;
    inv_fact[n]=qpow(fact[n],p-2,p);
    for(int i=n;i>=1;i--)inv_fact[i-1]=inv_fact[i]*i%p;
}
\end{lstlisting}
\note{O(\allowbreak{}1)\allowbreak{}，要求 0<\allowbreak{}=\allowbreak{}n<\allowbreak{}p O(\allowbreak{}1)\allowbreak{}，求模当前 mod 的 C(\allowbreak{}n,\allowbreak{}m)\allowbreak{}，n 是阶乘下标而非卡特兰参数。}
\begin{lstlisting}
ll C(ll n,ll m)
{
    if(m<0||m>n||n<0)return 0;
    return fact[n]*inv_fact[m]%mod*inv_fact[n-m]%mod;
}
\end{lstlisting}
\note{O(\allowbreak{}n\textasciicircum{}2)\allowbreak{}，出栈序列计数：g[\allowbreak{}i]\allowbreak{}[\allowbreak{}j]\allowbreak{} 表示还有 i 个待入栈、栈内 j 个的方案数 每步可以入栈（i>\allowbreak{}0 -\allowbreak{}>\allowbreak{} (\allowbreak{}i-\allowbreak{}1,\allowbreak{}j+\allowbreak{}1)\allowbreak{}）或出栈（j>\allowbreak{}0 -\allowbreak{}>\allowbreak{} (\allowbreak{}i,\allowbreak{}j-\allowbreak{}1)\allowbreak{}），最终停在 g[\allowbreak{}0]\allowbreak{}[\allowbreak{}0]\allowbreak{} 两个转移的 j 都变小，所以同一行 j 必须从大到小扫 O(\allowbreak{}n\textasciicircum{}2)\allowbreak{}，g[\allowbreak{}n]\allowbreak{}[\allowbreak{}0]\allowbreak{} 起步；i 递减接收入栈转移，固定行 j 递减接收出栈转移。}
\begin{lstlisting}
ll cat_dp(int n)
{
    static ll g[1005][1005];
    for(int i=0;i<=n;i++)
        for(int j=0;j<=n;j++)g[i][j]=0;
    g[n][0]=1;
    for(int i=n;i>=0;i--)
        for(int j=n;j>=0;j--)
        {
            if(!g[i][j])continue;
            if(i>0)g[i-1][j+1]=(g[i-1][j+1]+g[i][j])%mod;// 入栈
            if(j>0)g[i][j-1]=(g[i][j-1]+g[i][j])%mod;// 出栈
        }
    return g[0][0];
}
\end{lstlisting}
\note{O(\allowbreak{}log n)\allowbreak{}，Cat[\allowbreak{}n]\allowbreak{} 的组合公式：C(\allowbreak{}2n,\allowbreak{}n)\allowbreak{}/\allowbreak{}(\allowbreak{}n+\allowbreak{}1)\allowbreak{}，即 C(\allowbreak{}2n,\allowbreak{}n)\allowbreak{}*inv(\allowbreak{}n+\allowbreak{}1)\allowbreak{} O(\allowbreak{}log mod)\allowbreak{}，返回 Cat[\allowbreak{}n]\allowbreak{}；n+\allowbreak{}1 必须可逆且 C 表已覆盖 2*n。}
\begin{lstlisting}
ll cat_rec(int n)
{
    if(n==0)return 1;
    return C(2*n,n)*qpow(n+1,mod-2,mod)%mod;
}
\end{lstlisting}
% SOURCE: 05-数学/Stirling数.cpp
\subsection{Stirling数}

\note{n为上界，mod>\allowbreak{}0，first\_kind选择无符号第一类/\allowbreak{}第二类；返回s[\allowbreak{}i]\allowbreak{}[\allowbreak{}k]\allowbreak{}三角表，s[\allowbreak{}0]\allowbreak{}[\allowbreak{}0]\allowbreak{}=\allowbreak{}1。第一类数排列环，第二类数非空集合块，不能把两种递推系数混用。}
\note{O(\allowbreak{}n\ensuremath{^2})\allowbreak{} 两类 Stirling 三角：第一类无符号 c(\allowbreak{}n,\allowbreak{}k)\allowbreak{} 排列恰有 k 个环；第二类 S(\allowbreak{}n,\allowbreak{}k)\allowbreak{} 集合分成 k 个非空块。 若只求一行，可滚动压成 O(\allowbreak{}n)\allowbreak{} 空间；有符号第一类乘 (\allowbreak{}-\allowbreak{}1)\allowbreak{}\textasciicircum{}(\allowbreak{}n-\allowbreak{}k)\allowbreak{}。模 mod>\allowbreak{}0。}
\begin{lstlisting}
vector<vector<long long>> stirling(int n,long long mod,bool first_kind)
{
    assert(n>=0&&mod>0);vector<vector<long long>> s(n+1,vector<long long>(n+1));s[0][0]=1%mod;
    for(int i=1;i<=n;i++)for(int k=1;k<=i;k++)s[i][k]=(s[i-1][k-1]+(__int128)(first_kind?i-1:k)*s[i-1][k])%mod;
    return s;
}
\end{lstlisting}
% SOURCE: 05-数学/行列式与矩阵树定理.cpp
\subsection{行列式与矩阵树定理}

\note{n/\allowbreak{}m为图规模，mod为模数；矩阵树输入是拉普拉斯矩阵，删去同一行列求余子式行列式。矩阵参数保留以区分原矩阵和余子式，消元过程会改工作副本；有向入树/\allowbreak{}出树的度数方向必须一致。}
\note{适用：无向多重图生成树计数；拉普拉斯任意 n-\allowbreak{}1 阶主子式相等。 下标：矩阵 a/\allowbreak{}g 使用 1.\allowbreak{}.\allowbreak{}n，n<\allowbreak{}N；g[\allowbreak{}u]\allowbreak{}[\allowbreak{}v]\allowbreak{} 是边的重数，须对称。 参数：det\_mod 的 p 为素数，matrix\_tree 使用全局 mod；无向自环不计。 关键：度数放对角、邻接重数取负，删去同一编号行列；换行必须翻转行列式符号。 易错：det\_mod 原地消元，保留原矩阵要先复制；精确 det\_naive 和返回值会受 ll 范围限制。 复杂度：模行列式 O(\allowbreak{}n\textasciicircum{}3+\allowbreak{}n log p)\allowbreak{}，空间 O(\allowbreak{}n\textasciicircum{}2)\allowbreak{}；暴力仅供极小图核对。 秩一更新：det(\allowbreak{}I+\allowbreak{}u*v\textasciicircum{}T)\allowbreak{}=\allowbreak{}1+\allowbreak{}v\textasciicircum{}T*u；A 可逆时 det(\allowbreak{}A+\allowbreak{}u*v\textasciicircum{}T)\allowbreak{}=\allowbreak{}det(\allowbreak{}A)\allowbreak{}*(\allowbreak{}1+\allowbreak{}v\textasciicircum{}T*A\textasciicircum{}-\allowbreak{}1*u)\allowbreak{}。 已知精确 |det|=\allowbreak{}D 只判正负：选不整除 D 的奇素数 p，模行列式比较 ±D；D=\allowbreak{}0 单独处理。 必须是素数才可复用模高斯消元；仅“奇模数”不足以保证主元可逆。}
\begin{lstlisting}
typedef __int128 lll;
const int N=105;
const int E=1005;// 边的上限
int n,m;
ll mod=1000000007;
\end{lstlisting}
\note{O(\allowbreak{}log p)\allowbreak{}，快速幂 O(\allowbreak{}log n)\allowbreak{}，a 为底数、n>\allowbreak{}=\allowbreak{}0 为指数，p 为正模数，模乘先扩宽。}
\begin{lstlisting}
ll qpow(ll a,ll n,ll p)
{
    ll res=1%p;
    a%=p;
    while(n)
    {
        if(n&1)res=(ll)((lll)res*a%p);
        a=(ll)((lll)a*a%p);
        n>>=1;
    }
    return res;
}
\end{lstlisting}
\note{O(\allowbreak{}n\textasciicircum{}3)\allowbreak{}，模素数意义下高斯消元求行列式（不取模则要求 mod 为素数以便求逆） O(\allowbreak{}n\textasciicircum{}3+\allowbreak{}n log p)\allowbreak{}，a 是 n 阶方阵，p 为素数；找不到非零主元返回 0。}
\begin{lstlisting}
ll det_mod(ll a[N][N],int n,ll p)
{
    ll res=1;
    for(int c=1;c<=n;c++)
    {
        int pivot=-1;
        for(int i=c;i<=n;i++)
            if(a[i][c]%p){pivot=i;break;}
        if(pivot==-1)return 0;// 有一列全 0，行列式为 0
        if(pivot!=c)
        {
            for(int j=1;j<=n;j++)swap(a[c][j],a[pivot][j]);
            res=(p-res)%p;// 交换两行，行列式变号
        }
        res=(ll)((lll)res*a[c][c])%p;
        ll inv=qpow(a[c][c],p-2,p);// 主元的逆元
        for(int i=c+1;i<=n;i++)
        {
            if(!a[i][c])continue;
            ll t=(ll)((lll)a[i][c]*inv%p);
            for(int j=c;j<=n;j++)a[i][j]=(a[i][j]-(lll)t*a[c][j])%p;
        }
    }
    return (res%p<0?res%p+p:res%p);
}
\end{lstlisting}
\note{O(\allowbreak{}n! )\allowbreak{}，按第一行展开求行列式，只用于小矩阵对拍 O(\allowbreak{}n!)\allowbreak{} 递归展开，n>\allowbreak{}=\allowbreak{}1；临时矩阵在递归栈上，只用于小阶数。}
\note{O(\allowbreak{}n\textasciicircum{}2 +\allowbreak{} n\textasciicircum{}3)\allowbreak{}，矩阵树定理：邻接矩阵直接建拉普拉斯矩阵，求任意 n-\allowbreak{}1 阶主子式 返回 n 个点的无向（允许重边、自环不加）生成树个数，边权全为 1 O(\allowbreak{}n\textasciicircum{}3+\allowbreak{}n log mod)\allowbreak{}，g 为对称重数矩阵；返回模 mod 的生成树数，单点返回 1。}
\begin{lstlisting}
ll matrix_tree(int n,ll g[N][N])
{
    if(n==1)return 1;// 单点视为 1 棵树
    static ll lap[N][N],sub[N][N];
    for(int i=1;i<=n;i++)
        for(int j=1;j<=n;j++)lap[i][j]=0;
    for(int i=1;i<=n;i++)
        for(int j=1;j<=n;j++)
            if(i!=j&&g[i][j])
            {
                lap[i][j]-=g[i][j];// 邻接矩阵取负
                lap[i][i]+=g[i][j];// 度数放在对角
            }
    for(int i=1;i<=n-1;i++)
        for(int j=1;j<=n-1;j++)sub[i][j]=((lap[i][j]%mod)+mod)%mod;// 去掉第 n 行第 n 列
    return det_mod(sub,n-1,mod);
}
\end{lstlisting}
% SOURCE: 05-数学/矩阵求逆与多右端项.cpp
\subsection{矩阵求逆与多右端项}

\note{A为n\ensuremath{\times}n系数，B为n\ensuremath{\times}r右端项，编号0起；成功时B改成X，A以副本消元。求逆令B为单位矩阵，false表示奇异或接近奇异；模版p须素数，实数版EPS按尺度调整。}
\note{一次 Gauss–Jordan 解 AX=\allowbreak{}B（A 为 n*n，B 为 n*r），O(\allowbreak{}n\ensuremath{^3}+\allowbreak{}n\ensuremath{^2}r)\allowbreak{}，原地把 B 变成 X。 仅处理 A 可逆的情况；false 表示奇异/\allowbreak{}接近奇异。求逆令 B=\allowbreak{}I 即可，不重跑 n 次消元。 实数版本 eps 按尺度调整；模版本 p 必须素数，p>\allowbreak{}=\allowbreak{}2，允许负输入。编号从 0 开始。}
\begin{lstlisting}
bool solve_matrix(vector<vector<long double>> a,vector<vector<long double>> &b,long double eps=1e-12L)
{
    int n=a.size(),r=n?b[0].size():0;assert((int)b.size()==n);
    for(int c=0;c<n;c++)
    {
        int k=c;for(int i=c+1;i<n;i++)if(fabsl(a[i][c])>fabsl(a[k][c]))k=i;
        if(fabsl(a[k][c])<=eps)return false;swap(a[k],a[c]);swap(b[k],b[c]);
        long double z=a[c][c];for(auto &x:a[c])x/=z;for(auto &x:b[c])x/=z;
        for(int i=0;i<n;i++)if(i!=c){z=a[i][c];for(int j=c;j<n;j++)a[i][j]-=z*a[c][j];for(int j=0;j<r;j++)b[i][j]-=z*b[c][j];}
    }
    return true;
}
bool solve_matrix_mod(vector<vector<long long>> a,vector<vector<long long>> &b,long long p)
{
    using ll=long long;int n=a.size(),r=n?b[0].size():0;assert((int)b.size()==n&&p>=2);
    for(auto &v:a)for(ll &x:v)x=(x%p+p)%p;for(auto &v:b)for(ll &x:v)x=(x%p+p)%p;
    auto power=[&](ll x,ll e){ll z=1;for(;e;e>>=1,x=(__int128)x*x%p)if(e&1)z=(__int128)z*x%p;return z;};
    for(int c=0;c<n;c++)
    {
        int k=c;while(k<n&&!a[k][c])k++;if(k==n)return false;swap(a[k],a[c]);swap(b[k],b[c]);
        ll z=power(a[c][c],p-2);for(ll &x:a[c])x=(__int128)x*z%p;for(ll &x:b[c])x=(__int128)x*z%p;
        for(int i=0;i<n;i++)if(i!=c){z=a[i][c];for(int j=c;j<n;j++)a[i][j]=((a[i][j]-(__int128)z*a[c][j])%p+p)%p;for(int j=0;j<r;j++)b[i][j]=((b[i][j]-(__int128)z*b[c][j])%p+p)%p;}
    }
    return true;
}
\end{lstlisting}
% SOURCE: 05-数学/杜教筛.cpp
\subsection{杜教筛}

\note{M为小范围预筛界，sum\_phi/\allowbreak{}sum\_mu为其前缀和，mp\_phi/\allowbreak{}mp\_mu缓存大参数递归结果。prime/\allowbreak{}vis/\allowbreak{}phi/\allowbreak{}mu供预筛；换预筛或函数定义后清缓存，整除分块右端点用n/\allowbreak{}(\allowbreak{}n/\allowbreak{}l)\allowbreak{}计算。}
\note{适用：大 n 的 phi/\allowbreak{}mu 前缀和，筛不完所有数时用卷积恒等式递归。 参数：get\_pre 的 n 是筛上界 M，1<\allowbreak{}=\allowbreak{}M<\allowbreak{}N；du\_phi/\allowbreak{}du\_mu 的 n 是查询上界。 关键：整除商 n/\allowbreak{}l 相同的闭区间为 [\allowbreak{}l,\allowbreak{}n/\allowbreak{}(\allowbreak{}n/\allowbreak{}l)\allowbreak{}]\allowbreak{}，一次合并整个区间。 结论：sum\_phi[\allowbreak{}n]\allowbreak{} 为 1.\allowbreak{}.\allowbreak{}n 的 phi 和，sum\_mu[\allowbreak{}n]\allowbreak{} 为 Mertens 函数；n=\allowbreak{}0 返回前缀 0。 易错：多组重筛前清 vis、前缀和及缓存；n*(\allowbreak{}n+\allowbreak{}1)\allowbreak{}/\allowbreak{}2 和块乘积仍可能溢出 ll。 复杂度：M 取查询上界的约 2/\allowbreak{}3 次幂时总时间 O(\allowbreak{}n\textasciicircum{}(\allowbreak{}2/\allowbreak{}3)\allowbreak{})\allowbreak{}；筛空间 O(\allowbreak{}M)\allowbreak{}，缓存商值。}
\begin{lstlisting}
const int N=5000006;// 预处理上界，一般取 n^(2/3) 量级
int M;
int prime[N],cnt;// 素数表
bool vis[N];
int phi[N];// 欧拉函数
int mu[N];// 莫比乌斯函数
ll sum_phi[N];// phi 前缀和
ll sum_mu[N];// mu 前缀和
unordered_map<ll,ll> mp_phi,mp_mu;// 记忆化
\end{lstlisting}
\note{O(\allowbreak{}n)\allowbreak{}，线性筛 phi /\allowbreak{} mu 并求前缀和 O(\allowbreak{}n)\allowbreak{}，预处理 1.\allowbreak{}.\allowbreak{}n，M 记录可直接查表范围；先筛函数值再累加前缀。}
\begin{lstlisting}
void get_pre(int n)
{
    M=n,cnt=0;
    fill(vis,vis+n+1,false);
    mp_phi.clear(),mp_mu.clear(); sum_phi[0]=sum_mu[0]=0;
    vis[1]=true;
    phi[1]=1,mu[1]=1;
    for(int i=2;i<=n;i++)
    {
        if(!vis[i])
        {
            prime[++cnt]=i;
            phi[i]=i-1;// 素数的 phi
            mu[i]=-1;
        }
        for(int j=1;j<=cnt&&(ll)i*prime[j]<=n;j++)
        {
            int t=i*prime[j];
            vis[t]=true;
            if(i%prime[j]==0)
            {
                phi[t]=phi[i]*prime[j];
                mu[t]=0;
                break;// 每个合数只被最小质因子筛一次
            }
            phi[t]=phi[i]*(prime[j]-1);
            mu[t]=-mu[i];
        }
    }
    for(int i=1;i<=n;i++)
    {
        sum_phi[i]=sum_phi[i-1]+phi[i];
        sum_mu[i]=sum_mu[i-1]+mu[i];
    }
}
\end{lstlisting}
\note{O(\allowbreak{}n\textasciicircum{}(\allowbreak{}2/\allowbreak{}3)\allowbreak{})\allowbreak{}，杜教筛求 \ensuremath{\Sigma}\_\{i=\allowbreak{}1\}\textasciicircum{}\{n\} phi(\allowbreak{}i)\allowbreak{} 用恒等式 \ensuremath{\Sigma}\_\{i=\allowbreak{}1\}\textasciicircum{}\{n\} phi(\allowbreak{}i)\allowbreak{} =\allowbreak{} n*(\allowbreak{}n+\allowbreak{}1)\allowbreak{}/\allowbreak{}2 -\allowbreak{} \ensuremath{\Sigma}\_\{l=\allowbreak{}2\}\textasciicircum{}\{n\} (\allowbreak{}r-\allowbreak{}l+\allowbreak{}1)\allowbreak{}*S(\allowbreak{}n/\allowbreak{}l)\allowbreak{} 适当预筛时 O(\allowbreak{}n\textasciicircum{}(\allowbreak{}2/\allowbreak{}3)\allowbreak{})\allowbreak{}，缓存同一 n；从 l=\allowbreak{}2 起排除待求的自身项。}
\begin{lstlisting}
ll du_phi(ll n)
{
    if(n<=M)return sum_phi[n];
    if(mp_phi.count(n))return mp_phi[n];
    __int128 res=(__int128)n*(n+1)/2;
    for(ll l=2,r;l<=n;l=r+1)
    {
        r=n/(n/l);// 整除分块
        res-=(__int128)(r-l+1)*du_phi(n/l);
    }
    return mp_phi[n]=res;
}
\end{lstlisting}
\note{O(\allowbreak{}n\textasciicircum{}(\allowbreak{}2/\allowbreak{}3)\allowbreak{})\allowbreak{}，杜教筛求 \ensuremath{\Sigma}\_\{i=\allowbreak{}1\}\textasciicircum{}\{n\} mu(\allowbreak{}i)\allowbreak{} 用恒等式 1 =\allowbreak{} \ensuremath{\Sigma}\_\{l=\allowbreak{}1\}\textasciicircum{}\{n\} (\allowbreak{}r-\allowbreak{}l+\allowbreak{}1)\allowbreak{}*S(\allowbreak{}n/\allowbreak{}l)\allowbreak{}，即 S(\allowbreak{}n)\allowbreak{} =\allowbreak{} 1 -\allowbreak{} \ensuremath{\Sigma}\_\{l=\allowbreak{}2\}\textasciicircum{}\{n\} (\allowbreak{}r-\allowbreak{}l+\allowbreak{}1)\allowbreak{}*S(\allowbreak{}n/\allowbreak{}l)\allowbreak{} 适当预筛时 O(\allowbreak{}n\textasciicircum{}(\allowbreak{}2/\allowbreak{}3)\allowbreak{})\allowbreak{}，利用 mu*1 的前缀恒为 1，递归参数 n/\allowbreak{}l 严格减小。}
\begin{lstlisting}
ll du_mu(ll n)
{
    if(n<=M)return sum_mu[n];
    if(mp_mu.count(n))return mp_mu[n];
    ll res=1;
    for(ll l=2,r;l<=n;l=r+1)
    {
        r=n/(n/l);
        res-=(r-l+1)*du_mu(n/l);
    }
    return mp_mu[n]=res;
}
\end{lstlisting}
% SOURCE: 05-数学/Min25筛.cpp
\subsection{Min25筛}

\note{coeff=\allowbreak{}\{c0,\allowbreak{}c1,\allowbreak{}c2\}只定义素数处f(\allowbreak{}p)\allowbreak{}，prime\_power(\allowbreak{}p,\allowbreak{}e,\allowbreak{}p\textasciicircum{}e)\allowbreak{}定义素数幂值；两者须属于同一个积性函数。w为不同整除商，small/\allowbreak{}large映射下标，g筛素数幂贡献；返回模意义前缀和。}
\note{Min\_25：积性 f，f(\allowbreak{}1)\allowbreak{}=\allowbreak{}1，素数处 f(\allowbreak{}p)\allowbreak{}=\allowbreak{}c0+\allowbreak{}c1*p+\allowbreak{}c2*p\ensuremath{^2} mod 1e9+\allowbreak{}7。 prime\_power(\allowbreak{}p,\allowbreak{}e,\allowbreak{}p\textasciicircum{}e)\allowbreak{} 由题目提供，不是随意对普通整数套同一多项式；先筛素数贡献，再枚举最小素因子。 常见复杂度 O(\allowbreak{}n\textasciicircum{}(\allowbreak{}3/\allowbreak{}4)\allowbreak{}/\allowbreak{}log n)\allowbreak{}，空间 O(\allowbreak{}sqrt n)\allowbreak{}；n>\allowbreak{}=\allowbreak{}0，sqrt(\allowbreak{}n)\allowbreak{} 须能分配，建议先估算性能。 例 f(\allowbreak{}p\textasciicircum{}e)\allowbreak{}=\allowbreak{}p\textasciicircum{}e*(\allowbreak{}p\textasciicircum{}e-\allowbreak{}1)\allowbreak{}，coeff=\allowbreak{}\{0,\allowbreak{}-\allowbreak{}1,\allowbreak{}1\}。phi/\allowbreak{}mu 的单纯前缀仍优先杜教筛。 需要精确 pi(\allowbreak{}floor(\allowbreak{}n/\allowbreak{}i)\allowbreak{})\allowbreak{} 时用“整除商上的质数计数.\allowbreak{}cpp”；本接口返回模意义积性函数和。}
\begin{lstlisting}
template<class PrimePower>
ll min25(ll n,array<ll,3> coeff,PrimePower prime_power)
{
    const ll mod=1000000007,iv2=500000004,iv6=166666668;if(n<1)return 0;
    ll rt=sqrtl(n);while((rt+1)<=n/(rt+1))rt++;while(rt>n/rt)rt--;
    assert(rt<=INT_MAX-1);int r=rt;vector<int> prime={0};vector<bool> composite(r+1);
    for(int i=2;i<=r;i++){if(!composite[i])prime.push_back(i);for(int j=1;j<(int)prime.size()&&(ll)i*prime[j]<=r;j++){composite[i*prime[j]]=true;if(i%prime[j]==0)break;}}
    vector<array<ll,3>> pref(prime.size());for(int i=1;i<(int)prime.size();i++){ll p=prime[i];pref[i]={ (pref[i-1][0]+1)%mod,(pref[i-1][1]+p)%mod,(pref[i-1][2]+p*p)%mod};}
    vector<ll>w;vector<array<ll,3>>g;vector<int> small(r+1),large(r+1);
    for(ll l=1;l<=n;){ll x=n/l;int id=w.size();w.push_back(x);ll z=x%mod;g.push_back({(z-1+mod)%mod,(z*(z+1)%mod*iv2-1+mod)%mod,(z*(z+1)%mod*(2*z+1)%mod*iv6-1+mod)%mod});if(x<=r)small[x]=id;else large[n/x]=id;l=n/x+1;}
    auto id=[&](ll x){return x<=r?small[x]:large[n/x];};
    for(int i=1;i<(int)prime.size();i++){ll p=prime[i];for(int j=0;j<(int)w.size()&&p<=w[j]/p;j++){int k=id(w[j]/p);ll power=1;for(int d=0;d<3;d++){g[j][d]=(g[j][d]-power*(g[k][d]-pref[i-1][d])%mod+mod)%mod;power=power*p%mod;}}}
    for(ll &c:coeff)c=(c%mod+mod)%mod;
    auto sum=[&](auto&&sum,ll x,int j)->ll
    {
        if(x<=1||(j>0&&prime[j]>=x))return 0;ll ans=0;int k=id(x);
        for(int d=0;d<3;d++)ans=(ans+coeff[d]*(g[k][d]-pref[j][d]+mod))%mod;
        for(int i=j+1;i<(int)prime.size()&&prime[i]<=x/prime[i];i++)
        {ll p=prime[i],pw=p;for(int e=1;;e++){ll f=(prime_power(p,e,pw)%mod+mod)%mod;ans=(ans+f*(sum(sum,x/pw,i)+(e!=1)))%mod;if(pw>x/p)break;pw*=p;}}
        return ans;
    };
    return (sum(sum,n,0)+1)%mod;
}
\end{lstlisting}
% SOURCE: 05-数学/整除商上的质数计数.cpp
\subsection{整除商上的质数计数}

\note{状态只保存不同的floor(\allowbreak{}n/\allowbreak{}i)\allowbreak{}，小商与大商分别映射到紧凑下标；筛后读pi(\allowbreak{}x)\allowbreak{}。n为非负上界，精确计数使用整数，平方根及乘法边界按实现校正。}
\note{Min\_25 的计数阶段：一次预处理，得到全部 v=\allowbreak{}floor(\allowbreak{}n/\allowbreak{}i)\allowbreak{} 处的精确 pi(\allowbreak{}v)\allowbreak{}，不取模。 w 降序列出商值，g[\allowbreak{}j]\allowbreak{}=\allowbreak{}pi(\allowbreak{}w[\allowbreak{}j]\allowbreak{})\allowbreak{}。不能对任意大 x 查询；x<\allowbreak{}=\allowbreak{}sqrt(\allowbreak{}n)\allowbreak{} 均在表内。 常见时间 O(\allowbreak{}n\textasciicircum{}(\allowbreak{}3/\allowbreak{}4)\allowbreak{}/\allowbreak{}log n)\allowbreak{}，空间 O(\allowbreak{}sqrt n)\allowbreak{}；n>\allowbreak{}=\allowbreak{}0，先估算筛规模和运行时间。}
\begin{lstlisting}
struct PrimeCount
{
    ll n,rt;
    vector<ll> w,g;
    vector<int> small,large;
    PrimeCount(ll n):n(n),rt(sqrtl(n))
    {
        assert(n>=0);
        while(rt+1<=n/(rt+1))++rt;
        while(rt&&rt>n/rt)--rt;
        small.resize(rt+1),large.resize(rt+1);
        for(ll l=1;l<=n;)
        {
            ll x=n/l;int j=w.size();w.push_back(x),g.push_back(x-1);
            if(x<=rt)small[x]=j;else large[n/x]=j;
            l=n/x+1;
        }
        for(ll p=2;p<=rt;++p)
        {
            ll before=g[id(p-1)];if(g[id(p)]==before)continue;
            for(int j=0;j<(int)w.size()&&p<=w[j]/p;++j)g[j]-=g[id(w[j]/p)]-before;
        }
    }
    int id(ll x)const
    {
        assert(1<=x&&x<=n);
        int j=x<=rt?small[x]:large[n/x];assert(w[j]==x);return j;
    }
    ll count(ll x)const{return x==0?0:g[id(x)];}
};
\end{lstlisting}
\note{例：PrimeCount pc(\allowbreak{}n)\allowbreak{}; pc.\allowbreak{}count(\allowbreak{}n)\allowbreak{}；整除块 [\allowbreak{}l,\allowbreak{}r]\allowbreak{} 内素数个数为 pi(\allowbreak{}r)\allowbreak{}-\allowbreak{}pi(\allowbreak{}l-\allowbreak{}1)\allowbreak{}。 本模板的块端点 r=\allowbreak{}n/\allowbreak{}(\allowbreak{}n/\allowbreak{}l)\allowbreak{}，故 r 和前一块端点 l-\allowbreak{}1 都能查询；一般端点不保证在表内。}
% SOURCE: 05-数学/合数模组合数exLucas.cpp
\subsection{合数模组合数exLucas}

\note{n、k为组合数参数，mod为任意正模数；按素数幂拆分，分别处理阶乘中该素因子的指数与去因子部分，再CRT合并。缓存表与素数幂绑定，换模数须重建。}
\note{C(\allowbreak{}n,\allowbreak{}k)\allowbreak{} mod m，n,\allowbreak{}k 非负 64 位，1<\allowbreak{}=\allowbreak{}m<\allowbreak{}=\allowbreak{}2\textasciicircum{}31-\allowbreak{}1；分解素数幂 +\allowbreak{} 去 p 因子阶乘 +\allowbreak{} CRT。 每个素数幂 pk 需 O(\allowbreak{}pk)\allowbreak{} 预处理，故要求 pk<\allowbreak{}=\allowbreak{}2e6（可按内存调整）；查询 O(\allowbreak{}log\_p n * log pk)\allowbreak{}。 n 极大、pk 极大时不能直接用本版；普通素数且小 n 优先原组合数模板。}
\begin{lstlisting}
using ll=long long;using U=unsigned long long;
ll exgcd(ll a,ll b,ll &x,ll &y){if(!b){x=1;y=0;return a;}ll z=exgcd(b,a%b,y,x);y-=a/b*x;return z;}
ll inverse(ll a,ll m){ll x,y;assert(exgcd(a,m,x,y)==1);return (x%m+m)%m;}
ll power(ll a,U b,ll m){ll z=1;for(;b;b>>=1,a=a*a%m)if(b&1)z=z*a%m;return z;}
ll binomial_prime_power(U n,U k,int p,int pk)
{
    vector<ll> f(pk+1,1);for(int i=1;i<=pk;i++)f[i]=f[i-1]*(i%p?i:1)%pk;
    auto fact=[&](U x){ll ans=1;while(x){ans=ans*power(f[pk],x/pk,pk)%pk*f[x%pk]%pk;x/=p;}return ans;};
    auto exponent=[&](U x){U z=0;while(x)z+=x/=p;return z;};
    U e=exponent(n)-exponent(k)-exponent(n-k);int q=0;for(int x=pk;x>1;x/=p)q++;
    if(e>=(U)q)return 0;
    return fact(n)*inverse(fact(k),pk)%pk*inverse(fact(n-k),pk)%pk*power(p,e,pk)%pk;
}
ll exlucas(U n,U k,int m)
{
    assert(m>0);if(k>n||m==1)return 0;int rest=m;ll ans=0;
    for(int p=2;(ll)p*p<=rest;p++)if(rest%p==0)
    {int pk=1;while(rest%p==0)rest/=p,pk*=p;assert(pk<=2000000);ll v=binomial_prime_power(n,k,p,pk),M=m/pk;ans=(ans+(__int128)v*M*inverse(M%pk,pk))%m;}
    if(rest>1){assert(rest<=2000000);ll v=binomial_prime_power(n,k,rest,rest),M=m/rest;ans=(ans+(__int128)v*M*inverse(M%rest,rest))%m;}
    return ans;
}
\end{lstlisting}
% SOURCE: 05-数学/原根与阶.cpp
\subsection{原根与阶}

\note{prime/\allowbreak{}cnt保存素数表，spf是最小质因子，n为筛上界；阶与原根需要分解phi(\allowbreak{}模数)\allowbreak{}。先建立足够范围的质因子信息，查询模数的适用性单独判断，存在原根不是所有模数的性质。}
\note{适用：循环群生成元、周期、NTT 单位根；阶是使 a\textasciicircum{}k=\allowbreak{}1 的最小正 k。 参数：模数 p>\allowbreak{}0，阶与原根要求底数互素；get\_root 专用于素数模数。 下标：筛的 prime 从 1 开始，spf[\allowbreak{}x]\allowbreak{} 是最小质因子，筛界小于 N。 关键：阶整除 phi(\allowbreak{}p)\allowbreak{}，逐个试除其素因子；素数原根检查 (\allowbreak{}p-\allowbreak{}1)\allowbreak{}/\allowbreak{}q 次幂不为 1。 易错：试除循环 i*i、余数归一和朴素模乘仍须防溢出；seen 版要求模数<\allowbreak{}N。 复杂度：素因子试除 O(\allowbreak{}sqrt p)\allowbreak{}，每个候选做若干快速幂；枚举最小原根不能保证仅 O(\allowbreak{}log\textasciicircum{}2 p)\allowbreak{}。}
\begin{lstlisting}
typedef __int128 lll;
const int N=1000006;
int n;
int prime[N],cnt;// 素数表（顺便当最小质因子筛用）
int spf[N];// 最小质因子
\end{lstlisting}
\note{O(\allowbreak{}n)\allowbreak{}，线性筛，顺便记最小质因子 O(\allowbreak{}n)\allowbreak{}，n 为筛上界；spf 需初始清零，多组筛要重置表。}
\begin{lstlisting}
void get_prime(int n)
{
    cnt=0;
    fill(spf,spf+n+1,0);
    for(int i=2;i<=n;i++)
    {
        if(!spf[i])prime[++cnt]=i,spf[i]=i;// i 是素数
        for(int j=1;j<=cnt&&(ll)i*prime[j]<=n;j++)
        {
            spf[i*prime[j]]=prime[j];
            if(i%prime[j]==0)break;// 只被最小质因子筛一次
        }
    }
}
\end{lstlisting}
\note{O(\allowbreak{}log n)\allowbreak{}，快速幂（模乘用 \_\_int128 防溢出） O(\allowbreak{}log n)\allowbreak{}，模 mod 幂，n>\allowbreak{}=\allowbreak{}0；a 应为非负归一余数，乘法用 \_\_int128。}
\begin{lstlisting}
ll qpow(ll a,ll n,ll mod)
{
    ll res=1%mod;
    a%=mod;
    while(n)
    {
        if(n&1)res=(ll)((lll)res*a%mod);
        a=(ll)((lll)a*a%mod);
        n>>=1;
    }
    return res;
}
\end{lstlisting}
\note{O(\allowbreak{}sqrt n)\allowbreak{}，试除分解质因数（不要求 n 是素数），返回去重后的素因子 O(\allowbreak{}sqrt n)\allowbreak{}，n>\allowbreak{}=\allowbreak{}1，返回不同素因子，供阶的约简使用。}
\begin{lstlisting}
vector<ll> factor(ll n)
{
    vector<ll> v;
    for(ll i=2;i<=n/i;i++)
        if(n%i==0)
        {
            v.push_back(i);
            while(n%i==0)n/=i;
        }
    if(n>1)v.push_back(n);
    return v;
}
\end{lstlisting}
\note{O(\allowbreak{}sqrt n)\allowbreak{}，单个数的欧拉函数，顺序筛 phi(\allowbreak{}p)\allowbreak{} 时要用（p 不一定是素数） O(\allowbreak{}sqrt n)\allowbreak{}，n>\allowbreak{}=\allowbreak{}1，返回与 n 互素的 1.\allowbreak{}.\allowbreak{}n 整数个数。}
\begin{lstlisting}
ll phi_of(ll n)
{
    ll r=n;
    for(ll i=2;i<=n/i;i++)
        if(n%i==0)
        {
            r=r/i*(i-1);
            while(n%i==0)n/=i;
        }
    if(n>1)r=r/n*(n-1);
    return r;
}
\end{lstlisting}
\note{O(\allowbreak{}sqrt p log p)\allowbreak{}，求 a 模 p 的阶，要求 gcd(\allowbreak{}a,\allowbreak{}p)\allowbreak{}=\allowbreak{}1（p 可以是合数） 做法：ord | phi(\allowbreak{}p)\allowbreak{}，枚举 phi(\allowbreak{}p)\allowbreak{} 的素因子逐个除掉，除掉一个因子后接着试同一个因子 试除加 O(\allowbreak{}log p)\allowbreak{} 次幂测试，a 为底数、p 为模数；不互素返回 -\allowbreak{}1。}
\begin{lstlisting}
ll get_order(ll a,ll p)
{
    if(__gcd(a,p)!=1)return -1;// 不互素没有阶
    if(p==1)return 1;
    ll ph=phi_of(p);
    vector<ll> f=factor(ph);
    ll ord=ph;
    for(int i=0;i<(int)f.size();i++)
        while(ord%f[i]==0&&qpow(a,ord/f[i],p)==1)ord/=f[i];// 能除就除
    return ord;
}
\end{lstlisting}
\note{O(\allowbreak{}sqrt(\allowbreak{}p)\allowbreak{} +\allowbreak{} log\textasciicircum{}2 p)\allowbreak{}，求模 p 的最小原根（p 为素数，p=\allowbreak{}2 特判返回 1） 枚举候选乘每次幂测试成本，p 为素数；p=\allowbreak{}2 的原根为 1。}
\begin{lstlisting}
ll get_root(ll p)
{
    if(p==2)return 1;
    vector<ll> f=factor(p-1);
    for(ll g=2;g<p;g++)
    {
        bool ok=true;
        for(int i=0;i<(int)f.size();i++)
            if(qpow(g,(p-1)/f[i],p)==1)// 只要有一个为 1 就不是原根
            {
                ok=false;
                break;
            }
        if(ok)return g;
    }
    return -1;// p 不是素数时可能走到这里
}
\end{lstlisting}
\note{O(\allowbreak{}sqrt(\allowbreak{}p)\allowbreak{} log p)\allowbreak{}，原根判定：g 是模 p 原根 <\allowbreak{}=\allowbreak{}>\allowbreak{} ord\_p(\allowbreak{}g)\allowbreak{}=\allowbreak{}phi(\allowbreak{}p)\allowbreak{} 对素数 p 就是 ord=\allowbreak{}p-\allowbreak{}1；合数模数下用 phi(\allowbreak{}p)\allowbreak{}，只有 2,\allowbreak{}4,\allowbreak{}p\textasciicircum{}k,\allowbreak{}2p\textasciicircum{}k 才有原根 与求阶同量级，判断 g 的阶是否等于 phi(\allowbreak{}p)\allowbreak{}；接口另约定 p=\allowbreak{}1 返回 true。}
\begin{lstlisting}
bool is_root(ll g,ll p)
{
    if(__gcd(g,p)!=1)return false;
    if(p==1)return true;
    return get_order(g,p)==phi_of(p);
}
\end{lstlisting}
\note{O(\allowbreak{}p sqrt(\allowbreak{}p)\allowbreak{} log p)\allowbreak{}，求任意模数 n 的最小原根（n 为 2,\allowbreak{}4,\allowbreak{}p\textasciicircum{}k,\allowbreak{}2p\textasciicircum{}k 时才有解），无解返回 -\allowbreak{}1 至多 n 个候选，各自求阶；只在有原根的模数上有解，n=\allowbreak{}1 约定返回 0。}
\begin{lstlisting}
ll get_root_general(ll n)
{
    if(n==1)return 0;
    for(ll g=1;g<n;g++)
        if(is_root(g,n))return g;
    return -1;
}
\end{lstlisting}
% SOURCE: 05-数学/二次剩余.cpp
\subsection{二次剩余}

\note{P为当前素数模，W为Cipolla扩域非剩余；扩域乘法同时读取二者，不能在同一计算中改变模数。输入先归一，0单独处理，判无解后再求根；模2与奇素数条件按入口处理。}
\note{适用：求模素数平方根 x\textasciicircum{}2=\allowbreak{}a；有非零解时通常是 x 与 p-\allowbreak{}x 两个根。 参数：p 为素数；a 可负，入口归一化；模 2 特判，零的平方根只有零。 状态：P/\allowbreak{}W 是一次 Cipolla 运算使用的全局模数和非剩余，不能并发复用。 关键：找 b 使 b\textasciicircum{}2-\allowbreak{}a 非剩余，在扩域中求 (\allowbreak{}b+\allowbreak{}sqrt(\allowbreak{}W)\allowbreak{})\allowbreak{}\textasciicircum{}(\allowbreak{}p/\allowbreak{}2+\allowbreak{}1)\allowbreak{}。 易错：\_\_int128 保护乘法，但双乘积相加和 ll 归一化也要在范围内；返回根不保证最小。 复杂度：随机找非剩余期望常数轮，每轮 O(\allowbreak{}log p)\allowbreak{}，加一次扩域幂；朴素核对只用于小 p。}
\begin{lstlisting}
typedef __int128 lll;
\end{lstlisting}
\note{O(\allowbreak{}log p)\allowbreak{}，快速幂，模乘走 \_\_int128 O(\allowbreak{}log n)\allowbreak{}，a 是底数、n>\allowbreak{}=\allowbreak{}0 是指数，p>\allowbreak{}0，模乘先升为 \_\_int128。}
\begin{lstlisting}
ll qpow(ll a,ll n,ll p)
{
    ll res=1%p;
    a%=p;
    while(n)
    {
        if(n&1)res=(ll)((lll)res*a%p);
        a=(ll)((lll)a*a%p);
        n>>=1;
    }
    return res;
}
\end{lstlisting}
\note{O(\allowbreak{}log p)\allowbreak{}，勒让德符号：1 是二次剩余，-\allowbreak{}1 是非剩余，0 表示 a\ensuremath{\equiv}0 用欧拉判别法 a\textasciicircum{}(\allowbreak{}(\allowbreak{}p-\allowbreak{}1)\allowbreak{}/\allowbreak{}2)\allowbreak{} O(\allowbreak{}log p)\allowbreak{}，返回 0/\allowbreak{}1/\allowbreak{}-\allowbreak{}1；欧拉判别只对奇素数 p 成立。}
\begin{lstlisting}
int legendre(ll a,ll p)
{
    a=(a%p<0?a%p+p:a%p);
    if(a==0)return 0;
    ll t=qpow(a,(p-1)/2,p);
    return t==1?1:-1;
}
\end{lstlisting}
\note{Cipolla 用的扩域 F\_p[\allowbreak{}sqrt(\allowbreak{}w)\allowbreak{}]\allowbreak{} 上的数：real +\allowbreak{} imag*sqrt(\allowbreak{}w)\allowbreak{}}
\begin{lstlisting}
struct Cip{
    ll real,imag;
};
ll P,W;// P 是模数，W 是扩域里的非剩余
\end{lstlisting}
\note{O(\allowbreak{}1)\allowbreak{}，扩域乘法 (\allowbreak{}a+\allowbreak{}b√w)\allowbreak{}(\allowbreak{}c+\allowbreak{}d√w)\allowbreak{} =\allowbreak{} (\allowbreak{}ac+\allowbreak{}bdw)\allowbreak{} +\allowbreak{} (\allowbreak{}ad+\allowbreak{}bc)\allowbreak{}√w O(\allowbreak{}1)\allowbreak{}，x/\allowbreak{}y 是扩域数，以全局 P/\allowbreak{}W 相乘；虚部平方用 W 代替。}
\begin{lstlisting}
Cip mul(Cip x,Cip y)
{
    Cip r;
    r.real=((lll)x.real*y.real+(lll)x.imag*y.imag%P*W)%P;
    r.imag=((lll)x.real*y.imag+(lll)x.imag*y.real)%P;
    return r;
}
\end{lstlisting}
\note{O(\allowbreak{}log p)\allowbreak{}，扩域快速幂 O(\allowbreak{}log n)\allowbreak{}，a 是扩域底数，n 是非负指数，须先设置 P/\allowbreak{}W。}
\begin{lstlisting}
Cip cpw(Cip a,ll n)
{
    Cip res;
    res.real=1%P,res.imag=0;
    while(n)
    {
        if(n&1)res=mul(res,a);
        a=mul(a,a);
        n>>=1;
    }
    return res;
}
\end{lstlisting}
\note{O(\allowbreak{}log\textasciicircum{}2 p)\allowbreak{} 期望，Cipolla 算法：解 x\textasciicircum{}2 \ensuremath{\equiv} a (\allowbreak{}mod p)\allowbreak{}，p 为奇素数 返回一个解（另一个是 p-\allowbreak{}x），a 为非二次剩余时返回 -\allowbreak{}1 期望 O(\allowbreak{}log p)\allowbreak{} 次宽整数操作，返回一个根或 -\allowbreak{}1；b 的选择随机重试。}
\begin{lstlisting}
ll cipolla(ll a,ll p)
{
    a=(a%p<0?a%p+p:a%p);
    if(a==0)return 0;
    if(p==2)return a;// 模 2 时 x=a 就是解
    if(legendre(a,p)!=1)return -1;// 无解判定
    if(p%4==3)return qpow(a,p/4+1,p);// p(*@$\equiv$@*)3 (mod 4) 有显式公式
    P=p;
    mt19937_64 rnd(chrono::steady_clock::now().time_since_epoch().count());
    ll b;
    while(true)
    {
        b=rnd()%p;
        W=((lll)b*b-a)%p;// w = b^2 - a
        W=(W%p<0?W%p+p:W%p);
        if(W!=0&&legendre(W,p)==-1)break;// w 必须是非二次剩余
    }
    Cip x;
    x.real=b,x.imag=1;
    Cip r=cpw(x,p/2+1);// (b+√w)^(p/2+1) 的虚部必为 0
    return r.real;
}
\end{lstlisting}
% SOURCE: 05-数学/同余最短路.cpp
\subsection{同余最短路}

\note{coin为正整数硬币，取最小面值m作模；返回d[\allowbreak{}r]\allowbreak{}为可表示且余数r的最小数。x可表示当且仅当d[\allowbreak{}x\%m]\allowbreak{}<\allowbreak{}=\allowbreak{}x；INF表示该余数不可达，最小面值过大时状态数也过大。}
\note{正整数硬币无限使用：以最小硬币 m 为模，dist[\allowbreak{}r]\allowbreak{} 为能表示的最小且 \ensuremath{\equiv}r 的数。 x>\allowbreak{}=\allowbreak{}0 可表示 iff dist[\allowbreak{}x\%m]\allowbreak{}<\allowbreak{}=\allowbreak{}x；用已有 Dijkstra 思路，边 r -\allowbreak{}>\allowbreak{} (\allowbreak{}r+\allowbreak{}a)\allowbreak{}\%m，权 a。 O(\allowbreak{}m*k log m)\allowbreak{}，空间 O(\allowbreak{}m)\allowbreak{}，适合 m 小；不能用最大的硬币为模随意增大状态。 gcd>\allowbreak{}1 时不可表示数无限多；gcd=\allowbreak{}1 时最大不可表示非负数是 max(\allowbreak{}dist)\allowbreak{}-\allowbreak{}m，m=\allowbreak{}1 时为 -\allowbreak{}1。}
\begin{lstlisting}
vector<long long> residue_shortest_path(vector<long long> coin)
{
    using ll=long long;const ll INF=LLONG_MAX/4;assert(!coin.empty());sort(coin.begin(),coin.end());assert(coin[0]>0&&coin[0]<=INT_MAX);
    coin.erase(unique(coin.begin(),coin.end()),coin.end());int m=coin[0];vector<ll>d(m,INF);d[0]=0;
    priority_queue<pair<ll,int>,vector<pair<ll,int>>,greater<pair<ll,int>>>q;q.push({0,0});
    while(!q.empty()){auto [x,u]=q.top();q.pop();if(x!=d[u])continue;for(ll a:coin)if(a<=INF-x){int v=(u+a%m)%m;if(x+a<d[v])d[v]=x+a,q.push({d[v],v});}}
    return d;
}
\end{lstlisting}
% SOURCE: 05-数学/SternBrocot有理数二分.cpp
\subsection{SternBrocot有理数二分}

\note{上下界是相邻有理数，判定函数决定目标在中项哪侧，批量跳步避免逐个中项。分子分母上界及端点开闭由入口约定；交叉乘法用宽类型，约分后的分母为正。}
\note{在 0<\allowbreak{}=\allowbreak{}p<\allowbreak{}=\allowbreak{}P、1<\allowbreak{}=\allowbreak{}q<\allowbreak{}=\allowbreak{}Q 的既约分数中找单调判定边界，P>\allowbreak{}=\allowbreak{}0、Q>\allowbreak{}=\allowbreak{}1。 check(\allowbreak{}p,\allowbreak{}q)\allowbreak{} 为真是一个向下闭合的前缀，必须含 0/\allowbreak{}1；返回 \{最大真分数,\allowbreak{}最小假分数\}。 无合法假分数时右界为 1/\allowbreak{}0（正无穷哨兵），不调用 check(\allowbreak{}1,\allowbreak{}0)\allowbreak{}；P,\allowbreak{}Q<\allowbreak{}=\allowbreak{}1e18。 连续同方向走 k 步用倍增+\allowbreak{}二分加速，总 O(\allowbreak{}log(\allowbreak{}P+\allowbreak{}Q)\allowbreak{})\allowbreak{} 次判定，额外空间 O(\allowbreak{}1)\allowbreak{}。}
\begin{lstlisting}
template<class Check>
array<pair<ll,ll>,2> fraction_bounds(ll P,ll Q,Check check)
{
    assert(0<=P&&P<=1000000000000000000LL&&1<=Q&&Q<=1000000000000000000LL);
    assert(check(0,1));
    ll a=0,b=1,c=1,d=0;
    while((__int128)a+c<=P&&(__int128)b+d<=Q)
    {
        bool left=check(a+c,b+d);
        ll x=left?a:c,y=left?b:d,u=left?c:a,v=left?d:b;
        ll cap=LLONG_MAX;if(u)cap=min(cap,(P-x)/u);if(v)cap=min(cap,(Q-y)/v);
        auto same=[&](ll k){return check(x+k*u,y+k*v)==left;};
        ll good=1,bad=cap;
        while(good<cap)
        {
            ll k=good>cap/2?cap:good*2;
            if(!same(k)){bad=k;break;}
            good=k;
        }
        while(good+1<bad)
        {
            ll k=good+(bad-good)/2;
            if(same(k))good=k;else bad=k;
        }
        if(left)a=x+good*u,b=y+good*v;
        else c=x+good*u,d=y+good*v;
    }
    return {{{a,b},{c,d}}};
}
\end{lstlisting}
\note{邻界满足 cb-\allowbreak{}ad=\allowbreak{}1；mediant=\allowbreak{}(\allowbreak{}a+\allowbreak{}c)\allowbreak{}/\allowbreak{}(\allowbreak{}b+\allowbreak{}d)\allowbreak{}，在它之间的其他既约分数分子分母更大。 精确比较 p/\allowbreak{}q<\allowbreak{}=\allowbreak{}x/\allowbreak{}y：(\allowbreak{}\_\_int128)\allowbreak{}p*y<\allowbreak{}=\allowbreak{}(\allowbreak{}\_\_int128)\allowbreak{}x*q；不要转 double 判等。 check 很贵时可缓存结果；链内反复走一次 mediant 可能退化为 O(\allowbreak{}P+\allowbreak{}Q)\allowbreak{}。}
% SOURCE: 05-数学/整数拆分数.cpp
\subsection{整数拆分数}

\note{n为最大和，mod>\allowbreak{}0；返回p[\allowbreak{}0.\allowbreak{}.\allowbreak{}n]\allowbreak{}，p(\allowbreak{}0)\allowbreak{}=\allowbreak{}1，拆分不计顺序。五边形下标成对出现，符号按j奇偶交替，模数不要求素数。}
\note{p(\allowbreak{}n)\allowbreak{}：把 n 拆成正整数之和，不计顺序；五边形数定理，O(\allowbreak{}n sqrt n)\allowbreak{}、O(\allowbreak{}n)\allowbreak{}。 p(\allowbreak{}0)\allowbreak{}=\allowbreak{}1，递推使用 j(\allowbreak{}3j±1)\allowbreak{}/\allowbreak{}2，符号按 j 奇偶交替；mod>\allowbreak{}0，不要求素数。}
\begin{lstlisting}
vector<long long> partition_numbers(int n,long long mod)
{
    assert(n>=0&&mod>0);vector<long long> p(n+1);p[0]=1%mod;
    for(int i=1;i<=n;i++)for(long long j=1;j*(3*j-1)/2<=i;j++)
    {
        int a=j*(3*j-1)/2,b=j*(3*j+1)/2;__int128 v=p[a<=i?i-a:0];
        if(b<=i)v+=p[i-b];if(!(j&1))v=-v;
        p[i]=(p[i]+v)%mod;if(p[i]<0)p[i]+=mod;
    }
    return p;
}
\end{lstlisting}
% SOURCE: 05-数学/线性规划单纯形.cpp
\subsection{线性规划单纯形}

\note{A、b、c分别是约束矩阵、右端上界、目标系数，求max c*x且A*x<\allowbreak{}=\allowbreak{}b、x>\allowbreak{}=\allowbreak{}0；D为单纯形表，B/\allowbreak{}N为基/\allowbreak{}非基编号。Solve输出x和最优值，±INF区分无解/\allowbreak{}无界，重新求解须重新构造。}
\note{两阶段单纯形：max c*x，约束 A*x<\allowbreak{}=\allowbreak{}b 且 x>\allowbreak{}=\allowbreak{}0；实数 LP，不保证整数解。 Solve 返回 -\allowbreak{}INF 无解、+\allowbreak{}INF 无界，否则最优值并输出 x；重复调用须重新构建对象。 实践中较快，但最坏指数级；EPS、缩放和病态数据影响精度。long double 存表。 算法整理自 Stanford ACM /\allowbreak{} KACTL 两阶段单纯形；B/\allowbreak{}N 是基/\allowbreak{}非基变量编号。}
\begin{lstlisting}
struct LPSolver
{
    using Real=long double;static constexpr Real EPS=1e-10L;
    int m,n;vector<int>B,N;vector<vector<Real>> D;
    LPSolver(const vector<vector<Real>>&A,const vector<Real>&b,const vector<Real>&c):m(b.size()),n(c.size()),B(m),N(n+1),D(m+2,vector<Real>(n+2))
    {for(int i=0;i<m;i++){for(int j=0;j<n;j++)D[i][j]=A[i][j];B[i]=n+i;D[i][n]=-1;D[i][n+1]=b[i];}for(int j=0;j<n;j++)N[j]=j,D[m][j]=-c[j];N[n]=-1;D[m+1][n]=1;}
    void pivot(int r,int s)
    {Real inv=1/D[r][s];for(int i=0;i<m+2;i++)if(i!=r)for(int j=0;j<n+2;j++)if(j!=s)D[i][j]-=D[r][j]*D[i][s]*inv;for(int j=0;j<n+2;j++)if(j!=s)D[r][j]*=inv;for(int i=0;i<m+2;i++)if(i!=r)D[i][s]*=-inv;D[r][s]=inv;swap(B[r],N[s]);}
    bool simplex(int phase)
    {
        int row=phase==1?m+1:m;
        while(true)
        {int s=-1;for(int j=0;j<=n;j++)if(!(phase==2&&N[j]==-1))if(s<0||D[row][j]<D[row][s]-EPS||(fabsl(D[row][j]-D[row][s])<=EPS&&N[j]<N[s]))s=j;
            if(s<0||D[row][s]>=-EPS)return true;int r=-1;
            for(int i=0;i<m;i++)if(D[i][s]>EPS){if(r<0)r=i;else{Real a=D[i][n+1]/D[i][s],b=D[r][n+1]/D[r][s];if(a<b-EPS||(fabsl(a-b)<=EPS&&B[i]<B[r]))r=i;}}
            if(r<0)return false;pivot(r,s);}
    }
    Real Solve(vector<Real> &x)
    {
        const Real INF=numeric_limits<Real>::infinity();x.assign(n,0);
        if(!n){for(int i=0;i<m;i++)if(D[i][n+1]<-EPS)return -INF;return 0;}
        if(m){int r=0;for(int i=1;i<m;i++)if(D[i][n+1]<D[r][n+1])r=i;
            if(D[r][n+1]<-EPS){pivot(r,n);if(!simplex(1)||fabsl(D[m+1][n+1])>EPS)return -INF;
                for(int i=0;i<m;i++)if(B[i]==-1){int s=-1;for(int j=0;j<=n;j++)if(fabsl(D[i][j])>EPS&&(s<0||N[j]<N[s]))s=j;if(s>=0)pivot(i,s);}}}
        if(!simplex(2))return INF;for(int i=0;i<m;i++)if(0<=B[i]&&B[i]<n)x[B[i]]=D[i][n+1];return D[m][n+1];
    }
};
\end{lstlisting}
% SOURCE: 05-数学/自适应Simpson积分.cpp
\subsection{自适应Simpson积分}

\note{f为被积函数，l/\allowbreak{}r为积分区间，eps为误差目标；递归比较整段与两半Simpson估计。函数在区间内须可积且数值合理，奇点、振荡及递归深度仍需按题目处理。}
\note{f 在区间内取值有限，eps>\allowbreak{}0 是绝对误差目标；返回积分估计，允许 l>\allowbreak{}r。 误差估计依赖函数足够光滑；不连续点、奇点应先拆段，振荡函数可能漏采样。 两半 Simpson 与整段比较，误差约为差值/\allowbreak{}15；分给子区间 eps/\allowbreak{}2，复用端点值。 ok=\allowbreak{}false 表示深度耗尽或浮点中点无法细分，仍返回当前估计，不能当成精度达标。 每次细分只新增两个采样；时间 O(\allowbreak{}采样数)\allowbreak{}，栈空间 O(\allowbreak{}depth)\allowbreak{}。}
\begin{lstlisting}
template<class F>
double simpson(F f,double l,double r,double eps,bool &ok,int depth=25)
{
    assert(eps>0&&depth>=0&&isfinite(l)&&isfinite(r));
    ok=true;
    if(l==r)return 0;
    double sign=1;
    if(l>r)swap(l,r),sign=-1;
    auto area=[](double l,double r,double a,double b,double c){return (r-l)*(a+4*b+c)/6;};
    auto solve=[&](auto &&self,double l,double r,double a,double b,double c,double s,double e,int dep)->double
    {
        double m=l+(r-l)/2,x=l+(m-l)/2,y=m+(r-m)/2;
        if(x==l||x==m||y==m||y==r){ok=false;return s;}
        double d=f(x),v=f(y),left=area(l,m,a,d,b),right=area(m,r,b,v,c);
        double delta=left+right-s;
        if(fabs(delta)<=15*e)return left+right+delta/15;
        if(!dep){ok=false;return left+right+delta/15;}
        return self(self,l,m,a,d,b,left,e/2,dep-1)+self(self,m,r,b,v,c,right,e/2,dep-1);
    };
    double m=l+(r-l)/2,a=f(l),b=f(m),c=f(r);
    return sign*solve(solve,l,r,a,b,c,area(l,r,a,b,c),eps,depth);
}
\end{lstlisting}
\section{动态规划}
% SOURCE: 06-动态规划/递推与线性DP.cpp
\subsection{递推与线性DP}

\note{n/\allowbreak{}m为各例规模，a、s、t为数值/\allowbreak{}字符输入；f/\allowbreak{}g为滚动状态，dp为二维状态，best保存辅助最优量。tri为数字三角形，grid中1为障碍，dp2为路径计数；按所选例子初始化边界，其他例子的数组不必复制。}
\begin{lstlisting}
const int N=1005;
const int INF=0x3f3f3f3f;
const int NEG=-0x3f3f3f3f;
int n,m;
int a[N];
char s[N],t[N];
int f[N],g[N],dp[N][N],best[N];// f/g 是滚动的线性 dp 数组
int tri[15][15];// 数字三角形
ll dp2[15][15];
int grid[15][15];// 网格路径（1 表示障碍）
\end{lstlisting}
\note{O(\allowbreak{}n\textasciicircum{}2)\allowbreak{}，最长上升子序列（严格递增），f[\allowbreak{}i]\allowbreak{} 表示以 i 结尾的 LIS 长度 转移：f[\allowbreak{}i]\allowbreak{}=\allowbreak{}max(\allowbreak{}f[\allowbreak{}j]\allowbreak{})\allowbreak{}+\allowbreak{}1，其中 j<\allowbreak{}i 且 a[\allowbreak{}j]\allowbreak{}<\allowbreak{}a[\allowbreak{}i]\allowbreak{}}
\begin{lstlisting}
int lis_n2()
{
    int ans=0;
    for(int i=1;i<=n;i++)
    {
        f[i]=1;
        for(int j=1;j<i;j++)
            if(a[j]<a[i])f[i]=max(f[i],f[j]+1);
        ans=max(ans,f[i]);
    }
    return ans;
}
\end{lstlisting}
\note{O(\allowbreak{}n log n)\allowbreak{}，LIS 贪心+\allowbreak{}二分：g[\allowbreak{}len]\allowbreak{} 存长度为 len 的上升子序列的最小结尾}
\begin{lstlisting}
int lis_nlogn()
{
    int len=0;
    for(int i=1;i<=n;i++)
    {
        int p=lower_bound(g+1,g+len+1,a[i])-g;
        g[p]=a[i];
        len=max(len,p);
    }
    return len;
}
\end{lstlisting}
\note{O(\allowbreak{}n)\allowbreak{}，最大子段和：f[\allowbreak{}i]\allowbreak{} 表示以 i 结尾的最大子段和，f[\allowbreak{}i]\allowbreak{}=\allowbreak{}max(\allowbreak{}f[\allowbreak{}i-\allowbreak{}1]\allowbreak{},\allowbreak{}0)\allowbreak{}+\allowbreak{}a[\allowbreak{}i]\allowbreak{}}
\begin{lstlisting}
int max_subarray()
{
    f[0]=0;
    int ans=INT_MIN;
    for(int i=1;i<=n;i++)
    {
        f[i]=max(f[i-1],0)+a[i];
        ans=max(ans,f[i]);
    }
    return ans;
}
\end{lstlisting}
\note{O(\allowbreak{}n\textasciicircum{}2)\allowbreak{}，数字三角形：从顶走到底的最大路径和，顺推 转移：dp[\allowbreak{}i]\allowbreak{}[\allowbreak{}j]\allowbreak{}=\allowbreak{}max(\allowbreak{}dp[\allowbreak{}i-\allowbreak{}1]\allowbreak{}[\allowbreak{}j-\allowbreak{}1]\allowbreak{},\allowbreak{}dp[\allowbreak{}i-\allowbreak{}1]\allowbreak{}[\allowbreak{}j]\allowbreak{})\allowbreak{}+\allowbreak{}tri[\allowbreak{}i]\allowbreak{}[\allowbreak{}j]\allowbreak{} 边界必须判掉：j=\allowbreak{}1 没有左上方，j=\allowbreak{}i 没有正上方（否则会把没算过的 0 当答案）}
\begin{lstlisting}
int triangle_max()
{
    dp[0][1]=0;
    for(int i=1;i<=n;i++)
        for(int j=1;j<=i;j++)
        {
            if(j==1)dp[i][j]=dp[i-1][j]+tri[i][j];
            else if(j==i)dp[i][j]=dp[i-1][j-1]+tri[i][j];
            else dp[i][j]=max(dp[i-1][j-1],dp[i-1][j])+tri[i][j];
        }
    int ans=INT_MIN;
    for(int j=1;j<=n;j++)ans=max(ans,dp[n][j]);
    return ans;
}
\end{lstlisting}
\note{O(\allowbreak{}n\textasciicircum{}2)\allowbreak{}，网格路径数：只能往右/\allowbreak{}往下走，grid=\allowbreak{}1 是障碍}
\begin{lstlisting}
ll grid_paths()
{
    for(int i=0;i<=n;i++)
        for(int j=0;j<=m;j++)dp2[i][j]=0;
    dp2[1][1]=(grid[1][1]==1)?0:1;
    for(int i=1;i<=n;i++)
        for(int j=1;j<=m;j++)
        {
            if(i==1&&j==1)continue;
            if(grid[i][j]==1){dp2[i][j]=0;continue;}
            dp2[i][j]=dp2[i-1][j]+dp2[i][j-1];// 从上面或左面推过来
        }
    return dp2[n][m];
}
\end{lstlisting}
\note{O(\allowbreak{}n\textasciicircum{}2)\allowbreak{}，最长公共子序列长度}
\begin{lstlisting}
int lcs()
{
    for(int i=0;i<=n;i++)
        for(int j=0;j<=m;j++)
        {
            if(i==0||j==0){dp[i][j]=0;continue;}
            if(s[i]==t[j])dp[i][j]=dp[i-1][j-1]+1;
            else dp[i][j]=max(dp[i-1][j],dp[i][j-1]);
        }
    return dp[n][m];
}
\end{lstlisting}
\note{O(\allowbreak{}n)\allowbreak{}，爬楼梯方案数（每次 1 或 2 级），f[\allowbreak{}i]\allowbreak{}=\allowbreak{}f[\allowbreak{}i-\allowbreak{}1]\allowbreak{}+\allowbreak{}f[\allowbreak{}i-\allowbreak{}2]\allowbreak{}}
\begin{lstlisting}
ll climb(int n)
{
    ll x=1,y=1;// f[0]=1,f[1]=1
    for(int i=2;i<=n;i++)
    {
        ll z=x+y;
        x=y,y=z;
    }
    return y;
}
\end{lstlisting}
\note{O(\allowbreak{}n)\allowbreak{}，数字解码方案数：'1'\textasciitilde{}'9' 单独成一位，'10'\textasciitilde{}'26' 两位成一位}
\begin{lstlisting}
ll decode_ways()
{
    ll f0=1,f1=(n>=1&&s[1]!='0')?1:0;// f[i-2],f[i-1]
    if(n==0)return 0;
    if(n==1)return f1;
    for(int i=2;i<=n;i++)
    {
        ll cur=0;
        if(s[i]!='0')cur+=f1;
        int two=(s[i-1]-'0')*10+(s[i]-'0');
        if(two>=10&&two<=26)cur+=f0;
        f0=f1,f1=cur;
    }
    return f1;
}
\end{lstlisting}
% SOURCE: 06-动态规划/背包(01完全多重).cpp
\subsection{背包(01完全多重)}

\note{n为物品数，V为容量；w、v、c分别是重量、价值、数量，typ为0/\allowbreak{}1/\allowbreak{}多重类别，下标1起。f[\allowbreak{}j]\allowbreak{}是容量j的最优值，q、qb为多重背包的旧状态队列；要求恰好装满时用不可达负无穷初始化并读exact\_ans。}
\begin{lstlisting}
const int N=1005;// 物品数上限
const int M=100005;// 容量上限
const int INF=0x3f3f3f3f;
int n,V;
int w[N],v[N],c[N],typ[N];// 重量 价值 个数 类型(0:01 1:完全 2:多重)
int f[M];// 一维滚动数组，f[j] 表示容量 j 的最优价值
int q[M],qb[M];// 单调队列：存 (旧f[余数+j*w]-j*v) 和这个 j
int exact_ans;
\end{lstlisting}
\note{O(\allowbreak{}nV)\allowbreak{}，01 背包：容量不超过 V 的最大价值}
\begin{lstlisting}
int knap_01()
{
    for(int j=0;j<=V;j++)f[j]=0;
    for(int i=1;i<=n;i++)
        for(int j=V;j>=w[i];j--)// 倒序，保证每个物品最多用一次
            f[j]=max(f[j],f[j-w[i]]+v[i]);
    return f[V];
}
\end{lstlisting}
\note{O(\allowbreak{}nV)\allowbreak{}，01 背包恰好装满 V，装不满返回 -\allowbreak{}INF}
\begin{lstlisting}
int knap_01_exact()
{
    for(int j=0;j<=V;j++)f[j]=-INF;
    f[0]=0;// 只有容量 0 是合法起点，其余都不可达
    for(int i=1;i<=n;i++)
        for(int j=V;j>=w[i];j--)
            if(f[j-w[i]]!=-INF)f[j]=max(f[j],f[j-w[i]]+v[i]);
    return f[V];
}
\end{lstlisting}
\note{O(\allowbreak{}nV)\allowbreak{}，完全背包：容量正序，同一物品可重复选}
\begin{lstlisting}
int knap_complete()
{
    for(int j=0;j<=V;j++)f[j]=0;
    for(int i=1;i<=n;i++)
        for(int j=w[i];j<=V;j++)
            f[j]=max(f[j],f[j-w[i]]+v[i]);
    return f[V];
}
\end{lstlisting}
\note{O(\allowbreak{}V*\ensuremath{\Sigma}log(\allowbreak{}c[\allowbreak{}i]\allowbreak{}+\allowbreak{}1)\allowbreak{})\allowbreak{}，多重背包：二进制拆分把 c 个物品拆成 1,\allowbreak{}2,\allowbreak{}4,\allowbreak{}.\allowbreak{}.\allowbreak{}.\allowbreak{} 个 01 物品}
\begin{lstlisting}
int knap_multiple_binary()
{
    for(int j=0;j<=V;j++)f[j]=0;
    for(int i=1;i<=n;i++)
    {
        int k=1,cc=min(c[i],V/w[i]);
        while(cc>0)
        {
            int t=min(k,cc);
            for(int j=V;j>=t*w[i];j--)
                f[j]=max(f[j],f[j-t*w[i]]+t*v[i]);
            cc-=t,k<<=1;
        }
    }
    return f[V];
}
\end{lstlisting}
\note{O(\allowbreak{}nV)\allowbreak{}，多重背包：按余数分组 +\allowbreak{} 单调队列，窗口大小 lim+\allowbreak{}1 同余数下 j=\allowbreak{}r+\allowbreak{}k*w，f[\allowbreak{}j]\allowbreak{}=\allowbreak{}max(\allowbreak{}f[\allowbreak{}旧j']\allowbreak{}-\allowbreak{}k'*v)\allowbreak{}+\allowbreak{}k*v，k-\allowbreak{}k'<\allowbreak{}=\allowbreak{}c}
\begin{lstlisting}
int knap_multiple_deque()
{
    for(int j=0;j<=V;j++)f[j]=0;
    for(int i=1;i<=n;i++)
    {
        if(w[i]>V)continue;
        int lim=min(c[i],V/w[i]);// 实际最多能放几个
        for(int r=0;r<w[i]&&r<=V;r++)// 枚举余数类
        {
            int head=0,tail=0,k=0;
            for(int j=r;j<=V;j+=w[i],k++)
            {
                int cur=f[j]-k*v[i];// 先取旧值，再覆盖 f[j]
                while(head<tail&&qb[head]<k-lim)head++;// 队首超出 c 个
                while(head<tail&&q[tail-1]<=cur)tail--;// 队尾不优
                q[tail]=cur,qb[tail]=k,tail++;
                f[j]=q[head]+k*v[i];
            }
        }
    }
    return f[V];
}
\end{lstlisting}
\note{O(\allowbreak{}V*(\allowbreak{}n+\allowbreak{}\ensuremath{\Sigma}多重log(\allowbreak{}c[\allowbreak{}i]\allowbreak{}+\allowbreak{}1)\allowbreak{})\allowbreak{})\allowbreak{}，混合背包：0=\allowbreak{}01 1=\allowbreak{}完全 2=\allowbreak{}多重，多重用二进制拆分}
\begin{lstlisting}
int knap_mixed()
{
    for(int j=0;j<=V;j++)f[j]=0;
    for(int i=1;i<=n;i++)
    {
        if(typ[i]==0)
            for(int j=V;j>=w[i];j--)f[j]=max(f[j],f[j-w[i]]+v[i]);
        else if(typ[i]==1)
            for(int j=w[i];j<=V;j++)f[j]=max(f[j],f[j-w[i]]+v[i]);
        else
        {
            int k=1,cc=min(c[i],V/w[i]);
            while(cc>0)
            {
                int t=min(k,cc);
                for(int j=V;j>=t*w[i];j--)f[j]=max(f[j],f[j-t*w[i]]+t*v[i]);
                cc-=t,k<<=1;
            }
        }
    }
    return f[V];
}
\end{lstlisting}
% SOURCE: 06-动态规划/分组背包.cpp
\subsection{分组背包}

\note{cnt[\allowbreak{}i]\allowbreak{}、gw[\allowbreak{}i]\allowbreak{}[\allowbreak{}j]\allowbreak{}、gv[\allowbreak{}i]\allowbreak{}[\allowbreak{}j]\allowbreak{}是第i组数量、重量、价值；mw、mv为主件，ac、aw、av为附件数及属性。f[\allowbreak{}j]\allowbreak{}为容量j的最优值。组内选择须从同一旧状态转移；主件与附件下标均1起，重新求解清空f。}
\begin{lstlisting}
const int N=105;// 组数上限
const int K=1005;// 每组物品数上限
const int VMAX=20005;// 容量上限
const int KN=3;// 每个主件最多 2 个附件
int n,V,g;
int cnt[N],gw[N][K],gv[N][K];// 每组物品数 与 组内物品
int f[VMAX];// 一维滚动数组，f[j] 表示容量 j 的最大价值
int mw[N],mv[N],ac[N];// 主件重量 价值 附件数
int aw[N][KN],av[N][KN];// 附件重量 价值，1-indexed
\end{lstlisting}
\note{O(\allowbreak{}V*总物品数)\allowbreak{}，分组背包：每组至多选一个 必须先枚举容量(\allowbreak{}倒序)\allowbreak{}、再枚举组内物品，否则同组会选多个}
\begin{lstlisting}
int group_knap(int g,int V)
{
    for(int j=0;j<=V;j++)f[j]=0;
    for(int k=1;k<=g;k++)
        for(int j=V;j>=0;j--)
            for(int i=1;i<=cnt[k];i++)
                if(j>=gw[k][i])
                    f[j]=max(f[j],f[j-gw[k][i]]+gv[k][i]);
    return f[V];
}
\end{lstlisting}
\note{O(\allowbreak{}V*4n)\allowbreak{}，依赖背包：主件+\allowbreak{}最多 2 个附件，把每个主件的合法方案当成一组 想买附件必须先买主件，所以组内方案只有 4 种}
\begin{lstlisting}
int depend_knap(int n,int V)
{
    for(int j=0;j<=V;j++)f[j]=0;
    for(int i=1;i<=n;i++)
    {
        int cw[5],cv[5],tot=0;// 该主件的所有合法方案
        tot++,cw[tot]=mw[i],cv[tot]=mv[i];// 只买主件
        if(ac[i]>=1)tot++,cw[tot]=mw[i]+aw[i][1],cv[tot]=mv[i]+av[i][1];
        if(ac[i]>=2)
        {
            tot++,cw[tot]=mw[i]+aw[i][2],cv[tot]=mv[i]+av[i][2];
            if(mw[i]+aw[i][1]+aw[i][2]<=V)
                tot++,cw[tot]=mw[i]+aw[i][1]+aw[i][2],cv[tot]=mv[i]+av[i][1]+av[i][2];
        }
        for(int j=V;j>=0;j--)// 倒序，一个主件只取一种方案
            for(int t=1;t<=tot;t++)
                if(j>=cw[t])f[j]=max(f[j],f[j-cw[t]]+cv[t]);
    }
    return f[V];
}
\end{lstlisting}
% SOURCE: 06-动态规划/最长上升子序列LIS.cpp
\subsection{最长上升子序列LIS}

\note{a[\allowbreak{}1.\allowbreak{}.\allowbreak{}n]\allowbreak{}为序列；d维护各长度最小末尾，f、pre记录长度与前驱，st还原下标。b、tmp、mp、c供非降/\allowbreak{}排列转换等独立示例使用；严格递增用lower\_bound，允许相等用upper\_bound，不能混用。}
\begin{lstlisting}
const int N=100005;
const int INF=0x3f3f3f3f;
int n;
int a[N],b[N],d[N],f[N],pre[N],tmp[N],mp[N],c[N];
int st[N];// 还原方案用的栈
\end{lstlisting}
\note{O(\allowbreak{}n\textasciicircum{}2)\allowbreak{}，最长严格上升子序列：f[\allowbreak{}i]\allowbreak{} 表示以 a[\allowbreak{}i]\allowbreak{} 结尾的最长长度}
\begin{lstlisting}
int lis_n2()
{
    int ans=0;
    for(int i=1;i<=n;i++)
    {
        f[i]=1;
        for(int j=1;j<i;j++)
            if(a[j]<a[i])f[i]=max(f[i],f[j]+1);
        ans=max(ans,f[i]);
    }
    return ans;
}
\end{lstlisting}
\note{O(\allowbreak{}n log n)\allowbreak{}，最长严格上升子序列 d[\allowbreak{}len]\allowbreak{} 表示长度 len 的上升子序列的最小结尾，d 单调不减}
\begin{lstlisting}
int lis_nlogn()
{
    int len=0;
    for(int i=1;i<=n;i++)
    {
        int p=lower_bound(d+1,d+len+1,a[i])-d;
        d[p]=a[i];
        len=max(len,p);
    }
    return len;
}
\end{lstlisting}
\note{O(\allowbreak{}n log n)\allowbreak{}，最长不降子序列：只把 >\allowbreak{}=\allowbreak{} 改成 >\allowbreak{}，二分找第一个比 a[\allowbreak{}i]\allowbreak{} 大的}
\begin{lstlisting}
int lnds_nlogn()
{
    int len=0;
    for(int i=1;i<=n;i++)
    {
        int p=upper_bound(d+1,d+len+1,a[i])-d;
        d[p]=a[i];
        len=max(len,p);
    }
    return len;
}
\end{lstlisting}
\note{O(\allowbreak{}n log n)\allowbreak{}，最长严格下降子序列：用降序比较器，避免 INT\_MIN 取负溢出}
\begin{lstlisting}
int lds_nlogn()
{
    int len=0;
    for(int i=1;i<=n;i++)
    {
        int p=lower_bound(d+1,d+len+1,a[i],greater<int>())-d;
        d[p]=a[i];
        len=max(len,p);
    }
    return len;
}
\end{lstlisting}
\note{O(\allowbreak{}n log n)\allowbreak{}，两个排列(\allowbreak{}值 1.\allowbreak{}.\allowbreak{}n)\allowbreak{}的 LCS 转 LIS a 中每个值的位置记下来，按 b 的顺序排成 c，c 的 LIS 就是 LCS}
\begin{lstlisting}
int lcs_perm()
{
    for(int i=1;i<=n;i++)mp[a[i]]=i;// 值 -> 在 a 中的位置
    for(int i=1;i<=n;i++)c[i]=mp[b[i]];
    int len=0;
    for(int i=1;i<=n;i++)
    {
        int p=lower_bound(d+1,d+len+1,c[i])-d;
        d[p]=c[i],len=max(len,p);
    }
    return len;
}
\end{lstlisting}
\note{O(\allowbreak{}n\textasciicircum{}2)\allowbreak{}，记录前驱并输出一组最优解}
\begin{lstlisting}
void lis_scheme()
{
    int ans=0,best=0,top=0;
    for(int i=1;i<=n;i++)
    {
        f[i]=1,pre[i]=0;
        for(int j=1;j<i;j++)
            if(a[j]<a[i]&&f[j]+1>f[i])f[i]=f[j]+1,pre[i]=j;
        if(f[i]>ans)ans=f[i],best=i;
    }
    while(best)st[++top]=a[best],best=pre[best];// 逆着走前驱
    printf("长度 %d，一组方案：",ans);
    for(int i=top;i>=1;i--)printf("%d ",st[i]);
    printf("\n");
}
\end{lstlisting}
% SOURCE: 06-动态规划/最长公共子序列LCS.cpp
\subsection{最长公共子序列LCS}

\note{a[\allowbreak{}1.\allowbreak{}.\allowbreak{}n]\allowbreak{}、b[\allowbreak{}1.\allowbreak{}.\allowbreak{}m]\allowbreak{}是字符序列，f为滚动一行的长度，g为保留完整状态的二维表，st用于还原。求长度只需f；求方案用g，平局时按还原规则选一条；每次求解重置所用表。}
\begin{lstlisting}
const int N=1005;
int n,m;
char a[N],b[N];
int f[N];// 一维滚动数组
int g[N][N];// LCS 二维表，还原方案用
int st[N];// 还原方案用的栈
\end{lstlisting}
\note{O(\allowbreak{}nm)\allowbreak{}，最长公共子序列：g[\allowbreak{}i]\allowbreak{}[\allowbreak{}j]\allowbreak{} 表示 a 前 i 个与 b 前 j 个的 LCS 长度}
\begin{lstlisting}
int lcs_nm()
{
    for(int i=0;i<=n;i++)g[i][0]=0;
    for(int j=0;j<=m;j++)g[0][j]=0;
    for(int i=1;i<=n;i++)
        for(int j=1;j<=m;j++)
            if(a[i]==b[j])g[i][j]=g[i-1][j-1]+1;// 这一位匹配，直接要
            else g[i][j]=max(g[i-1][j],g[i][j-1]);// 否则丢 a 的尾或 b 的尾
    return g[n][m];
}
\end{lstlisting}
\note{O(\allowbreak{}nm)\allowbreak{}，LCS 一维滚动：f[\allowbreak{}j]\allowbreak{} 是当前行的值，pre 存上一行左上角}
\begin{lstlisting}
int lcs_roll()
{
    for(int j=0;j<=m;j++)f[j]=0;
    for(int i=1;i<=n;i++)
    {
        int pre=0;// f[i-1][j-1]
        for(int j=1;j<=m;j++)
        {
            int tmpj=f[j];// 旧的 f[i-1][j]
            if(a[i]==b[j])f[j]=pre+1;
            else f[j]=max(f[j],f[j-1]);// f[j-1] 已是 f[i][j-1]
            pre=tmpj;
        }
    }
    return f[m];
}
\end{lstlisting}
\note{O(\allowbreak{}nm)\allowbreak{}，还原一组 LCS 并输出（必须先调用 lcs\_nm 建好 g）}
\begin{lstlisting}
void lcs_scheme()
{
    int top=0,i=n,j=m;
    while(i>0&&j>0)
    {
        if(a[i]==b[j])st[++top]=a[i],i--,j--;
        else if(g[i-1][j]>=g[i][j-1])i--;
        else j--;
    }
    printf("长度 %d，一组方案：",g[n][m]);
    for(int k=top;k>=1;k--)putchar(st[k]);
    printf("\n");
}
\end{lstlisting}
\note{O(\allowbreak{}nm)\allowbreak{}，编辑距离：g[\allowbreak{}i]\allowbreak{}[\allowbreak{}j]\allowbreak{} 表示 a 前 i 个变成 b 前 j 个的最少操作数}
\begin{lstlisting}
int edit_dist()
{
    for(int i=0;i<=n;i++)g[i][0]=i;// 全删
    for(int j=0;j<=m;j++)g[0][j]=j;// 全插
    for(int i=1;i<=n;i++)
        for(int j=1;j<=m;j++)
            if(a[i]==b[j])g[i][j]=g[i-1][j-1];
            else g[i][j]=min(min(g[i-1][j],g[i][j-1]),g[i-1][j-1])+1;// 删 插 换
    return g[n][m];
}
\end{lstlisting}
\note{O(\allowbreak{}nm)\allowbreak{}，编辑距离一维滚动：f[\allowbreak{}j-\allowbreak{}1]\allowbreak{} 是 f[\allowbreak{}i]\allowbreak{}[\allowbreak{}j-\allowbreak{}1]\allowbreak{}，pre 是 f[\allowbreak{}i-\allowbreak{}1]\allowbreak{}[\allowbreak{}j-\allowbreak{}1]\allowbreak{}}
\begin{lstlisting}
int edit_dist_roll()
{
    for(int j=0;j<=m;j++)f[j]=j;// 第 0 行
    for(int i=1;i<=n;i++)
    {
        int pre=f[0];// f[i-1][0]
        f[0]=i;// f[i][0]=i
        for(int j=1;j<=m;j++)
        {
            int tmpj=f[j];// 旧的 f[i-1][j]
            if(a[i]==b[j])f[j]=pre;
            else f[j]=min(min(f[j],f[j-1]),pre)+1;
            pre=tmpj;
        }
    }
    return f[m];
}
\end{lstlisting}
\note{O(\allowbreak{}nm)\allowbreak{}，最长公共子串（必须连续）：不匹配就直接清 0}
\begin{lstlisting}
int lcsubstr()
{
    int ans=0;
    for(int j=0;j<=m;j++)f[j]=0;
    for(int i=1;i<=n;i++)
    {
        int pre=0;
        for(int j=1;j<=m;j++)
        {
            int tmpj=f[j];
            if(a[i]==b[j])f[j]=pre+1;
            else f[j]=0;
            pre=tmpj;
            ans=max(ans,f[j]);
        }
    }
    return ans;
}
\end{lstlisting}
% SOURCE: 06-动态规划/区间DP.cpp
\subsection{区间DP}

\note{a、sum为输入和前缀和；gmin/\allowbreak{}gmax是环形石子合并，emax是能量项链，pmin/\allowbreak{}pmax是含负数乘法的多边形状态。val/\allowbreak{}op为原输入，v/\allowbreak{}o为复制后的环序列；不同示例使用不同表，容量须覆盖展开到2n的下标。}
\begin{lstlisting}
const int N=505;// 断环成链后要开 2n
const int INF=0x3f3f3f3f;
const int NEG=-0x3f3f3f3f;
int n;
int a[N],sum[N];// a 存石子/项链，sum 是前缀和
int gmin[N][N],gmax[N][N];// 环形石子合并
int emax[N][N];// 能量项链
int pmax[N][N],pmin[N][N];// 多边形游戏：有负数乘法，必须同时存 max 和 min
int val[N],v[N];// 多边形顶点
char op[N],o[N];// 多边形边上的运算符，op[i] 连 val[i] 和 val[i+1]
\end{lstlisting}
\note{O(\allowbreak{}n\textasciicircum{}3)\allowbreak{}，环形石子合并最小代价 f[\allowbreak{}l]\allowbreak{}[\allowbreak{}r]\allowbreak{} 表示把 [\allowbreak{}l,\allowbreak{}r]\allowbreak{} 合并成一堆的最小代价，枚举断点 k 断环成链 a[\allowbreak{}i+\allowbreak{}n]\allowbreak{}=\allowbreak{}a[\allowbreak{}i]\allowbreak{}，答案是 min f[\allowbreak{}i]\allowbreak{}[\allowbreak{}i+\allowbreak{}n-\allowbreak{}1]\allowbreak{}}
\begin{lstlisting}
int stone_merge_min()
{
    for(int i=1;i<=n;i++)a[i+n]=a[i];
    sum[0]=0;
    for(int i=1;i<=2*n;i++)sum[i]=sum[i-1]+a[i];
    for(int i=1;i<=2*n;i++)
        for(int j=1;j<=2*n;j++)gmin[i][j]=(i==j)?0:INF;// 边界：一堆不用合并
    for(int len=2;len<=n;len++)// 必须按区间长度从小到大
        for(int l=1;l+len-1<=2*n;l++)
        {
            int r=l+len-1;
            for(int k=l;k<r;k++)
                gmin[l][r]=min(gmin[l][r],gmin[l][k]+gmin[k+1][r]+sum[r]-sum[l-1]);
        }
    int ans=INF;
    for(int i=1;i<=n;i++)ans=min(ans,gmin[i][i+n-1]);
    return ans;
}
\end{lstlisting}
\note{O(\allowbreak{}n\textasciicircum{}3)\allowbreak{}，环形石子合并最大代价}
\begin{lstlisting}
int stone_merge_max()
{
    for(int i=1;i<=n;i++)a[i+n]=a[i];
    sum[0]=0;
    for(int i=1;i<=2*n;i++)sum[i]=sum[i-1]+a[i];
    for(int i=1;i<=2*n;i++)
        for(int j=1;j<=2*n;j++)gmax[i][j]=0;
    for(int len=2;len<=n;len++)
        for(int l=1;l+len-1<=2*n;l++)
        {
            int r=l+len-1;
            for(int k=l;k<r;k++)
                gmax[l][r]=max(gmax[l][r],gmax[l][k]+gmax[k+1][r]+sum[r]-sum[l-1]);
        }
    int ans=0;
    for(int i=1;i<=n;i++)ans=max(ans,gmax[i][i+n-1]);
    return ans;
}
\end{lstlisting}
\note{O(\allowbreak{}n\textasciicircum{}3)\allowbreak{}，能量项链：emax[\allowbreak{}l]\allowbreak{}[\allowbreak{}r]\allowbreak{} 表示 [\allowbreak{}l,\allowbreak{}r]\allowbreak{} 合成一颗珠子的最大能量 合并 (\allowbreak{}l.\allowbreak{}.\allowbreak{}k)\allowbreak{} 和 (\allowbreak{}k.\allowbreak{}.\allowbreak{}r)\allowbreak{} 时释放 a[\allowbreak{}l]\allowbreak{}*a[\allowbreak{}k]\allowbreak{}*a[\allowbreak{}r]\allowbreak{}，注意端点共用}
\begin{lstlisting}
int energy_necklace()
{
    for(int i=1;i<=n+1;i++)a[i+n]=a[i];// 链要开到 2n+1
    for(int i=1;i<=2*n+1;i++)
        for(int j=1;j<=2*n+1;j++)emax[i][j]=0;
    for(int len=2;len<=n;len++)
        for(int l=1;l+len<=2*n+1;l++)
        {
            int r=l+len;
            for(int k=l+1;k<r;k++)
                emax[l][r]=max(emax[l][r],emax[l][k]+emax[k][r]+a[l]*a[k]*a[r]);
        }
    int ans=0;
    for(int i=1;i<=n;i++)ans=max(ans,emax[i][i+n]);
    return ans;
}
\end{lstlisting}
\note{O(\allowbreak{}n\textasciicircum{}4)\allowbreak{}，多边形游戏：枚举删哪条边(\allowbreak{}旋转)\allowbreak{}，再区间 dp 求最大值 乘法遇到负数会让"最小"翻成"最大"，所以 max 由 max*max /\allowbreak{} max*min /\allowbreak{} min*max /\allowbreak{} min*min 取}
\begin{lstlisting}
int polygon_game()
{
    int ans=NEG;
    for(int cut=1;cut<=n;cut++)// 删掉第 cut 条边，链从 val[cut+1] 开始
    {
        for(int i=1;i<=n;i++)v[i]=val[(cut+i-1)%n+1];
        for(int i=1;i<=n-1;i++)o[i]=op[(cut+i-1)%n+1];
        for(int i=1;i<=n;i++)
            for(int j=1;j<=n;j++)pmax[i][j]=NEG,pmin[i][j]=INF;
        for(int i=1;i<=n;i++)pmax[i][i]=pmin[i][i]=v[i];
        for(int len=2;len<=n;len++)
            for(int l=1;l+len-1<=n;l++)
            {
                int r=l+len-1;
                for(int k=l;k<r;k++)
                {
                    int cand[4],cnt=0;
                    if(o[k]=='+')
                    {
                        cand[cnt++]=pmax[l][k]+pmax[k+1][r];
                        cand[cnt++]=pmin[l][k]+pmin[k+1][r];
                    }
                    else
                    {
                        cand[cnt++]=pmax[l][k]*pmax[k+1][r];
                        cand[cnt++]=pmax[l][k]*pmin[k+1][r];
                        cand[cnt++]=pmin[l][k]*pmax[k+1][r];
                        cand[cnt++]=pmin[l][k]*pmin[k+1][r];
                    }
                    for(int t=0;t<cnt;t++)
                        pmax[l][r]=max(pmax[l][r],cand[t]),pmin[l][r]=min(pmin[l][r],cand[t]);
                }
            }
        ans=max(ans,pmax[1][n]);
    }
    return ans;
}
\end{lstlisting}
% SOURCE: 06-动态规划/树形DP.cpp
\subsection{树形DP}

\note{adj是无向树，son是有根子节点；h为欢乐值，score为学分，par/\allowbreak{}vis供建根。dp[\allowbreak{}u]\allowbreak{}[\allowbreak{}0/\allowbreak{}1]\allowbreak{}表示不选/\allowbreak{}选u，f[\allowbreak{}u]\allowbreak{}[\allowbreak{}j]\allowbreak{}表示子树选j门课，d1/\allowbreak{}d2为最长两条向下链；m为选课数，dia保存直径。}
\note{乘法换根不能默认模除可用：子项可能为 0，或模数非素；前后缀积能排除某个儿子而不求逆。 小边界连通 DP：三个不同边界点只有 5 种集合划分；边界重合先去重，再合并划分。 若统计连通块总数，维护(\allowbreak{}方案数,\allowbreak{}已封闭块数之和)\allowbreak{}，新封闭 t 块增加 t*两侧方案数乘积。}
\begin{lstlisting}
const int N=1005;
const int NEG=-0x3f3f3f3f;
int n,m,root,dia;
int h[N],score[N],par[N],vis[N];// 欢乐值 学分 父亲 是否当过儿子
vector<int> son[N],adj[N];
int dp[N][2];// dp[x][0] x 不来 dp[x][1] x 来
int f[N][N];// f[x][j] 以 x 为根的子树里选 j 门课
int d1[N],d2[N];// 最长/次长向下链
\end{lstlisting}
\note{O(\allowbreak{}n)\allowbreak{}，没有上司的舞会 dp[\allowbreak{}x]\allowbreak{}[\allowbreak{}0]\allowbreak{} 表示 x 不参加时子树的最大欢乐值(\allowbreak{}儿子可选可不选)\allowbreak{} dp[\allowbreak{}x]\allowbreak{}[\allowbreak{}1]\allowbreak{} 表示 x 参加时的最大值(\allowbreak{}儿子都不能参加)\allowbreak{}}
\begin{lstlisting}
void dfs_party(int x)
{
    dp[x][0]=0,dp[x][1]=h[x];
    for(int i=0;i<(int)son[x].size();i++)
    {
        int y=son[x][i];
        dfs_party(y);
        dp[x][0]+=max(dp[y][0],dp[y][1]);
        dp[x][1]+=dp[y][0];
    }
}
\end{lstlisting}
\note{O(\allowbreak{}n*m\textasciicircum{}2)\allowbreak{}，树上背包(\allowbreak{}选课)\allowbreak{}：f[\allowbreak{}x]\allowbreak{}[\allowbreak{}j]\allowbreak{} 表示以 x 为根的子树里选 j 门课的最大收益 先把所有儿子合并上来，最后再把自己塞进去，体现"选子必须先选父" 恰好选 j 门；不可达为 NEG，支持负学分；虚拟根 0 不占名额。}
\begin{lstlisting}
void dfs_knap(int x)
{
    for(int t=0;t<=m;t++)f[x][t]=NEG;
    f[x][0]=0;
    for(int i=0;i<(int)son[x].size();i++)
    {
        int y=son[x][i];
        dfs_knap(y);
        for(int t=m;t>=0;t--)// 倒序，每个儿子只贡献一次
            for(int j=t;j>=0;j--)
                if(f[x][t-j]!=NEG&&f[y][j]!=NEG)f[x][t]=max(f[x][t],f[x][t-j]+f[y][j]);
    }
    if(x!=0)// 虚拟根 0 没有学分
        for(int t=m;t>0;t--)f[x][t]=f[x][t-1]==NEG?NEG:f[x][t-1]+score[x];
}
\end{lstlisting}
\note{O(\allowbreak{}n)\allowbreak{}，树形 dp 求直径（边数）：d1[\allowbreak{}x]\allowbreak{} 是最长向下链，d2[\allowbreak{}x]\allowbreak{} 是次长}
\begin{lstlisting}
void dfs_dia(int x)
{
    d1[x]=0,d2[x]=0;
    for(int i=0;i<(int)son[x].size();i++)
    {
        int y=son[x][i];
        dfs_dia(y);
        int t=d1[y]+1;
        if(t>d1[x])d2[x]=d1[x],d1[x]=t;
        else if(t>d2[x])d2[x]=t;
    }
    dia=max(dia,d1[x]+d2[x]);// 拐点在自己身上的最长路
}
\end{lstlisting}
% SOURCE: 06-动态规划/状压DP.cpp
\subsection{状压DP}

\note{mp[\allowbreak{}i]\allowbreak{}以1表示禁放位置，S/\allowbreak{}cnt枚举合法行状态及数量，dp滚动保存相邻两行状态。另一例px/\allowbreak{}py为坐标、0号为原点，f[\allowbreak{}S]\allowbreak{}[\allowbreak{}u]\allowbreak{}为访问集合S且停在u的距离；两例状态含义独立，不能共用初始化。}
\begin{lstlisting}
const int N=110;
const int K=105;// 每行合法状态数上限(m<=10 时约 60 个)
const int NS=1<<15;// TSP 状态数上限
const int NEG=-0x3f3f3f3f;
int n,m,top;
int mp[N];// 地形压缩，1 表示山地不能放
int S[1<<10],cnt[1<<10];// 合法状态 和它里面 1 的个数
int dp[2][K][K];// 滚动两行，dp[i][j][k]：第 i 行状态 j、第 i-1 行状态 k
double px[20],py[20];// 吃奶酪的坐标，0 号是原点
double f[NS][16];// f[状态][当前点]，走完集合的最小距离
\end{lstlisting}
\note{O(\allowbreak{}log x)\allowbreak{}，二进制里 1 的个数}
\begin{lstlisting}
int popcnt(int x)
{
    int c=0;
    while(x)x&=x-1,c++;
    return c;
}
\end{lstlisting}
\note{O(\allowbreak{}n*top\textasciicircum{}3)\allowbreak{}，炮兵阵地：同行不能间隔 1 或 2 列，上下两行也不能同列}
\begin{lstlisting}
int artillery(int n,int m)
{
    top=0;
    for(int i=0;i<(1<<m);i++)// 预处理一行内的合法状态
        if((i&(i<<1))==0&&(i&(i<<2))==0)S[top]=i,cnt[top++]=popcnt(i);
    for(int j=0;j<top;j++)
        for(int k=0;k<top;k++)dp[0][j][k]=NEG;
    dp[0][0][0]=0;// 第 0 行只有空状态合法
    int ans=0;
    for(int i=1;i<=n;i++)
    {
        int cur=i&1,pre=cur^1;
        for(int j=0;j<top;j++)
            for(int k=0;k<top;k++)dp[cur][j][k]=NEG;
        for(int j=0;j<top;j++)// 本行状态
        {
            if(S[j]&mp[i])continue;// 山地不能放
            for(int k=0;k<top;k++)// 上一行状态
            {
                if(S[j]&S[k])continue;// 同列冲突
                int best=NEG;
                for(int l=0;l<top;l++)// 上上行状态
                    if(dp[pre][k][l]!=NEG&&(S[l]&S[j])==0)best=max(best,dp[pre][k][l]);
                if(best!=NEG)
                {
                    dp[cur][j][k]=best+cnt[j];
                    ans=max(ans,dp[cur][j][k]);
                }
            }
        }
    }
    return ans;
}
\end{lstlisting}
\note{O(\allowbreak{}1)\allowbreak{}，两点距离，i=\allowbreak{}0 表示原点}
\begin{lstlisting}
double dist(int i,int j)
{
    double dx=px[i]-px[j],dy=py[i]-py[j];
    return sqrt(dx*dx+dy*dy);
}
\end{lstlisting}
\note{O(\allowbreak{}2\textasciicircum{}n*n\textasciicircum{}2)\allowbreak{}，吃奶酪：从原点出发走遍所有点，最小总距离 f[\allowbreak{}s]\allowbreak{}[\allowbreak{}i]\allowbreak{} 表示已经走过集合 s、当前停在 i 的最小距离}
\begin{lstlisting}
double tsp(int n)
{
    if(!n)return 0;
    int full=1<<n;
    for(int s=0;s<full;s++)
        for(int i=0;i<n;i++)f[s][i]=1e18;
    for(int i=0;i<n;i++)f[1<<i][i]=dist(0,i+1);// 第一步从原点走到 i
    for(int s=1;s<full;s++)
        for(int i=0;i<n;i++)
            if(f[s][i]<1e18)
                for(int j=0;j<n;j++)
                    if(((s>>j)&1)==0)
                        f[s|(1<<j)][j]=min(f[s|(1<<j)][j],f[s][i]+dist(i+1,j+1));
    double ans=1e18;
    for(int i=0;i<n;i++)ans=min(ans,f[full-1][i]);
    return ans;
}
\end{lstlisting}
% SOURCE: 06-动态规划/数位DP.cpp
\subsection{数位DP}

\topic{最优计数：后续乘零}
\note{同一状态的历史，若后续目标乘非负系数 $w$，应同时维护最优值、最优方案数 \texttt{ways} 与全部可行方案数 \texttt{all}。$w>0$ 延用 \texttt{ways}，$w=0$ 改用 \texttt{all}：旧分数 2、5 各一种，乘零后有两种最优方案。负系数还需维护最小值。分支须无重复，且状态内历史的后续可行性相同；仅方案数取模，目标值不可取模后比较（2025 沈阳 D）。}

\note{dig[\allowbreak{}1]\allowbreak{}是个位，len为位数，tar是目标数字，K为数位和上限。ca/\allowbreak{}sa为计数与次数，g2/\allowbreak{}g3/\allowbreak{}g4供不同约束的记忆化；受上界限制的状态不能直接复用自由状态缓存，改变约束后清空对应缓存。}
\note{最优方案计数陷阱：若下一步将目标值乘非负 w，w=\allowbreak{}0 会让所有可行历史同分。 此时需另存 all（全部可行方案数），不能只用 ways（旧最优方案数）；w>\allowbreak{}0 才沿用 ways。 例如历史得分 2、5，各一种，乘0后最佳为0且有两种。若 w<\allowbreak{}0，还须维护旧最小值。 分支须无重复计数，且同状态历史的后续可行性一致；只对方案数取模，不能把目标值先取模比较。}
\begin{lstlisting}
const int N=25;// 位数
int dig[N],len;// 拆位，dig[1] 是个位
int tar,K;// 目标数字 与 数位和上限
ll ca[N],sa[N];// 示例1 的记搜：数的个数 与 数字出现次数
ll g2[N][12];// 示例2 不含 4
ll g3[N][12];// 示例3 相邻数字差 >= 2
ll g4[N][205];// 示例4 数位和 <= K
\end{lstlisting}
\note{拆位，dig[\allowbreak{}1]\allowbreak{} 是个位}
\begin{lstlisting}
void get_dig(ll x)
{
    len=0;
    if(x==0){dig[++len]=0;return;}
    while(x>0)dig[++len]=x%10,x/=10;
}
\end{lstlisting}
\note{示例1：统计数字 tar 在 1.\allowbreak{}.\allowbreak{}x 里出现的总次数（不把 0 记为一位数字） 返回 (\allowbreak{}合法数的个数,\allowbreak{} 这些数里 tar 的出现次数)\allowbreak{} 关键边界：记忆化只在 !limit \&\& !lead 时记录和取用}
\begin{lstlisting}
pair<ll,ll> dfs_cnt(int pos,bool limit,bool lead)
{
    if(pos==0)return make_pair(1,0);// 数完了，算 1 个补全
    if(!limit&&!lead&&ca[pos]!=-1)return make_pair(ca[pos],sa[pos]);
    int up=limit?dig[pos]:9;
    ll cnt=0,sum=0;
    for(int i=0;i<=up;i++)
    {
        pair<ll,ll> t=dfs_cnt(pos-1,limit&&i==up,lead&&i==0);
        cnt+=t.first;
        if(i==tar&&!(lead&&i==0))sum+=t.first;// 这一位贡献 t.first 次
        sum+=t.second;
    }
    if(!limit&&!lead)ca[pos]=cnt,sa[pos]=sum;// 贴上界/前导零的状态不能记
    return make_pair(cnt,sum);
}
\end{lstlisting}
\note{示例2：统计 0.\allowbreak{}.\allowbreak{}x 里十进制表示不含 4 的数的个数 不含 4 跟前导零无关，所以不用 lead 也可以直接记}
\begin{lstlisting}
ll dfs_no4(int pos,bool limit)
{
    if(pos==0)return 1;
    if(!limit&&g2[pos][0]!=-1)return g2[pos][0];
    int up=limit?dig[pos]:9;
    ll res=0;
    for(int i=0;i<=up;i++)
    {
        if(i==4)continue;
        res+=dfs_no4(pos-1,limit&&i==up);
    }
    if(!limit)g2[pos][0]=res;
    return res;
}
\end{lstlisting}
\note{示例3：统计 0.\allowbreak{}.\allowbreak{}x 里相邻两位数字差 >\allowbreak{}=\allowbreak{} 2 的数的个数(\allowbreak{}windy 数)\allowbreak{} pre=\allowbreak{}10 表示前面还没有有效数字(\allowbreak{}还在前导零阶段)\allowbreak{}}
\begin{lstlisting}
ll dfs_windy(int pos,int pre,bool limit,bool lead)
{
    if(pos==0)return 1;
    if(!limit&&!lead&&g3[pos][pre]!=-1)return g3[pos][pre];
    int up=limit?dig[pos]:9;
    ll res=0;
    for(int i=0;i<=up;i++)
    {
        if(!(lead&&i==0)&&pre!=10&&abs(i-pre)<2)continue;// 相邻差不够就剪掉
        res+=dfs_windy(pos-1,(lead&&i==0)?10:i,limit&&i==up,lead&&i==0);
    }
    if(!limit&&!lead)g3[pos][pre]=res;
    return res;
}
\end{lstlisting}
\note{示例4：统计 0.\allowbreak{}.\allowbreak{}x 里数位和 <\allowbreak{}=\allowbreak{} K 的数的个数 前导零对 digit sum 没有影响，所以状态只需要 (\allowbreak{}pos,\allowbreak{}sum)\allowbreak{}}
\begin{lstlisting}
ll dfs_sumk(int pos,int sum,bool limit)
{
    if(sum>K)return 0;
    if(pos==0)return 1;
    if(!limit&&g4[pos][sum]!=-1)return g4[pos][sum];
    int up=limit?dig[pos]:9;
    ll res=0;
    for(int i=0;i<=up;i++)res+=dfs_sumk(pos-1,sum+i,limit&&i==up);
    if(!limit)g4[pos][sum]=res;
    return res;
}
\end{lstlisting}
\note{求 [\allowbreak{}0,\allowbreak{}x]\allowbreak{}，x<\allowbreak{}0 时返回 0}
\begin{lstlisting}
ll solve_cnt(ll x,int dd)
{
    if(x<0)return 0;
    tar=dd,get_dig(x);
    memset(ca,-1,sizeof(ca)),memset(sa,-1,sizeof(sa));
    return dfs_cnt(len,true,true).second;
}
ll solve_no4(ll x)
{
    if(x<0)return 0;
    get_dig(x);
    memset(g2,-1,sizeof(g2));
    return dfs_no4(len,true);
}
ll solve_windy(ll x)
{
    if(x<0)return 0;
    get_dig(x);
    memset(g3,-1,sizeof(g3));
    return dfs_windy(len,10,true,true);
}
ll solve_sumk(ll x,int k)
{
    if(x<0)return 0;
    K=k,get_dig(x);
    memset(g4,-1,sizeof(g4));
    return dfs_sumk(len,0,true);
}
\end{lstlisting}
% SOURCE: 06-动态规划/概率期望DP.cpp
\subsection{概率期望DP}

\textbf{自环移项、线性性、尾和}
\[
E=c+pE+\sum_jq_jE_j
\quad\Longrightarrow\quad E=\frac{c+\sum_jq_jE_j}{1-p}.
\]
\[
\mathbb E\Bigl[\sum_iX_i\Bigr]=\sum_i\mathbb E[X_i],\qquad
\mathbb E[T]=\sum_{k\ge0}\Pr(T>k).
\]
\note{第一式要求 $p<1$ 且期望有限；尾和适用于非负整数随机变量。
无需独立性即可使用期望线性性。}
\note{正推代码是截断近似；当前存活概率小并不保证期望尾项小。若未终止状态的剩余期望统一不超过 $B$，截断到 $K$ 后的误差才可界为 $B\Pr(T>K)$。终点可能永不到达时，须先判断期望是否有限。}
\note{若 $S$ 是已取得的项目集合，每步等概率抽取 $n$ 项中的一项：}
\[
E[S]=\frac{n+\sum_{i\notin S}E[S\cup\{i\}]}{n-|S|},
\quad E[\text{完成状态}]=0.
\]
\note{稳态 move-to-front 模型中，若 $p_i+p_j>0$，
一对项目的期望逆序贡献为 $p_ip_j/(p_i+p_j)$；零概率对贡献为 0。}

\note{E为状态期望，P/\allowbreak{}np为本步/\allowbreak{}下一步概率；A、bvec、sol分别是系数矩阵、常数项和解，下标0起，gauss(\allowbreak{}sz)\allowbreak{}会改写前两者。st/\allowbreak{}ns是二维概率滚动表，行列不超过15。入口传规模和概率等标量，每次重建相应状态；正推结果是截断近似。}
\begin{lstlisting}
const int N=1005;
int n;
double p;// 每步向目标前进的概率
double E[N];// 期望步数
double P[N];// 当前时刻各位置的概率
double A[N][N],bvec[N],sol[N];// 线性方程组用
double np[N],st[15][15],ns[15][15];
\end{lstlisting}
\note{O(\allowbreak{}n\textasciicircum{}3)\allowbreak{}，高斯消元解线性方程组，用来和期望 dp 对拍}
\begin{lstlisting}
void gauss(int sz)
{
    for(int i=0;i<sz;i++)
    {
        int piv=i;
        for(int j=i;j<sz;j++)
            if(fabs(A[j][i])>fabs(A[piv][i]))piv=j;
        if(fabs(A[piv][i])<1e-14)continue;
        for(int j=0;j<sz;j++)swap(A[i][j],A[piv][j]);
        swap(bvec[i],bvec[piv]);
        double d=A[i][i];
        for(int j=i;j<sz;j++)A[i][j]/=d;
        bvec[i]/=d;
        for(int j=0;j<sz;j++)
            if(j!=i&&fabs(A[j][i])>1e-15)
            {
                double f=A[j][i];
                for(int k=i;k<sz;k++)A[j][k]-=f*A[i][k];
                bvec[j]-=f*bvec[i];
            }
    }
    for(int i=0;i<sz;i++)sol[i]=bvec[i];
}
\end{lstlisting}
\note{=\allowbreak{}=\allowbreak{}=\allowbreak{}=\allowbreak{}=\allowbreak{} 期望 dp 逆推 =\allowbreak{}=\allowbreak{}=\allowbreak{}=\allowbreak{}=\allowbreak{} 数轴上从 0 出发走到 n，每步以 p 前进、1-\allowbreak{}p 后退（在 0 处退不动） E[\allowbreak{}n]\allowbreak{}=\allowbreak{}0，E[\allowbreak{}i]\allowbreak{}=\allowbreak{}1+\allowbreak{}p*E[\allowbreak{}i+\allowbreak{}1]\allowbreak{}+\allowbreak{}(\allowbreak{}1-\allowbreak{}p)\allowbreak{}*E[\allowbreak{}i-\allowbreak{}1]\allowbreak{}，i=\allowbreak{}0 时后退为原地不动，最后一项为 (\allowbreak{}1-\allowbreak{}p)\allowbreak{}*E[\allowbreak{}0]\allowbreak{} 这个方程里 E[\allowbreak{}i]\allowbreak{} 左右都有，所以不能简单地从后往前推，要解三对角方程（p=\allowbreak{}0.\allowbreak{}5 时有公式 n(\allowbreak{}n+\allowbreak{}1)\allowbreak{}-\allowbreak{}i(\allowbreak{}i+\allowbreak{}1)\allowbreak{}） 有限反射区间上 p>\allowbreak{}0 时期望有限；p=\allowbreak{}0 且 n>\allowbreak{}0 时返回 -\allowbreak{}1}
\begin{lstlisting}
double exp_hit(int n,double p)
{
    if(n==0)return 0;
    if(p<=0)return -1;
    if(p>=1-1e-12)return (double)n;// 每步必进
    int sz=n;// 未知数 E[0..n-1]
    for(int i=0;i<sz;i++)
        for(int j=0;j<sz;j++)A[i][j]=0;
    for(int i=0;i<sz;i++)
    {
        A[i][i]=i==0?p:1;
        if(i+1<sz)A[i][i+1]-=p;// E[i] 里含 p*E[i+1]
        if(i-1>=0)A[i][i-1]-=(1-p);// E[i] 里含 (1-p)*E[i-1]
        bvec[i]=1;
    }
    gauss(sz);
    return sol[0];
}
\end{lstlisting}
\note{正推法：把"某个时刻还在游走的概率"记作 cost，期望 =\allowbreak{} sum\_\{t>\allowbreak{}=\allowbreak{}0\} P(\allowbreak{}T>\allowbreak{}t)\allowbreak{} 这里 T 是首次到达 n 的时刻，P(\allowbreak{}T>\allowbreak{}t)\allowbreak{}=\allowbreak{}P(\allowbreak{}X\_t<\allowbreak{}n)\allowbreak{} 下方正推函数均为截断近似：固定步数或当前存活概率小，不等于期望尾项误差已小。 剩余误差是 sum(\allowbreak{}t>\allowbreak{}=\allowbreak{}K)\allowbreak{}P(\allowbreak{}T>\allowbreak{}t)\allowbreak{}；若所有未终止状态的剩余期望<\allowbreak{}=\allowbreak{}B，才有尾项<\allowbreak{}=\allowbreak{}B*P(\allowbreak{}T>\allowbreak{}K)\allowbreak{}。 若终点可能永不到达，期望可为无穷，不能用有限截断结果冒充完整期望。}
\begin{lstlisting}
double exp_forward_walk(int n,double p)
{
    for(int i=0;i<=n;i++)P[i]=0;
    P[0]=1;
    double ans=0;
    int K=300000;
    for(int t=0;t<K;t++)
    {
        double cost=0;
        for(int i=0;i<n;i++)cost+=P[i];// 还没到 n
        ans+=cost;
        if(cost<1e-14)break;
        for(int i=0;i<=n;i++)np[i]=0;
        for(int i=0;i<n;i++)
        {
            np[i+1]+=P[i]*p;// 前进
            if(i-1>=0)np[i-1]+=P[i]*(1-p);// 后退，0 处退不动所以不加
            else np[0]+=P[i]*(1-p);
        }
        np[n]+=P[n];
        for(int i=0;i<=n;i++)P[i]=np[i];
    }
    return ans;
}
\end{lstlisting}
\note{正推法第二个例子：二维格子里走 每次 a 概率向右、b 概率向上、其余概率不动，走到 (\allowbreak{}n-\allowbreak{}1,\allowbreak{}m-\allowbreak{}1)\allowbreak{} 停，求期望步数}
\begin{lstlisting}
double exp_forward_grid(int n,int m,double a,double b)
{
    for(int i=0;i<n;i++)
        for(int j=0;j<m;j++)st[i][j]=0;
    st[0][0]=1;
    double ans=0;
    int K=200000;
    for(int t=0;t<K;t++)
    {
        double cost=0;
        for(int i=0;i<n;i++)
            for(int j=0;j<m;j++)
                if(!(i==n-1&&j==m-1))cost+=st[i][j];// 此刻还没走到终点
        ans+=cost;
        if(cost<1e-14)break;
        for(int i=0;i<n;i++)
            for(int j=0;j<m;j++)ns[i][j]=0;
        for(int i=0;i<n;i++)
            for(int j=0;j<m;j++)
            {
                if(i==n-1&&j==m-1){ns[i][j]+=st[i][j];continue;}
                double v=st[i][j];
                ns[i][j]+=v*(1-a-b);// 原地不动
                ns[min(i+1,n-1)][j]+=v*a;// 向右，到边界就停在终点方向
                ns[i][min(j+1,m-1)]+=v*b;// 向上
            }
        for(int i=0;i<n;i++)
            for(int j=0;j<m;j++)st[i][j]=ns[i][j];
    }
    return ans;
}
\end{lstlisting}
\note{顺带一个经典期望：n 种卡等概率抽，集齐 n 种的期望次数 =\allowbreak{} n*H\_n}
\begin{lstlisting}
double coupon_expected(int n)
{
    double ans=0;
    for(int i=1;i<=n;i++)ans+=(double)n/i;
    return ans;
}
\end{lstlisting}
\note{自环移项：E=\allowbreak{}c+\allowbreak{}p*E+\allowbreak{}sum(\allowbreak{}q[\allowbreak{}i]\allowbreak{}*E[\allowbreak{}i]\allowbreak{})\allowbreak{} =\allowbreak{}>\allowbreak{} E=\allowbreak{}(\allowbreak{}c+\allowbreak{}sum(\allowbreak{}q[\allowbreak{}i]\allowbreak{}*E[\allowbreak{}i]\allowbreak{})\allowbreak{})\allowbreak{}/\allowbreak{}(\allowbreak{}1-\allowbreak{}p)\allowbreak{}。 p<\allowbreak{}1、期望有限；后继互相依赖时仍需解方程组。抽 n 种物品、已有集合 S： E[\allowbreak{}S]\allowbreak{}=\allowbreak{}(\allowbreak{}n+\allowbreak{}sum(\allowbreak{}i不在S)\allowbreak{}E[\allowbreak{}S并\{i\}]\allowbreak{})\allowbreak{}/\allowbreak{}(\allowbreak{}n-\allowbreak{}|S|)\allowbreak{}，完成状态为 0，分母是“未抽过的种数”。 指示变量：E[\allowbreak{}sum X[\allowbreak{}i]\allowbreak{}]\allowbreak{}=\allowbreak{}sum E[\allowbreak{}X[\allowbreak{}i]\allowbreak{}]\allowbreak{}，无需独立；非负整数 T：E[\allowbreak{}T]\allowbreak{}=\allowbreak{}sum(\allowbreak{}t>\allowbreak{}=\allowbreak{}0)\allowbreak{}P(\allowbreak{}T>\allowbreak{}t)\allowbreak{}。 按概率 p[\allowbreak{}i]\allowbreak{} 独立抽元素并移到表头，稳态期望跨越数=\allowbreak{}sum(\allowbreak{}i!=\allowbreak{}j)\allowbreak{}p[\allowbreak{}i]\allowbreak{}*p[\allowbreak{}j]\allowbreak{}/\allowbreak{}(\allowbreak{}p[\allowbreak{}i]\allowbreak{}+\allowbreak{}p[\allowbreak{}j]\allowbreak{})\allowbreak{}。 两者概率都为 0 的项记 0；不要为整个排列建马尔可夫状态。}
% SOURCE: 06-动态规划/单调队列优化DP.cpp
\subsection{单调队列优化DP}

\note{n为长度，k或[\allowbreak{}L,\allowbreak{}R]\allowbreak{}给出合法前驱距离，a为输入，f为状态，res为输出缓冲，q存前驱下标。先按距离移除过期候选，再按转移值淘汰劣候选；每个示例分别初始化队首、队尾和不可达状态。}
\begin{lstlisting}
const int N=100005;
const int INF=0x3f3f3f3f;
int n,k,L,R;
int a[N],f[N],res[N];
int q[N];// 单调队列，存下标（也可以用 deque<int>，见 P1886 写法）
\end{lstlisting}
\note{O(\allowbreak{}n)\allowbreak{}，滑动窗口最大值 三步：入队前把队尾比它差的弹掉 -\allowbreak{}>\allowbreak{} 入队 -\allowbreak{}>\allowbreak{} 把越界的队首弹掉}
\begin{lstlisting}
void window_max()
{
    int head=0,tail=0;
    for(int i=1;i<=n;i++)
    {
        while(head<tail&&a[q[tail-1]]<=a[i])tail--;
        q[tail++]=i;
        while(q[head]<i-k+1)head++;// 窗口左端是 i-k+1
        if(i>=k)res[i-k+1]=a[q[head]];
    }
}
\end{lstlisting}
\note{O(\allowbreak{}n)\allowbreak{}，滑动窗口最小值，队列改成单调递增}
\begin{lstlisting}
void window_min()
{
    int head=0,tail=0;
    for(int i=1;i<=n;i++)
    {
        while(head<tail&&a[q[tail-1]]>=a[i])tail--;
        q[tail++]=i;
        while(q[head]<i-k+1)head++;
        if(i>=k)res[i-k+1]=a[q[head]];
    }
}
\end{lstlisting}
\note{O(\allowbreak{}n)\allowbreak{}，单调队列优化转移：f[\allowbreak{}i]\allowbreak{}=\allowbreak{}a[\allowbreak{}i]\allowbreak{}+\allowbreak{}max(\allowbreak{}f[\allowbreak{}j]\allowbreak{})\allowbreak{}，i-\allowbreak{}R<\allowbreak{}=\allowbreak{}j<\allowbreak{}=\allowbreak{}i-\allowbreak{}L f[\allowbreak{}i]\allowbreak{} 表示以 i 结尾的最大得分；j 也可以不选(\allowbreak{}原地起步)\allowbreak{} 处理 i 时先把下标 i-\allowbreak{}L 入队，再把 <\allowbreak{} i-\allowbreak{}R 的弹掉，窗口正好是 [\allowbreak{}i-\allowbreak{}R,\allowbreak{}i-\allowbreak{}L]\allowbreak{}}
\begin{lstlisting}
int jump_max_score()
{
    int head=0,tail=0,ans=-INF;
    for(int i=1;i<=n;i++)
    {
        int add=i-L;
        if(add>=1)
        {
            while(head<tail&&f[q[tail-1]]<=f[add])tail--;
            q[tail++]=add;
        }
        while(head<tail&&q[head]<i-R)head++;
        f[i]=a[i];
        if(head<tail)f[i]=max(f[i],f[q[head]]+a[i]);
        ans=max(ans,f[i]);
    }
    return ans;
}
\end{lstlisting}
% SOURCE: 06-动态规划/斜率优化DP.cpp
\subsection{斜率优化DP}

\note{a、s为原序列与前缀量，f是最优值；x、y是候选点坐标，A、B、C是展开后的系数，q存凸壳下标。L是具体模型参数。套用前验证插入斜率、查询横坐标的单调方向，并用足够宽的乘法比较。}
\begin{lstlisting}
const int N=100005;
const ll INF=4e18;
ll n,L;
ll a[N],s[N],f[N],x[N],y[N],A[N],B[N],C[N];
int q[N];// 下凸壳用的单调队列，存下标
\end{lstlisting}
\note{经典转移（P3195 玩具装箱）： f[\allowbreak{}i]\allowbreak{}=\allowbreak{}min\_\{j<\allowbreak{}i\} f[\allowbreak{}j]\allowbreak{}+\allowbreak{}(\allowbreak{}s[\allowbreak{}i]\allowbreak{}-\allowbreak{}s[\allowbreak{}j]\allowbreak{}-\allowbreak{}L)\allowbreak{}\textasciicircum{}2 ，其中 s[\allowbreak{}i]\allowbreak{} 是 c 的前缀和、L 是常数 展开成只和 j 有关的线性函数： 令 X=\allowbreak{}s[\allowbreak{}j]\allowbreak{}、Y=\allowbreak{}f[\allowbreak{}j]\allowbreak{}+\allowbreak{}s[\allowbreak{}j]\allowbreak{}\textasciicircum{}2、k=\allowbreak{}2*(\allowbreak{}s[\allowbreak{}i]\allowbreak{}-\allowbreak{}L)\allowbreak{}、C=\allowbreak{}(\allowbreak{}s[\allowbreak{}i]\allowbreak{}-\allowbreak{}L)\allowbreak{}\textasciicircum{}2 则 f[\allowbreak{}i]\allowbreak{}=\allowbreak{}min\_j (\allowbreak{}Y -\allowbreak{} k*X)\allowbreak{} +\allowbreak{} C ，即把直线 y=\allowbreak{}kx 从下方去切点集，取最小截距 a>\allowbreak{}=\allowbreak{}0；前缀和、平方、f/\allowbreak{}y/\allowbreak{}k 都须放入 ll。 本题 s 单调不减 -\allowbreak{}>\allowbreak{} X 递增（加点单调），k 递增（询问单调），可以单调队列 O(\allowbreak{}n)\allowbreak{}}
\begin{lstlisting}
ll sq(ll v){return v*v;}
\end{lstlisting}
\note{单调队列维护下凸壳，O(\allowbreak{}n)\allowbreak{} 斜率优化 叉积 cross(\allowbreak{}o,\allowbreak{}a,\allowbreak{}b)\allowbreak{}=\allowbreak{}(\allowbreak{}x[\allowbreak{}a]\allowbreak{}-\allowbreak{}x[\allowbreak{}o]\allowbreak{})\allowbreak{}*(\allowbreak{}y[\allowbreak{}b]\allowbreak{}-\allowbreak{}y[\allowbreak{}o]\allowbreak{})\allowbreak{}-\allowbreak{}(\allowbreak{}y[\allowbreak{}a]\allowbreak{}-\allowbreak{}y[\allowbreak{}o]\allowbreak{})\allowbreak{}*(\allowbreak{}x[\allowbreak{}b]\allowbreak{}-\allowbreak{}x[\allowbreak{}o]\allowbreak{})\allowbreak{} cross(\allowbreak{}o,\allowbreak{}a,\allowbreak{}b)\allowbreak{}<\allowbreak{}=\allowbreak{}0 表示 a 在 ob 连线之上，下凸壳要把 a 弹掉}
\begin{lstlisting}
ll slope_dp()
{
    s[0]=f[0]=0;
    for(int i=1;i<=n;i++)s[i]=s[i-1]+a[i];
    int head=0,tail=0;
    x[0]=0,y[0]=0;// f[0]=0
    q[tail++]=0;
    for(int i=1;i<=n;i++)
    {
        ll k=2*(s[i]-L);
        while(head+1<tail&&(__int128)y[q[head+1]]-y[q[head]]<=(__int128)k*(x[q[head+1]]-x[q[head]]))head++;
        int j=q[head];
        f[i]=(ll)((__int128)y[j]-(__int128)k*x[j]+(__int128)(s[i]-L)*(s[i]-L));
        x[i]=s[i],y[i]=f[i]+sq(s[i]);
        if(tail>head&&x[q[tail-1]]==x[i])
        {
            if(y[q[tail-1]]<=y[i])continue;
            --tail;
        }
        while(head+1<tail)
        {
            int o=q[tail-2],p=q[tail-1];
            if((__int128)(x[p]-x[o])*((__int128)y[i]-y[o])-((__int128)y[p]-y[o])*(x[i]-x[o])<=0)tail--;
            else break;
        }
        q[tail++]=i;
    }
    return f[n];
}
\end{lstlisting}
\note{把转移式乘开成 A[\allowbreak{}i]\allowbreak{}*B[\allowbreak{}j]\allowbreak{}+\allowbreak{}C[\allowbreak{}i]\allowbreak{} 的标准形式，方便套李超线段树/\allowbreak{}一般凸壳 f[\allowbreak{}i]\allowbreak{}=\allowbreak{}min\_j (\allowbreak{}f[\allowbreak{}j]\allowbreak{}+\allowbreak{}s[\allowbreak{}j]\allowbreak{}\textasciicircum{}2+\allowbreak{}2*L*s[\allowbreak{}j]\allowbreak{} -\allowbreak{} 2*s[\allowbreak{}i]\allowbreak{}*s[\allowbreak{}j]\allowbreak{})\allowbreak{} +\allowbreak{} s[\allowbreak{}i]\allowbreak{}\textasciicircum{}2 +\allowbreak{} L\textasciicircum{}2 -\allowbreak{} 2*L*s[\allowbreak{}i]\allowbreak{}}
\begin{lstlisting}
void build_ab()
{
    s[0]=f[0]=0;
    for(int i=1;i<=n;i++)s[i]=s[i-1]+a[i];
    for(int j=0;j<=n;j++)A[j]=-2*s[j],B[j]=f[j]+s[j]*s[j]+2*L*s[j];// 斜率、截距
    for(int i=1;i<=n;i++)C[i]=s[i]*s[i]+L*L-2*L*s[i];// 与 j 无关的常数
}
\end{lstlisting}
% SOURCE: 06-动态规划/矩阵优化DP.cpp
\subsection{矩阵优化DP}

\note{c[\allowbreak{}0.\allowbreak{}.\allowbreak{}k-\allowbreak{}1]\allowbreak{}是递推系数，f[\allowbreak{}0.\allowbreak{}.\allowbreak{}k-\allowbreak{}1]\allowbreak{}是初值，f(\allowbreak{}t)\allowbreak{}=\allowbreak{}sum(\allowbreak{}c[\allowbreak{}j]\allowbreak{}*f(\allowbreak{}t-\allowbreak{}1-\allowbreak{}j)\allowbreak{})\allowbreak{}；solve(\allowbreak{}n)\allowbreak{}返回模MOD的第n项，n>\allowbreak{}=\allowbreak{}0，k>\allowbreak{}0。矩阵乘法和快速幂保留对象参数，以区分底数、累积结果与乘积。}
\begin{lstlisting}
const ll MOD=1000000007;
struct Matrix
{
    int n;
    vector<vector<ll> > a;
    Matrix(int n_,bool unit=false):n(n_),a(n_,vector<ll>(n_,0))
    {
        if(unit)for(int i=0;i<n;i++)a[i][i]=1;
    }
};
Matrix mul(const Matrix &a,const Matrix &b)
{
    Matrix c(a.n);
    for(int i=0;i<a.n;i++)
        for(int k=0;k<a.n;k++)
            for(int j=0;j<a.n;j++)c.a[i][j]=(c.a[i][j]+a.a[i][k]*b.a[k][j])%MOD;
    return c;
}
Matrix qpow(Matrix a,ll n)
{
    Matrix r(a.n,true);
    while(n)
    {
        if(n&1)r=mul(r,a);
        a=mul(a,a),n>>=1;
    }
    return r;
}
\end{lstlisting}
\note{O(\allowbreak{}k\textasciicircum{}3 log n)\allowbreak{}，f(\allowbreak{}n)\allowbreak{}=\allowbreak{}c[\allowbreak{}0]\allowbreak{}f(\allowbreak{}n-\allowbreak{}1)\allowbreak{}+\allowbreak{}.\allowbreak{}.\allowbreak{}.\allowbreak{}+\allowbreak{}c[\allowbreak{}k-\allowbreak{}1]\allowbreak{}f(\allowbreak{}n-\allowbreak{}k)\allowbreak{}，初值 f(\allowbreak{}0.\allowbreak{}.\allowbreak{}k-\allowbreak{}1)\allowbreak{}}
\begin{lstlisting}
vector<ll> c,f;
ll solve(ll n)
{
    int k=c.size();
    assert(k>0&&f.size()==c.size()&&n>=0);
    if(n<k)return (f[n]%MOD+MOD)%MOD;
    Matrix a(k);
    for(int j=0;j<k;j++)a.a[0][j]=(c[j]%MOD+MOD)%MOD;
    for(int i=1;i<k;i++)a.a[i][i-1]=1;
    a=qpow(a,n-k+1);
    ll ans=0;
    for(int i=0;i<k;i++)ans=(ans+a.a[0][i]*((f[k-1-i]%MOD+MOD)%MOD))%MOD;
    return ans;
}
\end{lstlisting}
% SOURCE: 06-动态规划/SOS子集与超集和.cpp
\subsection{SOS子集与超集和}

\note{f[\allowbreak{}S]\allowbreak{}为每个掩码的原权，原地变换后为子集/\allowbreak{}超集聚合值；inverse撤销对应变换。两个输入与按集合大小分层的卷积缓冲是不同运算对象，保留参数；整数位宽与状态数量2\textasciicircum{}n分别检查。}
\note{全部掩码的聚合转移统一放在这里；枚举一个掩码的子集见第一章位运算。长度必须为2\textasciicircum{}k，普通求和的所有中间值须能容纳于ll。}
\begin{lstlisting}

\end{lstlisting}
\topic{子集和及逆变换}
\note{subset\_sum原地把f[\allowbreak{}S]\allowbreak{}变为所有T包含于S的原值之和，逐位从不含该位的状态传入。inverse=\allowbreak{}true逐位减回，恢复原数组。时间 $O(k2^k)$，额外空间常数；不要先取不适当的模再反推原整数值。}
\begin{lstlisting}
void subset_sum(vector<ll> &f,bool inverse=false)
{
    int n=f.size();
    assert(n>0&&(n&(n-1))==0);
    for(int b=1;b<n;b*=2)
        for(int s=0;s<n;s++)if(s&b)
            {if(inverse)f[s]-=f[s^b];else f[s]+=f[s^b];}
}
\end{lstlisting}
\topic{超集和及逆变换}
\note{superset\_sum把f[\allowbreak{}S]\allowbreak{}变为所有包含S的T的原值之和，从含该位的状态传向不含该位的状态。时间和空间同上，inverse=\allowbreak{}true恢复原数组。}
\begin{lstlisting}
void superset_sum(vector<ll> &f,bool inverse=false)
{
    int n=f.size();
    assert(n>0&&(n&(n-1))==0);
    for(int b=1;b<n;b*=2)
        for(int s=0;s<n;s++)if(!(s&b))
            {if(inverse)f[s]-=f[s|b];else f[s]+=f[s|b];}
}
\end{lstlisting}
\topic{不交子集卷积}
\note{目标 $h[S]=\sum_{T\subseteq S} f[T]g[S\setminus T]$。先按popcount分层，子集变换后逐层卷积，再逆变换提取第|S|层，保证两部分不相交。时间 $O(k^2 2^k)$，空间 $O(k2^k)$；输入等长，允许负值，模p>\allowbreak{}0，不要求p为素数。普通SOS只做聚合，不能直接代替这个乘积转移。}
\begin{lstlisting}
vector<ll> subset_convolution(const vector<ll> &a,const vector<ll> &b,ll p)
{
    int n=a.size();assert(n>0&&(n&(n-1))==0&&b.size()==a.size()&&p>0);
    int k=__builtin_ctz((unsigned)n);
    vector<vector<ll>> f(k+1,vector<ll>(n)),g=f,h=f;
    for(int s=0;s<n;s++){int c=__builtin_popcount((unsigned)s);f[c][s]=(a[s]%p+p)%p;g[c][s]=(b[s]%p+p)%p;}
    auto transform=[&](vector<vector<ll>> &v,bool inv)
    {
        for(int c=0;c<=k;c++)for(int bit=1;bit<n;bit*=2)for(int s=0;s<n;s++)if(s&bit)
        {__int128 z=(__int128)v[c][s]+(inv?-(__int128)v[c][s^bit]:v[c][s^bit]);v[c][s]=(z%p+p)%p;}
    };
    transform(f,false);transform(g,false);
    for(int c=0;c<=k;c++)for(int j=0;j<=c;j++)for(int s=0;s<n;s++)h[c][s]=(h[c][s]+(__int128)f[j][s]*g[c-j][s])%p;
    transform(h,true);vector<ll> ans(n);for(int s=0;s<n;s++)ans[s]=h[__builtin_popcount((unsigned)s)][s];return ans;
}
\end{lstlisting}
% SOURCE: 06-动态规划/四边形不等式与决策单调性.cpp
\subsection{四边形不等式与决策单调性}

\note{a、s为输入及前缀和，w是区间代价，fd是DP值，opt是最优决策。K为分组数，dq、beg、hd、tl维护在线决策区间。分治优化读上一层，Knuth读更短区间；只有证明各自单调条件后才能使用。}
\begin{lstlisting}
const int N=205,M=25;
const ll INF=4000000000000000000LL;
int n,K;
ll a[N],s[N],w[N][N],fd[N][N];
int opt[N][N],dq[N],beg[N],hd,tl;
\end{lstlisting}
\note{非负 a，w[\allowbreak{}l]\allowbreak{}[\allowbreak{}r]\allowbreak{}=\allowbreak{}(\allowbreak{}s[\allowbreak{}r]\allowbreak{}-\allowbreak{}s[\allowbreak{}l-\allowbreak{}1]\allowbreak{})\allowbreak{}\ensuremath{^2}；分成 k 个非空连续段。 fd[\allowbreak{}k]\allowbreak{}[\allowbreak{}i]\allowbreak{}=\allowbreak{}min\_\{k-\allowbreak{}1<\allowbreak{}=\allowbreak{}j<\allowbreak{}i\} fd[\allowbreak{}k-\allowbreak{}1]\allowbreak{}[\allowbreak{}j]\allowbreak{}+\allowbreak{}w[\allowbreak{}j+\allowbreak{}1]\allowbreak{}[\allowbreak{}i]\allowbreak{}。 要先证明决策单调性；不能仅凭样例或“区间代价”就套。 数值及加法须小于 INF；O(\allowbreak{}K n log n)\allowbreak{}。}
\begin{lstlisting}
void solve_layer(int k,int l,int r,int optl,int optr)
{
    if(l>r)return;
    int mid=(l+r)/2,best=optl;
    ll val=INF;
    for(int j=optl;j<=min(mid-1,optr);j++)
        if(fd[k-1][j]+w[j+1][mid]<val)
            val=fd[k-1][j]+w[j+1][mid],best=j;
    fd[k][mid]=val,opt[k][mid]=best;
    solve_layer(k,l,mid-1,optl,best);
    solve_layer(k,mid+1,r,best,optr);
}
void dnc_partition()
{
    for(int i=1;i<=n;i++)fd[1][i]=w[1][i],opt[1][i]=0;
    for(int k=2;k<=min(K,n);k++)solve_layer(k,k,n,k-1,n-1);
}
\end{lstlisting}
\note{单层：fd[\allowbreak{}2]\allowbreak{}[\allowbreak{}i]\allowbreak{}=\allowbreak{}min\_\{1<\allowbreak{}=\allowbreak{}j<\allowbreak{}i\} fd[\allowbreak{}1]\allowbreak{}[\allowbreak{}j]\allowbreak{}+\allowbreak{}w[\allowbreak{}j+\allowbreak{}1]\allowbreak{}[\allowbreak{}i]\allowbreak{}。 前提：新决策胜过旧决策的位置形成后缀，最优决策单调不减。 dq 存决策，beg 存它开始胜出的下标；每个决策进出一次，O(\allowbreak{}n log n)\allowbreak{}。}
\begin{lstlisting}
void mq_layer()
{
    if(n<2)return;
    auto value=[&](int j,int i){return fd[1][j]+w[j+1][i];};
    hd=0,tl=1,dq[0]=1,beg[0]=2;
    for(int i=2;i<=n;i++)
    {
        while(hd+1<tl&&beg[hd+1]<=i)++hd;
        fd[2][i]=value(dq[hd],i);
        if(i==n)break;
        while(hd<tl&&value(i,max(i+1,beg[tl-1]))<=value(dq[tl-1],max(i+1,beg[tl-1])))--tl;
        if(hd==tl){dq[tl]=i,beg[tl++]=i+1; continue;}
        int old=dq[tl-1];
        if(value(i,n)>value(old,n))continue;
        int l=max(i+1,beg[tl-1]),rr=n;
        while(l<rr)
        {
            int mid=(l+rr)/2;
            if(value(i,mid)<=value(old,mid))rr=mid;
            else l=mid+1;
        }
        dq[tl]=i,beg[tl++]=l;
    }
}
\end{lstlisting}
% SOURCE: 06-动态规划/SlopeTrick.cpp
\subsection{SlopeTrick}

\note{对象两堆保存凸分段线性函数的左右断点，min\_value为最小值，偏移量支持整体平移。add\_abs加入绝对值项，prefix\_min/\allowbreak{}suffix\_min做单侧最小化，shift(\allowbreak{}a,\allowbreak{}b)\allowbreak{}限制前驱区间；相同断点允许重复。}
\note{凸分段线性函数 f(\allowbreak{}x)\allowbreak{}=\allowbreak{}min\_value +\allowbreak{} \ensuremath{\Sigma} max(\allowbreak{}l-\allowbreak{}x,\allowbreak{}0)\allowbreak{} +\allowbreak{} \ensuremath{\Sigma} max(\allowbreak{}x-\allowbreak{}r,\allowbreak{}0)\allowbreak{}，整数断点、单位斜率。 两堆维护拐点，add\_abs(\allowbreak{}a)\allowbreak{} 加 |x-\allowbreak{}a|；add\_left/\allowbreak{}right 分别加 max(\allowbreak{}a-\allowbreak{}x,\allowbreak{}0)\allowbreak{}/\allowbreak{}max(\allowbreak{}x-\allowbreak{}a,\allowbreak{}0)\allowbreak{}。 prefix\_min：f(\allowbreak{}x)\allowbreak{}=\allowbreak{}min(\allowbreak{}y<\allowbreak{}=\allowbreak{}x)\allowbreak{}f(\allowbreak{}y)\allowbreak{}；suffix\_min 对称。shift(\allowbreak{}a,\allowbreak{}b)\allowbreak{}：min(\allowbreak{}x-\allowbreak{}b<\allowbreak{}=\allowbreak{}y<\allowbreak{}=\allowbreak{}x-\allowbreak{}a)\allowbreak{}f(\allowbreak{}y)\allowbreak{}，a<\allowbreak{}=\allowbreak{}b。 每次加项 O(\allowbreak{}log n)\allowbreak{}，平移 O(\allowbreak{}1)\allowbreak{}，eval O(\allowbreak{}n)\allowbreak{} 仅作查询/\allowbreak{}验证；数值须放入 ll。 例：dp\_i(\allowbreak{}x)\allowbreak{}=\allowbreak{}|x-\allowbreak{}a\_i|+\allowbreak{}min(\allowbreak{}y<\allowbreak{}=\allowbreak{}x)\allowbreak{}dp\_(\allowbreak{}i-\allowbreak{}1)\allowbreak{}(\allowbreak{}y)\allowbreak{}，每轮 prefix\_min(\allowbreak{})\allowbreak{}; add\_abs(\allowbreak{}a\_i)\allowbreak{}。}
\begin{lstlisting}
struct SlopeTrick
{
    using ll=long long;priority_queue<ll> L;priority_queue<ll,vector<ll>,greater<ll>> R;
    ll left_shift=0,right_shift=0,min_value=0;
    void add_left(ll a){if(!R.empty())min_value+=max(0LL,a-(R.top()+right_shift));R.push(a-right_shift);ll x=R.top()+right_shift;R.pop();L.push(x-left_shift);}
    void add_right(ll a){if(!L.empty())min_value+=max(0LL,(L.top()+left_shift)-a);L.push(a-left_shift);ll x=L.top()+left_shift;L.pop();R.push(x-right_shift);}
    void add_abs(ll a){add_left(a);add_right(a);}
    void add_constant(ll a){min_value+=a;}
    void prefix_min(){R={};}
    void suffix_min(){L={};}
    void shift(ll a,ll b){assert(a<=b);left_shift+=a;right_shift+=b;}
    // 此处 eval 只求凸函数数值，不解析或执行任何代码。
    ll eval(ll x)const{ll ans=min_value;auto l=L;auto r=R;while(!l.empty())ans+=max(0LL,l.top()+left_shift-x),l.pop();while(!r.empty())ans+=max(0LL,x-r.top()-right_shift),r.pop();return ans;}
};
\end{lstlisting}
% SOURCE: 06-动态规划/动态DP.cpp
\subsection{动态DP}

\note{adj为树，w为点权；fa/\allowbreak{}dep/\allowbreak{}sz/\allowbreak{}son是父亲、深度、子树大小、重儿子，head/\allowbreak{}pos\_of/\allowbreak{}dfn定位重链。g0/\allowbreak{}g1是轻子树贡献，chnode存各链点序，tr存链上转移矩阵；更新点权后沿链顶向上修正贡献，换树须重建全部链状态。}
\note{矩阵语义：计数用 (\allowbreak{}+\allowbreak{},\allowbreak{}*)\allowbreak{}；可达性用 (\allowbreak{}OR,\allowbreak{}AND)\allowbreak{}；最优值用 (\allowbreak{}max,\allowbreak{}+\allowbreak{})\allowbreak{} 或 (\allowbreak{}min,\allowbreak{}+\allowbreak{})\allowbreak{}。 空段单位矩阵：对角为乘法单位，其他为加法零（布尔 0、max-\allowbreak{}plus -\allowbreak{}INF、min-\allowbreak{}plus +\allowbreak{}INF）。 列向量先过 A 再过 B 为 B*A；不能只把普通矩阵的模数改掉。INF 项不得直接相加。 扫参数时，仿射式相等的条件只可能恒真/\allowbreak{}恒假/\allowbreak{}在唯一 k 成立；在 k 和 k+\allowbreak{}1 都安排更新。 轮廓线仅保留相对标号时，平移已处理前缀 delta 会改变目标 delta*前缀供需差，需结算。}
\begin{lstlisting}
const int N=2005;
const ll NEG=-4e18;
int n,q;
ll w[N];// 点权
vector<int> adj[N];
\end{lstlisting}
\note{=\allowbreak{}=\allowbreak{}=\allowbreak{}=\allowbreak{}=\allowbreak{} 转移矩阵 =\allowbreak{}=\allowbreak{}=\allowbreak{}=\allowbreak{}=\allowbreak{} g0[\allowbreak{}u]\allowbreak{}=\allowbreak{}轻儿子的 max(\allowbreak{}f0[\allowbreak{}v]\allowbreak{},\allowbreak{}f1[\allowbreak{}v]\allowbreak{})\allowbreak{} 之和，g1[\allowbreak{}u]\allowbreak{}=\allowbreak{}w[\allowbreak{}u]\allowbreak{}+\allowbreak{}轻儿子的 f0[\allowbreak{}v]\allowbreak{} 之和 f0[\allowbreak{}u]\allowbreak{}=\allowbreak{}g0[\allowbreak{}u]\allowbreak{}+\allowbreak{}max(\allowbreak{}f0[\allowbreak{}重儿子]\allowbreak{},\allowbreak{}f1[\allowbreak{}重儿子]\allowbreak{})\allowbreak{}，f1[\allowbreak{}u]\allowbreak{}=\allowbreak{}g1[\allowbreak{}u]\allowbreak{}+\allowbreak{}f0[\allowbreak{}重儿子]\allowbreak{} 写成矩阵： (\allowbreak{}f0[\allowbreak{}u]\allowbreak{};f1[\allowbreak{}u]\allowbreak{})\allowbreak{} =\allowbreak{} [\allowbreak{}g0 g0; g1 -\allowbreak{}inf]\allowbreak{} * (\allowbreak{}f0[\allowbreak{}son]\allowbreak{};f1[\allowbreak{}son]\allowbreak{})\allowbreak{} 答案 =\allowbreak{} max(\allowbreak{}f0[\allowbreak{}根]\allowbreak{},\allowbreak{}f1[\allowbreak{}根]\allowbreak{})\allowbreak{}}
\begin{lstlisting}
struct Matrix
{
    ll a[2][2];
    Matrix operator*(const Matrix &o)const
    {
        Matrix c;
        for(int i=0;i<2;i++)for(int j=0;j<2;j++)
        {
            c.a[i][j]=NEG;
            for(int k=0;k<2;k++)if(a[i][k]!=NEG&&o.a[k][j]!=NEG)
                c.a[i][j]=max(c.a[i][j],a[i][k]+o.a[k][j]);
        }
        return c;
    }
};
int fa[N],dep[N],sz[N],son[N],head[N],pos_of[N],dfn[N],tim;
ll g0[N],g1[N];
int nchain;
vector<int> chnode[N];// 每条链上的点，顺序：链头 -> 链尾
vector<int> chid[N];// 点 u 在链内的下标（链头是 0）
vector<Matrix> tr[N];
Matrix mt_of(int u)// 点 u 在重链上的转移矩阵
{
    Matrix m;
    m.a[0][0]=g0[u],m.a[0][1]=g0[u];
    m.a[1][0]=g1[u],m.a[1][1]=NEG;
    return m;
}
\end{lstlisting}
\note{线段树：叶子下标 0 放链尾（最深），下标越大越靠近链头 列向量转移要求浅 * 深，倒序叶子的右段矩阵乘左段矩阵}
\begin{lstlisting}
void build_seg(int c,int p,int l,int r)
{
    int len=chnode[c].size();
    if(l==r){tr[c][p]=mt_of(chnode[c][len-1-l]);return;}
    int mid=(l+r)>>1;
    build_seg(c,p<<1,l,mid);
    build_seg(c,p<<1|1,mid+1,r);
    tr[c][p]=tr[c][p<<1|1]*tr[c][p<<1];
}
void modify(int c,int p,int l,int r,int x,Matrix v)
{
    if(l==r){tr[c][p]=v;return;}
    int mid=(l+r)>>1;
    if(x<=mid)modify(c,p<<1,l,mid,x,v);
    else modify(c,p<<1|1,mid+1,r,x,v);
    tr[c][p]=tr[c][p<<1|1]*tr[c][p<<1];
}
Matrix query(int c,int p,int l,int r,int ql,int qr)
{
    if(ql<=l&&r<=qr)return tr[c][p];
    int mid=(l+r)>>1;
    if(qr<=mid)return query(c,p<<1,l,mid,ql,qr);
    if(ql>mid)return query(c,p<<1|1,mid+1,r,ql,qr);
    return query(c,p<<1|1,mid+1,r,ql,qr)*query(c,p<<1,l,mid,ql,qr);
}
\end{lstlisting}
\note{取点 u 的 f0,\allowbreak{}f1：查它所在链从 u 到链尾的一段 叶子下标 0 是链尾，所以链上 [\allowbreak{}chid[\allowbreak{}u]\allowbreak{},\allowbreak{} len-\allowbreak{}1]\allowbreak{} 对应叶子区间 [\allowbreak{}0,\allowbreak{} len-\allowbreak{}1-\allowbreak{}chid[\allowbreak{}u]\allowbreak{}]\allowbreak{}}
\begin{lstlisting}
void get_f(int u,ll &f0,ll &f1)
{
    int c=head[u],len=chnode[c].size();
    Matrix m=query(c,1,0,len-1,0,len-1-chid[u][0]);
    f0=m.a[0][0],f1=m.a[1][0];
}
\end{lstlisting}
\note{第一次 dfs：重儿子、子树大小}
\begin{lstlisting}
void dfs1(int u,int p)
{
    fa[u]=p,dep[u]=dep[p]+1,sz[u]=1,son[u]=0;
    for(int v:adj[u])
    {
        if(v==p)continue;
        dfs1(v,u);
        sz[u]+=sz[v];
        if(sz[v]>sz[son[u]])son[u]=v;
    }
}
\end{lstlisting}
\note{第二次 dfs：剖分重链，链内按链头-\allowbreak{}>\allowbreak{}链尾顺序放进 chnode}
\begin{lstlisting}
void dfs2(int u,int h)
{
    head[u]=h;
    dfn[u]=++tim;// 记录链头的访问顺序，用来决定建树顺序
    chid[u].push_back(chnode[h].size());
    chnode[h].push_back(u);
    for(int v:adj[u])
        if(v!=fa[u]&&v!=son[u])dfs2(v,v);
    if(son[u])dfs2(son[u],h);
}
\end{lstlisting}
\note{预处理：算 g0,\allowbreak{}g1 并建线段树}
\begin{lstlisting}
void build_all()
{
    dfs1(1,0);
    for(int i=1;i<=n;i++)chnode[i].clear(),chid[i].clear();
    tim=0;
    dfs2(1,1);
    // 按链头被访问的先后顺序倒着处理：保证轻儿子所在链已经算完
    vector<int> corder;
    for(int i=1;i<=n;i++)
        if(chnode[i].size())corder.push_back(i);
    sort(corder.begin(),corder.end(),[&](int x,int y){return dfn[x]>dfn[y];});
    for(int c:corder)
    {
        int L=chnode[c].size();
        for(int x=L-1;x>=0;x--)// 链内从链尾（最深）往链头算
        {
            int u=chnode[c][x];
            g0[u]=0,g1[u]=w[u];
            for(int v:adj[u])
                if(v!=fa[u]&&v!=son[u])
                {
                    ll v0,v1;
                    get_f(v,v0,v1);
                    g0[u]+=max(v0,v1),g1[u]+=v0;
                }
        }
        tr[c].assign(L*4+4,Matrix());
        build_seg(c,1,0,L-1);
    }
    nchain=corder.size();
    for(int u=1;u<=n;u++)pos_of[u]=chnode[head[u]].size()-1-chid[u][0];
}
\end{lstlisting}
\note{点权修改：沿着重链往上跳，每跳一条链就把这条链的链头 dp 结果并进父结点的轻儿子贡献}
\begin{lstlisting}
void update(int u,ll x)
{
    g1[u]+=x-w[u],w[u]=x;
    while(u)
    {
        int c=head[u],len=chnode[c].size();
        ll old0,old1,new0,new1;
        get_f(c,old0,old1);
        modify(c,1,0,len-1,pos_of[u],mt_of(u));
        get_f(c,new0,new1);
        u=fa[c];
        if(u)
        {
            g0[u]+=max(new0,new1)-max(old0,old1);
            g1[u]+=new0-old0;
        }
    }
}
\end{lstlisting}
% SOURCE: 06-动态规划/插头DP(轮廓线).cpp
\subsection{插头DP(轮廓线)}

\note{vis[\allowbreak{}i]\allowbreak{}[\allowbreak{}j]\allowbreak{}表示可走格，f[\allowbreak{}0/\allowbreak{}1]\allowbreak{}以压缩轮廓状态为键、方案数为值，n/\allowbreak{}m为行列。位段编码连通性而非仅“是否占用”；每格转移前清空下一层，换行移位并保持闭环只在允许的最后位置出现。}
\begin{lstlisting}
const int N=14;// 轮廓线长度（列数上限）
int n,m;// n 行 m 列，插头 dp 适合 n,m<=12
int vis[N][N];
\end{lstlisting}
\note{状态编码：2 bit 一个插头，从低位到高位对应轮廓线上第 0,\allowbreak{}1,\allowbreak{}.\allowbreak{}.\allowbreak{}.\allowbreak{},\allowbreak{}m 个位置 0 =\allowbreak{} 无插头，1 =\allowbreak{} 左括号 '(\allowbreak{}',\allowbreak{} 2 =\allowbreak{} 右括号 ')\allowbreak{}' 轮廓线上"还没接好的插头"一定是匹配好的括号序列，所以只需要记括号类型 第 j 个格子的左插头在位置 j，上插头在位置 j+\allowbreak{}1（j 从 0 起）}
\begin{lstlisting}
unordered_map<int,ll> f[2];// 滚动数组，按格转移
int getp(int st,int p)// 取第 p 个插头的类型
{
    return (st>>(p*2))&3;
}
int setp(int st,int p,int v)// 把第 p 个插头改成 v
{
    return (st&~(3<<(p*2)))|(v<<(p*2));
}
\end{lstlisting}
\note{找第 p 个插头的配对括号位置，找不到返回 -\allowbreak{}1}
\begin{lstlisting}
int match_p(int st,int p,int m)
{
    int v=getp(st,p);
    if(v==0)return -1;
    int cnt=0;
    if(v==1)// '(' 向右找配对的 ')'
    {
        for(int i=p;i<=m;i++)
        {
            int c=getp(st,i);
            if(c==1)cnt++;
            else if(c==2)cnt--;
            if(cnt==0)return i;
        }
    }
    else// ')' 向左找配对的 '('
    {
        for(int i=p;i>=0;i--)
        {
            int c=getp(st,i);
            if(c==2)cnt++;
            else if(c==1)cnt--;
            if(cnt==0)return i;
        }
    }
    return -1;
}
int has_plug(int st,int m)// 轮廓线上还有没有未闭合的插头（只用到 0..m-1 位）
{
    for(int i=0;i<m;i++)
        if(getp(st,i))return 1;
    return 0;
}
\end{lstlisting}
\note{求 n*m 网格的哈密顿回路条数（每个格子恰好走一次，最后回到起点） 逐格转移，四种情况：0 个插头 /\allowbreak{} 1 个插头(\allowbreak{}分左右括号)\allowbreak{} /\allowbreak{} 2 个插头(\allowbreak{}分三种括号组合)\allowbreak{} 第 j 位是左插头，第 j+\allowbreak{}1 位是上插头；输出下插头在 j，右插头在 j+\allowbreak{}1 复杂度 O(\allowbreak{}n*m*状态数)\allowbreak{}，本题模板适合 n,\allowbreak{}m<\allowbreak{}=\allowbreak{}12}
\begin{lstlisting}
ll plug_dp(int n,int m)
{
    if((n*m)%2==1)return 0;// 奇数个格子不可能有哈密顿回路
    f[0].clear(),f[1].clear();
    f[0][0]=1;
    int cur=0;
    for(int i=0;i<n;i++)
    {
        if(i>0)// 换行后下插头变为下一行的上插头
        {
            f[cur^1].clear();
            for(auto &pr:f[cur])
                if(!getp(pr.first,m))f[cur^1][pr.first<<2]+=pr.second;
            f[cur].clear();
            cur^=1;
        }
        for(int j=0;j<m;j++)
        {
            f[cur^1].clear();
            for(auto &pr:f[cur])
            {
                int st=pr.first;
                ll cnt=pr.second;
                int l=getp(st,j),u=getp(st,j+1);
                int base=setp(setp(st,j,0),j+1,0);
                if(!l&&!u)
                {
                    if(i+1<n&&j+1<m)
                        f[cur^1][setp(setp(base,j,1),j+1,2)]+=cnt;
                }
                else if(!l||!u)
                {
                    int v=l?l:u;
                    if(i+1<n)f[cur^1][setp(base,j,v)]+=cnt;
                    if(j+1<m)f[cur^1][setp(base,j+1,v)]+=cnt;
                }
                else if(l==1&&u==2)
                {
                    if(i==n-1&&j==m-1&&base==0)f[cur^1][0]+=cnt;
                }
                else if(l==2&&u==1)f[cur^1][base]+=cnt;
                else if(l==1&&u==1)
                {
                    int p=match_p(st,j+1,m);
                    if(p!=-1)f[cur^1][setp(base,p,1)]+=cnt;
                }
                else if(l==2&&u==2)
                {
                    int p=match_p(st,j,m);
                    if(p!=-1)f[cur^1][setp(base,p,2)]+=cnt;
                }
            }
            cur^=1;
        }
    }
    return f[cur][0];
}
\end{lstlisting}
\note{铺砖轮廓线：第 j 位表示当前列已被之前的骨牌占用}
\begin{lstlisting}
ll domino_tiling(int n,int m)
{
    vector<ll> g(1<<m,0),h;
    g[0]=1;
    for(int i=0;i<n;i++)
        for(int j=0;j<m;j++)
        {
            h.assign(1<<m,0);
            for(int st=0;st<(1<<m);st++)
            {
                if(!g[st])continue;
                if(st>>j&1)h[st^(1<<j)]+=g[st];
                else
                {
                    if(j+1<m&&!(st>>(j+1)&1))h[st|(1<<(j+1))]+=g[st];
                    if(i+1<n)h[st|(1<<j)]+=g[st];
                }
            }
            g.swap(h);
        }
    return g[0];
}
\end{lstlisting}
\section{搜索}
% SOURCE: 07-搜索/BFS最短路.cpp
\subsection{BFS最短路}

\note{n/\allowbreak{}m为网格大小，block为障碍，dis=\allowbreak{}-\allowbreak{}1表示未访问；dx/\allowbreak{}dy为四邻/\allowbreak{}马步方向。md、g、tim、dt用于多源、时序等独立例子；先传播危险到达时间，再检查人能否更早到达，重建队列与距离后开始搜索。}
\begin{lstlisting}
const int N=405;
\end{lstlisting}
\note{四方向偏移表（下标 0 空着，从 1 开始用）}
\begin{lstlisting}
int dx4[]={0,1,-1,0,0};
int dy4[]={0,0,0,1,-1};
\end{lstlisting}
\note{马走日，八方向}
\begin{lstlisting}
int dx8[]={0,-2,-2,-1,-1,1,1,2,2};
int dy8[]={0,-1,1,-2,2,-2,2,-1,1};
\end{lstlisting}
\note{-\allowbreak{}-\allowbreak{}-\allowbreak{}-\allowbreak{}-\allowbreak{}-\allowbreak{}-\allowbreak{}-\allowbreak{}-\allowbreak{}-\allowbreak{} 一、网格 BFS 求最少步数 -\allowbreak{}-\allowbreak{}-\allowbreak{}-\allowbreak{}-\allowbreak{}-\allowbreak{}-\allowbreak{}-\allowbreak{}-\allowbreak{}-\allowbreak{}}
\begin{lstlisting}
int n,m;
int dis[N][N];//-1 表示不可达
bool block[N][N];//障碍格
\end{lstlisting}
\note{单源 BFS，O(\allowbreak{}nm)\allowbreak{}}
\begin{lstlisting}
void bfs_grid(int sx,int sy)
{
    memset(dis,-1,sizeof(dis));
    queue<pair<int,int> > q;
    if(block[sx][sy])return;
    dis[sx][sy]=0;
    q.push(make_pair(sx,sy));
    while(!q.empty())
    {
        int x=q.front().first,y=q.front().second;
        q.pop();
        for(int i=1;i<=4;i++)
        {
            int xx=x+dx4[i],yy=y+dy4[i];
            if(xx<1||xx>n||yy<1||yy>m)continue;//边界
            if(block[xx][yy])continue;//障碍
            if(dis[xx][yy]!=-1)continue;//已访问
            dis[xx][yy]=dis[x][y]+1;
            q.push(make_pair(xx,yy));
        }
    }
}
\end{lstlisting}
\note{-\allowbreak{}-\allowbreak{}-\allowbreak{}-\allowbreak{}-\allowbreak{}-\allowbreak{}-\allowbreak{}-\allowbreak{}-\allowbreak{}-\allowbreak{} 二、多源 BFS -\allowbreak{}-\allowbreak{}-\allowbreak{}-\allowbreak{}-\allowbreak{}-\allowbreak{}-\allowbreak{}-\allowbreak{}-\allowbreak{}-\allowbreak{}}
\begin{lstlisting}
int md[N][N];
\end{lstlisting}
\note{多个起点同时入队，求每个点到最近起点的距离，O(\allowbreak{}nm)\allowbreak{}}
\begin{lstlisting}
void bfs_multi(vector<pair<int,int> > src)
{
    memset(md,-1,sizeof(md));
    queue<pair<int,int> > q;
    for(int i=0;i<(int)src.size();i++)
    {
        int x=src[i].first,y=src[i].second;
        md[x][y]=0;
        q.push(src[i]);
    }
    while(!q.empty())
    {
        int x=q.front().first,y=q.front().second;
        q.pop();
        for(int i=1;i<=8;i++)//马步
        {
            int xx=x+dx8[i],yy=y+dy8[i];
            if(xx<1||xx>n||yy<1||yy>m)continue;
            if(md[xx][yy]!=-1)continue;
            md[xx][yy]=md[x][y]+1;
            q.push(make_pair(xx,yy));
        }
    }
}
\end{lstlisting}
\note{-\allowbreak{}-\allowbreak{}-\allowbreak{}-\allowbreak{}-\allowbreak{}-\allowbreak{}-\allowbreak{}-\allowbreak{}-\allowbreak{}-\allowbreak{} 三、泛洪 /\allowbreak{} 连通块 -\allowbreak{}-\allowbreak{}-\allowbreak{}-\allowbreak{}-\allowbreak{}-\allowbreak{}-\allowbreak{}-\allowbreak{}-\allowbreak{}-\allowbreak{}}
\begin{lstlisting}
int g[N][N];
\end{lstlisting}
\note{从 (\allowbreak{}x,\allowbreak{}y)\allowbreak{} 出发把同一连通块染色；调用外面留一圈 0 保证边界的连通性}
\begin{lstlisting}
void flood(int x,int y,int c)
{
    assert(c!=0);
    if(g[x][y]!=0)return;
    queue<pair<int,int> > q;
    g[x][y]=c;
    q.push(make_pair(x,y));
    while(!q.empty())
    {
        int xx=q.front().first,yy=q.front().second;
        q.pop();
        for(int i=1;i<=4;i++)
        {
            int a=xx+dx4[i],b=yy+dy4[i];
            if(a<0||a>n+1||b<0||b>m+1)continue;//外面留了一圈
            if(g[a][b])continue;//非 0 都是非空
            g[a][b]=c;
            q.push(make_pair(a,b));
        }
    }
}
\end{lstlisting}
\note{-\allowbreak{}-\allowbreak{}-\allowbreak{}-\allowbreak{}-\allowbreak{}-\allowbreak{}-\allowbreak{}-\allowbreak{}-\allowbreak{}-\allowbreak{} 四、带时间限制的 BFS（流星雨）-\allowbreak{}-\allowbreak{}-\allowbreak{}-\allowbreak{}-\allowbreak{}-\allowbreak{}-\allowbreak{}-\allowbreak{}-\allowbreak{}-\allowbreak{}}
\begin{lstlisting}
const int INF=0x3f3f3f3f;
int tim[N][N],dt[N][N];
\end{lstlisting}
\note{tim[\allowbreak{}i]\allowbreak{}[\allowbreak{}j]\allowbreak{} =\allowbreak{} 该格被砸毁的时刻（INF 表示永不），求最早能到达「永不被砸」的格子的时刻 站在格子上要求到达时刻 t 满足 t <\allowbreak{} tim；注意被砸的格子在时刻 tim 就不能站了 起点 (\allowbreak{}0,\allowbreak{}0)\allowbreak{} 若在时刻 0 就被砸则无解}
\begin{lstlisting}
int bfs_meteor()
{
    memset(dt,-1,sizeof(dt));
    if(tim[0][0]<=0)return -1;
    queue<pair<int,int> > q;
    dt[0][0]=0;
    q.push(make_pair(0,0));
    while(!q.empty())
    {
        int x=q.front().first,y=q.front().second;
        q.pop();
        if(tim[x][y]==INF)return dt[x][y];//永久安全的格子，BFS 首次到即最早
        for(int i=1;i<=4;i++)
        {
            int xx=x+dx4[i],yy=y+dy4[i];
            if(xx<0||xx>N-1||yy<0||yy>N-1)continue;
            if(dt[xx][yy]!=-1)continue;
            if(dt[x][y]+1>=tim[xx][yy])continue;//到达时已被砸
            dt[xx][yy]=dt[x][y]+1;
            q.push(make_pair(xx,yy));
        }
    }
    return -1;
}
\end{lstlisting}
\note{-\allowbreak{}-\allowbreak{}-\allowbreak{}-\allowbreak{}-\allowbreak{}-\allowbreak{}-\allowbreak{}-\allowbreak{}-\allowbreak{}-\allowbreak{} 五、八数码状态哈希 -\allowbreak{}-\allowbreak{}-\allowbreak{}-\allowbreak{}-\allowbreak{}-\allowbreak{}-\allowbreak{}-\allowbreak{}-\allowbreak{}-\allowbreak{}}
\note{把 3*3 压成一个 36 位整数：第 i 格（0.\allowbreak{}.\allowbreak{}8）放在右起第 (\allowbreak{}8-\allowbreak{}i)\allowbreak{} 个 4 bit， 也就是「按行读成十六进制」的样子，目标态 123456780 编码后就是 0x123456780}
\begin{lstlisting}
ll encode(int b[3][3])
{
    ll s=0;
    for(int i=0;i<3;i++)
        for(int j=0;j<3;j++)s|=(ll)b[i][j]<<(4*(8-(i*3+j)));
    return s;
}
void decode(ll s,int b[3][3])
{
    for(int i=0;i<3;i++)
        for(int j=0;j<3;j++)b[i][j]=(s>>(4*(8-(i*3+j))))&15;
}
\end{lstlisting}
\note{八数码里空格向上下左右移动：pdx 是行号变化，pdy 是列号变化 行号增大是往下，所以「下」是 +\allowbreak{}1；这里写成 \{0,\allowbreak{}-\allowbreak{}1,\allowbreak{}1,\allowbreak{}0,\allowbreak{}0\} 表示 上、下、左、右}
\begin{lstlisting}
const int pdx[]={0,-1,1,0,0};
const int pdy[]={0,0,0,-1,1};
\end{lstlisting}
\note{交换第 p 格与第 q 格（p,\allowbreak{}q 是 0.\allowbreak{}.\allowbreak{}8 的格子编号），O(\allowbreak{}1)\allowbreak{} 用 decode/\allowbreak{}encode 保证索引口径和上面完全一致，不容易写错}
\begin{lstlisting}
ll flip(ll s,int p,int q)
{
    int b[3][3];
    decode(s,b);
    swap(b[p/3][p%3],b[q/3][q%3]);
    return encode(b);
}
\end{lstlisting}
\note{八数码最少步数（标准 123456780 为目标），无解返回 -\allowbreak{}1，O(\allowbreak{}状态数*4)\allowbreak{}}
\begin{lstlisting}
int bfs_puzzle(int st[3][3])
{
    int goal[3][3]={{1,2,3},{4,5,6},{7,8,0}};
    const ll T=encode(goal);//目标态编码 = 0x123456780
    ll S=encode(st);
    if(S==T)return 0;
    map<ll,int> dis;
    queue<pair<ll,int> > q;//(状态, 空格位置)
    int z=0;
    for(int i=0;i<3;i++)
        for(int j=0;j<3;j++)
            if(!st[i][j])z=i*3+j;
    dis[S]=0;
    q.push(make_pair(S,z));
    while(!q.empty())
    {
        ll u=q.front().first;
        int p=q.front().second;
        q.pop();
        int zx=p/3,zy=p%3;
        for(int i=1;i<=4;i++)
        {
            int xx=zx+pdx[i],yy=zy+pdy[i];
            if(xx<0||xx>2||yy<0||yy>2)continue;
            ll v=flip(u,zx*3+zy,xx*3+yy);
            if(dis.count(v))continue;
            dis[v]=dis[u]+1;
            if(v==T)return dis[v];
            q.push(make_pair(v,xx*3+yy));
        }
    }
    return -1;
}
\end{lstlisting}
% SOURCE: 07-搜索/DFS与剪枝.cpp
\subsection{DFS与剪枝}

\note{a为候选值，per为当前排列/\allowbreak{}选择，vis为使用标记；col/\allowbreak{}dg/\allowbreak{}udg记录皇后的列与对角线占用。best是当前上界，sum\_rest是未处理非负元素的后缀和；递归返回撤销本层修改，不把某例计数与另一例混用。}
\begin{lstlisting}
const int N=15;
int n,k;
int a[N];
\end{lstlisting}
\note{-\allowbreak{}-\allowbreak{}-\allowbreak{}-\allowbreak{}-\allowbreak{}-\allowbreak{}-\allowbreak{}-\allowbreak{}-\allowbreak{}-\allowbreak{} 一、枚举类 DFS -\allowbreak{}-\allowbreak{}-\allowbreak{}-\allowbreak{}-\allowbreak{}-\allowbreak{}-\allowbreak{}-\allowbreak{}-\allowbreak{}-\allowbreak{}}
\note{全排列，O(\allowbreak{}n!)\allowbreak{}}
\begin{lstlisting}
int per[N];
bool vis[N];
void dfs_perm(int dep)//dep 从 1 开始，填第 dep 个位置
{
    if(dep>n)
    {
        for(int i=1;i<=n;i++)printf("%d ",per[i]);
        printf("\n");
        return;
    }
    for(int i=1;i<=n;i++)
        if(!vis[i])
        {
            vis[i]=true,per[dep]=i;
            dfs_perm(dep+1);
            vis[i]=false;
        }
}
int comb_cnt;
\end{lstlisting}
\note{组合，从 [\allowbreak{}idx,\allowbreak{}n]\allowbreak{} 里选，已选 dep 个；O(\allowbreak{}C(\allowbreak{}n,\allowbreak{}k)\allowbreak{})\allowbreak{}}
\begin{lstlisting}
void dfs_comb(int dep,int idx)//必须递增选，天然去重
{
    if(dep==k)
    {
        comb_cnt++;
        return;
    }
    for(int i=idx;i<=n;i++)
        dfs_comb(dep+1,i+1);
}
int sub_cnt;
\end{lstlisting}
\note{子集枚举（DFS 版），O(\allowbreak{}2\textasciicircum{}n)\allowbreak{}}
\begin{lstlisting}
void dfs_subset(int idx)
{
    if(idx>n)
    {
        sub_cnt++;
        return;
    }
    dfs_subset(idx+1);//不选 a[idx]，优先不选可以保证子集不重不漏
    dfs_subset(idx+1);//选 a[idx]
}
\end{lstlisting}
\note{位掩码枚举子集，O(\allowbreak{}2\textasciicircum{}n)\allowbreak{}，注意 1<\allowbreak{}<\allowbreak{}n 在 n>\allowbreak{}=\allowbreak{}31 时溢出}
\begin{lstlisting}
void mask_subset()
{
    for(int mask=0;mask<(1<<n);mask++)
        for(int i=0;i<n;i++)
            if(mask>>i&1){}//i 在子集里
}
\end{lstlisting}
\note{-\allowbreak{}-\allowbreak{}-\allowbreak{}-\allowbreak{}-\allowbreak{}-\allowbreak{}-\allowbreak{}-\allowbreak{}-\allowbreak{}-\allowbreak{} 二、可行性剪枝 -\allowbreak{}-\allowbreak{}-\allowbreak{}-\allowbreak{}-\allowbreak{}-\allowbreak{}-\allowbreak{}-\allowbreak{}-\allowbreak{}-\allowbreak{}}
\note{组合数之和为素数个数；传入 sum 而不是在叶子处再求和，边界剪枝}
\begin{lstlisting}
int prime_cnt;
bool is_prime(int x)//试除法，O(sqrt(x))
{
    if(x<2)return false;
    for(int i=2;i<=x/i;i++)
        if(x%i==0)return false;
    return true;
}
void dfs_prime(int step,int sum,int idx)//step 已选个数，sum 已选之和
{
    if(step==k)
    {
        if(is_prime(sum))prime_cnt++;
        return;
    }
    if(n-idx<k-step)return;//可行性剪枝：剩下的数不够凑满 k 个
    for(int i=idx;i<n;i++)
        dfs_prime(step+1,sum+a[i],i+1);
}
\end{lstlisting}
\note{N 皇后计数，逐行放；列 /\allowbreak{} 主对角 r-\allowbreak{}c+\allowbreak{}n /\allowbreak{} 副对角 r+\allowbreak{}c 三个冲突数组}
\begin{lstlisting}
int nn,cnt_queen;
int col[N],dg[N<<1],udg[N<<1];
void dfs_queen(int r)
{
    if(r>nn)
    {
        cnt_queen++;
        return;
    }
    for(int c=1;c<=nn;c++)
    {
        if(col[c]||dg[r-c+nn]||udg[r+c])continue;
        col[c]=dg[r-c+nn]=udg[r+c]=1;
        dfs_queen(r+1);
        col[c]=dg[r-c+nn]=udg[r+c]=0;
    }
}
\end{lstlisting}
\note{-\allowbreak{}-\allowbreak{}-\allowbreak{}-\allowbreak{}-\allowbreak{}-\allowbreak{}-\allowbreak{}-\allowbreak{}-\allowbreak{}-\allowbreak{} 三、最优性剪枝（Branch and Bound）-\allowbreak{}-\allowbreak{}-\allowbreak{}-\allowbreak{}-\allowbreak{}-\allowbreak{}-\allowbreak{}-\allowbreak{}-\allowbreak{}-\allowbreak{}}
\begin{lstlisting}
int best;
int sum_rest[N+1];//sum_rest[i] = a[i]+...+a[n]（非负元素之和），必须是「还没放下去的元素」之和
\end{lstlisting}
\note{把 n 个数分成两组，最小化两组和的差；O(\allowbreak{}2\textasciicircum{}n)\allowbreak{}，剪枝后远小于 最优性剪枝：还没放的数最多把差值拉小 sum\_rest[\allowbreak{}idx]\allowbreak{}， 所以 |s1-\allowbreak{}s2|-\allowbreak{}sum\_rest[\allowbreak{}idx]\allowbreak{} >\allowbreak{} best 时这一支不可能更优，直接砍 注意 idx 这个位置的 a[\allowbreak{}idx]\allowbreak{} 还没选，sum\_rest[\allowbreak{}idx]\allowbreak{} 必须包含它}
\begin{lstlisting}
void dfs_split(int idx,int s1,int s2)
{
    if(idx>n)
    {
        best=min(best,abs(s1-s2));
        return;
    }
    if(abs(s1-s2)-sum_rest[idx]>best)return;//乐观下界剪枝
    dfs_split(idx+1,s1+a[idx],s2);
    dfs_split(idx+1,s1,s2+a[idx]);
}
\end{lstlisting}
% SOURCE: 07-搜索/折半搜索.cpp
\subsection{折半搜索}

\note{a是0起数值序列，sums[\allowbreak{}0/\allowbreak{}1]\allowbreak{}存两半的子集和；入口只传目标值等标量。求解时重新枚举并排序；负数允许，但剪枝不能假定剩余元素非负。}
\begin{lstlisting}
vector<ll> a,sums[2];
\end{lstlisting}
\note{O(\allowbreak{}2\textasciicircum{}(\allowbreak{}n/\allowbreak{}2)\allowbreak{} n)\allowbreak{}，含负数；统计和等于目标的子集个数，n<\allowbreak{}=\allowbreak{}40}
\begin{lstlisting}
void enum_sum(int l,int r,int id)
{
    auto &s=sums[id];s.assign(1,0);
    for(int i=l;i<r;i++)
    {
        int m=s.size();
        for(int j=0;j<m;j++)s.push_back(s[j]+a[i]);
    }
}
ll solve(ll target)
{
    int n=a.size();
    enum_sum(0,n/2,0);enum_sum(n/2,n,1);
    auto &l=sums[0],&r=sums[1];
    sort(r.begin(),r.end());
    ll ans=0;
    for(ll x:l)
    {
        __int128 want=(__int128)target-x;
        if(want<LLONG_MIN||want>LLONG_MAX)continue;
        auto p=equal_range(r.begin(),r.end(),(ll)want);
        ans+=p.second-p.first;
    }
    return ans;
}
\end{lstlisting}
% SOURCE: 07-搜索/双向BFS.cpp
\subsection{双向BFS}

\note{mp为网格，'\#'表示墙，df/\allowbreak{}db分别是两端到格子的距离；disf/\allowbreak{}disb为编码状态的两侧距离表，dig为编码位数。每次清空两侧访问状态，按完整层扩展再判断会合；反向搜索必须使用逆转移。}
\begin{lstlisting}
typedef unsigned long long ull;
const int N=1005;
\end{lstlisting}
\note{双向广搜：起点和终点同时 BFS，交替扩展「待扩展结点少」的那一侧。 第一次汇合时两边的距离之和就是最短路，扩展量约为单向的 2*b\textasciicircum{}(\allowbreak{}d/\allowbreak{}2)\allowbreak{}。 要求边权全为 1（或者等权），否则汇合时的答案不一定最优。}
\begin{lstlisting}
int dx4[]={0,1,-1,0,0};
int dy4[]={0,0,0,1,-1};
\end{lstlisting}
\note{-\allowbreak{}-\allowbreak{}-\allowbreak{}-\allowbreak{}-\allowbreak{}-\allowbreak{}-\allowbreak{}-\allowbreak{}-\allowbreak{}-\allowbreak{} 一、网格迷宫上的双向 BFS -\allowbreak{}-\allowbreak{}-\allowbreak{}-\allowbreak{}-\allowbreak{}-\allowbreak{}-\allowbreak{}-\allowbreak{}-\allowbreak{}-\allowbreak{}}
\begin{lstlisting}
int n,m;
char mp[N][N];//'#' 是墙，'.' 是空地
int df[N][N],db[N][N];//正向/反向到该格的距离
\end{lstlisting}
\note{返回最少步数，不可达返回 -\allowbreak{}1，O(\allowbreak{}nm)\allowbreak{} df 是正向距离，db 是反向距离；注意汇合时要扫一整层邻居（不能只看新点的距离）， 并且要等「两侧已扩展层数之和 >\allowbreak{}=\allowbreak{} 目前最好答案」才能停，否则会退出太早。}
\begin{lstlisting}
int bibfs_grid(int sx,int sy,int ex,int ey)
{
    if(mp[sx][sy]=='#'||mp[ex][ey]=='#')return -1;
    if(sx==ex&&sy==ey)return 0;
    memset(df,-1,sizeof(df));
    memset(db,-1,sizeof(db));
    queue<pair<int,int> > qf,qb;
    df[sx][sy]=0,qf.push(make_pair(sx,sy));
    db[ex][ey]=0,qb.push(make_pair(ex,ey));
    int depf=0,depb=0;//正向/反向「已经完整扩展过」的层数
    int best=0x3f3f3f3f;
    while(!qf.empty()&&!qb.empty())
    {
        if(depf+depb>=best)break;//再走也不会更短
        // 每次都挑小的一侧扩展，这是双向 BFS 的关键
        if(qf.size()<=qb.size())
        {
            int x=qf.front().first,y=qf.front().second;
            qf.pop();
            depf=max(depf,df[x][y]);
            for(int i=1;i<=4;i++)
            {
                int xx=x+dx4[i],yy=y+dy4[i];
                if(xx<1||xx>n||yy<1||yy>m)continue;
                if(mp[xx][yy]=='#')continue;
                if(df[xx][yy]==-1)
                {
                    df[xx][yy]=df[x][y]+1;
                    qf.push(make_pair(xx,yy));
                }
                if(db[xx][yy]!=-1)best=min(best,df[xx][yy]+db[xx][yy]);//两侧在 (xx,yy) 汇合
            }
        }
        else
        {
            int x=qb.front().first,y=qb.front().second;
            qb.pop();
            depb=max(depb,db[x][y]);
            for(int i=1;i<=4;i++)
            {
                int xx=x+dx4[i],yy=y+dy4[i];
                if(xx<1||xx>n||yy<1||yy>m)continue;
                if(mp[xx][yy]=='#')continue;
                if(db[xx][yy]==-1)
                {
                    db[xx][yy]=db[x][y]+1;
                    qb.push(make_pair(xx,yy));
                }
                if(df[xx][yy]!=-1)best=min(best,df[xx][yy]+db[xx][yy]);
            }
        }
    }
    return best==0x3f3f3f3f?-1:best;
}
\end{lstlisting}
\note{-\allowbreak{}-\allowbreak{}-\allowbreak{}-\allowbreak{}-\allowbreak{}-\allowbreak{}-\allowbreak{}-\allowbreak{}-\allowbreak{}-\allowbreak{} 二、数字变换上的双向 BFS -\allowbreak{}-\allowbreak{}-\allowbreak{}-\allowbreak{}-\allowbreak{}-\allowbreak{}-\allowbreak{}-\allowbreak{}-\allowbreak{}-\allowbreak{}}
\note{操作：把某一位数字 +\allowbreak{}1，逢 9 变 0（模 10）。 例：1234 -\allowbreak{}>\allowbreak{} 1235、1234 -\allowbreak{}>\allowbreak{} 1244、1999 -\allowbreak{}>\allowbreak{} 1990（每次只改一位） 位数少（<\allowbreak{}=\allowbreak{}6）时状态最多 10\textasciicircum{}6，用 ull 直接编码（每 4 bit 一位）， 每位 4 bit，反向搜索用 -\allowbreak{}1，枚举有向操作的前驱。}
\begin{lstlisting}
int dig;
unordered_map<ull,int> disf,disb;
ull step_digit(ull s,int k,int delta=1)//第 k 位（从低位起 0 开始）+1，O(1)
{
    int v=(s>>(4*k))&15;
    v=(v+delta+10)%10;
    return (s&~(15ULL<<(4*k)))|((ull)v<<(4*k));
}
\end{lstlisting}
\note{返回最少步数，不可达返回 -\allowbreak{}1，O(\allowbreak{}状态数*dig)\allowbreak{}}
\begin{lstlisting}
int bibfs_num(ull s,ull t)
{
    if(s==t)return 0;
    disf.clear(),disb.clear();
    queue<ull> qf,qb;
    disf[s]=0,qf.push(s);
    disb[t]=0,qb.push(t);
    int depf=0,depb=0,best=0x3f3f3f3f;
    while(!qf.empty()&&!qb.empty())
    {
        if(depf+depb>=best)break;
        if(qf.size()<=qb.size())
        {
            ull u=qf.front();qf.pop();
            depf=max(depf,disf[u]);
            for(int k=0;k<dig;k++)
            {
                ull v=step_digit(u,k);
                if(!disf.count(v))
                {
                    disf[v]=disf[u]+1;
                    qf.push(v);
                }
                if(disb.count(v))best=min(best,disf[v]+disb[v]);
            }
        }
        else
        {
            ull u=qb.front();qb.pop();
            depb=max(depb,disb[u]);
            for(int k=0;k<dig;k++)
            {
                ull v=step_digit(u,k,-1);
                if(!disb.count(v))
                {
                    disb[v]=disb[u]+1;
                    qb.push(v);
                }
                if(disf.count(v))best=min(best,disf[v]+disb[v]);
            }
        }
    }
    return best==0x3f3f3f3f?-1:best;
}
\end{lstlisting}
% SOURCE: 07-搜索/迭代加深与IDA星.cpp
\subsection{迭代加深与IDA星}

\note{id\_a/\allowbreak{}id\_vis供迭代加深例，ks与ks\_goal为当前/\allowbreak{}目标棋盘，ks\_sx/\allowbreak{}ks\_sy为空格位置，ks\_lim为深度界；cr\_s/\allowbreak{}cr\_t是块移动例的当前/\allowbreak{}目标序列。各例初始化答案和深度界；交换或移动后必须回溯恢复状态。}
\note{=\allowbreak{}=\allowbreak{}=\allowbreak{}=\allowbreak{}=\allowbreak{}=\allowbreak{}=\allowbreak{}=\allowbreak{}=\allowbreak{}=\allowbreak{} 一、迭代加深 IDDFS =\allowbreak{}=\allowbreak{}=\allowbreak{}=\allowbreak{}=\allowbreak{}=\allowbreak{}=\allowbreak{}=\allowbreak{}=\allowbreak{}=\allowbreak{} 适合「答案深度浅，但每层分支巨大」的搜索：逐层放宽深度上限， 空间只有 O(\allowbreak{}深度)\allowbreak{}，且第一次找到答案时深度即为最优（每层代价相同的前提）。 单次 DFS 为 O(\allowbreak{}b\textasciicircum{}lim)\allowbreak{}，总复杂度 O(\allowbreak{}b\textasciicircum{}lim)\allowbreak{}，比 BFS 省状态存储。}
\begin{lstlisting}
int id_n;
int id_a[25];
bool id_vis[25];
int id_ans;
\end{lstlisting}
\note{用互不相同的 a[\allowbreak{}i]\allowbreak{} 做加减，凑出 target，问最少用几个数}
\begin{lstlisting}
bool id_dfs(int dep,int lim,ll cur,ll target)
{
    if(cur==target)return true;
    if(dep==lim)return false;//到达深度上限，直接回溯
    for(int i=1;i<=id_n;i++)
    {
        if(id_vis[i])continue;
        id_vis[i]=true;
        bool ok=id_dfs(dep+1,lim,cur+id_a[i],target)||id_dfs(dep+1,lim,cur-id_a[i],target);
        id_vis[i]=false;
        if(ok)return true;
        id_vis[i]=false;
    }
    return false;
}
\end{lstlisting}
\note{返回最少用几个数；无解返回 -\allowbreak{}1}
\begin{lstlisting}
int solve_iddfs(ll target)
{
    fill(id_vis,id_vis+25,false);
    for(int lim=0;lim<=id_n;lim++)//上限从小到大
        if(id_dfs(0,lim,0,target))return lim;
    return -1;
}
\end{lstlisting}
\note{=\allowbreak{}=\allowbreak{}=\allowbreak{}=\allowbreak{}=\allowbreak{}=\allowbreak{}=\allowbreak{}=\allowbreak{}=\allowbreak{}=\allowbreak{} 二、IDA*：knight swap（洛谷 P2324 同型）=\allowbreak{}=\allowbreak{}=\allowbreak{}=\allowbreak{}=\allowbreak{}=\allowbreak{}=\allowbreak{}=\allowbreak{}=\allowbreak{}=\allowbreak{} 估价函数 h =\allowbreak{} 不在目标位置的棋子数，每次移动最多修正 1 个，可采纳。 f =\allowbreak{} g +\allowbreak{} h >\allowbreak{} maxdep 就剪枝。}
\begin{lstlisting}
const int KB=5;
int ks[KB][KB],ks_goal[KB][KB];
int ks_sx,ks_sy,cnt_ks,ks_lim,ks_sol;
int kdx[]={0,-2,-2,-1,-1,1,1,2,2};
int kdy[]={0,-1,1,-2,2,-2,2,-1,1};
int ks_h()//O(25)
{
    int c=0;
    for(int i=0;i<KB;i++)
        for(int j=0;j<KB;j++)
            if(ks[i][j]!=-1&&ks[i][j]!=ks_goal[i][j])c++;
    return c;
}
long long ks_nodes;//统计访问结点数，用来看搜索规模
void ks_dfs(int x,int y,int g,int pre)
{
    ks_nodes++;
    int h=ks_h();
    if(h==0){ks_sol=g;return;}
    if(g+h>ks_lim)return;//IDA* 核心剪枝
    for(int i=1;i<=8;i++)
    {
        int xx=x+kdx[i],yy=y+kdy[i];
        if(xx<0||xx>=KB||yy<0||yy>=KB)continue;
        if(xx*KB+yy==pre)continue;//不立刻走回上一步
        swap(ks[x][y],ks[xx][yy]);
        ks_dfs(xx,yy,g+1,x*KB+y);
        swap(ks[x][y],ks[xx][yy]);
        if(ks_sol!=-1)return;
    }
}
\end{lstlisting}
\note{返回最少步数，无解（超过 maxdepth）返回 -\allowbreak{}1}
\begin{lstlisting}
int ks_solve(int maxdepth)
{
    ks_sol=-1;
    for(int lim=0;lim<=maxdepth;lim++)
    {
        ks_lim=lim;
        ks_dfs(ks_sx,ks_sy,0,-1);
        if(ks_sol!=-1)return ks_sol;
    }
    return -1;
}
\end{lstlisting}
\note{暴力 BFS 求真实最少步数，用于对拍（状态空间小的时候可以用）}
\note{=\allowbreak{}=\allowbreak{}=\allowbreak{}=\allowbreak{}=\allowbreak{}=\allowbreak{}=\allowbreak{}=\allowbreak{}=\allowbreak{}=\allowbreak{} 三、IDA*：九连环型旋钮（洛谷 P5507 思路）=\allowbreak{}=\allowbreak{}=\allowbreak{}=\allowbreak{}=\allowbreak{}=\allowbreak{}=\allowbreak{}=\allowbreak{}=\allowbreak{}=\allowbreak{} n 个旋钮各 4 个状态，一次操作把旋钮 i 和它指向的旋钮 t[\allowbreak{}i]\allowbreak{} 都 +\allowbreak{}1（模 4）。 估价：每个旋钮到 1 的最少次数之和的一半上取整。 每步总共 +\allowbreak{}1 两次（允许 t[\allowbreak{}i]\allowbreak{}=\allowbreak{}=\allowbreak{}i）；连续操作同一旋钮不能跳过。}
\begin{lstlisting}
int cr_n;
int cr_s[20],cr_t[20];
int cr_sol,cr_lim;
const int DIST1[4]={1,0,3,2};//状态 0/1/2/3 变到 1 的最少次数
int cr_h()
{
    int sum=0;
    for(int i=1;i<=cr_n;i++)sum+=DIST1[cr_s[i]];
    return (sum+1)/2;
}
void cr_dfs(int g,int pre)
{
    int h=cr_h();
    if(h==0){cr_sol=g;return;}
    if(g+h>cr_lim)return;
    for(int i=1;i<=cr_n;i++)
    {
        cr_s[i]=(cr_s[i]+1)&3;
        cr_s[cr_t[i]]=(cr_s[cr_t[i]]+1)&3;
        cr_dfs(g+1,i);
        cr_s[cr_t[i]]=(cr_s[cr_t[i]]+3)&3;
        cr_s[i]=(cr_s[i]+3)&3;
        if(cr_sol!=-1)return;
    }
}
int cr_solve(int maxdepth)
{
    cr_sol=-1;
    for(int lim=cr_h();lim<=maxdepth;lim++)//下界从 h 开始，别从 0 浪费
    {
        cr_lim=lim;
        cr_dfs(0,0);
        if(cr_sol!=-1)return cr_sol;
    }
    return -1;
}
\end{lstlisting}
% SOURCE: 07-搜索/A星与K短路.cpp
\subsection{A星与K短路}

\note{g/\allowbreak{}rg分别为正图和反图，st/\allowbreak{}en为源汇，k为排名；h[\allowbreak{}u]\allowbreak{}在反图中求得，是u到en的最短距离，vis记录各点弹出次数。边权须满足对应Dijkstra前提，重建图及计数后运行；此实现允许重复经过点。}
\begin{lstlisting}
const int N=1005;
const ll INF=4000000000000000000LL;
\end{lstlisting}
\note{=\allowbreak{}=\allowbreak{}=\allowbreak{}=\allowbreak{}=\allowbreak{}=\allowbreak{}=\allowbreak{}=\allowbreak{}=\allowbreak{}=\allowbreak{}=\allowbreak{}=\allowbreak{}=\allowbreak{}=\allowbreak{}=\allowbreak{}=\allowbreak{}=\allowbreak{} 一、A* 求第 k 短路 =\allowbreak{}=\allowbreak{}=\allowbreak{}=\allowbreak{}=\allowbreak{}=\allowbreak{}=\allowbreak{}=\allowbreak{}=\allowbreak{}=\allowbreak{}=\allowbreak{}=\allowbreak{}=\allowbreak{}=\allowbreak{}=\allowbreak{}=\allowbreak{}=\allowbreak{} h(\allowbreak{}x)\allowbreak{} =\allowbreak{} x 到终点的最短路（在反图上从终点 Dijkstra 一次预处理得到），可采纳 每次从堆里弹出「到终点第 i 次」的路径就是第 i 短路；O(\allowbreak{}k*E*log)\allowbreak{} 注意：与点 k 短路不同，这里按路径长度非递减计数，同长度不同路径要分别算 边权非负；s=\allowbreak{}=\allowbreak{}t 时计入长度 0 的空走法（排除它就把 k 加 1）。 本模板允许路径重复经过同一个点（标准 k 短路定义），所以不能用 vis 砍点}
\begin{lstlisting}
struct Edge
{
    int to,w;
};
struct Graph
{
    vector<Edge> adj[N];
    void init(int n){for(int i=0;i<=n;i++)adj[i].clear();}
    void add_edge(int u,int v,int w){adj[u].push_back({v,w});}
};
Graph g,rg;
int n,m,k,st,en;
ll h[N];//h[i] = i 到 en 的最短路
int vis[N];
\end{lstlisting}
\note{反图 Dijkstra 预处理 h，O(\allowbreak{}m log n)\allowbreak{}}
\begin{lstlisting}
void dijkstra_rev(int src)
{
    for(int i=1;i<=n;i++)h[i]=INF,vis[i]=0;
    priority_queue<pair<ll,int>,vector<pair<ll,int> >,greater<pair<ll,int> > > pq;
    h[src]=0;
    pq.push(make_pair(0,src));
    while(!pq.empty())
    {
        int u=pq.top().second;
        pq.pop();
        if(vis[u])continue;
        vis[u]=1;
        for(Edge e:rg.adj[u])
        {
            int v=e.to,w=e.w;
            if(h[v]>h[u]+w)
            {
                h[v]=h[u]+w;
                pq.push(make_pair(h[v],v));
            }
        }
    }
}
struct State
{
    int u; ll d;//d 是已经走过的实际距离
    bool operator>(const State &o)const{return d+h[u]>o.d+h[o.u];}//小根堆按 f=g+h
};
\end{lstlisting}
\note{求第 k 短路的长度；不足 k 条返回 -\allowbreak{}1}
\begin{lstlisting}
ll kth_shortest(int s,int t,int kk)
{
    if(kk<1)return -1;
    dijkstra_rev(t);
    if(h[s]==INF)return -1;//根本不连通
    priority_queue<State,vector<State>,greater<State> > pq;
    int cnt[N]={0};
    pq.push({s,0});
    while(!pq.empty())
    {
        State cur=pq.top();
        pq.pop();
        int u=cur.u;
        cnt[u]++;
        if(u==t&&cnt[u]==kk)return cur.d;
        if(cnt[u]>kk)continue;//同一个点被弹出超过 k 次，后面的一定不会更优
        for(Edge e:g.adj[u])
        {
            int v=e.to,w=e.w;
            if(h[v]<INF)pq.push({v,cur.d+w});
        }
    }
    return -1;
}
\end{lstlisting}
\note{=\allowbreak{}=\allowbreak{}=\allowbreak{}=\allowbreak{}=\allowbreak{}=\allowbreak{}=\allowbreak{}=\allowbreak{}=\allowbreak{}=\allowbreak{}=\allowbreak{}=\allowbreak{}=\allowbreak{}=\allowbreak{}=\allowbreak{}=\allowbreak{}=\allowbreak{} 三、A* 解八数码（洛谷 P1379）=\allowbreak{}=\allowbreak{}=\allowbreak{}=\allowbreak{}=\allowbreak{}=\allowbreak{}=\allowbreak{}=\allowbreak{}=\allowbreak{}=\allowbreak{}=\allowbreak{}=\allowbreak{}=\allowbreak{}=\allowbreak{}=\allowbreak{}=\allowbreak{}=\allowbreak{} 状态压成 64 位整数（每格 4 bit），h 用「不在目标位置的格子数」}
\begin{lstlisting}
ll pz_encode(int b[3][3])
{
    ll s=0;
    for(int i=0;i<3;i++)
        for(int j=0;j<3;j++)s|=(ll)b[i][j]<<(4*(8-(i*3+j)));
    return s;
}
void pz_decode(ll s,int b[3][3])
{
    for(int i=0;i<3;i++)
        for(int j=0;j<3;j++)b[i][j]=(s>>(4*(8-(i*3+j))))&15;
}
\end{lstlisting}
\note{八数码里空格向上下左右移动}
\begin{lstlisting}
const int pdx[]={0,-1,1,0,0};
const int pdy[]={0,0,0,-1,1};
\end{lstlisting}
\note{交换第 p 格与第 q 格（p,\allowbreak{}q 是 0.\allowbreak{}.\allowbreak{}8 的格子编号），O(\allowbreak{}1)\allowbreak{}}
\begin{lstlisting}
ll pz_flip(ll s,int p,int q)
{
    int b[3][3];
    pz_decode(s,b);
    swap(b[p/3][p%3],b[q/3][q%3]);
    return pz_encode(b);
}
struct PzState
{
    ll s;
    int g,h,z;
    bool operator>(const PzState &o)const{return g+h>o.g+o.h;}
};
int pz_h(ll s)//不在目标位置的格子数，可采纳
{
    const ll t=0x123456780ll;//encode(123456780) 的结果，空格的 0 在最低 4 bit
    int c=0;
    for(int i=0;i<9;i++)
    {
        int a=(s>>(4*i))&15,b=(t>>(4*i))&15;
        if(a&&a!=b)c++;
    }
    return c;
}
\end{lstlisting}
\note{A* 解八数码，h 用「不在目标位置的格子数」，O(\allowbreak{}状态数 log)\allowbreak{}}
\begin{lstlisting}
int astar_puzzle(int st[3][3])
{
    int goal[3][3]={{1,2,3},{4,5,6},{7,8,0}};
    const ll GOAL=pz_encode(goal);//目标态编码 = 0x123456780
    ll s0=pz_encode(st);
    if(s0==GOAL)return 0;
    int z0=0;
    for(int i=0;i<3;i++)
        for(int j=0;j<3;j++)
            if(!st[i][j])z0=i*3+j;
    unordered_map<ll,int> best;
    priority_queue<PzState,vector<PzState>,greater<PzState> > pq;
    best[s0]=0;
    PzState s1;
    s1.s=s0,s1.g=0,s1.h=pz_h(s0),s1.z=z0;
    pq.push(s1);
    while(!pq.empty())
    {
        PzState cur=pq.top();
        pq.pop();
        if(cur.s==GOAL)return cur.g;
        if(best[cur.s]<cur.g)continue;//过时状态
        int zx=cur.z/3,zy=cur.z%3;
        for(int i=1;i<=4;i++)
        {
            int xx=zx+pdx[i],yy=zy+pdy[i];
            if(xx<0||xx>2||yy<0||yy>2)continue;
            ll v=pz_flip(cur.s,zx*3+zy,xx*3+yy);
            int g2=cur.g+1;
            unordered_map<ll,int>::iterator it=best.find(v);
            if(it!=best.end()&&it->second<=g2)continue;
            best[v]=g2;
            PzState t2;
            t2.s=v,t2.g=g2,t2.h=pz_h(v),t2.z=xx*3+yy;
            pq.push(t2);
        }
    }
    return -1;
}
\end{lstlisting}
% SOURCE: 07-搜索/模拟退火与爬山.cpp
\subsection{模拟退火与爬山}

\note{当前状态与历史最佳状态分别保存，温度控制接受较差解概率；评价函数把状态映射为题目目标。随机种子、时间预算和降温系数按规模设置，返回最好记录，不能只输出最后状态。}
\begin{lstlisting}
struct Point
{
    double x,y;
};
\end{lstlisting}
\note{示例目标：有界区域中最小化凸二次函数；通用启发式不保证全局最优 更换目标函数后需同步调整扰动尺度、温度和迭代次数}
\begin{lstlisting}
struct Optimizer
{
    double cx,cy;
    mt19937 rng;
    Optimizer(double x,double y,unsigned seed):cx(x),cy(y),rng(seed){}
    double value(Point p)
    {
        return (p.x-cx)*(p.x-cx)+2*(p.y-cy)*(p.y-cy);
    }
    double rnd()
    {
        return uniform_real_distribution<double>(0,1)(rng);
    }
    Point clamp(Point p)
    {
        p.x=max(-10.0,min(10.0,p.x)),p.y=max(-10.0,min(10.0,p.y));
        return p;
    }
    Point anneal(Point cur)
    {
        Point best=cur;
        for(double temp=10;temp>1e-5;temp*=0.98)
            for(int i=1;i<=20;i++)
            {
                Point p=clamp({cur.x+(rnd()*2-1)*temp,cur.y+(rnd()*2-1)*temp});
                double delta=value(p)-value(cur);
                if(delta<=0||rnd()<exp(-delta/temp))cur=p;
                if(value(cur)<value(best))best=cur;
            }
        return best;
    }
    Point climb(Point cur)
    {
        for(double step=10;step>1e-8;step*=0.5)
        {
            bool change=true;
            while(change)
            {
                change=false;
                for(int dx=-1;dx<=1;dx++)
                    for(int dy=-1;dy<=1;dy++)
                    {
                        Point p=clamp({cur.x+dx*step,cur.y+dy*step});
                        if(value(p)+1e-18<value(cur))cur=p,change=true;
                    }
            }
        }
        return cur;
    }
};
\end{lstlisting}
% SOURCE: 07-搜索/DLX舞蹈链.cpp
\subsection{DLX舞蹈链}

\note{节点的左右上下链接构成稀疏0/\allowbreak{}1矩阵，列头计当前可用行数，行号对应方案项。add\_row给一行的非零列，dance输出所选行；覆盖列后必须逆序恢复，精确覆盖与可重复覆盖不能混用。}
\note{精确覆盖：每列恰好选一次；稀疏链表，最少候选列优先，最坏指数复杂度}
\begin{lstlisting}
struct Dlx
{
    vector<int> l,r,u,d,col,row,sz,ans;
    int m;
    void init(int n)
    {
        m=n;
        l.resize(n+1),r.resize(n+1),u.resize(n+1),d.resize(n+1);
        col.assign(n+1,0),row.assign(n+1,0),sz.assign(n+1,0),ans.clear();
        for(int i=0;i<=n;i++)l[i]=i-1,r[i]=i+1,u[i]=d[i]=i;
        l[0]=n,r[n]=0;
    }
    void add_row(int id,const vector<int> &a)
    {
        int first=0;
        for(int c:a)
        {
            assert(c>=1&&c<=m);
            int p=l.size();
            l.push_back(p),r.push_back(p),u.push_back(u[c]),d.push_back(c);
            col.push_back(c),row.push_back(id);
            d[u[c]]=p,u[c]=p,sz[c]++;
            if(!first)first=p;
            else l[p]=l[first],r[p]=first,r[l[first]]=p,l[first]=p;
        }
    }
    void remove(int c)
    {
        r[l[c]]=r[c],l[r[c]]=l[c];
        for(int i=d[c];i!=c;i=d[i])
            for(int j=r[i];j!=i;j=r[j])d[u[j]]=d[j],u[d[j]]=u[j],sz[col[j]]--;
    }
    void resume(int c)
    {
        for(int i=u[c];i!=c;i=u[i])
            for(int j=l[i];j!=i;j=l[j])sz[col[j]]++,d[u[j]]=j,u[d[j]]=j;
        r[l[c]]=c,l[r[c]]=c;
    }
    bool dance()
    {
        if(r[0]==0)return true;
        int c=r[0];
        for(int j=r[c];j;j=r[j])if(sz[j]<sz[c])c=j;
        remove(c);
        bool found=false;
        for(int i=d[c];i!=c&&!found;i=d[i])
        {
            ans.push_back(row[i]);
            for(int j=r[i];j!=i;j=r[j])remove(col[j]);
            found=dance();
            for(int j=l[i];j!=i;j=l[j])resume(col[j]);
            if(!found)ans.pop_back();
        }
        resume(c);// 返回后恢复链表，答案另外保存
        return found;
    }
    bool solve()
    {
        ans.clear();// 链表已恢复，允许重复求解
        return dance();
    }
};
\end{lstlisting}
\section{计算几何}
% SOURCE: 08-计算几何/向量基础.cpp
\subsection{向量基础}

\note{Point的x/\allowbreak{}y为坐标；多个点、线段或多边形是独立几何对象，保留参数。多边形顶点0起、首点不重复，判定接口区分内部/\allowbreak{}边界/\allowbreak{}外部；EPS仅用于几何判定，排序须严格弱序。}
\note{二维相似变换用复数：保持定向 z-\allowbreak{}>\allowbreak{}a*z+\allowbreak{}b，反转定向 z-\allowbreak{}>\allowbreak{}a*conj(\allowbreak{}z)\allowbreak{}+\allowbreak{}b。 p0!=\allowbreak{}p1 映到 q0,\allowbreak{}q1：a=\allowbreak{}(\allowbreak{}q1-\allowbreak{}q0)\allowbreak{}/\allowbreak{}(\allowbreak{}p1-\allowbreak{}p0)\allowbreak{}，b=\allowbreak{}q0-\allowbreak{}a*p0；退化点对不能除。}
\begin{lstlisting}
const double EPS=1e-9;
\end{lstlisting}
\note{点的编号从 0 开始（几何用 0-\allowbreak{}indexed 更自然：多边形顶点 0.\allowbreak{}.\allowbreak{}n-\allowbreak{}1）}
\begin{lstlisting}
struct Point
{
    double x,y;
    Point():x(0),y(0){}
    Point(double x,double y):x(x),y(y){}
    Point operator+(const Point &b)const{return Point(x+b.x,y+b.y);}
    Point operator-(const Point &b)const{return Point(x-b.x,y-b.y);}
    Point operator*(double k)const{return Point(x*k,y*k);}
    bool operator<(const Point &b)const{return x!=b.x?x<b.x:y<b.y;}
    bool operator==(const Point &b)const{return fabs(x-b.x)<EPS&&fabs(y-b.y)<EPS;}
};
int sgn(double x)//三态符号函数，所有判断都经过它，避免浮点直接比大小
{
    if(x>EPS)return 1;
    if(x<-EPS)return -1;
    return 0;
}
double dot(Point a,Point b){return a.x*b.x+a.y*b.y;}//点积，>0 锐角
double cross(Point a,Point b){return a.x*b.y-a.y*b.x;}//叉积，>0 表示 b 在 a 逆时针侧
double len(Point a){return sqrt(a.x*a.x+a.y*a.y);}
double len2(Point a){return a.x*a.x+a.y*a.y;}
\end{lstlisting}
\note{-\allowbreak{}-\allowbreak{}-\allowbreak{}-\allowbreak{}-\allowbreak{}-\allowbreak{}-\allowbreak{}-\allowbreak{}-\allowbreak{}-\allowbreak{} 关系判定 -\allowbreak{}-\allowbreak{}-\allowbreak{}-\allowbreak{}-\allowbreak{}-\allowbreak{}-\allowbreak{}-\allowbreak{}-\allowbreak{}-\allowbreak{}}
\begin{lstlisting}
bool is_parallel(Point a,Point b){return sgn(cross(a,b))==0;}//平行（含反向）
bool is_vertical(Point a,Point b){return sgn(dot(a,b))==0;}//垂直
\end{lstlisting}
\note{点 p 是否在线段 ab 上（含端点），O(\allowbreak{}1)\allowbreak{}}
\begin{lstlisting}
bool on_segment(Point a,Point b,Point p)
{
    if(sgn(cross(b-a,p-a))!=0)return false;//不在直线上
    return sgn(dot(p-a,p-b))<=0;//投影落在线段内
}
\end{lstlisting}
\note{直线 ab 与直线 cd 的交点；平行/\allowbreak{}重合时 ok=\allowbreak{}false}
\begin{lstlisting}
Point line_intersect(Point a,Point b,Point c,Point d,bool &ok)
{
    double s1=cross(b-a,c-a),s2=cross(b-a,d-a);
    if(sgn(s1-s2)==0&&sgn(s1)==0){ok=false;return Point();}//共线，无穷多交点
    if(sgn(s1-s2)==0){ok=false;return Point();}//平行
    ok=true;
    return c+(d-c)*(s1/(s1-s2));
}
\end{lstlisting}
\note{线段 ab 与线段 cd 是否相交（含端点、含共线重叠），O(\allowbreak{}1)\allowbreak{}}
\begin{lstlisting}
bool segment_cross(Point a,Point b,Point c,Point d)
{
    int d1=sgn(cross(b-a,c-a)),d2=sgn(cross(b-a,d-a));
    int d3=sgn(cross(d-c,a-c)),d4=sgn(cross(d-c,b-c));
    if(d1*d2<0&&d3*d4<0)return true;//规范相交
    if(d1==0&&on_segment(a,b,c))return true;//端点落在另一条线段上
    if(d2==0&&on_segment(a,b,d))return true;
    if(d3==0&&on_segment(c,d,a))return true;
    if(d4==0&&on_segment(c,d,b))return true;
    return false;
}
\end{lstlisting}
\note{点到线段的距离，O(\allowbreak{}1)\allowbreak{}}
\begin{lstlisting}
double point_seg_dis(Point a,Point b,Point p)
{
    if(sgn(dot(p-a,b-a))<=0)return len(p-a);//垂足在 a 外侧
    if(sgn(dot(p-b,a-b))<=0)return len(p-b);
    return fabs(cross(b-a,p-a))/len(b-a);
}
\end{lstlisting}
\note{点与简单多边形位置关系，射线法。 返回：0 外部，1 边界上，2 内部}
\begin{lstlisting}
int point_in_polygon(vector<Point> &p,Point q)
{
    int n=p.size();
    for(int i=0;i<n;i++)
        if(on_segment(p[i],p[(i+1)%n],q))return 1;
    int cnt=0;
    for(int i=0;i<n;i++)
    {
        Point a=p[i],b=p[(i+1)%n];
        if((a.y>q.y)!=(b.y>q.y))//跨过 q 的水平线
        {
            double x=a.x+(q.y-a.y)*(b.x-a.x)/(b.y-a.y);
            if(x>q.x)cnt++;
        }
    }
    return cnt&1?2:0;
}
\end{lstlisting}
\note{仅适用于凸多边形的内测：叉积同号，O(\allowbreak{}n)\allowbreak{}}
\begin{lstlisting}
bool in_convex_polygon(vector<Point> &p,Point q)//顶点逆时针
{
    int n=p.size();
    if(!n)return false;
    if(n==1)return q==p[0];
    if(n==2)return on_segment(p[0],p[1],q);
    for(int i=0;i<n;i++)
        if(sgn(cross(p[(i+1)%n]-p[i],q-p[i]))<0)return false;
    return true;
}
\end{lstlisting}
% SOURCE: 08-计算几何/凸包Andrew.cpp
\subsection{凸包Andrew}

\note{p为0起点集，返回逆时针凸包且不重复首点；普通版删除共线中间点，keep\_col版保留边界共线点。排序在副本上完成，面积是有向值，退化线段不表示二维区域。}
\begin{lstlisting}
const double EPS=1e-9;
const double PI=acos(-1.0);
struct Point
{
    double x,y;
    Point(){}
    Point(double x,double y):x(x),y(y){}
    Point operator+(const Point &b)const{return Point(x+b.x,y+b.y);}
    Point operator-(const Point &b)const{return Point(x-b.x,y-b.y);}
    Point operator*(double k)const{return Point(x*k,y*k);}
    bool operator<(const Point &b)const{return x!=b.x?x<b.x:y<b.y;}
    bool operator==(const Point &b)const{return fabs(x-b.x)<EPS&&fabs(y-b.y)<EPS;}
};
int sgn(double x)
{
    if(x>EPS)return 1;
    if(x<-EPS)return -1;
    return 0;
}
double cross(Point a,Point b){return a.x*b.y-a.y*b.x;}
double dot(Point a,Point b){return a.x*b.x+a.y*b.y;}
double len(Point a){return sqrt(a.x*a.x+a.y*a.y);}
double len2(Point a){return a.x*a.x+a.y*a.y;}
\end{lstlisting}
\note{-\allowbreak{}-\allowbreak{}-\allowbreak{}-\allowbreak{}-\allowbreak{}-\allowbreak{}-\allowbreak{}-\allowbreak{}-\allowbreak{}-\allowbreak{} 极角排序 -\allowbreak{}-\allowbreak{}-\allowbreak{}-\allowbreak{}-\allowbreak{}-\allowbreak{}-\allowbreak{}-\allowbreak{}-\allowbreak{}-\allowbreak{} atan2 精度差、常数大；用半平面 +\allowbreak{} 叉积排序，O(\allowbreak{}log)\allowbreak{} 常数极小 顺序：先按 y>\allowbreak{}=\allowbreak{}0 /\allowbreak{} y<\allowbreak{}0 分上下半平面，同半平面内按极角逆时针}
\note{排序比较不使用 EPS，避免破坏严格弱序；零向量无极角，实际应用应先去掉。}
\begin{lstlisting}
bool cmp_polar(Point a,Point b)
{
    int ha=a.y<0||(a.y==0&&a.x<0),hb=b.y<0||(b.y==0&&b.x<0);
    if(ha!=hb)return ha<hb;
    double v=cross(a,b);
    return v!=0?v>0:len2(a)<len2(b);
}
\end{lstlisting}
\note{也可以用 atan2，写法最短但有精度损失}
\begin{lstlisting}
bool cmp_atan2(Point a,Point b){return atan2(a.y,a.x)<atan2(b.y,b.x);}
\end{lstlisting}
\note{-\allowbreak{}-\allowbreak{}-\allowbreak{}-\allowbreak{}-\allowbreak{}-\allowbreak{}-\allowbreak{}-\allowbreak{}-\allowbreak{}-\allowbreak{} Andrew 单调链求凸包 -\allowbreak{}-\allowbreak{}-\allowbreak{}-\allowbreak{}-\allowbreak{}-\allowbreak{}-\allowbreak{}-\allowbreak{}-\allowbreak{}-\allowbreak{} 返回逆时针凸包顶点，无重复点；点集共线时返回两端点 O(\allowbreak{}n log n)\allowbreak{}，瓶颈是排序}
\begin{lstlisting}
vector<Point> convex_hull(vector<Point> p)
{
    int n=p.size();
    sort(p.begin(),p.end());
    n=unique(p.begin(),p.end())-p.begin();//先去重，否则共线点会出错
    p.resize(n);
    if(n<3)return p;
    vector<Point> h(2*n);
    int k=0;
    for(int i=0;i<n;i++)//下凸壳，叉积 <=0 弹出（共线点只留两端）
    {
        while(k>=2&&sgn(cross(h[k-1]-h[k-2],p[i]-h[k-2]))<=0)k--;
        h[k++]=p[i];
    }
    for(int i=n-2,t=k+1;i>=0;i--)//上凸壳
    {
        while(k>=t&&sgn(cross(h[k-1]-h[k-2],p[i]-h[k-2]))<=0)k--;
        h[k++]=p[i];
    }
    h.resize(k-1);//最后一点是起点，去掉
    return h;
}
\end{lstlisting}
\note{保留共线点的版本：把 <\allowbreak{}=\allowbreak{}0 改成 <\allowbreak{}0，凸包边上会留下中间点}
\begin{lstlisting}
vector<Point> convex_hull_keep_col(vector<Point> p)
{
    int n=p.size();
    sort(p.begin(),p.end());
    n=unique(p.begin(),p.end())-p.begin();
    p.resize(n);
    if(n<3)return p;
    bool all_line=true;
    for(Point x:p)if(sgn(cross(p.back()-p.front(),x-p.front()))!=0)all_line=false;
    if(all_line)return p;
    vector<Point> h(2*n);
    int k=0;
    for(int i=0;i<n;i++)
    {
        while(k>=2&&sgn(cross(h[k-1]-h[k-2],p[i]-h[k-2]))<0)k--;
        h[k++]=p[i];
    }
    for(int i=n-2,t=k+1;i>=0;i--)
    {
        while(k>=t&&sgn(cross(h[k-1]-h[k-2],p[i]-h[k-2]))<0)k--;
        h[k++]=p[i];
    }
    h.resize(k-1);
    return h;
}
double polygon_area(vector<Point> &p)//有向面积，逆时针为正
{
    double s=0;
    int n=p.size();
    for(int i=0;i<n;i++)s+=cross(p[i],p[(i+1)%n]);
    return s/2;
}
bool on_seg(Point a,Point b,Point p)
{
    if(sgn(cross(b-a,p-a))!=0)return false;
    return sgn(dot(p-a,p-b))<=0;
}
\end{lstlisting}
% SOURCE: 08-计算几何/旋转卡壳与多边形面积.cpp
\subsection{旋转卡壳与多边形面积}

\note{p为逆时针凸多边形，无重复首点和共线中间点；直径返回距离，宽度为两支撑线间距，面积有符号。点集先求凸包再卡壳，单点/\allowbreak{}线段退化按接口处理。}
\note{整体直径卡壳不直接给出每个顶点各自最远点；需要另证候选单调性或使用对应查询算法。}
\begin{lstlisting}
const double EPS=1e-9;
struct Point
{
    double x,y;
    Point(){}
    Point(double x,double y):x(x),y(y){}
    Point operator+(const Point &b)const{return Point(x+b.x,y+b.y);}
    Point operator-(const Point &b)const{return Point(x-b.x,y-b.y);}
    Point operator*(double k)const{return Point(x*k,y*k);}
    bool operator<(const Point &b)const{return x!=b.x?x<b.x:y<b.y;}
    bool operator==(const Point &b)const{return fabs(x-b.x)<EPS&&fabs(y-b.y)<EPS;}
};
int sgn(double x)
{
    if(x>EPS)return 1;
    if(x<-EPS)return -1;
    return 0;
}
double cross(Point a,Point b){return a.x*b.y-a.y*b.x;}
double dot(Point a,Point b){return a.x*b.x+a.y*b.y;}
double len(Point a){return sqrt(a.x*a.x+a.y*a.y);}
double len2(Point a){return a.x*a.x+a.y*a.y;}
\end{lstlisting}
\note{极角排序，半平面 +\allowbreak{} 叉积，O(\allowbreak{}1)\allowbreak{} 比较 排序比较不使用 EPS，避免破坏严格弱序；零向量无极角，实际应用应先去掉。}
\begin{lstlisting}
bool cmp_polar(Point a,Point b)
{
    int ha=a.y<0||(a.y==0&&a.x<0),hb=b.y<0||(b.y==0&&b.x<0);
    if(ha!=hb)return ha<hb;
    double v=cross(a,b);
    return v!=0?v>0:len2(a)<len2(b);
}
\end{lstlisting}
\note{Andrew 单调链，返回逆时针凸包，O(\allowbreak{}n log n)\allowbreak{}}
\begin{lstlisting}
vector<Point> convex_hull(vector<Point> p)
{
    int n=p.size();
    sort(p.begin(),p.end());
    n=unique(p.begin(),p.end())-p.begin();
    p.resize(n);
    if(n<3)return p;
    vector<Point> h(2*n);
    int k=0;
    for(int i=0;i<n;i++)
    {
        while(k>=2&&sgn(cross(h[k-1]-h[k-2],p[i]-h[k-2]))<=0)k--;
        h[k++]=p[i];
    }
    for(int i=n-2,t=k+1;i>=0;i--)
    {
        while(k>=t&&sgn(cross(h[k-1]-h[k-2],p[i]-h[k-2]))<=0)k--;
        h[k++]=p[i];
    }
    h.resize(k-1);
    return h;
}
double polygon_area(vector<Point> &p)//有向面积，逆时针为正
{
    double s=0;
    int n=p.size();
    for(int i=0;i<n;i++)s+=cross(p[i],p[(i+1)%n]);
    return s/2;
}
\end{lstlisting}
\note{凸多边形周长}
\begin{lstlisting}
double polygon_perimeter(vector<Point> &p)
{
    double s=0;
    int n=p.size();
    for(int i=0;i<n;i++)s+=len(p[(i+1)%n]-p[i]);
    return s;
}
\end{lstlisting}
\note{=\allowbreak{}=\allowbreak{}=\allowbreak{}=\allowbreak{}=\allowbreak{}=\allowbreak{}=\allowbreak{}=\allowbreak{}=\allowbreak{}=\allowbreak{}=\allowbreak{}=\allowbreak{}=\allowbreak{}=\allowbreak{}=\allowbreak{}=\allowbreak{}=\allowbreak{} 旋转卡壳：凸包直径 =\allowbreak{}=\allowbreak{}=\allowbreak{}=\allowbreak{}=\allowbreak{}=\allowbreak{}=\allowbreak{}=\allowbreak{}=\allowbreak{}=\allowbreak{}=\allowbreak{}=\allowbreak{}=\allowbreak{}=\allowbreak{}=\allowbreak{}=\allowbreak{}=\allowbreak{} 单调旋转，O(\allowbreak{}n)\allowbreak{}。i 是起点，j 是当前离直线 p[\allowbreak{}i]\allowbreak{}-\allowbreak{}p[\allowbreak{}i+\allowbreak{}1]\allowbreak{} 最远的点， 「面积」用叉积的绝对值表示，不必真的除以 2。}
\begin{lstlisting}
double rotating_diameter(vector<Point> &p)
{
    int n=p.size();
    if(n<=1)return 0;
    if(n==2)return len(p[1]-p[0]);
    double ans=0;
    for(int i=0,j=1;i<n;i++)
    {
        Point a=p[i],b=p[(i+1)%n];
        while(fabs(cross(b-a,p[(j+1)%n]-a))>fabs(cross(b-a,p[j]-a)))j=(j+1)%n;
        ans=max(ans,len2(p[j]-p[i]));
        ans=max(ans,len2(p[j]-p[(i+1)%n]));
    }
    return sqrt(ans);
}
\end{lstlisting}
\note{凸包宽度：所有方向上的最小「平行支撑线间距」，也是旋转卡壳，O(\allowbreak{}n)\allowbreak{}}
\begin{lstlisting}
double min_width(vector<Point> &p)
{
    int n=p.size();
    if(n<=2)return 0;
    double ans=1e100;
    for(int i=0,j=1;i<n;i++)
    {
        Point a=p[i],b=p[(i+1)%n];
        while(fabs(cross(b-a,p[(j+1)%n]-a))>fabs(cross(b-a,p[j]-a)))j=(j+1)%n;
        ans=min(ans,fabs(cross(b-a,p[j]-a))/len(b-a));
    }
    return ans;
}
\end{lstlisting}
\note{=\allowbreak{}=\allowbreak{}=\allowbreak{}=\allowbreak{}=\allowbreak{}=\allowbreak{}=\allowbreak{}=\allowbreak{}=\allowbreak{}=\allowbreak{}=\allowbreak{}=\allowbreak{}=\allowbreak{}=\allowbreak{}=\allowbreak{}=\allowbreak{}=\allowbreak{} 点在凸包内判定 =\allowbreak{}=\allowbreak{}=\allowbreak{}=\allowbreak{}=\allowbreak{}=\allowbreak{}=\allowbreak{}=\allowbreak{}=\allowbreak{}=\allowbreak{}=\allowbreak{}=\allowbreak{}=\allowbreak{}=\allowbreak{}=\allowbreak{}=\allowbreak{}=\allowbreak{} 顶点必须逆时针且无重复；O(\allowbreak{}log n)\allowbreak{}：先二分出极角所在的三角形，再用叉积判 返回 1 内部/\allowbreak{}边界，0 外部}
\begin{lstlisting}
int in_convex(vector<Point> &p,Point q)
{
    int n=p.size();
    if(!n)return 0;
    if(n==1)return q==p[0];
    if(n==2)
    {
        if(sgn(cross(p[1]-p[0],q-p[0]))!=0)return 0;
        return sgn(dot(q-p[0],q-p[1]))<=0;
    }
    if(sgn(cross(p[1]-p[0],q-p[0]))==0)return sgn(dot(q-p[0],q-p[1]))<=0;
    if(sgn(cross(p[n-1]-p[0],q-p[0]))==0)return sgn(dot(q-p[0],q-p[n-1]))<=0;
    if(sgn(cross(p[1]-p[0],q-p[0]))<0)return 0;//在 p0p1 外侧
    if(sgn(cross(p[n-1]-p[0],q-p[0]))>0)return 0;//在 p0p(n-1) 另一侧
    int lo=1,hi=n-1;
    while(hi-lo>1)//找夹住 q 的两个相邻顶点
    {
        int mid=(lo+hi)>>1;
        if(sgn(cross(p[mid]-p[0],q-p[0]))>=0)lo=mid;
        else hi=mid;
    }
    return sgn(cross(p[(lo+1)%n]-p[lo],q-p[lo]))>=0;
}
\end{lstlisting}
% SOURCE: 08-计算几何/半平面交.cpp
\subsection{半平面交}

\note{每条有向直线保留左侧半平面，方向与角度用于排序；队列存当前边界及相邻交点。返回多边形顶点，无界交需另加题目边框，不能把人为截断结果当原区域。}
\begin{lstlisting}
const double EPS=1e-9;
struct Point
{
    double x,y;
    Point operator+(Point b)const{return {x+b.x,y+b.y};}
    Point operator-(Point b)const{return {x-b.x,y-b.y};}
    Point operator*(double k)const{return {x*k,y*k};}
};
double cross(Point a,Point b)
{
    return a.x*b.y-a.y*b.x;
}
struct Line
{
    Point p,v;
    double ang;
    Line(Point a,Point b):p(a),v(b-a),ang(atan2(v.y,v.x)){if(ang<0)ang+=2*acos(-1.0);}
};
bool inside(Line a,Point p)
{
    return cross(a.v,p-a.p)>=-EPS;
}
Point meet(Line a,Line b)
{
    return a.p+a.v*(cross(b.v,b.p-a.p)/cross(b.v,a.v));
}
\end{lstlisting}
\note{O(\allowbreak{}n log n)\allowbreak{}，保留有向直线左侧；方向非零，交集须有界（可显式加题目边界）。 空集及线段/\allowbreak{}点交集返回空多边形；不以任意大方框冒充无界交集。}
\begin{lstlisting}
vector<Point> half_plane(vector<Line> a)
{
    sort(a.begin(),a.end(),[](Line x,Line y){return x.ang<y.ang;});
    vector<Line> b;
    for(Line l:a)
    {
        if(!b.empty()&&fabs(cross(b.back().v,l.v))<EPS&&b.back().v.x*l.v.x+b.back().v.y*l.v.y>0)
        {
            if(cross(b.back().v,l.p-b.back().p)>0)b.back()=l;
        }
        else b.push_back(l);
    }
    if(b.size()>1&&fabs(cross(b.front().v,b.back().v))<EPS&&
       b.front().v.x*b.back().v.x+b.front().v.y*b.back().v.y>0)
    {
        if(cross(b.front().v,b.back().p-b.front().p)>0)b.front()=b.back();
        b.pop_back();
    }
    deque<Line> q;
    for(Line l:b)
    {
        while(q.size()>1&&!inside(l,meet(q[q.size()-2],q.back())))q.pop_back();
        while(q.size()>1&&!inside(l,meet(q[0],q[1])))q.pop_front();
        if(!q.empty()&&fabs(cross(q.back().v,l.v))<EPS)return {};
        q.push_back(l);
    }
    while(q.size()>2&&!inside(q.front(),meet(q[q.size()-2],q.back())))q.pop_back();
    while(q.size()>2&&!inside(q.back(),meet(q[0],q[1])))q.pop_front();
    if(q.size()<3||fabs(cross(q.front().v,q.back().v))<EPS)return {};
    vector<Point> p;
    for(int i=0;i<(int)q.size();i++)p.push_back(meet(q[i],q[(i+1)%q.size()]));
    double twice=0;
    for(int i=0;i<(int)p.size();i++)twice+=cross(p[i],p[(i+1)%p.size()]);
    if(fabs(twice)<EPS)return {};
    return p;
}
double area(vector<Point> p)
{
    double s=0;
    for(int i=0;i<(int)p.size();i++)s+=cross(p[i],p[(i+1)%p.size()]);
    return fabs(s)/2;
}
\end{lstlisting}
% SOURCE: 08-计算几何/圆的交点与公切线.cpp
\subsection{圆的交点与公切线}

\note{圆由圆心与非负半径确定，交点/\allowbreak{}切线返回点或直线对象；切线对应的切点分别属于两个圆。重合圆、内含、相切和零半径先按接口分类，不能用交点数概括无限交点。}
\begin{lstlisting}
const double EPS=1e-10;
struct Point
{
    double x,y;
    Point operator+(Point b)const{return {x+b.x,y+b.y};}
    Point operator-(Point b)const{return {x-b.x,y-b.y};}
    Point operator*(double k)const{return {x*k,y*k};}
};
struct Circle{Point p;double r;};// 所有接口要求 r>=0，EPS 按坐标尺度调整。
double dot(Point a,Point b){return a.x*b.x+a.y*b.y;}
Point turn(Point a){return {-a.y,a.x};}
double norm(Point a){return hypot(a.x,a.y);}
bool same(Point a,Point b){return norm(a-b)<=EPS;}
\end{lstlisting}
\note{圆与无限直线 ab 的交点，0/\allowbreak{}1/\allowbreak{}2 个；a=\allowbreak{}=\allowbreak{}b 时按一个点处理。 r>\allowbreak{}=\allowbreak{}0，坐标尺度适中；EPS 是容差，极端尺度需要按题调整。}
\begin{lstlisting}
vector<Point> circle_line(Circle c,Point a,Point b)
{
    Point v=b-a;
    double len=norm(v);
    if(len<=EPS)return fabs(norm(a-c.p)-c.r)<=EPS?vector<Point>{a}:vector<Point>{};
    v=v*(1/len);
    Point foot=a+v*dot(c.p-a,v);
    double dist=norm(foot-c.p);
    if(dist>c.r+EPS)return {};
    if(fabs(dist-c.r)<=EPS)return {foot};
    double h=sqrt(max(0.0,(c.r-dist)*(c.r+dist)));
    return {foot+v*h,foot-v*h};
}
\end{lstlisting}
\note{圆圆交点：相离/\allowbreak{}内含 0 个，相切 1 个，相交 2 个；重合时 infinite=\allowbreak{}true。 两个重合的零半径圆只有一个交点；其余重合圆有无穷交点。}
\begin{lstlisting}
vector<Point> circle_circle(Circle a,Circle b,bool &infinite)
{
    infinite=false;
    Point v=b.p-a.p;
    double d=norm(v);
    if(d<=EPS)
    {
        if(fabs(a.r-b.r)>EPS)return {};
        if(a.r<=EPS)return {a.p};
        infinite=true;return {};
    }
    if(d>a.r+b.r+EPS||d<fabs(a.r-b.r)-EPS)return {};
    v=v*(1/d);
    double x=(d+(a.r-b.r)*(a.r+b.r)/d)/2;
    Point p=a.p+v*x;
    double h2=(a.r-x)*(a.r+x);
    if(fabs(d-a.r-b.r)<=EPS||fabs(d-fabs(a.r-b.r))<=EPS)return {p};
    double h=sqrt(max(0.0,h2));
    return {p+turn(v)*h,p-turn(v)*h};
}
\end{lstlisting}
\note{所有公切线，返回各线在两圆上的切点对；外切线与内切线统一计算。 切点重合时，切线过该点，方向垂直圆心连线；不能直接用两切点作方向。 相离两圆 4 条，外切 3 条，相交 2 条，内切 1 条，严格内含 0 条。 重合圆 infinite=\allowbreak{}true；不同圆心的两点圆返回连接它们的直线。}
\begin{lstlisting}
vector<pair<Point,Point>> common_tangents(Circle a,Circle b,bool &infinite)
{
    Point d=b.p-a.p;
    double len=norm(d);
    infinite=false;
    if(len<=EPS){infinite=fabs(a.r-b.r)<=EPS;return {};}
    Point u=d*(1/len);
    vector<pair<Point,Point>> ans;
    for(int s:{1,-1})
    {
        double r=s*b.r,h=(a.r-r)/len;
        if(fabs(h)>1+EPS)continue;
        h=clamp(h,-1.0,1.0);
        double z=sqrt(max(0.0,1-h*h));
        for(int sign:{1,-1})
        {
            Point v=u*h+turn(u)*(z*sign);
            pair<Point,Point> line={a.p+v*a.r,b.p+v*r};
            bool duplicate=false;
            for(auto e:ans)if(same(e.first,line.first)&&same(e.second,line.second))duplicate=true;
            if(!duplicate)ans.push_back(line);
        }
    }
    return ans;
}
\end{lstlisting}
\note{两圆交面积，O(\allowbreak{}1)\allowbreak{}；相离/\allowbreak{}相切、包含、重合、零半径均可。角度用弧度。}
\begin{lstlisting}
double circle_intersection_area(Circle a,Circle b)
{
    double d=norm(a.p-b.p),pi=acos(-1.0);
    if(d>=a.r+b.r)return 0;
    if(d<=fabs(a.r-b.r))return pi*min(a.r,b.r)*min(a.r,b.r);
    double x=acos(clamp((d*d+a.r*a.r-b.r*b.r)/(2*d*a.r),-1.0,1.0));
    double y=acos(clamp((d*d+b.r*b.r-a.r*a.r)/(2*d*b.r),-1.0,1.0));
    return a.r*a.r*(x-sin(2*x)/2)+b.r*b.r*(y-sin(2*y)/2);
}
\end{lstlisting}
\note{圆与简单多边形交面积，O(\allowbreak{}n)\allowbreak{}：切分每条边，内部用三角形、外部用扇形。 顶点按边界顺序，不要求凸；逆时针返回正有向面积，顺时针为负，通常取 fabs。}
\begin{lstlisting}
double circle_polygon_area(Circle c,const vector<Point> &p)
{
    auto cross=[](Point a,Point b){return a.x*b.y-a.y*b.x;};double ans=0;
    if(c.r<=0)return 0;
    for(int i=0;i<(int)p.size();i++)
    {
        Point a=p[i]-c.p,b=p[(i+1)%p.size()]-c.p,d=b-a;
        double A=dot(d,d);if(A==0)continue;
        vector<double> ts={0,1};double B=dot(a,d),C=dot(a,a)-c.r*c.r,D=B*B-A*C;
        if(D>0)for(double t:{(-B-sqrt(D))/A,(-B+sqrt(D))/A})if(t>0&&t<1)ts.push_back(t);
        sort(ts.begin(),ts.end());
        for(int j=1;j<(int)ts.size();j++)
        {
            Point u=a+d*ts[j-1],v=a+d*ts[j],mid=a+d*((ts[j-1]+ts[j])/2);
            if(dot(mid,mid)<=c.r*c.r)ans+=cross(u,v)/2;
            else ans+=c.r*c.r*atan2(cross(u,v),dot(u,v))/2;
        }
    }
    return ans;
}
\end{lstlisting}
% SOURCE: 08-计算几何/平面最近点对.cpp
\subsection{平面最近点对}

\note{a存0起平面点，tmp为等长归并工作区；closest\_pair(\allowbreak{})\allowbreak{}返回距离，并把a最终改为按y排序。solve(\allowbreak{}l,\allowbreak{}r)\allowbreak{}处理半开区间，入口按x排序、出口按y排序；少于两点返回正无穷，重合点返回0，每次入口重排a并调整tmp。}
\begin{lstlisting}
struct Point
{
    double x,y;
};
vector<Point> a,tmp;
double dis2(Point a,Point b)
{
    return (a.x-b.x)*(a.x-b.x)+(a.y-b.y)*(a.y-b.y);
}
bool cmp_y(Point a,Point b)
{
    return a.y<b.y;
}
\end{lstlisting}
\note{递归入口按 x 排序，出口按 y 排序；归并与带内扫描 O(\allowbreak{}n)\allowbreak{}。}
\begin{lstlisting}
double solve(int l,int r)
{
    if(r-l<=3)
    {
        double ans=numeric_limits<double>::infinity();
        for(int i=l;i<r;i++)for(int j=l;j<i;j++)ans=min(ans,dis2(a[i],a[j]));
        sort(a.begin()+l,a.begin()+r,cmp_y);
        return ans;
    }
    int mid=(l+r)>>1;
    double x=a[mid].x;
    double ans=min(solve(l,mid),solve(mid,r));
    merge(a.begin()+l,a.begin()+mid,a.begin()+mid,a.begin()+r,tmp.begin()+l,cmp_y);
    copy(tmp.begin()+l,tmp.begin()+r,a.begin()+l);
    vector<Point> p;
    for(int i=l;i<r;i++)
        if((a[i].x-x)*(a[i].x-x)<ans)
        {
            for(int j=(int)p.size()-1;j>=0&&(a[i].y-p[j].y)*(a[i].y-p[j].y)<ans;j--)ans=min(ans,dis2(a[i],p[j]));
            p.push_back(a[i]);
        }
    return ans;
}
\end{lstlisting}
\note{O(\allowbreak{}n log n)\allowbreak{}，最近点对距离；少于两点返回正无穷，重合点返回 0。}
\begin{lstlisting}
double closest_pair()
{
    sort(a.begin(),a.end(),[](Point x,Point y){return x.x!=y.x?x.x<y.x:x.y<y.y;});
    tmp.resize(a.size());
    return sqrt(solve(0,a.size()));
}
\end{lstlisting}
% SOURCE: 08-计算几何/最小圆覆盖.cpp
\subsection{最小圆覆盖}

\note{a是点集，Circle\{p,\allowbreak{}r\}为圆心及半径，返回覆盖圆；随机打乱在副本上，原输入不变。two为两点直径圆，three为三点外接圆或共线退化圆，不能将非共线three直接当三点最小覆盖圆。}
\begin{lstlisting}
const double EPS=1e-8;
struct Point
{
    double x,y;
};
struct Circle
{
    Point p;
    double r;
};
double dis(Point a,Point b)
{
    return hypot(a.x-b.x,a.y-b.y);
}
bool inside(Circle c,Point p)
{
    return dis(c.p,p)<=c.r+EPS;
}
Circle two(Point a,Point b)
{
    Point p={(a.x+b.x)/2,(a.y+b.y)/2};
    return {p,dis(a,b)/2};
}
Circle three(Point a,Point b,Point c)
{
    double x=b.x-a.x,y=b.y-a.y,u=c.x-a.x,v=c.y-a.y;
    double d=2*(x*v-y*u);
    if(fabs(d)<1e-12)
    {
        Circle z=two(a,b);
        if(dis(a,c)>2*z.r)z=two(a,c);
        if(dis(b,c)>2*z.r)z=two(b,c);
        return z;// 共线时最远两点作直径
    }
    double s=x*x+y*y,t=u*u+v*v;
    Point p={a.x+(s*v-t*y)/d,a.y+(x*t-u*s)/d};
    return {p,dis(p,a)};
}
\end{lstlisting}
\note{随机增量，期望 O(\allowbreak{}n)\allowbreak{}；空集半径 0，重复/\allowbreak{}共线点可用。}
\begin{lstlisting}
Circle min_circle(vector<Point> a)
{
    mt19937 rng(19260817);
    shuffle(a.begin(),a.end(),rng);
    Circle c={{0,0},0};
    for(int i=0;i<(int)a.size();i++)
        if(!inside(c,a[i]))
        {
            c={a[i],0};
            for(int j=0;j<i;j++)
                if(!inside(c,a[j]))
                {
                    c=two(a[i],a[j]);
                    for(int k=0;k<j;k++)
                        if(!inside(c,a[k]))c=three(a[i],a[j],a[k]);
                }
        }
    return c;
}
\end{lstlisting}
% SOURCE: 08-计算几何/闵可夫斯基和.cpp
\subsection{闵可夫斯基和}

\topic{随机平移下的交面积期望}
\note{非退化凸多边形 $P,Q$，在交面积为正的平移向量域内均匀采样 $t$。该域为 $P+(-Q)$ 的内部，边界零测度，因此}
\[\mathbb E\,|P\cap(Q+t)|=\frac{|P|\,|Q|}{|P+(-Q)|}.\]
\note{交换积分：固定 $x\in P$，使 $x\in Q+t$ 的 $t$ 集合面积为 $|Q|$，总积分即 $|P||Q|$。将 $Q$ 的坐标取负后求闵可夫斯基和，逆时针顺序仍保留。不适用于任意矩形内均匀采样或随机旋转（2025 沈阳 G）。}
\note{平面凸区域面积 $S$、周长 $L$，先加半径 $r$ 的圆盘，再加半径 $R$ 的三维球：}
\begin{align*}
S'&=S+Lr+\pi r^2,\qquad L'=L+2\pi r,\\
V&=2RS'+\frac\pi2 R^2L'+\frac{4\pi}3R^3.
\end{align*}
\note{二维圆盘和三维球的半径不能直接相加。}

\note{a/\allowbreak{}b为两个独立逆时针凸多边形，边方向合并得到和多边形；先规范起点并处理共线边。差集用于相交判定时取a+\allowbreak{}(\allowbreak{}-\allowbreak{}b)\allowbreak{}，平移不改变边方向，空集与低维退化另按入口处理。}
\note{凸平面区域面积 S、周长 L，加半径 r 圆盘后：S'=\allowbreak{}S+\allowbreak{}L*r+\allowbreak{}pi*r\ensuremath{^2}，L'=\allowbreak{}L+\allowbreak{}2*pi*r。 单独将该平面区域加半径 R 的三维球，体积 V=\allowbreak{}2*R*S+\allowbreak{}(\allowbreak{}pi/\allowbreak{}2)\allowbreak{}*R\ensuremath{^2}*L+\allowbreak{}(\allowbreak{}4*pi/\allowbreak{}3)\allowbreak{}*R\ensuremath{^3}。 两步同时出现时先更新 S、L，再代入球半径；二维圆盘和三维球的半径不能直接相加。}
\note{随机平移交面积：非退化凸多边形 P、Q，在有正交面积的平移向量域内均匀采样 t。 域为 P+\allowbreak{}(\allowbreak{}-\allowbreak{}Q)\allowbreak{} 的内部，E[\allowbreak{}area(\allowbreak{}P交(\allowbreak{}Q+\allowbreak{}t)\allowbreak{})\allowbreak{}]\allowbreak{}=\allowbreak{}area(\allowbreak{}P)\allowbreak{}*area(\allowbreak{}Q)\allowbreak{}/\allowbreak{}area(\allowbreak{}P+\allowbreak{}(\allowbreak{}-\allowbreak{}Q)\allowbreak{})\allowbreak{}。 对每个 x属于P，积分指示[\allowbreak{}x属于Q+\allowbreak{}t]\allowbreak{}，可行 t 的面积均为 area(\allowbreak{}Q)\allowbreak{}；交换积分即得分子。 需先把 Q 各点取负（不是反转点序），再调用 minkowski；不适用于任意矩形内采样或随机旋转。}
\begin{lstlisting}
struct Point
{
    ll x,y;
    Point operator+(Point b)const{return {x+b.x,y+b.y};}
    Point operator-(Point b)const{return {x-b.x,y-b.y};}
    bool operator<(Point b)const{return x!=b.x?x<b.x:y<b.y;}
    bool operator==(Point b)const{return x==b.x&&y==b.y;}
};
__int128 cross(Point a,Point b)
{
    return (__int128)a.x*b.y-(__int128)a.y*b.x;
}
\end{lstlisting}
\note{O(\allowbreak{}n log n)\allowbreak{}，去重并删除共线中间点；坐标加减须不溢出 ll，叉积用 \_\_int128。}
\begin{lstlisting}
vector<Point> hull(vector<Point> a)
{
    sort(a.begin(),a.end());
    a.erase(unique(a.begin(),a.end()),a.end());
    if(a.size()<3)return a;
    vector<Point> p;
    for(Point x:a)
    {
        while(p.size()>1&&cross(p.back()-p[p.size()-2],x-p.back())<=0)p.pop_back();
        p.push_back(x);
    }
    int k=p.size();
    for(int i=(int)a.size()-2;i>=0;i--)
    {
        while((int)p.size()>k&&cross(p.back()-p[p.size()-2],a[i]-p.back())<=0)p.pop_back();
        p.push_back(a[i]);
    }
    p.pop_back();
    return p;
}
void rotate_low(vector<Point> &a)
{
    int k=0;
    for(int i=1;i<(int)a.size();i++)if(a[i].y<a[k].y||(a[i].y==a[k].y&&a[i].x<a[k].x))k=i;
    rotate(a.begin(),a.begin()+k,a.end());
}
\end{lstlisting}
\note{O(\allowbreak{}n+\allowbreak{}m)\allowbreak{}，输入逆时针严格凸包（无重复首点）；单点/\allowbreak{}线段也支持。 任意点集先调用 hull；空集的和仍为空集。}
\begin{lstlisting}
vector<Point> minkowski(vector<Point> a,vector<Point> b)
{
    if(a.empty()||b.empty())return {};
    rotate_low(a),rotate_low(b);
    int n=a.size(),m=b.size(),i=0,j=0;
    Point p=a[0]+b[0];
    vector<Point> c={p};
    while(i<n||j<m)
    {
        Point x=i<n?a[(i+1)%n]-a[i]:Point{0,0};
        Point y=j<m?b[(j+1)%m]-b[j]:Point{0,0};
        __int128 v=cross(x,y);
        if(j==m||(i<n&&v>0))p=p+x,i++;
        else if(i==n||v<0)p=p+y,j++;
        else p=p+x+y,i++,j++;
        c.push_back(p);
    }
    c.pop_back();
    // 线性删除退化零边与共线点，正常输出已按边极角有序。
    vector<Point> d;
    for(Point x:c)
    {
        if(!d.empty()&&x==d.back())continue;
        while(d.size()>1&&cross(d.back()-d[d.size()-2],x-d.back())==0)d.pop_back();
        d.push_back(x);
    }
    while(d.size()>2&&cross(d.back()-d[d.size()-2],d.front()-d.back())==0)d.pop_back();
    if(d.size()>2&&cross(d.front()-d.back(),d[1]-d.front())==0)d.erase(d.begin());
    return d;
}
\end{lstlisting}
% SOURCE: 08-计算几何/多边形面积并.cpp
\subsection{多边形面积并}

\note{poly是多边形列表，每个顶点序列逆时针、简单且无零长边；可非凸。算法按边记录其他多边形遮蔽区间，编号用于共线同向边去重；返回并面积，原输入不变。}
\note{多个简单多边形的面积并，允许非凸，要求每个多边形顶点逆时针、无自交且无零长度边。 每条边按遮蔽区间的进入/\allowbreak{}离开事件扫描；只累加露在外面的边段，做 Green 积分。 重合且同向边只保留一份；反向边积分相消。O(\allowbreak{}N\ensuremath{^2} log N)\allowbreak{}，浮点 EPS 按尺度调整。 用有向边边界积分求面积并时，共线重叠边只能计一次；固定多边形编号作为去重优先级。}
\begin{lstlisting}
struct Point{double x,y;Point operator+(Point b)const{return {x+b.x,y+b.y};}Point operator-(Point b)const{return {x-b.x,y-b.y};}Point operator*(double t)const{return {x*t,y*t};}};
double cross(Point a,Point b){return a.x*b.y-a.y*b.x;}
double dot(Point a,Point b){return a.x*b.x+a.y*b.y;}
const double EPS=1e-9;
int inside(Point p,const vector<Point> &a)// 1 内部，0 边界，-1 外部
{
    bool in=false;for(int i=0;i<(int)a.size();i++){Point u=a[i],v=a[(i+1)%a.size()];if(fabs(cross(v-u,p-u))<=EPS&&dot(p-u,p-v)<=EPS)return 0;if((u.y>p.y)!=(v.y>p.y)){double x=u.x+(v.x-u.x)*(p.y-u.y)/(v.y-u.y);if(x>p.x)in=!in;}}
    return in?1:-1;
}
\end{lstlisting}
\note{沿每条边收集其他多边形遮蔽区间的进入/\allowbreak{}离开事件，O(\allowbreak{}N\ensuremath{^2} log N)\allowbreak{}。 算法参考 KACTL PolygonUnion.\allowbreak{}h（black\_horse2014/\allowbreak{}chilli）；EPS 按坐标尺度调整。}
\begin{lstlisting}
double polygon_union_area(const vector<vector<Point>> &poly)
{
    auto sign=[](double x){return (x>EPS)-(x<-EPS);};double ans=0;
    for(int i=0;i<(int)poly.size();i++)for(int k=0;k<(int)poly[i].size();k++)
    {
        Point a=poly[i][k],b=poly[i][(k+1)%poly[i].size()],v=b-a;double len=dot(v,v);if(len==0)continue;
        vector<pair<double,int>> events={{0,0},{1,0}};
        for(int j=0;j<(int)poly.size();j++)if(j!=i)for(int h=0;h<(int)poly[j].size();h++)
        {
            Point c=poly[j][h],d=poly[j][(h+1)%poly[j].size()];int sc=sign(cross(v,c-a)),sd=sign(cross(v,d-a));
            if(sc!=sd)
            {
                double sa=cross(d-c,a-c),sb=cross(d-c,b-c);
                if(min(sc,sd)<0&&fabs(sa-sb)>EPS)events.push_back({sa/(sa-sb),sign(sc-sd)});
            }
            else if(sc==0&&j<i&&dot(v,d-c)>0)
            {events.push_back({dot(c-a,v)/len,1});events.push_back({dot(d-a,v)/len,-1});}
        }
        for(auto &e:events)e.first=clamp(e.first,0.0,1.0);sort(events.begin(),events.end());
        int count=0;double last=0,visible=0;
        for(auto [at,delta]:events){if(count==0)visible+=at-last;count+=delta;last=at;}
        ans+=cross(a,b)*visible/2;
    }
    return ans;
}
\end{lstlisting}
% SOURCE: 08-计算几何/圆并面积.cpp
\subsection{圆并面积}

\note{c存Circle\{x,\allowbreak{}y,\allowbreak{}r\}，r>\allowbreak{}=\allowbreak{}0；每圆按角度合并遮蔽区间，累加可见圆弧积分。包含/\allowbreak{}重合圆自动去除，返回面积而非周长；角度跨越2\ensuremath{\pi}时拆区间，EPS按坐标尺度设置。}
\note{n 个圆的面积并，O(\allowbreak{}n\ensuremath{^2} log n)\allowbreak{}：对每个圆合并被其他圆遮住的角区间，积分可见圆弧。 自动去除重合/\allowbreak{}被包含的圆；r>\allowbreak{}=\allowbreak{}0，坐标尺度适中，EPS 按题调整。只返回面积并。}
\begin{lstlisting}
struct Circle{double x,y,r;};
double circle_union_area(const vector<Circle> &c)
{
    const double EPS=1e-10,PI=acos(-1.0);double ans=0;int n=c.size();
    for(int i=0;i<n;i++)
    {
        auto a=c[i];if(a.r<=0)continue;bool covered=false;vector<pair<double,double>> seg;
        for(int j=0;j<n;j++)if(i!=j)
        {
            auto b=c[j];double dx=b.x-a.x,dy=b.y-a.y,d=hypot(dx,dy);
            if(d+a.r<=b.r+EPS)
            {if(d>EPS||a.r<b.r-EPS||j<i){covered=true;break;}else continue;}
            if(d>=a.r+b.r-EPS||d+b.r<=a.r+EPS)continue;
            double theta=atan2(dy,dx);if(theta<0)theta+=2*PI;
            double h=acos(clamp((d*d+a.r*a.r-b.r*b.r)/(2*d*a.r),-1.0,1.0));
            double l=theta-h,r=theta+h;
            if(l<0)seg.push_back({l+2*PI,2*PI}),l=0;
            if(r>2*PI)seg.push_back({0,r-2*PI}),r=2*PI;
            seg.push_back({l,r});
        }
        if(covered)continue;sort(seg.begin(),seg.end());
        auto arc=[&](double l,double r){return (a.r*a.r*(r-l)+a.r*a.x*(sin(r)-sin(l))+a.r*a.y*(cos(l)-cos(r)))/2;};
        double end=0;for(auto [l,r]:seg){if(l>end)ans+=arc(end,l);end=max(end,r);}ans+=arc(end,2*PI);
    }
    return max(0.0,ans);
}
\end{lstlisting}
% SOURCE: 08-计算几何/三维几何基础.cpp
\subsection{三维几何基础}

\note{Point3为三维坐标，dot为点积，cross为叉积向量；直线、平面各由对应点和方向确定。平行、共面等退化先判定，法向量未归一时距离计算需除其长度。}
\begin{lstlisting}
const double EPS=1e-9;
struct Point
{
    double x,y,z;
    Point operator+(Point b)const{return {x+b.x,y+b.y,z+b.z};}
    Point operator-(Point b)const{return {x-b.x,y-b.y,z-b.z};}
    Point operator*(double k)const{return {x*k,y*k,z*k};}
};
double dot(Point a,Point b)
{
    return a.x*b.x+a.y*b.y+a.z*b.z;
}
Point cross(Point a,Point b)
{
    return {a.y*b.z-a.z*b.y,a.z*b.x-a.x*b.z,a.x*b.y-a.y*b.x};
}
double len(Point a)
{
    return sqrt(dot(a,a));
}
struct Plane
{
    Point p,n;// 平面上一点与非零法向量
};
\end{lstlisting}
\note{O(\allowbreak{}1)\allowbreak{}，三点建平面；共线或重复点返回 false。}
\begin{lstlisting}
bool make_plane(Point a,Point b,Point c,Plane &s)
{
    s={a,cross(b-a,c-a)};
    return len(s.n)>EPS;
}
\end{lstlisting}
\note{O(\allowbreak{}1)\allowbreak{}，点到平面距离及垂足；退化平面返回 false。}
\begin{lstlisting}
bool project(Point p,Plane s,Point &q,double &d)
{
    double v=dot(s.n,s.n);
    if(v<EPS*EPS)return false;
    double t=dot(p-s.p,s.n)/v;
    q=p-s.n*t,d=fabs(t)*sqrt(v);
    return true;
}
\end{lstlisting}
\note{O(\allowbreak{}1)\allowbreak{}，直线 a+\allowbreak{}t*v 与平面求交；0 无交，1 唯一交，2 直线在平面内，-\allowbreak{}1 退化。}
\begin{lstlisting}
int line_plane(Point a,Point v,Plane s,Point &q)
{
    if(len(v)<EPS||len(s.n)<EPS)return -1;
    double x=dot(s.n,v),y=dot(s.n,s.p-a);
    if(fabs(x)<=EPS*len(s.n)*len(v))return fabs(y)<=EPS*len(s.n)?2:0;
    q=a+v*(y/x);
    return 1;
}
\end{lstlisting}
\note{O(\allowbreak{}1)\allowbreak{}，点到直线距离；零方向退化成点。}
\begin{lstlisting}
double point_line(Point p,Point a,Point v)
{
    if(len(v)<EPS)return len(p-a);
    return len(cross(p-a,v))/len(v);
}
\end{lstlisting}
\note{绕坐标轴正方向按右手法则转 times*90 度，axis=\allowbreak{}0/\allowbreak{}1/\allowbreak{}2 为 x/\allowbreak{}y/\allowbreak{}z；避免 sin/\allowbreak{}cos 误差。}
\begin{lstlisting}
Point rotate_axis(Point p,int axis,int times)
{
    assert(0<=axis&&axis<3);times=(times%4+4)%4;
    while(times--)
        if(axis==0)p={p.x,-p.z,p.y};
        else if(axis==1)p={p.z,p.y,-p.x};
        else p={-p.y,p.x,p.z};
    return p;
}
\end{lstlisting}
\note{非原点轴先减轴上基点，再旋转再加回；维护旋转矩阵时，列向量先 A 后 B 为 B*A。 90 度矩阵：Rx=\allowbreak{}[\allowbreak{}1 0 0;0 0 -\allowbreak{}1;0 1 0]\allowbreak{}，Ry=\allowbreak{}[\allowbreak{}0 0 1;0 1 0;-\allowbreak{}1 0 0]\allowbreak{}，Rz=\allowbreak{}[\allowbreak{}0 -\allowbreak{}1 0;1 0 0;0 0 1]\allowbreak{}。}
% SOURCE: 08-计算几何/三维凸包.cpp
\subsection{三维凸包}

\note{p为0起三维点集，Hull3D的faces存三角面原点下标，dim为仿射维数，area/\allowbreak{}volume为面积/\allowbreak{}体积。三维面朝外；二维退化只给单面剖分及单面面积，不能当封闭表面积。}
\note{增量三维凸包，O(\allowbreak{}n\ensuremath{^2})\allowbreak{}；返回朝外三角面、维数、表面积、体积。允许重复/\allowbreak{}共面点。 低维退化：dim=\allowbreak{}0/\allowbreak{}1 面为空，dim=\allowbreak{}2 返回平面凸包的一份三角剖分，area 是单面面积、volume=\allowbreak{}0。 用固定初始四面体内部点定向；EPS 是尺度相关容差，浮点接近退化须按题调整。}
\begin{lstlisting}
struct Point3{long double x,y,z;Point3 operator+(Point3 b)const{return {x+b.x,y+b.y,z+b.z};}Point3 operator-(Point3 b)const{return {x-b.x,y-b.y,z-b.z};}Point3 operator*(long double t)const{return {x*t,y*t,z*t};}};
Point3 cross(Point3 a,Point3 b){return {a.y*b.z-a.z*b.y,a.z*b.x-a.x*b.z,a.x*b.y-a.y*b.x};}
long double dot(Point3 a,Point3 b){return a.x*b.x+a.y*b.y+a.z*b.z;}
long double norm3(Point3 a){return sqrtl(dot(a,a));}
struct Hull3D{int dim=0;vector<array<int,3>> faces;long double area=0,volume=0;};
Hull3D convex_hull_3d(const vector<Point3> &p,long double eps=1e-10L)
{
    Hull3D out;int n=p.size();if(n<2)return out;
    int a=0,b=0,c=0,d=0;for(int i=1;i<n;i++)if(norm3(p[i]-p[a])>norm3(p[b]-p[a]))b=i;
    if(norm3(p[b]-p[a])<=eps)return out;out.dim=1;
    for(int i=0;i<n;i++)if(norm3(cross(p[b]-p[a],p[i]-p[a]))>norm3(cross(p[b]-p[a],p[c]-p[a])))c=i;
    Point3 normal=cross(p[b]-p[a],p[c]-p[a]);if(norm3(normal)<=eps)return out;out.dim=2;
    normal=normal*(1/norm3(normal));for(int i=0;i<n;i++)if(fabsl(dot(normal,p[i]-p[a]))>fabsl(dot(normal,p[d]-p[a])))d=i;
    if(fabsl(dot(normal,p[d]-p[a]))<=eps)
    {
        Point3 ex=(p[b]-p[a])*(1/norm3(p[b]-p[a])),ey=cross(normal,ex);vector<tuple<long double,long double,int>> v;
        for(int i=0;i<n;i++)v.push_back({dot(p[i]-p[a],ex),dot(p[i]-p[a],ey),i});sort(v.begin(),v.end());
        v.erase(unique(v.begin(),v.end(),[&](auto x,auto y){return fabsl(get<0>(x)-get<0>(y))<=eps&&fabsl(get<1>(x)-get<1>(y))<=eps;}),v.end());
        auto turn=[&](int i,int j,int k){auto [x,y,id]=v[i];auto [X,Y,jd]=v[j];auto [u,w,kd]=v[k];return (X-x)*(w-y)-(Y-y)*(u-x);};vector<int> h;
        for(int i=0;i<(int)v.size();i++){while(h.size()>1&&turn(h[h.size()-2],h.back(),i)<=eps)h.pop_back();h.push_back(i);}int lower=h.size();
        for(int i=(int)v.size()-2;i>=0;i--){while((int)h.size()>lower&&turn(h[h.size()-2],h.back(),i)<=eps)h.pop_back();h.push_back(i);}h.pop_back();
        for(int i=1;i+1<(int)h.size();i++)out.faces.push_back({get<2>(v[h[0]]),get<2>(v[h[i]]),get<2>(v[h[i+1]])});
    }
    else
    {
        out.dim=3;Point3 inner=(p[a]+p[b]+p[c]+p[d])*.25L;
        auto face=[&](int u,int v,int w){if(dot(cross(p[v]-p[u],p[w]-p[u]),inner-p[u])>0)swap(v,w);return array<int,3>{u,v,w};};
        out.faces={face(a,b,c),face(a,b,d),face(a,c,d),face(b,c,d)};
        for(int i=0;i<n;i++)if(i!=a&&i!=b&&i!=c&&i!=d)
        {
            vector<array<int,3>> keep;map<pair<int,int>,int> horizon;
            for(auto f:out.faces)
            {auto [u,v,w]=f;Point3 no=cross(p[v]-p[u],p[w]-p[u]);if(dot(no,p[i]-p[u])>eps*norm3(no))
                {for(auto [x,y]:{pair(u,v),pair(v,w),pair(w,u)}){auto it=horizon.find({y,x});if(it!=horizon.end())horizon.erase(it);else horizon[{x,y}]=1;}}
                else keep.push_back(f);}
            for(auto [edge,count]:horizon)keep.push_back(face(edge.first,edge.second,i));out.faces.swap(keep);
        }
    }
    for(auto [u,v,w]:out.faces){out.area+=norm3(cross(p[v]-p[u],p[w]-p[u]))/2;if(out.dim==3)out.volume+=dot(p[u]-p[a],cross(p[v]-p[a],p[w]-p[a]))/6;}
    out.volume=fabsl(out.volume);return out;
}
\end{lstlisting}
% SOURCE: 08-计算几何/Delaunay与Voronoi.cpp
\subsection{Delaunay与Voronoi}

\note{input为0起点集，Delaunay三角返回原点编号；重复点仅保留第一份。Voronoi另传凸逆时针边界boundary，返回各点在边框内的区域，重复后续点为空；截断边框不是无界区域的真实边。}
\note{抛物面提升 Delaunay：三维下凸包投影，O(\allowbreak{}n\ensuremath{^2})\allowbreak{}，0-\allowbreak{}indexed 原点编号，重合点保留第一份。 共圆时任选合法对角线；全共线返回空。用 long double 与坐标归一化改善数值，不用于病态极近点。 Voronoi 用半平面裁剪：返回给定凸 CCW 边界内每个点的区域，重合后出现的点区域为空。 裁剪版总 O(\allowbreak{}n\ensuremath{^3})\allowbreak{} 最坏，适合中小规模；无界 Voronoi 必须另表示射线，不能直接把截断边界当真实边。}
\begin{lstlisting}
#include "三维凸包.cpp" // 复用三维增量凸包；打印时查同章，避免重复实现。
struct Point{long double x,y;Point operator+(Point b)const{return {x+b.x,y+b.y};}Point operator-(Point b)const{return {x-b.x,y-b.y};}Point operator*(long double t)const{return {x*t,y*t};}};
long double cross(Point a,Point b){return a.x*b.y-a.y*b.x;}
long double dot(Point a,Point b){return a.x*b.x+a.y*b.y;}
long double incircle(Point a,Point b,Point c,Point p)
{a=a-p;b=b-p;c=c-p;return dot(a,a)*cross(b,c)-dot(b,b)*cross(a,c)+dot(c,c)*cross(a,b);}
vector<array<int,3>> delaunay(const vector<Point> &input)
{
    vector<Point> p;vector<int> id;set<pair<long double,long double>> seen;
    for(int i=0;i<(int)input.size();i++)if(seen.insert({input[i].x,input[i].y}).second)p.push_back(input[i]),id.push_back(i);
    int n=p.size();if(n<3)return {};long double xmin=p[0].x,xmax=xmin,ymin=p[0].y,ymax=ymin;
    for(auto q:p)xmin=min(xmin,q.x),xmax=max(xmax,q.x),ymin=min(ymin,q.y),ymax=max(ymax,q.y);
    Point center={(xmin+xmax)/2,(ymin+ymax)/2};long double scale=max(xmax-xmin,ymax-ymin);if(scale==0)return {};
    vector<Point3> lifted;for(auto &q:p){q=(q-center)*(1/scale);lifted.push_back({q.x,q.y,dot(q,q)});}
    auto hull=convex_hull_3d(lifted,1e-18L);vector<array<int,3>> ans;
    for(auto [a,b,c]:hull.faces)
    {
        long double turn=cross(p[b]-p[a],p[c]-p[a]);
        if(fabsl(turn)<=1e-18L)continue;
        if(hull.dim==3&&turn>=0)continue;
        if(turn<0)swap(b,c);ans.push_back({id[a],id[b],id[c]});
    }
    return ans;
}
vector<vector<Point>> voronoi_cells(const vector<Point> &p,const vector<Point> &boundary)
{
    const long double EPS=1e-12L;int n=p.size();vector<vector<Point>> cell(n);
    for(int i=0;i<n;i++)
    {
        vector<Point>a=boundary;
        for(int j=0;j<n&&!a.empty();j++)if(i!=j)
        {
            Point normal=p[j]-p[i];long double norm=dot(normal,normal);if(norm==0){if(j<i)a.clear();continue;}
            auto side=[&](Point x){return dot(normal,x-p[i])-norm/2;};vector<Point>b;
            for(int k=0;k<(int)a.size();k++){Point u=a[k],v=a[(k+1)%a.size()];long double x=side(u),y=side(v);bool in=x<=EPS,next=y<=EPS;if(in)b.push_back(u);if(in!=next)b.push_back(u+(v-u)*(x/(x-y)));}a.swap(b);
        }
        cell[i]=a;
    }
    return cell;
}
\end{lstlisting}
\section{实现与调试}
% SOURCE: 09-其他/卡常技巧.cpp
\subsection{卡常技巧}

\note{性能例子的规模、数组和计时结果只用于说明连续访问、缓冲与常数；按实际数据测量。改变类型或压位前检查范围和内存，不能把演示机器的耗时当跨平台保证。}
\note{先定位时间和内存瓶颈，再比较相同输入、相同编译选项下的总耗时。复杂度和数据布局优先；快读写见第一章，位函数见位运算，字并行见bitset，避免在本节重复模板。}
\begin{lstlisting}

\end{lstlisting}
\topic{归一化模加}
\note{若 $0\le a,b<mod$ 且mod>\allowbreak{}0，a+\allowbreak{}b最多减一次mod。先扩宽再相加，避免int中间溢出；批量累加不满足这一范围时不能只减一次。}
\begin{lstlisting}
inline int add_mod(int a,int b,int mod){ll c=1LL*a+b;return c>=mod?c-mod:c;}
\end{lstlisting}
\topic{连续访问与循环展开}
\note{二维数组的最后一维连续，遍历时尽量让它作内层。结构体填充以sizeof实测，内存布局往往比局部算术改写更影响吞吐。下面保留四路展开求和作为可对照实现，输入a[\allowbreak{}1.\allowbreak{}.\allowbreak{}n]\allowbreak{}，总和必须容纳于ll；是否更快由编译器与实际输入决定。}
\begin{lstlisting}
ll sum_unrolled(int *a,int n)
{
    ll s=0;
    int i=1;
    for(;i+3<=n;i+=4)
    {
        s+=a[i],s+=a[i+1],s+=a[i+2],s+=a[i+3];
    }
    for(;i<=n;i++)s+=a[i];
    return s;
}
\end{lstlisting}
\topic{测量与编译边界}
\note{测量代码放在测试目录，结果必须被消费，避免整段计算被优化掉。常量乘除通常由编译器处理，inline不保证强制内联；不要把register、重复求值的MIN/\allowbreak{}MAX宏当作优化模板。编译选项按比赛环境确定，不默认启用本机特有指令集。}
% SOURCE: 09-其他/交互题模板.cpp
\subsection{交互题模板}

\note{询问参数按题目协议填写，读到的答复先检查失败/\allowbreak{}结束哨兵；每次完整输出后flush。缓冲与查询次数约束由协议规定，测试假评测器不进入手册，通信技巧见独立正文。}
\note{交互题输出通常先进入缓冲区，只写换行不保证评测机马上收到。每次询问后用endl或显式flush；关闭cin与cout的绑定后，更不能依赖读入时自动刷新。printf对应fflush(\allowbreak{}stdout)\allowbreak{}。}
\topic{输出与读回}
\note{询问格式、次数上限、最终答案和多测结束方式按题面写。读入失败或错误返回就退出；-\allowbreak{}1只是下面示例的约定。调试信息写cerr。}
\begin{lstlisting}
// 发一次询问并读回回答
int ask(const string &s)
{
    cout<<s<<endl;//endl 会 flush，等价于 cout<<s<<"\n"<<flush;
    int r;
    if(!(cin>>r)||r==-1)exit(0); // 本示例协议以 -1 表示错误；其他题按协议改
    return r;
}
void answer(int v)
{
    cout<<"! "<<v<<endl;
}
\end{lstlisting}
\topic{通信技巧}
\note{先数接收方必须区分多少种情况：若有M种，固定长度二进制编码至少需要ceil(\allowbreak{}log2 M)\allowbreak{}位；若接收方已有部分信息，只需区分该信息下仍可能的情况。}
\note{常用做法是只发送差异或缺失部分；用位掩码打包小范围字段；已知总和或异或时省掉最后一项。发送集合可利用有序差值、组合排名，发送排列可用排列排名。编解码顺序、字段位宽、取整规则必须一致；压缩收益按最坏情况算。}
\note{若题目允许接收方先询问，可利用双方已有信息缩小候选集，再发送候选编号。允许通信错误时才考虑校验或纠错；哈希仅能在允许概率错误且能控制碰撞概率时使用。}
% SOURCE: 09-其他/随机数与对拍.cpp
\subsection{随机数与对拍}

\note{rng维护随机状态，种子决定复现；小型gen/\allowbreak{}solve/\allowbreak{}brute示例仅展示生成、候选解和参考解的连接方式。失败时保存种子与输入，实际维护先覆盖手造边界，再做少量独立对照。}
\note{需要随机数时用下面的生成器；偶尔对拍，写一个小规模暴力作参考即可。}
\begin{lstlisting}

\end{lstlisting}
\topic{随机生成}
\note{randint(\allowbreak{}l,\allowbreak{}r)\allowbreak{}生成闭区间整数，要求l\ensuremath{\le}r。固定seed可复现；排列用iota后shuffle。}
\begin{lstlisting}
unsigned long long seed=chrono::steady_clock::now().time_since_epoch().count();
mt19937_64 rnd(seed);
ll randint(ll l,ll r){return uniform_int_distribution<ll>(l,r)(rnd);}
\end{lstlisting}
\topic{简单对拍}
\note{按题意补gen、solve、brute，发现不同就输出完整输入。先试几个手造边界，不必机械地跑很多轮；浮点结果按允许误差比较。}
\begin{lstlisting}
vector<ll> gen(); // 小数据生成器，主动覆盖重复、负数、极小规模等边界
ll solve(vector<ll> a);
ll brute(vector<ll> a);
bool stress(int rounds)
{
    rnd.seed(seed);
    for(int t=1;t<=rounds;t++)
    {
        vector<ll> a=gen();
        ll x=solve(a),y=brute(a);
        if(x!=y)
        {
            cerr<<"seed="<<seed<<" case="<<t<<" got="<<x<<" expected="<<y<<'\n';
            cerr<<a.size()<<'\n';
            for(ll v:a)cerr<<v<<' ';
            cerr<<'\n';
            return false; // 首次不同就停，保留种子和完整输入；别只打印答案
        }
    }
    return true;
}
\end{lstlisting}
\section{模型判据与常用结论}
\topic{基础技巧}
\par\textbf{交换论证贪心}\quad
先比较相邻两项交换前后的目标值，再用不增代价的次序排序\par
\textbf{信号}：任务排序、等待时间、乘积排序；\textbf{易错}：只按单项大小排序而未验证交换差；比较乘积溢出\par
\par\textbf{双指针前提}\quad
只有右端扩张及左端收缩对可行性的影响单调时，才能让指针不回退\par
\textbf{信号}：有序两数配对、非负区间和、覆盖子串；\textbf{易错}：带负数的和直接滑窗；重复字符计数漏减\par
\par\textbf{反悔贪心}\quad
单位耗时且收益非负的任务按截止时间排序，超过可选数量时撤销已选收益最小项\par
\textbf{信号}：每件耗时 1、有截止时间、最大总收益；\textbf{易错}：把最大收益问题误套最长耗时堆；本库工作调度只支持单位耗时\par
\par\textbf{单调栈}\quad
求左侧最近严格较小值时弹掉所有大于等于当前值的栈顶\par
\textbf{信号}：最近更大或更小、柱状图、可见关系；\textbf{易错}：严格与非严格混淆；相等高度需按计数语义处理\par
\topic{数据结构}
\par\textbf{结合律与区间聚合}\quad
区间结构先定义可结合的 merge 与单位元，再让查询按原序合并\par
\textbf{信号}：区间统计、修改与查询混合；\textbf{易错}：把不满足结合律的信息直接拼接；非交换操作调换左右\par
\par\textbf{区间乘加标记}\quad
新变换 x*m+\allowbreak{}a 合成到旧标记时 mul*=\allowbreak{}m、add=\allowbreak{}add*m+\allowbreak{}a\par
\textbf{信号}：同时区间乘和区间加、模意义区间和；\textbf{易错}：合成顺序反写；乘法更新忘乘旧 add；模乘中间溢出\par
\par\textbf{静态 ST 表}\quad
结合且幂等的操作可用两段重叠区间 O(\allowbreak{}1)\allowbreak{} 查询；区间和不满足幂等性。\par
\textbf{信号}：静态区间最大最小、无修改；\textbf{易错}：忘记区间含两端；将重叠方案用于求和\par
\par\textbf{颜色段均摊}\quad
split(\allowbreak{}r+\allowbreak{}1)\allowbreak{} 后 split(\allowbreak{}l)\allowbreak{}，赋值删除中间段并插入新段，碎片遍历不保证最坏快\par
\textbf{信号}：区间赋值、极长同色段、随机数据；\textbf{易错}：碎片遍历无最坏对数保证；赋值删除的迭代器不能再用，先分右端再分左端\par
\par\textbf{扫描线与离线}\quad
扫描坐标排序后维护上一条带的聚合状态，同坐标事件按查询语义成组处理\par
\textbf{信号}：矩形并面积、二维偏序、所有询问已知；\textbf{易错}：把离散排名当长度；同坐标处理时机不一致\par
\topic{字符串}
\par\textbf{KMP 与循环节}\quad
非空串候选周期 p=\allowbreak{}n-\allowbreak{}nxt[\allowbreak{}n]\allowbreak{}，p 整除 n 才是完整重复块，否则完整循环节为 n\par
\textbf{信号}：最小循环节、重复字符串、单模式匹配；\textbf{易错}：本库 nxt 按前缀长度索引且 nxt[\allowbreak{}0]\allowbreak{}=\allowbreak{}-\allowbreak{}1；匹配后不回退会漏重叠\par
\par\textbf{Manacher}\quad
本库插分隔符后的 rr 等于原串最长回文长度，回文出现数贡献为 (\allowbreak{}rr+\allowbreak{}1)\allowbreak{}/\allowbreak{}2\par
\textbf{信号}：最长回文子串、全部回文出现数；\textbf{易错}：把 rr 再乘二；哨兵出现在输入中；复用未清旧半径\par
\par\textbf{AC 自动机覆盖}\quad
BFS 建 fail 后把 fail 的终止信息并入当前状态，转移即可识别后缀匹配\par
\textbf{信号}：多个模式、全部包含、禁止子串；\textbf{易错}：只检查当前 Trie 终点；重复模式是否分开计数不明确\par
\par\textbf{有限状态与 DP}\quad
只保存决定未来合法性的有限历史，令 dp[\allowbreak{}i+\allowbreak{}1]\allowbreak{}[\allowbreak{}go(\allowbreak{}s,\allowbreak{}c)\allowbreak{}]\allowbreak{}+\allowbreak{}=\allowbreak{}dp[\allowbreak{}i]\allowbreak{}[\allowbreak{}s]\allowbreak{}\par
\textbf{信号}：历史可压成有限状态、按长度统计、禁止模式；\textbf{易错}：状态不含足够历史；ABC418G 仅作同标签入口，禁串示例不是该题题解\par
\topic{图论}
\par\textbf{非负权最短路}\quad
非负边权用 Dijkstra，堆项距离不等于当前 dis 时跳过过期项\par
\textbf{信号}：路径费用非负、稀疏图；\textbf{易错}：负边仍用 Dijkstra；INF 加权溢出\par
\par\textbf{Floyd 中转顺序}\quad
中转点 k 在最外层，dis[\allowbreak{}i]\allowbreak{}[\allowbreak{}j]\allowbreak{}=\allowbreak{}min(\allowbreak{}dis[\allowbreak{}i]\allowbreak{}[\allowbreak{}j]\allowbreak{},\allowbreak{}dis[\allowbreak{}i]\allowbreak{}[\allowbreak{}k]\allowbreak{}+\allowbreak{}dis[\allowbreak{}k]\allowbreak{}[\allowbreak{}j]\allowbreak{})\allowbreak{}\par
\textbf{信号}：n 小、多源距离、增量开放中转点；\textbf{易错}：循环次序换错；不可达 INF 参与加法\par
\par\textbf{强连通缩点}\quad
有向图的 SCC 缩点图必为 DAG，可在缩点后按拓扑序 DP\par
\textbf{信号}：互相可达、环内收益合并、压缩依赖；\textbf{易错}：把不在栈中的点用于 low 回退；缩点后自环未剔除\par
\par\textbf{二分图与奇环}\quad
无向图可二染色当且仅当不含奇环，每个连通块都要启动染色\par
\textbf{信号}：敌对双方、分组、相邻必须不同色；\textbf{易错}：仅遍历一个块；同色自环漏判\par
\par\textbf{最大流与二分匹配}\quad
二分图选互不冲突配对可建单位容量网络，最大流等于最大匹配\par
\textbf{信号}：最多配对、一人一项、行列互斥；\textbf{易错}：反向残量边没维护；二分答案容量未重建\par
\par\textbf{树遍历与子树区间}\quad
DFS 首次进入编号后，子树恰为 [\allowbreak{}in[\allowbreak{}u]\allowbreak{},\allowbreak{}out[\allowbreak{}u]\allowbreak{}]\allowbreak{}，整条任意路径通常不是单区间\par
\textbf{信号}：子树加、子树查询、祖先判定；\textbf{易错}：混用欧拉序与首次进入序；把路径也当连续区间\par
\topic{数学}
\par\textbf{扩欧与不定方程}\quad
a、b 不全为零时 ax+\allowbreak{}by=\allowbreak{}c 有整数解当且仅当 gcd(\allowbreak{}a,\allowbreak{}b)\allowbreak{} 整除 c\par
\textbf{信号}：同余方程、整数系数线性组合；\textbf{易错}：exgcd 对负参数可能返回负 gcd；最小非负解未规范\par
\par\textbf{模逆元}\quad
a 模 m 可逆当且仅当 gcd(\allowbreak{}a,\allowbreak{}m)\allowbreak{}=\allowbreak{}1，费马 a\textasciicircum{}(\allowbreak{}p-\allowbreak{}2)\allowbreak{} 仅用于素数 p\par
\textbf{信号}：模除法、组合公式；\textbf{易错}：对合数照搬费马；分母为零还求逆\par
\par\textbf{矩阵与线性递推}\quad
固定维线性递推写成状态乘转移矩阵，再二进制幂跳过大量相同步数\par
\textbf{信号}：Fibonacci、大 n、固定转移；\textbf{易错}：行列向量约定混淆；零次幂不是零矩阵\par
\par\textbf{高斯消元}\quad
逐列选非零主元消元，零系数行非零常数则无解，否则自由元数为未知数减秩\par
\textbf{信号}：联立线性方程、秩、唯一解判断；\textbf{易错}：浮点主元不取较大绝对值；模合数下主元不一定可逆\par
\par\textbf{容斥}\quad
多个禁条件的并集按非空子集交集计数，奇数子集加、偶数子集减\par
\textbf{信号}：至少一个、一个都不、条件数量小；\textbf{易错}：交集计算误用乘积；总数减并集又弄反符号\par
\par\textbf{整除分块}\quad
l\ensuremath{\le}n 时令 q=\allowbreak{}n/\allowbreak{}l、r=\allowbreak{}n/\allowbreak{}q，则 [\allowbreak{}l,\allowbreak{}r]\allowbreak{} 上 floor(\allowbreak{}n/\allowbreak{}i)\allowbreak{} 相同，不同商块共 O(\allowbreak{}sqrt(\allowbreak{}n)\allowbreak{})\allowbreak{} 个\par
\textbf{信号}：求和项只依赖 n/\allowbreak{}i、n 很大；\textbf{易错}：l 超过 n 后除零；同时两个商需取两者右端最小值\par
\par\textbf{无偏博弈与 SG}\quad
有限无环正常规则无偏游戏用 SG=\allowbreak{} mex(\allowbreak{}后继 SG)\allowbreak{}，SG=\allowbreak{}0 必败，独立子局 SG 取异或\par
\textbf{信号}：两人轮流、无法行动者输、独立子局；\textbf{易错}：误用于反常规则或有偏游戏；简单 Nim 结论套任意游戏\par
\topic{动态规划}
\par\textbf{01 与完全背包}\quad
01 背包容量倒序避免重复选，完全背包正序允许再次选；恰好装满须按目标将不可达置为相应无穷值\par
\textbf{信号}：每件一次、无限件、容量限制；\textbf{易错}：正逆序写反；至多容量与恰好容量边界混淆\par
\par\textbf{区间 DP}\quad
按长度递增计算 f[\allowbreak{}l]\allowbreak{}[\allowbreak{}r]\allowbreak{}=\allowbreak{}min\_k(\allowbreak{}f[\allowbreak{}l]\allowbreak{}[\allowbreak{}k]\allowbreak{}+\allowbreak{}f[\allowbreak{}k+\allowbreak{}1]\allowbreak{}[\allowbreak{}r]\allowbreak{}+\allowbreak{}cost(\allowbreak{}l,\allowbreak{}r)\allowbreak{})\allowbreak{}，环复制后取 n 长度窗口最优\par
\textbf{信号}：区间合并、最后一次划分、环形石子；\textbf{易错}：长度顺序错；环答案只取从 1 开始；max 与 min 初始化混淆\par
\par\textbf{树形 DP}\quad
树最大权独立集有 f[\allowbreak{}u]\allowbreak{}[\allowbreak{}1]\allowbreak{}=\allowbreak{}w[\allowbreak{}u]\allowbreak{}+\allowbreak{}\ensuremath{\Sigma}f[\allowbreak{}v]\allowbreak{}[\allowbreak{}0]\allowbreak{}，f[\allowbreak{}u]\allowbreak{}[\allowbreak{}0]\allowbreak{}=\allowbreak{}\ensuremath{\Sigma}max(\allowbreak{}f[\allowbreak{}v]\allowbreak{}[\allowbreak{}0]\allowbreak{},\allowbreak{}f[\allowbreak{}v]\allowbreak{}[\allowbreak{}1]\allowbreak{})\allowbreak{}\par
\textbf{信号}：父子不能同时选、树上最优选择；\textbf{易错}：先算父亲后算孩子；森林漏根；负权强制选点\par
\par\textbf{LIS 与 Dilworth}\quad
序列严格 LIS 长度等于最少不升子序列划分数，求严格 LIS 用 lower\_bound 更新最小尾值；最长不下降子序列改用 upper\_bound\par
\textbf{信号}：导弹拦截、偏序链划分、最少轨道；\textbf{易错}：严格上升与不下降混用；二维排序同值方向不处理\par
\par\textbf{期望线性性}\quad
E[\allowbreak{}\ensuremath{\Sigma}X\_i]\allowbreak{}=\allowbreak{}\ensuremath{\Sigma}E[\allowbreak{}X\_i]\allowbreak{} 不要求独立，但乘积期望一般不能直接拆\par
\textbf{信号}：随机贡献和、每位置概率、期望次数；\textbf{易错}：把期望的非线性函数当函数的期望；模概率分母不可逆\par
\par\textbf{分治优化的单调决策}\quad
分层转移 f[\allowbreak{}t]\allowbreak{}[\allowbreak{}i]\allowbreak{}=\allowbreak{}min\_j(\allowbreak{}f[\allowbreak{}t-\allowbreak{}1]\allowbreak{}[\allowbreak{}j]\allowbreak{}+\allowbreak{}w(\allowbreak{}j,\allowbreak{}i)\allowbreak{})\allowbreak{} 的代价矩阵满足 Monge 不等式时可选一致的单调最优决策\par
\textbf{信号}：同层决策随 i 不回退；先证明 Monge 或直接证明单调；\textbf{易错}：必须证明 Monge 或决策单调；随机无反例不能代替证明。\par
\par\textbf{Knuth 区间 DP 优化}\quad
区间合并型 DP 在 a\ensuremath{\le}b\ensuremath{\le}c\ensuremath{\le}d 时若 w(\allowbreak{}b,\allowbreak{}c)\allowbreak{}\ensuremath{\le}w(\allowbreak{}a,\allowbreak{}d)\allowbreak{} 且 w(\allowbreak{}a,\allowbreak{}c)\allowbreak{}+\allowbreak{}w(\allowbreak{}b,\allowbreak{}d)\allowbreak{}\ensuremath{\le}w(\allowbreak{}a,\allowbreak{}d)\allowbreak{}+\allowbreak{}w(\allowbreak{}b,\allowbreak{}c)\allowbreak{}，可取 opt[\allowbreak{}l]\allowbreak{}[\allowbreak{}r-\allowbreak{}1]\allowbreak{}\ensuremath{\le}opt[\allowbreak{}l]\allowbreak{}[\allowbreak{}r]\allowbreak{}\ensuremath{\le}opt[\allowbreak{}l+\allowbreak{}1]\allowbreak{}[\allowbreak{}r]\allowbreak{}\par
\textbf{信号}：非负石子合并等区间代价；不是所有区间 DP；\textbf{易错}：只适用于相应区间合并递推；统一并列最优决策的取法。\par
\topic{搜索}
\par\textbf{迭代加深}\quad
单位代价下逐增深度上限，首次找到的解深度最小，IDA* 以 g+\allowbreak{}h 超界剪枝\par
\textbf{信号}：解深浅、分支大、内存紧；\textbf{易错}：h 高估剩余最小步数；跨轮保留不合适的 visited\par
\par\textbf{A* 下界}\quad
非负边权下按 g+\allowbreak{}h 取最小；h 不高估且 h(\allowbreak{}终点)\allowbreak{}=\allowbreak{}0 时允许重开状态可保最优，永久关闭状态还需一致性。\par
\textbf{信号}：第 k 短路、代价搜索、可靠下界；\textbf{易错}：不一致的 h 配永久 visited 会错；K 短路不能把每点只出队一次。\par
\par\textbf{折半与双向搜索}\quad
子集分半枚举再匹配；限深可逆变换从起点正向、终点反向各搜半深度，再合并相同中间状态\par
\textbf{信号}：n 约 40 的子集、可逆变换且深度小；\textbf{易错}：反搜须反转规则；相同子集和的频次应相乘，不能只存存在标记。\par
\topic{计算几何}
\par\textbf{向量与行走可达}\quad
给定可任意转向的步长，终点距 d 可达当且仅当 max(\allowbreak{}0,\allowbreak{}2*最大步长-\allowbreak{}总长)\allowbreak{}\ensuremath{\le}d\ensuremath{\le}总长\par
\textbf{信号}：若干线段任意方向、能否到坐标；\textbf{易错}：只判三角上界漏下界；距离平方与距离混比\par
\topic{图论}
\par\textbf{竞赛图哈密顿路径}\quad
任意竞赛图都有一条经过每个顶点恰一次的有向哈密顿路径，但不保证哈密顿回路\par
\textbf{信号}：任意两点间恰有一个方向的边、求访问所有点的次序\par
\par\textbf{强连通竞赛图}\quad
至少 3 点的强连通竞赛图存在有向哈密顿回路\par
\textbf{信号}：完整有向对局关系且整体强连通\par
\par\textbf{竞赛图三角形}\quad
竞赛图若无有向三角形则无有向环，可按拓扑顺序形成传递次序\par
\textbf{信号}：检查竞赛图是否传递；不能说任意三点都有环\par
\par\textbf{DAG 最小路径覆盖}\quad
DAG 的不交顶点有向路径覆盖数等于 n-\allowbreak{}二分拆点图最大匹配数\par
\textbf{信号}：所有顶点恰好覆盖一次、允许单点路径；若要求可达链须先考虑传递闭包\par
\par\textbf{最大流最小割}\quad
有限非负容量网络中最大 s-\allowbreak{}t 流值等于最小 s-\allowbreak{}t 割容量\par
\textbf{信号}：瓶颈、最少代价切断联系、容量限制\par
\par\textbf{残量图最小割}\quad
最大流结束后，残量图从 s 可达点集与其补集构成一个最小割\par
\textbf{信号}：需要输出最小割方案\par
\par\textbf{割边方向}\quad
有向割容量只加从源侧到汇侧的原始边，反方向不计\par
\textbf{信号}：枚举小图最小割、对拍流网络\par
\par\textbf{容量积分性}\quad
整数容量网络存在整数最大流，单位容量的两层网络可解释为离散配对\par
\par\textbf{Kőnig 定理}\quad
二分图最小顶点覆盖大小等于最大匹配大小\par
\textbf{信号}：用最少点覆盖全部边\par
\par\textbf{二分图最大独立集}\quad
二分图最大独立集大小等于顶点数减最大匹配大小\par
\textbf{信号}：最多选点且任意两选点不相邻\par
\par\textbf{Hall 定理}\quad
二分图能匹配左部全部顶点当且仅当每个左部子集 S 都满足邻集大小至少为 S 的大小\par
\textbf{信号}：每个任务须获得不同资源、证明完美匹配存在\par
\par\textbf{一般图匹配边界}\quad
存在奇环时不能直接套二分图匈牙利算法，需一般图匹配或证明特殊结构\par
\textbf{信号}：非二分关系的最多配对\par
\par\textbf{平面图欧拉式}\quad
连通平面嵌入满足 V-\allowbreak{}E+\allowbreak{}F=\allowbreak{}2，F 包含外部面；c 个连通块时为 1+\allowbreak{}c\par
\par\textbf{简单平面图边数}\quad
简单平面图 n\ensuremath{\ge}3 时 m\ensuremath{\le}3n-\allowbreak{}6\par
\textbf{信号}：证明图不平面、估算边数\par
\par\textbf{二分平面图边数}\quad
简单二分平面图 n\ensuremath{\ge}3 时 m\ensuremath{\le}2n-\allowbreak{}4\par
\textbf{信号}：所有面无三角形、平面二染色\par
\par\textbf{无向欧拉路}\quad
忽略孤立点后图须连通，奇度点为 0 或 2 才存在欧拉路\par
\textbf{信号}：每条无向边恰走一次\par
\par\textbf{有向欧拉路}\quad
非零度顶点在无向底图连通，入出度全相等或仅起终点差为 ±1 才存在有向欧拉路\par
\par\textbf{Hierholzer 逆序}\quad
走到无未用边的点才退栈记录，所得序列逆序后才是欧拉遍历\par
\par\textbf{树重心块大小}\quad
树重心删除后每个连通块大小不超过 n/\allowbreak{}2，重心至多两个且若两个则相邻\par
\textbf{信号}：分治、最小化最大子树块\par
\par\textbf{树距离和最小点}\quad
无权树上所有点到某点的距离总和在重心处取最小\par
\textbf{信号}：无权树上设施选址；人口权不同改用加权重心\par
\par\textbf{加权点距离换根}\quad
边长 w 的父子 u-\allowbreak{}v，子树点权和 s、总点权 W 时，距离和 F[\allowbreak{}v]\allowbreak{}=\allowbreak{}F[\allowbreak{}u]\allowbreak{}+\allowbreak{}w*(\allowbreak{}W-\allowbreak{}2s)\allowbreak{}\par
\par\textbf{树直径两遍法}\quad
非负边权树中从任意点找最远点 a，再从 a 找最远点 b，可得直径\par
\textbf{信号}：树上最远距离、无负权边\par
\par\textbf{树上偏心距}\quad
非负边权树直径端点 a、b 满足每点最远距离等于 max(\allowbreak{}dist(\allowbreak{}u,\allowbreak{}a)\allowbreak{},\allowbreak{}dist(\allowbreak{}u,\allowbreak{}b)\allowbreak{})\allowbreak{}\par
\textbf{信号}：所有点到最远点的距离\par
\par\textbf{树距离 LCA}\quad
树上 dist(\allowbreak{}u,\allowbreak{}v)\allowbreak{}=\allowbreak{}D[\allowbreak{}u]\allowbreak{}+\allowbreak{}D[\allowbreak{}v]\allowbreak{}-\allowbreak{}2D[\allowbreak{}lca(\allowbreak{}u,\allowbreak{}v)\allowbreak{}]\allowbreak{}，D 是根距离\par
\par\textbf{树点路径差分}\quad
路径 u-\allowbreak{}v 的点计数加一可做 d[\allowbreak{}u]\allowbreak{}+\allowbreak{}+\allowbreak{}、d[\allowbreak{}v]\allowbreak{}+\allowbreak{}+\allowbreak{}、d[\allowbreak{}lca]\allowbreak{}-\allowbreak{}-\allowbreak{}、d[\allowbreak{}parent(\allowbreak{}lca)\allowbreak{}]\allowbreak{}-\allowbreak{}-\allowbreak{} 后子树汇总\par
\textbf{信号}：统计每点被多少条路径经过\par
\par\textbf{树边路径差分}\quad
路径 u-\allowbreak{}v 的边计数加一可做 d[\allowbreak{}u]\allowbreak{}+\allowbreak{}+\allowbreak{}、d[\allowbreak{}v]\allowbreak{}+\allowbreak{}+\allowbreak{}、d[\allowbreak{}lca]\allowbreak{}-\allowbreak{}=\allowbreak{}2 后子树汇总\par
\textbf{信号}：统计每条父子边经过次数\par
\par\textbf{桥判据}\quad
无向 DFS 树边 u-\allowbreak{}v 为桥当且仅当 low[\allowbreak{}v]\allowbreak{}>\allowbreak{}dfn[\allowbreak{}u]\allowbreak{}，父边按边号排除\par
\textbf{信号}：删除一条边断连、重边图\par
\par\textbf{割点判据}\quad
非根 u 若有孩子 v 满足 low[\allowbreak{}v]\allowbreak{}\ensuremath{\ge}dfn[\allowbreak{}u]\allowbreak{} 则为割点，根须至少两个 DFS 孩子\par
\textbf{信号}：删点连通性\par
\par\textbf{差分约束}\quad
x[\allowbreak{}v]\allowbreak{}-\allowbreak{}x[\allowbreak{}u]\allowbreak{}\ensuremath{\le}w 建 u→v 权 w，所有点从超级源可达后存在负环当且仅当无解\par
\textbf{信号}：大量相对上界、排程不等式\par
\par\textbf{DAG 最长路}\quad
DAG 按拓扑序松弛可求最长路，负边也可；不可达初始化为负无穷\par
\par\textbf{0-1 BFS}\quad
边权仅 0 和 1 时按松弛把零边放队首、一边放队尾，可在线性规模内求最短路\par
\textbf{信号}：免费操作与一次代价操作\par
\par\textbf{超源多源最短路}\quad
所有起点初始距离置零等价于添加零权超级源，统一求离各源最近距离\par
\textbf{信号}：多出口、多个补给站\par
\par\textbf{最短路计数零环}\quad
最短路计数不能无条件按距离递增 DP，零权等距边可能形成环\par
\textbf{信号}：边权允许零、路径数量\par
\par\textbf{MST 割性质}\quad
任意割的一条最小权跨割边属于某棵 MST，若唯一最小则属于每棵 MST\par
\par\textbf{MST 环性质}\quad
某环中唯一最大权边不属于任何 MST，存在并列最大时只能保证可避开其中一条\par
\par\textbf{MST 瓶颈性质}\quad
任意 MST 都是最小瓶颈生成树，反向不一定成立\par
\par\textbf{函数图结构}\quad
每点出度恰一的有限有向图，每个弱连通块有一个有向环及指向环的树\par
\par\textbf{2-SAT 判据}\quad
蕴含图中任何变量与其否定同 SCC 则无解，否则有解\par
\textbf{信号}：每条限制涉及两个布尔文字\par
\topic{数学}
\par\textbf{裴蜀多元}\quad
不全为零的整数 a\_i 所有整数线性组合恰为 gcd(\allowbreak{}a\_i)\allowbreak{} 的倍数\par
\par\textbf{同余方程解数}\quad
g=\allowbreak{}gcd(\allowbreak{}a,\allowbreak{}m)\allowbreak{}>\allowbreak{}0，ax\ensuremath{\equiv}b(\allowbreak{}mod m)\allowbreak{} 在 [\allowbreak{}0,\allowbreak{}m)\allowbreak{} 有解时恰有 g 个解，相邻相差 m/\allowbreak{}g\par
\par\textbf{CRT 互质}\quad
两两互质模数的同余组在模各模数乘积意义下存在唯一解\par
\textbf{信号}：多个互质周期同时对齐\par
\par\textbf{扩展 CRT}\quad
x\ensuremath{\equiv}a(\allowbreak{}mod m)\allowbreak{}、x\ensuremath{\equiv}b(\allowbreak{}mod n)\allowbreak{} 可合并当且仅当 gcd(\allowbreak{}m,\allowbreak{}n)\allowbreak{} 整除 b-\allowbreak{}a，合并模数是 lcm(\allowbreak{}m,\allowbreak{}n)\allowbreak{}\par
\textbf{信号}：模数不互质、周期碰面\par
\par\textbf{欧拉定理}\quad
m>\allowbreak{}1 且 gcd(\allowbreak{}a,\allowbreak{}m)\allowbreak{}=\allowbreak{}1 时 a\textasciicircum{}\ensuremath{\varphi}(\allowbreak{}m)\allowbreak{}\ensuremath{\equiv}1(\allowbreak{}mod m)\allowbreak{}，指数可按 \ensuremath{\varphi}(\allowbreak{}m)\allowbreak{} 化简。\par
\textbf{信号}：大指数取模且底数互质\par
\par\textbf{非互质欧拉降幂}\quad
m>\allowbreak{}1、b\ensuremath{\ge}\ensuremath{\varphi}(\allowbreak{}m)\allowbreak{} 时 a\textasciicircum{}b\ensuremath{\equiv}a\textasciicircum{}(\allowbreak{}b mod \ensuremath{\varphi}(\allowbreak{}m)\allowbreak{}+\allowbreak{}\ensuremath{\varphi}(\allowbreak{}m)\allowbreak{})\allowbreak{}(\allowbreak{}mod m)\allowbreak{}，小指数保留原值\par
\textbf{信号}：指数塔、底数不互质；必须记录指数是否达到阈值\par
\par\textbf{Lucas}\quad
素数 p 下 C(\allowbreak{}n,\allowbreak{}k)\allowbreak{} mod p 等于各 p 进制位 C(\allowbreak{}n\_i,\allowbreak{}k\_i)\allowbreak{} 的乘积\par
\par\textbf{阶乘逆元边界}\quad
素数 p 下阶乘逆元预处理要求最大下标小于 p，否则阶乘为零不可逆\par
\par\textbf{Catalan 递推与闭式}\quad
C\_0=\allowbreak{}1，C\_n=\allowbreak{}\ensuremath{\Sigma}C\_i*C\_(\allowbreak{}n-\allowbreak{}1-\allowbreak{}i)\allowbreak{}，计数合法括号或二叉树形状；C\_n=\allowbreak{}binom(\allowbreak{}2n,\allowbreak{}n)\allowbreak{}-\allowbreak{}binom(\allowbreak{}2n,\allowbreak{}n+\allowbreak{}1)\allowbreak{}=\allowbreak{}binom(\allowbreak{}2n,\allowbreak{}n)\allowbreak{}/\allowbreak{}(\allowbreak{}n+\allowbreak{}1)\allowbreak{}，模除法须检查逆元\par
\par\textbf{第二类 Stirling}\quad
S(\allowbreak{}n,\allowbreak{}k)\allowbreak{}=\allowbreak{}S(\allowbreak{}n-\allowbreak{}1,\allowbreak{}k-\allowbreak{}1)\allowbreak{}+\allowbreak{}k*S(\allowbreak{}n-\allowbreak{}1,\allowbreak{}k)\allowbreak{}，S(\allowbreak{}0,\allowbreak{}0)\allowbreak{}=\allowbreak{}1，计数无标号非空组划分\par
\textbf{信号}：把 n 个不同物品分成 k 个非空集合\par
\par\textbf{第一类 Stirling}\quad
无符号第一类 c(\allowbreak{}n,\allowbreak{}k)\allowbreak{}=\allowbreak{}c(\allowbreak{}n-\allowbreak{}1,\allowbreak{}k-\allowbreak{}1)\allowbreak{}+\allowbreak{}(\allowbreak{}n-\allowbreak{}1)\allowbreak{}c(\allowbreak{}n-\allowbreak{}1,\allowbreak{}k)\allowbreak{}，计数有 k 个循环的排列\par
\textbf{信号}：排列循环数、置换分解\par
\par\textbf{前缀和同余鸽巢}\quad
任意 n 个整数必有一个非空连续段，其和被 n 整除\par
\textbf{信号}：求存在性、前缀和模 n\par
\par\textbf{错排}\quad
D\_0=\allowbreak{}1、D\_1=\allowbreak{}0，D\_n=\allowbreak{}(\allowbreak{}n-\allowbreak{}1)\allowbreak{}*(\allowbreak{}D\_(\allowbreak{}n-\allowbreak{}1)\allowbreak{}+\allowbreak{}D\_(\allowbreak{}n-\allowbreak{}2)\allowbreak{})\allowbreak{}\par
\textbf{信号}：排列没有固定点\par
\par\textbf{隔板法的非负与正数解}\quad
k\ensuremath{\ge}1 时 k 个非负整数和为 n 的解数为 C(\allowbreak{}n+\allowbreak{}k-\allowbreak{}1,\allowbreak{}k-\allowbreak{}1)\allowbreak{}；k 个正整数和为 n 的有序解数为 C(\allowbreak{}n-\allowbreak{}1,\allowbreak{}k-\allowbreak{}1)\allowbreak{}，要求 n\ensuremath{\ge}k\ensuremath{\ge}1。\par
\textbf{信号}：相同物品放不同盒子、允许空盒\par
\par\textbf{Vandermonde}\quad
\ensuremath{\Sigma}\_k C(\allowbreak{}a,\allowbreak{}k)\allowbreak{}C(\allowbreak{}b,\allowbreak{}r-\allowbreak{}k)\allowbreak{}=\allowbreak{}C(\allowbreak{}a+\allowbreak{}b,\allowbreak{}r)\allowbreak{}，越界组合数视为零\par
\textbf{信号}：两组选择合计 r 个\par
\par\textbf{约数和}\quad
n=\allowbreak{}\ensuremath{\prod}p\textasciicircum{}e 时正约数和为 \ensuremath{\prod}(\allowbreak{}1+\allowbreak{}p+\allowbreak{}.\allowbreak{}.\allowbreak{}.\allowbreak{}+\allowbreak{}p\textasciicircum{}e)\allowbreak{}，模合数可直接等比求和避免除法\par
\textbf{信号}：所有约数总和\par
\par\textbf{\ensuremath{\varphi} 乘积公式}\quad
正整数 n 满足 \ensuremath{\varphi}(\allowbreak{}n)\allowbreak{}=\allowbreak{}n*\ensuremath{\prod}\_(\allowbreak{}p整除n)\allowbreak{}(\allowbreak{}1-\allowbreak{}1/\allowbreak{}p)\allowbreak{}\par
\textbf{信号}：互质数量、欧拉降幂\par
\par\textbf{\ensuremath{\mu} 取值}\quad
含平方质因子时 \ensuremath{\mu}(\allowbreak{}n)\allowbreak{}=\allowbreak{}0，否则 \ensuremath{\mu}(\allowbreak{}n)\allowbreak{}=\allowbreak{}(\allowbreak{}-\allowbreak{}1)\allowbreak{}\textasciicircum{}不同质因子数，\ensuremath{\mu}(\allowbreak{}1)\allowbreak{}=\allowbreak{}1\par
\par\textbf{\ensuremath{\mu} 约数和}\quad
\ensuremath{\Sigma}\_(\allowbreak{}d整除n)\allowbreak{} \ensuremath{\mu}(\allowbreak{}d)\allowbreak{} 为 n=\allowbreak{}1 时的 1，否则为 0\par
\par\textbf{欧拉约数和}\quad
\ensuremath{\Sigma}\_(\allowbreak{}d整除n)\allowbreak{} \ensuremath{\varphi}(\allowbreak{}d)\allowbreak{}=\allowbreak{}n\par
\par\textbf{Legendre}\quad
n! 中素数 p 的指数是 \ensuremath{\Sigma}\_(\allowbreak{}k\ensuremath{\ge}1)\allowbreak{} floor(\allowbreak{}n/\allowbreak{}p\textasciicircum{}k)\allowbreak{}\par
\textbf{信号}：阶乘质因数、尾零数量\par
\par\textbf{随机哈希概率}\quad
素数域中对不同且非恒等的长度\ensuremath{\le}L 哈希多项式随机选底数，单次碰撞概率\ensuremath{\le}(\allowbreak{}L-\allowbreak{}1)\allowbreak{}/\allowbreak{}可选底数数；Q 次比较再用并合界\par
\textbf{信号}：子串比较、需估计错误概率；固定底数没有这一保证\par
\par\textbf{独立双哈希}\quad
两个独立抽样哈希同时碰撞的概率上界可相乘，固定的相关参数不能直接如此估算\par
\textbf{信号}：一次碰撞不能接受、比较次数巨大\par
\par\textbf{线性基秩}\quad
GF(\allowbreak{}2)\allowbreak{} 线性基秩为 r 时能表示 2\textasciicircum{}r 个不同异或值，n 个输入的每个可达值有 2\textasciicircum{}(\allowbreak{}n-\allowbreak{}r)\allowbreak{} 个子集方案\par
\textbf{信号}：子集异或计数\par
\par\textbf{条件期望}\quad
E[\allowbreak{}X]\allowbreak{}=\allowbreak{}\ensuremath{\Sigma}P(\allowbreak{}A\_i)\allowbreak{}*E[\allowbreak{}X given A\_i]\allowbreak{}，其中 A\_i 是互斥穷尽事件\par
\textbf{信号}：先分第一步、随机状态转移\par
\par\textbf{几何等待}\quad
每次独立成功概率 p>\allowbreak{}0 时，直到首次成功的试验次数期望为 1/\allowbreak{}p\par
\textbf{信号}：重复独立试验、第一次成功\par
\topic{动态规划}
\par\textbf{中位数 L1}\quad
\ensuremath{\Sigma} abs(\allowbreak{}x-\allowbreak{}a\_i)\allowbreak{} 在中位数处最小；偶数个时两个中位数之间都最优\par
\textbf{信号}：数轴聚点、总绝对距离最小\par
\par\textbf{加权中位数}\quad
非负权下加权绝对距离和在左右权重都不超过总权一半的位置取最小\par
\textbf{信号}：带人数的数轴选址\par
\par\textbf{平方距离均值}\quad
实数域下 \ensuremath{\Sigma}(\allowbreak{}x-\allowbreak{}a\_i)\allowbreak{}\textasciicircum{}2 在算术平均数处最小，整数域只需查平均数两邻整数\par
\textbf{信号}：平方损失、整数选点\par
\par\textbf{区间调度}\quad
最多不重叠区间按右端点从小到大选取，端点相接合法性须按题意定义\par
\par\textbf{区间调度带权}\quad
带权区间调度不能直接贪心取最早结束，排序后 DP 对接前一兼容区间\par
\textbf{信号}：每活动价值不同\par
\par\textbf{多重背包的同余类单调队列}\quad
多重背包固定容量模重量的余数后，数量限制成为滑动窗口，可用单调队列\par
\textbf{信号}：某物品数量巨大、需要 O(\allowbreak{}nV)\allowbreak{} 转移\par
\par\textbf{二进制拆多重物品}\quad
数量 c 拆为若干 1、2、4 和剩余块，使所有 0.\allowbreak{}.\allowbreak{}c 数量可由块选取表示\par
\textbf{信号}：多重背包数量上限\par
\par\textbf{分组背包顺序}\quad
每组至多一项须从上一组状态统一转移，不能让同组新值再选另一项\par
\textbf{信号}：套餐组选一、主附件组合\par
\par\textbf{LCS 与排列 LIS}\quad
两个无重复序列的 LCS 可把一列映射到另一列位置，再求映射序列 LIS\par
\textbf{信号}：两个排列公共子序列\par
\par\textbf{二维偏序防同值}\quad
要求两个维度都严格递增时，第一维升序且同值第二维降序，再求第二维严格 LIS\par
\par\textbf{Dilworth 一般形式}\quad
有限偏序最少链划分数等于最大反链大小；序列 LIS 与不升划分是其对偶特例\par
\par\textbf{斜率优化展开}\quad
f[\allowbreak{}i]\allowbreak{}=\allowbreak{}min\_j(\allowbreak{}f[\allowbreak{}j]\allowbreak{}+\allowbreak{}(\allowbreak{}S\_i-\allowbreak{}S\_j-\allowbreak{}L)\allowbreak{}\textasciicircum{}2)\allowbreak{} 可拆为 S\_i\ensuremath{^2}-\allowbreak{}2LS\_i+\allowbreak{}L\ensuremath{^2} 加直线查询 (\allowbreak{}-\allowbreak{}2S\_j)\allowbreak{}*S\_i+\allowbreak{}f[\allowbreak{}j]\allowbreak{}+\allowbreak{}S\_j\ensuremath{^2}+\allowbreak{}2LS\_j\par
\textbf{信号}：平方前缀和代价\par
\par\textbf{单调凸壳前提}\quad
候选直线斜率按序插入且查询横坐标单调，才可用相应方向的队尾维护和队首弹出\par
\textbf{信号}：前缀和严格单调；重复斜率保留更优截距\par
\par\textbf{凸壳交点判据}\quad
斜率有序时若中间直线两侧交点次序破坏下包络凸性，则删掉中线；比较用交叉乘法并用宽类型\par
\textbf{信号}：避免浮点交点误差、斜率优化\par
\par\textbf{非单调直线查询}\quad
查询横坐标无序时不能永久弹队首，可凸壳上二分；插入斜率也无序时考虑李超树\par
\textbf{信号}：插入斜率无序的一般线性 DP 转移\par
\par\textbf{数位 DP 上界}\quad
统计 [\allowbreak{}l,\allowbreak{}r]\allowbreak{} 用 F(\allowbreak{}r)\allowbreak{}-\allowbreak{}F(\allowbreak{}l-\allowbreak{}1)\allowbreak{}，tight 状态限制本位上界，前导零是否成数字依题意\par
\textbf{信号}：大区间数字性质、不能逐个枚举\par
\par\textbf{背包复杂度伪多项式}\quad
O(\allowbreak{}nV)\allowbreak{} 依赖容量的数值而非输入位数，V 很大时转向价值维、折半或稀疏状态\par
\textbf{信号}：重量达 1e9、n 或总价值小\par
\par\textbf{树背包合并}\quad
合并孩子时容量拆成父侧与孩子侧，两层容量循环必须限制在实际子树大小\par
\textbf{信号}：子树选 k 点、连通选择\par
\par\textbf{概率转移归一}\quad
转移概率从某状态出去的总和应为 1，吸收状态单独处理\par
\textbf{信号}：概率 DP、期望状态方程\par
\par\textbf{变长任务反悔}\quad
每任务耗时可变且只求最多完成件数时，按截止时间排序，超时撤销已选耗时最大的任务\par
\topic{序列与字符串}
\par\textbf{border 链}\quad
一个前缀所有非空真 border 的长度可沿 pi 或 nxt 链逐级跳出\par
\par\textbf{前缀函数增幅}\quad
pi[\allowbreak{}i]\allowbreak{}\ensuremath{\le}pi[\allowbreak{}i-\allowbreak{}1]\allowbreak{}+\allowbreak{}1，失配时沿前一 border 跳转而不回退文本指针\par
\par\textbf{重叠匹配}\quad
匹配完整模式后回到其最长真 border，可继续找到重叠出现\par
\par\textbf{Z 函数匹配}\quad
在 pattern+\allowbreak{}分隔符+\allowbreak{}text 上计算 Z，text 段 Z\ensuremath{\ge}模式长处就是匹配起点\par
\textbf{信号}：求所有匹配与前缀 LCP\par
\par\textbf{自动机禁状态}\quad
禁止模式 DP 的坏状态需沿 fail 传播，一旦到坏状态就不再转移或计数\par
\textbf{信号}：避免多个禁串、按长度计数\par
\par\textbf{自动机矩阵}\quad
有限状态计数转移不随位置改变时可把转移次数写为矩阵，长度 n 用矩阵幂\par
\textbf{信号}：状态数小、n 很大\par
\topic{复杂度与套路}
\par\textbf{根号分治}\quad
把参数按 B 分界，小参数预处理、大参数按总量/\allowbreak{}参数限制枚举，常取 B 约 sqrt(\allowbreak{}总量)\allowbreak{}\par
\textbf{信号}：两种代价一增一减\par
\par\textbf{根号值域分治}\quad
频次>\allowbreak{}B 的不同值最多 n/\allowbreak{}B 种，低频值可按出现列表枚举\par
\textbf{信号}：按颜色、权值、出现次数查询\par
\par\textbf{入出栈均摊}\quad
每项最多入栈一次出栈一次，即使单次循环很多，总弹出次数仍为 O(\allowbreak{}n)\allowbreak{}\par
\par\textbf{小并大均摊}\quad
每次把小集合元素搬进至少同样大的集合，每个元素所在集合大小至少翻倍，总搬运 O(\allowbreak{}n log n)\allowbreak{}\par
\textbf{信号}：子树集合合并\par
\par\textbf{势能法}\quad
摊还代价为实际代价加势能变化；势能始终非负且初值可控时，摊还和界住总实际代价\par
\textbf{信号}：一次很慢但总量可控、段删除\par
\par\textbf{几何增长数组}\quad
容量按固定比例扩张时复制总次数为线性量，逐次只加一容量会二次退化\par
\textbf{信号}：动态数组、自制缓冲区\par
\par\textbf{莫队移动}\quad
静态区间询问按左端块、右端顺序离线，使总指针移动约 O(\allowbreak{}n\ensuremath{^2}/\allowbreak{}B+\allowbreak{}qB)\allowbreak{}，不是单次 O(\allowbreak{}1)\allowbreak{}\par
\textbf{信号}：可增删端点的区间统计、所有询问已知\par
\par\textbf{修改莫队}\quad
加入时间维可处理离线点修改，回退修改须保存旧值并同步当前区间贡献\par
\textbf{信号}：带点修改的区间统计\par
\par\textbf{扫描线事件先后}\quad
二维条件 x\ensuremath{\le}查询值时同 x 的插入须先于查询，x<\allowbreak{} 时顺序反过来\par
\par\textbf{整体二分}\quad
答案落有序值域且贡献可按中点累计时，可按判定把询问分到左右并回滚临时修改\par
\textbf{信号}：离线第 k 小、值域二分\par
\par\textbf{平行二分}\quad
多个询问各维护答案区间，把相同中点的检查按一次扫描批量处理\par
\textbf{信号}：添加操作后首次达阈值\par
\par\textbf{分治时间轴}\quad
边或约束的生效时间为区间时分配到时间线段树，DFS 配回滚结构\par
\textbf{信号}：动态连通、加入删除边离线\par
\par\textbf{回滚并查集}\quad
回滚通常按大小合并而不路径压缩，记录根父亲与大小改动后按栈撤销\par
\textbf{信号}：时间分治、撤销连边\par
\par\textbf{调和复杂度}\quad
\ensuremath{\Sigma}\_(\allowbreak{}i=\allowbreak{}1.\allowbreak{}.\allowbreak{}n)\allowbreak{} floor(\allowbreak{}n/\allowbreak{}i)\allowbreak{}=\allowbreak{}O(\allowbreak{}n log n)\allowbreak{}，多重枚举倍数总量不能按 n\ensuremath{^2} 估计\par
\textbf{信号}：筛法、枚举因子倍数\par
\par\textbf{树链剖分段数}\quad
沿路径每跨一条轻边子树大小至少翻倍，任意路径只跨 O(\allowbreak{}log n)\allowbreak{} 条重链\par
\textbf{信号}：路径转区间、链上数据结构\par
\par\textbf{虚树结点数}\quad
k 个不同关键点与相邻 DFS 序点的 LCA 闭包构成虚树，结点数不超过 2k-\allowbreak{}1\par
\textbf{信号}：每次只问少量关键点\par

\topic{WQS：数量限制变罚项}
令 $F(k)$ 为恰选 $k$ 次的最小代价，
$G(\lambda)=\min_k\{F(k)+\lambda k\}$。
若 $F$ 离散凸，且 $\lambda$ 支持目标 $K$（或 $K$ 位于并列最优数量之间），则
\[
F(K)=G(\lambda)-\lambda K.
\]
罚项增大时最优次数不增；并列时统一取更多次数。
仅次数单调不够：$F(0)=F(2)=0,F(1)=10$ 时只能恢复下凸包。
若单次 DP 为 $O(T)$，二分罚项为 $O(T\log W)$；须证明范围与凸性。
\topic{Burnside / P\'olya：对称下计数}
有限群 $G$ 作用于合法方案集合：
\[
\#\text{轨道}=\frac1{|G|}\sum_{g\in G}\operatorname{Fix}(g).
\]
自由用 $c$ 色染位置时，$\operatorname{Fix}(g)=c^{\#\text{轮换}}$；
有颜色数量限制时须重新数不动方案。$n\ge1$，令
\[
R=\sum_{i=0}^{n-1}c^{\gcd(n,i)}
=\sum_{d\mid n}\varphi(d)c^{n/d}.
\]
只认旋转的项链数为 $R/n$；允许翻转的手链数为 $(R+H)/(2n)$，其中
\[
H=\begin{cases}
nc^{(n+1)/2},&n\text{ 奇},\\
\frac n2\bigl(c^{n/2+1}+c^{n/2}\bigr),&n\text{ 偶}.
\end{cases}
\]
模 $M$ 除以 $|G|$：互素时用逆元；否则可先模 $M|G|$ 计算分子，再整数除以 $|G|$。
例如 $n=6,c=2$：项链 14，手链 13。

\topic{扩展欧拉、LGV 与建模补充}
扩展欧拉降幂：m>\allowbreak{}1，\ensuremath{\varphi}=\allowbreak{}\ensuremath{\varphi}(\allowbreak{}m)\allowbreak{}。gcd(\allowbreak{}a,\allowbreak{}m)\allowbreak{}=\allowbreak{}1 时 a\textasciicircum{}b \ensuremath{\equiv} a\textasciicircum{}(\allowbreak{}b mod \ensuremath{\varphi})\allowbreak{}；不互素时只有 b>\allowbreak{}=\allowbreak{}\ensuremath{\varphi} 才能用 a\textasciicircum{}(\allowbreak{}b mod \ensuremath{\varphi}+\allowbreak{}\ensuremath{\varphi})\allowbreak{}。b<\allowbreak{}\ensuremath{\varphi} 必须保留原指数，不能无条件加 \ensuremath{\varphi}。超长十进制指数同时维护余数和“是否达到 \ensuremath{\varphi}”的饱和值标记；m=\allowbreak{}1 结果为 0。\par\smallskip
LGV：DAG 中 M[\allowbreak{}i]\allowbreak{}[\allowbreak{}j]\allowbreak{} 为起点 s\_i 到终点 t\_j 的路径数，则 det(\allowbreak{}M)\allowbreak{} 是按终点排列符号加权的、两两点不相交路径族数量。只有起终点满足兼容顺序、交叉排列必相交时，才能把行列式当作固定配对的非负计数。边权版本用路径权积；求路径数用拓扑 DP，行列式复用现有模板。\par\smallskip
Prüfer：新增编码/\allowbreak{}解码实现。n 个标号顶点的树为 n\textasciicircum{}(\allowbreak{}n-\allowbreak{}2)\allowbreak{}（n>\allowbreak{}=\allowbreak{}2）；度数 d\_i>\allowbreak{}=\allowbreak{}1 且和为 2n-\allowbreak{}2 时为 (\allowbreak{}n-\allowbreak{}2)\allowbreak{}!/\allowbreak{}\ensuremath{\prod}(\allowbreak{}d\_i-\allowbreak{}1)\allowbreak{}!。模除法仍要求分母可逆，合数模不能直接费马求逆。\par\smallskip
同余最短路：正硬币无限使用，取最小硬币 m，dist[\allowbreak{}r]\allowbreak{} 是可表示的最小 r 类数；x 可表示当且仅当 x>\allowbreak{}=\allowbreak{}dist[\allowbreak{}x\%m]\allowbreak{}。gcd>\allowbreak{}1 时无有限的最大不可表示数；gcd=\allowbreak{}1 时 max(\allowbreak{}dist)\allowbreak{}-\allowbreak{}m 是最大不可表示数。新增模板按隐式边做 Dijkstra，不另建 m*k 条边。\par\smallskip
\section{赛事建模与技巧}
\note{从已核对题面与题解中提炼可复用部分。复用已有模板，以下条件与公式不代表整道原题的全部解法；不重复收录基础操作。}
\topic{可合并信息与贡献}
\par\textbf{点集只留直径端点}\quad
\note{非负边权树中，非空点集 A 的直径端点 a,\allowbreak{}b 满足 $\max_{x\in A}d(v,x)=\max(d(v,a),d(v,b))$。合并两个点集，只比较原来两条直径和四组跨集合端点距离；空集为单位元。并列时其他点也可能最远，不能说最远点只能是端点。复用 LCA 距离（2021 澳门 J、2022 沈阳 G）。}
\par\textbf{整体加下的 gcd 摘要}\quad
\note{非空整数集合存任一元素 f 及 $g=\gcd\{|x-f|\}$。整体加只改 f；合并两集合得 $(f_1,\gcd(g_1,g_2,|f_2-f_1|))$，实际元素 gcd 为 $\gcd(|f|,g)$。单点 g=\allowbreak{}0，空集另标记；差与绝对值须防溢出（2022 杭州 M）。}
\par\textbf{截零累加的可结合变换}\quad
\note{$s\gets\max(s+t_i,0)$，s 非负。一段存总和 S 与含空后缀的最大后缀和 M，作用为 $T(s)=\max(s+S,M)$。先 A 后 B：$S=S_A+S_B$，$M=\max(M_A+S_B,M_B)$；空段 (\allowbreak{}0,\allowbreak{}0)\allowbreak{}。重复 k 次（k\ensuremath{\ge}1）：$T^k(s)=\max(s+kS,M+\max(0,(k-1)S))$。可作线段树摘要（2022 济南 B）。}
\par\textbf{段生命周期与离线矩形}\quad
\note{只有端点 split 和区间赋值创建段时，总创建数 O(\allowbreak{}n+\allowbreak{}q)\allowbreak{}，赋值删除段的遍历可按生命周期收费。每段对应下标区间\ensuremath{\times}活跃时间半开区间，可转离线二维扫描。add、查询反复遍历却不删段，不适用此均摊，最坏仍可 O(\allowbreak{}nq)\allowbreak{}（2022 广州 B）。}
\par\textbf{置换分环转循环相关}\quad
\note{环长 L，$R_k=\sum_{i=0}^{L-1}a_i a_{(i+k)\bmod L}$。令 $b_i=a_{(-i)\bmod L}$，普通卷积 c=\allowbreak{}b*a，则 $R_k=c_k+c_{k+L}$，越界项记0。同长度环先累加；所有环长度和为 n，不同长度仅 O(\allowbreak{}√n)\allowbreak{}，查询枚举长度取 k mod L。精确整数结果需控制 FFT 误差或用足够多模数（2021 澳门 E）。}
\par\textbf{多项式权重换二项式基}\quad
\note{维护 $\sum_{\text{方案}}\binom tj$，一次增加段数由 $\binom{t+1}j=\binom tj+\binom t{j-1}$ 转移。$t^2=2\binom t2+\binom t1$；一般 $t^d=\sum_j S(d,j)j!\binom tj$，无须模除法（2023 CCPC 预选 A）。}
\par\textbf{绝对差拆成阈值}\quad
\note{整数 x,\allowbreak{}y 有 $|x-y|=\sum_{t\in\mathbb Z}|[x>t]-[y>t]|$。按阈值统计01贡献，值域大时按相邻不同值间距压缩。连通无向图上 $\sum_{uv\in E}|a_u-a_v|\ge\max a-\min a$；等号要求每个非平凡阈值割恰一条边，故边双连通块内同值，重边按条计（2021 澳门 H、2023 EC I）。}
\par\textbf{积水量化为前后缀最大值}\quad
\note{$P_i=\max_{j\le i}a_j$、$S_i=\max_{j\ge i}a_j$、$M=\max a$，则积水为 $\sum P_i+\sum S_i-nM-\sum a_i$。仅单点增高时，被更高纪录覆盖的旧纪录不再复活，可用 set 删段均摊；允许降低时不成立（2023 南京 M）。}
\topic{贪心、逆推与构造}
\par\textbf{全部资源任务的安全次序}\quad
\note{每任务先消耗 a\ensuremath{\ge}0，再返还 b\ensuremath{\ge}0，执行前资源至少 a。先做 a\ensuremath{\le}b 的任务，按 a 升序；再做 a>\allowbreak{}b 的任务，按 b 降序。此顺序最小初始资源为 $\max(0,\max_i(a_i+\sum_{j<i}(a_j-b_j)))$。不直接推广到任选任务或收益最大化（2023 CCPC 预选 I）。}
\par\textbf{免费名额与背包边界}\quad
\note{每件至多一次，免费 k 件，其他所选重量和\ensuremath{\le}W。若免费件比付费件轻，交换身份不增加付费重量，因此存在免费件均不轻于付费件的最优解。重量降序后枚举分界：前缀取价值最大的 k 件，后缀做01背包。免费件不必是全局最重 k 件（2023 南京 G）。}
\par\textbf{除当前物品外的背包}\quad
\note{分治候选下标 [\allowbreak{}l,\allowbreak{}r]\allowbreak{}，入口 DP 已加入区间外所有物品。下传左边时加入右半，反之亦然；到叶 i 就得到排除 i 的 DP。容量 V 的01背包总 O(\allowbreak{}nV log n)\allowbreak{}，按 DFS 复制数组需要 O(\allowbreak{}V log n)\allowbreak{} 额外空间；保存全部分支会更多。不能对非单调特殊收益再套贪心（2022 杭州 C）。}
\par\textbf{逆推 max 约束与位置计数}\quad
\note{从 $f_i=\max(f_{i-1}+c_i,g_i)$ 的上界反推，可消除正向等待决策；每个分支的上界都必须满足。若最终得到 n 个不同对象必须放在不同位置、且第 i 个允许的前缀位置数为 b\_i，排序 b 后方案数为 $\prod_{i=1}^n\max(0,b_i-i+1)$。位置同质、限制确为前缀是前提；概率由计数除 n! 得到，尾和从 k=\allowbreak{}0 起（2025 成都 C、沈阳 E；2026 预选 II D）。}
\par\textbf{先规范最优形态，再枚举}\quad
\note{先通过交换证明操作可收拢为层、后缀、主路径加完整挂树等形态，再对少量形态做 DP 或贪心。分别排序两组量前要证明配对关系无关；不能凭样例推断。先手工检查零代价、空方案和端点，避免把题目特有结论当通用模板（2025 沈阳 B、武汉 I、杭州 G）。}
\par\textbf{特殊角色的线性不变量}\quad
\note{活跃位置 i 反射到 $P'_i=2P_j-P_i$，随后活跃角色交给 j，其他位置不变，则 $\sum P-2P_i$ 保持。终态可用 $\sum P-2P_s=\sum Q-2Q_t$ 筛选 t；这是必要条件，充分性还需对应规则证明，不可任意推广（2025 沈阳 K）。}
\topic{图与树}
\par\textbf{最小瓶颈路径}\quad
\note{非负边权，最小化路径最大边：Dijkstra 松弛改为 $nd=\max(dis_u,w)$，源距离0，仍用小根堆与过期项检查。无向图也可用最小生成树路径；最大生成树对应的方向不同，不能混用（2022 威海 F、2023 合肥 J）。}
\par\textbf{最大匹配与可增益边}\quad
\note{已求最大二分匹配，在残量网络中，左点 u 从源可达、右点 v 能达汇，当且仅当新增边 (\allowbreak{}u,\allowbreak{}v)\allowbreak{} 可增大匹配。因此增益端点对数是两集合大小之积；已有边不同时满足此条件，重边不产生新端点对（2023 济南 E）。}
\note{无向无重复访问的行走博弈：双方轮流把同一棋子沿边移到未访问点，不能走者输。起点 s 必胜，当且仅当每个最大匹配都覆盖 s，等价于删除 s 后匹配数下降1。取匹配伙伴作为应对；若某最大匹配不覆盖 s，对手可按配对回应。一般无向图需一般图匹配，不仅限二分图（2025 武汉 D）。}
\par\textbf{连通点集与路径：点减边}\quad
\note{树上非空连通点集 C 与路径 P 的交仍连通，所以 $[C\cap P\ne\varnothing]=|C\cap V(P)|-|E(C)\cap E(P)|$。对各连通点集求和，查询变成路径点贡献减边贡献。消息 $f_{u\to v}=\prod_{w\in N(u),w\ne v}(1+f_{w\to u})$；点贡献为全邻居积，边 uv 贡献为 $f_{u\to v}f_{v\to u}$。模意义用前后缀积，不因因子零而求逆（2024 成都 E）。}
\par\textbf{动态根路径并}\quad
\note{关键点按 DFS 序排，加入 x 时与旧集合的最深 LCA 可在前驱、后继取得。根永久作为辅助关键点，新增覆盖为 (\allowbreak{}L,\allowbreak{}x]\allowbreak{}；删点在去掉 x 后找邻居反向更新。根路径并边权总和为 $\frac12\sum_{\text{循环相邻 }u,v}d(u,v)$，集合计重及根点是否计数须单独处理（Metropolis J）。}
\par\textbf{SCC 中检查非零和环}\quad
\note{一个 SCC 内，所有有向环权和为0，当且仅当存在整数势能 h，全部内部边满足 $h_v-h_u=w_{uv}$。沿遍历树赋值，再检查所有内部边，冲突即存在非零和环；不是只判负环。缩点后可反向传播能到达冲突块的起点。势能范围须足够，负余数规范到 [\allowbreak{}0,\allowbreak{}n)\allowbreak{}（2024 沈阳 M，原题按 n 取模）。}
\par\textbf{路径公共边的二阶矩}\quad
\note{n\ensuremath{\ge}2，两条独立均匀无向非空路径，$K=n(n-1)/2$。边 e 两侧大小 s,\allowbreak{}n-\allowbreak{}s，$\lambda_e=s(n-s)$；不同边 e,\allowbreak{}f 删去后，两端分量大小 a,\allowbreak{}c，$\lambda_{ef}=ac$。公共边数 X 满足 $\mathbb E[X^2]=(\sum_e\lambda_e^2+2\sum_{e<f}\lambda_{ef}^2)/K^2$。不同边项不含对角；计算可按 LCA 分类，不直接 O(\allowbreak{}n\ensuremath{^2})\allowbreak{} 枚举大树（2024 哈尔滨 L）。}
\par\textbf{小 k 匹配核化}\quad
\note{简单二分图，任意边权，求所有恰好 t\ensuremath{\le}k 条的最优匹配。先每个左点只留最重 k 边，再在所得图每个右点留最重 k 边；最后留全局最重 $(2k-1)(k-1)+1$ 边，保持这些最优值。顺序不能换成双侧各自独立筛边，同端点重边先取最大。费用流按每次增广分别记录，不能只取最大流结果（Metropolis K）。}
\topic{字符串、矩阵与位运算}
\par\textbf{自定义字母序的逆序计数}\quad
\note{固定字符串序列，预统计 C[\allowbreak{}a]\allowbreak{}[\allowbreak{}b]\allowbreak{}：i<\allowbreak{}j 的两串首次不同字符分别为 a,\allowbreak{}b；另计后串为前串真前缀的对数 P。字符顺序 rank 下答案 $P+\sum_{rank(a)>rank(b)}C[a][b]$。相同串无贡献，前缀项不随字母序变化。Trie 顺序插入可 O(\allowbreak{}\ensuremath{\sigma}·总字长)\allowbreak{} 统计，每问 O(\allowbreak{}\ensuremath{\sigma}\ensuremath{^2})\allowbreak{}，计数用64位（2022 杭州 K）。}
\par\textbf{自动机长字符段与递归串}\quad
\note{固定自动机，消费一个字符 c 后统计到达状态权：$G_0(u,c)=go(u,c)$、$H_0(u,c)=val(G_0(u,c))$；$G_{j+1}=G_j(G_j(u,c),c)$，$H_{j+1}=H_j(u,c)+H_j(G_j(u,c),c)$。长同字符段按位倍增，O(\allowbreak{}log L)\allowbreak{}（2021 哈尔滨 L）。}
\note{子序列匹配目标 T，行向量状态 j 表示已匹配前 j 字符。字符矩阵对角为1，若 T[\allowbreak{}j]\allowbreak{}=\allowbreak{}c，则 j→j+\allowbreak{}1 加1；单字符不能跳过多个位置。递归串 $U_i=U_{i-1}+c_i+U_{i-1}$ 的矩阵为 $G_0=I$，$G_i=G_{i-1}F_{c_i}G_{i-1}$。答案 $G_n[0][|T|]$，空目标恰1种（2024 CCPC 网络 D）。}
\par\textbf{固定掩码操作的仿射复合}\quad
\note{在 W 位 GF(\allowbreak{}2)\allowbreak{} 向量上，AND 是对角矩阵，OR m 是 $diag(\sim m)x+m$，XOR 是 $x+c$，循环移位是置换矩阵；补码只限 W 位。若 f=\allowbreak{}Ax+\allowbreak{}b、g=\allowbreak{}Bx+\allowbreak{}c，则 $g\circ f=(BA,Bb+c)$。列压整数后作用 O(\allowbreak{}W)\allowbreak{}，复合 O(\allowbreak{}W\ensuremath{^2})\allowbreak{}，非交换顺序要和线段树合并方向一致（2023 深圳 H）。}
\note{$0\le a\le x,0\le b\le y$ 时最大 a OR b：若 x AND y=\allowbreak{}0，为 x OR y；否则 h 为 x AND y 的最高置位值，答案 $(x\mathbin{\mathrm{OR}}y)\mathbin{\mathrm{OR}}(h-1)$。最高公共位以下可补满，h 不能取 OR 的最高位（2023 深圳 E）。}
\par\textbf{空间交、包含与共同 xor}\quad
\note{固定 W 维 GF(\allowbreak{}2)\allowbreak{} 标准点积，$(A\cap B)^\perp=A^\perp+B^\perp$，$(A^\perp)^\perp=A$，$\dim A^\perp=W-\dim A$。正交补用模2消元求零空间，仿射集合不能直接套用（2022 沈阳 M）。}
\note{两个线性子空间的并还是子空间，当且仅当一者包含另一者。W 维空间中，超过半数的向量必张满，因为真子空间最多一半；等于一半不够（2026 预选 II A）。}
\note{同时给 A,\allowbreak{}B xor 一个可表示 x，A xor B 不变，不能各自独立最小化。若 A\ensuremath{\ne}B，设 h 为差的最高位：先最小化 h 以上共同前缀，再分 h 位主元用或不用，确定较大者，低位基最小化它，取两支较优；A=\allowbreak{}B 直接求最小 xor（2024 CCPC 网络 J）。}
\topic{数学与概率}
\par\textbf{路径计数诱导均匀采样}\quad
\note{DAG 中 cnt[\allowbreak{}u]\allowbreak{} 为到终点的路径数，沿 u→v 以 cnt[\allowbreak{}v]\allowbreak{}/\allowbreak{}cnt[\allowbreak{}u]\allowbreak{} 选分支，才使完整路径均匀；每条出边均分通常只使下一步均匀。重边按独立边是否计为不同路线须统一；不可达点 cnt=\allowbreak{}0 不可除（2026 CCPC 预选 K）。}
\par\textbf{零频率叶与编码长度}\quad
\note{Huffman 模型若要求给每个名册对象编码，零频率对象仍是叶，不能因不影响直接代价就删除，因为会改变树结构。单符号若规定编码非空，须至少一位，经典根深度0的约定不适用。相同频率可批量配对，但奇数余项仍要与全局下一最小项合并（2026 CCPC 预选 L）。}
\par\textbf{等长排列前缀集合不用哈希}\quad
\note{两排列中，将第一排列前 k 项映射到第二排列的位置。因为恰有 k 个不同位置，其最大值等于 k 当且仅当两者前 k 项集合相同。真排列与等长是前提，存在重复值时不成立；可扫描前缀最大值 O(\allowbreak{}n)\allowbreak{} 判全部边界（2026 CCPC 预选 E）。}
\par\textbf{可撤销 Bernoulli 因子}\quad
\note{独立 Bernoulli 和的生成多项式 $F(z)=\prod_i((1-p_i)+p_i z)$。撤销一个因子后 $F_k=(1-p)G_k+pG_{k-1}$；1-\allowbreak{}p 非零时从低次递推，p=\allowbreak{}1 时 $G_k=F_{k+1}$，p=\allowbreak{}0 不变。模意义要求实际除数可逆，不等于实数非零就可用逆元；更新一项再乘新因子 O(\allowbreak{}次数)\allowbreak{}（2024 哈尔滨 E）。}
\par\textbf{布尔卷积的精确性}\quad
\note{0/\allowbreak{}1 数组长度 L\ensuremath{_1},\allowbreak{}L\ensuremath{_2}，普通卷积每系数\ensuremath{\le}min(\allowbreak{}L\ensuremath{_1},\allowbreak{}L\ensuremath{_2})\allowbreak{}。NTT 模数严格大于该界时，结果非零恰为可达，之后恢复0/\allowbreak{}1。连续计数卷积不二值化时不能保证无模后假零；截断指数范围必须另证可达路径不会被截掉（2024 昆明 I）。}
\par\textbf{小质因子掩码，大质因子分组}\quad
\note{数值\ensuremath{\le}V 时，每个数最多含一个大于√V的质因子。对小质因子的出现做掩码，按唯一大质因子分组。对子集乘积求 \ensuremath{\varphi}，先乘选中的数，再对出现的不同质因子各乘一次 (\allowbreak{}p-\allowbreak{}1)\allowbreak{}/\allowbreak{}p，空集贡献 \ensuremath{\varphi}(\allowbreak{}1)\allowbreak{}=\allowbreak{}1；模数须让相关 p 可逆。小质因子数量也必须能承受指数状态（2024 重庆 A）。}
\par\textbf{整数区间选不同数的和}\quad
\note{l\ensuremath{\le}r，$1\le k\le r-l+1$，选 k 个不同整数的可能和恰为无空洞区间 $[k(2l+k-1)/2,\ k(2r-k+1)/2]$，因此只需上下界判断。k=\allowbreak{}0 仅有和0，乘法先扩宽（2022 威海 K）。}
\par\textbf{下中位数的符号和}\quad
\note{排列区间含唯一值 x，a\ensuremath{\ge}x 记+\allowbreak{}1，其余-\allowbreak{}1。x 为排序第 ceil(\allowbreak{}L/\allowbreak{}2)\allowbreak{} 个，当且仅当奇长度和为1、偶长度和为2。上中位数的偶长条件不同；存在相等重复值时不能用此唯一 x 判据（2023 EC D）。}
\par\textbf{参数消耗证明根号种类数}\quad
\note{若每种不同正整数参数值 d 都至少消耗 d 个互不重叠单位，总消耗\ensuremath{\le}N，则不同值 k 满足 $1+\cdots+k\le N$，故 k=\allowbreak{}O(\allowbreak{}√N)\allowbreak{}。只能在消耗不重复收费且参数为不同正整数时使用，不能把“值很快变”当证明（2026 预选 II F）。}
\par\textbf{单位群中阶的 lcm 分解}\quad
\note{对模 n 的单位群，CRT 拆质数幂；奇素数幂单位群循环，$2^e$ 在 e\ensuremath{\ge}3 时为 $C_2\times C_{2^{e-2}}$。分解为循环因子后，对每个质数 q，若其在各因子阶中的指数为 e\_i，则最大阶指数\ensuremath{\le}t 的选择数为 $F_q(t)=q^{\sum_i\min(t,e_i)}$。阶和该质因子贡献 $1+\sum_{t\ge1}q^t(F_q(t)-F_q(t-1))$，最后乘各质因子贡献；只对单位成立，n=\allowbreak{}1 单独约定（2026 预选 II E）。}
\topic{搜索与几何}
\par\textbf{奇偶层 BFS 的等待条件}\quad
\note{无向无孤点图且必须每步走边，往返可延迟任意偶数步。对 (\allowbreak{}点,\allowbreak{}时间奇偶)\allowbreak{} 多源 BFS 得最早危险时间 \ensuremath{\tau}；若危险只活 d 步，仅 \ensuremath{\tau}\ensuremath{\le}d 生效。玩家时刻 t 到达 v 须 $t<\tau[v][t\bmod2]$，同时到达失败。是否允许对向过边、停留、孤点均会改变模型，不套到任意有向图（2024 预选 II E）。}
\par\textbf{最短路 DAG 与自适应折半}\quad
\note{从源、汇分别 BFS，仅保留 $d_s(u)+1+d_t(v)=d_s(t)$ 的边，所有最少边路径落在 DAG。非负费用可剪枝；以路径分支量而非固定层数选折半界，跨边配对时起终点费用各算一次。n=\allowbreak{}1 单点路线单独处理，路径计数需精确上界，不能拿溢出后的数分类（2026 预选 I G）。}
\par\textbf{二次曲线包络消共同项}\quad
\note{圆心在 x 轴的上半圆：$y^2=R_i^2-(x-c_i)^2$。令 $z=x^2+y^2$，比较变为直线 $z=2c_i x+R_i^2-c_i^2$ 的上包络；仍需限制圆弧有效 x 区间及 y\ensuremath{\ge}0。全局直径的旋转卡壳不能直接求每个顶点的最远点（2026 预选 II B）。}
\par\textbf{整数 incircle 的符号与容量}\quad
\note{平移查询点 p 到原点，a=\allowbreak{}A-\allowbreak{}p、b=\allowbreak{}B-\allowbreak{}p、c=\allowbreak{}C-\allowbreak{}p：$D=|a|^2\operatorname{cross}(b,c)-|b|^2\operatorname{cross}(a,c)+|c|^2\operatorname{cross}(a,b)$。A,\allowbreak{}B,\allowbreak{}C 非共线时，D 与 cross(\allowbreak{}B-\allowbreak{}A,\allowbreak{}C-\allowbreak{}A)\allowbreak{} 同号表示圆内，D=\allowbreak{}0 共圆。平移坐标绝对值\ensuremath{\le}M 时保守界为 $12M^4$；减法先扩宽，界与中间积均能放128位才用该类型，不能只看原坐标是64位（Metropolis M）。}
\topic{配对补证后的推导技巧}
\par\textbf{最后一次操作给资源阈值}\quad
\note{n 堆石头，取1入袋、从袋放n到一堆，初值 a、终值 b，袋可剩石。设 $S=\sum a_i$、$r_i=b_i\bmod n$、$R=\sum r_i$、$P=\max r_i$。可达恰为 $\sum b_i\le S$ 且（所有 $r_i\le a_i$，或 $S-R\ge n-P$）。必要性看最后一次放n，此时其他堆至少为最终值；充分性先清空，再将最大余数堆留到最后建。等号合法，不能写严格小于（香港 Stonebag）。}
\par\textbf{杨图反对角线与轮廓计数}\quad
\note{格子 (\allowbreak{}r,\allowbreak{}c)\allowbreak{} 属于反对角线 $k=r+c-1$，数量 $t_k$，$t_0=0$、越界为0。正确差分 $d_k=t_k-t_{k+1}$：最初一段为-\allowbreak{}1，之后非负。同反对角线计数的杨图在合法右下子图转置下恰好互达；轮廓条带排列给 $\prod_{k\ge1,\ d_k\ge0}\binom{d_{k-1}+d_k+1}{d_k}$。加入第r行长a，只改 $d_{r-1}-=1$、$d_{r+a-1}+=1$，最多四个因子；负差分项跳过，不做负阶乘（2025 沈阳 H）。}
\par\textbf{高位扫描反推低位进位}\quad
\note{k个数每个\ensuremath{\le}m，所有数对xor之和为n。第b位有s个1，贡献 $s(k-s)2^b$。加法等式 $s(k-s)+c_b=n_b+2c_{b+1}$；为处理tight从高到低扫，已知 $c_{b+1}$ 后反推 $c_b=n_b+2c_{b+1}-s(k-s)$，保留非负且不超过 $\lfloor k^2/4\rfloor$。初始高位进位0，终止低位进位0。tight只存数量u：m位0权 $\binom{k-u}{s}$；位1若j个tight取1，权 $\binom uj\binom{k-u}{s-j}$、新u=\allowbreak{}j（2022 广州 M）。}
\par\textbf{相邻不等且总 xor 为零}\quad
\note{原串长度n\ensuremath{\ge}1，值域 $[0,2^m)$，仅线性相邻不等，令 $t=2^m-1$。答案奇n为 $t^{n-1}$，偶n为 $t^{n-1}+(-1)^{n/2}t^{n/2}$。xor字符筛选下，非零字符加权串满足 $S_{n+2}=-tS_n$、$S_1=0$、$S_2=-2^m$；零字符贡献 $2^m t^{n-1}$。两次幂即可，无须模逆；m=\allowbreak{}0、n=\allowbreak{}1按 $0^0=1$。不能把题解的配对数误当原串长度（2025 预选 II E）。}
\par\textbf{余数转移的反转卷积下标}\quad
\note{随机排列互异模数取模，只考虑\ensuremath{\le}当前值的下一有效模数，均匀选择；令 $w_i=f_i/c_i$。目标块余数j\ensuremath{\in}[\allowbreak{}l,\allowbreak{}r]\allowbreak{}，大模数a>\allowbreak{}r的来源 $i=j+qa>r$ 已知。设 $U[x]=w_{V-x}$（仅来源>\allowbreak{}r）、$K[d]=\#\{a>r:a\mid d,\ d>0\}$，普通卷积取 $(U*K)[V-j]$。不能加入0倍数；小模数仅枚举l<\allowbreak{}a\ensuremath{\le}r，块内按值递减处理。终止状态c\_j=\allowbreak{}0的概率乘n!才是排列数（2026 预选 I B）。}
\par\textbf{参数只在单点成立的动态 DP}\quad
\note{村民前缀狼数 $c_i=\alpha_i k+\beta_i$，村民间距1或2。令 $A_i=[c_i=c_{i-1}]$、$B_i=[c_i=c_{i-2}+1]$，明确状态 $(dp_i,dp_{i-1})^\mathsf T$，叶矩阵 $\left(\begin{smallmatrix}A_i&B_i\\1&0\end{smallmatrix}\right)$。区间乘积为右积乘左积；仿射等式非恒定时只在一个k成立，扫描参数需在k启用、k+\allowbreak{}1撤销。同叶多事件先合并，不相互覆盖（2026 预选 I K）。}
\par\textbf{局部合流保证唯一合并规范形}\quad
\note{相邻强连通块A、B在区间出边图 $a_i\le i\le b_i$ 中可合并恰为 $\max_{i\in A}b_i\ge\min B$ 且 $\min_{i\in B}a_i\le\max A$。不交合并可交换；重叠的两次合并都可继续得到同一三块并，且每步减少块数，故终局唯一。计数可用总集合减去至少两块的稳定链，按上一块最大点/\allowbreak{}最大b接下一块，避免枚举切法重计（2024 决赛 M）。}
\par\textbf{st 序给带权连通切分}\quad
\note{点双中指定s、t，开耳分解可给s首t尾的序，每个其他点均有较早、较晚邻点，因此任意非空前缀与后缀都连通。外接分支权随原点归属；总权1 mod3时，有权2点可单切，否则只剩0/\allowbreak{}1，至少四个权1点时取含恰两个的前缀。只判存在无需输出st序。原图不连通时恰两分量且各权2也可行（2024 决赛 K）。}
\par\textbf{饱和匹配的模约束交换界}\quad
\note{基若可表示为左对象匹配到全部右位置，两个基的饱和匹配之并分成交替路径/\allowbreak{}环；路径端点把基差元素一一配对，任意路径子集都可双向交换。无约束最优基A与同余最优且离A最近的B，若差至少q对，q个费用差的q+\allowbreak{}1个前缀和模q重复，得到非空零和交换；A最优使B换后不差且更近，矛盾。故差\ensuremath{\le}q-\allowbreak{}1。这个证明依赖强交换双射，不泛化到任意拟阵（2024 香港 M，q=\allowbreak{}4）。}
\par\textbf{单点修改的新增方案乘积}\quad
\note{区间合法性只因降低a\_x而放宽，旧划分全部保留。新增划分有唯一包含x的新合法段[\allowbreak{}l,\allowbreak{}r]\allowbreak{}，其余段仍用原前缀数f、后缀数g，所以新答案为旧答案加 $\sum f_{l-1}g_{r+1}$，仅枚举旧不合法/\allowbreak{}新合法且含x的段。区间最大值降低须x为唯一最大值；不能用f+\allowbreak{}g，不能把旧合法段再加一次（2023 济南 F）。}
\par\textbf{最大边加次大边的最短路}\quad
\note{非负权图，空路径瓶颈0；两次瓶颈Dijkstra得 $d_s,d_t$。枚举边(\allowbreak{}u,\allowbreak{}v,\allowbreak{}w)\allowbreak{}及反向，$B=\max(d_s[u],d_t[v])$，仅B\ensuremath{\le}w时取w+\allowbreak{}B。最优路最大边分开后两侧瓶颈\ensuremath{\le}次大边；候选拼接删环又能得到代价不超过w+\allowbreak{}B的简单路，两向证明最优。直接s-\allowbreak{}t边候选w+\allowbreak{}0，不能无条件把w当最大边（2023 合肥 J）。}
\par\textbf{跨两行平方串的配对比较}\quad
\note{两行长度N，向右/\allowbreak{}下且起终任意，路径最多N+\allowbreak{}1字符。补不同哨兵写成 $S[i,j)T[j,k)$；半长d的配对应比较 $S[x]=T[x+d]$。转折在对齐块[\allowbreak{}a,\allowbreak{}a+\allowbreak{}d]\allowbreak{}中时，用 $L=\mathrm{lcs}(S[0,a),T[0,a+d))$、$R=\mathrm{lcp}(S[a,\cdot),T[a+d,\cdot))$，存在平方串恰为L+\allowbreak{}R\ensuremath{\ge}d。半长上限 $\lfloor(N+1)/2\rfloor$；不是S/\allowbreak{}S，也不是阈值N（2024 哈尔滨 D）。}
\par\textbf{自环期望与辅助初值分开}\quad
\note{不可分排列块重洗后按新最大块继续：g\_n为不可分排列数，C\_n为所有排列后续总代价，真实 $f_1=C_1=0$、n>\allowbreak{}1有 $f_n=1+C_n/n!$。按首块j拆，$C_n=g_nf_n+\sum_{j<n}(g_jf_j(n-j)!+g_jC_{n-j})$，隔离自环分母 $n!-g_n$。为统一卷积可另设辅助h\_1=\allowbreak{}1、n>\allowbreak{}1设h\_n=\allowbreak{}n!f\_n，从而C\_n=\allowbreak{}h\_n-\allowbreak{}n!；h\_1不是f\_1，不能混用（Metropolis F）。}
\par\textbf{凸包上的预算查询与优化 oracle}\quad
\note{子集S形成 $(x,y)=(|S|,|N(S)|)$，操作系数和\ensuremath{\le}1且空集允许，所有组合恰为这些点的凸包。下包Y(\allowbreak{}x)\allowbreak{}非减，总成本ax+\allowbreak{}bY(\allowbreak{}x)\allowbreak{}，收益ax。a>\allowbreak{}0时二分成本相邻顶点再线性插值；a=\allowbreak{}0答0，b=\allowbreak{}0答min(\allowbreak{}k,\allowbreak{}an)\allowbreak{}。恢复包的支持目标为Ax-\allowbreak{}By，用闭合子图；右端点是 $(n,|N(\text{全部左点})|)$，非必(\allowbreak{}n,\allowbreak{}n)\allowbreak{}。方向整数边长总和\ensuremath{\le}2n给顶点数O(\allowbreak{}n\textasciicircum{}(\allowbreak{}2/\allowbreak{}3)\allowbreak{})\allowbreak{}（香港 E）。}
\par\textbf{相邻消去变非交叉匹配}\quad
\note{洞H、机器人B最初交替为HBH…BH；激活B总与剩余序列相邻方向的H匹配并删除，仍交替，故本题不可能出界消失。最终匹配无交叉，且唯一未配H不能被弧包围；反过来按内层到外层删对即可实现。逐符号括号DP只存高度h、是否跳过空H的e、前一步是否入栈的o。B入栈须向右、出栈须向左；H只在h=\allowbreak{}0且e=\allowbreak{}0时可跳过。配对若前步入栈则原相邻是bad，否则good加1；顶层B由 $(h-1+e)\bmod2=1$ 判定，末态h=\allowbreak{}0/\allowbreak{}e=\allowbreak{}1，O(\allowbreak{}n\ensuremath{^2})\allowbreak{}时间、O(\allowbreak{}n)\allowbreak{}空间（2023 决赛 M）。}
\par\textbf{二分图单位势的奇偶舍入}\quad
\note{整数最优单位Lipschitz势x、二染色c，将奇偶不匹配点分别+\allowbreak{}1/\allowbreak{}-\allowbreak{}1得到x\ensuremath{^+}、x\ensuremath{^-}；原差±1时两端同时改，原差0时恰一端改，所以所有新边差恰±1。两新势平均为原势，线性最优目标也平均，故都最优。窄网格轮廓内部差±1，唯一对角间隙差-\allowbreak{}2/\allowbreak{}0/\allowbreak{}2，状态O(\allowbreak{}2\textasciicircum{}m)\allowbreak{}。归一化全部旧势减k时，累计分数必须减 $k\sum_{\text{已处理}v}(A_v-B_v)$，不能只删偏移（2025 决赛 A）。}
\par\textbf{最优盒体积界只用于覆盖}\quad
\note{四变量非负同余、正费用、素数p，字典序最小最优解满足 $\prod(x_i+1)\le p$，否则盒内两向量同余，正反差一者改善费用或字典序。枚举24个大小顺序，前两坐标组合O(\allowbreak{}√p)\allowbreak{}，第三上界U\ensuremath{\le}√p。设 $u=a_3/a_4$、$t=(m-a_1x_1-a_2x_2)/a_4$（模p），预存 $v_x=(-ux)\bmod p$、$c_x=b_3x+b_4v_x$。查询为 $b_4t+\min(\min_{x\le U,v_x<p-t}c_x,\ \min_{x\le U,v_x\ge p-t}c_x-b_4p)$，按U离线插入、按v范围min，O(\allowbreak{}√p logp)\allowbreak{}。候选可放宽大小顺序和盒界，但必须满足原同余/\allowbreak{}非负；覆盖最优后额外合法候选不会更优（2025 决赛 C）。}
\par\textbf{观测背包的唯一默认候选}\quad
\note{预先固定监控和搜索叶子，要求所有叶的观测签名不同；全零观测可留一个默认叶，根处也能靠排除法识别。未查父监控时，所有孩子合计至多一个默认叶；查父监控能辨孩子，每个孩子可各留一个。搜索k叶费用t\_k在树DP后统一付，t\_0=\allowbreak{}0；单叶树答案0。flag“最多一个”应包含无默认叶的状态，不能因状态简写漏掉min分支（2022 广州 A）。}
\par\textbf{双端分段回文的残串方向}\quad
\note{每段和转十进制（0是单个字符0），从两端比对。状态存未分段区间及一侧按外向内阅读的残串p；仅从另一侧选段，右段要反转，比较较短长度后留长者残余。残串为空时统一从左开始，避免一个分割有两种开始顺序。区间空时接受p自身回文，不能只接受空p；全零数组应计 $2^{n-1}$。残串用源段端点/\allowbreak{}偏移存，正确状态不代表朴素枚举转移已满足原时限（2022 广州 D）。}
\end{multicols}
\end{document}
